\section{Contextual editorial arguments}

Each passage below is an explicitly identified editorial treatment. The original statement and its exact source location remain available in the source records. The retained context supplies the definitions, statement and surrounding proof needed to read the argument. Source corrections, completed proofs and stronger conclusions keep their individual decisions in the accompanying ledger. No novelty or official endorsement is claimed.

\subsection{Editorial treatment of Algebra lines 1000--1031}

\hypertarget{algebra-context-001}{}

\noindent\href{https://github.com/stacks/stacks-project/blob/a04446e57ec1fbc252a871afcec7752fb2807b14/algebra.tex\#L1000}{Original source, lines 1000--1031}. Complete original and reviewed TeX passages accompany this note. Every intervention is separately recorded; the following treatment is editorial.


\begin{lemma}
\label{editorial-source-lemma-almost-directed-colimit-exact}
Let $R$ be a ring. Let $\mathcal{I}$ be an index category satisfying the assumptions of
Categories, Lemma \ref{categories-lemma-split-into-directed}.
Then taking colimits of diagrams of $R$-modules over $\mathcal{I}$
is exact (i.e., the analogue of
Lemma \ref{algebra-lemma-directed-colimit-exact}
holds in this situation).
\end{lemma}

\begin{proof}
By
Categories, Lemma \ref{categories-lemma-split-into-directed}
we may write $\mathcal{I} = \coprod_{j \in J} \mathcal{I}_j$ with each
$\mathcal{I}_j$ a filtered category, and $J$ possibly empty. By
Categories, Lemma \ref{categories-lemma-directed-category-system}
taking colimits over the index categories $\mathcal{I}_j$ is
the same as taking the colimit over some directed set. Hence
Lemma \ref{algebra-lemma-directed-colimit-exact}
applies to these colimits. This reduces the problem to showing that
coproducts in the category of $R$-modules over the set $J$ are exact.
In other words, given exact sequences
$L_j \to M_j \to N_j$ of $R$-modules, we have to show that
$$
\bigoplus\nolimits_{j \in J} L_j
\longrightarrow
\bigoplus\nolimits_{j \in J} M_j
\longrightarrow
\bigoplus\nolimits_{j \in J} N_j
$$
is exact. Indeed, denote its component maps by
$\alpha_j : L_j \to M_j$ and $\beta_j : M_j \to N_j$.
An element $(m_j)_{j \in J}$ of $\bigoplus_{j \in J} M_j$
has finite support. It is in the kernel of $\bigoplus_j\beta_j$
if and only if every $m_j$ is in $\Ker(\beta_j)$.
For each index in its finite support, exactness gives
$l_j \in L_j$ such that $\alpha_j(l_j) = m_j$.
Set $l_j = 0$ at every other index. Then $(l_j)_{j \in J}$
has finite support and maps to $(m_j)_{j \in J}$.
Conversely, $\beta_j\alpha_j = 0$ at every index, so the image
of $\bigoplus_j\alpha_j$ is contained in this kernel.

\medskip\noindent
More generally, for complexes $L_j \xrightarrow{\alpha_j} M_j
\xrightarrow{\beta_j} N_j$, the canonical $R$-linear map
$$
\bigoplus_{j \in J}
\bigl(\Ker(\beta_j)/\Im(\alpha_j)\bigr)
\longrightarrow
\Ker\bigl(\bigoplus_j\beta_j\bigr)
/\Im\bigl(\bigoplus_j\alpha_j\bigr),
\qquad ([m_j])_j \longmapsto [(m_j)_j]
$$
is an isomorphism. Choose representatives at the finitely many
nonzero components and zero representatives elsewhere. Changing these
choices changes $(m_j)_j$ by the image of a finitely supported
family in $\bigoplus_j L_j$, so the map is well defined.
Every element of the target has a finitely supported representative
whose components lie in $\Ker(\beta_j)$; this proves surjectivity.
If a representative lies in $\Im(\bigoplus_j\alpha_j)$, each
$m_j$ lies in $\Im(\alpha_j)$, proving injectivity.
All maps used are $R$-linear. Applying this to the colimits over
the components $\mathcal{I}_j$ proves the homology assertion as well.
If $J$ is empty, all the displayed direct sums and homology modules
are zero, and the same assertions hold.
\end{proof}


\subsection{Editorial treatment of Algebra lines 1301--1324}

\hypertarget{algebra-context-002}{}

\noindent\href{https://github.com/stacks/stacks-project/blob/a04446e57ec1fbc252a871afcec7752fb2807b14/algebra.tex\#L1301}{Original source, lines 1301--1324}. Complete original and reviewed TeX passages accompany this note. Every intervention is separately recorded; the following treatment is editorial.


\begin{proposition}
\label{editorial-source-proposition-localize-twice-module}
View $S'^{-1}M$ as an $A$-module, then $S^{-1}(S'^{-1}M)$ is
isomorphic to $(SS')^{-1}M$.
\end{proposition}

\begin{proof}
Note that given an $A$-module $M$, the universal property of
Lemma \ref{algebra-lemma-universal-property-localization-module} is available.
We give an explicit construction of the isomorphism.

\medskip\noindent
We define the maps as follows
\begin{align*}
& f : S^{-1}(S'^{-1}M) \longrightarrow (SS')^{-1}M, \quad \frac{x/s'}{s}\mapsto
x/ss'\\
& g : (SS')^{-1}M \longrightarrow S^{-1}(S'^{-1}M), \quad x/t\mapsto
\frac{x/s'}{s}\ \text{for some }s\in S, s'\in S', \text{ and }
t = ss'
\end{align*}
We have to check that these homomorphisms are well-defined, that is,
independent of the choice of the fraction. The universal property proves that these maps are well defined:
every element of $S'$ and of $S$ acts invertibly on $(SS')^{-1}M$,
so the map $M \to (SS')^{-1}M$ extends first to $S'^{-1}M$
and then to $S^{-1}(S'^{-1}M)$, giving $f$ with the displayed formula.
On $S^{-1}(S'^{-1}M)$ every $s' \in S'$ acts invertibly already
on the inner localization, and every $s \in S$ acts invertibly
on the outer localization. Thus every $ss' \in SS'$ acts invertibly,
and $M \to S^{-1}(S'^{-1}M)$ extends uniquely to $g$.
The displayed formula for $g$ follows by inverting the actions of
$s'$ and $s$; it therefore does not depend on the chosen factorization
of the denominator. Both composites restrict to the original
localization map on $M$. Uniqueness in the universal property,
applied once for $(SS')^{-1}M$ and successively for the two
localizations, gives $f \circ g = \mathrm{id}_{(SS')^{-1}M}$ and
$g \circ f = \mathrm{id}_{S^{-1}(S'^{-1}M)}$.
\end{proof}


\subsection{Editorial treatment of Algebra lines 1420--1433}

\hypertarget{algebra-context-003}{}

\noindent\href{https://github.com/stacks/stacks-project/blob/a04446e57ec1fbc252a871afcec7752fb2807b14/algebra.tex\#L1420}{Original source, lines 1420--1433}. Complete original and reviewed TeX passages accompany this note. Every intervention is separately recorded; the following treatment is editorial.


\begin{lemma}
\label{editorial-source-lemma-submodule-localization}
Any submodule $N'$ of $S^{-1}M$ is of the form $S^{-1}N$ for some
$N\subset M$. Indeed one can take $N$ to be the inverse image of
$N'$ in $M$.
\end{lemma}

\begin{proof}
Let $N$ be the inverse image of $N'$ in $M$. Then $S^{-1}N\supset N'$:
if $x/s \in N'$, then $x/1 = s(x/s) \in N'$, so $x \in N$
and hence $x/s \in S^{-1}N$. To show they are equal, take $x/s$ in
$S^{-1}N$, where $s\in S$ and $x\in N$. This yields that $x/1\in
N'$. Since $N'$ is an $S^{-1}A$-submodule we have
$x/s = x/1\cdot 1/s\in N'$. This finishes the proof.
\end{proof}


\subsection{Editorial treatment of Algebra lines 1588--1614}

\hypertarget{algebra-context-004}{}

\noindent\href{https://github.com/stacks/stacks-project/blob/a04446e57ec1fbc252a871afcec7752fb2807b14/algebra.tex\#L1588}{Original source, lines 1588--1614}. Complete original and reviewed TeX passages accompany this note. Every intervention is separately recorded; the following treatment is editorial.


\begin{lemma}
\label{editorial-source-lemma-characterize-finite-module-hom}
Let $R$ be a ring. Let $N$ be an $R$-module. The following are equivalent
\begin{enumerate}
\item $N$ is a finite $R$-module,
\item for any filtered colimit $M = \colim M_i$ of $R$-modules the map
$\colim \Hom_R(N, M_i) \to \Hom_R(N, M)$ is injective.
\end{enumerate}
\end{lemma}

\begin{proof}
Assume (1) and choose generators $x_1, \ldots, x_m$ for $N$.
If $N \to M_i$ is a module map and the composition
$N \to M_i \to M$ is zero, then because $M = \colim_{i' \geq i} M_{i'}$
for each $j \in \{1, \ldots, m\}$ we can find an $i' \geq i$ such that
$x_j$ maps to zero in $M_{i'}$. Since there are finitely many
$x_j$ we can find a single $i'$ which works for all of them.
Then the composition $N \to M_i \to M_{i'}$ is zero and we conclude
the map is injective, i.e., part (2) holds.

\medskip\noindent
Assume (2). For a finite subset $E \subset N$ denote $N_E \subset N$
the $R$-submodule generated by the elements of $E$. Then
$0 = \colim_E N/N_E$ is a filtered colimit. At the stage
$E = \emptyset$ we have $N/N_E = N$, and the class of
$\mathrm{id}_N$ in $\colim_E\Hom_R(N, N/N_E)$ maps to zero in
$\Hom_R(N, 0)$. By (2) this class is zero. Thus, at some finite
stage $E$, its image is the zero map. That image is the quotient
map $q_E : N \to N/N_E$, so $q_E = 0$ and $N = N_E$.
Equivalently, $\mathrm{id}_N$ factors through the inclusion
$N_E \hookrightarrow N$. Hence $N$ is finitely generated.
\end{proof}


\subsection{Editorial treatment of Algebra lines 1715--1776}

\hypertarget{algebra-context-005}{}

\noindent\href{https://github.com/stacks/stacks-project/blob/a04446e57ec1fbc252a871afcec7752fb2807b14/algebra.tex\#L1715}{Original source, lines 1715--1776}. Complete original and reviewed TeX passages accompany this note. Every intervention is separately recorded; the following treatment is editorial.


\begin{lemma}
\label{editorial-source-lemma-tensor-product}
Let $M, N$ be $R$-modules. Then there exists a pair $(T, g)$
where $T$ is an $R$-module, and
$g : M \times N \to T$ an $R$-bilinear
mapping, with the following universal property:
For any $R$-module $P$ and any $R$-bilinear mapping
$f : M \times N \to P$, there
exists a unique $R$-linear
mapping $\tilde{f} : T \to P$ such that $f = \tilde{f} \circ g$.
In other words, the following diagram commutes:
$$
\xymatrix{
M \times N \ar[rr]^f \ar[dr]_g & & P\\
& T \ar[ur]_{\tilde f}
}
$$
Moreover, if $(T, g)$ and $(T', g')$
are two pairs with this property, then there
exists a unique isomorphism
$j : T \to T'$ such that $j\circ g = g'$.
\end{lemma}

\noindent
The $R$-module $T$ which satisfies the above universal property is called
the  {\it tensor product} of $R$-modules $M$ and $N$, denoted as
$M \otimes_R N$.

\begin{proof}
We first prove the existence of such $R$-module $T$.
Let $M, N$ be $R$-modules.
Let $T$ be the quotient module
$P/Q$, where $P$ is the free $R$-module $R^{(M \times N)}$ and $Q$ is the
$R$-module generated by all elements of
the following types: ($a\in R$, $x,x'\in M$, $y,y'\in N$)
\begin{align*}
(x + x', y) - (x, y) - (x', y), \\
(x, y + y') - (x, y) - (x, y'), \\
(ax, y) - a(x, y), \\
(x, ay) - a(x, y)
\end{align*}
Let $\pi : M \times N \to T$ denote the natural map.
This map is $R$-bilinear, as
implied by the above relations
when we check the bilinearity conditions. Denote the image by
$\pi(x, y) = x \otimes
y$. Then these elements generate
$T$. Now let $P'$ be any $R$-module and let
$f : M \times N \to P'$ be an $R$-bilinear map. The unique
$R$-linear map from the free module $P$ to $P'$ that sends the
basis element $(x,y)$ to $f(x,y)$ kills every displayed generator
of $Q$, by bilinearity. It therefore induces an $R$-linear map
$f' : T = P/Q \to P'$ with $f'(x\otimes y) = f(x,y)$.
Thus $f = f'\circ\pi$. The elements $x\otimes y$ generate $T$,
so these values determine $f'$ uniquely.
Suppose there is another pair $(T', g')$ satisfying the same properties.
Then there is a unique $j : T \to T'$ and
also $j' : T' \to T$ such that $g' = j\circ g$, $g = j'\circ g'$.
But then both the maps $(j'\circ j) \circ g$ and $g$
satisfy the universal properties, so by uniqueness they are equal,
and hence $j'\circ j$ is identity on $T$.
Similarly $(j\circ j') \circ g' = g'$ and $j\circ j'$ is identity on $T'$.
So $j$ is an isomorphism.
\end{proof}


\subsection{Editorial treatment of Algebra lines 1933--1991}

\hypertarget{algebra-context-006}{}

\noindent\href{https://github.com/stacks/stacks-project/blob/a04446e57ec1fbc252a871afcec7752fb2807b14/algebra.tex\#L1933}{Original source, lines 1933--1991}. Complete original and reviewed TeX passages accompany this note. Every intervention is separately recorded; the following treatment is editorial.


\begin{lemma}[Tensor products commute with colimits]
\label{editorial-source-lemma-tensor-products-commute-with-limits}
Let $(M_i, \mu_{ij})$ be a system over the preordered set $I$.
Let $N$ be an $R$-module. Then
$$
\colim (M_i \otimes N) \cong (\colim M_i)\otimes N.
$$
Moreover, the isomorphism is induced by the homomorphisms
$\mu_i \otimes 1: M_i \otimes N \to M \otimes N$
where $M = \colim_i M_i$ with natural maps $\mu_i : M_i \to M$.
More generally, for any small index category $\mathcal{I}$ and any
diagram $i\mapsto M_i$ of $R$-modules, the same canonical map
$$
\colim_{i\in\mathcal{I}}(M_i\otimes_R N)
\longrightarrow
\left(\colim_{i\in\mathcal{I}}M_i\right)\otimes_R N,
\qquad [m_i\otimes n]\longmapsto\mu_i(m_i)\otimes n,
$$
is an isomorphism.
\end{lemma}

\begin{proof}
First proof. The functor $M' \mapsto M' \otimes_R N$ is left adjoint
to the functor $N' \mapsto \Hom_R(N, N')$ by
Lemma \ref{algebra-lemma-hom-from-tensor-product}. Thus $M' \mapsto M' \otimes_R N$
commutes with all colimits, see
Categories, Lemma \ref{categories-lemma-adjoint-exact}.

\medskip\noindent
Second direct proof. Let $P = \colim (M_i \otimes N)$ with coprojections
$\lambda_i : M_i \otimes N \to P$. Let $M = \colim M_i$ with coprojections
$\mu_i : M_i \to M$. Then for all $i\leq j$, the following diagram commutes:
$$
\xymatrix{
M_i \otimes N \ar[r]_{\mu_i \otimes 1} \ar[d]_{\mu_{ij} \otimes 1} &
M \otimes N \ar[d]^{\text{id}} \\
M_j \otimes N \ar[r]^{\mu_j \otimes 1} &
M \otimes N
}
$$
By Lemma \ref{algebra-lemma-homomorphism-limit} these maps induce a unique homomorphism
$\psi : P \to M \otimes N$ such that
$\mu_i \otimes 1 = \psi \circ \lambda_i$.

\medskip\noindent
To construct the inverse map, for each $i\in I$, there is the canonical
$R$-bilinear mapping $g_i : M_i \times N \to
M_i \otimes N$. This induces a unique mapping
$\widehat{\phi} : M \times N \to P$
such that $\widehat{\phi} \circ (\mu_i \times 1) = \lambda_i \circ g_i$.
It is $R$-bilinear. Thus it induces an
$R$-linear mapping $\phi : M \otimes N \to P$.
From the commutative diagram below:
$$
\xymatrix{
M_i \times N \ar[r]^{g_i} \ar[d]^{\mu_i \times \text{id}} &
M_i \otimes N\ar[r]_{\text{id}} \ar[d]_{\lambda_i} &
M_i \otimes N \ar[d]_{\mu_i \otimes \text{id}} \ar[rd]^{\lambda_i} \\
M \times N \ar[r]^{\widehat{\phi}} &
P \ar[r]^{\psi} & M \otimes N \ar[r]^{\phi} & P
}
$$
we see that $\psi\circ\widehat{\phi} = g$, the canonical $R$-bilinear mapping
$g : M \times N \to M \otimes N$. So
$\psi\circ\phi$ is identity on $M \otimes N$. From the right-hand square and
triangle, $\phi\circ\psi$ is also
identity on $P$.
\end{proof}


\subsection{Editorial treatment of Algebra lines 2087--2121}

\hypertarget{algebra-context-007}{}

\noindent\href{https://github.com/stacks/stacks-project/blob/a04446e57ec1fbc252a871afcec7752fb2807b14/algebra.tex\#L2087}{Original source, lines 2087--2121}. Complete original and reviewed TeX passages accompany this note. Every intervention is separately recorded; the following treatment is editorial.


\begin{lemma}
\label{editorial-source-lemma-tensor-localization}
Let $M$ be an $R$-module. Then the $S^{-1}R$-modules $S^{-1}M$
and $S^{-1}R \otimes_R M$ are canonically isomorphic, and the
canonical isomorphism $f : S^{-1}R \otimes_R M \to S^{-1}M$
is given by
$$
f((a/s) \otimes m) = am/s, \forall a \in R, m \in M, s \in S
$$
\end{lemma}

\begin{proof}
Obviously, the map
$f' : S^{-1}R \times M \to S^{-1}M$ given by $f'(a/s, m) = am/s$ is
bilinear, and thus by the
universal property, this map induces a unique $S^{-1}R$-module homomorphism
$f : S^{-1}R \otimes_R M \to S^{-1}M$ as in the statement of the lemma.
Actually every element in $S^{-1}M$ is of the form $m/s$, $m\in M, s\in S$ and
every element in
$S^{-1}R \otimes_R M$ is of the form $1/s \otimes m$. To see the latter fact,
write an element in $S^{-1}R\otimes_R M$ as a finite sum
indexed by a finite set $K$. Put
$s = \prod_{k\in K}s_k$ and
$t_k = \prod_{\ell\in K\setminus\{k\}}s_\ell$ for $k\in K$.
Thus $s=s_kt_k$ and all these elements belong to $S$.
Every sum in the next display is over $k\in K$.
For the empty sum take $s=1$, so its tensor representative is
$1\otimes 0=0$. We obtain
$$
\sum_k \frac{a_k}{s_k} \otimes m_k =
\sum_k \frac{a_k t_k}{s} \otimes m_k =
\frac{1}{s} \otimes \sum_k {a_k t_k}m_k = \frac{1}{s} \otimes m
$$
Here $m = \sum_k {a_k t_k}m_k$. Then it is obvious that $f$ is surjective,
and if $f(\frac{1}{s} \otimes m) = m/s = 0$ then there exists $t \in S$ with
$tm = 0$ in $M$. Then we have
$$
\frac{1}{s} \otimes m = \frac{1}{st} \otimes tm = \frac{1}{st} \otimes 0 = 0
$$
Therefore $f$ is injective.
\end{proof}


\subsection{Editorial treatment of Algebra lines 2240--2268}

\hypertarget{algebra-context-008}{}

\noindent\href{https://github.com/stacks/stacks-project/blob/a04446e57ec1fbc252a871afcec7752fb2807b14/algebra.tex\#L2240}{Original source, lines 2240--2268}. Complete original and reviewed TeX passages accompany this note. Every intervention is separately recorded; the following treatment is editorial.


\begin{lemma}
\label{editorial-source-lemma-presentation-sym-exterior}
Let $R$ be a ring.
Let $M_2 \to M_1 \to M \to 0$ be an exact sequence of $R$-modules.
For every integer $n \geq 1$ there are exact sequences
$$
M_2 \otimes_R \text{Sym}^{n - 1}(M_1)
\to
\text{Sym}^n(M_1)
\to
\text{Sym}^n(M)
\to
0
$$
and similarly
$$
M_2 \otimes_R \wedge^{n - 1}(M_1)
\to
\wedge^n(M_1)
\to
\wedge^n(M)
\to
0
$$
In degree zero both induced maps are the identity map $R \to R$.
\end{lemma}

\begin{proof}
Denote the maps by $u : M_2 \to M_1$ and $v : M_1 \to M$.
We prove both assertions, writing $C(M_1)$ for the symmetric algebra
or the exterior algebra in the respective case. Let $J$ be the
homogeneous two sided ideal of $C(M_1)$ generated by all
$u(y)$ in degree one, for $y \in M_2$.
The map $v$ induces a graded $R$-algebra map
$C(M_1)/J \to C(M)$. Conversely, the assignment
$v(x) \mapsto x \bmod J$ defines an $R$-linear map
$M \to (C(M_1)/J)^1$: it is independent of the lift $x$
because $\Ker(v) = \Im(u)$. Sending a pure tensor in $M$
to the product of these images defines a map from its tensor algebra,
with the identity map on the degree-zero part $R$.
In the symmetric case it kills the commutator relations.
In the exterior case it kills $m\otimes m$, since a lift $x$ of $m$
satisfies $x\wedge x = 0$ in $C(M_1)$.
Thus it induces a graded $R$-algebra map
$C(M) \to C(M_1)/J$. The two maps are inverse on $R$ and on
the degree-one generators, and hence on the entire algebras.

\medskip\noindent
It follows that the kernel of $C^n(M_1) \to C^n(M)$ is $J^n$.
For $n \geq 1$ every homogeneous term generating $J^n$ has the form
$a\,u(y)\,b$, with $a$ of degree $p$, $b$ of degree $q$ and
$p+1+q=n$. In the symmetric algebra this equals $u(y)ab$.
In the exterior algebra it equals $(-1)^p u(y)ab$:
on a pure degree-$p$ tensor this is obtained by $p$ successive
interchanges of degree-one factors, and the assertion extends by
$R$-linearity. Therefore $J^n$ is exactly the image of the map
$$
M_2\otimes_R C^{n-1}(M_1) \longrightarrow C^n(M_1),
\qquad y\otimes z \longmapsto u(y)z.
$$
The reverse inclusion holds because every $u(y)$ belongs to $J$.
This proves the two displayed exact sequences, including their
surjective last maps. In degree zero $J^0=0$, and both maps
$C^0(M_1)=R \to C^0(M)=R$ are the identity.
\end{proof}


\subsection{Editorial treatment of Algebra lines 2834--2858}

\hypertarget{algebra-context-009}{}

\noindent\href{https://github.com/stacks/stacks-project/blob/a04446e57ec1fbc252a871afcec7752fb2807b14/algebra.tex\#L2834}{Original source, lines 2834--2858}. Complete original and reviewed TeX passages accompany this note. Every intervention is separately recorded; the following treatment is editorial.


\begin{lemma}
\label{editorial-source-lemma-charpoly-module}
Let $R$ be a ring.
Let $M$ be a finite $R$-module.
Let $\varphi : M \to M$ be an endomorphism.
Then there exists a monic polynomial $P \in R[T]$ such that
$P(\varphi) = 0$ as an endomorphism of $M$.
\end{lemma}

\begin{proof}
Choose a surjective $R$-module map $R^{\oplus n} \to M$, given by
$(a_1, \ldots, a_n) \mapsto \sum a_ix_i$ for some generators $x_i \in M$.
Choose $(a_{i1}, \ldots, a_{in}) \in R^{\oplus n}$ such that
$\varphi(x_i) = \sum a_{ij} x_j$. In other words the diagram
$$
\xymatrix{
R^{\oplus n} \ar[d]_{A^{\mathrm{t}}} \ar[r] & M \ar[d]^\varphi \\
R^{\oplus n} \ar[r] & M
}
$$
is commutative where $A = (a_{ij})$ and $A^{\mathrm{t}}$
denotes its transpose: the lift sends the $i$th basis vector to
$\sum_j a_{ij}e_j$, which maps to $\varphi(x_i)$.
By
Lemma \ref{algebra-lemma-charpoly}
there exists a monic polynomial $P$ such that $P(A) = 0$.
Taking transposes gives $P(A^{\mathrm{t}})=P(A)^{\mathrm{t}}=0$,
since the coefficient matrices are scalar and transpose preserves
powers of this one matrix.
Then it follows that $P(\varphi) = 0$.
\end{proof}


\subsection{Editorial treatment of Algebra lines 2860--2887}

\hypertarget{algebra-context-010}{}

\noindent\href{https://github.com/stacks/stacks-project/blob/a04446e57ec1fbc252a871afcec7752fb2807b14/algebra.tex\#L2860}{Original source, lines 2860--2887}. Complete original and reviewed TeX passages accompany this note. Every intervention is separately recorded; the following treatment is editorial.


\begin{lemma}
\label{editorial-source-lemma-charpoly-module-ideal}
Let $R$ be a ring. Let $I \subset R$ be an ideal.
Let $M$ be a finite $R$-module.
Let $\varphi : M \to M$ be an endomorphism such
that $\varphi(M) \subset IM$.
Then there exists a monic polynomial
$P = T^n + a_1 T^{n - 1} + \ldots + a_n \in R[T]$
such that $a_j \in I^j$ and $P(\varphi) = 0$ as an endomorphism of $M$.
\end{lemma}

\begin{proof}
Choose a surjective $R$-module map $R^{\oplus n} \to M$, given by
$(a_1, \ldots, a_n) \mapsto \sum a_ix_i$ for some generators $x_i \in M$.
Choose $(a_{i1}, \ldots, a_{in}) \in I^{\oplus n}$ such that
$\varphi(x_i) = \sum a_{ij} x_j$. In other words the diagram
$$
\xymatrix{
R^{\oplus n} \ar[d]_{A^{\mathrm{t}}} \ar[r] & M \ar[d]^\varphi \\
I^{\oplus n} \ar[r] & M
}
$$
is commutative where $A = (a_{ij})$: its transpose sends
$e_i$ to $\sum_j a_{ij}e_j \in I^{\oplus n}$, whose image in
$M$ is $\varphi(x_i)$. By
Lemma \ref{algebra-lemma-charpoly}
the polynomial
$P(t) = \det(t\text{id}_{n \times n} - A)$
is monic of degree $n$ and satisfies $P(A)=0$.
Transposing yields $P(A^{\mathrm{t}})=0$, so the commutative diagram
and the surjection $R^{\oplus n}\to M$ give $P(\varphi)=0$.
In the determinant expansion, a term contributing to the coefficient
of $t^{n-j}$ uses exactly $j$ entries of $-A$ and $n-j$ copies of $t$
from diagonal entries. Each such term is therefore a product of
$j$ elements of $I$, with its determinant sign, and lies in $I^j$.
Their sum is the coefficient $a_j$, so $a_j\in I^j$ as required.
When $n=0$, $M=0$ and the empty determinant is $P=1$;
its action is the zero endomorphism of the zero module.
\end{proof}


\subsection{Editorial treatment of Algebra lines 2892--2912}

\hypertarget{algebra-context-011}{}

\noindent\href{https://github.com/stacks/stacks-project/blob/a04446e57ec1fbc252a871afcec7752fb2807b14/algebra.tex\#L2892}{Original source, lines 2892--2912}. Complete original and reviewed TeX passages accompany this note. Every intervention is separately recorded; the following treatment is editorial.


\begin{lemma}
\label{editorial-source-lemma-fun}
Let $R$ be a ring.
Let $M$ be a finite $R$-module.
Let $\varphi : M \to M$ be a surjective $R$-module map.
Then $\varphi$ is an isomorphism. Moreover, its inverse belongs
to the subalgebra $R[\varphi]\subset\operatorname{End}_R(M)$:
there is a polynomial $Q\in R[T]$ with $\varphi^{-1}=Q(\varphi)$.
\end{lemma}

\begin{proof}[First proof]
Write $R' = R[x]$ and think of $M$ as a finite $R'$-module with
$x$ acting via $\varphi$. Set $I = (x) \subset R'$. By our assumption that
$\varphi$ is surjective we have $IM = M$. Hence we may apply
Lemma \ref{editorial-source-lemma-charpoly-module-ideal}
to $M$ as an $R'$-module, the ideal $I$ and the endomorphism $\text{id}_M$.
We conclude that
$(1 + a_1 + \ldots + a_n)\text{id}_M = 0$ with $a_j \in I$.
Write $a_j = b_j(x)x$ for some $b_j(x) \in R[x]$.
Translating back into $\varphi$ we see that
$\text{id}_M = -(\sum_{j = 1}^{n} b_j(\varphi)) \varphi$, and hence
$\varphi$ is invertible. More explicitly, put
$Q(T)=-\sum_{j=1}^{n}b_j(T)\in R[T]$, where $b_j(T)$ is
obtained from $b_j(x)$ by replacing its indeterminate $x$ by $T$.
The displayed identity gives $Q(\varphi)\varphi=\mathrm{id}_M$.
Every power of $\varphi$ and every scalar multiplication commute
with $\varphi$, so also $\varphi Q(\varphi)=\mathrm{id}_M$.
Thus the same polynomial gives the two-sided inverse.
If $M=0$, the zero polynomial gives its unique endomorphism,
which is also its identity and inverse.
\end{proof}


\subsection{Editorial treatment of Algebra lines 2986--3102}

\hypertarget{algebra-context-012}{}

\noindent\href{https://github.com/stacks/stacks-project/blob/a04446e57ec1fbc252a871afcec7752fb2807b14/algebra.tex\#L2986}{Original source, lines 2986--3102}. Complete original and reviewed TeX passages accompany this note. Every intervention is separately recorded; the following treatment is editorial.


\begin{lemma}
\label{editorial-source-lemma-Zariski-topology}
Let $R$ be a ring.
\begin{enumerate}
\item The spectrum of a ring $R$ is empty if and only if $R$
is the zero ring.
\item Every nonzero ring has a maximal ideal.
\item Every nonzero ring has a minimal prime ideal.
\item Given an ideal $I \subset R$ and a prime ideal
$I \subset \mathfrak p$ there exists a prime
$I \subset \mathfrak q \subset \mathfrak p$ such
that $\mathfrak q$ is minimal over $I$.
\item If $T \subset R$, and if $(T)$ is the ideal generated by
$T$ in $R$, then $V((T)) = V(T)$.
\item If $I$ is an ideal and $\sqrt{I}$ is its radical,
see basic notion (\ref{algebra-item-radical-ideal}), then $V(I) = V(\sqrt{I})$.
\item Given an ideal $I$ of $R$ we have $\sqrt{I} =
\bigcap_{I \subset \mathfrak p} \mathfrak p$.
\item If $I$ is an ideal then $V(I) = \emptyset$ if and only
if $I$ is the unit ideal.
\item If $I$, $J$ are ideals of $R$ then $V(I) \cup V(J) =
V(I \cap J)$.
\item If $(I_a)_{a\in A}$ is a set of ideals of $R$ then
$\bigcap_{a\in A} V(I_a) = V(\bigcup_{a\in A} I_a)$.
\item If $f \in R$, then $D(f) \amalg V(f) = \Spec(R)$.
\item If $f \in R$ then $D(f) = \emptyset$ if and only if $f$
is nilpotent.
\item If $f = u f'$ for some unit $u \in R$, then $D(f) = D(f')$.
\item If $I \subset R$ is an ideal, and $\mathfrak p$ is a prime of
$R$ with $\mathfrak p \not\in V(I)$, then there exists an $f \in I$
such that $\mathfrak p \in D(f)$, and $D(f) \cap V(I) = \emptyset$.
\item If $f, g \in R$, then $D(fg) = D(f) \cap D(g)$.
\item If $f_i \in R$ for $i \in I$, then
$\bigcup_{i\in I} D(f_i)$ is the complement of $V(\{f_i \}_{i\in I})$
in $\Spec(R)$.
\item If $f \in R$ and $D(f) = \Spec(R)$, then $f$ is a unit.
\end{enumerate}
\end{lemma}

\begin{proof}
We address each part in the corresponding item below.
\begin{enumerate}
\item This is a direct consequence of (2) or (3).
\item Let $\mathfrak{A}$ be the set of all proper ideals of $R$. This set is
ordered by inclusion and is non-empty, since $(0) \in \mathfrak{A}$ is a proper
ideal. An empty totally ordered subset has the upper bound $(0)$.
Let $A$ be a nonempty totally ordered subset of $\mathfrak A$.
Then $\bigcup_{I \in A} I$ is in
fact an ideal. Since $1 \notin I$ for all $I \in A$, the union does not contain
$1$ and thus is proper. Hence $\bigcup_{I \in A} I$ is in $\mathfrak{A}$ and is
an upper bound for the set $A$. Thus by Zorn's lemma $\mathfrak{A}$ has a
maximal element, which is the sought-after maximal ideal.
\item Since $R$ is nonzero, it contains a maximal ideal which is a prime ideal.
Thus the set $\mathfrak{A}$ of all prime ideals of $R$ is nonempty.
$\mathfrak{A}$ is ordered by reverse-inclusion. An empty totally
ordered subset has any one of these prime ideals as an upper bound.
Let $A$ be a nonempty totally ordered subset of $\mathfrak{A}$.
The intersection $J = \bigcap_{I \in A} I$ is an ideal, and it is proper
because it is contained in any chosen member of $A$.
To prove that it is prime, let $xy \in J$.
Then $xy \in I$ for all $I \in A$. Now let $B = \{I \in A | y \in I\}$. Let $K
= \bigcap_{I \in B} I$. Since $A$ is totally ordered, either $K = J$ (and we're
done, since then $y \in J$) or $K \supset J$ and for all $I \in A$ such that
$I$ is properly contained in $K$, we have $y \notin I$. But that means that for
all those $I$, we have $x \in I$, since they are prime.
Here is why this implies $x \in J$. Since $K\ne J$, choose
$a\in K\setminus J$ and $I_0\in A$ with $a\notin I_0$.
For each $H\in B$, total ordering forces $I_0\subset H$:
the other possibility $H\subset I_0$ would put $a\in K\subset H$
inside $I_0$. Consequently $I_0\subset K$, strictly because
$a\in K\setminus I_0$. This also holds when $B$ is empty,
with its intersection interpreted as $R$.
For any $L\in A$, either $L\subset I_0$, in which case
$L\subsetneq K$ and the preceding primality argument gives $x\in L$,
or $I_0\subset L$, in which case $x\in I_0\subset L$.
Thus $x$ belongs to every $L\in A$ and hence to $J$. In either case,
$J$ is prime as desired. Hence by Zorn's lemma we get a maximal element which
in this case is a minimal prime ideal.
\item This is the same exact argument as (3) except you only consider prime
ideals contained in $\mathfrak{p}$ and containing $I$.
This collection is nonempty because it contains $\mathfrak p$,
which also supplies an upper bound for the empty chain in reverse inclusion.
For a nonempty chain its intersection still contains $I$, is contained
in $\mathfrak p$, and is prime by the argument in (3).
The resulting minimal member $\mathfrak q$ is minimal over $I$
among all primes: any prime between $I$ and $\mathfrak q$ also
lies inside $\mathfrak p$ and therefore belongs to the same collection.
\item $(T)$ is the smallest ideal containing $T$. Hence if $T \subset I$, some
ideal, then $(T) \subset I$ as well. Hence if $I \in V(T)$, then $I \in V((T))$
as well. The other inclusion is obvious.
\item Since $I \subset \sqrt{I}, V(\sqrt{I}) \subset V(I)$. Now let
$\mathfrak{p} \in V(I)$. Let $x \in \sqrt{I}$. Then $x^n \in I$ for some $n$.
Hence $x^n \in \mathfrak{p}$. But since $\mathfrak{p}$ is prime, a boring
induction argument gets you that $x \in \mathfrak{p}$. Hence $\sqrt{I} \subset
\mathfrak{p}$ and $\mathfrak{p} \in V(\sqrt{I})$.
\item Let $f \in R \setminus \sqrt{I}$. Then $f^n \notin I$ for all $n$. Hence
$S = \{1, f, f^2, \ldots\}$ is a multiplicative subset, not containing $0$.
Moreover, $S\cap I=\emptyset$ and $S^{-1}I$ is proper.
Indeed, if $1=a/s$ with $a\in I$ and $s\in S$, equality of fractions
gives $u(s-a)=0$ for some $u\in S$. Then $us=ua\in I\cap S$,
a contradiction. Thus $(S^{-1}R)/(S^{-1}I)$ is a nonzero ring.
By (2) it has a maximal ideal; its inverse image supplies a
prime ideal $\bar{\mathfrak{p}} \subset S^{-1}R$ containing $S^{-1}I$. Then the
pull-back $\mathfrak{p}$ in $R$ of $\bar{\mathfrak{p}}$ is a prime ideal
containing $I$ that does not intersect $S$. This shows that $\bigcap_{I \subset
\mathfrak p} \mathfrak p \subset \sqrt{I}$. Now if $a \in \sqrt{I}$, then $a^n
\in I$ for some $n$. Hence if $I \subset \mathfrak{p}$, then $a^n \in
\mathfrak{p}$. But since $\mathfrak{p}$ is prime, we have $a \in \mathfrak{p}$.
Thus the equality is shown. If $I=R$, both sides are $R$,
with the intersection over the empty set of prime ideals interpreted as $R$.
\item $I$ is not the unit ideal if and only if $I$
is contained in some maximal ideal (to
see this, apply (2) to the ring $R/I$) which is therefore prime.
\item If $\mathfrak{p} \in V(I) \cup V(J)$, then $I \subset \mathfrak{p}$ or $J
\subset \mathfrak{p}$ which means that $I \cap J \subset \mathfrak{p}$. Now if
$I \cap J \subset \mathfrak{p}$, then $IJ \subset \mathfrak{p}$ and hence
either $I \subset \mathfrak{p}$ or $J \subset \mathfrak{p}$, since $\mathfrak{p}$ is prime.
\item $\mathfrak{p} \in \bigcap_{a \in A} V(I_a) \Leftrightarrow
I_a \subset \mathfrak{p}, \forall a \in A \Leftrightarrow
\mathfrak{p} \in V(\bigcup_{a\in A} I_a)$
\item If $\mathfrak{p}$ is a prime ideal and $f \in R$, then either $f \in
\mathfrak{p}$ or $f \notin \mathfrak{p}$ (strictly) which is what the disjoint
union says.
\item If $a \in R$ is nilpotent, then $a^n = 0$ for some $n$. Hence $a^n \in
\mathfrak{p}$ for any prime ideal. Thus $a \in \mathfrak{p}$ as can be shown by
induction and $D(a) = \emptyset$. Now, as shown in (7), if $a \in R$ is not
nilpotent, then there is a prime ideal that does not contain it.
\item $f \in \mathfrak{p} \Leftrightarrow uf \in \mathfrak{p}$, since $u$ is
invertible.
\item If $\mathfrak{p} \notin V(I)$, then $\exists f \in I \setminus
\mathfrak{p}$. Then $f \notin \mathfrak{p}$ so $\mathfrak{p} \in D(f)$. Also if
$\mathfrak{q} \in D(f)$, then $f \notin \mathfrak{q}$ and thus $I$ is not
contained in $\mathfrak{q}$. Thus $D(f) \cap V(I) = \emptyset$.
\item If $fg \in \mathfrak{p}$, then $f \in \mathfrak{p}$ or $g \in
\mathfrak{p}$. Hence if $f \notin \mathfrak{p}$ and $g \notin \mathfrak{p}$,
then $fg \notin \mathfrak{p}$. Since $\mathfrak{p}$ is an ideal, if $fg \notin
\mathfrak{p}$, then $f \notin \mathfrak{p}$ and $g \notin \mathfrak{p}$.
\item $\mathfrak{p} \in \bigcup_{i \in I} D(f_i) \Leftrightarrow \exists i \in
I, f_i \notin \mathfrak{p} \Leftrightarrow \mathfrak{p} \in \Spec(R)
\setminus V(\{f_i\}_{i \in I})$
\item If $D(f) = \Spec(R)$, then $V(f) = \emptyset$ and
hence $fR = R$, so $f$ is a unit.
\end{enumerate}
\end{proof}


\subsection{Editorial treatment of Algebra lines 3162--3204}

\hypertarget{algebra-context-013}{}

\noindent\href{https://github.com/stacks/stacks-project/blob/a04446e57ec1fbc252a871afcec7752fb2807b14/algebra.tex\#L3162}{Original source, lines 3162--3204}. Complete original and reviewed TeX passages accompany this note. Every intervention is separately recorded; the following treatment is editorial.


\begin{lemma}
\label{editorial-source-lemma-spec-localization}
Let $R$ be a ring. Let $S \subset R$ be a multiplicative subset.
The map $R \to S^{-1}R$ induces via the functoriality of $\Spec$
a homeomorphism
$$
\Spec(S^{-1}R)
\longrightarrow
\{\mathfrak p \in \Spec(R) \mid S \cap \mathfrak p = \emptyset \}
$$
where the topology on the right hand side is that induced from the
Zariski topology on $\Spec(R)$. The inverse map is given
by $\mathfrak p \mapsto S^{-1}\mathfrak p = \mathfrak p(S^{-1}R)$.
\end{lemma}

\begin{proof}
Denote the right hand side of the arrow of the lemma by $D$.
Choose a prime $\mathfrak p' \subset S^{-1}R$ and let $\mathfrak p$ be
the inverse image of $\mathfrak p'$ in $R$. Since $\mathfrak p'$
does not contain $1$ we see that $\mathfrak p$ does not contain
any element of $S$. Hence $\mathfrak p \in D$ and we see that
the image is contained in $D$. Let $\mathfrak p \in D$.
By assumption the image $\overline{S}$ does not contain $0$.
By basic notion (\ref{algebra-item-localization-zero})
$\overline{S}^{-1}(R/\mathfrak p)$ is not the zero ring.
By basic notion (\ref{algebra-item-localize-ideal}) we see
$S^{-1}R / S^{-1}\mathfrak p = \overline{S}^{-1}(R/\mathfrak p)$
is a domain, and hence $S^{-1}\mathfrak p$ is a prime.
The equality of rings also shows that the inverse image of
$S^{-1}\mathfrak p$ in $R$ is equal to $\mathfrak p$,
because $R/\mathfrak p \to \overline{S}^{-1}(R/\mathfrak p)$
is injective by basic notion (\ref{algebra-item-localize-nonzerodivisors}).
In the other direction, start with $\mathfrak p'\subset S^{-1}R$
and put $\mathfrak p$ equal to its inverse image.
For $a\in R$ and $s\in S$, the element $s/1$ is a unit, so
$a/s\in\mathfrak p'$ if and only if $a/1\in\mathfrak p'$,
if and only if $a\in\mathfrak p$.
It follows that $S^{-1}\mathfrak p=\mathfrak p'$.
Together with the preceding contraction identity, this proves that
the map $\Spec(S^{-1}R) \to \Spec(R)$
is bijective onto $D$ with inverse as given.
It is continuous by Lemma \ref{algebra-lemma-spec-functorial}.
Finally, let $D(g) \subset \Spec(S^{-1}R)$ be a standard
open. Write $g = h/s$ for some $h\in R$ and $s\in S$.
Since $g$ and $h/1$ differ by a unit we have $D(g) =
D(h/1)$ in $\Spec(S^{-1}R)$.
Hence by Lemma \ref{algebra-lemma-spec-functorial} and the bijectivity
above the image of $D(g) = D(h/1)$ is $D \cap D(h)$.
This proves the map is open as well.
\end{proof}


\subsection{Editorial treatment of Algebra lines 3366--3397}

\hypertarget{algebra-context-014}{}

\noindent\href{https://github.com/stacks/stacks-project/blob/a04446e57ec1fbc252a871afcec7752fb2807b14/algebra.tex\#L3366}{Original source, lines 3366--3397}. Complete original and reviewed TeX passages accompany this note. Every intervention is separately recorded; the following treatment is editorial.


\begin{lemma}
\label{editorial-source-lemma-characterize-local-ring}
Let $R$ be a ring. The following are equivalent:
\begin{enumerate}
\item $R$ is a local ring,
\item $\Spec(R)$ has exactly one closed point,
\item $R$ has a maximal ideal $\mathfrak m$
and every element of $R \setminus \mathfrak m$
is a unit, and
\item $R$ is not the zero ring and for every $x \in R$ either $x$
or $1 - x$ is invertible or both.
\end{enumerate}
\end{lemma}

\begin{proof}
Let $R$ be a ring, and $\mathfrak m$ a maximal ideal.
If $x \in R \setminus \mathfrak m$, and $x$ is not a unit
then there is a maximal ideal $\mathfrak m'$ containing $x$.
Hence $R$ has at least two maximal ideals. Conversely,
if $\mathfrak m'$ is another maximal ideal, then choose
$x \in \mathfrak m'$, $x \not \in \mathfrak m$. Clearly
$x$ is not a unit. This proves the equivalence of (1) and (3).
The equivalence of (1) and (2) follows because the closure of $\{\mathfrak p\}$ is
$V(\mathfrak p)$. Indeed, this is a closed set containing
$\mathfrak p$, and any closed set $V(T)$ containing $\mathfrak p$
has $T\subset\mathfrak p$, hence contains $V(\mathfrak p)$.
A maximal ideal therefore gives a closed point. Conversely, if
$\{\mathfrak p\}$ is closed, apply (2) of
Lemma \ref{editorial-source-lemma-Zariski-topology} to $R/\mathfrak p$.
The inverse image of a maximal ideal of that quotient is a maximal
ideal $\mathfrak m$ of $R$ containing $\mathfrak p$.
Then $\mathfrak m\in V(\mathfrak p)=\{\mathfrak p\}$,
so $\mathfrak m=\mathfrak p$.
If $R$ is local then (4) holds since $x$ is either in $\mathfrak m$
or not. If (4) holds, and $\mathfrak m$, $\mathfrak m'$ are distinct
maximal ideals then we may choose $x \in R$ such that
$x \bmod \mathfrak m' = 0$ and $x \bmod \mathfrak m = 1$
by the Chinese remainder theorem
(Lemma \ref{algebra-lemma-chinese-remainder}).
This element $x$ is not invertible and neither is $1 - x$ which is
a contradiction. Thus (4) and (1) are equivalent.
\end{proof}


\subsection{Editorial treatment of Algebra lines 3499--3517}

\hypertarget{algebra-context-015}{}

\noindent\href{https://github.com/stacks/stacks-project/blob/a04446e57ec1fbc252a871afcec7752fb2807b14/algebra.tex\#L3499}{Original source, lines 3499--3517}. Complete original and reviewed TeX passages accompany this note. Every intervention is separately recorded; the following treatment is editorial.


\begin{lemma}
\label{editorial-source-lemma-in-image}
Let $\varphi : R \to S$ be a ring map. Let $\mathfrak p$
be a prime of $R$. The following are equivalent
\begin{enumerate}
\item $\mathfrak p$ is in the image of
$\Spec(S) \to \Spec(R)$,
\item $S \otimes_R \kappa(\mathfrak p) \not = 0$,
\item $S_{\mathfrak p}/\mathfrak p S_{\mathfrak p} \not = 0$,
\item $(S/\mathfrak pS)_{\mathfrak p} \not = 0$, and
\item $\mathfrak p = \varphi^{-1}(\mathfrak pS)$.
\end{enumerate}
\end{lemma}

\begin{proof}
We have already seen the equivalence of the first two
in Remark \ref{algebra-remark-fundamental-diagram}.
Put $T=R\setminus\mathfrak p$ and $Q=S/\mathfrak pS$, with
$T$ acting on $Q$ through $\varphi$.
The identifications in that remark identify the three rings in
(2), (3), and (4) with $T^{-1}Q$. In particular the first map is
$$
S\otimes_R\kappa(\mathfrak p)\longrightarrow T^{-1}Q,\qquad
s\otimes\overline{a/u}\longmapsto
\frac{\varphi(a)s+\mathfrak pS}{u},
$$
where $a\in R$, $u\in T$, and the bar denotes the class in
$R_{\mathfrak p}/\mathfrak pR_{\mathfrak p}$.
The formula is $R$-balanced, respects multiplication, kills
$\mathfrak pR_{\mathfrak p}$, and inverts every member of $T$.
Its inverse sends $(s+\mathfrak pS)/u$ to
$s\otimes\overline{1/u}$: the relations defining the quotient
and localization hold because $\mathfrak p$ acts as zero and
each $u\in T$ as a unit in the second tensor factor.
Both composites are the identity on the displayed generators.
The corresponding quotient-localization map is
$s/\varphi(u)+\mathfrak pS_{\mathfrak p}\mapsto
(s+\mathfrak pS)/u$, with the reverse formula as inverse.
Finally, $T^{-1}Q$ is the zero ring if and only if there is
$u\in T$ with $\varphi(u)\in\mathfrak pS$, by the equality
criterion for $1/1=0/1$ in a localization.
The ideal $\varphi^{-1}(\mathfrak pS)$ always contains
$\mathfrak p$. Hence it misses $T$ if and only if it equals
$\mathfrak p$. This proves the equivalence with (5).
\end{proof}


\subsection{Editorial treatment of Algebra lines 3519--3538}

\hypertarget{algebra-context-016}{}

\noindent\href{https://github.com/stacks/stacks-project/blob/a04446e57ec1fbc252a871afcec7752fb2807b14/algebra.tex\#L3519}{Original source, lines 3519--3538}. Complete original and reviewed TeX passages accompany this note. Every intervention is separately recorded; the following treatment is editorial.


\begin{remark}
\label{editorial-source-remark-local-ring-fibre}
Let $R \to S$ be a ring map. Let $\mathfrak q$ be a prime
ideal of $S$ lying over the prime ideal $\mathfrak p$ of $R$.
According to Remark \ref{algebra-remark-fundamental-diagram}
the prime $\mathfrak q$ corresponds to a unique prime
$\overline{\mathfrak q}$ of the fibre ring
$F = S \otimes_R \kappa(\mathfrak p)$.
Then we have
$$
F_{\overline{\mathfrak q}}
\cong S_\mathfrak q \otimes_{R_\mathfrak p} \kappa(\mathfrak p)
\cong S_\mathfrak q/\mathfrak p S_\mathfrak q
$$
Namely, there is an obvious ring map
$F \to S_\mathfrak q \otimes_{R_\mathfrak p} \kappa(\mathfrak p)$
which, under the displayed isomorphism, identifies with $F \to F_{\overline{\mathfrak q}}$.
The second isomorphism follows from the fact that $\kappa(\mathfrak p)$
is the quotient of $R_\mathfrak p$ by $\mathfrak pR_\mathfrak p$.
Explicitly, write $T=R\setminus\mathfrak p$. The fibre-ring
description is $F=T^{-1}(S/\mathfrak pS)$, and the prime in it is
$\overline{\mathfrak q}=T^{-1}(\mathfrak q/\mathfrak pS)$.
For $a\in S$ and $u\in T$, the fraction
$(a+\mathfrak pS)/u$ belongs to $\overline{\mathfrak q}$
if and only if $a\in\mathfrak q$; here every $\varphi(u)$
is outside $\mathfrak q$.
Thus the mutually inverse maps between $F_{\overline{\mathfrak q}}$
and $S_{\mathfrak q}/\mathfrak pS_{\mathfrak q}$ are
$$
\frac{(a+\mathfrak pS)/u}{(b+\mathfrak pS)/v}
\longmapsto
\frac{a\varphi(v)}{\varphi(u)b}+\mathfrak pS_{\mathfrak q},
\qquad
\frac{a}{b}+\mathfrak pS_{\mathfrak q}
\longmapsto
\frac{(a+\mathfrak pS)/1}{(b+\mathfrak pS)/1},
$$
where $u,v\in T$ and $b\in S\setminus\mathfrak q$.
The first map is induced by $S\to S_{\mathfrak q}/\mathfrak pS_{\mathfrak q}$:
it kills $\mathfrak pS$, inverts $T$, and inverts exactly the
additional denominators displayed above. The second is induced by
$S\to F_{\overline{\mathfrak q}}$, which inverts
$S\setminus\mathfrak q$ and kills $\mathfrak pS$.
The formulas show that both composites fix every fraction.
The tensor comparison sends
$$
\frac{a}{b}\otimes\overline{r/t}
\longmapsto
\frac{a\varphi(r)}{b\varphi(t)}+\mathfrak pS_{\mathfrak q},
$$
for $r\in R$, $t\in T$, and has inverse
$a/b+\mathfrak pS_{\mathfrak q}\mapsto(a/b)\otimes\overline{1}$.
Balancing over $R_{\mathfrak p}$ and the vanishing of
$\mathfrak pR_{\mathfrak p}$ prove well-definedness and both
inverse identities. These maps identify the original map out of
$F$ with localization at $\overline{\mathfrak q}$.
\end{remark}


\subsection{Editorial treatment of Algebra lines 3586--3601}

\hypertarget{algebra-context-017}{}

\noindent\href{https://github.com/stacks/stacks-project/blob/a04446e57ec1fbc252a871afcec7752fb2807b14/algebra.tex\#L3586}{Original source, lines 3586--3601}. Complete original and reviewed TeX passages accompany this note. Every intervention is separately recorded; the following treatment is editorial.


\begin{lemma}
\label{editorial-source-lemma-surjective-on-spec-units}
Let $\varphi : R \to S$ be a ring map such that the induced map
$\Spec(S) \to \Spec(R)$ is surjective. Then an element $x \in R$
is a unit if and only if $\varphi(x) \in S$ is a unit.
The same conclusion holds under the weaker assumption that the image
of $\Spec(S)\to\Spec(R)$ contains every closed point of $\Spec(R)$.
\end{lemma}

\begin{proof}
If $x$ is a unit, then so is $\varphi(x)$. Conversely, if $\varphi(x)$
is a unit, then $\varphi(x) \not \in \mathfrak q$ for all
$\mathfrak q \in \Spec(S)$. Hence
$x \not \in \varphi^{-1}(\mathfrak q) = \Spec(\varphi)(\mathfrak q)$
for all $\mathfrak q \in \Spec(S)$. Since $\Spec(\varphi)$ is surjective
we conclude that $x$ is a unit by
part (17) of Lemma \ref{editorial-source-lemma-Zariski-topology}.
For the weaker assumption, suppose that $x$ is not a unit.
Then $xR$ is proper, and a maximal ideal of the nonzero ring $R/xR$
pulls back to a maximal ideal $\mathfrak m$ of $R$ containing $x$.
This is a closed point, so the assumption supplies
$\mathfrak q\in\Spec(S)$ with
$\varphi^{-1}(\mathfrak q)=\mathfrak m$.
Consequently $\varphi(x)\in\mathfrak q$, so $\varphi(x)$ is not a unit.
The forward implication still follows by applying $\varphi$ to an inverse.
The assumption is strictly weaker: for a prime integer $p$, the residue map
$\mathbf Z_{(p)}\to\mathbf F_p$ hits the unique closed point $(p)$
of $\Spec(\mathbf Z_{(p)})$, but misses its prime ideal $(0)$.
Indeed, localization of the integer prime ideals leaves exactly $(0)$
and $(p)$, and the only prime of the field $\mathbf F_p$ contracts to $(p)$.
\end{proof}


\subsection{Editorial treatment of Algebra lines 3698--3733}

\hypertarget{algebra-context-018}{}

\noindent\href{https://github.com/stacks/stacks-project/blob/a04446e57ec1fbc252a871afcec7752fb2807b14/algebra.tex\#L3698}{Original source, lines 3698--3733}. Complete original and reviewed TeX passages accompany this note. Every intervention is separately recorded; the following treatment is editorial.


\begin{lemma}
\label{editorial-source-lemma-NAK-localization}
Let $R$ be a ring, let $S \subset R$ be a multiplicative subset,
let $I \subset R$ be an ideal, and let $M$ be a finite $R$-module.
If $x_1, \ldots, x_r \in M$ generate $S^{-1}(M/IM)$
as an $S^{-1}(R/I)$-module, then there exists an $f \in S + I$
such that $x_1, \ldots, x_r$ generate $M_f$ as an
$R_f$-module.\footnote{Special cases: (I) $I = 0$. The lemma says
if $x_1, \ldots, x_r$ generate $S^{-1}M$, then $x_1, \ldots, x_r$
generate $M_f$ for some $f \in S$. (II) $I = \mathfrak p$ is
a prime ideal and $S = R \setminus \mathfrak p$. The lemma says if
$x_1, \ldots, x_r$ generate $M \otimes_R \kappa(\mathfrak p)$
then $x_1, \ldots, x_r$ generate $M_f$ for some
$f \in R$, $f \not \in \mathfrak p$.}
\end{lemma}

\begin{proof}
Special case $I = 0$. Let $y_1, \ldots, y_s$ be generators for $M$ over $R$.
Since $S^{-1}M$ is generated by $x_1, \ldots, x_r$, for each $i$
we can write $y_i = \sum (a_{ij}/s_{ij})x_j$ in $S^{-1}M$
for some $a_{ij} \in R$ and $s_{ij} \in S$. Multiplying by the product
$s=\prod_{k=1}^{s_0}\prod_{\ell=1}^r s_{k\ell}\in S$,
where $s_0$ denotes the number of chosen generators $y_i$, we see that
$sy_i = \sum_{j=1}^r a'_{ij}x_j$ in $S^{-1}M$, with
$$
a'_{ij}=a_{ij}
\prod_{\substack{1\leq k\leq s_0,\ 1\leq\ell\leq r\\(k,\ell)\ne(i,j)}}
s_{k\ell}\in R.
$$
Empty products are $1$ and empty sums are $0$.
This in turn means there exist $t_i \in S$ such that
$t_isy_i = \sum t_ia'_{ij}x_j$ in $M$. Thus if $t \in S$
is the product of the $t_i$, put
$t^{(i)}=\prod_{k\ne i}t_k$, so $t=t_it^{(i)}$.
Multiplying the equality in $M$ by $t^{(i)}$ gives
$sty_i=\sum_{j=1}^r t a'_{ij}x_j$.
In $M_{st}$ it follows that
$y_i=\sum_{j=1}^r(t a'_{ij}/(st))x_j$.
Thus $y_i$ is in the $R_{st}$-submodule generated by
$x_1,\ldots,x_r$. Hence $x_1, \ldots, x_r$ generate $M_{st}$.

\medskip\noindent
General case. By the special case, we can find an $s \in S$
such that $x_1, \ldots, x_r$ generate $(M/IM)_s$ over $(R/I)_s$.
By Lemma \ref{algebra-lemma-NAK} we can find a $g \in 1 + I_s \subset R_s$
such that $x_1, \ldots, x_r$ generate $(M_s)_g$ over $(R_s)_g$.
Write $g=1+i/s'$ with $i\in I$ and $s'=s^n$ for some
$n\geq0$. Put $f=ss'+is=s(s'+i)$. Since $ss'=s^{n+1}\in S$
and $is\in I$, we have $f\in S+I$.
In $(R_s)_g$, the image of $f$ equals $ss'g$ and is a unit.
Conversely, in $R_f$ the element $s$ has inverse $(s'+i)/f$,
and the image of $g$ has inverse $ss'/f$:
$$
s\frac{s'+i}{f}=1,\qquad
\left(1+\frac{i}{s'}\right)\frac{ss'}{f}=1.
$$
The two universal properties therefore give mutually inverse maps
$R_f\longleftrightarrow(R_s)_g$ which agree with the original
maps from $R$. More explicitly, the first sends $r/f^k$ to
$(r/1)/(ss'g)^k$. The reverse map sends
$r/s^a$ to $r(s'+i)^a/f^a$ and $g^{-1}$ to $ss'/f$.
Their composites fix the image of every $r\in R$ and each inverse
adjoined in the respective localizations, so both composites are identities.
The corresponding module maps are also inverse: they send
$m/f^k$ to $(m/1)/(ss'g)^k$ and send
$(m/s^a)/g^b$ to $m(s'+i)^a(ss')^b/f^{a+b}$.
They preserve the images of all the chosen $x_j$.
Consequently those images generate $M_f$ over $R_f$, as required.
These localization identities remain valid when a localization is
the zero ring and when the original lists of generators are empty.
\end{proof}


\subsection{Editorial treatment of Algebra lines 4007--4031}

\hypertarget{algebra-context-019}{}

\noindent\href{https://github.com/stacks/stacks-project/blob/a04446e57ec1fbc252a871afcec7752fb2807b14/algebra.tex\#L4007}{Original source, lines 4007--4031}. Complete original and reviewed TeX passages accompany this note. Every intervention is separately recorded; the following treatment is editorial.


\begin{lemma}
\label{editorial-source-lemma-connected-component}
Let $R$ be a ring.
A connected component of
$\Spec(R)$ is of the form $V(I)$,
where $I$ is an ideal generated by idempotents
such that every idempotent of $R$ either maps to
$0$ or $1$ in $R/I$.
Conversely, if $I$ is a proper ideal generated by idempotents and
every idempotent of $R$ maps to $0$ or $1$ in $R/I$, then $V(I)$
is a connected component of $\Spec(R)$.
\end{lemma}

\begin{proof}
Let $\mathfrak p$ be a prime of $R$. By
Lemma \ref{algebra-lemma-topology-spec}
we see that the hypotheses of
Topology, Lemma \ref{topology-lemma-connected-component-intersection}
are satisfied for the topological space $\Spec(R)$.
Hence the connected component of $\mathfrak p$ in $\Spec(R)$
is the intersection of open and closed subsets of $\Spec(R)$
containing $\mathfrak p$. Hence it equals $V(I)$ where
$I$ is generated by the idempotents $e \in R$ such that $e$ maps to $0$
in $\kappa(\mathfrak p)$, see
Lemma \ref{algebra-lemma-disjoint-decomposition}.
Any idempotent $e$ which is not in this collection maps to $1$
in $\kappa(\mathfrak p)$. Its complementary idempotent $1-e$
therefore belongs to $I$, so $e$ maps to $1$ in $R/I$.

\medskip\noindent
For the converse, we first show that idempotents lift modulo any ideal
$I$ generated by idempotents. Let $\overline a\in R/I$ be idempotent
and choose a lift $a\in R$. Write
$a^2-a=\sum_{j=1}^m r_je_j$, where $r_j\in R$ and the $e_j\in I$
are idempotent. Define
$$
c=\prod_{j=1}^m(1-e_j),\qquad
e=1-c=
\sum_{\emptyset\ne J\subset\{1,\ldots,m\}}
(-1)^{|J|+1}\prod_{j\in J}e_j.
$$
Every factor $1-e_j$ is idempotent, so $c^2=c$ and $e^2=e$.
The expanded sum shows $e\in I$. For each $j$ the product $ce_j$
contains the factor $(1-e_j)e_j=0$, hence
$c(a^2-a)=\sum_{j=1}^m r_jce_j=0$.
It follows that $b=ca$ satisfies
$b^2-b=c^2a^2-ca=c(a^2-a)=0$ and $b-a=-ea\in I$.
Thus $b$ is an idempotent lift of $\overline a$.
This also covers $m=0$, when $c=1$, $e=0$, and $a$ is already idempotent.
Under the hypothesis of the converse, every such $b$ maps to $0$ or $1$,
so $R/I$ has no nontrivial idempotents. Since $I$ is proper, $R/I$ is
nonzero. Lemmas \ref{algebra-lemma-characterize-spec-connected} and
\ref{algebra-lemma-spec-closed} show that $V(I)$ is nonempty and connected.
By Lemma \ref{algebra-lemma-closed-union-connected-components}, it is also
a union of connected components of $\Spec(R)$. A nonempty connected
subset lies in a single connected component; since $V(I)$ is a union
of such components, it is exactly that component.
\end{proof}


\subsection{Editorial treatment of Algebra lines 4133--4217}

\hypertarget{algebra-context-020}{}

\noindent\href{https://github.com/stacks/stacks-project/blob/a04446e57ec1fbc252a871afcec7752fb2807b14/algebra.tex\#L4133}{Original source, lines 4133--4217}. Complete original and reviewed TeX passages accompany this note. Every intervention is separately recorded; the following treatment is editorial.


\begin{lemma}
\label{editorial-source-lemma-cover}
\begin{slogan}
Zariski-local properties of modules and algebras
\end{slogan}
Let $R$ be a ring. Let $M$ be an $R$-module. Let $S$ be an $R$-algebra.
Suppose that $f_1, \ldots, f_n$ is a finite list of
elements of $R$ such that $\bigcup D(f_i) = \Spec(R)$,
in other words $(f_1, \ldots, f_n) = R$.
\begin{enumerate}
\item If each $M_{f_i} = 0$ then $M = 0$.
\item If each $M_{f_i}$ is a finite $R_{f_i}$-module,
then $M$ is a finite $R$-module.
\item If each $M_{f_i}$ is a finitely presented $R_{f_i}$-module,
then $M$ is a finitely presented $R$-module.
\item Let $M \to N$ be a map of $R$-modules. If $M_{f_i} \to N_{f_i}$
is an isomorphism for each $i$ then $M \to N$ is an isomorphism.
\item Let $0 \to M'' \to M \to M' \to 0$ be a complex of $R$-modules.
If $0 \to M''_{f_i} \to M_{f_i} \to M'_{f_i} \to 0$ is exact for each $i$,
then $0 \to M'' \to M \to M' \to 0$ is exact.
\item If each $R_{f_i}$ is Noetherian, then $R$ is Noetherian.
\item If each $S_{f_i}$ is a finite type $R_{f_i}$-algebra, then
$S$ is a finite type $R$-algebra.
\item If each $S_{f_i}$ is of finite presentation over $R_{f_i}$, then
$S$ is a finitely presented $R$-algebra.
\end{enumerate}
\end{lemma}

\begin{proof}
We prove each of the parts in turn.
\begin{enumerate}
\item By Proposition \ref{algebra-proposition-localize-twice}
this implies $M_\mathfrak p = 0$ for all $\mathfrak p \in \Spec(R)$,
so we conclude by Lemma \ref{algebra-lemma-characterize-zero-local}.
\item For each $i$ take a finite generating set $X_i$ of $M_{f_i}$.
For each $x\in X_i$, write $x=y_x/f_i^{a_x}$ with
$y_x\in M$ and $a_x\geq0$, and put $Y_i=\{y_x\mid x\in X_i\}$.
Keep the original elements $x$: the image of $y_x$ is $f_i^{a_x}x$,
and $x=f_i^{-a_x}(y_x/1)$ in $M_{f_i}$.
Since $f_i$ is a unit in $R_{f_i}$, the images of $Y_i$ and the
original set $X_i$ generate the same $R_{f_i}$-submodule. Let $Y$
be the union of these sets. This is still a finite set.
Consider the obvious $R$-linear map $R^Y \rightarrow M$ sending the basis
element $e_y$ to $y$. By assumption this map is surjective after localizing
at an arbitrary prime ideal $\mathfrak p$ of $R$, so it is surjective by
Lemma \ref{algebra-lemma-characterize-zero-local}
and $M$ is finitely generated.
\item By (2) we have a short exact sequence
$$
0 \rightarrow K \rightarrow R^m \rightarrow M \rightarrow 0
$$
Since localization is an exact functor and $M_{f_i}$ is finitely
presented we see that $K_{f_i}$ is finitely generated for all
$1 \leq i \leq n$ by Lemma \ref{algebra-lemma-extension}.
By (2) this implies that $K$ is a finite $R$-module and therefore
$M$ is finitely presented.
\item By Proposition \ref{algebra-proposition-localize-twice}
the assumption implies that the induced morphism
on localizations at all prime ideals is an isomorphism, so we conclude
by Lemma \ref{algebra-lemma-characterize-zero-local}.
\item By Proposition \ref{algebra-proposition-localize-twice} the assumption
implies that the induced
sequence of localizations at all prime ideals is short exact, so we
conclude by Lemma \ref{algebra-lemma-characterize-zero-local}.
\item We will show that every ideal of $R$ has a finite generating set:
For this, let $I \subset R$ be an arbitrary ideal. By
Proposition \ref{algebra-proposition-localization-exact}
each $I_{f_i} \subset R_{f_i}$ is an ideal. These are all
finitely generated by assumption, so we conclude by (2).
\item For each $i$ take a finite generating set $X_i$ of $S_{f_i}$.
For each $x\in X_i$, write $x=y_x/f_i^{a_x}$ with
$y_x\in S$ and $a_x\geq0$, and put $Y_i=\{y_x\mid x\in X_i\}$.
The equalities $y_x/1=f_i^{a_x}x$ and
$x=f_i^{-a_x}(y_x/1)$ show that the images of $Y_i$ generate
the same $R_{f_i}$-subalgebra as $X_i$, because both scalar factors
belong to the base ring $R_{f_i}$. Let $Y$ be the union of the
finite sets $Y_i$; it is finite.
Consider the algebra homomorphism $R[X_y]_{y \in Y} \rightarrow S$
induced by $Y$. Since it is an algebra homomorphism, the image $T$
is an $R$-submodule of the $R$-module $S$, so we can consider the
quotient module $S/T$. By assumption, this is zero if we localize
at the $f_i$, so it is zero by (1) and therefore $S$ is an
$R$-algebra of finite type.
\item By the previous item, there exists a surjective $R$-algebra
homomorphism $R[X_1, \ldots, X_n] \rightarrow S$. Let $K$ be the kernel
of this map. This is an ideal in $R[X_1, \ldots, X_n]$, finitely generated
in each localization at $f_i$. Since the $f_i$ generate the unit ideal
in $R$, they also generate the unit ideal in $R[X_1, \ldots, X_n]$, so an
application of (2) finishes the proof.
\end{enumerate}
\end{proof}


\subsection{Editorial treatment of Algebra lines 4306--4351}

\hypertarget{algebra-context-021}{}

\noindent\href{https://github.com/stacks/stacks-project/blob/a04446e57ec1fbc252a871afcec7752fb2807b14/algebra.tex\#L4306}{Original source, lines 4306--4351}. Complete original and reviewed TeX passages accompany this note. Every intervention is separately recorded; the following treatment is editorial.


\begin{lemma}
\label{editorial-source-lemma-cover-module}
Let $R$ be a ring. Let $f_1, \ldots, f_n$ be elements of $R$
generating the unit ideal. Let $M$ be an $R$-module.
The sequence
$$
0 \to
M \xrightarrow{\alpha}
\bigoplus\nolimits_{i = 1}^n M_{f_i} \xrightarrow{\beta}
\bigoplus\nolimits_{i, j = 1}^n M_{f_i f_j}
$$
is exact, where $\alpha(m) = (m/1, \ldots, m/1)$
and $\beta(m_1/f_1^{e_1}, \ldots, m_n/f_n^{e_n})
= (m_i/f_i^{e_i} - m_j/f_j^{e_j})_{(i, j)}$.
\end{lemma}

\begin{proof}
It suffices to prove exactness after localization at each maximal
ideal $\mathfrak m$, by Lemma \ref{algebra-lemma-characterize-zero-local}.
The finite direct sums commute with localization. For each original
index $i$, the comparison of the two localization orders is
$$
(M_{f_i})_{\mathfrak m}\longrightarrow
(M_{\mathfrak m})_{f_i/1},\qquad
\frac{x/f_i^a}{s}\longmapsto\frac{x/s}{(f_i/1)^a},
$$
where $s\in R\setminus\mathfrak m$. Its reverse formula is the
inverse, by Proposition \ref{editorial-source-proposition-localize-twice-module}.
The same formulas with $f_if_j$ give the comparisons on the last term.
They commute with the original maps $\alpha$ and $\beta$, including
the difference in each ordered pair $(i,j)$.

\medskip\noindent
Choose an original index $k$ with $f_k\notin\mathfrak m$, which exists
because $(f_1,\ldots,f_n)=R$. Keep all indices and all elements $f_i$.
The natural map
$\iota_k:M_{\mathfrak m}\to(M_{f_k})_{\mathfrak m}$
has the explicit inverse
$$
\theta_k\left(\frac{x/f_k^a}{s}\right)=\frac{x}{sf_k^a}
\quad\text{in }M_{\mathfrak m}.
$$
The denominator $sf_k^a$ lies outside $\mathfrak m$.
For each $i$, the natural map
$(M_{f_i})_{\mathfrak m}\to(M_{f_kf_i})_{\mathfrak m}$
is also an isomorphism, with inverse
$$
\theta_{k,i}\left(\frac{x/(f_kf_i)^a}{s}\right)
=\frac{x/f_i^a}{sf_k^a}.
$$
Its forward formula sends $(x/f_i^a)/s$ to
$(f_k^ax/(f_kf_i)^a)/s$. The two composites are identities by
the equality-of-fractions relations. These maps come from localization
at elements which already act invertibly, so the formulas are
well-defined and preserve the original module structures.

\medskip\noindent
Projection to the $k$th summand followed by $\theta_k$ is a left
inverse of $\alpha_{\mathfrak m}$, proving injectivity.
Now let $(z_i)_{i=1}^n$ be in the kernel of $\beta_{\mathfrak m}$
and put $z=\theta_k(z_k)\in M_{\mathfrak m}$.
The $(k,i)$ coordinate says that the images of $z_k$ and $z_i$
in $(M_{f_kf_i})_{\mathfrak m}$ agree. Applying $\theta_{k,i}$
shows that $z_i$ is the image of $z$ in $(M_{f_i})_{\mathfrak m}$.
For completeness, on a representative $z_k=(x/f_k^a)/s$ the image
in the pair localization is $(f_i^ax/(f_kf_i)^a)/s$; applying
$\theta_{k,i}$ gives $(f_i^ax/f_i^a)/(sf_k^a)$, which is precisely
the image of $x/(sf_k^a)=z$. Thus
$\alpha_{\mathfrak m}(z)=(z_i)_{i=1}^n$.
The reverse inclusion follows from the original identity
$\beta\alpha=0$, so the localized sequence is exact.
The local exactness criterion proves the original sequence exact.
If the covering list is empty, the unit-ideal hypothesis gives $R=0$,
and every unital $R$-module is zero, so the same assertion holds.
\end{proof}


\subsection{Editorial treatment of Algebra lines 4438--4523}

\hypertarget{algebra-context-022}{}

\noindent\href{https://github.com/stacks/stacks-project/blob/a04446e57ec1fbc252a871afcec7752fb2807b14/algebra.tex\#L4438}{Original source, lines 4438--4523}. Complete original and reviewed TeX passages accompany this note. Every intervention is separately recorded; the following treatment is editorial.


\begin{lemma}
\label{editorial-source-lemma-glue-modules}
Let $R$ be a ring. Let $f_1, \ldots, f_n \in R$. Suppose we are given
the following data:
\begin{enumerate}
\item For each $i$ an $R_{f_i}$-module $M_i$.
\item For each pair $i, j$ an $R_{f_if_j}$-module isomorphism
$\psi_{ij} : (M_i)_{f_j} \to (M_j)_{f_i}$.
\end{enumerate}
which satisfy the ``cocycle condition'' that all the diagrams
$$
\xymatrix{
(M_i)_{f_jf_k}
\ar[rd]_{\psi_{ij}}
\ar[rr]^{\psi_{ik}}
& &
(M_k)_{f_if_j} \\
&
(M_j)_{f_if_k} \ar[ru]_{\psi_{jk}}
}
$$
commute (for all triples $i, j, k$). Given this data define
$$
M = \Ker\left(
\bigoplus\nolimits_{1 \leq i \leq n} M_i
\longrightarrow
\bigoplus\nolimits_{1 \leq i, j \leq n} (M_i)_{f_j}
\right)
$$
where $(m_1, \ldots, m_n)$ maps to the element whose
$(i, j)$th entry is $m_i/1 - \psi_{ji}(m_j/1)$.
Then the natural map $M \to M_i$ induces an isomorphism
$M_{f_i} \to M_i$. Moreover $\psi_{ij}(m/1) = m/1$
for all $m \in M$ (with obvious notation).
\end{lemma}

\begin{proof}
Fix an original index $k\in\{1,\ldots,n\}$.
Localizing the defining kernel at $f_k$ and using exactness of
localization gives the isomorphism
$$
M_{f_k}\longrightarrow H_k:=
\Ker\left(
\bigoplus_{i=1}^n(M_i)_{f_k}\longrightarrow
\bigoplus_{i,j=1}^n((M_i)_{f_j})_{f_k}
\right),\qquad
\frac{(m_i)_{i=1}^n}{f_k^a}\longmapsto
\left(\frac{m_i}{f_k^a}\right)_{i=1}^n.
$$
The map defining $H_k$ keeps the original $(i,j)$ component
$z_i/1-\psi_{ji}(z_j/1)$, after the indicated localizations.
The comparison between localization orders sends
$(x/f_j^a)/f_k^b$ to $(x/f_k^b)/f_j^a$ and has the reverse
formula as inverse, by Proposition \ref{editorial-source-proposition-localize-twice-module}.

\medskip\noindent
Since $M_k$ is already an $R_{f_k}$-module, the canonical map
$M_k\to(M_k)_{f_k}$ has inverse
$x/f_k^a\mapsto f_k^{-a}x$.
The cocycle condition for $(k,k,k)$ gives $\psi_{kk}^2=\psi_{kk}$
under these comparisons; since $\psi_{kk}$ is invertible, it is the
identity. Likewise the conditions with repeated indices give
$\psi_{ik}\psi_{ki}=\mathrm{id}$ on each overlap.
Define
$$
\eta_k:M_k\longrightarrow\bigoplus_{i=1}^n(M_i)_{f_k},\qquad
a\longmapsto\bigl(\psi_{ki}(a/1)\bigr)_{i=1}^n,
$$
where $a/1$ is the image of $a$ in $(M_k)_{f_i}$ for the $i$th entry.
The cocycle for $(k,j,i)$ says
$\psi_{ji}\psi_{kj}=\psi_{ki}$ after localizing at $f_if_j$,
so every defining difference for $H_k$ vanishes on this tuple.
Thus $\eta_k$ takes values in $H_k$.
Projection onto the $k$th summand, followed by
$(M_k)_{f_k}\to M_k$, defines $\pi_k:H_k\to M_k$.
The identity $\psi_{kk}=\mathrm{id}$ proves $\pi_k\eta_k=\mathrm{id}$.
For $z=(z_i)\in H_k$, its $(i,k)$ equation identifies $z_i$
with $\psi_{ki}(\pi_k(z)/1)$; hence $\eta_k\pi_k=\mathrm{id}$.

\medskip\noindent
There is also an explicit fraction representing this inverse in $M_{f_k}$.
Write $\psi_{ki}(a/1)=b_i/f_k^{r_i}$ with $b_i\in M_i$ and $r_i\geq0$.
Choose $N\geq r_i$ for all $i$, and put $b'_i=f_k^{N-r_i}b_i$.
For each pair $(i,j)$ the element
$d_{ij}=b'_i/1-\psi_{ji}(b'_j/1)$ of $(M_i)_{f_j}$
vanishes after localization at $f_k$.
Choose $s_{ij}\geq0$ with $f_k^{s_{ij}}d_{ij}=0$, and choose
$L\geq s_{ij}$ for all pairs. The tuple
$m=(f_k^Lb'_i)_{i=1}^n$ belongs to $M$, and the fraction
$m/f_k^{N+L}$ maps to $\eta_k(a)$.
Independence of these choices follows from the injectivity of the
displayed kernel comparison, since their images are the same tuple.
Consequently the original projection induces $M_{f_k}\cong M_k$.
Finally, for $m=(m_i)\in M$, its $(j,i)$ equation reads
$m_j/1=\psi_{ij}(m_i/1)$, which is the asserted compatibility.
The proof uses no cover assumption on the $f_i$ and changes none of them.
If $n=0$, then $M=0$ and there are no indexed assertions to prove.
\end{proof}


\subsection{Editorial treatment of Algebra lines 4592--4604}

\hypertarget{algebra-context-023}{}

\noindent\href{https://github.com/stacks/stacks-project/blob/a04446e57ec1fbc252a871afcec7752fb2807b14/algebra.tex\#L4592}{Original source, lines 4592--4604}. Complete original and reviewed TeX passages accompany this note. Every intervention is separately recorded; the following treatment is editorial.


\begin{lemma}
\label{editorial-source-lemma-total-ring-fractions}
Let $R$ be a ring.
Let $S \subset R$ be a multiplicative subset consisting of nonzerodivisors.
Then $Q(R) \cong Q(S^{-1}R)$.
In particular $Q(R) \cong Q(Q(R))$.
\end{lemma}

\begin{proof}
Let $T\subset R$ be the set of nonzerodivisors.
Since $S\subset T$, the map $R\to S^{-1}R$ is injective.
An element $x=r/f$ of $S^{-1}R$, with $r\in R$ and $f\in S$,
is a nonzerodivisor if and only if $r\in T$.
Indeed, if $x$ is a nonzerodivisor and $ra=0$ in $R$, then
$x(a/1)=0$, hence $a/1=0$, and injectivity gives $a=0$.
Conversely, if $r\in T$ and $x(a/s)=0$, then
$ura=0$ for some $u\in S$. Both $u$ and $r$ are nonzerodivisors,
so $a=0$ and $a/s=0$.
The two ring maps are
$$
Q(R)\longrightarrow Q(S^{-1}R),\qquad
\frac{r}{t}\longmapsto\frac{r/1}{t/1},
$$
and
$$
Q(S^{-1}R)\longrightarrow Q(R),\qquad
\frac{r/s}{a/u}\longmapsto\frac{ru}{sa},
\quad s,u\in S,\quad a\in T.
$$
The first is induced by $R\to S^{-1}R$ because each $t\in T$
remains a nonzerodivisor. The second is induced by $S^{-1}R\to Q(R)$:
it sends every nonzerodivisor $a/u$ to a unit with inverse $u/a$.
Thus both formulas are well-defined by the localization universal
properties. Their composites fix each element of $R$ and every
adjoined inverse; explicitly, $r/t$ returns to $r/t$, and
$(r/s)/(a/u)$ returns to $(ru/1)/(sa/1)$, the original fraction.
Taking $S=T$ gives the final assertion, with the same maps.
\end{proof}


\subsection{Editorial treatment of Algebra lines 4664--4695}

\hypertarget{algebra-context-024}{}

\noindent\href{https://github.com/stacks/stacks-project/blob/a04446e57ec1fbc252a871afcec7752fb2807b14/algebra.tex\#L4664}{Original source, lines 4664--4695}. Complete original and reviewed TeX passages accompany this note. Every intervention is separately recorded; the following treatment is editorial.


\begin{lemma}
\label{editorial-source-lemma-irreducible}
Let $R$ be a ring.
\begin{enumerate}
\item For a prime $\mathfrak p \subset R$ the closure
of $\{\mathfrak p\}$ in the Zariski topology is $V(\mathfrak p)$.
In a formula $\overline{\{\mathfrak p\}} = V(\mathfrak p)$.
\item The irreducible closed subsets of $\Spec(R)$ are
exactly the subsets $V(\mathfrak p)$, with $\mathfrak p \subset R$
a prime.
\item The irreducible components (see Topology,
Definition \ref{topology-definition-irreducible-components})
of $\Spec(R)$ are  exactly the subsets $V(\mathfrak p)$,
with $\mathfrak p \subset R$ a minimal prime.
\end{enumerate}
\end{lemma}

\begin{proof}
Note that if $ \mathfrak p \in V(I)$, then
$I \subset \mathfrak p$. Hence,
clearly $\overline{\{\mathfrak p\}} = V(\mathfrak p)$.
In particular $V(\mathfrak p)$ is the closure of
a singleton and hence irreducible.
The second assertion implies the third.
To show the second, let
$V(I) \subset \Spec(R)$ with $I$ a radical ideal.
If $I=R$, then $V(I)=\emptyset$, which is not irreducible.
We may therefore suppose that $I$ is proper.
If $I$ is not prime, then choose $a, b\in R$, $a, b\not \in I$
with $ab\in I$. In this case $V(I, a) \cup V(I, b) = V(I)$,
but neither $V(I, b) = V(I)$ nor $V(I, a) = V(I)$, by
Lemma \ref{editorial-source-lemma-Zariski-topology}. Hence $V(I)$ is not
irreducible.
\end{proof}


\subsection{Editorial treatment of Algebra lines 4863--4887}

\hypertarget{algebra-context-025}{}

\noindent\href{https://github.com/stacks/stacks-project/blob/a04446e57ec1fbc252a871afcec7752fb2807b14/algebra.tex\#L4863}{Original source, lines 4863--4887}. Complete original and reviewed TeX passages accompany this note. Every intervention is separately recorded; the following treatment is editorial.


\begin{example}
\label{editorial-source-example-spec-Zx}
In this example we describe $X = \Spec(\mathbf{Z}[x])$.
Fix $\mathfrak p \in X$.
Let $\phi : \mathbf{Z} \to \mathbf{Z}[x]$ and notice
that $\phi^{-1}(\mathfrak p) \in \Spec(\mathbf{Z})$.
If $\phi^{-1}(\mathfrak p) = (q)$ for $q$ a prime number $q > 0$,
then $\mathfrak p$ corresponds to a prime in $(\mathbf{Z}/(q))[x]$.
The zero prime of that polynomial ring gives $\mathfrak p=(q)$.
Every nonzero prime is generated by an irreducible polynomial
$\overline{f_q}\in(\mathbf{Z}/(q))[x]$.
Choose any coefficient lift $f_q\in\mathbf{Z}[x]$ of the same degree.
The inverse image of $(\overline{f_q})$ is $(q,f_q)$: if the reduction
of $F\in\mathbf Z[x]$ is $\overline A\,\overline{f_q}$, choose a lift
$A\in\mathbf Z[x]$; then $F-Af_q\in q\mathbf Z[x]$.
The converse containment follows by reduction.
Thus the primes over $(q)$ are $(q)$ and the ideals $(q,f_q)$ with
$\overline{f_q}$ irreducible. The lift $f_q$ itself need not be
irreducible in $\mathbf Z[x]$: for $q=5$, the degree-one reduction
of $2x+2$ is irreducible, while $2x+2=2(x+1)$ in $\mathbf Z[x]$.
Now assume that $\phi^{-1}(\mathfrak p)=(0)$.
In this case, if $\mathfrak p = (0)$ there is nothing more to prove. Otherwise, $\mathfrak p$ contains nonconstant polynomials
and it contains no nonzero integer. Factoring any nonzero member
in the unique factorization domain $\mathbf Z[x]$ and using primality
gives an irreducible $f\in\mathfrak p$. It is nonconstant because
$\mathfrak p\cap\mathbf Z=(0)$, and it is primitive because a
nonunit integer content would be a nontrivial factor.
Every other nonconstant irreducible $g\in\mathfrak p$ is likewise primitive.
Gauss' lemma makes $f$ and $g$ irreducible in $\mathbf Q[x]$.
Moreover, if $g=(d/e)f$ with coprime nonzero integers $d,e$,
the equality $eg=df$ has contents $|e|$ and $|d|$, respectively.
Thus $|e|=|d|=1$, so association over $\mathbf Q[x]$ implies
association over $\mathbf Z[x]$ for these original polynomials.
If $f$ and $g$ were not associates, the Euclidean algorithm in
$\mathbf Q[x]$ would give $1=af+bg$ with $a,b\in\mathbf Q[x]$.
Choose a positive common denominator $m$ of every coefficient of
$a$ and $b$, and put $\bar a=ma$, $\bar b=mb$ in $\mathbf Z[x]$.
Then $m=\bar a f+\bar b g\in\mathfrak p\cap\mathbf Z$, a contradiction.
Therefore all irreducibles belonging to $\mathfrak p$ are associates of $f$.
For every nonzero $F\in\mathfrak p$, a factorization of $F$ has
an irreducible factor in $\mathfrak p$ by primality, and that factor
is a unit multiple of $f$. Thus $f$ divides every $F\in\mathfrak p$,
while $f\in\mathfrak p$, proving $\mathfrak p=(f)$.
Conversely, every nonconstant irreducible $f\in\mathbf Z[x]$
generates a prime ideal, because $\mathbf Z[x]$ is a unique
factorization domain. Its contraction to $\mathbf Z$ is zero:
a nonzero constant cannot be divisible by a positive-degree polynomial.
Together with the separately retained zero prime, this describes
all primes with zero contraction.
\end{example}


\subsection{Editorial treatment of Algebra lines 4889--4929}

\hypertarget{algebra-context-026}{}

\noindent\href{https://github.com/stacks/stacks-project/blob/a04446e57ec1fbc252a871afcec7752fb2807b14/algebra.tex\#L4889}{Original source, lines 4889--4929}. Complete original and reviewed TeX passages accompany this note. Every intervention is separately recorded; the following treatment is editorial.


\begin{example}
\label{editorial-source-example-spec-kxy}
In this example we describe $X = \Spec(k[x, y])$
when $k$ is an arbitrary field.
Clearly $(0)$ is prime, and any principal ideal generated by an
irreducible polynomial will also be a prime since $k[x, y]$ is a
unique factorization domain. Now assume $\mathfrak p$ is an
element of $X$ that is not principal. Since $k[x, y]$ is a
Noetherian UFD, the prime ideal $\mathfrak p$ can be generated
by a finite number of irreducible polynomials $(f_1, \ldots, f_n)$.
Now, I claim that if $f, g$ are irreducible polynomials in $k[x, y]$
that are not associates, then $(f, g) \cap k[x] \neq 0$. To do this,
if either $f$ or $g$ lies in $k[x]$, that nonzero polynomial
itself belongs to $(f,g)\cap k[x]$ and proves the assertion.
Otherwise both have positive degree in $y$. Each is primitive as
a polynomial in $y$ over $k[x]$, since a nonunit content would
factor the original irreducible polynomial in $k[x,y]$.
Gauss' lemma therefore makes both irreducible in $k(x)[y]$.
In this case it is enough to show that $f$ and $g$ are relatively prime when
viewed in $k(x)[y]$. In this case, $k(x)[y]$ is a Euclidean domain,
so by applying the Euclidean algorithm and clearing denominators, we
obtain $p = af + bg$ for $0 \ne p \in k[x]$ and $a, b \in k[x, y]$. Thus, assume this is not
the case, that is, that some nonunit $h \in k(x)[y]$ divides both
$f$ and $g$. Then, by Gauss's lemma, for some $a, b \in k(x)$ we
have $ah | f$ and $bh | g$ for $ah, bh \in k[x, y]$. By
irreducibility, $ah = f$ and
$bh = g$ (since $h \notin k(x)$). So, back in $k(x)[y]$, $f, g $
are associates, as $\frac{a}{b} g = f$. Since
$k(x)$ is the fraction field of $k[x]$, we can write $g = \frac{r}{s} f $
for elements $r , s \in k[x]$ sharing no common factors. This
implies that $sg = rf$ in $k[x, y]$ and so $s$ must divide $f$
since $k[x,y]$ is a UFD. If $s$ is a unit in $k[x]$, write
$s=c\in k^\times$. Then $g=(r/c)f$; irreducibility of $g$
forces $r/c$ to be a unit, contradicting that $f$ and $g$ are
not associates. If $s$ is not a unit, the divisibility $s\mid f$
and irreducibility of $f$ give $f=us$ for a unit $u\in k^\times$.
This would put $f$ in $k[x]$, contrary to its positive $y$-degree.
Both alternatives contradict the existence of the common divisor.
Thus $f$ and $g$ are relatively prime in $k(x)[y]$, and the
cleared Bezout identity gives a nonzero polynomial in $(f,g)\cap k[x]$.
To distinguish it from a selected irreducible factor below, denote
this resulting polynomial by $P$.
The irreducible generators of the nonprincipal prime $\mathfrak p$
include two which are not associates, since otherwise they generate
a principal ideal. Applying the argument to those two puts
$P\in\mathfrak p\cap k[x]$. It cannot be a nonzero constant,
because $\mathfrak p$ is proper. Write its full factorization as
$P=c\prod_{j=1}^d p_j$, with $c\in k^\times$ and irreducible
$p_j\in k[x]$, repeating factors according to their multiplicities.
Primality gives some $p_j\in\mathfrak p$; denote this chosen
factor by $p$. Thus $\mathfrak p$ contains the irreducible
polynomial $p\in k[x]\subset k[x,y]$.
By the above, $\mathfrak p$ corresponds to a prime in the ring
$k[x,y]/(p)\cong(k[x]/(p))[y]=k(\alpha)[y]$, where
$\alpha=x+(p)$ in the field $k[x]/(p)$.
The isomorphism sends $F(x,y)+(p)$ to $F(\alpha,y)$:
coefficientwise reduction is surjective and its kernel consists
exactly of the polynomials whose coefficients are multiples of $p$.
Every element of $k[x]/(p)$ is a polynomial in $\alpha$, and this
quotient is a field, so it is $k(\alpha)$. No leading coefficient
of the chosen $p$ is changed or omitted. This is a
PID, and so any prime ideal corresponds to $(0)$ or an
irreducible polynomial in $k(\alpha)[y]$. Thus, $\mathfrak p$
is of the form $(p)$ or $(p, f)$ where $f$ is a
polynomial in $k[x, y]$ that is irreducible in the quotient
$k[x, y]/(p)$.
\end{example}


\subsection{Editorial treatment of Algebra lines 4931--5083}

\hypertarget{algebra-context-027}{}

\noindent\href{https://github.com/stacks/stacks-project/blob/a04446e57ec1fbc252a871afcec7752fb2807b14/algebra.tex\#L4931}{Original source, lines 4931--5083}. Complete original and reviewed TeX passages accompany this note. Every intervention is separately recorded; the following treatment is editorial.


\begin{example}
\label{editorial-source-example-affine-open-not-standard}
Consider the ring
$$
R = \{ f \in \mathbf{Q}[z]\text{ with }f(0) = f(1) \}.
$$
Consider the map
$$
\varphi : \mathbf{Q}[A, B] \to R
$$
defined by $\varphi(A)=A(z)=z^2-z$ and
$\varphi(B)=B(z)=z^3-z^2=zA(z)$.
Put $F=A^3-B^2+AB$. The exact identity
$$
F(A(z),B(z))=A(z)^3-z^2A(z)^2+zA(z)^2
=A(z)^2\bigl(A(z)-z^2+z\bigr)=0
$$
shows $(F)\subset\Ker(\varphi)$.
Since the leading coefficient of $F$ as a polynomial in $B$ is $-1$,
division in $\mathbf Q[A][B]$ writes any polynomial as
$Q(A,B)F+C(A)+BD(A)$.
Suppose its remainder evaluates to zero, so
$C(A(z))+zA(z)D(A(z))=0$.
Substituting $1-z$ for $z$ fixes $A(z)$ and sends $B(z)$ to
$(1-z)A(z)$; thus also $C(A(z))+(1-z)A(z)D(A(z))=0$.
Subtracting gives $(2z-1)A(z)D(A(z))=0$ in the domain $\mathbf Q[z]$.
The first two factors are nonzero, hence $D(A(z))=0$.
Substitution $\mathbf Q[A]\to\mathbf Q[z]$, $A\mapsto z^2-z$,
is injective: a nonzero polynomial of degree $d$ has after substitution
degree $2d$ with the same nonzero leading coefficient.
Consequently $D=0$ and then $C=0$, proving $\Ker(\varphi)=(F)$.
This ideal is prime because the image is a subring of $\mathbf Q[z]$.
It follows that $F$ is irreducible: in a factorization $F=GH$, primality
puts one factor, say $G$, in $(F)$, and writing $G=FK$ and cancelling
the nonzero $F$ gives $KH=1$.
The surjectivity proved next therefore gives the induced isomorphism
$$
\overline\varphi:\mathbf Q[A,B]/(A^3-B^2+AB)\longrightarrow R.
$$

\medskip\noindent
To see that $\varphi$ is surjective, we must express any
$f\in R$ as a $\mathbf{Q}$-coefficient polynomial in $A(z) = z^2-z$
and $B(z) = z^3-z^2$. Note the relation $zA(z) = B(z)$. Let
$a = f(0) = f(1)$. Then $z(z-1)$ must divide $f(z)-a$, so we can write
$f(z) = z(z-1)g(z)+a = A(z)g(z)+a$.  If $\deg(g) < 2$, then
$g(z) = c_1z + c_0$ and $f(z) = A(z)(c_1z + c_0)+a = c_1B(z)+c_0A(z)+a$, so we
are done.  If $\deg(g)\geq 2$, then by the polynomial division
algorithm, we can write $g(z) = A(z)h(z)+b_1z + b_0$
($\deg(h)\leq\deg(g)-2$), so $f(z) = A(z)^2h(z)+b_1B(z)+b_0A(z)+a$.
Applying division to $h(z)$ and iterating, we obtain an expression
for $f(z)$ as a polynomial in $A(z)$ and $B(z)$; hence $\varphi$ is
surjective.

\medskip\noindent
Now let $a \in \mathbf{Q}$, $a \neq 0, \frac{1}{2}, 1$ and
consider
$$
R_a = \{ f \in \mathbf{Q}[z, \frac{1}{z-a}]\text{ with }f(0) = f(1)
\}.
$$
This is a finitely generated $\mathbf Q$-algebra. Here is a proof
with the stated generators. Put
$$
\delta=a^2-a,\qquad
w=z+\frac{\delta}{z-a},\qquad
D=\mathbf Q[z,1/(z-a)],\qquad
E=\mathbf Q[z^2-z,z^3-z,w].
$$
The values of the first two generators at $0$ and $1$ are $0$;
the values of $w$ are both $1-a$. Thus $E\subset R_a$.
Since $\delta\ne0$ and $(z-a)^{-1}=(w-z)/\delta$, we have
$D=E[z]$. The original relation $z^2=z+A(z)$, with $A(z)\in E$,
shows $D=E+Ez$ as an $E$-module.
Also $R_a=\mathbf Q+A(z)D$: if $f\in R_a$ and
$\lambda=f(0)=f(1)$, write $f-\lambda=P(z)/(z-a)^N$.
Its numerator vanishes at both $0$ and $1$, so $P(z)=z(z-1)H(z)$
for a polynomial $H$, giving $f-\lambda\in A(z)D$.
The reverse inclusion follows by evaluation.
Finally $A(z)\in E$ and
$zA(z)=z^3-z^2=(z^3-z)-(z^2-z)\in E$, hence
$A(z)D\subset E+E$ and $R_a\subset E$.
This proves $R_a=E$, with all three original generators retained.
We have the following inclusions:
$$
R\subset R_a\subset\mathbf{Q}[z, \frac{1}{z-a}], \quad
R\subset\mathbf{Q}[z]\subset\mathbf{Q}[z, \frac{1}{z-a}].
$$
Recall (Lemma \ref{editorial-source-lemma-spec-localization}) that for a ring $T$ and a
multiplicative subset $S\subset T$, the ring map $T \to S^{-1}T$
induces a map on spectra $\Spec(S^{-1}T) \to \Spec(T)$
which is a homeomorphism onto the subset
$$
\{\mathfrak p \in \Spec(T) \mid S \cap \mathfrak p = \emptyset\}
\subset \Spec(T).
$$
When $S = \{ 1, f, f^2, \ldots\}$ for some $f\in T$, this is
the open set $D(f)\subset \Spec(T)$.  We now verify a corresponding
property for the ring map $R \to R_a$: we will show that the map
$\theta : \Spec(R_a) \to \Spec(R)$ induced by inclusion
$R\subset R_a$ is a homeomorphism onto an open subset of
$\Spec(R)$ by verifying that $\theta$ is an injective local
homeomorphism.  We do so with respect to an open cover of
$\Spec(R_a)$ by two distinguished opens, as we now describe.
For any $r\in\mathbf{Q}$, let $\text{ev}_r : R \to \mathbf{Q}$ be the
homomorphism given by evaluation at $r$.  Note that for $r = 0$ and
$r = 1-a$, this can be extended to a homomorphism
$\text{ev}_r' : R_a \to \mathbf{Q}$ (the latter because $\frac{1}{z-a}$
is well-defined at $z = 1-a$, since $a\neq\frac{1}{2}$).  However,
$\text{ev}_a$ does not extend to $R_a$.  Write
$\mathfrak{m}_r = \Ker(\text{ev}_r)$. We have
$$
\mathfrak{m}_0 = (z^2-z, z^3-z),
$$
$$
\mathfrak{m}_a = ((z-1 + a)(z-a), (z^2-1 + a)(z-a)), \text{ and}
$$
$$
\mathfrak{m}_{1-a} = ((z-1 + a)(z-a), (z-1 + a)(z^2-a)).
$$
For every $r\in\mathbf Q$, the presentation already proved gives
$\mathfrak m_r=(A-r(r-1),B-r^2(r-1))$.
Indeed, evaluating $A$ and $B$ at those two constants gives a quotient
equal to $\mathbf Q$, so this ideal is exactly the evaluation kernel.
Write
$$
u=(z-1+a)(z-a)=A-\delta,\qquad
v=(z^2+z+2a-2)(z-a)=B+(2-a)A-2\delta.
$$
The displayed generators above are verified by the identities
$$
\begin{aligned}
z^3-z&=B+A,\\
(z^2-1+a)(z-a)&=(B-a\delta)+(1-a)(A-\delta),\\
(z-1+a)(z^2-a)&=(B-(1-a)\delta)+a(A-\delta).
\end{aligned}
$$
Each formula is reversible by subtracting the indicated multiple
of $A-\delta$, proving the asserted ideal equalities.
In particular $(u,v)=(A-\delta,B-a\delta)=\mathfrak m_a$,
since $v=(B-a\delta)+(2-a)(A-\delta)$.
Note that
$\mathfrak{m}_a$ is not in the image of $\theta$: we have
$$
(z^2 - z)^2(z - a)\left(\frac{a^2 - a}{z - a} + z\right) =
(z^2 - z)^2(a^2 - a) + (z^2 - z)^2(z - a)z
$$
The left hand side is in $\mathfrak m_a R_a$ because
$(z^2 - z)(z - a)$ is in $\mathfrak m_a$ and because
$(z^2 - z)(\frac{a^2 - a}{z - a} + z)$ is in $R_a$. Similarly
the element $(z^2 - z)^2(z - a)z$ is in $\mathfrak m_a R_a$
because $(z^2 - z)$ is in $R_a$ and $(z^2 - z)(z - a)$ is in $\mathfrak m_a$.
As $a \not \in \{0, 1\}$ we conclude that
$(z^2 - z)^2 \in \mathfrak m_a R_a$. Hence
no ideal $I$ of $R_a$ can satisfy $I \cap R = \mathfrak m_a$, as such
an $I$ would have to contain $(z^2 - z)^2$, which is in $R$ but not in
$\mathfrak m_a$. The distinguished open set
$D((z-1 + a)(z-a))\subset\Spec(R)$ is equal to the complement of
the closed set $\{\mathfrak{m}_a, \mathfrak{m}_{1-a}\}$.
Then check that $R_{(z-1 + a)(z-a)} = (R_a)_{(z-1 + a)(z-a)}$; calling
this localized ring $R'$, then, it follows that the map $R \to R'$
factors as $R \to R_a \to R'$.  By Lemma
\ref{editorial-source-lemma-spec-localization}, then, these maps express
$\Spec(R') \subset \Spec(R_a)$ and
$\Spec(R') \subset \Spec(R)$ as open subsets; hence
$\theta : \Spec(R_a) \to \Spec(R)$, when restricted to
$D((z-1 + a)(z-a))$, is a homeomorphism onto an open subset.
Similarly, $\theta$ restricted to
$D((z^2 + z + 2a-2)(z-a)) \subset \Spec(R_a)$ is a homeomorphism
onto the open subset $D((z^2 + z + 2a-2)(z-a)) \subset \Spec(R)$.
In $\Spec(R)$, the complement of this distinguished open consists
of $\mathfrak m_a$ and the one or two closed points associated to the
irreducible factors of $z^2+z+2a-2$; a repeated factor gives only one point.
In $\Spec(R_a)$ the point $\mathfrak m_a$ is absent.
The verification below identifies these zero sets as well.
Furthermore, the ideal in $R_a$ generated by the
elements $(z^2 + z + 2a-2)(z-a)$ and $(z-1 + a)(z-a)$ is all of $R_a$, so
these two distinguished open sets cover $\Spec(R_a)$. Hence in
order to show that $\theta$ is a homeomorphism onto
$\Spec(R)-\{\mathfrak{m}_a\}$, it suffices to show
that these one or two points can never equal $\mathfrak{m}_{1-a}$.
And this is indeed the case, since $1-a$ is a root of $z^2 + z + 2a-2$
if and only if $a = 0$ or $a = 1$, both of which do not occur.

\medskip\noindent
Here are the exact localization maps and the cover certificate.
As identities in the original rational-function field,
$$
uw=B+(a-1)\delta\in R,
\qquad
vw=zv+\delta(z^2+z+2a-2)\in R.
$$
The second numerator is polynomial and has value
$2\delta(a-1)$ at both $z=0$ and $z=1$, proving the stated membership.
Thus $w=(B+(a-1)\delta)/u$ belongs to $R_u$, and
$w=(zv+\delta(z^2+z+2a-2))/v$ belongs to $R_v$.
Since $R_a$ is generated by $A$, $B+A$ and $w$, inclusion induces
isomorphisms
$R_u\to(R_a)_u$ and $R_v\to(R_a)_v$.
Their inverse maps send the original $w$ to the respective fractions
just displayed, send each element of $R$ to itself, and send
$u^{-1}$ or $v^{-1}$ to that same inverse in the rational-function field.
These formulas prove both composites are identities; on the overlap
they agree as maps into $\mathbf Q(z)$.
The full unit-ideal identity is
$$
u(w-a+2)-v=\delta(2a-1).
$$
Its right hand side is a nonzero rational constant under the present
hypotheses. Division by this displayed constant gives
$1=((w-a+2)/(\delta(2a-1)))u-(1/(\delta(2a-1)))v$ in $R_a$.
Hence $D(u)$ and $D(v)$ cover $\Spec(R_a)$.
Their images are respectively $D_R(u)$ and $D_R(v)$, and those images
have union $\Spec(R)\setminus V(u,v)=\Spec(R)\setminus\{\mathfrak m_a\}$.
Contraction of prime ideals gives
$\theta^{-1}(D_R(u))=D_{R_a}(u)$ and the analogous equality for $v$.
If two source points have the same image, choose one of these target
opens containing it; both points belong to its inverse image and
are equal by that chart isomorphism. Thus $\theta$ is injective.
The two chart maps also prove surjectivity onto the stated union and
openness there, establishing the asserted homeomorphism.

\medskip\noindent
For the zero-set descriptions, the presentation of $R$ shows that the
only prime containing $A$ is $\mathfrak m_0$: modulo $A$ the relation
gives $B^2=0$. The values of $u$ and $v$ at this point are respectively
$-\delta$ and $-2\delta$, both nonzero.
Moreover $R_A=\mathbf Q[z,1/A(z)]$ inside $\mathbf Q(z)$, since
$z=B/A$ and the converse inclusion preserves $A(z)^{-1}$.
Thus all zeros under consideration can be read in this polynomial
localization. The two roots of $u$ are $a$ and $1-a$, distinct and
different from $0,1$. The other factors of $v$ come from
$q(z)=z^2+z+2a-2$, with $q(0)=2a-2\ne0$ and $q(1)=2a\ne0$.
Also $q(a)=a^2+3a-2\ne0$ for rational $a$: equality would give
$(2a+3)^2=17$, whereas $17$ is not a square of a rational number.
Consequently the irreducible factors of $q$ give one or two distinct
closed points in addition to $\mathfrak m_a$ in the target spectrum.
Finally $q(1-a)=a^2-a=\delta\ne0$, so none is $\mathfrak m_{1-a}$.
Deleting $\mathfrak m_a$ through the proved chart correspondence
gives exactly the claimed source zero set.

\medskip\noindent
Despite this homeomorphism which mimics the behavior of a
localization at an element of $R$, while
$\mathbf{Q}[z, \frac{1}{z-a}]$ is the localization of $\mathbf{Q}[z]$
at the element $z-a$, the ring $R_a$ is {\it not} a
localization of $R$: Any proper localization $S^{-1}R$ results in more
units than the original ring $R$.  The units of $R$ are
$\mathbf{Q}^\times$, the units of $\mathbf{Q}$.  In fact, it is
easy to see that the units of $R_a$ are $\mathbf{Q}^*$.
Namely, the units of $\mathbf{Q}[z, \frac{1}{z - a}]$ are
$c (z - a)^n$ for $c \in \mathbf{Q}^*$ and $n \in \mathbf{Z}$
and such a unit belongs to $R_a$ if and only if
$c(-a)^n=c(1-a)^n$. Its inverse then has equal values too, so
this is exactly the condition for it to be a unit of $R_a$.
For $n\ne0$ the equality implies
$\bigl((-a)/(1-a)\bigr)^n=1$.
A rational number with a nonzero integer power equal to $1$ is $1$ or $-1$:
writing it as a quotient of coprime integers shows that their absolute
values must agree. The value $1$ would give $-a=1-a$, which is impossible;
the value $-1$ gives $a=1/2$, which is excluded here.
Thus $n=0$ and $R_a^\times=\mathbf Q^\times=R^\times$.
The ambient description of units follows as well from polynomial
factorization: if $P(z)/(z-a)^r$ has inverse $Q(z)/(z-a)^s$, then
$P(z)Q(z)=(z-a)^{r+s}$, so $P(z)=c(z-a)^j$ and the original unit
is $c(z-a)^{j-r}$, with the coefficient and exponent retained.
Moreover $R_a\ne R$: the element
$w=z+(a^2-a)/(z-a)$ has a pole at $z=a$, since its numerator
$z(z-a)+(a^2-a)$ has nonzero value $a^2-a$ there, whereas all
elements of $R$ are polynomials.
If the inclusion $R\to R_a$ were a localization at a multiplicative
subset $S$, every $s\in S$ would become a unit of $R_a$ and hence
would already be a nonzero rational constant, a unit of $R$.
Then $S^{-1}R=R$ under the original map, contradicting $w\notin R$.
Therefore there is no $R$-algebra isomorphism $S^{-1}R\cong R_a$
compatible with the specified inclusion.

\medskip\noindent
We used the fact that $a\neq 0, 1$ to ensure that
$\frac{1}{z-a}$ makes sense at $z = 0, 1$.  We used the fact that
$a\neq 1/2$ in a few places: (1) In order to be able to talk about
the kernel of $\text{ev}_{1-a}$ on $R_a$, which ensures that
$\mathfrak{m}_{1-a}$ is a point of $R_a$ (i.e., that $R_a$ is
missing just one point of $R$). (2) At the end in order to conclude
that $(z-a)^n$ can belong to $R_a$ only for $n = 0$; indeed, if
$a = 1/2$, then it belongs to $R_a$ whenever $n$ is even. Hence
there are indeed more units in the ring defined by the same formula
at $a=1/2$. In this exceptional case we can decide the localization
question exactly. Set $h=(z-1/2)^2$. Its values at $0$ and $1$ are
both $1/4$, so $h\in R$ and $h^{-1}\in R_{1/2}$.
The original inclusion induces
$$
R_h\longrightarrow
R_{1/2}=\{f\in\mathbf Q[z,1/(z-1/2)]\mid f(0)=f(1)\}.
$$
It is injective because both rings are subrings of $\mathbf Q(z)$.
For surjectivity, write any $f\in R_{1/2}$ as
$P(z)/(z-1/2)^N$ and choose $m\geq0$ with $2m\geq N$.
Then $h^mf=P(z)(z-1/2)^{2m-N}$ is a polynomial, and its values at
$0$ and $1$ are equal because $h(0)=h(1)=1/4$ and $f(0)=f(1)$.
Thus $h^mf\in R$ and $f=(h^mf)/h^m\in R_h$.
This proves $R_{1/2}=R_{(z-1/2)^2}$ via the specified map.
The same endpoint calculation gives its exact unit group
$R_{1/2}^{\times}=\{c(z-1/2)^{2j}\mid c\in\mathbf Q^{\times},\ j\in\mathbf Z\}$,
including negative exponents. Its spectrum is the standard open
$D_R((z-1/2)^2)=\Spec(R)\setminus\{\mathfrak m_{1/2}\}$:
its only zero lies in $D_R(A)$ and corresponds there to the factor
$z-1/2$, while its value at $\mathfrak m_0$ is $1/4\ne0$.
\end{example}


\subsection{Editorial treatment of Algebra lines 5182--5190}

\hypertarget{algebra-context-028}{}

\noindent\href{https://github.com/stacks/stacks-project/blob/a04446e57ec1fbc252a871afcec7752fb2807b14/algebra.tex\#L5182}{Original source, lines 5182--5190}. Complete original and reviewed TeX passages accompany this note. Every intervention is separately recorded; the following treatment is editorial.


\begin{example}
\label{editorial-source-example-oka-family-bound-cardinality}
Let $R$ be a ring.
Let $\kappa$ be an infinite cardinal.
The family of ideals which can be generated by at most $\kappa$ elements
is an Oka family. To verify this, suppose $(I:a)$ has a
generating family $(u_\lambda)_{\lambda\in L}$ and $(I,a)$ has
a generating family $(v_\mu)_{\mu\in P}$, with $|L|,|P|\leq\kappa$.
Write $v_\mu=b_\mu+c_\mu a$, with $b_\mu\in I$ and $c_\mu\in R$.
Then $(I,a)=(a,b_\mu\mid\mu\in P)$.
For $x\in I$, an expression
$x=ra+\sum_{\mu\in P_0}s_\mu b_\mu$ with $P_0$ finite gives
$ra\in I$, hence $r\in(I:a)$. Writing
$r=\sum_{\lambda\in L_0}t_\lambda u_\lambda$ with $L_0$ finite yields
$$
x=\sum_{\lambda\in L_0}t_\lambda a u_\lambda+
\sum_{\mu\in P_0}s_\mu b_\mu.
$$
Thus $I=(a u_\lambda,b_\mu\mid\lambda\in L,\mu\in P)$, with
at most $\kappa+\kappa=\kappa$ generators. The ideal $R$ is
generated by $1$, so it belongs to this family as well.
\end{example}


\subsection{Editorial treatment of Algebra lines 5237--5262}

\hypertarget{algebra-context-029}{}

\noindent\href{https://github.com/stacks/stacks-project/blob/a04446e57ec1fbc252a871afcec7752fb2807b14/algebra.tex\#L5237}{Original source, lines 5237--5262}. Complete original and reviewed TeX passages accompany this note. Every intervention is separately recorded; the following treatment is editorial.


\begin{lemma}
\label{editorial-source-lemma-cohen}
Let $R$ be a ring.
\begin{enumerate}
\item An ideal $I \subset R$ maximal with respect to not being
finitely generated is prime.
\item If every prime ideal of $R$ is
finitely generated, then
every ideal of $R$ is finitely generated\footnote{Later we will say
that $R$ is Noetherian.}.
\end{enumerate}
\end{lemma}

\begin{proof}
The first assertion is an immediate consequence of
Example \ref{algebra-example-oka-family-finitely-generated} and
Proposition \ref{algebra-proposition-oka}. For the second,
suppose that there exists an ideal $I \subset R$ which is not finitely
generated. In the nonempty set of such ideals, the empty chain
has the chosen $I$ as an upper bound. The union of a nonempty totally ordered chain $\left\{I_\alpha\right\}$
of ideals that are not finitely generated is not finitely generated;
indeed, if $I = \bigcup I_\alpha$ were generated by
$a_1, \ldots, a_n$, then all the generators would belong to some
$I_\alpha $ and would consequently generate it.
By Zorn's lemma, there is an ideal maximal with respect to being not finitely
generated. By the first part this ideal is prime.
\end{proof}


\subsection{Editorial treatment of Algebra lines 5264--5286}

\hypertarget{algebra-context-030}{}

\noindent\href{https://github.com/stacks/stacks-project/blob/a04446e57ec1fbc252a871afcec7752fb2807b14/algebra.tex\#L5264}{Original source, lines 5264--5286}. Complete original and reviewed TeX passages accompany this note. Every intervention is separately recorded; the following treatment is editorial.


\begin{lemma}
\label{editorial-source-lemma-primes-principal}
Let $R$ be a ring.
\begin{enumerate}
\item An ideal $I \subset R$ maximal with respect to not being
principal is prime.
\item If every prime ideal of $R$ is principal, then
every ideal of $R$ is principal.
\end{enumerate}
\end{lemma}

\begin{proof}
The first part follows from
Example \ref{algebra-example-oka-family-principal} and
Proposition \ref{algebra-proposition-oka}.
For the second, suppose that there exists an ideal $I \subset R$
which is not principal. In the nonempty set of nonprincipal ideals,
the empty chain has the chosen $I$ as an upper bound.
The union of a nonempty totally ordered chain
$\left\{I_\alpha\right\}$ of ideals that are not principal is not principal;
indeed, if $I = \bigcup I_\alpha$ were generated by
$a$, then $a$ would belong to some $I_\alpha $ and $a$ would generate it.
By Zorn's lemma, there is an ideal maximal with respect to not being
principal. This ideal is necessarily prime by the first part.
\end{proof}


\subsection{Editorial treatment of Algebra lines 5288--5306}

\hypertarget{algebra-context-031}{}

\noindent\href{https://github.com/stacks/stacks-project/blob/a04446e57ec1fbc252a871afcec7752fb2807b14/algebra.tex\#L5288}{Original source, lines 5288--5306}. Complete original and reviewed TeX passages accompany this note. Every intervention is separately recorded; the following treatment is editorial.


\begin{lemma}
\label{editorial-source-lemma-characterize-domain}
Let $R$ be a ring.
\begin{enumerate}
\item An ideal maximal among the ideals which do not contain a
nonzerodivisor is prime.
\item If $R$ is nonzero and every nonzero prime ideal in $R$
contains a nonzerodivisor, then $R$ is a domain.
\end{enumerate}
\end{lemma}

\begin{proof}
Consider the set $S$ of nonzerodivisors. It is a multiplicative
subset of $R$. Hence any ideal maximal with respect to not intersecting
$S$ is prime, see
Lemma \ref{algebra-lemma-simple}.
For the second assertion, assume $R\ne0$. Then $0\notin S$,
so $(0)$ is an ideal disjoint from $S$. In the set of ideals disjoint
from $S$, every nonempty totally ordered chain has its union as an
upper bound: an element of that union belongs to some member, so
the union still misses $S$. The empty chain is bounded by $(0)$.
Zorn's lemma supplies a maximal member $\mathfrak p$, and the first
assertion shows it is prime. By hypothesis every nonzero prime
meets $S$, so this $\mathfrak p$ must equal $(0)$.
Thus $(0)$ is prime, and the nonzero ring $R$ is a domain.
\end{proof}


\subsection{Editorial treatment of Algebra lines 5308--5327}

\hypertarget{algebra-context-032}{}

\noindent\href{https://github.com/stacks/stacks-project/blob/a04446e57ec1fbc252a871afcec7752fb2807b14/algebra.tex\#L5308}{Original source, lines 5308--5327}. Complete original and reviewed TeX passages accompany this note. Every intervention is separately recorded; the following treatment is editorial.


\begin{remark}
\label{editorial-source-remark-cohen-bound-cardinality}
Let $R$ be a ring. Let $\kappa$ be an infinite cardinal.
By applying
Example \ref{editorial-source-example-oka-family-bound-cardinality} and
Proposition \ref{algebra-proposition-oka}
we see that any ideal maximal with respect to the property of not being
generated by $\kappa$ elements is prime. This result is not so
useful because there exists a ring for which every prime ideal
of $R$ can be generated by $\aleph_0$ elements, but some
ideal cannot. Namely, let $k$ be a field, let $T$ be a set whose
cardinality is greater than $\aleph_0$ and let
$$
R = k[\{x_n\}_{n \geq 1}, \{z_{t, n}\}_{t \in T, n \geq 0}]/
\bigl(x_n^2\ (n\geq1),\ z_{t,n}^2\ (t\in T,n\geq0),\
x_nz_{t,n}-z_{t,n-1}\ (t\in T,n\geq1)\bigr)
$$
This is a local ring with unique prime ideal
$\mathfrak m=(x_n\mid n\geq1)$. Indeed, every prime contains all
$x_n$ and $z_{t,n}$ because their squares vanish, and the quotient
by all these elements is the field $k$. The relations
$z_{t,n}=x_{n+1}z_{t,n+1}$ put every $z_{t,n}$ in $(x_n)$.
Every element of this ideal is nilpotent: it uses only finitely many
of the square-zero generators, and a product of more factors than
the number of those generators has a repeated factor and vanishes.
The augmentation to $k$ shows the ring is nonzero.
Thus $\mathfrak m$ is the unique prime and is countably generated.

\medskip\noindent
The ideal $J=(z_{t,n}\mid t\in T,n\geq0)$ in fact requires exactly
$|T|$ generators. To prove the lower bound we construct maps detecting
each original $T$-label. For $N\geq0$ let
$$
C_N=k[X_1,\ldots,X_N,W_N]/(X_1^2,\ldots,X_N^2,W_N^2).
$$
Define $C_N\to C_{N+1}$ by $X_i\mapsto X_i$ for $i\leq N$
and $W_N\mapsto X_{N+1}W_{N+1}$.
The square-free monomials
$X_E=\prod_{i\in E}X_i$ and $X_EW_N$, for $E\subset\{1,\ldots,N\}$,
form a $k$-basis of $C_N$.
Their images are respectively $X_E$ and $X_EX_{N+1}W_{N+1}$,
distinct members of the corresponding basis of $C_{N+1}$.
Thus all these maps are injective. Put $C=\colim_N C_N$.
Concretely, take classes of pairs $(N,c)$ with $c\in C_N$, identifying
two pairs when their images agree at a common later stage. Addition
and multiplication are computed at a common stage; compatibility of
the ring maps makes them independent of the stage and representatives.
The injectivity of every transition map makes each $C_N\to C$ injective.
In $C$ the original classes satisfy
$W_{n-1}=X_nW_n$ for $n\geq1$, and
$W_0=X_1\cdots X_NW_N\ne0$ at every stage and in the colimit.

\medskip\noindent
For each $t_0\in T$, the assignments
$$
\rho_{t_0}:R\longrightarrow C,\qquad
x_n\longmapsto X_n,\qquad
z_{t,n}\longmapsto
\begin{cases}W_n,&t=t_0,\\0,&t\ne t_0\end{cases}
$$
respect every displayed defining relation, including its index range,
and hence define a ring map. In particular
$\rho_{t_0}(z_{t_0,0})=W_0\ne0$.
Suppose a family $(j_\lambda)_{\lambda\in L}$ with $|L|<|T|$
generated $J$. Each $j_\lambda$ is a finite sum of multiples
of generators $z_{t,n}$, so the labels occurring in all these
finite expressions form a set $T'\subset T$ of cardinality at most
$\max(|L|,\aleph_0)<|T|$; here $T$ is uncountable.
Choose $t_0\in T\setminus T'$. Then $\rho_{t_0}$ kills every
$j_\lambda$, irrespective of its coefficients, but does not kill
$z_{t_0,0}\in J$, a contradiction.
This proves the lower bound. The displayed generating set for $J$
has cardinality $|T\times\mathbf Z_{\geq0}|=|T|$, proving equality.
In particular $J$ is not countably generated.
More generally, for any prescribed infinite cardinal $\kappa$, choosing
$|T|>\kappa$ gives a ring whose only prime ideal is countably
generated, hence generated by at most $\kappa$ elements, while $J$
cannot be generated by $\kappa$ elements.
\end{remark}


\subsection{Editorial treatment of Algebra lines 5329--5362}

\hypertarget{algebra-context-033}{}

\noindent\href{https://github.com/stacks/stacks-project/blob/a04446e57ec1fbc252a871afcec7752fb2807b14/algebra.tex\#L5329}{Original source, lines 5329--5362}. Complete original and reviewed TeX passages accompany this note. Every intervention is separately recorded; the following treatment is editorial.


\begin{example}
\label{editorial-source-example-noetherian-topology-on-spec}
\begin{reference}
Comment by Lukas Heger of November 12, 2020.
\end{reference}
Let $R$ be a ring and $X = \Spec(R)$. Since closed subsets of $X$ correspond to
radical ideals of $R$ (Lemma \ref{editorial-source-lemma-Zariski-topology}) we see that $X$ is a
Noetherian topological space if and only if we have ACC for radical ideals.
This holds if and only if every radical ideal is the radical of a finitely
generated ideal. For completeness, under ACC the set of radicals
of finitely generated subideals of a radical ideal $K$ has a maximal
member $\sqrt{F}$. Otherwise successive strict enlargements would
give an infinite ascending chain. For every $x\in K$ the ideal
$\sqrt{F+(x)}$ is another such member containing $\sqrt F$,
so maximality gives $x\in\sqrt F$ and hence $K=\sqrt F$.
Conversely, the union $K$ of an ascending sequence of radical ideals
$K_1\subset K_2\subset\cdots$ is radical: if a power of an element
lies in the union, it lies in one $K_i$, which then contains the element.
Write $K=\sqrt{(f_1,\ldots,f_n)}$. Each $f_j$ belongs to some
$K_i$, and a single $K_N$ contains all of them. Therefore
$K=\sqrt{(f_1,\ldots,f_n)}\subset K_N\subset K$, so the sequence
stabilizes. This proves ACC. Let
$$
\mathcal{F} = \{I \subset R \mid
\sqrt{I} = \sqrt{(f_1, \ldots, f_n)}\text{ for some }n
\text{ and }f_1, \ldots, f_n \in R\}.
$$
We now prove that $\mathcal F$ is an Oka family. First, $R\in\mathcal F$
using the generator $1$. For any $I$ and $a$, the product
$(I,a)(I:a)$ is contained in $I$. Conversely, for $x\in I$ we have
$x\in(I,a)$ and $x\in(I:a)$, since $xa\in I$; hence
$x^2\in(I,a)(I:a)$. These two containments prove the identity
$$
\sqrt{I} = \sqrt{(I, a)(I : a)}
$$
for every ideal $I\subset R$ and element $a\in R$.
If $\sqrt{(I,a)}=\sqrt{(f_1,\ldots,f_r)}$ and
$\sqrt{(I:a)}=\sqrt{(g_1,\ldots,g_s)}$, it follows that
$$
\sqrt I=\sqrt{(f_i g_j\mid 1\leq i\leq r,\ 1\leq j\leq s)}.
$$
Indeed, a prime contains a product of two ideals if and only if it
contains one of them, and replacing each ideal by another with the
same radical preserves precisely those primes. The radical is the
intersection of its containing primes by Lemma \ref{editorial-source-lemma-Zariski-topology}.
The displayed finite set of products therefore proves the Oka condition.
If the complement of $\mathcal F$ is nonempty, choose an ideal
in it as an upper bound for the empty chain. On the other hand,
if we have a nonempty totally ordered chain of ideals
$\{I_\alpha\}$ none of which are in $\mathcal{F}$, then the union
$I = \bigcup I_\alpha$ cannot be in $\mathcal{F}$ either. Otherwise
$\sqrt{I} = \sqrt{(f_1, \ldots, f_n)}$, then $f_i^e \in I$ for some $e$,
then $f_i^e \in I_\alpha$ for some $\alpha$ independent of $i$, then
$\sqrt{I_\alpha} = \sqrt{(f_1, \ldots, f_n)}$, contradiction.
Thus if the set of ideals not in $\mathcal{F}$ is nonempty, then it
has maximal elements and
exactly as in Lemma \ref{editorial-source-lemma-cohen} we conclude that $X$ is a
Noetherian topological space if and only if every prime ideal of $R$
is equal to $\sqrt{(f_1, \ldots, f_n)}$ for some $f_1, \ldots, f_n \in R$.
In detail, if every prime belongs to $\mathcal F$ but its complement
were nonempty, the chain argument and Zorn's lemma would give a maximal
excluded ideal, which is prime by Proposition \ref{algebra-proposition-oka},
a contradiction. Thus every ideal belongs to $\mathcal F$, giving
ACC for radical ideals by the equivalence just proved.
Conversely, ACC gives the finite radical presentation of every radical
ideal, in particular every prime. This proves the stated criterion
for the spectrum to be Noetherian.
\end{example}


\subsection{Editorial treatment of Algebra lines 5433--5451}

\hypertarget{algebra-context-034}{}

\noindent\href{https://github.com/stacks/stacks-project/blob/a04446e57ec1fbc252a871afcec7752fb2807b14/algebra.tex\#L5433}{Original source, lines 5433--5451}. Complete original and reviewed TeX passages accompany this note. Every intervention is separately recorded; the following treatment is editorial.


\begin{lemma}
\label{editorial-source-lemma-constructible}
Let $R$ be a ring. A subset of $\Spec(R)$ is constructible if and only
if it can be written as a finite union of subsets of the form
$D(f) \cap V(g_1, \ldots, g_m)$ for $f, g_1, \ldots, g_m \in R$.
\end{lemma}

\begin{proof}
By Lemma \ref{algebra-lemma-qc-open} the subset $D(f)$ and the complement of
$V(g_1, \ldots, g_m)$ are retro-compact open. Hence
$D(f) \cap V(g_1, \ldots, g_m)$ is a constructible subset and so is
any finite union of such. Conversely, by Topology, Definition
\ref{topology-definition-constructible}, a constructible subset has an
expression $T=\bigcup_{\alpha=1}^k(U_\alpha\cap V_\alpha^c)$ with
$U_\alpha,V_\alpha$ retrocompact open. By Lemma \ref{algebra-lemma-qc-open}, write
$U_\alpha=\bigcup_{i=1}^{n_\alpha}D(f_{\alpha i})$ and
$V_\alpha=\bigcup_{j=1}^{m_\alpha}D(g_{\alpha j})$. Then
$$
T=\bigcup_{\alpha=1}^k\bigcup_{i=1}^{n_\alpha}
\left(D(f_{\alpha i})\cap
V(g_{\alpha1},\ldots,g_{\alpha m_\alpha})\right).
$$
These are finite unions; an empty inner union is empty, and
$V()=\Spec(R)$ when $m_\alpha=0$.
\end{proof}


\subsection{Editorial treatment of Algebra lines 5453--5472}

\hypertarget{algebra-context-035}{}

\noindent\href{https://github.com/stacks/stacks-project/blob/a04446e57ec1fbc252a871afcec7752fb2807b14/algebra.tex\#L5453}{Original source, lines 5453--5472}. Complete original and reviewed TeX passages accompany this note. Every intervention is separately recorded; the following treatment is editorial.


\begin{lemma}
\label{editorial-source-lemma-constructible-is-image}
Let $R$ be a ring and let $T \subset \Spec(R)$
be constructible. Then there exists a ring map $R \to S$ of
finite presentation such that $T$ is the image of
$\Spec(S)$ in $\Spec(R)$.
\end{lemma}

\begin{proof}
Write
$T=\bigcup_{i=1}^k(D(f_i)\cap V(g_{i1},\ldots,g_{im_i}))$
as in Lemma \ref{editorial-source-lemma-constructible}, and set
$S_i=(R/(g_{i1},\ldots,g_{im_i}))_{f_i}$.
The canonical maps on spectra identify $\Spec(S_i)$ with the indicated
subset, by Lemmas \ref{algebra-lemma-standard-open} and \ref{algebra-lemma-spec-closed}.
For $k=0$, take $S=R/(1)$, whose spectrum is empty and which is finitely
presented over $R$. For $k>0$, take $S=\prod_{i=1}^kS_i$, with its
diagonal $R$-algebra structure. Lemma \ref{algebra-lemma-spec-product} identifies
its spectrum with the disjoint union of the component spectra; the
contraction map to $\Spec(R)$ therefore has image exactly $T$.

\medskip\noindent
Here is a finite presentation, including all component relations. Let
$C=R[E_1,\ldots,E_k,Y_1,\ldots,Y_k]/J$, where $J$ is generated by
$$
\begin{gathered}
E_i^2-E_i\quad(1\leq i\leq k),\qquad
E_iE_j\quad(1\leq i<j\leq k),\qquad
\sum_{i=1}^kE_i-1,\\
(1-E_i)Y_i,\quad f_iY_i-E_i\quad(1\leq i\leq k),\qquad
g_{ij}E_i\quad(1\leq i\leq k,\ 1\leq j\leq m_i).
\end{gathered}
$$
Map $C$ to $S$ by sending $E_i$ to the component idempotent and $Y_i$
to the element whose $i$th component is $1/f_i$ and whose other
components are zero. Every displayed relation vanishes. In $C$, the
orthogonal idempotents $E_i$ sum to $1$, so
$C\longrightarrow\prod_iE_iC$, $c\longmapsto(E_ic)_i$, has inverse
$(c_i)_i\longmapsto\sum_i c_i$. The ring $E_iC$, with identity $E_i$,
has the $R$-algebra map $r\mapsto rE_i$; all $g_{ij}$ vanish and
$f_iE_i$ has inverse $Y_i$. This defines a map $S_i\to E_iC$ sending
$\bar r/f_i^a$ to $rY_i^a$ for $a>0$, and $\bar r$ to $rE_i$.
The localization relations make the map well defined. The map in the
other direction sends $E_i$ to $1$ and $Y_i$ to $1/f_i$.
The composites fix the coefficient images and these generators; all
other $E_j,Y_j$ vanish in $E_iC$. Thus these maps are inverse, and
$C\cong S$ as $R$-algebras. This proves finite presentation without
any additional finiteness hypothesis on $R$.
\end{proof}


\subsection{Editorial treatment of Algebra lines 5554--5595}

\hypertarget{algebra-context-036}{}

\noindent\href{https://github.com/stacks/stacks-project/blob/a04446e57ec1fbc252a871afcec7752fb2807b14/algebra.tex\#L5554}{Original source, lines 5554--5595}. Complete original and reviewed TeX passages accompany this note. Every intervention is separately recorded; the following treatment is editorial.


\begin{lemma}
\label{editorial-source-lemma-characteristic-polynomial-prime}
Let $R \to A$ be a ring homomorphism.
Assume $A \cong R^{\oplus n}$ as an $R$-module.
Let $f \in A$. The multiplication map $m_f: A
\to A$ is $R$-linear and hence
has a characteristic polynomial
$P(T) = T^n + r_{n-1}T^{n-1} + \ldots + r_0 \in R[T]$.
For any prime
$\mathfrak{p} \in \Spec(R)$, $f$ acts nilpotently on $A
\otimes_R \kappa(\mathfrak{p})$ if and only if $\mathfrak p \in
V(r_0, \ldots, r_{n-1})$.
Moreover, the image of $\Spec(A_f)$ in $\Spec(R)$ is
$\bigcup_{i=0}^{n-1}D(r_i)$. This formula commutes with every ring
base change $R\to R'$, and $\Spec(A)\to\Spec(R)$ is open after every
such base change.
\end{lemma}

\begin{proof}
This follows quite easily once we prove that the characteristic
polynomial $\bar P(T) \in \kappa(\mathfrak p)[T]$ of the
multiplication map $m_{\bar f}: A \otimes_R \kappa(\mathfrak p) \to
A \otimes_R \kappa(\mathfrak p)$ which multiplies elements of $A
\otimes_R \kappa(\mathfrak p)$ by $\bar f$, the image of $f$ viewed in
$A\otimes_R\kappa(\mathfrak p)$, is just the image of $P(T)$ in
$\kappa(\mathfrak p)[T]$. Let $(a_{ij})$ be the matrix of the map
$m_f$ with entries in $R$, using a basis $e_1, \ldots, e_n$
of $A$ as an $R$-module.
Then, $A \otimes_R \kappa(\mathfrak p) \cong (R \otimes_R
\kappa(\mathfrak p))^{\oplus n} = \kappa(\mathfrak p)^n$, which is
an $n$-dimensional vector space over $\kappa(\mathfrak p)$ with
basis $e_1 \otimes 1, \ldots, e_n \otimes 1$. The image $\bar f = f
\otimes 1$, and so the multiplication map $m_{\bar f}$ has matrix
$(a_{ij} \otimes 1)$. Thus, the characteristic polynomial is
precisely the image of $P(T)$.

\medskip\noindent
For a linear endomorphism $N$ of a finite dimensional vector space,
nilpotence implies that the filtration
$0=\ker(N^0)\subset\ker(N)\subset\cdots\subset\ker(N^a)$
eventually reaches the whole space. A basis extending bases along this
filtration gives a strictly triangular matrix, since
$N(\ker(N^j))\subset\ker(N^{j-1})$; its characteristic polynomial
is $T^n$. Conversely, if the characteristic polynomial is $T^n$,
Cayley--Hamilton gives $N^n=0$ for $n>0$.
For $n=0$ the space is zero, its only endomorphism is zero and its
characteristic polynomial is $1=T^0$. Applying this to $m_{\bar f}$
gives nilpotence exactly when all the images of
$r_0,\ldots,r_{n-1}$ in $\kappa(\mathfrak p)$ vanish.
The kernel of $R\to\kappa(\mathfrak p)$ is $\mathfrak p$, proving
the first assertion, including $V()=\Spec(R)$ when $n=0$.

\medskip\noindent
Put $B=A\otimes_R\kappa(\mathfrak p)$. Multiplication by $\bar f$
is nilpotent if and only if $\bar f$ is nilpotent: one implication
follows by applying the power of the multiplication map to $1$, and
the other follows by multiplication by $\bar f^a=0$.
The localization $B_{\bar f}$ is zero if and only if
$\bar f^a=0$ for some $a\geq0$, by the equality criterion for
$1/1=0/1$. Also
$A_f\otimes_R\kappa(\mathfrak p)\cong B_{\bar f}$:
the map sends $(a/f^j)\otimes\lambda$ to
$(a\otimes\lambda)/\bar f^j$, and its inverse is induced by
$a\otimes\lambda\mapsto(a/1)\otimes\lambda$ and the inverse
$(1/f)\otimes1$ of $\bar f$. These maps fix the generators and
the inverse of $f$, so their composites are identities.
Lemma \ref{editorial-source-lemma-in-image} now gives
$$
\mathfrak p\in\operatorname{Im}(\Spec(A_f)\to\Spec(R))
\ \Longleftrightarrow\ B_{\bar f}\ne0
\ \Longleftrightarrow\ \mathfrak p\in\bigcup_{i=0}^{n-1}D(r_i).
$$
For $n=0$, $A=0$ and the union and image are both empty.

\medskip\noindent
For an arbitrary ring map $R\to R'$, the images of the chosen basis
give a basis of $A'=A\otimes_RR'$. Its multiplication matrix for
$f'=f\otimes1$ is the image $(a'_{ij})$ of $(a_{ij})$. The determinant
formula
$$
\det(TI-(a_{ij}))=
\sum_{\sigma\in S_n}\operatorname{sgn}(\sigma)
\prod_{i=1}^n(T\delta_{i,\sigma(i)}-a_{i,\sigma(i)})
$$
shows that every coefficient of its characteristic polynomial is the
image of the corresponding original coefficient, with all signs and
terms retained. Consequently the image of $\Spec(A'_{f'})$ is
$\bigcup_{i=0}^{n-1}D(r'_i)$, the inverse image of
$\bigcup_{i=0}^{n-1}D(r_i)$ in $\Spec(R')$.
Finally, the same image formula applies to every element of $A'$,
whether or not it comes from an element of $A$. Standard opens form
a basis of $\Spec(A')$, and arbitrary unions of their open images
are open. Hence $\Spec(A')\to\Spec(R')$ is open.
\end{proof}


\subsection{Editorial treatment of Algebra lines 5597--5643}

\hypertarget{algebra-context-037}{}

\noindent\href{https://github.com/stacks/stacks-project/blob/a04446e57ec1fbc252a871afcec7752fb2807b14/algebra.tex\#L5597}{Original source, lines 5597--5643}. Complete original and reviewed TeX passages accompany this note. Every intervention is separately recorded; the following treatment is editorial.


\begin{lemma}
\label{editorial-source-lemma-affineline-special}
Let $R$ be a ring. Let $f, g \in R[x]$ be polynomials.
Assume the leading coefficient of $g$ is a unit of $R$.
There exist elements $r_i\in R$, $i = 1\ldots, n$ such that
the image of $D(f) \cap V(g)$ in $\Spec(R)$ is
$\bigcup_{i = 1, \ldots, n} D(r_i)$.
\end{lemma}

\begin{proof}
Write $g = ux^d + a_{d-1}x^{d-1} + \ldots + a_0$, where
$d$ is the degree of $g$, and hence $u \in R^*$.
If $d=0$, then $g=u$ is a unit, so $V(g)$ and its image are empty;
the required union may be taken to have no terms. Suppose $d>0$ and
consider the original ring $A=R[x]/(g)$. It is finite free with basis
the images of $1,x,\ldots,x^{d-1}$. Indeed, if a polynomial has leading
term $bx^e$ with $e\geq d$, subtracting
$bu^{-1}x^{e-d}g$ cancels that term and lowers its degree. Iteration
gives a remainder of degree less than $d$, keeping $g$ and all its
coefficients unchanged. If a nonzero multiple $qg$ had degree less
than $d$, its leading coefficient would be the product of the nonzero
leading coefficient of $q$ with the unit $u$, hence nonzero, and its
degree would be $\deg(q)+d\geq d$, a contradiction. This proves both
spanning and independence. Consider multiplication by the image of
$f$ on $A$. This is an $R$-module map.
Hence we can let $P(T) \in R[T]$ be the characteristic polynomial
of this map. Write $P(T) = T^d + r_{d-1} T^{d-1} + \ldots + r_0$.
We claim that $r_0, \ldots, r_{d-1}$ have the desired property.
We will use below the property of characteristic polynomials
that
$$
\mathfrak p \in V(r_0, \ldots, r_{d-1})
\Leftrightarrow
\text{multiplication by }f\text{ is nilpotent on }
A \otimes_R \kappa(\mathfrak p).
$$
This was proved in Lemma \ref{editorial-source-lemma-characteristic-polynomial-prime}.

\medskip\noindent
Suppose $\mathfrak q\in D(f) \cap V(g)$, and let
$\mathfrak p = \mathfrak q \cap R$. Then there is a nonzero map
$A \otimes_R \kappa(\mathfrak p) \to \kappa(\mathfrak q)$ which
is compatible with multiplication by $f$.
And $f$ acts as a unit on $\kappa(\mathfrak q)$.
Thus we conclude $\mathfrak p \not \in  V(r_0, \ldots, r_{d-1})$.

\medskip\noindent
On the other hand, suppose that $r_i \not\in \mathfrak p$ for some
prime $\mathfrak p$ of $R$ and some $0 \leq i \leq d - 1$.
Then multiplication by $f$ is not nilpotent on the algebra
$A \otimes_R \kappa(\mathfrak p)$.
Hence there exists a prime ideal $\overline{\mathfrak q} \subset
A \otimes_R \kappa(\mathfrak p)$ not containing the image of $f$.
The inverse image of $\overline{\mathfrak q}$ in $R[x]$
is an element of $D(f) \cap V(g)$ mapping to $\mathfrak p$.
\end{proof}


\subsection{Editorial treatment of Algebra lines 5645--5709}

\hypertarget{algebra-context-038}{}

\noindent\href{https://github.com/stacks/stacks-project/blob/a04446e57ec1fbc252a871afcec7752fb2807b14/algebra.tex\#L5645}{Original source, lines 5645--5709}. Complete original and reviewed TeX passages accompany this note. Every intervention is separately recorded; the following treatment is editorial.


\begin{theorem}[Chevalley's Theorem]
\label{editorial-source-theorem-chevalley}
Suppose that $R \to S$ is of finite presentation.
The image of a constructible subset of
$\Spec(S)$ in $\Spec(R)$ is constructible.
\end{theorem}

\begin{proof}
Write $S = R[x_1, \ldots, x_n]/(f_1, \ldots, f_m)$.
We may factor $R \to S$ as $R \to R[x_1] \to R[x_1, x_2]
\to \ldots \to R[x_1, \ldots, x_{n-1}] \to S$. Hence
we may assume that $S = R[x]/(f_1, \ldots, f_m)$.
In this case we factor the map as $R \to R[x] \to S$,
and by Lemma \ref{algebra-lemma-closed-fp} we reduce to
the case $S = R[x]$. By Lemma \ref{editorial-source-lemma-constructible} it suffices
to show that if
$T = (\bigcup_{i = 1\ldots n} D(f_i)) \cap V(g_1, \ldots, g_m)$
for $f_i , g_j \in R[x]$ then the image in $\Spec(R)$ is
constructible. Since finite unions of constructible sets
are constructible, it suffices to deal with the case $n = 1$,
i.e., when $T = D(f) \cap V(g_1, \ldots, g_m)$.

\medskip\noindent
For $c \in R$, the following are disjoint decompositions of sets:
$$
\Spec(R) =
V(c) \amalg D(c) =
\Spec(R/(c)) \amalg \Spec(R_c),
$$
and correspondingly $\Spec(R[x]) =
V(c) \amalg D(c) = \Spec(R/(c)[x]) \amalg
\Spec(R_c[x])$. The intersection of $T = D(f) \cap V(g_1, \ldots, g_m)$
with each part still has the same shape, with $f$, $g_i$ replaced
by their images in $R/(c)[x]$, respectively $R_c[x]$.
Note that the image of $T$
in $\Spec(R)$ is the union of the image of
$T \cap V(c)$ and $T \cap D(c)$. Using Lemmas \ref{algebra-lemma-open-fp}
and \ref{algebra-lemma-closed-fp} it suffices to prove the images of both
parts are constructible in $\Spec(R/(c))$, respectively
$\Spec(R_c)$.

\medskip\noindent
We prove the assertion for
$T=D(f)\cap V(g_1,\ldots,g_m)$ over every ring by induction on
$$
N=\sum_{j=1}^m w(g_j),\qquad
w(g)=\begin{cases}0&g=0,\\ \deg(g)+1&g\ne0.\end{cases}
$$
All original polynomial labels are retained, including zero relations.
For $N=0$ all $g_j$ are zero, so $T=D(f)$ and
Lemma \ref{algebra-lemma-affineline-open} applies. If the coefficient ring is
zero, the spectrum is empty and the assertion also holds. For $N>0$
choose a nonzero $g_i$ of least degree among the nonzero relations and
write its full polynomial as
$g_i=c x^{d_i}+\sum_{a=0}^{d_i-1}c_{ia}x^a$, with $c\ne0$.

\medskip\noindent
Use the preceding decomposition into $R/(c)$ and $R_c$. Over $R/(c)$,
the weight of the image of $g_i$ strictly decreases, including when
$d_i=0$ and that image is zero. Every other weight is at most its
original value, so induction applies on this piece.
If $R_c=0$, its piece is empty. Otherwise the image of $c$ is a unit
and the image of $g_i$ retains degree $d_i$. If $g_i$ was the only
nonzero relation over $R$, Lemma \ref{editorial-source-lemma-affineline-special}
applies over $R_c$. If there was another nonzero relation, choose
one, say
$g_j=c_jx^{d_j}+\sum_{a=0}^{d_j-1}c_{ja}x^a$, where $d_j\geq d_i$
by the choice of $i$. In $R_c[x]$ replace only this relation by
$$
g'_j=g_j-(c_j/c)x^{d_j-d_i}g_i.
$$
Here every coefficient denotes its image in $R_c$; in particular
$c_j/c$ may be zero. The coefficient of $x^{d_j}$ cancels exactly,
so $g'_j=0$ or $\deg(g'_j)<d_j$. The other relations have weight
no larger than before base change, and the total weight is therefore
strictly less than the original $N$. Moreover,
$$
g'_j=g_j-(c_j/c)x^{d_j-d_i}g_i,\qquad
g_j=g'_j+(c_j/c)x^{d_j-d_i}g_i
$$
prove both inclusions between the original ideal and the ideal with
$g_j$ replaced by $g'_j$. Thus the two lists define exactly the same
subset $T$ over $R_c$, and induction proves its image constructible
there. In the notation $i=1,j=2$ this is precisely
$V(g_1,g'_2,g_3,\ldots,g_m)$, with every separator present.
The two induction base cases are a standard open and a standard open
intersected with the zero set of one polynomial having unit leading
coefficient, as in Lemmas \ref{algebra-lemma-affineline-open} and
\ref{editorial-source-lemma-affineline-special}. Finally,
Lemmas \ref{algebra-lemma-open-fp} and \ref{algebra-lemma-closed-fp} carry the two
constructible images back to $\Spec(R)$, and their finite union is
the desired image.
\end{proof}


\subsection{Editorial treatment of Algebra lines 5729--5767}

\hypertarget{algebra-context-039}{}

\noindent\href{https://github.com/stacks/stacks-project/blob/a04446e57ec1fbc252a871afcec7752fb2807b14/algebra.tex\#L5729}{Original source, lines 5729--5767}. Complete original and reviewed TeX passages accompany this note. Every intervention is separately recorded; the following treatment is editorial.


\begin{lemma}
\label{editorial-source-lemma-generic-finite-presentation}
Let $R \subset S$ be an inclusion of domains.
Assume that $R \to S$ is of finite type.
There exists a nonzero $f \in R$, and a nonzero $g \in S$
such that $R_f \to S_{fg}$ is of finite presentation.
\end{lemma}

\begin{proof}
We use induction on the number of generators of the original
$R$-algebra $S$. For no generators, $S=R$ and we take $f=g=1$.
For one generator write $S=R[x]/\mathfrak q$, where
$\mathfrak q\cap R=(0)$. If $\mathfrak q=(0)$, again take $f=g=1$.
Otherwise choose a nonzero polynomial $h\in\mathfrak q$ of least
degree and retain its full expression
$$
h=f x^d+a_{d-1}x^{d-1}+\cdots+a_0,
\qquad f\ne0,\quad f,a_0,\ldots,a_{d-1}\in R.
$$
Here $d\geq1$ since $R\to S$ is injective. Over $R_f$ its leading
coefficient $f$ is a unit, but the polynomial $h$ is unchanged.
Division by $h$, cancelling a leading term $b x^e$ by subtracting
$bf^{-1}x^{e-d}h$, gives a unique remainder of degree less than $d$.
Uniqueness holds because the leading coefficient of a nonzero
multiple of $h$ is a nonzero coefficient times the unit $f$.
The canonical map
$$
R_f[x]/(h)\longrightarrow S_f
$$
is surjective. If a remainder $v\in R_f[x]$ of degree less than $d$
maps to zero, choose $a\geq0$ with $f^av\in R[x]$.
Its image in $S_f$ is zero, so $f^b f^av$ maps to zero in $S$ for
some $b\geq0$. Since $S$ is a domain and $f\ne0$ in $S$, already
$f^av$ maps to zero in $S$. If $v\ne0$, the domain property of $R$
implies $f^av\ne0$ of degree less than $d$, contrary to the choice
of $h$. Thus $v=0$, proving injectivity and finite presentation.
This proves the one-generator case with the original $f$ and $g=1$.

\medskip\noindent
For $n>1$ let $x_1,\ldots,x_n$ be the original generators and
$S'=R[x_1,\ldots,x_{n-1}]\subset S$. By the induction hypothesis
there are nonzero $f\in R$ and $g\in S'$ for which
$R_f\to S'_{fg}$ is finitely presented. The extension
$S'_{fg}\subset S_{fg}$ is an inclusion of domains generated by
$x_n$. Apply the one-generator case to obtain nonzero
$f'\in S'_{fg}$ and $g'\in S_{fg}$ with
$$
(S'_{fg})_{f'}\longrightarrow (S_{fg})_{f'g'}
$$
finitely presented. Retain both elements and write their actual
fractions as $f'=a/(fg)^r$ and $g'=b/(fg)^s$, with
$a\in S'$, $b\in S$, $r,s\geq0$ and $a,b\ne0$.
Set $G=gab\in S$, which is nonzero. There are canonical inverse
$S$-algebra maps
$$
(S_{fg})_{f'g'}\ \longleftrightarrow\ S_{fG}.
$$
Indeed, in $S_{fG}$ the product $fgab$ is a unit, so each of
$f,g,a,b$ is a unit. Thus $fg,f',g'$ are units and the map from
the left ring exists. Conversely, in the left ring $fg$ and $f'g'$
are units, so $f,g,f',g'$ are units, and then
$a=f'(fg)^r$, $b=g'(fg)^s$ are units. Hence $fG=fgab$ is a unit
and the reverse map exists. The two composites fix $S$ and every
inverse adjoined in the indicated localizations, so they are identities.

\medskip\noindent
Finally $S'_{fg}\to(S'_{fg})_{f'}$ has the finite presentation
$S'_{fg}[z]/(f'z-1)$. Finite presentations compose: if
$P=A[X_1,\ldots,X_l]/(F_1,\ldots,F_v)$ and
$Q=P[Y_1,\ldots,Y_t]/(H_1,\ldots,H_w)$, lift each of the finitely
many coefficients in the $H_j$ to $A[X_1,\ldots,X_l]$ to obtain
$\widetilde H_j$. The canonical map identifies
$Q$ with $A[X_1,\ldots,X_l,Y_1,\ldots,Y_t]/(F_i,\widetilde H_j)$:
maps in both directions send every variable to its displayed class,
and their composites fix the coefficient ring and all variables.
Applying this twice proves that
$R_f\to(S_{fg})_{f'g'}\cong S_{fG}$ is finitely presented, with
the required elements $f\in R$ and $G\in S$.
\end{proof}


\subsection{Editorial treatment of Algebra lines 5769--5814}

\hypertarget{algebra-context-040}{}

\noindent\href{https://github.com/stacks/stacks-project/blob/a04446e57ec1fbc252a871afcec7752fb2807b14/algebra.tex\#L5769}{Original source, lines 5769--5814}. Complete original and reviewed TeX passages accompany this note. Every intervention is separately recorded; the following treatment is editorial.


\begin{lemma}
\label{editorial-source-lemma-characterize-image-finite-type}
Let $R \to S$ be a finite type ring map.
Denote $X = \Spec(R)$ and $Y = \Spec(S)$.
Write $f : Y \to X$ for the induced
map of spectra. Let $E \subset Y = \Spec(S)$ be a
constructible set.
If a point $\xi \in X$ is in $f(E)$, then
$\overline{\{\xi\}} \cap f(E)$ contains an open
dense subset of $\overline{\{\xi\}}$.
\end{lemma}

\begin{proof}
Let $\xi \in X$ be a point of $f(E)$. Choose a point $\eta \in E$
mapping to $\xi$. Let $\mathfrak p \subset R$ be the prime
corresponding to $\xi$ and let $\mathfrak q \subset S$ be the
prime corresponding to $\eta$. Consider the diagram
$$
\xymatrix{
\eta \ar[r] \ar@{|->}[d] & E \cap Y' \ar[r] \ar[d] &
Y' = \Spec(S/\mathfrak q) \ar[r] \ar[d] &
Y \ar[d] \\
\xi \ar[r] & f(E) \cap X' \ar[r] &
X' = \Spec(R/\mathfrak p) \ar[r] &
X
}
$$
By Lemma \ref{algebra-lemma-affine-map-quasi-compact} the set $E \cap Y'$
is constructible in $Y'$.
The original ring map induces an injection of domains
$R/\mathfrak p\hookrightarrow S/\mathfrak q$, still of finite type:
the original finite set of generators maps to a generating set of the
quotient. Let $\bar f:Y'\to X'$ be its map of spectra. The image
$\bar f(E\cap Y')$ is contained in $f(E)\cap X'$, so it suffices
to find a dense open subset of $X'$ contained in this smaller image.
The points $\xi,\eta$ are the generic points of $X',Y'$ respectively.
Lemma \ref{editorial-source-lemma-generic-finite-presentation} supplies nonzero
$a\in R/\mathfrak p$ and $b\in S/\mathfrak q$ with
$(R/\mathfrak p)_a\to(S/\mathfrak q)_{ab}$ finitely presented.
Thus $U=D(a)\subset X'$ and $V=D(ab)\subset Y'$ are dense opens,
contain $\xi$ and $\eta$, and $\bar f(V)\subset U$.
The set $E\cap Y'\cap V$ is constructible in $V$, by Topology,
Lemma \ref{topology-lemma-open-immersion-constructible-inverse-image}.
Chevalley's theorem (Theorem \ref{editorial-source-theorem-chevalley}) implies its
image $C$ is constructible in $U$ and contains $\xi$.
By Topology, Lemma \ref{topology-lemma-generic-point-in-constructible},
$C$ contains a dense open $U'\subset U$.
For completeness, a finite expression
$C=\bigcup_j(U_j\cap V_j^c)$ in the irreducible space $U$ has a
term containing $\xi$. Its open set $V_j$ cannot be nonempty:
every nonempty open of an irreducible space contains its generic
point. Hence $V_j=\emptyset$ and the nonempty open $U_j$ lies in
$C$. Such an open is dense in $U$. Since $U$ is open dense in $X'$,
$U'$ is open dense in $X'$ and is contained in $f(E)\cap X'$.
Under the original closed immersion
$X'=V(\mathfrak p)=\overline{\{\xi\}}\subset X$, this is exactly
the required assertion.
\end{proof}


\subsection{Editorial treatment of Algebra lines 5821--5869}

\hypertarget{algebra-context-041}{}

\noindent\href{https://github.com/stacks/stacks-project/blob/a04446e57ec1fbc252a871afcec7752fb2807b14/algebra.tex\#L5821}{Original source, lines 5821--5869}. Complete original and reviewed TeX passages accompany this note. Every intervention is separately recorded; the following treatment is editorial.


\begin{lemma}
\label{editorial-source-lemma-surjective-spec-radical-ideal}
Let $\varphi : R \to S$ be a ring map.
The following are equivalent:
\begin{enumerate}
\item The map $\Spec(S) \to \Spec(R)$ is surjective.
\item For any ideal $I \subset R$
the inverse image of $\sqrt{IS}$ in $R$ is equal to $\sqrt{I}$.
\item For any radical ideal $I \subset R$ the inverse image
of $IS$ in $R$ is equal to $I$.
\item For every prime $\mathfrak p$ of $R$ the inverse
image of $\mathfrak p S$ in $R$ is $\mathfrak p$.
\end{enumerate}
In this case the same is true after any base change: Given a ring map
$R \to R'$ the ring map $R' \to R' \otimes_R S$ has the equivalent
properties (1), (2), (3), and (4) as well.
\end{lemma}

\begin{proof}
For any ideal $J\subset S$, an element $x\in R$ belongs to
$\sqrt{\varphi^{-1}(J)}$ if and only if
$\varphi(x)^a\in J$ for some $a\geq1$. Thus
$\sqrt{\varphi^{-1}(J)}=\varphi^{-1}(\sqrt J)$.
Assume (1). For any ideal $I\subset R$, the inclusion
$\sqrt I\subset\varphi^{-1}(\sqrt{IS})$ follows by applying
$\varphi$ to a power lying in $I$. If $x\notin\sqrt I$, there is
a prime $\mathfrak p\supset I$ with $x\notin\mathfrak p$, by
Lemma \ref{editorial-source-lemma-Zariski-topology}. Choose a prime $\mathfrak q$
of $S$ over $\mathfrak p$. Then $IS\subset\mathfrak q$ but
$\varphi(x)\notin\mathfrak q$, so $\varphi(x)\notin\sqrt{IS}$.
This proves (2). If $I$ is radical, (2) gives
$I\subset\varphi^{-1}(IS)\subset\varphi^{-1}(\sqrt{IS})=I$,
proving (3), which immediately implies (4).
Finally assume (4). The quotient map $R/\mathfrak p\to S/\mathfrak pS$
is injective for every prime $\mathfrak p$. Inverting the nonzero
elements of the domain $R/\mathfrak p$ gives a nonzero ring:
otherwise some such element would map to zero in $S/\mathfrak pS$
by the equality criterion for $1/1=0/1$, contradicting injectivity.
This ring is $S\otimes_R\kappa(\mathfrak p)$, so
Lemma \ref{editorial-source-lemma-in-image} places $\mathfrak p$ in the image.
This proves (1), and hence the equivalence of all four properties.

\medskip\noindent
Assume (1) holds. Let $R \to R'$ be a ring map. Let
$\mathfrak p' \subset R'$ be a prime ideal lying over the prime
$\mathfrak p$ of $R$. To see that $\mathfrak p'$ is in the image
of $\Spec(R' \otimes_R S) \to \Spec(R')$ we have to show
that $(R' \otimes_R S) \otimes_{R'} \kappa(\mathfrak p')$ is not zero, see
Lemma \ref{editorial-source-lemma-in-image}.
Write $K=\kappa(\mathfrak p)$ and $K'=\kappa(\mathfrak p')$.
The given map induces an injection $\iota:K\hookrightarrow K'$:
its map $R/\mathfrak p\to R'/\mathfrak p'$ is injective and every
nonzero denominator remains nonzero in $K'$. There is an isomorphism
$$
(R'\otimes_RS)\otimes_{R'}K'
\longrightarrow (S\otimes_RK)\otimes_KK'
$$
sending $(r'\otimes s)\otimes\lambda$ to
$(s\otimes1)\otimes\overline{r'}\lambda$, where
$\overline{r'}$ denotes the image in $K'$.
Its inverse sends $(s\otimes a)\otimes\lambda$ to
$(1\otimes s)\otimes\iota(a)\lambda$.
The balancing relations over $R$, $R'$ and $K$ hold because the
coefficient maps into $K'$ agree; the maps preserve products and
units. Their composites are identities: in one direction move
$r'$ across the tensor product over $R'$, and in the other move
$a$ across the tensor product over $K$.
By (1) and Lemma \ref{editorial-source-lemma-in-image}, $S\otimes_RK$ is nonzero.
A $K$-basis containing its nonzero element $1$ gives a direct sum
of copies of $K$, whose tensor product with $K'$ is the corresponding
direct sum of copies of $K'$. It is nonzero. Thus every
$\mathfrak p'$ is in the image, proving (1) after base change and
therefore all four equivalent properties there.
\end{proof}


\subsection{Editorial treatment of Algebra lines 5920--5944}

\hypertarget{algebra-context-042}{}

\noindent\href{https://github.com/stacks/stacks-project/blob/a04446e57ec1fbc252a871afcec7752fb2807b14/algebra.tex\#L5920}{Original source, lines 5920--5944}. Complete original and reviewed TeX passages accompany this note. Every intervention is separately recorded; the following treatment is editorial.


\begin{lemma}
\label{editorial-source-lemma-image-dense-generic-points}
Let $R \to S$ be a ring map. The following are equivalent:
\begin{enumerate}
\item The kernel of $R \to S$ consists of nilpotent elements.
\item The minimal primes of $R$ are in the image of
$\Spec(S) \to \Spec(R)$.
\item The image of $\Spec(S) \to \Spec(R)$ is dense
in $\Spec(R)$.
\end{enumerate}
\end{lemma}

\begin{proof}
Write $\varphi:R\to S$ for the original map and $I=\ker\varphi$.
By Lemma \ref{editorial-source-lemma-Zariski-topology} and the power criterion for a
radical,
$$
\bigcap_{\mathfrak q\in\Spec(S)}\varphi^{-1}(\mathfrak q)
=\varphi^{-1}(\sqrt{(0)_S})=\sqrt I.
$$
Indeed $\varphi(x)$ is nilpotent exactly when $x^a\in I$ for some
$a\geq1$. A closed subset $V(J)\subset\Spec(R)$ contains the image
exactly when $J$ lies in this intersection. Hence its closure is
$V(\sqrt I)=V(I)$, proving (1) $\Leftrightarrow$ (3).
For (2) $\Rightarrow$ (3), every prime of $R$ contains a minimal
prime, and lies in the closure of that point, by
Lemma \ref{editorial-source-lemma-Zariski-topology}.
For (1) $\Rightarrow$ (2), the original map factors as
$$
R\longrightarrow R/I\lhook\joinrel\longrightarrow S.
$$
Every prime of $R$ contains $I$ under (1), and the quotient map
identifies $\Spec(R/I)$ homeomorphically with $\Spec(R)$ by
contraction, with inverse $\mathfrak p\mapsto\mathfrak p/I$.
These inverse order-preserving maps carry minimal primes to minimal
primes. The injection $R/I\hookrightarrow S$ hits all those points
by Lemma \ref{algebra-lemma-injective-minimal-primes-in-image}. Its contraction
composed with the quotient contraction is the contraction for
$\varphi$, proving (2) for the original map.
\end{proof}


\subsection{Editorial treatment of Algebra lines 5961--5975}

\hypertarget{algebra-context-043}{}

\noindent\href{https://github.com/stacks/stacks-project/blob/a04446e57ec1fbc252a871afcec7752fb2807b14/algebra.tex\#L5961}{Original source, lines 5961--5975}. Complete original and reviewed TeX passages accompany this note. Every intervention is separately recorded; the following treatment is editorial.


\begin{lemma}
\label{editorial-source-lemma-intersection-not-zero}
Let $A \subset B$ be an inclusion of domains inducing an algebraic extension
of fraction fields. If $J \subset B$ is a nonzero ideal, then $A \cap J$
is nonzero too. Thus the image of a proper closed subset of $\Spec(B)$
is not dense in $\Spec(A)$.
For a specified ideal $J$, the same conclusion $A\cap J\ne0$ already
follows if $J$ contains one nonzero element algebraic over
$\operatorname{Frac}(A)$; the entire fraction-field extension need
not be algebraic for this conclusion.
\end{lemma}

\begin{proof}
Let $x\in J$ be nonzero and algebraic over $K=\operatorname{Frac}(A)$.
Choose a nonzero polynomial $\sum_{i=0}^dc_iT^i\in K[T]$ of least
degree vanishing at $x$, with $c_d\ne0$. Then $d\geq1$.
Also $c_0\ne0$: otherwise dividing the equation by the nonzero
$x\in\operatorname{Frac}(B)$ gives a nonzero polynomial relation
of smaller degree. Write every coefficient as $c_i=b_i/t_i$ with
$b_i,t_i\in A$, $t_i\ne0$, keeping zero coefficients as well, and set
$$
D=\prod_{i=0}^dt_i,\qquad
a_i=b_i\prod_{\substack{0\leq j\leq d\\j\ne i}}t_j.
$$
Then $D\ne0$, $a_i=Dc_i$, $a_0,a_d\ne0$, and the complete equation is
$$
a_dx^d+a_{d-1}x^{d-1}+\cdots+a_0=0.
$$
Consequently $a_0=-\sum_{i=1}^d a_ix^i$ is a nonzero element of
$A\cap J$. This uses only the stated algebraicity of $x$.
For a proper closed subset $V(J)\subset\Spec(B)$, the domain property
ensures $J\ne0$, since $(0)$ is a prime. Its image is contained in
$V(a_0)\subset\Spec(A)$, a proper closed subset because
$a_0\ne0$ and $(0)$ is a prime of $A$. Thus that image is not dense.
\end{proof}


\subsection{Editorial treatment of Algebra lines 5980--6238}

\hypertarget{algebra-context-044}{}

\noindent\href{https://github.com/stacks/stacks-project/blob/a04446e57ec1fbc252a871afcec7752fb2807b14/algebra.tex\#L5980}{Original source, lines 5980--6238}. Complete original and reviewed TeX passages accompany this note. Every intervention is separately recorded; the following treatment is editorial.


\subsubsection{Noetherian rings}
\label{editorial-source-section-Noetherian}

\noindent
A ring $R$ is {\it Noetherian} if any ideal of $R$ is finitely generated.
This is equivalent to the ascending chain condition for ideals of $R$.
Indeed, the union of an ascending sequence of ideals is an ideal; if
it has a finite generating set, all those generators lie in one term,
and the sequence stabilizes there. Conversely, if an ideal is not
finitely generated, start with $(0)$ and successively choose an element
of the ideal outside the ideal generated by all the preceding choices.
This gives a strictly increasing sequence of ideals, contradicting
the ascending chain condition.
By
Lemma \ref{editorial-source-lemma-cohen}
it suffices to check that every prime ideal of $R$ is finitely generated.

\begin{lemma}
\label{editorial-source-lemma-Noetherian-permanence}
\begin{slogan}
Noetherian property is stable by passage to finite type extension
and localization.
\end{slogan}
Any finitely generated ring over a Noetherian ring
is Noetherian. Any localization of a Noetherian ring
is Noetherian.
\end{lemma}

\begin{proof}
The statement on localizations follows from the fact
that any ideal $J \subset S^{-1}R$ is of the form
$I \cdot S^{-1}R$. Any quotient $R/I$ of a Noetherian
ring $R$ is Noetherian because any ideal $\overline{J} \subset R/I$
is of the form $J/I$ for some ideal $I \subset J \subset R$.
Thus it suffices to show that if $R$ is Noetherian so
is $R[X]$. Suppose $J_1\subset J_2\subset\cdots$ is an ascending
chain of ideals in $R[X]$. For $i\geq1$ and $d\geq0$ define
$$
I_{i,d}=\{a\in R\mid aX^d+\sum_{e=0}^{d-1}b_eX^e\in J_i
\text{ for some }b_0,\ldots,b_{d-1}\in R\}.
$$
This contains zero, and addition and scalar multiplication of the
displayed polynomials show that it is an ideal. For a nonzero $a$,
the displayed polynomial has degree exactly $d$ and leading
coefficient $a$. Multiplication by $X^{d'-d}$ and the inclusion
$J_i\subset J_{i'}$ prove $I_{i,d}\subset I_{i',d'}$ when
$i\leq i'$ and $d\leq d'$.

\medskip\noindent
Put $K_d=\bigcup_{i\geq1}I_{i,d}$, an ideal because the union is
ascending in $i$. The sequence $K_0\subset K_1\subset\cdots$
stabilizes at some $d_0$. Each $K_d$ for $0\leq d\leq d_0$ has a
finite generating set. Every generator lies in some $I_{i,d}$;
taking the largest of these finitely many indices gives $i_d$ with
$I_{i_d,d}=K_d$ (take $i_d=1$ if the generating set is empty).
Let $i_0=\max_{0\leq d\leq d_0}i_d$. For $d\leq d_0$ and
$i\geq i_0$ we have $I_{i,d}=I_{i_0,d}=K_d$.
For $d>d_0$ the inclusions
$$
K_{d_0}=I_{i_0,d_0}\subset I_{i_0,d}\subset I_{i,d}
\subset K_d=K_{d_0}
$$
give the same equality. Thus the required $i_0$ works simultaneously
for every degree.

\medskip\noindent
For $i\geq i_0$ and $f\in J_i$, prove $f\in J_{i_0}$ by induction
on its degree, with $f=0$ treated separately. If $f\ne0$ has degree
$d$ and leading coefficient $a\ne0$, the equality
$I_{i,d}=I_{i_0,d}$ supplies $g\in J_{i_0}$ of degree $d$ with
the same coefficient $a$. Then $f-g\in J_i$ is zero or has degree
less than $d$. For $d=0$ it is zero; otherwise induction applies.
Hence $f-g\in J_{i_0}$ and $f\in J_{i_0}$. All $J_i$ with
$i\geq i_0$ are therefore equal, as required.
\end{proof}

\begin{lemma}
\label{editorial-source-lemma-Noetherian-power-series}
If $R$ is a Noetherian ring, then so is the formal power
series ring $R[[x_1, \ldots, x_n]]$.
\end{lemma}

\begin{proof}
The isomorphism
$R[[x_1,\ldots,x_{n+1}]]\cong R[[x_1,\ldots,x_n]][[x_{n+1}]]$
sends a coefficient family $(a_{\alpha,e})$ to
$\sum_{e\geq0}(\sum_{\alpha\in\mathbf Z_{\geq0}^n}a_{\alpha,e}
x_1^{\alpha_1}\cdots x_n^{\alpha_n})x_{n+1}^e$.
Its inverse reads off the same coefficients. For each fixed
multi-index, the Cauchy product has only finitely many contributing
pairs, so these mutually inverse maps preserve addition, multiplication
and the identity. It therefore suffices to prove that $R[[x]]$ is
Noetherian when $R$ is Noetherian. Let $I\subset R[[x]]$ be an ideal.
We have to show that $I$ is a finitely generated ideal.
For each nonnegative integer
$d$ denote $I_d = \{a \in R \mid ax^d + \text{h.o.t.} \in I\}$.
Then we see that $I_0 \subset I_1 \subset \ldots$ stabilizes as $R$
is Noetherian. Choose $d_0$ such that $I_{d_0} = I_{d_0 + 1} = \ldots$.
For each $d \leq d_0$ choose elements $f_{d, j} \in I \cap (x^d)$,
$j = 1, \ldots, n_d$ such that if we write
$f_{d, j} = a_{d, j}x^d + \text{h.o.t}$ then $I_d = (a_{d, j})$.
Denote $I' = (\{f_{d, j}\}_{d = 0, \ldots, d_0, j = 1, \ldots, n_d})$.
Then it is clear that $I' \subset I$. Pick $f \in I$.
Choose the coefficients $c_{d,i}$ successively for $d=0,\ldots,d_0$.
Before the step for $d$, the current remainder lies in $I\cap(x^d)$.
Its coefficient of $x^d$ belongs to $I_d$, so express it as
$\sum_{i=1}^{n_d}c_{d,i}a_{d,i}$ and subtract
$\sum_{i=1}^{n_d}c_{d,i}f_{d,i}$. The new remainder lies in
$I\cap(x^{d+1})$. Thus
$$
h_0=f-\sum_{d=0}^{d_0}\sum_{i=1}^{n_d}c_{d,i}f_{d,i}
\in I\cap(x^{d_0+1}).
$$
For every $e\geq1$, having constructed
$h_{e-1}\in I\cap(x^{d_0+e})$, its coefficient of $x^{d_0+e}$
lies in $I_{d_0+e}=I_{d_0}$. Choose $c_{i,e}\in R$ such that this
coefficient is $\sum_{i=1}^{n_{d_0}}c_{i,e}a_{d_0,i}$, and put
$$
h_e=h_{e-1}-\sum_{i=1}^{n_{d_0}}c_{i,e}x^e f_{d_0,i}
\in I\cap(x^{d_0+e+1}).
$$
These choices define actual formal power series
$C_i(x)=\sum_{e=1}^\infty c_{i,e}x^e\in R[[x]]$.
For any fixed degree $N$, the coefficient of $x^N$ in
$C_i(x)f_{d_0,i}$ receives contributions only from
$1\leq e\leq N-d_0$, a finite set. The remainder formula with
$d_0+e+1>N$ therefore proves, coefficient by coefficient, the identity
$$
f=\sum_{d=0}^{d_0}\sum_{i=1}^{n_d}c_{d,i}f_{d,i}
 +\sum_{i=1}^{n_{d_0}}\left(\sum_{e=1}^\infty c_{i,e}x^e\right)f_{d_0,i}.
$$
This is a finite $R[[x]]$-linear combination of the original chosen
generators, so $f$ belongs to $I'$. No assertion that $I'$ is closed
under limits is needed: the equality is an equality of all formal
coefficients with finitely many generator terms.
\end{proof}

\noindent
The following lemma, although easy, is useful because
finite type $\mathbf{Z}$-algebras come up quite often in
a technique called ``absolute Noetherian reduction''.

\begin{lemma}
\label{editorial-source-lemma-obvious-Noetherian}
Any finite type algebra over a field is Noetherian.
Any finite type algebra over $\mathbf{Z}$ is Noetherian.
\end{lemma}

\begin{proof}
This is immediate from Lemma \ref{editorial-source-lemma-Noetherian-permanence}
and the fact that fields are Noetherian rings and that
$\mathbf{Z}$ is a Noetherian ring (because it is a
principal ideal domain).
\end{proof}

\begin{lemma}
\label{editorial-source-lemma-Noetherian-finite-type-is-finite-presentation}
Let $R$ be a Noetherian ring.
\begin{enumerate}
\item Any finite $R$-module is of finite presentation.
\item Any submodule of a finite $R$-module is finite.
\item Any finite type $R$-algebra is of finite presentation over $R$.
\end{enumerate}
\end{lemma}

\begin{proof}
Let $M$ be a finite $R$-module. By
Lemma \ref{algebra-lemma-trivial-filter-finite-module}
we can find a finite filtration of $M$ whose successive quotients are
of the form $R/I$. Since any ideal is finitely generated, each of
the quotients $R/I$ is finitely presented. Hence $M$ is finitely
presented by
Lemma \ref{algebra-lemma-extension}.
This proves (1).

\medskip\noindent
Let $N \subset M$ be a submodule. As $M$ is finite, the quotient
$M/N$ is finite. Thus $M/N$ is of finite presentation by part (1).
Thus we see that $N$ is finite by Lemma \ref{algebra-lemma-extension} part (5).
This proves part (2).

\medskip\noindent
To see (3) note that any ideal of
$R[x_1, \ldots, x_n]$ is finitely generated by
Lemma \ref{editorial-source-lemma-Noetherian-permanence}.
\end{proof}

\begin{lemma}
\label{editorial-source-lemma-Noetherian-topology}
If $R$ is a Noetherian ring then $\Spec(R)$
is a Noetherian topological space, see Topology,
Definition \ref{topology-definition-noetherian}.
\end{lemma}

\begin{proof}
This is because any closed subset of $\Spec(R)$
is uniquely of the form $V(I)$ with $I$ a radical ideal,
see Lemma \ref{editorial-source-lemma-Zariski-topology}.
And this correspondence is inclusion reversing.
Thus the result follows from the definitions.
\end{proof}

\begin{lemma}
\label{editorial-source-lemma-Noetherian-irreducible-components}
\begin{slogan}
A Noetherian affine scheme has finitely many generic points.
\end{slogan}
If $R$ is a Noetherian ring then $\Spec(R)$
has finitely many irreducible components. In other words
$R$ has finitely many minimal primes.
\end{lemma}

\begin{proof}
By Lemma \ref{editorial-source-lemma-Noetherian-topology} and
Topology, Lemma \ref{topology-lemma-Noetherian}
we see there are finitely many irreducible components.
By Lemma \ref{editorial-source-lemma-irreducible} these correspond to
minimal primes of $R$.
\end{proof}

\begin{lemma}
\label{editorial-source-lemma-Noetherian-base-change-finite-type}
Let $R \to S$ be a ring map. Let $R \to R'$ be of finite type.
If $S$ is Noetherian, then the base change $S' = R' \otimes_R S$
is Noetherian.
\end{lemma}

\begin{proof}
By Lemma \ref{algebra-lemma-base-change-finiteness} finite type is stable under
base change. Thus $S \to S'$ is of finite type. Since $S$ is Noetherian we
can apply Lemma \ref{editorial-source-lemma-Noetherian-permanence}.
\end{proof}

\begin{lemma}
\label{editorial-source-lemma-Noetherian-field-extension}
Let $k$ be a field and let $R$ be a Noetherian $k$-algebra.
If $K/k$ is a finitely generated field extension then
$K \otimes_k R$ is Noetherian.
\end{lemma}

\begin{proof}
Since $K/k$ is a finitely generated field extension, there exists
a finitely generated $k$-algebra $B \subset K$ such that $K$ is
the fraction field of $B$. In other words, $K = S^{-1}B$
with $S = B \setminus \{0\}$. Then $K \otimes_k R = S^{-1}(B \otimes_k R)$.
Then $B \otimes_k R$ is Noetherian by
Lemma \ref{editorial-source-lemma-Noetherian-base-change-finite-type}.
Finally, $K \otimes_k R = S^{-1}(B \otimes_k R)$ is Noetherian by
Lemma \ref{editorial-source-lemma-Noetherian-permanence}.
\end{proof}

\noindent
Here are some fun lemmas that are sometimes useful.

\begin{lemma}
\label{editorial-source-lemma-subring-of-local-ring}
Let $R$ be a ring and $\mathfrak p \subset R$ be a prime.
There exists an $f \in R$, $f \not \in \mathfrak p$ such
that $R_f \to R_\mathfrak p$ is injective in each of the
following cases
\begin{enumerate}
\item $R$ is a domain,
\item $R$ is Noetherian, or
\item $R$ is reduced and has finitely many minimal primes.
\end{enumerate}
In (2), it suffices more generally that
$\ker(R\to R_{\mathfrak p})$ is finitely generated.
\end{lemma}

\begin{proof}
If $R$ is a domain, then $R \subset R_\mathfrak p$, hence $f = 1$ works.
More generally, suppose the kernel $I=\ker(R\to R_{\mathfrak p})$
is generated by $a_1,\ldots,a_t$. For each $a_j$, the localization
criterion gives $s_j\notin\mathfrak p$ with $s_ja_j=0$.
Set $f=\prod_{j=1}^ts_j$, using $f=1$ when $t=0$.
Then $f\notin\mathfrak p$ and $fI=0$.
If $r/f^e$ maps to zero in $R_{\mathfrak p}$, the image of $r$
is zero, so $r\in I$ and $fr=0$. Hence $r/f^e=0$ in $R_f$.
This proves injectivity, and applies in particular when $R$ is Noetherian.

\medskip\noindent
For (3), retain all the minimal primes
$\mathfrak p_1,\ldots,\mathfrak p_n$. For each $i$ with
$\mathfrak p_i\not\subset\mathfrak p$, choose
$f_i\in\mathfrak p_i\setminus\mathfrak p$, and put
$$
f=\prod_{\mathfrak p_i\not\subset\mathfrak p}f_i.
$$
The empty product is $1$, and primality of $\mathfrak p$ gives
$f\notin\mathfrak p$. If $r\in\ker(R\to R_{\mathfrak p})$,
some $s\notin\mathfrak p$ satisfies $sr=0$. Whenever
$\mathfrak p_i\subset\mathfrak p$, we have $s\notin\mathfrak p_i$,
so $r\in\mathfrak p_i$ by primality. For every other $i$ the
chosen factor $f_i$ gives $f\in\mathfrak p_i$.
It follows that $fr$ belongs to every minimal prime. The intersection
of the minimal primes is the nilradical, which is zero since $R$
is reduced. Thus $fr=0$. The same fraction argument as above proves
that the original canonical map $R_f\to R_{\mathfrak p}$ is injective.
\end{proof}

\begin{lemma}
\label{editorial-source-lemma-surjective-endo-noetherian-ring-is-iso}
Any surjective endomorphism of a Noetherian ring is an isomorphism.
\end{lemma}

\begin{proof}
Suppose $f:R\to R$ is surjective and not injective, and choose
$0\ne a\in\ker(f)$. For each $n\geq1$, surjectivity of $f^n$
provides $b_n\in R$ with $f^n(b_n)=a$. Then
$f^{n+1}(b_n)=f(a)=0$ but $f^n(b_n)\ne0$.
Thus every inclusion in
$$
\ker(f)\subset\ker(f^2)\subset\ker(f^3)\subset\cdots
$$
is strict, witnessed by $b_n$ at the $n$th inclusion. This contradicts
the ascending chain condition on ideals of the Noetherian ring $R$.
Consequently $f$ is injective as well as surjective, and is an isomorphism.
\end{proof}










\subsection{Editorial treatment of Algebra lines 6253--6265}

\hypertarget{algebra-context-045}{}

\noindent\href{https://github.com/stacks/stacks-project/blob/a04446e57ec1fbc252a871afcec7752fb2807b14/algebra.tex\#L6253}{Original source, lines 6253--6265}. Complete original and reviewed TeX passages accompany this note. Every intervention is separately recorded; the following treatment is editorial.


\begin{example}
\label{editorial-source-example-locally-nilpotent-not-nilpotent}
Let $R = k[x_n | n \in \mathbf{N}]$ be the polynomial ring in infinitely
many variables over a field $k$. Let $I$ be the ideal generated by
the elements $x_n^n$ for $n \in \mathbf{N}$ and $S = R/I$. Then the ideal
$J \subset S$ generated by the images of $x_n$, $n \in \mathbf{N}$
is locally nilpotent, but not nilpotent. The convention here is
$\mathbf N=\{1,2,3,\ldots\}$, so no relation with exponent zero occurs.
An element $y\in J$ has a finite expression
$y=\sum_{i=1}^t b_i\overline{x}_{n_i}$ with $b_i\in S$ and $n_i\geq1$.
Put $N=1+\sum_{i=1}^t(n_i-1)$. In the full multinomial expansion
of $y^N$, every term has an exponent at least $n_i$ on one factor
$\overline{x}_{n_i}$, and hence vanishes. The empty sum gives $y=0$.
Thus $J$ is locally nilpotent.
For each $n\geq1$ define a $k$-algebra map
$$
R\longrightarrow k[t]/(t^{n+1}),\qquad
x_{n+1}\longmapsto t,\qquad x_j\longmapsto0\quad(j\ne n+1).
$$
Every generator $x_j^j$ of $I$ maps to zero, so this map factors
through $S$. The element $\overline{x}_{n+1}^n\in J^n$ maps to
$t^n\ne0$, since $1,t,\ldots,t^n$ form a $k$-basis of the target.
Consequently $J^n\ne0$ for every $n\geq1$ and $J$ is not nilpotent.
\end{example}


\subsection{Editorial treatment of Algebra lines 6267--6277}

\hypertarget{algebra-context-046}{}

\noindent\href{https://github.com/stacks/stacks-project/blob/a04446e57ec1fbc252a871afcec7752fb2807b14/algebra.tex\#L6267}{Original source, lines 6267--6277}. Complete original and reviewed TeX passages accompany this note. Every intervention is separately recorded; the following treatment is editorial.


\begin{lemma}
\label{editorial-source-lemma-locally-nilpotent}
Let $R \to R'$ be a ring map and let $I \subset R$ be a locally nilpotent
ideal. Then $IR'$ is a locally nilpotent ideal of $R'$.
\end{lemma}

\begin{proof}
Write $\varphi:R\to R'$ for the map. Every element of $IR'$ has a
finite expression $y=\sum_{i=1}^t b_i\varphi(x_i)$ with $x_i\in I$
and $b_i\in R'$. Choose positive integers $d_i$ with $x_i^{d_i}=0$
and put $N=1+\sum_{i=1}^t(d_i-1)$. The complete expansion is
$$
y^N=\sum_{\substack{\alpha_1,\ldots,\alpha_t\geq0\\
\alpha_1+\cdots+\alpha_t=N}}
\frac{N!}{\alpha_1!\cdots\alpha_t!}
\prod_{i=1}^t b_i^{\alpha_i}\varphi(x_i)^{\alpha_i}.
$$
The displayed multinomial coefficients are integers, acting through
the canonical map $\mathbf Z\to R'$. Every summand has
$\alpha_i\geq d_i$ for some $i$, since otherwise its total degree
would be at most $N-1$. Thus every summand is zero. If $t=0$ then
$y=0$ and $N=1$ suffices. In particular the case of two nilpotents
$x^n=y^m=0$ gives $(x+y)^{n+m-1}=0$, with both original exponents
retained. This proves local nilpotence after extension.
\end{proof}


\subsection{Editorial treatment of Algebra lines 6279--6297}

\hypertarget{algebra-context-047}{}

\noindent\href{https://github.com/stacks/stacks-project/blob/a04446e57ec1fbc252a871afcec7752fb2807b14/algebra.tex\#L6279}{Original source, lines 6279--6297}. Complete original and reviewed TeX passages accompany this note. Every intervention is separately recorded; the following treatment is editorial.


\begin{lemma}
\label{editorial-source-lemma-locally-nilpotent-unit}
Let $R$ be a ring and let $I \subset R$ be a locally nilpotent
ideal.
An element $x$ of $R$ is a unit if and only if the image of $x$
in $R/I$ is a unit.
\end{lemma}

\begin{proof}
If $x$ is a unit in $R$, then its image is clearly a unit in $R/I$.
It remains to prove the converse.
Assume the image of $y \in R$ in $R/I$ is the inverse of the image of $x$.
Then $xy=1-z$ for some $z\in I$, and $z^N=0$ for a positive
integer $N$. Retain all terms in the finite geometric sum. They give
$$
x\left(y\sum_{j=0}^{N-1}z^j\right)
=(1-z)\sum_{j=0}^{N-1}z^j=1-z^N=1.
$$
Commutativity gives the reverse product as well, so the displayed
element $y\sum_{j=0}^{N-1}z^j$ is the inverse of the original $x$.
In particular, $1+I$ is a group under multiplication: it is closed
under products, and if $z\in I$ with $z^N=0$ then
$(1+z)^{-1}=\sum_{j=0}^{N-1}(-z)^j\in1+I$.
\end{proof}


\subsection{Editorial treatment of Algebra lines 6299--6316}

\hypertarget{algebra-context-048}{}

\noindent\href{https://github.com/stacks/stacks-project/blob/a04446e57ec1fbc252a871afcec7752fb2807b14/algebra.tex\#L6299}{Original source, lines 6299--6316}. Complete original and reviewed TeX passages accompany this note. Every intervention is separately recorded; the following treatment is editorial.


\begin{lemma}
\label{editorial-source-lemma-Noetherian-power}
\begin{slogan}
An ideal in a Noetherian ring is nilpotent if each element
of the ideal is nilpotent.
\end{slogan}
Let $R$ be a Noetherian ring. Let $I, J$ be ideals of $R$.
Suppose $J \subset \sqrt{I}$. Then $J^n \subset I$ for some $n$.
In particular, in a Noetherian ring the notions of
``locally nilpotent ideal''
and ``nilpotent ideal'' coincide.
More generally, the conclusion $J^n\subset I$ needs only finite
generation of $J$, not Noetherianity of $R$. If
$J=(f_1,\ldots,f_s)$ and $f_i^{d_i}\in I$ with $d_i\geq1$, one may
take $n=1+\sum_{i=1}^s(d_i-1)$.
\end{lemma}

\begin{proof}
Choose generators $J=(f_1,\ldots,f_s)$ and positive exponents
$d_i$ such that $f_i^{d_i}\in I$. Put $N=1+\sum_{i=1}^s(d_i-1)$.
The ideal $J^N$ is generated by the full list of monomials
$\prod_i f_i^{\alpha_i}$ with $\alpha_i\geq0$ and $\sum_i\alpha_i=N$.
Every such list has $\alpha_i\geq d_i$ for some $i$, so that monomial
belongs to $I$. Hence $J^N\subset I$. For $s=0$, $J=0$ and $N=1$
works. The original larger choice $d_1+\cdots+d_s+1$ also works,
since every higher power of $J$ is contained in $J^N$.
Only the existence of a finite generating set for $J$ was used.

\medskip\noindent
The bound cannot be decreased using only the displayed exponents.
In $B=\mathbf Z[T_1,\ldots,T_s]/(T_1^{d_1},\ldots,T_s^{d_s})$,
the residue classes of monomials with $0\leq\alpha_i<d_i$ form a
$\mathbf Z$-basis: the defining monomial ideal is the free subgroup
spanned by all the other monomials. Therefore
$\prod_i T_i^{d_i-1}$ is a nonzero element of
$(T_1,\ldots,T_s)^{N-1}$. Together with the upper bound this proves
sharpness, including the empty product in degree zero.
\end{proof}


\subsection{Editorial treatment of Algebra lines 6318--6383}

\hypertarget{algebra-context-049}{}

\noindent\href{https://github.com/stacks/stacks-project/blob/a04446e57ec1fbc252a871afcec7752fb2807b14/algebra.tex\#L6318}{Original source, lines 6318--6383}. Complete original and reviewed TeX passages accompany this note. Every intervention is separately recorded; the following treatment is editorial.


\begin{lemma}
\label{editorial-source-lemma-lift-idempotents}
Let $R$ be a ring. Let $I \subset R$ be a locally nilpotent ideal.
Then $R \to R/I$ induces a bijection on idempotents.
\end{lemma}

\begin{proof}[First proof of Lemma \ref{editorial-source-lemma-lift-idempotents}]
As $I$ is locally nilpotent it is contained in every prime ideal.
Hence $\Spec(R/I) = V(I) = \Spec(R)$. Hence the
lemma follows from Lemma \ref{algebra-lemma-disjoint-decomposition}.
\end{proof}

\begin{proof}[Second proof of Lemma \ref{editorial-source-lemma-lift-idempotents}]
Suppose $\overline{e} \in R/I$ is an idempotent.
We have to lift $\overline{e}$ to an idempotent of $R$.

\medskip\noindent
First, choose any lift $f \in R$ of $\overline{e}$, and set
$x = f^2 - f$. Then, $x \in I$, so $x$ is nilpotent (since $I$
is locally nilpotent). Let now $J$ be the ideal of $R$ generated
by $x$. Then, $J$ is nilpotent (not just locally nilpotent),
since it is generated by the nilpotent $x$.

\medskip\noindent
Now, assume that we have found a lift $e \in R$ of $\overline{e}$
such that $e^2 - e \in J^k$ for some $k \geq 1$.
Let $e' = e - (2e - 1)(e^2 - e) = 3e^2 - 2e^3$, which is another
lift of $\overline{e}$ (since the idempotency of $\overline{e}$
yields $e^2 - e \in I$). Then
$$
(e')^2 - e' = (4e^2 - 4e - 3)(e^2 - e)^2 \in J^{2k}
$$
because expansion of either side gives
$4e^6-12e^5+9e^4+2e^3-3e^2$.
In particular every coefficient and sign in the original update is retained.

\medskip\noindent
Choose $N\geq1$ with $x^N=0$, so $J^N=0$.
Starting with $e_0=f$, define for every $j\geq0$
$$
e_{j+1}=e_j-(2e_j-1)(e_j^2-e_j)=3e_j^2-2e_j^3.
$$
The identity just proved gives inductively
$e_j^2-e_j\in J^{2^j}$, and every $e_j$ still maps to
$\overline e$ since each update differs by an element of $J$.
Any integer $j$ with $2^j\geq N$ therefore gives
$e_j^2-e_j=0$. This constructs the required lift after finitely
many explicitly specified updates.

\medskip\noindent
We thus have seen that if $\overline{e} \in R/I$ is any
idempotent, then there exists a lift of $\overline{e}$
which is an idempotent of $R$.
It remains to prove that this lift is unique. Indeed, let
$e_1$ and $e_2$ be two such lifts. We
need to show that $e_1 = e_2$.

\medskip\noindent
By definition of $e_1$ and $e_2$, we have $e_1 \equiv e_2
\mod I$, and both $e_1$ and $e_2$ are idempotent. From
$e_1 \equiv e_2 \mod I$, we see that $e_1 - e_2 \in I$,
so that $e_1 - e_2$ is nilpotent (since $I$ is locally nilpotent).
A straightforward
computation (using the idempotency of $e_1$ and $e_2$)
reveals that $(e_1 - e_2)^3 = e_1 - e_2$. Using this and
induction, we obtain $(e_1 - e_2)^k = e_1 - e_2$ for any
positive odd integer $k$. Since all high enough $k$ satisfy
$(e_1 - e_2)^k = 0$ (since $e_1 - e_2$ is nilpotent),
this shows $e_1 - e_2 = 0$, so that $e_1 = e_2$, which
completes our proof.
\end{proof}


\subsection{Editorial treatment of Algebra lines 6385--6399}

\hypertarget{algebra-context-050}{}

\noindent\href{https://github.com/stacks/stacks-project/blob/a04446e57ec1fbc252a871afcec7752fb2807b14/algebra.tex\#L6385}{Original source, lines 6385--6399}. Complete original and reviewed TeX passages accompany this note. Every intervention is separately recorded; the following treatment is editorial.


\begin{lemma}
\label{editorial-source-lemma-lift-idempotents-noncommutative}
Let $A$ be a possibly noncommutative algebra.
Let $e \in A$ be an element such that $x = e^2 - e$ is nilpotent.
Then there exists an idempotent of the form
$e' = e + x(\sum a_{i, j}e^ix^j) \in A$
with $a_{i, j} \in \mathbf{Z}$.
\end{lemma}

\begin{proof}
Choose $n\geq1$ with $x^n=0$, and distinguish the indeterminate
$E$ from the original element $e\in A$. Evaluation gives a unital
ring map
$$
\mathbf Z[E]/((E^2-E)^n)\longrightarrow A,\qquad E\longmapsto e.
$$
This map exists even if $A$ is noncommutative: integer multiples of
$1_A$ are central and every polynomial in the single element $e$
commutes with every other such polynomial. The relation vanishes
because $(e^2-e)^n=x^n=0$.
In the commutative source ring let $I=(E^2-E)$ and apply
Lemma \ref{editorial-source-lemma-lift-idempotents} to the class of $E$ modulo $I$.
Its idempotent lift $P$ differs from $E$ by an element of $I$,
so $P=E+(E^2-E)Q(E)$ for a polynomial $Q\in\mathbf Z[E]$.
Its image is the idempotent $e'=e+xQ(e)$ in $A$.
Writing $Q(E)=\sum_i a_{i,0}E^i$ gives the required finite sum,
with $a_{i,j}=0$ for $j>0$. Alternatively the same explicit cubic
updates in the preceding proof construct $P$. No uniqueness of
arbitrary idempotent lifts in the whole noncommutative algebra is asserted.
\end{proof}


\subsection{Editorial treatment of Algebra lines 6401--6438}

\hypertarget{algebra-context-051}{}

\noindent\href{https://github.com/stacks/stacks-project/blob/a04446e57ec1fbc252a871afcec7752fb2807b14/algebra.tex\#L6401}{Original source, lines 6401--6438}. Complete original and reviewed TeX passages accompany this note. Every intervention is separately recorded; the following treatment is editorial.


\begin{lemma}
\label{editorial-source-lemma-lift-nth-roots}
Let $R$ be a ring. Let $I \subset R$ be a locally nilpotent ideal.
Let $n \geq 1$ be an integer which is invertible in $R/I$. Then
\begin{enumerate}
\item the $n$th power map $1 + I \to 1 + I$, $1 + x \mapsto (1 + x)^n$
is a bijection,
\item a unit of $R$ is an $n$th power if and only if its image in $R/I$
is an $n$th power.
\end{enumerate}
More precisely, for each unit $a\in R$, reduction induces a bijection
between the solutions $u\in R^*$ of $u^n=a$ and the solutions
$\overline u\in(R/I)^*$ of $\overline u^{,n}=\overline a$.
\end{lemma}

\begin{proof}
Let $a \in R$ be a unit whose image in $R/I$ is the same as the image
of $b^n$ with $b \in R$. Then $b$ is a unit
(Lemma \ref{editorial-source-lemma-locally-nilpotent-unit}) and
$ab^{-n} = 1 + x$ for some $x \in I$. Hence $ab^{-n} = c^n$ by
part (1). Thus (2) follows from (1).

\medskip\noindent
Proof of (1). This is true because there is an inverse to the
map $1 + x \mapsto (1 + x)^n$. Namely, we can consider the map
which sends $1 + x$ to
\begin{align*}
(1 + x)^{1/n}
& =
1 + {1/n \choose 1}x +
{1/n \choose 2}x^2 +
{1/n \choose 3}x^3 + \ldots \\
& =
1 + \frac{1}{n} x + \frac{1 - n}{2n^2}x^2 +
\frac{(1 - n)(1 - 2n)}{6n^3}x^3 + \ldots
\end{align*}
We justify the coefficients and both inverse properties over the original ring.
For $k\geq1$ retain the full coefficient formula
$$
c_k={1/n\choose k}
=\frac{\prod_{j=0}^{k-1}(1-jn)}{n^k k!},\qquad c_0=1.
$$
For each prime $p$ not dividing $n$ and each $a\geq1$, exactly one
residue class of $j$ modulo $p^a$ solves $1-jn\equiv0\pmod{p^a}$.
Among $k$ consecutive integers $j=0,\ldots,k-1$ there are at least
$\lfloor k/p^a\rfloor$ elements of this class. If the numerator is
nonzero, its $p$-adic valuation is therefore at least
$\sum_{a\geq1}\lfloor k/p^a\rfloor$, which equals the valuation
of $k!$: each multiple of $p^a$ among $1,\ldots,k$ contributes once
to that sum. The denominator $n^k$ contributes nothing at $p$.
If the numerator is zero, the coefficient is zero and the conclusion
is immediate. Thus the denominator of $c_k$ has no prime divisor
outside those dividing $n$, and $c_k\in\mathbf Z[1/n]$.
This statement does not require $k!$ to be invertible in $R$.
The integer $n$ is invertible in $R$ by
Lemma \ref{editorial-source-lemma-locally-nilpotent-unit}, so every $c_k$ has a
well-defined image in $R$ via $\mathbf Z[1/n]\to R$.

\medskip\noindent
Set $F(X)=\sum_{k\geq0}c_kX^k\in\mathbf Z[1/n][[X]]$.
The full product formula gives
$n(k+1)c_{k+1}=(1-kn)c_k$; hence in $\mathbf Q[[X]]$ it gives
$n(1+X)F'(X)=F(X)$ and $F(0)=1$.
For $H=F^n$, the product rule yields $(1+X)H'=H$, so the derivative
of $H/(1+X)$ is zero. A formal series over $\mathbf Q$ with zero
derivative is constant, and its constant term here is $1$.
Therefore $F(X)^n=1+X$. The coefficientwise injection
$\mathbf Z[1/n][[X]]\hookrightarrow\mathbf Q[[X]]$ proves the same
identity over $\mathbf Z[1/n]$.
For $x\in I$ with $x^N=0$, reduction modulo $X^N$ followed by
$X\mapsto x$ makes the finite expression
$c=\sum_{k=0}^{N-1}c_kx^k\in1+I$ satisfy $c^n=1+x$.
Every larger truncation gives the same element, since the additional
powers of $x$ vanish. This proves surjectivity in (1).

\medskip\noindent
For injectivity, if $u,v\in1+I$ and $u^n=v^n$, use the complete
factorization
$$
0=u^n-v^n=(u-v)\sum_{j=0}^{n-1}u^{n-1-j}v^j.
$$
The sum maps to the unit $n$ in $R/I$, hence is a unit of $R$ by
Lemma \ref{editorial-source-lemma-locally-nilpotent-unit}. Thus $u=v$.
The constructed map is consequently the inverse in (1).

\medskip\noindent
Finally fix a unit $a$ and a root $\overline b$ of $\overline a$.
Choose a lift $b\in R$; it is a unit, and $ab^{-n}\in1+I$.
The unique root $c\in1+I$ furnished by (1) gives $u=bc$ with
$u^n=a$ and $\overline u=\overline b$.
If another such lift is $v$, then $uv^{-1}\in1+I$ and
$(uv^{-1})^n=1$, so injectivity in (1) gives $u=v$.
This proves the asserted bijection of root sets and also (2).
\end{proof}


\subsection{Editorial treatment of Algebra lines 6455--6477}

\hypertarget{algebra-context-052}{}

\noindent\href{https://github.com/stacks/stacks-project/blob/a04446e57ec1fbc252a871afcec7752fb2807b14/algebra.tex\#L6455}{Original source, lines 6455--6477}. Complete original and reviewed TeX passages accompany this note. Every intervention is separately recorded; the following treatment is editorial.


\begin{lemma}
\label{editorial-source-lemma-invert-closed-quotient}
Let $R$ be a ring. Let $S \subset R$ be a multiplicative subset.
Assume the image of the map $\Spec(S^{-1}R) \to \Spec(R)$
is closed. Then $S^{-1}R \cong R/I$ for some ideal $I \subset R$.
\end{lemma}

\begin{proof}
Let $I = \Ker(R \to S^{-1}R)$ so that $V(I)$ contains the image.
Say the image is the closed subset $V(I') \subset \Spec(R)$ for
some ideal $I' \subset R$. So $V(I') \subset V(I)$.
For $f \in I'$ we see that $f/1 \in S^{-1}R$
belongs to every prime ideal. Hence $f^n$ maps to zero in $S^{-1}R$
for some $n \geq 1$ (Lemma \ref{editorial-source-lemma-Zariski-topology}).
Hence $V(I') = V(I)$.
Every $g\in S$ avoids all primes in the image, hence every prime
containing $I$. Its class $\overline g$ in $R/I$ belongs to no maximal
ideal and is therefore a unit. The quotient and localization properties
give ring maps
$$
\alpha:R/I\longrightarrow S^{-1}R,\quad \overline r\longmapsto r/1,
\qquad
\beta:S^{-1}R\longrightarrow R/I,\quad r/s\longmapsto
\overline r\,\overline s^{-1}.
$$
For every $\overline r$, $\beta\alpha(\overline r)=\overline r$.
For every $r/s$, $\alpha\beta(r/s)=(r/1)(s/1)^{-1}=r/s$.
Thus the maps are inverse, and their composite with $R\to R/I$
is the original localization map. In particular the localization map
is surjective. Conversely, if $R\to S^{-1}R$ is surjective, it is
the quotient by its kernel and its spectrum has closed image by
Lemma \ref{algebra-lemma-spec-closed}.
\end{proof}


\subsection{Editorial treatment of Algebra lines 6479--6505}

\hypertarget{algebra-context-053}{}

\noindent\href{https://github.com/stacks/stacks-project/blob/a04446e57ec1fbc252a871afcec7752fb2807b14/algebra.tex\#L6479}{Original source, lines 6479--6505}. Complete original and reviewed TeX passages accompany this note. Every intervention is separately recorded; the following treatment is editorial.


\begin{lemma}
\label{editorial-source-lemma-invert-closed-split}
Let $R$ be a ring. Let $S \subset R$ be a multiplicative subset.
Assume the image of the map $\Spec(S^{-1}R) \to \Spec(R)$
is closed. If $R$ is Noetherian, or $\Spec(R)$ is a
Noetherian topological space, or $S$ is finitely generated as a monoid,
then $R\cong S^{-1}R\times R'$ for some ring $R'$.
More generally, the same conclusion holds whenever the complement
of the closed image is quasi-compact.
\end{lemma}

\begin{proof}
By Lemma \ref{editorial-source-lemma-invert-closed-quotient} we have $S^{-1}R \cong R/I$
for some ideal $I \subset R$. By Lemma \ref{algebra-lemma-disjoint-implies-product}
it suffices to show that $V(I)$ is open.
If $R$ is Noetherian then $\Spec(R)$ is a Noetherian
topological space, see Lemma \ref{editorial-source-lemma-Noetherian-topology}.
If $\Spec(R)$ is Noetherian, the complement of $V(I)$ is
quasi-compact by Topology, Lemma
\ref{topology-lemma-Noetherian-quasi-compact}. More generally assume
only that this complement is quasi-compact. Its cover by the opens
$D(f)$, $f\in I$, has a finite subcover $D(f_1),\ldots,D(f_n)$,
so $V(I)=V(f_1,\ldots,f_n)$. For each $i$ choose $s_i\in S$ with
$s_if_i=0$, using the localization kernel criterion, and set
$g=\prod_{i=1}^n s_i\in S$ (with empty product $1$).
Then $gf_i=0$ for every $i$. If a prime does not contain $g$, it
contains every $f_i$, so $D(g)\subset V(I)$. Conversely, every prime
in $V(I)$ lies in the localization image and avoids all of $S$,
including $g$. Hence $V(I)=D(g)$, which is open.

\medskip\noindent
If instead $S$ is generated as a monoid by $g_1,\ldots,g_m$, put
$t=g_1\cdots g_m$, with $t=1$ when the list is empty.
Each $g_i$ is a unit in $R_t$, because a factor of a unit in a
commutative ring is a unit. Conversely $t$ is a unit in $S^{-1}R$.
The resulting canonical maps $S^{-1}R\rightleftarrows R_t$ are
inverse: their composites fix $R$ and the unique inverses adjoined.
Thus their spectrum image is $D(t)$, again proving $V(I)$ open.

\medskip\noindent
To identify the original localization as the first product factor,
choose an idempotent $e$ with $D(e)=V(I)$ by
Lemma \ref{algebra-lemma-disjoint-decomposition}. Each $s\in S$ avoids every
prime of $R_e$, because those primes correspond to $D(e)=V(I)$;
thus it is a unit in $R_e$. Likewise $e/1$ avoids every prime of
$S^{-1}R$, so it is a unit there. Being idempotent, it equals $1$.
The canonical maps $S^{-1}R\rightleftarrows R_e$ therefore exist and
are inverse by the same identity-on-$R$ and inverse-element argument.
Finally
$$
R\longrightarrow R_e\times R_{1-e},\qquad
r\longmapsto(r/1,r/1)
$$
has inverse $(a/e^u,b/(1-e)^v)\longmapsto ea+(1-e)b$.
This formula is independent of fraction representatives: equality
in $R_e$ implies $e(a-a')=0$, and equality in $R_{1-e}$ implies
$(1-e)(b-b')=0$, by the localization equality criterion and
idempotency. The two composites are identities, since $e$ maps to
$1$ in $R_e$ and $0$ in $R_{1-e}$, and $e+(1-e)=1$ in $R$.
The map is a ring map, and hence so is its inverse. Taking
$R'=R_{1-e}$ gives the claimed decomposition, whose first projection
is exactly the original map $R\to S^{-1}R$.
\end{proof}


\subsection{Editorial treatment of Algebra lines 6523--6608}

\hypertarget{algebra-context-054}{}

\noindent\href{https://github.com/stacks/stacks-project/blob/a04446e57ec1fbc252a871afcec7752fb2807b14/algebra.tex\#L6523}{Original source, lines 6523--6608}. Complete original and reviewed TeX passages accompany this note. Every intervention is separately recorded; the following treatment is editorial.


\begin{theorem}[Hilbert Nullstellensatz]
\label{editorial-source-theorem-nullstellensatz}
Let $k$ be a field.
\begin{enumerate}
\item
\label{editorial-source-item-finite-kappa}
For any maximal ideal $\mathfrak m \subset k[x_1, \ldots, x_n]$
the field extension $\kappa(\mathfrak m)/k$ is finite.
\item
\label{editorial-source-item-polynomial-ring-Jacobson}
Any radical ideal $I \subset k[x_1, \ldots, x_n]$
is the intersection of maximal ideals containing it.
\end{enumerate}
The same is true in any finite type $k$-algebra.
\end{theorem}

\begin{proof}
It is enough to prove part (\ref{editorial-source-item-finite-kappa}) of
the theorem for the case of a polynomial
algebra $k[x_1, \ldots, x_n]$, because any finitely generated
$k$-algebra is a quotient of such a polynomial algebra.
We prove this by induction on $n$. The case $n = 0$ is clear.
Suppose that $\mathfrak m$ is a maximal ideal in $k[x_1, \ldots, x_n]$.
Let $\mathfrak p \subset k[x_n]$ be the intersection
of $\mathfrak m$ with $k[x_n]$.

\medskip\noindent
If $\mathfrak p \not = (0)$,
then $\mathfrak p$ is maximal and generated by an irreducible
monic polynomial $P$ (because of the Euclidean algorithm
in $k[x_n]$). Then
$k' = k[x_n]/\mathfrak p$ is a finite field extension of $k$
and contained in $\kappa(\mathfrak m)$. In this case
we get a surjection
$$
k'[x_1, \ldots, x_{n-1}]
\to
k'[x_1, \ldots, x_n] =
k' \otimes_k k[x_1, \ldots, x_n]
\longrightarrow
\kappa(\mathfrak m)
$$
where the last map sends the coefficient field $k'$ into
$\kappa(\mathfrak m)$ and each variable to its original residue
class. The image of $x_n$ already belongs to this coefficient field:
it is the class of $x_n$ in $k'=k[x_n]/\mathfrak p$. Hence the
composite from $k'[x_1,\ldots,x_{n-1}]$ is surjective. Its kernel
is maximal, and the induction hypothesis makes
$\kappa(\mathfrak m)/k'$ finite. The extension $k'/k$ is finite,
with basis $1,\overline{x}_n,\ldots,\overline{x}_n^{d-1}$ where
$d=\deg(P)$. Products of this basis with a $k'$-basis of
$\kappa(\mathfrak m)$ span the latter over $k$, proving finiteness
over the original field.

\medskip\noindent
If $\mathfrak p = (0)$ we consider the ring
extension $k[x_n] \subset k[x_1, \ldots, x_n]/\mathfrak m$.
This is a finitely generated ring extension, hence
of finite presentation by
Lemmas \ref{editorial-source-lemma-obvious-Noetherian} and
\ref{editorial-source-lemma-Noetherian-finite-type-is-finite-presentation}.
Thus the image of $\Spec(k[x_1, \ldots, x_n]/\mathfrak m)$
in $\Spec(k[x_n])$ is constructible by
Theorem \ref{editorial-source-theorem-chevalley}. Since the image
contains $(0)$ we conclude that it contains a standard
open $D(f)$ for a nonzero $f\in k[x_n]$.
This open contains infinitely many closed points, over finite fields
as well as infinite fields. Indeed, given any finite list of
irreducible polynomials $p_1,\ldots,p_t$ defining points already
chosen in $D(f)$, form the original polynomial
$$
q=1+x_n f\prod_{i=1}^t p_i.
$$
It has positive degree and hence an irreducible divisor $h$.
Such a divisor exists by induction on degree, factoring any reducible
nonconstant polynomial until an irreducible factor is reached.
No irreducible divisor of $f$ or of any $p_i$ can divide $q$, since
$q$ is congruent to $1$ modulo each such divisor. Thus $(h)$ is a
new maximal ideal not containing $f$. The empty list is allowed,
and this construction proves infinitude. But the field
$k[x_1,\ldots,x_n]/\mathfrak m$ has just one point in its spectrum;
its image cannot contain $D(f)$. This is the required contradiction.

\medskip\noindent
Proof of (\ref{editorial-source-item-polynomial-ring-Jacobson}). Let
$I \subset R$ be a radical ideal, with $R$ of finite type over $k$.
Let $f \in R$, $f \not \in I$. We have to find a maximal ideal
$\mathfrak m \subset R$ with $I \subset \mathfrak m$ and
$f \not \in \mathfrak m$. The ring $(R/I)_f$ is nonzero, since
$1 = 0$ in this ring would mean $f^n \in I$ and since $I$ is
radical this would mean $f \in I$ contrary to our assumption on $f$.
Thus we may choose a maximal ideal $\mathfrak m'$
in $(R/I)_f$, see Lemma \ref{editorial-source-lemma-Zariski-topology}.
Let $\mathfrak m \subset R$
be the inverse image of $\mathfrak m'$ in $R$. We see that
$I \subset \mathfrak m$
and $f \not \in \mathfrak m$. If we show that $\mathfrak m$ is a maximal
ideal of $R$, then we are done. We clearly have
$$
k \subset R/\mathfrak m \subset \kappa(\mathfrak m').
$$
The algebra $(R/I)_f$ is of finite type over $k$: it is the quotient
of $R[t]$ by the ideal generated by $I$ and $ft-1$, with $t$ mapping
to the inverse of the class of $f$. The mutually inverse maps send
$r/f^j$ to $rt^j$ and $t$ to $1/f$.
Thus part (\ref{editorial-source-item-finite-kappa}) applies and
$\kappa(\mathfrak m')/k$ is finite. The original injection
$R/\mathfrak m\hookrightarrow\kappa(\mathfrak m')$ makes
$R/\mathfrak m$ a finite dimensional $k$-vector space and a domain.
For each nonzero $a\in R/\mathfrak m$, multiplication by $a$ is
an injective linear endomorphism of this finite dimensional space,
hence surjective. Its image contains $1$, giving the inverse of $a$.
Consequently $R/\mathfrak m$ is a field, as also stated in Fields,
Lemma \ref{fields-lemma-subalgebra-algebraic-extension-field}.
Thus $\mathfrak m$ is maximal and the proof is complete.
\end{proof}


\subsection{Editorial treatment of Algebra lines 6610--6627}

\hypertarget{algebra-context-055}{}

\noindent\href{https://github.com/stacks/stacks-project/blob/a04446e57ec1fbc252a871afcec7752fb2807b14/algebra.tex\#L6610}{Original source, lines 6610--6627}. Complete original and reviewed TeX passages accompany this note. Every intervention is separately recorded; the following treatment is editorial.


\begin{lemma}
\label{editorial-source-lemma-field-finite-type-over-domain}
Let $R$ be a ring. Let $K$ be a field.
If $R \subset K$ and $K$ is of finite type over $R$,
then there exists an $f \in R$ such that $R_f$ is a field,
and $K/R_f$ is a finite field extension.
\end{lemma}

\begin{proof}
By Lemma \ref{editorial-source-lemma-characterize-image-finite-type} there
exist a nonempty open $U \subset \Spec(R)$
contained in the image $\{(0)\}$ of $\Spec(K) \to \Spec(R)$.
Choose $f \in R$, $f \not = 0$ such that $D(f) \subset U$, i.e.,
$D(f)=\{(0)\}$. The localization $R_f$ is a domain whose only
prime ideal is $(0)$. Every nonunit belongs to a maximal ideal,
so its only nonunit is zero and $R_f$ is a field. The original
inclusion $R\subset K$ extends to $R_f\subset K$, and the original
finite set of $R$-algebra generators of $K$ is also a set of
$R_f$-algebra generators. The Hilbert Nullstellensatz
(Theorem \ref{editorial-source-theorem-nullstellensatz}) therefore makes $K/R_f$ finite.
In fact the canonical inclusion $R_f\subset\operatorname{Frac}(R)$
is an equality: for $r/s\in\operatorname{Frac}(R)$ with $s\ne0$,
the element $s/1$ is nonzero in the field $R_f$ and has an inverse
there, so $r/s$ belongs to $R_f$. Thus the same $f$ realizes the
original fraction field as a principal localization of $R$.
\end{proof}


\subsection{Editorial treatment of Algebra lines 6727--6776}

\hypertarget{algebra-context-056}{}

\noindent\href{https://github.com/stacks/stacks-project/blob/a04446e57ec1fbc252a871afcec7752fb2807b14/algebra.tex\#L6727}{Original source, lines 6727--6776}. Complete original and reviewed TeX passages accompany this note. Every intervention is separately recorded; the following treatment is editorial.


\begin{lemma}
\label{editorial-source-lemma-characterize-jacobson}
Let $R$ be a ring. If $R$ is not Jacobson there exist
a prime $\mathfrak p \subset R$, an element $f \in R$
such that the following hold
\begin{enumerate}
\item $\mathfrak p$ is not a maximal ideal,
\item $f \not \in \mathfrak p$,
\item $V(\mathfrak p) \cap D(f) = \{\mathfrak p\}$, and
\item $(R/\mathfrak p)_f$ is a field.
\end{enumerate}
On the other hand, if $R$ is Jacobson, then for any pair $(\mathfrak p, f)$
such that (1) and (2) hold, $V(\mathfrak p)\cap D(f)$ contains
infinitely many maximal ideals of $R$; in particular it is infinite.
\end{lemma}

\begin{proof}
Assume $R$ is not Jacobson.
By Lemma \ref{algebra-lemma-jacobson} this means there exists a
closed subset $T \subset \Spec(R)$
whose set $T_0 \subset T$ of closed points is not dense in $T$.
Choose an $f \in R$ such that $T_0 \subset V(f)$ but
$T \not \subset V(f)$. Note that $T \cap D(f)$
is homeomorphic to $\Spec((R/I)_f)$ if $T = V(I)$, see
Lemmas \ref{algebra-lemma-spec-closed} and \ref{algebra-lemma-standard-open}.
The set $T\cap D(f)$ is nonempty, so its displayed coordinate ring
is nonzero. Every nonzero ring has a maximal ideal
(Lemma \ref{editorial-source-lemma-Zariski-topology}), so choose a closed point $t$
of the space $T\cap D(f)$. Then $t$ corresponds to a prime ideal
$\mathfrak p \subset R$ which is not maximal (as $t \not \in T_0$).
Thus (1) holds. By construction $f \not \in \mathfrak p$, hence (2).
As $t$ is a closed point of $T \cap D(f)$ we see that
$V(\mathfrak p) \cap D(f) = \{\mathfrak p\}$, i.e., (3) holds. Hence we
conclude that $(R/\mathfrak p)_f$ is a domain whose
spectrum has one point, hence (4) holds
(for example combine Lemmas \ref{editorial-source-lemma-characterize-local-ring} and
\ref{algebra-lemma-minimal-prime-reduced-ring}).

\medskip\noindent
Conversely, suppose $R$ is Jacobson and $(\mathfrak p,f)$ satisfies
(1) and (2). Suppose only finitely many maximal ideals
$\mathfrak m_1,\ldots,\mathfrak m_t$ lie in
$V(\mathfrak p)\cap D(f)$. Because $\mathfrak p$ is not maximal,
choose $g_i\in\mathfrak m_i\setminus\mathfrak p$ and let
$g=\prod_{i=1}^t g_i$, with $g=1$ when $t=0$.
Primality gives $fg\notin\mathfrak p$, so the locally closed subset
$V(\mathfrak p)\cap D(fg)$ is nonempty. Since $\Spec(R)$ is Jacobson,
it contains a point closed in $\Spec(R)$. That point would be one of
the listed maximal ideals, since $D(fg)\subset D(f)$, but $g$ belongs
to every listed ideal, excluding each of them. This contradiction
proves that there are infinitely many such maximal ideals.
\end{proof}


\subsection{Editorial treatment of Algebra lines 6778--6807}

\hypertarget{algebra-context-057}{}

\noindent\href{https://github.com/stacks/stacks-project/blob/a04446e57ec1fbc252a871afcec7752fb2807b14/algebra.tex\#L6778}{Original source, lines 6778--6807}. Complete original and reviewed TeX passages accompany this note. Every intervention is separately recorded; the following treatment is editorial.


\begin{lemma}
\label{editorial-source-lemma-pid-jacobson}
The ring $\mathbf{Z}$ is a Jacobson ring.
More generally, let $R$ be a ring such that
\begin{enumerate}
\item $R$ is a domain,
\item $R$ is Noetherian,
\item any nonzero prime ideal is a maximal ideal, and
\item $R$ has infinitely many maximal ideals.
\end{enumerate}
Then $R$ is a Jacobson ring. In (2), it suffices more generally
that $\Spec(R)$ is a Noetherian topological space.
\end{lemma}

\begin{proof}
Assume (1), (3), (4) and that $\Spec(R)$ is Noetherian. The statement
means that $(0) = \bigcap_{\mathfrak m \subset R} \mathfrak m$.
Since $R$ has infinitely many maximal ideals it suffices to
show that any nonzero $x\in R$ belongs to only finitely many
maximal ideals, or equivalently that $V(x)$ is finite.
By Lemma \ref{algebra-lemma-spec-closed}
we see that $V(x)$ is homeomorphic to $\Spec(R/xR)$.
By assumption (3) every prime of $R/xR$ is minimal and hence
corresponds to an irreducible component of $\Spec(R/xR)$
(Lemma \ref{editorial-source-lemma-irreducible}).
If $R$ is Noetherian, then $\Spec(R)$ is Noetherian by
Lemma \ref{editorial-source-lemma-Noetherian-topology}. More generally assume only
this topological condition. Its closed subspace $V(x)$ is Noetherian,
and therefore has finitely many irreducible components by Topology,
Lemma \ref{topology-lemma-Noetherian}. Every prime in $V(x)$ is
nonzero and hence maximal by (3), so each irreducible component is
a singleton. Thus $V(x)$ is finite. The infinitely many maximal
ideals cannot all contain this $x$, proving their intersection is zero.
The nonzero primes are already maximal by (3), so every prime is an
intersection of maximal ideals and Lemma \ref{algebra-lemma-jacobson-prime}
proves that $R$ is Jacobson.
\end{proof}


\subsection{Editorial treatment of Algebra lines 6809--6827}

\hypertarget{algebra-context-058}{}

\noindent\href{https://github.com/stacks/stacks-project/blob/a04446e57ec1fbc252a871afcec7752fb2807b14/algebra.tex\#L6809}{Original source, lines 6809--6827}. Complete original and reviewed TeX passages accompany this note. Every intervention is separately recorded; the following treatment is editorial.


\begin{example}
\label{editorial-source-example-infinite-product-fields-jacobson}
Let $A$ be an infinite set.
For each $\alpha \in A$, let $k_\alpha$ be a field.
We claim that $R = \prod_{\alpha\in A} k_\alpha$ is Jacobson.
For the original element $f=(f_\alpha)_\alpha$, define
$$
e_\alpha=\begin{cases}1&f_\alpha\ne0,\\0&f_\alpha=0,\end{cases}
\qquad
u_\alpha=\begin{cases}f_\alpha&f_\alpha\ne0,\\1&f_\alpha=0.\end{cases}
$$
Then $e=(e_\alpha)$ is idempotent, $u=(u_\alpha)$ is a unit with
inverse $(u_\alpha^{-1})$, and $f=ue$ coordinate by coordinate.
Thus $D(f)=D(e)$ and the canonical maps identify
$R_f\cong R_e\cong R/(1-e)$. The first pair of inverse maps comes
from inverting $f=ue$ with $u$ already a unit. For the second, an
invertible idempotent is $1$, so $R\to R_e$ kills $1-e$; conversely
$e$ becomes $1$ in $R/(1-e)$. These universal maps are inverse
because their composites fix the image of $R$ and the adjoined inverse.

In fact every prime of this product ring is maximal. In the domain
$R/\mathfrak p$, the image of an idempotent is either $0$ or $1$.
For any $f\notin\mathfrak p$, its decomposition $f=ue$ forces the
image of $e$ to be $1$, so the image of $f$ is the unit given by
the image of $u$. Every nonzero element of $R/\mathfrak p$ is
therefore a unit. All primes are maximal, and each radical ideal
is the intersection of primes containing it, proving that $R$ is Jacobson.
This also covers finite index sets and the empty product, whose
spectrum is empty. More generally, the following original argument
proves the Jacobson property whenever each canonical localization
$R\to R_f$ is surjective.
Namely, say $\mathfrak p \subset R$ is a prime ideal
and $f \in R$, $f \not \in \mathfrak p$. We have to find
a maximal ideal $\mathfrak m$ of $R$ such that
$\mathfrak p \subset \mathfrak m$ and $f \not\in \mathfrak m$.
Because $R_f$ is a quotient of $R$ we see that any maximal
ideal of $R_f$ corresponds to a maximal ideal of $R$
not containing $f$. Hence the result follows
by choosing a maximal ideal of $R_f$ containing $\mathfrak p R_f$.
\end{example}


\subsection{Editorial treatment of Algebra lines 6829--6842}

\hypertarget{algebra-context-059}{}

\noindent\href{https://github.com/stacks/stacks-project/blob/a04446e57ec1fbc252a871afcec7752fb2807b14/algebra.tex\#L6829}{Original source, lines 6829--6842}. Complete original and reviewed TeX passages accompany this note. Every intervention is separately recorded; the following treatment is editorial.


\begin{example}
\label{editorial-source-example-not-jacobson}
A domain $R$ with finitely many maximal ideals
$\mathfrak m_i$, $i = 1, \ldots, n$ is not a
Jacobson ring, except when it is a field.
If this domain is not a field, every maximal ideal $\mathfrak m_i$
is nonzero. Choose $0\ne a_i\in\mathfrak m_i$ for $1\leq i\leq n$.
The nonzero domain has at least one maximal ideal, so $n\geq1$.
The product $a_1\cdots a_n$ is nonzero and belongs to
$$
\mathfrak m_1\cdots\mathfrak m_n
\subset\mathfrak m_1\cap\cdots\cap\mathfrak m_n.
$$
Hence the radical ideal $(0)$ is not the intersection of the maximal
ideals containing it, proving the assertion. In particular a discrete
valuation ring is not Jacobson. More generally, in any local ring
having a prime $\mathfrak p$ distinct from its unique maximal ideal
$\mathfrak m$, the intersection of maximal ideals containing the
radical ideal $\mathfrak p$ is $\mathfrak m\ne\mathfrak p$.
Thus any local ring with at least two prime ideals is not Jacobson,
without a domain hypothesis in this last assertion.
\end{example}


\subsection{Editorial treatment of Algebra lines 6872--6886}

\hypertarget{algebra-context-060}{}

\noindent\href{https://github.com/stacks/stacks-project/blob/a04446e57ec1fbc252a871afcec7752fb2807b14/algebra.tex\#L6872}{Original source, lines 6872--6886}. Complete original and reviewed TeX passages accompany this note. Every intervention is separately recorded; the following treatment is editorial.


\begin{lemma}
\label{editorial-source-lemma-dimension}
Suppose that $k$ is a field and suppose that $V$ is a nonzero vector
space over $k$. Assume the dimension of $V$ (which is a cardinal number)
is smaller than the cardinality of $k$. Then for any linear operator
$T : V \to V$ there exists some monic polynomial $P(t) \in k[t]$ such that
$P(T)$ is not invertible.
\end{lemma}

\begin{proof}
Suppose every monic polynomial $P$ gives an invertible $P(T)$.
For an arbitrary nonzero $Q\in k[t]$ with leading coefficient $c$,
retain the exact comparison
$Q(T)=c\,(c^{-1}Q)(T)$, where $c^{-1}Q$ is monic.
Its inverse is $((c^{-1}Q)(T))^{-1}c^{-1}$, so every such $Q(T)$
is invertible, with its scalar factor retained.
Define the action of the original rational-function field by
$$
\rho:k(t)\longrightarrow\operatorname{End}_k(V),\qquad
p/q\longmapsto p(T)q(T)^{-1}.
$$
All polynomial operators commute, and their inverses commute with
them and with each other. If $p/q=p'/q'$, then $pq'=p'q$;
evaluation at $T$ and multiplication by $q(T)^{-1}q'(T)^{-1}$
prove equality of the proposed operators. The common-denominator
formulas prove additivity and multiplicativity, and $\rho(1)$ is
the identity. This is a $k(t)$-vector-space structure extending
the original $k$-structure on $V$. For $0\ne v\in V$, the map
$k(t)\to V$, $a\mapsto\rho(a)v$, is injective: if $a\ne0$,
multiplication by $\rho(a^{-1})$ recovers $v$.

\medskip\noindent
The elements $1/(t-\lambda)$ for $\lambda\in k$ are linearly
independent over $k$. Indeed, multiplying a finite relation
$\sum_{i=1}^r c_i/(t-\lambda_i)=0$ for distinct $\lambda_i$ by
the full product $\prod_{i=1}^r(t-\lambda_i)$ gives
$\sum_{i=1}^r c_i\prod_{j\ne i}(t-\lambda_j)=0$.
Evaluation at $t=\lambda_i$ gives
$c_i\prod_{j\ne i}(\lambda_i-\lambda_j)=0$, and every factor
in that product is nonzero. Thus every $c_i=0$.
Their images in $V$ are linearly independent, so
$\dim_k(V)\geq\#k$, contradicting the hypothesis.
\end{proof}


\subsection{Editorial treatment of Algebra lines 6891--6943}

\hypertarget{algebra-context-061}{}

\noindent\href{https://github.com/stacks/stacks-project/blob/a04446e57ec1fbc252a871afcec7752fb2807b14/algebra.tex\#L6891}{Original source, lines 6891--6943}. Complete original and reviewed TeX passages accompany this note. Every intervention is separately recorded; the following treatment is editorial.


\begin{theorem}
\label{editorial-source-theorem-uncountable-nullstellensatz}
Let $k$ be a field. Let $S$ be a $k$-algebra generated over $k$
by the elements $\{x_i\}_{i \in I}$. Assume the cardinality of $I$
is smaller than the cardinality of $k$. Then
\begin{enumerate}
\item for all maximal ideals $\mathfrak m \subset S$ the field
extension $\kappa(\mathfrak m)/k$
is algebraic, and
\item $S$ is a Jacobson ring.
\end{enumerate}
If $I$ is finite, both conclusions hold for every field $k$, without
the cardinality restriction on the finite generating set.
\end{theorem}

\begin{proof}
For finite $I$, both assertions follow from the Hilbert
Nullstellensatz (Theorem \ref{editorial-source-theorem-nullstellensatz}) over every
field, regardless of the size of that finite set. Now assume $I$
is infinite and put $\kappa=\#I<\#k$.
Let $P=k[\{X_i\}_{i\in I}]$ be the polynomial ring mapping onto
the original algebra $S$ by $X_i\mapsto x_i$. Every monomial has
finite support and is specified by a finite list of pairs
$(i,e)$ with $i\in I$, $e\in\mathbf Z_{\geq0}$.
For each finite length $r$, the set of such lists has cardinality
at most $\kappa$, since $\kappa$ is infinite; the union over all
$r\geq0$ still has cardinality $\kappa$.
Distinct variables supply $\kappa$ distinct monomials, so the
monomial basis of $P$ has cardinality exactly $\kappa$.
The images of these monomials span $S$ and every quotient $S/\mathfrak m$.
In particular, for a maximal ideal $\mathfrak m\subset S$, its
residue field is exactly $S/\mathfrak m$ and has $k$-dimension
at most $\kappa$. These are maps from the same original polynomial
ring and algebra, without replacing $S$ by $P$.

\medskip\noindent
To arrive at a contradiction pick $T \in \kappa(\mathfrak m)$ transcendental
over $k$. Note that the $k$-linear map
$T : \kappa(\mathfrak m) \to \kappa(\mathfrak m)$
given by multiplication by $T$ has the property that $P(T)$ is invertible
for all monic polynomials $P(t) \in k[t]$.
Also, $\kappa(\mathfrak m)$ has dimension at most the cardinality of $I$
over $k$ since it is a quotient of the vector space
$k[\{x_i\}_{i \in I}]$ over $k$ (whose dimension is $\# I$ as we saw above).
This is impossible by Lemma \ref{editorial-source-lemma-dimension}.

\medskip\noindent
To show that $S$ is Jacobson we argue as follows. If not then
there exists a prime $\mathfrak q \subset S$ and an element $f \in S$,
$f \not \in \mathfrak q$ such that $\mathfrak q$ is not maximal
and $(S/\mathfrak q)_f$ is a field, see
Lemma \ref{editorial-source-lemma-characterize-jacobson}.
The field $L=(S/\mathfrak q)_f$ is generated over $k$ by the
images of the original $x_i$ and the inverse of the image of $f$.
Thus it has a generating set of cardinality at most
$\kappa+1=\kappa<\#k$. Apply the first assertion already proved
to the maximal ideal $(0)$ of the field $L$: the extension $L/k$
is algebraic. The canonical inclusion
$S/\mathfrak q\subset L\subset\operatorname{Frac}(S/\mathfrak q)$
identifies $L$ with the entire fraction field, since a field containing
$S/\mathfrak q$ contains the inverse of every nonzero element of
that domain. Hence $L=\kappa(\mathfrak q)$ canonically.
The prime $\mathfrak q$ lies over the maximal ideal $(0)$ of $k$,
so Lemma \ref{algebra-lemma-finite-residue-extension-closed} makes
$\mathfrak q$ maximal in $S$, a contradiction.
\end{proof}


\subsection{Editorial treatment of Algebra lines 6945--6963}

\hypertarget{algebra-context-062}{}

\noindent\href{https://github.com/stacks/stacks-project/blob/a04446e57ec1fbc252a871afcec7752fb2807b14/algebra.tex\#L6945}{Original source, lines 6945--6963}. Complete original and reviewed TeX passages accompany this note. Every intervention is separately recorded; the following treatment is editorial.


\begin{lemma}
\label{editorial-source-lemma-base-change-Jacobson}
Let $k$ be a field. Let $S$ be a $k$-algebra.
For any field extension $K/k$ whose cardinality is larger
than the cardinality of $S$ we have
\begin{enumerate}
\item for every maximal ideal $\mathfrak m$ of $S_K$ the field
$\kappa(\mathfrak m)$ is algebraic over $K$, and
\item $S_K$ is a Jacobson ring.
\end{enumerate}
It suffices more generally that $S$ have a $k$-algebra generating
set of cardinality less than $\#K$. If this set is finite, no
cardinality restriction on $K$ is needed.
\end{lemma}

\begin{proof}
Fix the given field extension $K/k$ and write
$S_K=K\otimes_kS$. If $(s_i)_{i\in I}$ generates $S$ as a
$k$-algebra, then $(1\otimes s_i)_{i\in I}$ generates $S_K$ as a
$K$-algebra. Indeed every tensor is a finite sum of terms
$\lambda\otimes s$, and an expression for $s$ as a polynomial
in finitely many $s_i$ gives the corresponding expression for
$\lambda\otimes s$, with all coefficients carried by $k\to K$.
Theorem \ref{editorial-source-theorem-uncountable-nullstellensatz} therefore applies
whenever $\#I<\#K$, and its finite-generator case applies without
that restriction. Taking the generating family indexed by every
element of $S$ gives the hypotheses and conclusion stated originally.
\end{proof}


\subsection{Editorial treatment of Algebra lines 6965--6978}

\hypertarget{algebra-context-063}{}

\noindent\href{https://github.com/stacks/stacks-project/blob/a04446e57ec1fbc252a871afcec7752fb2807b14/algebra.tex\#L6965}{Original source, lines 6965--6978}. Complete original and reviewed TeX passages accompany this note. Every intervention is separately recorded; the following treatment is editorial.


\begin{example}
\label{editorial-source-example-countable-trick-does-not-work}
The trick in the proof of
Theorem \ref{editorial-source-theorem-uncountable-nullstellensatz}
really does not work if $k$ is a countable field and $I$ is countable too.
Let $k$ be a countable field. Let $x$ be a variable,
and let $k(x)$ be the field of rational functions in $x$.
Consider the polynomial algebra $R = k[x, \{x_f\}_{f \in k[x]-\{0\}}]$.
Let $I = (\{fx_f - 1\}_{f\in k[x] - \{0\}})$. Note that
Evaluation defines a map
$$
R\longrightarrow k(x),\qquad x\longmapsto x,\qquad
x_f\longmapsto f^{-1}\quad(0\ne f\in k[x]),
$$
which kills every original relation $fx_f-1$ and factors through
$R/I$. Conversely the map $k[x]\to R/I$ sends every nonzero
$f$ to a unit, with inverse the class of $x_f$. The localization
property therefore gives
$$
k(x)\longrightarrow R/I,\qquad p/q\longmapsto
\overline p\,\overline{x}_q\quad(q\ne0).
$$
These maps are inverse. On $k(x)$ their composite fixes all fractions;
on $R/I$ it fixes $\overline x$ and every $\overline{x}_f$, since
each is the unique inverse of $\overline f$. Thus $R/I\cong k(x)$
as $k$-algebras, so $I$ itself is maximal and any maximal ideal
containing $I$ is equal to $I$. This extension is transcendental.

The same presentation and inverse maps work over every field $k$.
For countable $k$, the polynomial set $k[x]$ is countably infinite,
so the original algebra has a countable generating family.
For infinite $k$, that family has cardinality exactly $\#k$, since
$\#k[x]=\max(\#k,\aleph_0)=\#k$ and the nonzero constants already
give $\#k$ distinct indices. Thus equality at the cardinal threshold
can fail to force algebraic residue extensions over every infinite
field, with the same explicit quotient.
\end{example}


\subsection{Editorial treatment of Algebra lines 7050--7087}

\hypertarget{algebra-context-064}{}

\noindent\href{https://github.com/stacks/stacks-project/blob/a04446e57ec1fbc252a871afcec7752fb2807b14/algebra.tex\#L7050}{Original source, lines 7050--7087}. Complete original and reviewed TeX passages accompany this note. Every intervention is separately recorded; the following treatment is editorial.


\begin{proposition}
\label{editorial-source-proposition-Jacobson-permanence}
Let $R$ be a Jacobson ring. Let $R \to S$ be a
ring map of finite type. Then
\begin{enumerate}
\item The ring $S$ is Jacobson.
\item The map $\Spec(S) \to \Spec(R)$ transforms
closed points to closed points.
\item For $\mathfrak m' \subset S$ maximal lying over $\mathfrak m \subset R$
the field extension $\kappa(\mathfrak m')/\kappa(\mathfrak m)$
is finite.
\end{enumerate}
\end{proposition}

\begin{proof}
Let $\mathfrak m' \subset S$ be a maximal ideal and
let $\varphi:R\to S$ be the original map and set
$\mathfrak m=\varphi^{-1}(\mathfrak m')$.
Then $R/\mathfrak m \to S/\mathfrak m'$ satisfies
the conditions of Lemma \ref{algebra-lemma-silly-jacobson}
by Lemma \ref{algebra-lemma-Jacobson-mod-ideal}.
Hence $R/\mathfrak m$ is a field and
$\mathfrak m$ a maximal ideal and the induced
residue field extension is finite. This proves (2) and (3).  

\medskip\noindent
If $S$ is not Jacobson, then by Lemma \ref{editorial-source-lemma-characterize-jacobson} there
exists a non-maximal prime ideal $\mathfrak q$ of $S$ and an
$g \in S$, $g \not\in \mathfrak q$ such that $(S/\mathfrak q)_g$ is a field.
To arrive at a contradiction we show that $\mathfrak q$ is a maximal ideal.
Let $\mathfrak p=\varphi^{-1}(\mathfrak q)$.
The map $R/\mathfrak p\to S/\mathfrak q$ is injective and of finite
type, using the images of the original finite set of algebra generators.
Adjoining the inverse of the image of $g$ adds one generator, so
$R/\mathfrak p\to(S/\mathfrak q)_g$ is also of finite type.
The target is a field and embeds canonically in
$\operatorname{Frac}(S/\mathfrak q)$, so it equals that fraction
field: it already contains the domain and all inverses of its nonzero
elements. Thus
$R/\mathfrak p \to (S/\mathfrak q)_g$ satisfies the conditions of
Lemma \ref{algebra-lemma-silly-jacobson} by
Lemma \ref{algebra-lemma-Jacobson-mod-ideal}.
Hence $R/\mathfrak p$ is a field and the field extension
$\kappa(\mathfrak p) \to (S/\mathfrak q)_g = \kappa(\mathfrak q)$ is
finite, thus algebraic. Then $\mathfrak q$ is a maximal ideal of $S$ by
Lemma \ref{algebra-lemma-finite-residue-extension-closed}. Contradiction.
\end{proof}


\subsection{Editorial treatment of Algebra lines 7099--7173}

\hypertarget{algebra-context-065}{}

\noindent\href{https://github.com/stacks/stacks-project/blob/a04446e57ec1fbc252a871afcec7752fb2807b14/algebra.tex\#L7099}{Original source, lines 7099--7173}. Complete original and reviewed TeX passages accompany this note. Every intervention is separately recorded; the following treatment is editorial.


\begin{lemma}
\label{editorial-source-lemma-image-finite-type-map-Jacobson-rings}
Let $R \to S$ be a finite type ring map of Jacobson rings.
Denote $X = \Spec(R)$ and $Y = \Spec(S)$.
Write $f : Y \to X$ for the induced
map of spectra. Let $E \subset Y = \Spec(S)$ be a
constructible set. Denote with a subscript ${}_0$ the set
of closed points of a topological space.
\begin{enumerate}
\item We have $f(E)_0 = f(E_0) = X_0 \cap f(E)$.
\item A point $\xi \in X$ is in $f(E)$ if and only if
$\overline{\{\xi\}} \cap f(E_0)$ is dense in $\overline{\{\xi\}}$.
\end{enumerate}
\end{lemma}

\begin{proof}
We have a commutative diagram of continuous maps
$$
\xymatrix{
E \ar[r] \ar[d] & Y \ar[d] \\
f(E) \ar[r] & X
}
$$
Suppose $x \in f(E)$ is closed in $f(E)$. Then $f^{-1}(\{x\})\cap E$
is nonempty and closed in $E$. Applying
Topology, Lemma \ref{topology-lemma-jacobson-inherited}
to both inclusions
$$
f^{-1}(\{x\}) \cap E \subset E \subset Y
$$
we find there exists a point $y \in f^{-1}(\{x\}) \cap E$ which is
closed in $Y$. In other words, there exists $y \in Y_0$ and $y \in E_0$
mapping to $x$. Hence $x \in f(E_0)$.
This proves that $f(E)_0 \subset f(E_0)$.
Proposition \ref{editorial-source-proposition-Jacobson-permanence} implies that
$f(E_0) \subset X_0 \cap f(E)$. The inclusion
$X_0 \cap f(E) \subset f(E)_0$ is trivial. This proves the
first assertion.

\medskip\noindent
Suppose that $\xi \in f(E)$. According to
Lemma \ref{editorial-source-lemma-characterize-image-finite-type}
the set $f(E) \cap \overline{\{\xi\}}$ contains a dense
open subset of $\overline{\{\xi\}}$. Since $X$ is Jacobson
we conclude that $f(E) \cap \overline{\{\xi\}}$ contains a
dense set of closed points, see Topology,
Lemma \ref{topology-lemma-jacobson-inherited}.
We conclude by part (1) of the lemma.

\medskip\noindent
Conversely, suppose $\overline{\{\xi\}}\cap f(E_0)$ is dense
in $\overline{\{\xi\}}$. By Lemma \ref{editorial-source-lemma-constructible-is-image}
choose a finitely presented map $S\to S'$ whose spectrum map
$g:Y'=\Spec(S')\to Y$ has image exactly $E$.
The ring $S'$ is Jacobson by Proposition \ref{editorial-source-proposition-Jacobson-permanence}.
Part (1) of the present lemma, already proved and applied to $S\to S'$,
gives $g(Y'_0)=E_0$. Put $h=f\circ g$, so
$h(Y'_0)=f(E_0)$ and $h(Y')=f(E)$, retaining the original $S$ and $E$.
If $\xi$ corresponds to $\mathfrak p\subset R$, use the exact diagram
$$
\xymatrix{
S'\ar[r] & S'/\mathfrak pS'\\
R\ar[r]\ar[u] & R/\mathfrak p\ar[u]
}
$$
Every point of $Y'_0$ whose image lies in $V(\mathfrak p)$ corresponds
to a prime of $S'$ containing $\mathfrak pS'$, hence to a point of
$\Spec(S'/\mathfrak pS')$. The spectrum image of
$R/\mathfrak p\to S'/\mathfrak pS'$ therefore contains the given
dense subset $h(Y'_0)\cap V(\mathfrak p)$ of
$V(\mathfrak p)=\Spec(R/\mathfrak p)$.
By Lemma \ref{algebra-lemma-domain-image-dense-set-points-generic-point},
there is a prime $\overline{\mathfrak q}'$ of $S'/\mathfrak pS'$
contracting to $(0)$ in $R/\mathfrak p$. Its inverse image
$\mathfrak q'\subset S'$ contracts to $\mathfrak p$ in $R$.
The point $g(\mathfrak q')$ belongs to $E$ and its image under the
original $f$ is $\xi$. Consequently $\xi\in f(E)$, as required.
\end{proof}


\subsection{Editorial treatment of Algebra lines 7189--7215}

\hypertarget{algebra-context-066}{}

\noindent\href{https://github.com/stacks/stacks-project/blob/a04446e57ec1fbc252a871afcec7752fb2807b14/algebra.tex\#L7189}{Original source, lines 7189--7215}. Complete original and reviewed TeX passages accompany this note. Every intervention is separately recorded; the following treatment is editorial.


\begin{lemma}
\label{editorial-source-lemma-conclude-jacobson-Noetherian}
With notation as above, assume that $R$ is a Noetherian Jacobson ring.
Further assume $R \to S$ is of finite type.
There is a commutative diagram
$$
\xymatrix{
\text{Constr}(Y) \ar[r]^{E \mapsto E_0} \ar[d]^{E \mapsto f(E)} &
\text{Constr}(Y_0) \ar[d]^{E \mapsto f(E)} \\
\text{Constr}(X) \ar[r]^{E \mapsto E_0} &
\text{Constr}(X_0)
}
$$
where the horizontal arrows are the bijections from
Topology, Lemma \ref{topology-lemma-jacobson-equivalent-constructible}.
The same diagram and conclusions hold more generally when $R$ is
Jacobson and $R\to S$ is of finite presentation, without requiring
$R$ to be Noetherian.
\end{lemma}

\begin{proof}
Under the original hypotheses, the finite type map $R\to S$ is of
finite presentation by
Lemma \ref{editorial-source-lemma-Noetherian-finite-type-is-finite-presentation}.
For the more general assertion, finite presentation is assumed directly.
In both cases it implies finite type, so $S$ is Jacobson by
Proposition \ref{editorial-source-proposition-Jacobson-permanence}. Therefore $X$ and
$Y$ are Jacobson spaces and the two horizontal maps are bijections
by Topology, Lemma \ref{topology-lemma-jacobson-equivalent-constructible},
which requires no Noetherian hypothesis.
Chevalley's theorem (Theorem \ref{editorial-source-theorem-chevalley}) states that
the image of a constructible subset of the source $Y$ is constructible
in the target $X$. For each such $E\subset Y$,
Lemma \ref{editorial-source-lemma-image-finite-type-map-Jacobson-rings} gives
$f(E_0)=f(E)_0$, proving commutativity.
The right vertical arrow is well defined as a direct-image map:
given a constructible subset $C\subset Y_0$, the top horizontal
bijection supplies a unique constructible $E\subset Y$ with $E_0=C$,
and then $f(C)=f(E)_0$ is constructible in $X_0$ by the bottom
horizontal bijection. This proves every arrow and the square in
the stated greater generality.
\end{proof}


\subsection{Editorial treatment of Algebra lines 7220--7361}

\hypertarget{algebra-context-067}{}

\noindent\href{https://github.com/stacks/stacks-project/blob/a04446e57ec1fbc252a871afcec7752fb2807b14/algebra.tex\#L7220}{Original source, lines 7220--7361}. Complete original and reviewed TeX passages accompany this note. Every intervention is separately recorded; the following treatment is editorial.


\begin{example}
\label{editorial-source-example-product-matrices-zero}
Let $k$ be a field. The space $\Spec(k[x, y]/(xy))$
has two irreducible components: namely the $x$-axis and the
$y$-axis. As a generalization, let
$$
R = k[x_{11}, x_{12}, x_{21}, x_{22}, y_{11}, y_{12}, y_{21}, y_{22}]/
\mathfrak a,
$$
where $\mathfrak a$ is the ideal in
$k[x_{11}, x_{12}, x_{21}, x_{22}, y_{11}, y_{12}, y_{21}, y_{22}]$
generated by the entries of the $2 \times 2$ product matrix
$$
\left(
\begin{matrix}
x_{11} & x_{12}\\
x_{21} & x_{22}
\end{matrix}
\right)
 \left(
\begin{matrix}
y_{11} & y_{12}\\
y_{21} & y_{22}
\end{matrix}
\right).
$$
In this example we will describe $\Spec(R)$.

\medskip\noindent
To prove the statement about $\Spec(k[x, y]/(xy))$ we argue as follows.
If $\mathfrak p\subset k[x,y]$ is any prime ideal containing $xy$,
then either $x$ or $y$ belongs to $\mathfrak p$. Hence the minimal such
prime ideals are just $(x)$ and $(y)$. In case $k$ is
algebraically closed, the $\text{max-Spec}$ of these components
can then be visualized as the point sets of $y$- and $x$-axis.

\medskip\noindent
For the generalization, note that we may identify the closed
points of the spectrum of
$k[x_{11}, x_{12}, x_{21}, x_{22}, y_{11}, y_{12}, y_{21}, y_{22}]$
with the space of matrices
$$
\left\{ (X, Y) \in \text{Mat}(2, k)\times \text{Mat}(2, k) \mid
X = \left(
\begin{matrix}
x_{11} & x_{12}\\
x_{21} & x_{22}
\end{matrix}
\right),
Y= \left(
\begin{matrix}
y_{11} & y_{12}\\
y_{21} & y_{22}
\end{matrix}
\right)
\right\}
$$
if $k$ is algebraically closed. We assume this for the following
description in terms of closed $k$-points. The coordinate-ring proof
below establishes the component statement over every field $k$.
Now define a group action of
$\text{GL}(2, k)\times \text{GL}(2, k)\times \text{GL}(2, k)$
on the space of matrices $\{(X, Y)\}$ by
$$
(g_1, g_2, g_3) \times (X, Y) \mapsto ((g_1Xg_2^{-1}, g_2Yg_3^{-1})).
$$
Here, also observe that the algebraic set
$$
\text{GL}(2, k)\times \text{GL}(2, k)\times \text{GL}(2, k) \subset
\text{Mat}(2, k)\times \text{Mat}(2, k) \times \text{Mat}(2, k)
$$
is irreducible since it is the max spectrum of the domain
$$
k[x_{11}, x_{12}, \ldots, z_{21}, z_{22}, (x_{11}x_{22}-x_{12}x_{21})^{-1}
, (y_{11}y_{22}-y_{12}y_{21})^{-1}, (z_{11}z_{22}-z_{12}z_{21})^{-1}].
$$
Since the image of an irreducible algebraic set is still
irreducible, the orbit closures below are irreducible. There are
six orbits, indexed by the pairs of ranks $(r,s)$ with $r+s\leq2$:
$$
(2,0),\quad(1,1),\quad(0,2),\quad(1,0),\quad(0,1),\quad(0,0).
$$
Indeed $XY=0$ says $\operatorname{Im}(Y)\subset\ker(X)$.
Choose a basis of the middle space starting with a basis of
$\operatorname{Im}(Y)$, extend it to $\ker(X)$ and then to the
whole middle space. Lift a basis of $\operatorname{Im}(Y)$ through
$Y$ and extend those lifts by a basis of $\ker(Y)$ in the source.
In the target extend the images under $X$ of the remaining middle
basis vectors to a basis. These choices show that the ranks determine
the orbit. If their basis matrices in the original coordinates are
$C_1,C_2,C_3$, respectively, then
$g_1=C_1^{-1}$, $g_2=C_2^{-1}$, $g_3=C_3^{-1}$ give exactly
the displayed action, with transformed matrices
$C_1^{-1}XC_2$ and $C_2^{-1}YC_3$. It preserves the original
equations because the transformed product is $g_1XYg_3^{-1}$,
and the inverse transformation uses $(g_1^{-1},g_2^{-1},g_3^{-1})$.
The three cases according to the rank of $X$ are:
\begin{enumerate}
\item $\exists (g_1, g_2)$ such that $g_1Xg_2^{-1} = I_{2\times 2}$.
Then $Y$ is necessarily $0$, which as an algebraic set is
invariant under the group action. It follows that this orbit is
contained in the irreducible algebraic set defined by the prime
ideal $(y_{11}, y_{12}, y_{21}, y_{22})$. Taking the closure, we see
that $(y_{11}, y_{12}, y_{21}, y_{22})$ is actually a component.
\item $\exists (g_1, g_2)$ such that
$$
g_1Xg_2^{-1} = \left(
\begin{matrix}
1 & 0 \\
0 & 0
\end{matrix}
\right).
$$
This case occurs if and only if $X$ is a rank 1 matrix,
and furthermore, $Y$ is killed by such an $X$ if and only if
$$
x_{11}y_{11}+x_{12}y_{21} = 0; \quad x_{11}y_{12}+x_{12}y_{22} = 0;
$$
$$
x_{21}y_{11}+x_{22}y_{21} = 0; \quad x_{21}y_{12}+x_{22}y_{22} = 0.
$$
Fix a rank $1$ matrix $X$. The nonzero matrices $Y$ satisfying these
equations form an irreducible locally closed set, for the following
reason (the pair $(X,0)$ belongs to the closure of the first orbit):
$0 = g_1Xg_2^{-1}g_2Y$ implies that
$$
g_2Y = \left(
\begin{matrix}
0 & 0 \\
y_{21}' & y_{22}'
\end{matrix}
\right).
$$
With a further $\text{GL}(2, k)$-action on the right by $g_3$,
$g_2Y$ can be brought into
$$
g_2Yg_3^{-1} = \left(
\begin{matrix}
0 & 0 \\
0 & 1
\end{matrix}
\right),
$$
and thus such $Y$'s form an irreducible algebraic set
isomorphic to the image of $\text{GL}(2, k)$ under this action. Finally,
notice that the ``rank 1" condition for $X$'s forms an open dense
subset of the irreducible algebraic set
$\det X = x_{11}x_{22} - x_{12}x_{21} = 0$.
On the open set where both matrices are nonzero, the four displayed
product equations and $\det X=0$ describe the rank-$(1,1)$ orbit.
To identify its closure, retain the polynomial ring
$B=k[x_{11},x_{12},x_{21},x_{22},y_{11},y_{12},y_{21},y_{22}]$
and put
$$
\begin{aligned}
\mathfrak q_M=(
&x_{11}y_{11}+x_{12}y_{21},\ x_{11}y_{12}+x_{12}y_{22},\\
&x_{21}y_{11}+x_{22}y_{21},\ x_{21}y_{12}+x_{22}y_{22},\\
&x_{11}x_{22}-x_{12}x_{21},\ y_{11}y_{22}-y_{12}y_{21}).
\end{aligned}
$$
The proof below shows that this is prime and cuts out that closure.
The sixth equation is already a consequence of the first four on
the open set where $X$ is nonzero: writing $a_{ij}=(XY)_{ij}$,
the exact matrix identity $X\det(Y)=(XY)\operatorname{adj}(Y)$ gives
$$
\begin{aligned}
x_{11}\det Y&=a_{11}y_{22}-a_{12}y_{21},&
x_{12}\det Y&=-a_{11}y_{12}+a_{12}y_{11},\\
x_{21}\det Y&=a_{21}y_{22}-a_{22}y_{21},&
x_{22}\det Y&=-a_{21}y_{12}+a_{22}y_{11}.
\end{aligned}
$$
Thus on each standard open where an entry of $X$ is invertible, the
original five-equation ideal and $\mathfrak q_M$ give the same quotient.
No sixth equation is discarded in the closure.
\item $\exists (g_1, g_2)$ such that $g_1Xg_2^{-1} = 0$, or
equivalently, $X = 0$. Then $Y$ can be arbitrary and this component
is thus defined by $(x_{11}, x_{12}, x_{21}, x_{22})$.
\end{enumerate}

\medskip\noindent
Here is the coordinate-ring proof, valid over the original arbitrary
field $k$. Keep $B$ and the four-entry ideal $\mathfrak a$ of the
original product $XY$, so $R=B/\mathfrak a$. Put
$$
\mathfrak q_X=(x_{11},x_{12},x_{21},x_{22}),\qquad
\mathfrak q_Y=(y_{11},y_{12},y_{21},y_{22}),
$$
and retain $\mathfrak q_M=\mathfrak a+(\det X,\det Y)$ as displayed
above. The first two ideals are prime, since their quotient rings
are polynomial rings in the four remaining entries.

For the middle ideal use the exact rearrangement of the original entries
$$
M=\begin{pmatrix}
x_{11}&x_{21}&y_{21}&y_{22}\\
x_{12}&x_{22}&-y_{11}&-y_{12}
\end{pmatrix}.
$$
Its six $2\times2$ minors, in column order
$(12),(13),(14),(23),(24),(34)$, are
$$
\det X,\quad-a_{11},\quad-a_{12},\quad-a_{21},\quad-a_{22},\quad\det Y.
$$
The entries of $M$ are eight independent polynomial coordinates:
the inverse rearrangement recovers each original entry, including
$y_{11}=-M_{23}$ and $y_{12}=-M_{24}$. Define
$$
\Phi:B\longrightarrow k[\lambda,\mu,c_1,c_2,c_3,c_4]
$$
by
$$
\begin{aligned}
x_{11}&\mapsto\lambda c_1,&x_{12}&\mapsto\mu c_1,&
x_{21}&\mapsto\lambda c_2,&x_{22}&\mapsto\mu c_2,\\
y_{11}&\mapsto-\mu c_3,&y_{12}&\mapsto-\mu c_4,&
y_{21}&\mapsto\lambda c_3,&y_{22}&\mapsto\lambda c_4.
\end{aligned}
$$
Thus the two rows of $M$ map to $(\lambda c_i)_i$ and
$(\mu c_i)_i$. Every minor maps to zero, proving
$\mathfrak q_M\subset\ker\Phi$.

To prove equality, write $s_i=M_{1i}$ and $t_i=M_{2i}$ for
$1\leq i\leq4$. The image of a monomial
$\prod_i s_i^{u_i}t_i^{v_i}$ is
$\lambda^{\sum_i u_i}\mu^{\sum_i v_i}\prod_i c_i^{u_i+v_i}$.
Two monomials have the same image exactly when all their column
totals $u_i+v_i$ agree and their first-row totals $\sum_i u_i$ agree.
For two such monomials with first-row exponents $u_i,u'_i$, if
$u\ne u'$ choose $i,j$ with $u_i>u'_i$ and $u_j<u'_j$.
Then $s_i$ and $t_j$ divide the first monomial. Replacing the factor
$s_it_j$ by $t_is_j$ changes it by a multiple of the minor
$s_it_j-t_is_j$, preserves all row and column totals, and decreases
$\sum_i|u_i-u'_i|$ by $2$. Repetition reaches the second monomial,
proving their difference belongs to the minor ideal.
For an arbitrary polynomial in $\ker\Phi$, group its finitely many
terms according to their image monomial. Distinct image monomials
are linearly independent over $k$, so the sum of coefficients in
each group is zero. Choosing one monomial in each group expresses
that whole group as a $k$-linear combination of the differences
just proved to lie in the minor ideal. Consequently
$\ker\Phi=\mathfrak q_M$. The induced map
$B/\mathfrak q_M\hookrightarrow k[\lambda,\mu,c_1,c_2,c_3,c_4]$
is injective, so $\mathfrak q_M$ is prime over every field,
with all original signs retained also in characteristic $2$.

Now let $\mathfrak p\subset B$ be any prime containing
$\mathfrak a$. In the field $\operatorname{Frac}(B/\mathfrak p)$,
the original matrices satisfy $XY=0$. If $\det X\notin\mathfrak p$,
then $X$ is invertible over this field and $Y=0$, so
$\mathfrak q_Y\subset\mathfrak p$. If
$\det Y\notin\mathfrak p$, the same argument gives
$\mathfrak q_X\subset\mathfrak p$. If both determinants belong
to $\mathfrak p$, then $\mathfrak q_M\subset\mathfrak p$.
These three prime ideals are pairwise incomparable.
For example $x_{11}\in\mathfrak q_X\setminus\mathfrak q_M$
since $\Phi(x_{11})=\lambda c_1\ne0$, while
$\det Y\in\mathfrak q_M\setminus\mathfrak q_X$ since it remains
a nonzero polynomial in the quotient by $\mathfrak q_X$.
The analogous facts with $X,Y$ exchanged prove incomparability
with $\mathfrak q_Y$, and the two coordinate ideals are
incomparable by their distinct sets of variables.
Therefore these are exactly the three minimal primes over
$\mathfrak a$. Their images in $R$ define exactly the three
irreducible components of $\Spec(R)$.

Finally, over an algebraically closed field the rank-$(2,0)$ orbit
is the nonempty open $\det X\ne0$ in $V(\mathfrak q_Y)$ and is
dense there; the rank-$(0,2)$ orbit is dense in $V(\mathfrak q_X)$
by the identical argument. In the irreducible $V(\mathfrak q_M)$,
the open where both $X$ and $Y$ are nonzero is precisely the
rank-$(1,1)$ orbit. It is nonempty, as witnessed by
$$
X=\begin{pmatrix}1&0\\0&0\end{pmatrix},\qquad
Y=\begin{pmatrix}0&0\\0&1\end{pmatrix},
$$
and hence dense. Density of closed points in these affine finite
type spectra follows from the Jacobson results above, so these
are the same closures in the stated description by closed points.
The other three rank orbits lie in these components and give no
additional irreducible component.
\end{example}


\subsection{Editorial treatment of Algebra lines 7364--7396}

\hypertarget{algebra-context-068}{}

\noindent\href{https://github.com/stacks/stacks-project/blob/a04446e57ec1fbc252a871afcec7752fb2807b14/algebra.tex\#L7364}{Original source, lines 7364--7396}. Complete original and reviewed TeX passages accompany this note. Every intervention is separately recorded; the following treatment is editorial.


\begin{example}
\label{editorial-source-example-idempotent-matrices}
For another example, consider
$R = k[\{t_{ij}\}_{i, j = 1}^{n}]/\mathfrak a$, where $\mathfrak a$ is
the ideal generated by the entries of the product matrix $T^2-T$,
$T = (t_{ij})$. From linear algebra, we know that under the
$GL(n, k)$-action defined by $g, T \mapsto gTg^{-1}$, $T$ is
classified by its rank and each $T$ is conjugate to some
$\text{diag}(1, \ldots, 1, 0, \ldots, 0)$, which has $r$ 1's and $n-r$ 0's.
Every idempotent matrix with entries in $k$ belongs to exactly one
of these rank orbits. Indeed each vector has the decomposition
$v=Tv+(v-Tv)$ into $\operatorname{Im}(T)$ and $\ker(T)$;
their intersection is zero since $Tw=w$ on the image and $Tw=0$
on the kernel. Bases of these two subspaces give the stated
conjugacy, and conjugation preserves rank. To deduce the irreducible
components of $\operatorname{Spec}(R)$, including its nonclosed
points, we give a coordinate-ring proof below. It proves that the
rank loci are disjoint open and closed irreducible components over
every field $k$. To separate ranks in characteristic zero one may use
$f(t_{ij}) = trace(T) = \sum_{i = 1}^nt_{ii}$. In positive
characteristic cases, things are slightly more tricky since we
might have $trace(T) = 0$ even if $T \neq 0$. For instance, if $k$ has characteristic $3$,
$$
trace\left(
\begin{matrix}
1 & & \\
& 1 & \\
& & 1
\end{matrix}
\right) = 3 = 0
$$
For every $0\leq j\leq n$ put
$e_j(T)=\operatorname{tr}(\bigwedge^jT)$, with $e_0(T)=1$.
On a rank-$r$ idempotent this equals $\binom rj$ in $k$:
in a basis of the image followed by the kernel, the exterior map
has eigenvalue $1$ on the $j$-fold wedges of image basis vectors
and $0$ on all other basis wedges. The exterior basis change is
invertible, so it preserves the trace of the original exterior map.
Thus if $r<s$, $e_s$ is zero on rank $r$ and one on rank $s$.
The complete tuple $(e_1,\ldots,e_n)$ separates all ranks in every
characteristic. A single chosen entry need not do so: in
characteristic $3$, ranks $3$ and $4$ both give $e_3=1$, since
$\binom33=1$ and $\binom43=4=1$ in $k$.
The displayed function $e_3$ isolates rank $3$ when $n=3$;
it does not isolate that rank for arbitrary $n$.
\medskip\noindent
We now prove the component assertions in the original ring $R$,
without assuming that $k$ is algebraically closed or that $R$ is
reduced. Keep the original matrix $T$ and define its coefficients
$d_0,\ldots,d_n\in R$ by
$$
F(z)=\det(I_n-T+zT)=\sum_{r=0}^n d_rz^r.
$$
The identity $T^2=T$ gives, in matrices over $R[x,y]$,
$$
(I_n-T+xT)(I_n-T+yT)=I_n-T+xyT.
$$
Hence $F(x)F(y)=F(xy)$. Comparing the coefficient of each
$x^ry^s$ proves $d_r^2=d_r$ and $d_rd_s=0$ for $r\ne s$;
evaluation at $z=1$ gives $\sum_r d_r=1$.
Put $R_r=R/(1-d_r)$ and $X_r=\operatorname{Spec}(R_r)$.
Inverting the idempotent $d_r$ is the same as imposing $d_r=1$,
so $X_r=D(d_r)=V(1-d_r)$ in $\operatorname{Spec}(R)$.
In particular these loci are open and closed. The exact ring maps
$$
R\longrightarrow\prod_{r=0}^n R_r,
\qquad a\longmapsto(a\bmod(1-d_r))_r,
\qquad
(\overline{a_r})_r\longmapsto\sum_{r=0}^n d_ra_r
$$
are inverse: the second expression does not depend on lifts since
$d_r(1-d_r)=0$, and orthogonality and the sum identity verify both
composites on every element.

For any prime $\mathfrak p\subset R$, apply the image/kernel
decomposition above to the original matrix over $\kappa(\mathfrak p)$.
If its rank is $r$, its determinant polynomial is exactly $z^r$.
Thus $d_s$ has residue $1$ for $s=r$ and $0$ for $s\ne r$.
It follows that $X_r$ is precisely the residue-field rank-$r$
locus, including all nonclosed points. Each $X_r$ is nonempty:
evaluation at $P_r=\operatorname{diag}(1,\ldots,1,0,\ldots,0)$,
with exactly $r$ ones, is a $k$-point.

For irreducibility take
$A=k[g_{ij}:1\leq i,j\leq n][\det(G)^{-1}]$,
where $G=(g_{ij})$. This is an integral domain. The formula
$T\mapsto GP_rG^{-1}$ gives the coordinate-ring homomorphism
$$
R_r\longrightarrow A,\qquad
t_{ij}\longmapsto\sum_{a=1}^r g_{ia}(G^{-1})_{aj}.
$$
The inverse entries are in $A$ by the adjugate formula; the
substituted matrix is idempotent and has determinant polynomial
$z^r$, so the displayed homomorphism is well defined. Its spectrum
map $\operatorname{Spec}(A)\to X_r$ is surjective. In fact for
any $\mathfrak p\in X_r$, bases of the image and kernel over
$\kappa(\mathfrak p)$ provide an invertible matrix $G_{\mathfrak p}$
with $T_{\mathfrak p}=G_{\mathfrak p}P_rG_{\mathfrak p}^{-1}$.
Evaluation $A\to\kappa(\mathfrak p)$ at these entries has a
prime kernel whose contraction to $R_r$ is the given prime,
because its composite sends every original $t_{ij}$ to its
residue at $\mathfrak p$.
The continuous image of an irreducible space is irreducible:
a union of two relatively closed proper subsets would pull back
to two closed subsets covering the source, and surjectivity makes
both proper. Therefore each $X_r$ is irreducible. An irreducible
subset cannot meet two members of the disjoint open-and-closed
cover $(X_r)$, since that would separate it into two nonempty
relatively closed subsets. The same separation excludes a connected
subset from meeting two members. Each $X_r$ is itself irreducible
and connected. Hence the irreducible components and the connected
components of $\operatorname{Spec}(R)$ are exactly $X_0,\ldots,X_n$.

The same original coordinates give useful explicit charts.
For a subset $I\subseteq\{1,\ldots,n\}$ let $P_I$ be the diagonal
matrix with entry $1$ at indices in $I$ and $0$ elsewhere, and set
$$
Q_I=TP_I+(I_n-T)(I_n-P_I),\qquad q_I=\det Q_I.
$$
Multilinearity in the columns of $F(z)$, retaining their original
order, gives
$$
d_r=\sum_{|I|=r}q_I:
$$
the term indexed by $I$ chooses the column $Te_j$ for $j\in I$
and $(I_n-T)e_j$ for $j\notin I$.
Thus $X_r$ is covered by the opens $D(q_I)$ with $|I|=r$.
To check that each such open lies in $X_r$, compute
$TQ_I=Q_IP_I$. After inverting $q_I$ this gives
$T=Q_IP_IQ_I^{-1}$, so $F(z)=z^r$ there.

Write $K=Q_I^{-1}$ in this localization. From $KT=P_IK$,
the definition of $Q_I$ gives the exact identity
$$
I_n=KQ_I=P_IKP_I+(I_n-P_I)K(I_n-P_I).
$$
Consequently the two diagonal blocks of $K$, indexed by $I$ and
its complement without changing the original row and column labels,
are identity matrices. Introduce variables $u_{ab}$ for the ordered
pairs $(a,b)$ with exactly one of $a,b$ in $I$. Let $K_I(u)$ have
these off-diagonal entries and the two identity diagonal blocks.
Put $h_I(u)=\det K_I(u)$ and
$$
B_I=k[u_{ab}: \text{exactly one of }a,b\text{ belongs to }I]
[h_I(u)^{-1}].
$$
There are $2r(n-r)$ such variables and $h_I(0)=1$.
The matrix $T'=K_I(u)^{-1}P_IK_I(u)$ is idempotent, and
$$
\begin{aligned}
Q_I(T')
&=K_I(u)^{-1}\bigl(P_IK_I(u)P_I
 +(I_n-P_I)K_I(u)(I_n-P_I)\bigr)\\
&=K_I(u)^{-1}.
\end{aligned}
$$
Thus $q_I(T')=h_I(u)^{-1}$ is invertible. Sending each original
$t_{ab}$ to the corresponding entry of $T'$ defines
$R[q_I^{-1}]\to B_I$. Conversely, sending each $u_{ab}$ to the
corresponding entry of $Q_I^{-1}$ defines $B_I\to R[q_I^{-1}]$:
the diagonal blocks agree by the identity just proved, and
$h_I$ maps to $q_I^{-1}$. These homomorphisms are inverse.
On the original matrix the composite is $Q_IP_IQ_I^{-1}=T$;
on the matrix of chart variables it is $Q_I(T')^{-1}=K_I(u)$.
The inverted determinants have the corresponding inverse images.
We have therefore proved the exact chart isomorphism
$$
R[q_I^{-1}]\cong B_I.
$$

In particular every chart ring is an integral domain and is reduced.
These charts prove that $R$ itself is reduced. If $a\in R$ is
nilpotent, its image is zero in each chart, so for each $I$ some
positive power $q_I^{N_I}$ annihilates $a$. The finitely many
$q_I$ cover the spectrum, since their sum over all $I$ is
$\sum_r d_r=1$. The powers $q_I^{N_I}$ also generate the unit
ideal: a maximal ideal containing them would contain every $q_I$,
contrary to that sum identity. Taking a linear combination equal
to $1$ proves $a=0$.
Each factor $R_r$ is reduced and has irreducible spectrum, hence is
a domain: if $ab=0$ and neither factor were zero, reducedness would
make both $D(a)$ and $D(b)$ nonempty, whereas their intersection is
$D(ab)=\varnothing$, contradicting irreducibility.
All constructions commute with every field extension $L/k$:
$R\otimes_k L=L[t_{ij}]/((T^2-T)_{ij})$, with the same coefficient
polynomials $d_r$ and chart maps. Repeating the argument over $L$
shows that each $R_r\otimes_k L$ is a domain. Thus each component
is geometrically integral, in the sense of being integral after
every extension of the original base field.

We can also identify its defining ideal exactly. Let $J_\ell(T)$
be the ideal generated by the $\ell\times\ell$ minors, with
$J_0(T)=R$ and $J_{n+1}(T)=0$. The diagonal entries of
$\bigwedge^jT$ in its indexed exterior basis are the principal
$j\times j$ minors. Expanding the determinant therefore gives
$$
\det(I_n+uT)=\sum_{j=0}^n e_j(T)u^j,
\qquad
d_s=\sum_{j=s}^n(-1)^{j-s}\binom js e_j(T).
$$
For $s\geq\ell$, every minor appearing in the latter sum belongs
to $J_\ell(T)$, by Laplace expansion along any fixed $\ell$ of
its rows. Hence $(d_\ell,\ldots,d_n)\subseteq J_\ell(T)$.
Conversely, on a chart of $X_s$ with $s<\ell$, the identity
$T=Q_IP_IQ_I^{-1}$ gives
$$
\bigwedge^\ell T
=(\bigwedge^\ell Q_I)(\bigwedge^\ell P_I)
 (\bigwedge^\ell Q_I^{-1})=0,
$$
since $P_I$ has only $s$ nonzero diagonal entries.
All $\ell\times\ell$ minors therefore vanish on each chart of
$X_s$, hence in $R_s$ by the same finite-cover annihilator argument.
In the proved product decomposition of $R$, the ideal of elements
whose coordinates for $s<\ell$ are zero is exactly
$(d_\ell,\ldots,d_n)$. This proves
$$
J_\ell(T)=(d_\ell,\ldots,d_n)
\quad(0\leq\ell\leq n+1).
$$
For the complementary idempotent matrix the Laurent-polynomial
identity
$$
\det(T+z(I_n-T))
=z^n\det(I_n-T+z^{-1}T)
$$
shows that $d_s(I_n-T)=d_{n-s}(T)$. Applying the minor formula
to both matrices gives the exact component ideal
$$
(1-d_r)=J_{r+1}(T)+J_{n-r+1}(I_n-T).
$$
Indeed the right side is generated by all $d_s$ with $s\ne r$;
their sum is $1-d_r$ and each satisfies $d_s(1-d_r)=d_s$.
These are also precisely the minimal prime ideals of $R$, since
each quotient $R_r$ is a domain and the $X_r$ are all components.
Finally $F(1+u)=\det(I_n+uT)$ gives
$$
e_j(T)=\sum_{r=0}^n\binom rj d_r,
$$
an equality in the original ring, strengthening the residue-field
trace calculation. Every integer coefficient is mapped to $k$;
no characteristic restriction has been introduced. If $n=0$,
empty determinants are $1$, $R=k$, $d_0=1$, and the same statements
hold with their empty matrices and lists.
\end{example}


\subsection{Editorial treatment of Algebra lines 7455--7467}

\hypertarget{algebra-context-069}{}

\noindent\href{https://github.com/stacks/stacks-project/blob/a04446e57ec1fbc252a871afcec7752fb2807b14/algebra.tex\#L7455}{Original source, lines 7455--7467}. Complete original and reviewed TeX passages accompany this note. Every intervention is separately recorded; the following treatment is editorial.


\begin{lemma}
\label{editorial-source-lemma-characterize-integral-element}
Let $\varphi : R \to S$ be a ring map. Let $y \in S$. If there exists a
finite $R$-submodule $M$ of $S$ such that $1 \in M$ and $yM \subset M$,
then $y$ is integral over $R$.
Conversely, if $y$ is integral, the original subalgebra
$\varphi(R)[y]\subset S$ is a finite $R$-submodule containing $1$
and stable under multiplication by $y$.
\end{lemma}

\begin{proof}
Keep $\varphi:R\to S$ for the original coefficient map and write
$m_y:M\to M$, $x\mapsto yx$, for the $R$-linear endomorphism.
By Lemma \ref{editorial-source-lemma-charpoly-module} there is a monic
$P(T)=T^d+\sum_{i=1}^d a_iT^{d-i}\in R[T]$ with $P(m_y)=0$.
Evaluation at the original element $1\in M$ gives in $S$
$$
P^\varphi(y)=y^d+\sum_{i=1}^d\varphi(a_i)y^{d-i}
=P(m_y)(1)=0.
$$
Conversely, such a monic equation expresses $y^d$ as
$-\sum_{i=1}^d\varphi(a_i)y^{d-i}$. Multiplying this equation by
each needed power of $y$ and using induction expresses every power
of $y$ in the $R$-span of $1,y,\ldots,y^{d-1}$.
That span is exactly $\varphi(R)[y]$, contains $1$ and is stable
under $m_y$. If $S$ is the zero ring the same conclusion holds with
the zero submodule; a degree-one equation may be used.
\end{proof}


\subsection{Editorial treatment of Algebra lines 7480--7500}

\hypertarget{algebra-context-070}{}

\noindent\href{https://github.com/stacks/stacks-project/blob/a04446e57ec1fbc252a871afcec7752fb2807b14/algebra.tex\#L7480}{Original source, lines 7480--7500}. Complete original and reviewed TeX passages accompany this note. Every intervention is separately recorded; the following treatment is editorial.


\begin{lemma}
\label{editorial-source-lemma-characterize-integral}
Let $\varphi : R \to S$ be a ring map. Let $s_1, \ldots, s_n$
be a finite set of elements of $S$.
In this case $s_i$ is integral over $R$ for all $i = 1, \ldots, n$
if and only if
there exists an $R$-subalgebra $S' \subset S$ finite over $R$
containing all of the $s_i$.
\end{lemma}

\begin{proof}
For each original $s_i$ choose its monic equation
$s_i^{d_i}+\sum_{j=0}^{d_i-1}\varphi(a_{ij})s_i^j=0$ with
$d_i\geq1$. Consider the finite list of original monomials
$s_1^{e_1}\cdots s_n^{e_n}$ with $0\leq e_i<d_i$.
Their $R$-span contains $1$ and is stable under each $s_i$:
if multiplication reaches exponent $d_i$, substitute the full equation
$s_i^{d_i}=-\sum_{j=0}^{d_i-1}\varphi(a_{ij})s_i^j$, keeping every
coefficient and all the other exponents. Thus this span is precisely
$\varphi(R)[s_1,\ldots,s_n]$ and is generated by at most
$\prod_{i=1}^n d_i$ elements. For $n=0$ the list is the single
empty monomial $1$ and the subalgebra is $\varphi(R)$.
Conversely, suppose given a finite $R$-subalgebra
$S'$ containing all the $s_i$. Then all of the
$s_i$ are integral by Lemma \ref{algebra-lemma-finite-is-integral}.
\end{proof}


\subsection{Editorial treatment of Algebra lines 7502--7515}

\hypertarget{algebra-context-071}{}

\noindent\href{https://github.com/stacks/stacks-project/blob/a04446e57ec1fbc252a871afcec7752fb2807b14/algebra.tex\#L7502}{Original source, lines 7502--7515}. Complete original and reviewed TeX passages accompany this note. Every intervention is separately recorded; the following treatment is editorial.


\begin{lemma}
\label{editorial-source-lemma-characterize-finite-in-terms-of-integral}
Let $R \to S$ be a ring map. The following are equivalent
\begin{enumerate}
\item $R \to S$ is finite,
\item $R \to S$ is integral and of finite type, and
\item there exist $x_1, \ldots, x_n \in S$ which generate $S$ as an
algebra over $R$ such that each $x_i$ is integral over $R$.
\end{enumerate}
\end{lemma}

\begin{proof}
If $S$ is finite as an $R$-module, a finite module generating
family also generates it as an $R$-algebra: the latter subalgebra
already contains every linear combination of those elements.
Lemma \ref{algebra-lemma-finite-is-integral} proves integrality, giving
(1) $\Rightarrow$ (2). Under (2), choose the finite algebra
generators; integrality holds for each, proving (3).
Under (3), Lemma \ref{editorial-source-lemma-characterize-integral} shows that the
subalgebra generated by those same elements is finite over $R$.
It is the original $S$, so (1) follows.
\end{proof}


\subsection{Editorial treatment of Algebra lines 7541--7554}

\hypertarget{algebra-context-072}{}

\noindent\href{https://github.com/stacks/stacks-project/blob/a04446e57ec1fbc252a871afcec7752fb2807b14/algebra.tex\#L7541}{Original source, lines 7541--7554}. Complete original and reviewed TeX passages accompany this note. Every intervention is separately recorded; the following treatment is editorial.


\begin{lemma}
\label{editorial-source-lemma-integral-closure-is-ring}
Let $R \to S$ be a ring homomorphism.
The set
$$
S' = \{s \in S \mid s\text{ is integral over }R\}
$$
is an $R$-subalgebra of $S$.
\end{lemma}

\begin{proof}
The image of each $r\in R$ satisfies $T-r$, hence belongs to $S'$;
this includes $0$ and $1$. If $u,v\in S'$,
Lemma \ref{editorial-source-lemma-characterize-integral} places both in the finite
$R$-subalgebra generated by these two original elements. It contains
$u+v$, $uv$ and $-u$, each of which is integral over $R$ by
Lemma \ref{algebra-lemma-finite-is-integral}. Therefore $S'$ is a subring
containing the image of $R$, with its original coefficient map.
\end{proof}


\subsection{Editorial treatment of Algebra lines 7556--7566}

\hypertarget{algebra-context-073}{}

\noindent\href{https://github.com/stacks/stacks-project/blob/a04446e57ec1fbc252a871afcec7752fb2807b14/algebra.tex\#L7556}{Original source, lines 7556--7566}. Complete original and reviewed TeX passages accompany this note. Every intervention is separately recorded; the following treatment is editorial.


\begin{lemma}
\label{editorial-source-lemma-finite-product-integral}
Let $R_i\to S_i$ be ring maps $i = 1, \ldots, n$.
Let $R$ and $S$ denote the product of the $R_i$ and $S_i$ respectively.
Then an element $s = (s_1, \ldots, s_n) \in S$ is integral over $R$
if and only if each $s_i$ is integral over $R_i$.
For a product indexed by an arbitrary set, the exact corresponding
condition is the existence of monic equations for all $s_i$ whose
degrees have one common finite bound.
\end{lemma}

\begin{proof}
Let $\varphi_i:R_i\to S_i$ be the original maps. If the tuple $s$
satisfies a monic equation of degree $D$ over $R=\prod_iR_i$,
projection to the $i$th coordinate gives an equation of degree $D$
for $s_i$ over $R_i$, with each coefficient projected separately.
Conversely choose monic equations
$P_i(T)=T^{d_i}+\sum_{j=0}^{d_i-1}a_{ij}T^j$ and a finite bound
$D\geq d_i$ for every $i$. Keep these polynomials and form
$Q_i(T)=T^{D-d_i}P_i(T)$; its coefficient of $T^j$ is
$a_{i,j-(D-d_i)}$ for $D-d_i\leq j<D$ and is zero for
$0\leq j<D-d_i$. Let $b_j\in\prod_iR_i$ have precisely these
coordinates. Then
$$
Q(T)=T^D+\sum_{j=0}^{D-1}b_jT^j
$$
is monic and its evaluation at the original tuple has $i$th coordinate
$s_i^{D-d_i}P_i^{\varphi_i}(s_i)=0$. This proves the assertion for
arbitrary products with that uniform bound. For a nonempty finite
product take $D=\max_i d_i$; for an empty product the rings are zero
and the degree-one polynomial $T$ works.

The bound is necessary for an infinite product. Fix a field $k$ and
use $R_i=k[v_i]$, $S_i=k[u_i]$ with $v_i\mapsto u_i^i$ for $i\geq1$.
The original element $s_i=u_i$ satisfies $T^i-v_i$.
No nonzero polynomial of degree less than $i$ in $T$ with
coefficients in $k[v_i]$ vanishes at $u_i$: after substitution,
the terms $a_j(u_i^i)u_i^j$ for $0\leq j<i$ have powers of $u_i$
in distinct residue classes modulo $i$, so they cannot cancel.
Thus the least monic degree is $i$. Each $s_i$ is integral but
no uniform finite bound exists. The preceding equivalence shows
that $(s_i)_i$ is not integral over $\prod_{i\geq1}R_i$.
\end{proof}


\subsection{Editorial treatment of Algebra lines 7582--7594}

\hypertarget{algebra-context-074}{}

\noindent\href{https://github.com/stacks/stacks-project/blob/a04446e57ec1fbc252a871afcec7752fb2807b14/algebra.tex\#L7582}{Original source, lines 7582--7594}. Complete original and reviewed TeX passages accompany this note. Every intervention is separately recorded; the following treatment is editorial.


\begin{lemma}
\label{editorial-source-lemma-finite-product-integral-closure}
Let $R_i\to S_i$ be ring maps $i = 1, \ldots, n$.
Denote the integral closure of $R_i$ in $S_i$ by $S'_i$.
Further let $R$ and $S$ denote the product of the $R_i$ and $S_i$ respectively.
Then the integral closure of $R$ in $S$
is the product of the $S'_i$. In particular it equals the image
of $R\to S$ if and only if each $S'_i$ equals the image of
$R_i\to S_i$. When all the given maps are injective, this says
that $R$ is integrally closed in $S$ if and only if every $R_i$
is integrally closed in $S_i$, using Definition \ref{algebra-definition-integral-closure}.
\end{lemma}

\begin{proof}
By Lemma \ref{editorial-source-lemma-finite-product-integral}, the integral tuples
are exactly $\prod_iS'_i$, with coordinatewise operations and the
original product coefficient map. The image of that map is
$\prod_i\varphi_i(R_i)$: each tuple of image elements has a tuple
of preimages in the finite product. Equality of these two products
is equivalent to equality in each coordinate, by projection and
by filling the other coordinates with zero. If the maps are
injective, their product is injective and the image equality is
exactly the subring property defined above. The empty product
case gives the zero ring and the assertions hold vacuously.
\end{proof}


\subsection{Editorial treatment of Algebra lines 7596--7627}

\hypertarget{algebra-context-075}{}

\noindent\href{https://github.com/stacks/stacks-project/blob/a04446e57ec1fbc252a871afcec7752fb2807b14/algebra.tex\#L7596}{Original source, lines 7596--7627}. Complete original and reviewed TeX passages accompany this note. Every intervention is separately recorded; the following treatment is editorial.


\begin{lemma}
\label{editorial-source-lemma-integral-closure-localize}
Integral closure commutes with localization: If $A \to B$ is a ring
map, and $S \subset A$ is a multiplicative subset, then the integral
closure of $S^{-1}A$ in $S^{-1}B$ is $S^{-1}B'$, where $B' \subset B$
is the integral closure of $A$ in $B$.
\end{lemma}

\begin{proof}
Since localization is exact we see that $S^{-1}B' \subset S^{-1}B$.
Suppose $x \in B'$ and $f \in S$. Then
$x^d + \sum_{i = 1, \ldots, d} a_i x^{d - i} = 0$
in $B$ for some $a_i \in A$. Hence also
$$
(x/f)^d + \sum\nolimits_{i = 1, \ldots, d} a_i/f^i (x/f)^{d - i} = 0
$$
in $S^{-1}B$. In this way we see that $S^{-1}B'$ is contained in
the integral closure of $S^{-1}A$ in $S^{-1}B$. Conversely, suppose
that $x/f \in S^{-1}B$ is integral over $S^{-1}A$. Choosing a monic equation of positive degree, we have
$$
(x/f)^d + \sum\nolimits_{i = 1, \ldots, d} (a_i/f_i) (x/f)^{d - i} = 0
$$
in $S^{-1}B$ for the original $a_i\in A$ and $f_i\in S$.
Put $F=\prod_{j=1}^d f_j\in S$ and define the actual element of $B$
$$
H=Fx^d+\sum_{i=1}^d a_i f^i
\left(\prod_{j\ne i}f_j\right)x^{d-i},
$$
where all coefficients from $A$ act through the given map $A\to B$.
Multiplying the original equality by $f^dF$ shows that $H/1=0$
in $S^{-1}B$. Therefore there is $f'\in S$ with $f'H=0$ in $B$.
No injectivity of $B\to S^{-1}B$ is assumed. Multiplying this
last equality by $(f')^{d-1}F^{d-1}$ gives exactly
$$
(f'Fx)^d+
\sum_{i=1}^d f^i(f')^i
\left(f_i^{i-1}\prod_{j\ne i}f_j^i\right)a_i
(f'Fx)^{d-i}=0.
$$
Indeed every term equals its original term in $H$ multiplied by
$(f')^dF^{d-1}$, with no denominator remaining. This is a monic
equation for the original scaled element $y=f'Fx$, so $y\in B'$.
The equality of fractions
$$
x/f=y/(ff'F)
$$
holds in $S^{-1}B$ by multiplying numerator and denominator by
$f'F\in S$. The right side is the image of an element of
$S^{-1}B'$, proving the reverse inclusion with all original
denominators and the annihilating factor $f'$ retained.
\end{proof}


\subsection{Editorial treatment of Algebra lines 7629--7662}

\hypertarget{algebra-context-076}{}

\noindent\href{https://github.com/stacks/stacks-project/blob/a04446e57ec1fbc252a871afcec7752fb2807b14/algebra.tex\#L7629}{Original source, lines 7629--7662}. Complete original and reviewed TeX passages accompany this note. Every intervention is separately recorded; the following treatment is editorial.


\begin{lemma}
\label{editorial-source-lemma-integral-closure-stalks}
\begin{slogan}
An element of an algebra over a ring is integral over the ring
if and only if it is locally integral at every prime ideal of the ring.
\end{slogan}
Let $\varphi : R \to S$ be a ring map.
Let $x \in S$. The following are equivalent:
\begin{enumerate}
\item $x$ is integral over $R$,
\item for every prime ideal $\mathfrak p \subset R$ the element
$x/1\in S_{\mathfrak p}$ is integral over $R_{\mathfrak p}$, and
\item for every maximal ideal $\mathfrak m\subset R$, the element
$x/1\in S_{\mathfrak m}$ is integral over $R_{\mathfrak m}$.
\end{enumerate}
\end{lemma}

\begin{proof}
The monic equation for (1) remains a monic equation after every
localization, so (1) implies (2), and (2) implies (3).
For the reverse implication let $B\subset S$ be the integral
closure of $R$ in $S$, with its original inclusion. It is an
$R$-submodule because it is an $R$-subalgebra. Lemma
\ref{editorial-source-lemma-integral-closure-localize} identifies its localization
$B_{\mathfrak m}\subset S_{\mathfrak m}$ with the integral closure
of $R_{\mathfrak m}$ there. Under (3), the class $\overline x$
in the quotient $R$-module $S/B$ therefore vanishes after localization
at every maximal ideal. This uses the exact quotient comparison
$(S/B)_{\mathfrak m}\cong S_{\mathfrak m}/B_{\mathfrak m}$,
which sends $(s+B)/u$ to $s/u+B_{\mathfrak m}$; the inverse is
induced by $s/u\mapsto(s+B)/u$, and the fraction equality criterion
and quotient relation prove both maps well defined and inverse.
If $\overline x\ne0$, its annihilator in $R$ is a proper ideal
and lies in a maximal ideal $\mathfrak m$. Vanishing of
$\overline x/1$ would give $u\notin\mathfrak m$ with
$u\overline x=0$, contradicting that containment. Thus
$\overline x=0$ and $x\in B$, proving (1). If $R$ is zero,
the unital map forces $S$ zero and the conclusion holds directly.

For completeness, retain also the original finite-subalgebra
calculation. Let $S'\subset S$ be generated by $\varphi(R)$ and
the original $x$. At a prime $\mathfrak p$ under (2), write its
monic equation of positive degree with coefficients
$a_i=b_i/u_i\in R_{\mathfrak p}$,
where $b_i\in R$ and $u_i\notin\mathfrak p$ for $1\leq i\leq d$.
Take $U=\prod_{i=1}^d u_i$ and the element
$$
H=\varphi(U)x^d+
\sum_{i=1}^d\varphi\left(b_i\prod_{j\ne i}u_j\right)x^{d-i}
\quad\text{of }S.
$$
Its image in $S_{\mathfrak p}$ is zero. Hence some
$v\notin\mathfrak p$ satisfies $\varphi(v)H=0$ in $S$.
Set $f=vU\notin\mathfrak p$. In $R_f$, each $u_i$ and $v$
is a unit because each is a factor of $f$. Dividing the last
equality by $vU$ in $S_f$ recovers the complete original equation
$$
x^d+\sum_{i=1}^d\varphi_f(b_i/u_i)x^{d-i}=0,
$$
where $\varphi_f:R_f\to S_f$ is the original localized map.
Thus $S'_f$, still the localization of that specified subalgebra,
is spanned over $R_f$ by $1,x,\ldots,x^{d-1}$.
The opens $D(f)$ so obtained cover $\operatorname{Spec}(R)$;
quasi-compactness (Lemma \ref{algebra-lemma-quasi-compact}) gives finitely
many $f_1,\ldots,f_n$ whose opens cover. They generate the unit
ideal, since no prime contains all of them. Lemma \ref{editorial-source-lemma-cover}
then shows that the original $S'$ is finite over $R$. Finally
Lemma \ref{editorial-source-lemma-characterize-integral} proves $x$ integral.
The factors $u_i$, $v$, $U$ and $f$ explicitly account both for
coefficient denominators and for the kernel of localization.
\end{proof}


\subsection{Editorial treatment of Algebra lines 7664--7686}

\hypertarget{algebra-context-077}{}

\noindent\href{https://github.com/stacks/stacks-project/blob/a04446e57ec1fbc252a871afcec7752fb2807b14/algebra.tex\#L7664}{Original source, lines 7664--7686}. Complete original and reviewed TeX passages accompany this note. Every intervention is separately recorded; the following treatment is editorial.


\begin{lemma}
\label{editorial-source-lemma-base-change-integral}
\begin{slogan}
Integrality and finiteness are preserved under base change.
\end{slogan}
Let $R \to S$ and $R \to R'$ be ring maps.
Set $S' = R' \otimes_R S$.
\begin{enumerate}
\item If $R \to S$ is integral so is $R' \to S'$.
\item If $R \to S$ is finite so is $R' \to S'$.
\end{enumerate}
\end{lemma}

\begin{proof}
For (1), write an arbitrary element of the original tensor algebra
as a finite sum $z=\sum_{i=1}^m r'_i\otimes s_i$.
Each $s_i$ satisfies its monic equation
$s_i^{d_i}+\sum_{j=0}^{d_i-1}\varphi(a_{ij})s_i^j=0$.
Tensoring gives the same equation for $1\otimes s_i$, with the
coefficients $a_{ij}$ mapped into $R'$. The finite-subalgebra
argument of Lemma \ref{editorial-source-lemma-characterize-integral} places all
$1\otimes s_i$, and hence their $R'$-linear combination $z$,
in a finite $R'$-subalgebra of $S'$. Lemma
\ref{algebra-lemma-finite-is-integral} proves $z$ integral.
For (2), choose original $R$-module generators $s_1,\ldots,s_m$
of $S$. An expression $s=\sum_i a_i s_i$ gives
$r'\otimes s=\sum_i r'\varphi'(a_i)(1\otimes s_i)$, where
$\varphi':R\to R'$ is the given map. Every tensor is a finite
sum of such terms, so $1\otimes s_1,\ldots,1\otimes s_m$
generate $S'$ as an $R'$-module.
\end{proof}


\subsection{Editorial treatment of Algebra lines 7688--7707}

\hypertarget{algebra-context-078}{}

\noindent\href{https://github.com/stacks/stacks-project/blob/a04446e57ec1fbc252a871afcec7752fb2807b14/algebra.tex\#L7688}{Original source, lines 7688--7707}. Complete original and reviewed TeX passages accompany this note. Every intervention is separately recorded; the following treatment is editorial.


\begin{lemma}
\label{editorial-source-lemma-integral-local}
Let $R \to S$ be a ring map.
Let $f_1, \ldots, f_n \in R$ generate the unit ideal.
\begin{enumerate}
\item If each $R_{f_i} \to S_{f_i}$ is integral, so is $R \to S$.
\item If each $R_{f_i} \to S_{f_i}$ is finite, so is $R \to S$.
\end{enumerate}
\end{lemma}

\begin{proof}
For (1), fix the original $s\in S$ and keep
$I=\{P\in R[x]:P^\varphi(s)=0\}$, the kernel of evaluation.
Let $J\subset R$ consist of zero and the leading coefficients
of the nonzero polynomials in $I$. This is an ideal. To add two
nonzero leading coefficients, multiply the corresponding polynomials
by powers of $x$ to make their degrees equal. If the sum of the
leading coefficients is nonzero, it is the leading coefficient of
the sum polynomial; if it is zero, it belongs to $J$ by definition.
For multiplication by $r\in R$, a nonzero product of leading
coefficients remains the leading coefficient of $rP$; a zero
product again belongs to $J$. Thus all ideal axioms hold.

For each original $f_i$, choose the local monic equation
$$
s^{d_i}+\sum_{j=0}^{d_i-1}
\varphi_{f_i}(a_{ij}/f_i^{c_{ij}})s^j=0
\quad\text{in }S_{f_i},
$$
where $d_i\geq1$, $a_{ij}\in R$ and $c_{ij}\geq0$.
Choose $N_i\geq c_{ij}$ for all $j$. The element
$\varphi(f_i^{N_i})s^{d_i}+
\sum_j\varphi(a_{ij}f_i^{N_i-c_{ij}})s^j$ becomes zero in
$S_{f_i}$, so some $M_i\geq0$ makes it zero upon multiplication
by $\varphi(f_i^{M_i})$. Consequently the original polynomial
$$
f_i^{N_i+M_i}x^{d_i}
+\sum_{j=0}^{d_i-1}a_{ij}f_i^{N_i+M_i-c_{ij}}x^j
$$
belongs to $I$. If $f_i^{N_i+M_i}\ne0$ it is its leading
coefficient; if it is zero it belongs to $J$ by definition.
In either case $f_i^{N_i+M_i}\in J$. A proper $J$ would lie
in a maximal ideal containing every $f_i$ (or would already
contain $1$ when an exponent is zero), contrary to the given
unit-ideal hypothesis. Therefore $1\in J$, so $I$ contains a
monic polynomial and $s$ is integral over $R$.

The maximal-ideal criterion above gives the same conclusion:
every maximal ideal omits some $f_i$ and the corresponding local
map is a further localization of $R_{f_i}\to S_{f_i}$.

For (2), choose finite $R_{f_i}$-module generators of $S_{f_i}$.
Write each as $s_{ij}/f_i^{a_{ij}}$, keeping its numerator
$s_{ij}\in S$ and exponent $a_{ij}\geq0$.
Let $N\subset S$ be the $R$-submodule generated by all these
finitely many numerators. Since $f_i$ is a unit after localization,
$N_{f_i}=S_{f_i}$. Thus the quotient $Q=S/N$ satisfies $Q_{f_i}=0$
for each $i$. For $q\in Q$ choose $b_i\geq1$ with $f_i^{b_i}q=0$.
The ideal generated by the $f_i^{b_i}$ is the unit ideal: if it were
proper, a maximal ideal containing it would contain all $f_i$,
contrary to the original unit-ideal assumption. Choose $c_i\in R$
with $\sum_i c_if_i^{b_i}=1$. Then
$q=\sum_i c_if_i^{b_i}q=0$. Hence $Q=0$ and the original numerators
$s_{ij}$ generate $S$ over $R$. The empty-cover case forces $R=0$
and hence $S=0$, where the assertion also holds.
\end{proof}


\subsection{Editorial treatment of Algebra lines 7709--7720}

\hypertarget{algebra-context-079}{}

\noindent\href{https://github.com/stacks/stacks-project/blob/a04446e57ec1fbc252a871afcec7752fb2807b14/algebra.tex\#L7709}{Original source, lines 7709--7720}. Complete original and reviewed TeX passages accompany this note. Every intervention is separately recorded; the following treatment is editorial.


\begin{lemma}
\label{editorial-source-lemma-integral-permanence}
Let $A \to B \to C$ be ring maps.
\begin{enumerate}
\item If $A \to C$ is integral so is $B \to C$.
\item If $A \to C$ is finite so is $B \to C$.
\end{enumerate}
\end{lemma}

\begin{proof}
Write $\alpha:A\to B$ and $\beta:B\to C$ for the original maps.
If $c\in C$ satisfies
$c^d+\sum_{i=0}^{d-1}(\beta\alpha)(a_i)c^i=0$, the same monic
polynomial with coefficients $\alpha(a_i)\in B$ proves integrality
over $B$. This proves (1) for every original $c$.
If $c_1,\ldots,c_m$ generate $C$ as an $A$-module, every expression
$c=\sum_i(\beta\alpha)(a_i)c_i$ is also a $B$-linear expression,
with coefficients $\alpha(a_i)$. The same family generates $C$
as a $B$-module, proving (2).
\end{proof}


\subsection{Editorial treatment of Algebra lines 7722--7732}

\hypertarget{algebra-context-080}{}

\noindent\href{https://github.com/stacks/stacks-project/blob/a04446e57ec1fbc252a871afcec7752fb2807b14/algebra.tex\#L7722}{Original source, lines 7722--7732}. Complete original and reviewed TeX passages accompany this note. Every intervention is separately recorded; the following treatment is editorial.


\begin{lemma}
\label{editorial-source-lemma-integral-closure-transitive}
Let $A \to B \to C$ be ring maps.
Let $B'$ be the integral closure of $A$ in $B$,
let $C'$ be the integral closure of $B'$ in $C$. Then
$C'$ is the integral closure of $A$ in $C$.
\end{lemma}

\begin{proof}
The coefficient map $A\to B$ lands in $B'$ by the degree-one
equations for its image elements, and $A\to B'$ is integral by
definition. If $c\in C'$ is integral over $B'$, transitivity of
integrality (Lemma \ref{algebra-lemma-integral-transitive}) shows that $c$
is integral over $A$. Conversely, every monic equation for $c$
over $A$ becomes a monic equation over $B'$ under the original
map $A\to B'$, since its composite into $C$ is unchanged.
Thus the two sets of integral elements of the same ring $C$ agree,
and their subalgebra structures and coefficient maps agree as well.
\end{proof}


\subsection{Editorial treatment of Algebra lines 7734--7762}

\hypertarget{algebra-context-081}{}

\noindent\href{https://github.com/stacks/stacks-project/blob/a04446e57ec1fbc252a871afcec7752fb2807b14/algebra.tex\#L7734}{Original source, lines 7734--7762}. Complete original and reviewed TeX passages accompany this note. Every intervention is separately recorded; the following treatment is editorial.


\begin{lemma}
\label{editorial-source-lemma-integral-overring-surjective}
Suppose that $R \to S$ is an integral
ring extension with $R \subset S$.
Then $\varphi : \Spec(S) \to \Spec(R)$
is surjective. More generally, for an arbitrary integral ring map
$\varphi:R\to S$ its spectrum image is $V(\ker\varphi)$.
\end{lemma}

\begin{proof}
Let $\mathfrak p \subset R$ be a prime ideal.
We have to show $\mathfrak pS_{\mathfrak p} \not = S_{\mathfrak p}$, see
Lemma \ref{editorial-source-lemma-in-image}.
The localization $R_{\mathfrak p} \to S_{\mathfrak p}$ is injective
(as localization is exact) and integral by
Lemma \ref{editorial-source-lemma-integral-closure-localize} or
\ref{editorial-source-lemma-base-change-integral}.
Keep these localizations and put $\mathfrak m=\mathfrak pR_{\mathfrak p}$.
Suppose, for contradiction, that
$1=\sum_i f_is_i$ in $S_{\mathfrak p}$ with
$f_i\in\mathfrak m$ and $s_i\in S_{\mathfrak p}$, using the
original localized coefficient map. Lemma
\ref{editorial-source-lemma-characterize-integral} gives an
$R_{\mathfrak p}$-subalgebra $S'\subset S_{\mathfrak p}$
containing all these $s_i$ and finite over $R_{\mathfrak p}$.
The same equation for $1$ holds in $S'$. Multiplying it by every
element of $S'$ proves $S'=\mathfrak mS'$. Nakayama's
Lemma \ref{algebra-lemma-NAK} gives $S'=0$, which contradicts the injection
of the nonzero local ring $R_{\mathfrak p}$ into $S'$.
Thus the original fibre has a prime and the desired preimage exists.

For arbitrary $\varphi$, the induced map
$R/\ker\varphi\to S$, $r+\ker\varphi\mapsto\varphi(r)$,
is injective and integral: reduce each coefficient of a monic
equation modulo the kernel. The injective case just proved makes
its spectrum map surjective. Composing with the original quotient
map identifies its image in $\operatorname{Spec}(R)$ with
$V(\ker\varphi)$, because primes of the quotient are exactly
the primes of $R$ containing that kernel.
\end{proof}


\subsection{Editorial treatment of Algebra lines 7764--7779}

\hypertarget{algebra-context-082}{}

\noindent\href{https://github.com/stacks/stacks-project/blob/a04446e57ec1fbc252a871afcec7752fb2807b14/algebra.tex\#L7764}{Original source, lines 7764--7779}. Complete original and reviewed TeX passages accompany this note. Every intervention is separately recorded; the following treatment is editorial.


\begin{lemma}
\label{editorial-source-lemma-integral-under-field}
Let $R$ be a ring. Let $K$ be a field.
If $R \subset K$ and $K$ is integral over $R$,
then $R$ is a field and $K$ is an algebraic extension.
If $R \subset K$ and $K$ is finite over $R$,
then $R$ is a field and $K$ is a finite algebraic extension.
\end{lemma}

\begin{proof}
Assume that $R \subset K$ is integral.
By Lemma \ref{editorial-source-lemma-integral-overring-surjective} we see that
$\Spec(R)$ has $1$ point. Since clearly $R$ is a domain we see
that $R = R_{(0)}$ is a field (Lemma \ref{algebra-lemma-minimal-prime-reduced-ring}).
For an explicit inverse calculation, if $0\ne r\in R$, the
element $r^{-1}\in K$ satisfies a monic equation
$(r^{-1})^d+\sum_{i=1}^d a_i(r^{-1})^{d-i}=0$ with $a_i\in R$
and $d\geq1$. Multiplication by $r^{d-1}$ in the original field
gives the full formula
$$
r^{-1}=-\sum_{i=1}^d a_i r^{i-1}\in R.
$$
Every element of $K$ is then algebraic over the field $R$, since
its monic integral equation is a nonzero polynomial over $R$.
In the finite case, Lemma \ref{algebra-lemma-finite-is-integral} first
gives integrality, and the original finite module generators
become a finite vector-space spanning family over this same field
$R$. Hence the algebraic extension has finite degree.
\end{proof}


\subsection{Editorial treatment of Algebra lines 7781--7810}

\hypertarget{algebra-context-083}{}

\noindent\href{https://github.com/stacks/stacks-project/blob/a04446e57ec1fbc252a871afcec7752fb2807b14/algebra.tex\#L7781}{Original source, lines 7781--7810}. Complete original and reviewed TeX passages accompany this note. Every intervention is separately recorded; the following treatment is editorial.


\begin{lemma}
\label{editorial-source-lemma-integral-over-field}
Let $k$ be a field. Let $S$ be a $k$-algebra over $k$.
\begin{enumerate}
\item If $S$ is a domain and finite dimensional over $k$,
then $S$ is a field.
\item If $S$ is integral over $k$ and a domain,
then $S$ is a field.
\item If $S$ is integral over $k$ then every prime of
$S$ is a maximal ideal (see
Lemma \ref{algebra-lemma-ring-with-only-minimal-primes}
for more consequences).
\end{enumerate}
\end{lemma}

\begin{proof}
For (3), a quotient $S/\mathfrak q$ by any prime is an integral
$k$-algebra and a domain. Its coefficient map from $k$ is injective,
since the quotient is nonzero and $k$ is a field. Assertion (2)
makes this quotient a field, hence $\mathfrak q$ is maximal.
Let $S$ be integral over $k$ and assume $S$ is a domain.
Take a nonzero $s\in S$. By Lemma
\ref{editorial-source-lemma-characterize-integral} we may find a
finite dimensional $k$-subalgebra $k \subset S' \subset S$
containing $s$. This subalgebra is a domain since it is a subring
of $S$. Once (1) is proved, $S'$ is a field and the inverse of
the original nonzero $s$ belongs to $S'\subset S$, proving (2).
For (1), assume that the original $S$ is finite dimensional over
$k$ and is a domain. Pick a nonzero $s\in S$. The $k$-linear
map $S\to S$, $a\mapsto sa$, is injective because $S$ is a
domain and $s\ne0$. Equality of the finite source and target
dimensions makes it surjective. A preimage $s_1$ of $1$ satisfies
$ss_1=1$. Thus every nonzero element is invertible and $S$ is a field.
\end{proof}


\subsection{Editorial treatment of Algebra lines 7812--7829}

\hypertarget{algebra-context-084}{}

\noindent\href{https://github.com/stacks/stacks-project/blob/a04446e57ec1fbc252a871afcec7752fb2807b14/algebra.tex\#L7812}{Original source, lines 7812--7829}. Complete original and reviewed TeX passages accompany this note. Every intervention is separately recorded; the following treatment is editorial.


\begin{lemma}
\label{editorial-source-lemma-integral-no-inclusion}
Suppose $R \to S$ is integral.
Let $\mathfrak q, \mathfrak q' \in \Spec(S)$
be distinct primes
having the same image in $\Spec(R)$.
Then neither $\mathfrak q \subset \mathfrak q'$
nor $\mathfrak q' \subset \mathfrak q$.
\end{lemma}

\begin{proof}
Let $\mathfrak p=\varphi^{-1}(\mathfrak q)=
\varphi^{-1}(\mathfrak q')$ and set $U=R\setminus\mathfrak p$.
Retain the fibre algebra
$A=S\otimes_R\kappa(\mathfrak p)\cong U^{-1}S/\mathfrak pU^{-1}S$.
The two ideals are sent to
$U^{-1}\mathfrak q/\mathfrak pU^{-1}S$ and
$U^{-1}\mathfrak q'/\mathfrak pU^{-1}S$. Localization and quotient
preserve inclusion, and contraction recovers each original prime:
$s/1\in U^{-1}\mathfrak q$ means
$\varphi(u)s\in\mathfrak q$ for some $u\in U$, which by primality
and $\varphi(u)\notin\mathfrak q$ is equivalent to $s\in\mathfrak q$.
Thus the two resulting primes are distinct. Lemma
\ref{editorial-source-lemma-base-change-integral} makes $A$ integral over
$\kappa(\mathfrak p)$, so Lemma \ref{editorial-source-lemma-integral-over-field}
makes both primes maximal. An inclusion between them would force
equality, and contraction would contradict the original distinction.
\end{proof}


\subsection{Editorial treatment of Algebra lines 7831--7850}

\hypertarget{algebra-context-085}{}

\noindent\href{https://github.com/stacks/stacks-project/blob/a04446e57ec1fbc252a871afcec7752fb2807b14/algebra.tex\#L7831}{Original source, lines 7831--7850}. Complete original and reviewed TeX passages accompany this note. Every intervention is separately recorded; the following treatment is editorial.


\begin{lemma}
\label{editorial-source-lemma-finite-finite-fibres}
Suppose $R \to S$ is finite.
Then the fibres of $\Spec(S)\to\Spec(R)$ are finite.
More precisely, if $s_1,\ldots,s_m$ generate $S$ as an $R$-module,
then for every $\mathfrak p\in\operatorname{Spec}(R)$,
$$
\begin{aligned}
\#\{\mathfrak q:\varphi^{-1}(\mathfrak q)=\mathfrak p\}
&\leq\sum_{\varphi^{-1}(\mathfrak q)=\mathfrak p}
[\kappa(\mathfrak q):\kappa(\mathfrak p)]
\\
&\leq\dim_{\kappa(\mathfrak p)}(S\otimes_R\kappa(\mathfrak p))
\leq m.
\end{aligned}
$$
The middle difference is exactly the dimension of the nilradical
of this fibre algebra. In particular the residue-degree sum equals
the fibre dimension exactly when this fibre is reduced.
\end{lemma}

\begin{proof}
By the discussion in
Remark \ref{algebra-remark-fundamental-diagram}
the fibres are the spectra of the rings $S \otimes_R \kappa(\mathfrak p)$.
As $R \to S$ is finite, these fibre rings are finite over
$\kappa(\mathfrak p)$ hence Noetherian by
Lemma \ref{editorial-source-lemma-Noetherian-permanence}.
By
Lemma \ref{editorial-source-lemma-integral-no-inclusion}
every prime of $S \otimes_R \kappa(\mathfrak p)$ is a minimal
prime. Hence by
Lemma \ref{editorial-source-lemma-Noetherian-irreducible-components}
there are at most finitely many.
\medskip\noindent
For the quantitative assertions, keep the original map
$\varphi:R\to S$ and the given module generators $s_1,\ldots,s_m$.
Fix $\mathfrak p\in\operatorname{Spec}(R)$ and put
$K=\kappa(\mathfrak p)$, $U=R\setminus\mathfrak p$,
$A=S\otimes_R K$ and $d=\dim_K A$.
The original surjection $R^m\to S$ induces the surjection
$$
K^m\longrightarrow A,\qquad
(\lambda_j)_j\longmapsto\sum_{j=1}^m s_j\otimes\lambda_j.
$$
Indeed $s=\sum_j\varphi(a_j)s_j$ gives
$s\otimes\lambda=\sum_j s_j\otimes(\overline{a_j}\lambda)$,
and every tensor is a finite sum of pure tensors. Thus $d\leq m$.

The exact fibre-ring isomorphism and its inverse are
$$
\begin{aligned}
A&\longrightarrow U^{-1}S/\mathfrak pU^{-1}S,&
s\otimes\frac{r+\mathfrak p}{u+\mathfrak p}
&\longmapsto\frac{s\varphi(r)}{\varphi(u)}+\mathfrak pU^{-1}S,\\
\frac{s}{\varphi(u)}+\mathfrak pU^{-1}S
&\longmapsto s\otimes\frac1{u+\mathfrak p}.&&
\end{aligned}
$$
Here $U^{-1}S$ means localization at the image $\varphi(U)$.
Tensor balancing respects the original coefficient map; a change
of fraction representatives is killed by a denominator from $U$
on both sides, and an element of $\mathfrak p$ acts as zero on
$K$. These facts verify the defining relations of both maps.
Their composites on the displayed fractions and pure tensors
are identities, proving the isomorphism.
Write $\beta:S\to A$, $s\mapsto s\otimes1$.
A prime $\mathfrak q$ of $S$ with
$\varphi^{-1}(\mathfrak q)=\mathfrak p$ corresponds to
$$
\mathfrak n_{\mathfrak q}
=U^{-1}\mathfrak q/\mathfrak pU^{-1}S\subset A.
$$
It is proper because $\mathfrak q$ avoids $\varphi(U)$, and
contraction recovers $\mathfrak q$: the condition
$s/1\in U^{-1}\mathfrak q$ is equivalent to
$\varphi(u)s\in\mathfrak q$ for some $u\in U$, hence to
$s\in\mathfrak q$ by primality. Conversely, a prime
$\mathfrak n\subset A$ contracts to a prime $\mathfrak q\subset S$
whose contraction in $R$ is $\mathfrak p$: elements of
$\mathfrak p$ act as zero and elements of $U$ act as units.
Every element of $A$ is represented by a fraction $s/\varphi(u)$
modulo $\mathfrak p$, and it belongs to $\mathfrak n$ precisely
when $\beta(s)$ does. This verifies the remaining composite of
the two prime correspondences. In particular the correspondence
is a bijection for the original, possibly noninjective map.

For every prime $\mathfrak n$ of $A$, the quotient $A/\mathfrak n$
is a finite dimensional domain over $K$, hence a field by
Lemma \ref{editorial-source-lemma-integral-over-field}. Therefore all these primes
are maximal. If $\mathfrak n=\mathfrak n_{\mathfrak q}$,
its field is canonically the original residue field:
$$
A/\mathfrak n_{\mathfrak q}
\cong U^{-1}(S/\mathfrak q)
\xrightarrow{\ \sim\ }\kappa(\mathfrak q),\qquad
\frac{s+\mathfrak q}{\varphi(u)+\mathfrak q}
\longmapsto\frac{s+\mathfrak q}{\varphi(u)+\mathfrak q}.
$$
The last map is injective since $S/\mathfrak q$ is a domain.
Its source is already a field containing $S/\mathfrak q$,
so it contains the inverse of each nonzero element of that domain;
it therefore contains every fraction and the map is surjective.
This also identifies the coefficient field $K$ in both targets.

Here is the Chinese remainder map with its actual preimages.
For any finite family of distinct maximal ideals
$\mathfrak n_1,\ldots,\mathfrak n_t$ of $A$, comaximality gives
$b_{ij}\in\mathfrak n_j$ with $1-b_{ij}\in\mathfrak n_i$
when $i\ne j$. Put $e_i=\prod_{j\ne i}b_{ij}$.
It is $1$ modulo $\mathfrak n_i$ and $0$ modulo each other
$\mathfrak n_j$. Hence
$$
\rho:A\longrightarrow\prod_{i=1}^t A/\mathfrak n_i,
\qquad a\longmapsto(a+\mathfrak n_i)_i
$$
is surjective: representatives $a_i$ of a prescribed tuple have
preimage $\sum_i e_ia_i$. Its kernel is exactly
$\bigcap_i\mathfrak n_i$. Comparing $K$-dimensions gives
$\sum_i\dim_K(A/\mathfrak n_i)\leq d$. Every field factor has
positive dimension, so there cannot be more than $d$ maximal
ideals: any $d+1$ distinct ones would contradict the inequality.
If $A\ne0$ it has a maximal ideal, for example a proper ideal of
greatest $K$-dimension. Thus its primes are a finite family and
the preceding construction applies to all of them at once.
Let $L_i=A/\mathfrak n_i\cong\kappa(\mathfrak q_i)$ and
$f_i=[L_i:K]$. With $J=\bigcap_i\mathfrak n_i$, it gives
$$
0\longrightarrow J\longrightarrow A\xrightarrow{\rho}
\prod_iL_i\longrightarrow0,
\qquad
d=\sum_i f_i+\dim_KJ.
$$

The ideal $J$ is precisely the nilradical, and is nilpotent.
To prove the latter assertion directly, the decreasing chain
$J,J^2,J^3,\ldots$ stabilizes since these are subspaces of the
finite dimensional $K$-space $A$. Choose $N\geq1$ with
$J^N=J^{N+1}$. A finite $K$-basis of $J^N$ generates it as an
$A$-module. Write each generator as a $J$-linear combination
of all generators, giving a matrix $C$ with entries in $J$.
Multiplication by the adjugate of $I-C$ gives
$\det(I-C)J^N=0$. This determinant belongs to $1+J$ and is a
unit: if it lay in a maximal ideal, that ideal would also contain
$J$, and hence would contain $1$. Therefore $J^N=0$.
Every nilpotent element lies in all prime ideals, and conversely
every element of $J$ is now nilpotent, proving the assertion.
Consequently the middle inequality is an equality exactly when
$A$ is reduced. Together with $f_i\geq1$ and $d\leq m$, this
proves all the displayed inequalities. For $A=0$ there are no
prime ideals, its dimension is zero, and the formulas use empty sums.

There is a more precise local formula, retaining every multiplicity.
For $A\ne0$ and the above $N$, the ideals $\mathfrak n_i^N$ are
pairwise comaximal. In fact, if $a+b=1$ with
$a\in\mathfrak n_i$, $b\in\mathfrak n_j$, every term of
$(a+b)^{2N-1}=1$ belongs to $\mathfrak n_i^N$ or $\mathfrak n_j^N$.
For two comaximal ideals their intersection is their product:
if $x$ is in both and $a+b=1$ with one element in each ideal,
then $x=xa+xb$ belongs to the product. Induction gives the finite
version; the product of the first ideals stays comaximal with
the next since multiplying equations congruent to $1$ modulo
the next ideal gives another such equation. It follows that
$$
\bigcap_i\mathfrak n_i^N
=\prod_i\mathfrak n_i^N
=\left(\prod_i\mathfrak n_i\right)^N
=J^N=0.
$$
The same Chinese remainder construction therefore gives
$$
A\xrightarrow{\ \sim\ }\prod_i C_i,
\qquad C_i=A/\mathfrak n_i^N.
$$
Each $C_i$ has unique maximal ideal
$\mathfrak m_i=\mathfrak n_i/\mathfrak n_i^N$: a prime
containing $\mathfrak n_i^N$ contains every element of
$\mathfrak n_i$, by taking its $N$th power, and thus equals
that maximal ideal. The residue field is $L_i$.

The identification $C_i\cong A_{\mathfrak n_i}$ has explicit maps.
Choose the Chinese remainder element $e_i$ that is $1$ modulo
$\mathfrak n_i^N$ and $0$ modulo the other powers. For
$a\in\mathfrak n_i^N$, the product $e_ia$ is in the intersection
of all powers and is zero. Since $e_i\notin\mathfrak n_i$,
this gives $a/1=0$ in $A_{\mathfrak n_i}$. Thus
$\overline a\mapsto a/1$ defines $C_i\to A_{\mathfrak n_i}$.
Its inverse sends $a/u$ to $\overline a\,\overline u^{-1}$;
the inverse exists since $u\notin\mathfrak n_i$ and $C_i$ is
local. The two composites on elements and fractions are identities.
There is also the exact original-ring comparison
$$
A_{\mathfrak n_i}\cong S_{\mathfrak q_i}/\mathfrak pS_{\mathfrak q_i}.
$$
One direction sends $s/v$ modulo $\mathfrak p$ to
$\beta(s)/\beta(v)$ for $v\notin\mathfrak q_i$.
The other is induced by
$s\otimes((r+\mathfrak p)/(u+\mathfrak p))
\mapsto s\varphi(r)/\varphi(u)$ modulo $\mathfrak p$.
It extends to localization at $\mathfrak n_i$, since an element
outside that maximal ideal maps to an element with nonzero
residue in the identified field $\kappa(\mathfrak q_i)$ and
hence to a unit of the target local ring. Tensor, fraction and
quotient relations are respected as in the preceding fibre map,
and the composites fix every such fraction, proving the comparison.

Use the finite filtration
$C_i\supset\mathfrak m_i\supset\mathfrak m_i^2\supset\cdots\supset0$.
Each quotient $\mathfrak m_i^j/\mathfrak m_i^{j+1}$ is annihilated
by $\mathfrak m_i$ and is a finite dimensional $L_i$-space.
Put
$$
\mu_i=\sum_{j\geq0}
\dim_{L_i}(\mathfrak m_i^j/\mathfrak m_i^{j+1}).
$$
This positive integer is the length of $C_i$ over itself:
refine each quotient by a vector-space flag, and note that every
simple module over the local ring $C_i$ is $L_i$, because the
annihilator of a nonzero simple-module generator is its unique
maximal ideal. Every composition series has the same number of
factors, since each such factor has $K$-dimension $f_i$.
Taking dimensions of the filtration and of the product yields
the exact formulas
$$
\begin{aligned}
\dim_K(S\otimes_RK)&=\sum_i\mu_i[\kappa(\mathfrak q_i):K],\\
\dim_KJ&=\sum_i(\mu_i-1)[\kappa(\mathfrak q_i):K],\\
\mu_i&=\operatorname{length}_{S_{\mathfrak q_i}/\mathfrak pS_{\mathfrak q_i}}
 (S_{\mathfrak q_i}/\mathfrak pS_{\mathfrak q_i}).
\end{aligned}
$$
Thus no local multiplicity has been absorbed into a residue degree.

The empty fibre also has an exact criterion:
$A=0$ if and only if $\ker\varphi\not\subseteq\mathfrak p$.
If some $u\in U$ is in the kernel, its zero image is inverted
in $U^{-1}S$, giving the zero ring. Conversely $A=0$ implies
$U^{-1}S=\mathfrak pR_{\mathfrak p}(U^{-1}S)$.
This is a finite $R_{\mathfrak p}$-module generated by the images
of the specified $s_j$. Nakayama's Lemma \ref{algebra-lemma-NAK} gives
$U^{-1}S=0$, so $1/1=0$ implies $\varphi(u)=0$ for some
$u\in U$. This proves the criterion without assuming injectivity.

For clarity, all equality cases follow from these formulas.
The number of primes equals the residue-degree sum precisely
when every $f_i=1$. That sum equals $d$ precisely when $J=0$,
equivalently all $\mu_i=1$. The upper equality $d=m$ holds
precisely when the original images $s_j\otimes1$ are a $K$-basis.
In particular the number of primes equals $d$ exactly when
$A\cong K^t$, where $t$ is the number of primes: the equalities
give $J=0$, $f_i=1$ and the proved Chinese remainder isomorphism;
the converse follows from that product's coordinate maximal ideals.
The number of primes equals $m$ exactly when $A\cong K^m$.
These statements include the empty product, interpreted as the zero ring.
In terms of the original relation module
$H=\ker(R^m\to S)$, the kernel of $K^m\to A$ is exactly
$$
\frac{H_{\mathfrak p}+\mathfrak pR_{\mathfrak p}^m}
{\mathfrak pR_{\mathfrak p}^m},
$$
because localization gives $S_{\mathfrak p}=R_{\mathfrak p}^m/H_{\mathfrak p}$
and reduction gives that quotient. Thus $d=m$ exactly when
$H_{\mathfrak p}\subseteq\mathfrak pR_{\mathfrak p}^m$, equivalently
$H\subseteq\mathfrak pR^m$. The last equivalence is coordinatewise:
$a/1\in\mathfrak pR_{\mathfrak p}$ if and only if $a\in\mathfrak p$,
by the fraction equality criterion and primality.

The bounds and their limitations can be seen in explicit examples.
For $R=k$, $S=k^m$ and the coordinate idempotent generators,
the fibre has $m$ primes and all inequalities are equalities.
For $R=k[t]$, $S=k[x]$ and $\varphi(t)=x^m$ with $m\geq1$,
the original generators $1,x,\ldots,x^{m-1}$ form a free
$R$-basis: collecting powers by their residue modulo $m$ gives
spanning and uniqueness of the coefficients in $k[t]$.
At $\mathfrak p=(t)$ the fibre is $k[x]/(x^m)$, so it has
one prime, residue degree $1$, multiplicity $m$ and nilradical
dimension $m-1$, although both original rings are domains.
For the noninjective map $k[t]\to k$, $t\mapsto0$, with generator
$1$, the fibre at $(t-1)$ is zero, while the fibre at $(t)$ is $k$.
Repeating the generator $1$ for $R=S=k$ gives $d=1<m$ whenever
the supplied family has $m>1$ entries.

Equality in the residue-degree bound does not imply separability.
In characteristic $\ell>0$, let $K=\mathbb F_\ell(t)$ and
$L=\mathbb F_\ell(u)$ with $t\mapsto u^\ell$.
The elements $1,u,\ldots,u^{\ell-1}$ are linearly independent
over $K$: clearing denominators in any relation gives
$\sum_{j=0}^{\ell-1}u^jP_j(u^\ell)=0$, whose powers in distinct
residue classes modulo $\ell$ cannot cancel. Their span is closed
under multiplication since $u^\ell=t$, and is a finite dimensional
domain over $K$, hence a field. It contains $u$ and therefore is
the whole $L$, proving $[L:K]=\ell$ and
$L\cong K[X]/(X^\ell-t)$. Its unique residue field has degree
equal to its dimension and it is reduced. But
$$
L\otimes_KL\cong L[X]/(X^\ell-u^\ell)
=L[X]/((X-u)^\ell)
$$
is not reduced: $X-u$ is a nonzero element of nilpotence order
$\ell$, as its degree is smaller than that of the defining monic
polynomial. Thus the equality asserts reducedness of the specified
fibre, not reducedness after every field extension.

Finally every finite choice of positive integers $f_i,\mu_i$
is realized, showing the numerical formula is sharp. Take
$K=F(t_0)$ for a field $F$ and
$L_i=F(u_i)$ with $t_0\mapsto u_i^{f_i}$.
The same distinct-residue-class argument proves
$[L_i:K]=f_i$, with basis $1,u_i,\ldots,u_i^{f_i-1}$;
the spanning ring is a finite dimensional domain and hence a field
containing $u_i$, so the argument applies to the full rational
function field. Set
$$
S=\prod_i L_i[\varepsilon_i]/(\varepsilon_i^{\mu_i}).
$$
Its $K$-basis consists of the elements supported in factor $i$
with entries $u_i^a\varepsilon_i^b$ for
$0\leq a<f_i$, $0\leq b<\mu_i$. The nilpotent ideal in each
factor is $(\varepsilon_i)$, its quotient is the field $L_i$,
and its successive powers have one dimensional quotients over
$L_i$ until exponent $\mu_i$. Thus its unique prime has residue
degree $f_i$ and local length $\mu_i$. A prime of the product
selects exactly one factor because the coordinate idempotents
are orthogonal and sum to $1$. Consequently the fibre over the
zero prime of $K$ has exactly these degrees and multiplicities.
The displayed basis has size $\sum_i\mu_if_i$; appending zeros
gives a specified generating family of any larger size $m$.
Thus positivity and $\sum_i\mu_if_i\leq m$ are the complete
universal numerical restrictions on these data.
\end{proof}


\subsection{Editorial treatment of Algebra lines 7852--7868}

\hypertarget{algebra-context-086}{}

\noindent\href{https://github.com/stacks/stacks-project/blob/a04446e57ec1fbc252a871afcec7752fb2807b14/algebra.tex\#L7852}{Original source, lines 7852--7868}. Complete original and reviewed TeX passages accompany this note. Every intervention is separately recorded; the following treatment is editorial.


\begin{lemma}
\label{editorial-source-lemma-integral-going-up}
Let $R \to S$ be a ring map such that
$S$ is integral over $R$.
Let $\mathfrak p \subset \mathfrak p' \subset R$
be primes. Let $\mathfrak q$ be a prime of $S$ mapping
to $\mathfrak p$. Then there exists a prime $\mathfrak q'$
with $\mathfrak q \subset \mathfrak q'$
mapping to $\mathfrak p'$.
\end{lemma}

\begin{proof}
Keep the original map $\varphi:R\to S$ and form
$$
\overline\varphi:R/\mathfrak p\longrightarrow S/\mathfrak q,
\qquad r+\mathfrak p\longmapsto\varphi(r)+\mathfrak q.
$$
It is well defined and injective because
$\varphi^{-1}(\mathfrak q)=\mathfrak p$; both rings are domains.
It is integral because the monic equation for each original
$s\in S$ descends coefficient by coefficient to its class.
Lemma \ref{editorial-source-lemma-integral-overring-surjective} therefore gives
a prime $\overline{\mathfrak q'}$ of $S/\mathfrak q$ contracting
to $\mathfrak p'/\mathfrak p$. Let $\mathfrak q'$ be its inverse
image in $S$. This is prime, contains the original $\mathfrak q$,
and its contraction is exactly $\mathfrak p'$: an element $r\in R$
maps into it if and only if $r+\mathfrak p\in\mathfrak p'/\mathfrak p$,
which is equivalent to $r\in\mathfrak p'$.
\end{proof}


\subsection{Editorial treatment of Algebra lines 7875--7925}

\hypertarget{algebra-context-087}{}

\noindent\href{https://github.com/stacks/stacks-project/blob/a04446e57ec1fbc252a871afcec7752fb2807b14/algebra.tex\#L7875}{Original source, lines 7875--7925}. Complete original and reviewed TeX passages accompany this note. Every intervention is separately recorded; the following treatment is editorial.


\begin{lemma}
\label{editorial-source-lemma-finite-finitely-presented-extension}
Let $R \to S$ be a finite and finitely presented ring map.
Let $M$ be an $S$-module.
Then $M$ is finitely presented as an $R$-module if and only if
$M$ is finitely presented as an $S$-module.
\end{lemma}

\begin{proof}
If $M$ is finitely presented over $R$, Lemma
\ref{algebra-lemma-finitely-presented-over-subring} applies: the given finite
map is of finite type, because its module generators also generate
its algebra. Hence $M$ is finitely presented over $S$.
To see the other assume that $M$ is finitely presented as an $S$-module.
Pick a presentation
$$
S^{\oplus m} \longrightarrow
S^{\oplus n} \longrightarrow
M \longrightarrow 0
$$
The kernel of $S^{\oplus n}\to M$ is the image of the displayed
$S^{\oplus m}$. A finite $R$-generating family of $S$ in each of
these $m$ summands maps to a finite $R$-generating family of that
kernel. Thus from
Lemma \ref{algebra-lemma-extension}
we see that it suffices to prove that $S$ is finitely presented as an
$R$-module.

\medskip\noindent
Pick $y_1, \ldots, y_n \in S$ such that $y_1, \ldots, y_n$ generate $S$
as an $R$-module. By Lemma \ref{editorial-source-lemma-characterize-integral-element}
each $y_i$ is integral over $R$. Choose monic polynomials
$P_i(x) \in R[x]$ with $P_i(y_i) = 0$. Consider the ring
$$
S' = R[x_1, \ldots, x_n]/(P_1(x_1), \ldots, P_n(x_n))
$$
Then we see that $S$ is of finite presentation as an $S'$-algebra
by Lemma \ref{algebra-lemma-compose-finite-type}. Since $S' \to S$ is surjective,
the kernel $J = \Ker(S' \to S)$ is finitely generated as an ideal by
Lemma \ref{algebra-lemma-finite-presentation-independent}. Hence $J$ is a finite
$S'$-module (immediate from the definitions).
Thus $S = \Coker(J \to S')$ is  of finite presentation as an $S'$-module
by Lemma \ref{algebra-lemma-extension}.
Hence, arguing as in the first paragraph, it suffices to show that $S'$ is
of finite presentation as an $R$-module. Actually, $S'$ is free as an
$R$-module with basis the monomials $x_1^{e_1} \ldots x_n^{e_n}$
for $0 \leq e_i < \deg(P_i)$. Namely, write $R \to S'$ as the composition
$$
R \to R[x_1]/(P_1(x_1)) \to R[x_1, x_2]/(P_1(x_1), P_2(x_2)) \to
\ldots \to S'
$$
This shows that the $i$th ring in this sequence is free as a module over the
$(i - 1)$st one with basis $1, x_i, \ldots, x_i^{\deg(P_i) - 1}$. The result
follows by induction, with the following division argument proving
both spanning and independence over every intermediate coefficient
ring. If $P_i(x_i)=x_i^{d_i}+\sum_{j<d_i}a_{ij}x_i^j$ is monic,
subtract the leading coefficient times $x_i^{e-d_i}P_i$ from
a polynomial of degree $e\geq d_i$. Repeating decreases the degree
and leaves a remainder of degree less than $d_i$, retaining each
coefficient $a_{ij}$. Uniqueness holds because for nonzero $Q$,
the product $P_iQ$ has leading coefficient equal to that of $Q$
and degree $d_i+\deg Q$, even if the coefficient ring has
zero divisors. Thus no nonzero remainder of degree less than
$d_i$ belongs to $(P_i)$. Choosing $d_i\geq1$ is always possible;
a monic equation of degree zero can be multiplied by $x_i$.
Applying the unique remainder successively in $x_1,\ldots,x_n$
gives exactly the displayed monomial basis of $S'$ over $R$,
of size $D=\prod_i d_i$ (the empty product is $1$).

To spell out the final presentations, let $h_1,\ldots,h_a$
generate the original ideal $J\subset S'$. The $aD$ elements
$h_j x_1^{e_1}\cdots x_n^{e_n}$, with the displayed exponent
bounds, generate $J$ as an $R$-module. Indeed expand each
coefficient in an ideal expression $\sum_j c_jh_j$ uniquely
in that monomial basis. Their coordinate vectors give a map
$R^{\oplus aD}\to R^{\oplus D}$ with cokernel the original
$S=S'/J$. This is an explicit finite $R$-presentation of $S$.
For the original $S^{\oplus n}$ of the first paragraph, take
one copy of this presentation per summand. That number of summands
is independent of the number of chosen $y_i$ in this paragraph.
Lift the finite $R$-generating family
of the kernel there to this finite free $R$-module. The lifted
vectors together with its existing finitely many relation vectors
generate the full kernel onto $M$: subtracting the lifts from
any vector mapping to zero leaves an old relation. This supplies
a finite $R$-presentation of the same $M$ and finishes the proof.
\end{proof}


\subsection{Editorial treatment of Algebra lines 7927--7963}

\hypertarget{algebra-context-088}{}

\noindent\href{https://github.com/stacks/stacks-project/blob/a04446e57ec1fbc252a871afcec7752fb2807b14/algebra.tex\#L7927}{Original source, lines 7927--7963}. Complete original and reviewed TeX passages accompany this note. Every intervention is separately recorded; the following treatment is editorial.


\begin{lemma}
\label{editorial-source-lemma-silly-normal}
Let $R$ be a ring. Let $x, y \in R$ be nonzerodivisors.
Let $R[x/y] \subset R_{xy}$ be the $R$-subalgebra generated
by $x/y$, and similarly for the subalgebras $R[y/x]$ and $R[x/y, y/x]$.
If $R$ is integrally closed in $R_x$ or $R_y$, then the sequence
$$
0 \to R \xrightarrow{(-1, 1)} R[x/y] \oplus R[y/x] \xrightarrow{(1, 1)}
R[x/y, y/x] \to 0
$$
is a short exact sequence of $R$-modules.
\end{lemma}

\begin{proof}
Because $x$ and $y$ are nonzerodivisors, their product is a
nonzerodivisor. Thus $R\to R_{xy}$ is injective, and so are the
maps $R_x\to R_{xy}$ and $R_y\to R_{xy}$: a fraction that becomes
zero has its numerator killed by a power of $xy$, hence has zero
numerator. All subrings in the statement therefore have their
specified inclusions in $R_{xy}$. The product of the original
elements $x/y$ and $y/x$ is $1$. Consequently each monomial
$(x/y)^a(y/x)^b$ with $a,b\geq0$ equals $(x/y)^{a-b}$ when
$a\geq b$, and equals $(y/x)^{b-a}$ otherwise. Splitting a
polynomial into these terms expresses it as a sum of elements
of $R[x/y]$ and $R[y/x]$, proving surjectivity of the map
$(u,v)\mapsto u+v$. Its kernel consists exactly of pairs
$(-\alpha,\alpha)$ with $\alpha\in R[x/y]\cap R[y/x]$.
The first map in the statement is injective and sends
$a\mapsto(-a,a)$, with the original signs.
Let $\alpha \in R[x/y] \cap R[y/x]$. To show exactness in the middle
we have to prove that $\alpha \in R$. By assumption we may write
$$
\alpha = a_0 + a_1 x/y + \ldots + a_n (x/y)^n =
b_0 + b_1 y/x + \ldots + b_m(y/x)^m
$$
for some $n, m \geq 0$ and $a_i, b_j \in R$.
Pick some $N > \max(n, m)$.
Consider the finite $R$-submodule $M$ of $R_{xy}$ generated by the elements
$$
(x/y)^N, (x/y)^{N - 1}, \ldots, x/y, 1, y/x, \ldots, (y/x)^{N - 1}, (y/x)^N
$$
We check $\alpha M\subset M$ using the full original expressions.
For $0\leq i\leq N$, multiply $(x/y)^i$ by
$\sum_{j=0}^m b_j(y/x)^j$. If $i\geq j$, the resulting
monomial is $(x/y)^{i-j}$ with $0\leq i-j\leq N$; if $i<j$,
it is $(y/x)^{j-i}$ with $0<j-i\leq m<N$.
Thus every summand, with its original coefficient $b_j$, belongs
to $M$. Multiplying $(y/x)^i$ instead by
$\sum_{j=0}^n a_j(x/y)^j$ gives $(y/x)^{i-j}$ for $i\geq j$
and $(x/y)^{j-i}$ for $i<j$, with the same bounds using $n<N$.
This proves stability for every listed generator. Since $1\in M$,
Lemma \ref{editorial-source-lemma-characterize-integral-element} proves $\alpha$
integral over $R$. The expression with coefficients $b_j$ places
$\alpha$ in the original $R_x$, and the expression with $a_j$
places it in $R_y$. Either of the two integral-closedness
hypotheses therefore gives $\alpha\in R$, proving middle exactness.

The calculation also identifies precisely what is needed without
either integral-closedness assumption. With the same original
nonzerodivisors, let $C=R[x/y]\cap R[y/x]\subset R_{xy}$.
It is a subring containing $R$, and the preceding finite-submodule
argument proves every element of $C$ integral over $R$.
There is always an exact sequence with its actual intersection
$$
0\longrightarrow C\xrightarrow{\alpha\mapsto(-\alpha,\alpha)}
R[x/y]\oplus R[y/x]\xrightarrow{(u,v)\mapsto u+v}
R[x/y,y/x]\longrightarrow0.
$$
Surjectivity and the kernel were proved above, and injection uses
the original inclusions. In the sequence with $R$ at the left,
the middle homology is exactly $C/R$: the map sends
$\alpha+R$ to the class of $(-\alpha,\alpha)$, is well defined
by the original map $R\to R[x/y]\oplus R[y/x]$, and has inverse
given by the second coordinate of a kernel pair modulo $R$.
Both composites are identities. Thus the original sequence is
exact precisely when this specified integral extension $C$ equals
$R$; the two stated hypotheses are sufficient conditions for that
equality, and the formula retains the precise defect otherwise.
\end{proof}


\subsection{Editorial treatment of Algebra lines 7971--8339}

\hypertarget{algebra-context-089}{}

\noindent\href{https://github.com/stacks/stacks-project/blob/a04446e57ec1fbc252a871afcec7752fb2807b14/algebra.tex\#L7971}{Original source, lines 7971--8339}. Complete original and reviewed TeX passages accompany this note. Every intervention is separately recorded; the following treatment is editorial.


\subsubsection{Normal rings}
\label{editorial-source-section-normal-rings}

\noindent
We first introduce the notion of a normal domain, and then we
introduce the (very general) notion of a normal ring.

\begin{definition}
\label{editorial-source-definition-domain-normal}
A domain $R$ is called {\it normal} if it is integrally
closed in its field of fractions.
\end{definition}

\begin{lemma}
\label{editorial-source-lemma-integral-closure-in-normal}
Let $R \to S$ be a ring map.
If $S$ is a normal domain, then the integral closure of $R$
in $S$ is a normal domain. It also equals the integral closure
of $R$ in the fraction field of the original $S$.
\end{lemma}

\begin{proof}
Write $B\subset S$ for that integral closure, retaining the
original map $R\to S$. The ring $B$ is a domain, since it is
a subring of the domain $S$. The inclusion induces the injective
map $\operatorname{Frac}(B)\to\operatorname{Frac}(S)$,
$a/b\mapsto a/b$; denominators remain nonzero and equality of
fractions is the same equality $ab'=a'b$ in either domain.
If $z\in\operatorname{Frac}(B)$ is integral over $B$, its monic
equation is also a monic equation over $S$. Normality of $S$
therefore places the image of $z$ in $S$. The subalgebra $B[z]$
is integral over $B$, and $R\to B$ is integral by the definition
of $B$. Transitivity (Lemma \ref{algebra-lemma-integral-transitive}) makes
$z$ integral over $R$, so $z\in B$. This proves normality of $B$.
If instead $z\in\operatorname{Frac}(S)$ is integral over $R$,
its same equation, with coefficients mapped into $S$, makes it
integral over $S$. Thus $z\in S$ and then $z\in B$.
The reverse inclusion follows from the original inclusion
$B\subset S\subset\operatorname{Frac}(S)$ and its monic equations,
proving the asserted equality of integral closures.
\end{proof}

\noindent
The following notion is occasionally useful when
studying normality.

\begin{definition}
\label{editorial-source-definition-almost-integral}
Let $R$ be a domain.
\begin{enumerate}
\item An element $g$ of the fraction
field of $R$ is called {\it almost integral over $R$}
if there exists an element $r \in R$, $r\not = 0$
such that $rg^n \in R$ for all $n \geq 0$.
\item The domain $R$ is called {\it completely normal} if every
almost integral element of the fraction field of $R$ is
contained in $R$.
\end{enumerate}
\end{definition}

\noindent
The following lemma shows that a Noetherian domain is normal
if and only if it is completely normal.

\begin{lemma}
\label{editorial-source-lemma-almost-integral}
Let $R$ be a domain with fraction field $K$.
If $u, v \in K$ are almost integral over $R$, then so are
$u + v$ and $uv$. Any element $g \in K$ which is integral over $R$
is almost integral over $R$. If $R$ is Noetherian
then the converse holds as well. More precisely, for an almost
integral $g$ and any original witness $0\ne r\in R$, integrality
of $g$ is equivalent to finite generation of the specified ideal
$I_g=\sum_{n\geq0}R\,rg^n\subset R$.
\end{lemma}

\begin{proof}
Keep witnesses $0\ne r,r'\in R$ with $ru^n\in R$ and
$r'v^n\in R$ for all $n\geq0$. Their product is nonzero, and
the full identities
$$
rr'(uv)^n=(ru^n)(r'v^n),\qquad
rr'(u+v)^n=\sum_{j=0}^n\binom nj(ru^j)(r'v^{n-j})
$$
place both expressions in $R$ for every $n\geq0$. The binomial
coefficients are integers acting through $\mathbf Z\to R$;
no division in $R$ is used. Also $r(-u)^n=(-1)^nru^n\in R$,
and every element of $R$ has witness $1$. Consequently the
almost integral elements form an $R$-subalgebra of the original $K$.
Suppose $g \in K$ is integral over $R$.
In this case there exists a $d > 0$ such that
the ring $R[g]$ is generated by $1, g, \ldots, g^d$ as an $R$-module.
Write each original $g^i=a_i/b_i$ for $0\leq i\leq d$,
with $a_i,b_i\in R$ and $b_i\ne0$, and choose the actual common
denominator $r=\prod_{i=0}^d b_i\ne0$. Then
$rg^i=a_i\prod_{j\ne i}b_j\in R$ for each listed power.
Every element of $R[g]$ is an $R$-linear combination of these
powers, so multiplying that combination by $r$ gives an element
of $R$. In particular $rg^n\in R$ for all $n\geq0$,
which proves that the original $g$ is almost integral.

\medskip\noindent
Let $g\in K$ be almost integral with its original witness
$0\ne r\in R$. The set
$I_g=\sum_{n\geq0}R\,rg^n=rR[g]\subset R$ is an ideal.
Multiplication by $r$ gives the exact $R$-module isomorphism
$$
R[g]\longrightarrow I_g,\quad z\longmapsto rz,
\qquad I_g\longrightarrow R[g],\quad w\longmapsto w/r.
$$
The second map lands in $R[g]$ by the displayed description of
$I_g$, and both composites are identities in the original field $K$.
If $I_g$ is finitely generated, so is $R[g]$, and Lemma
\ref{algebra-lemma-finite-is-integral} proves $g$ integral. Conversely,
if $g$ is integral, Lemma \ref{editorial-source-lemma-characterize-integral}
makes $R[g]$ finite, so the same isomorphism makes $I_g$ finite
for every choice of witness. This proves the exact criterion.
When $R$ is Noetherian the ideal $I_g$ is finitely generated,
giving the claimed converse. Equivalently, the original inclusion
$R[g]\subset(1/r)R$ identifies it with a submodule of the finite
$R$-module $(1/r)R$; multiplication by $r$ identifies this
submodule with the very same ideal $I_g$.
\end{proof}

\begin{lemma}
\label{editorial-source-lemma-localize-normal-domain}
If $R$ is a normal domain and $S\subset R$ is multiplicative
with $0\notin S$, then $S^{-1}R$ is a normal domain.
If $0\in S$, the localization is the zero ring, which is a normal
ring in the sense of Definition \ref{editorial-source-definition-ring-normal}.
\end{lemma}

\begin{proof}
If $0\in S$, every fraction is zero and the spectrum is empty,
so the later definition of a normal ring applies vacuously.
Suppose $0\notin S$ and keep $K=\operatorname{Frac}(R)$.
The original localization injects into $K$ by $a/s\mapsto a/s$.
Its fraction field is identified with $K$ by the inverse maps
$$
\frac{a/s}{b/t}\longmapsto\frac{at}{sb},\qquad
\frac ab\longmapsto\frac{a/1}{b/1}.
$$
All denominators are nonzero; fraction equality in a domain
verifies well-definedness, and multiplication of the original
numerators and denominators verifies both composites.
Let $g\in K$ be integral over $S^{-1}R$, with its monic equation
$P(g)=0$ for $P(x)=x^d+\sum_{0\leq j<d}a_jx^j$, $d\geq1$.
Write $a_j=c_j/s_j$ with $c_j\in R$, $s_j\in S$, and keep
$s=\prod_{j=0}^{d-1}s_j\in S$.
Then $sa_j=c_j\prod_{h\ne j}s_h\in R$, and
$s^{d-j}a_j=s^{d-j-1}(sa_j)\in R$ for every $0\leq j<d$.
The full scaled equation is
$$
(sg)^d+\sum_{j=0}^{d-1}s^{d-j}a_j(sg)^j
=s^d\left(g^d+\sum_{j=0}^{d-1}a_jg^j\right)=0.
$$
It is monic with coefficients in $R$, so normality gives
$sg\in R$. Thus the original $g$ is the image of $(sg)/s$
in $S^{-1}R$, proving the assertion.
\end{proof}

\begin{lemma}
\label{editorial-source-lemma-PID-normal}
A principal ideal domain is normal.
\end{lemma}

\begin{proof}
Let $R$ be a principal ideal domain and retain the original
$g=a/b\in\operatorname{Frac}(R)$, where $b\ne0$, and the full
monic equation
$$
(a/b)^n+\sum_{i=0}^{n-1}r_i(a/b)^i=0,\qquad r_i\in R.
$$
Multiplication by the original $b^n$ gives
$a^n+\sum_{i=0}^{n-1}r_i a^i b^{n-i}=0$ in $R$.
Choose $c$ generating the original ideal $(a,b)$, and write
$a=ca'$, $b=cb'$, $c=ua+vb$ for $a',b',u,v\in R$.
Since $b\ne0$, both $c$ and $b'$ are nonzero. Cancellation in
$c=c(ua'+vb')$ gives $ua'+vb'=1$.
The full original equation, after these substitutions, is
$$
c^n\left((a')^n+\sum_{i=0}^{n-1}r_i(a')^i(b')^{n-i}\right)=0.
$$
The nonzero factor $c^n$ in this domain makes the parenthesized
expression zero. Expanding the original identity
$1=(ua'+vb')^n$ and substituting that expression gives $1=b'w$,
where the entire inverse is
$$
w=-u^n\sum_{i=0}^{n-1}r_i(a')^i(b')^{n-i-1}
+\sum_{j=0}^{n-1}\binom nj(ua')^jv^{n-j}(b')^{n-j-1}\in R.
$$
Indeed the omitted top term of the latter binomial sum is
$u^n(a')^n$, exactly the term replaced using the monic equation.
Returning to the original numerator and denominator gives
$b(a'w)=cb'a'w=ca'=a$. Hence $g=a/b=a'w\in R$.
Every factor from the original fraction and every coefficient
of its monic equation is retained in this comparison.
\end{proof}

\begin{lemma}
\label{editorial-source-lemma-prepare-polynomial-ring-normal}
Let $R$ be a domain with fraction field $K$.
Suppose $f = \sum \alpha_i x^i$ is an
element of $K[x]$.
\begin{enumerate}
\item $f$ is integral over $R[x]$ if and only if
all $\alpha_i$ are integral over $R$, and
\item $f$ is almost integral over $R[x]$ if and only if
all $\alpha_i$ are almost integral over $R$.
\end{enumerate}
\end{lemma}

\begin{proof}
If $f=0$, all its coefficients are zero and the assertions hold
directly. Suppose $f\ne0$. We first prove the forward implication
in the second statement, keeping all coefficients, including zeros.
Write $f = \alpha_0 + \alpha_1 x + \ldots + \alpha_r x^r$
with $\alpha_r \not = 0$.
By assumption there exists $h = b_0 + b_1 x + \ldots + b_s x^s \in R[x]$,
$b_s \not = 0$ such that $f^n h \in R[x]$ for all
$n \geq 0$. This implies that $b_s \alpha_r^n \in R$
for all $n \geq 0$. Hence $\alpha_r$ is almost
integral over $R$. Since the set of almost integral
elements forms a subring (Lemma \ref{editorial-source-lemma-almost-integral}) we deduce that
$f - \alpha_r x^r = \alpha_0 + \alpha_1 x + \ldots + \alpha_{r - 1} x^{r - 1}$
is almost integral over $R[x]$. Here $\alpha_rx^r$ has the
actual witness $b_s$, since
$b_s(\alpha_rx^r)^n=(b_s\alpha_r^n)x^{rn}\in R[x]$.
The product $b_sh$ is consequently a witness for their difference
by the binomial calculation in Lemma \ref{editorial-source-lemma-almost-integral}.
The induction can be made fully explicit. For $0\leq t\leq r$
keep $f_t=\sum_{i=0}^t\alpha_ix^i$ and set
$$
h_t=b_s^{2^{r-t}-1}h,\qquad c_t=b_s^{2^{r-t}}.
$$
The original hypothesis says $h_rf_r^n\in R[x]$ for all $n\geq0$.
If $h_tf_t^n\in R[x]$, its coefficient of $x^{tn+s}$ is
$c_t\alpha_t^n$, also when $\alpha_t=0$. Hence this coefficient
belongs to $R$. The full binomial formula for
$f_{t-1}=f_t-\alpha_tx^t$ shows
$$
h_tc_t f_{t-1}^n
=\sum_{j=0}^n\binom nj(h_tf_t^j)
 \bigl(c_t(-\alpha_tx^t)^{n-j}\bigr)\in R[x].
$$
Since $h_tc_t=h_{t-1}$, this proves the descending induction.
In particular each original $\alpha_t$ has nonzero witness $c_t$,
and the single nonzero element $b_s^{2^r}$ is a witness for
all coefficients, by multiplying their equations by the remaining
power of $b_s$.

Conversely, if $\alpha_i$ has witness $0\ne w_i\in R$ for
$0\leq i\leq r$, put $w=\prod_iw_i\ne0$. For every $n\geq0$
the complete multinomial expansion gives
$$
w f^n=\sum_{k_0+\cdots+k_r=n}
\binom{n}{k_0,\ldots,k_r}
\left(\prod_{i=0}^r w_i\alpha_i^{k_i}\right)
x^{\sum_{i=0}^r i k_i}\in R[x].
$$
The indices $k_i$ are nonnegative integers and the multinomial
coefficients are integers, not inverses of factorials in $R$.
This proves the reverse almost-integral implication.

\medskip\noindent
In order to prove the first statement we will use absolute Noetherian
reduction. Namely, write $\alpha_i = a_i / b_i$ and
let $P(t) = t^d + \sum_{j < d} f_j t^j$ be a polynomial
with coefficients $f_j \in R[x]$ such that $P(f) = 0$.
Let $f_j = \sum f_{ji}x^i$. Consider the subring
$R_0 \subset R$ generated by the finite list of elements
$a_i, b_i, f_{ji}$ of $R$. It is a domain; let
$K_0$ be its field of fractions. Since $R_0$ is a finite type
$\mathbf{Z}$-algebra it is Noetherian, see
Lemma \ref{editorial-source-lemma-obvious-Noetherian}. It is still
the case that $f \in K_0[x]$ is integral over $R_0[x]$,
because the inclusion $R_0\subset R$ induces the injective maps
$K_0\to K$, $a/b\mapsto a/b$, and $K_0[x]\to K[x]$.
Each original denominator $b_i$ is nonzero in $R_0$, all the
original coefficients $f_{ji}$ belong to it, and the same monic
polynomial $P$ and same $f$ have $P(f)\in K_0[x]$.
Its image is zero in $K[x]$, so injectivity gives $P(f)=0$
already in $K_0[x]$.
By Lemma \ref{editorial-source-lemma-almost-integral} the element
$f$ is almost integral over $R_0[x]$. By the second statement of
the lemma, the elements $\alpha_i$ are almost integral
over $R_0$. And since $R_0$ is Noetherian, they are
integral over $R_0$, see Lemma \ref{editorial-source-lemma-almost-integral}.
Each resulting monic equation over $R_0$ maps coefficient by
coefficient to a monic equation for the same $\alpha_i$ over $R$.
For the reverse implication, if all the finitely many $\alpha_i$
are integral over $R$, the original algebra
$B=R[\alpha_0,\ldots,\alpha_r]\subset K$ is finite over $R$
by Lemma \ref{editorial-source-lemma-characterize-integral}. A finite module
generating family of $B$ also generates $B[x]$ over $R[x]$:
express each coefficient of a polynomial in that family.
Thus the original $f\in B[x]$ is integral over $R[x]$ by
Lemma \ref{algebra-lemma-finite-is-integral}, proving the other implication.
\end{proof}

\begin{lemma}
\label{editorial-source-lemma-polynomial-domain-normal}
Let $R$ be a normal domain.
Then $R[x]$ is a normal domain. For any domain $R$, the polynomial
ring $R[x]$ is completely normal if and only if $R$ is completely normal.
\end{lemma}

\begin{proof}
The result is true if $R$ is a field $K$ because
$K[x]$ is a euclidean domain and hence a principal ideal
domain and hence normal by Lemma \ref{editorial-source-lemma-PID-normal}.
The inclusions $R[x]\subset K[x]$ identify their fraction fields
with $K(x)$. More explicitly a fraction of polynomials over $R$
maps to the same fraction over $K$. Conversely, for $A/B$ with
$A,B\in K[x]$, choose a common nonzero denominator $c\in R$
of every coefficient of both polynomials. Its inverse image is
$(cA)/(cB)$, with both original scaled polynomials in $R[x]$.
The field equality criterion verifies independence of the chosen
common denominator and both composites.
Let $g$ be an element of the fraction field of
$R[x]$ which is integral over $R[x]$. Because $g$
is integral over $K[x]$ where $K$ is the fraction
field of $R$ we may write $g = \alpha_d x^d + \alpha_{d-1}x^{d-1} +
\ldots + \alpha_0$ with $\alpha_i \in K$.
By Lemma \ref{editorial-source-lemma-prepare-polynomial-ring-normal}
the elements $\alpha_i$ are integral over $R$ and
hence are in $R$. Thus the original $g$ belongs to $R[x]$.

For the second assertion, suppose $R$ is completely normal and
$g\in K(x)$ is almost integral over $R[x]$. Its same nonzero
polynomial witness also makes it almost integral over $K[x]$.
The ring $K[x]$ is Noetherian, since it is a principal ideal
domain, so Lemma \ref{editorial-source-lemma-almost-integral} makes $g$ integral
over $K[x]$. Normality of $K[x]$ then gives $g\in K[x]$.
The proved coefficientwise almost-integral equivalence shows
that every coefficient of this polynomial is almost integral
over $R$, hence lies in $R$. Therefore $g\in R[x]$.
Conversely, an almost integral element $g\in K$ and its witness
$0\ne r\in R$ remain the same constant fraction and witness
over $R[x]$. Complete normality of $R[x]$ places $g$ in
$R[x]\cap K=R$, where equality follows from uniqueness of
polynomial coefficients. This proves complete normality of $R$.
\end{proof}

\begin{lemma}
\label{editorial-source-lemma-power-series-over-Noetherian-normal-domain}
Let $R$ be a Noetherian normal domain. Then $R[[x]]$ is
a Noetherian normal domain. More generally, for any domain $R$,
$R[[x]]$ is completely normal if and only if $R$ is completely normal.
\end{lemma}

\begin{proof}
First prove the second assertion in the forward direction.
Assume that $R$ is completely normal, and retain nonzero
$f,g\in R[[x]]$ and the original $w=f/g$ in its fraction field,
now assumed almost integral. The power series ring is a domain:
the first nonzero coefficients of two nonzero series multiply
to the first nonzero coefficient of their product. Coefficient
inclusion gives an injection $R[[x]]\to K[[x]]$, where
$K=\operatorname{Frac}(R)$, hence an injection of its fraction
field into $K((x))$. A nonzero series over $K$ is $x^a$ times
a series with nonzero constant coefficient. For a series
$x^a\sum_{j\geq0}c_jx^j$ with $c_0\ne0$, its inverse is
$x^{-a}\sum_{j\geq0}d_jx^j$, where
$d_0=c_0^{-1}$ and
$d_j=-c_0^{-1}\sum_{i=1}^j c_id_{j-i}$ for $j\geq1$.
The coefficient of degree zero in the product is $1$ and every
higher coefficient is zero by this recurrence, proving the
inverse formula in $K((x))$. Thus the
original $w$ has the unique Laurent expansion
$$
w=a_nx^n+a_{n+1}x^{n+1}+\cdots,
\qquad n\in\mathbf Z,\quad a_n\ne0.
$$
Keep its original nonzero almost-integrality witness
$$
h=b_mx^m+b_{m+1}x^{m+1}+\cdots\in R[[x]],
\qquad b_m\ne0,
$$
so $hw^e\in R[[x]]$ for every $e\geq0$.
The first nonzero term of $hw^e$ has degree $m+en$ and
coefficient $b_ma_n^e$. If $n<0$, choose $e$ with $m+en<0$;
that nonzero negative-degree term contradicts membership in
$R[[x]]$. Hence $n\geq0$. For every $e\geq0$ the displayed
coefficient belongs to $R$, so complete normality gives $a_n\in R$.

Suppose all original coefficients $a_n,\ldots,a_{N-1}$ have
been proved to belong to $R$, retaining zeros in this list, and
put $P_N=\sum_{i=n}^{N-1}a_ix^i\in R[x]$.
The same original $h$ is a witness for $w-P_N$, since the full
identity
$$
h(w-P_N)^e=\sum_{j=0}^e\binom ej(hw^j)(-P_N)^{e-j}
$$
has every term in $R[[x]]$. The coefficient of $x^{m+eN}$
on the left is $b_ma_N^e$, including the case $a_N=0$:
$w-P_N$ has no terms of degree below $N$, so only these leading
terms can contribute at that degree. Thus $b_ma_N^e\in R$
for every $e\geq0$. The original nonzero $b_m$ witnesses
almost integrality of $a_N$ over $R$, giving $a_N\in R$.
Induction proves every coefficient belongs to $R$, hence the
original fraction $w$ belongs to $R[[x]]$. The zero fraction
belongs to the ring directly. This proves complete normality
of $R[[x]]$ without any Noetherian assumption.

Conversely, if $R[[x]]$ is completely normal and $z\in K$ is
almost integral over $R$, its original constant numerator,
denominator and witness define the same almost integral constant
fraction in the fraction field of $R[[x]]$. Thus it lies in
$R[[x]]\cap K=R$: uniqueness of the Laurent expansion forces
all coefficients but the constant one to be zero. Hence $R$
is completely normal, proving the equivalence.

Finally if the original $R$ is Noetherian and normal, Lemma
\ref{editorial-source-lemma-almost-integral} says that it is completely normal.
The argument just given makes $R[[x]]$ completely normal, hence
normal since every integral element is almost integral.
Its Noetherian property follows from
Lemma \ref{editorial-source-lemma-Noetherian-power-series}, proving the first assertion.
\end{proof}

\begin{lemma}
\label{editorial-source-lemma-normality-is-local}
Let $R$ be a domain. The following are equivalent:
\begin{enumerate}
\item The domain $R$ is a normal domain,
\item for every prime $\mathfrak p \subset R$ the local ring
$R_{\mathfrak p}$ is a normal domain, and
\item for every maximal ideal $\mathfrak m$ the ring $R_{\mathfrak m}$
is a normal domain.
\end{enumerate}
\end{lemma}

\begin{proof}
We deduce (1) $\Rightarrow$ (2) from Lemma \ref{editorial-source-lemma-localize-normal-domain}.
The implication (2) $\Rightarrow$ (3) is immediate. The implication
(3) $\Rightarrow$ (1) follows from the fact that for any domain $R$ we have
$$
R = \bigcap\nolimits_{\mathfrak m} R_{\mathfrak m}
$$
inside the original fraction field of $R$. To prove this equality,
let $g$ belong to every $R_{\mathfrak m}$ and retain
$I=\{x\in R:xg\in R\}$. It is an ideal: multiplication by
$g$ respects sums and multiplication by elements of $R$.
For each maximal $\mathfrak m$, the expression $g=a/s$ in
$R_{\mathfrak m}\subset\operatorname{Frac}(R)$ has
$s\notin\mathfrak m$ and gives $sg=a\in R$. Thus $s\in I$
and $I\not\subseteq\mathfrak m$. If $I$ were proper it would
lie in a maximal ideal, a contradiction. Therefore $1\in I$
and $g\in R$. Finally if $g$ is integral over $R$, its original
monic equation remains monic over each $R_{\mathfrak m}$.
Under (3), each of those rings contains $g$, and the proved
intersection identity gives $g\in R$, establishing (1).
\end{proof}

\noindent
Lemma \ref{editorial-source-lemma-normality-is-local} shows that the following definition
is compatible with Definition \ref{editorial-source-definition-domain-normal}. (It is the
definition from EGA -- see \cite[IV, 5.13.5 and 0, 4.1.4]{EGA}.)

\begin{definition}
\label{editorial-source-definition-ring-normal}
A ring $R$ is called {\it normal} if for every prime
$\mathfrak p \subset R$ the localization $R_{\mathfrak p}$ is
a normal domain (see Definition \ref{editorial-source-definition-domain-normal}).
\end{definition}

\noindent
It suffices in this definition to test maximal ideals. Indeed,
if each $R_{\mathfrak m}$ is a normal domain and
$\mathfrak p\subseteq\mathfrak m$, there is the exact isomorphism
$$
(R_{\mathfrak m})_{\mathfrak pR_{\mathfrak m}}
\longrightarrow R_{\mathfrak p},\qquad
\frac{a/s}{b/t}\longmapsto\frac{at}{sb},
$$
with inverse $a/u\mapsto(a/1)/(u/1)$.
Here $s,t\notin\mathfrak m$ and $b\notin\mathfrak p$, so
$sb\notin\mathfrak p$. The prime correspondence for localization
identifies $\mathfrak pR_{\mathfrak m}$ as a proper prime with
contraction $\mathfrak p$. Fraction equality witnesses map to
fraction equality witnesses in these localizations; the displayed
formulas fix every original fraction in both composites.
Lemma \ref{editorial-source-lemma-localize-normal-domain} then makes
$R_{\mathfrak p}$ a normal domain. The converse is part of the
definition, and every prime lies in a maximal ideal.

A normal ring is reduced. More explicitly the natural map
$R\to\prod_{\mathfrak m}R_{\mathfrak m}$ is injective: if
$a\ne0$, its annihilator is a proper ideal and is contained in
a maximal ideal $\mathfrak m$; then $a/1$ cannot be zero there,
since a denominator outside $\mathfrak m$ killing $a$ would
belong to that annihilator. A nilpotent element maps to zero
in every domain $R_{\mathfrak m}$ and hence is zero in $R$.
This also proves the stated subring assertion using all primes
(see Lemma \ref{algebra-lemma-characterize-zero-local}).
For the zero ring the products are empty and all assertions
hold directly.

\begin{lemma}
\label{editorial-source-lemma-normal-ring-integrally-closed}
A normal ring is integrally closed in its total ring of fractions.
\end{lemma}

\begin{proof}
Let $R$ be normal, let $\Sigma\subset R$ be its original
multiplicative set of nonzerodivisors, and keep
$Q(R)=\Sigma^{-1}R$. The map $R\to Q(R)$ is injective,
since any element in its kernel is killed by a member of $\Sigma$.
Let $x\in Q(R)$ be integral over $R$, and retain the ideal
$I=\{f\in R\mid fx\in R\}$.
Fix a prime $\mathfrak p$ and put $A=R_{\mathfrak p}$, a
normal domain. Every $s\in\Sigma$ has nonzero image in $A$:
an equality $s/1=0$ would give $u\notin\mathfrak p$ with
$us=0$, forcing $u=0$, a contradiction.
There is the exact comparison
$$
Q(R)\otimes_RA\cong\Sigma^{-1}A
\lhook\joinrel\longrightarrow\operatorname{Frac}(A).
$$
The first map sends $(r/s)\otimes c$ to $(rc)/(s/1)$;
its inverse sends $c/(s/1)$ to $(1/s)\otimes c$.
The original coefficient maps agree on $R$, so tensor balancing
gives the first map. The second is defined because each $s/1$
maps to the unit $(s/1)\otimes1$ with inverse $(1/s)\otimes1$.
The formulas verify both composites. The last map is injective
since $A$ is a domain and every image of $\Sigma$ is nonzero.
For the original fractions its full formula is
$$
(r/s)\otimes(a/u)\longmapsto
\frac{(ra)/1}{(su)/1},\qquad
s\in\Sigma,\quad u\notin\mathfrak p.
$$
This embeds the indicated localization in the fraction field;
it does not require that every nonzero element of $A$ come from
an original nonzerodivisor of $R$.

The image of $x\otimes1$ satisfies the same monic equation
with coefficients mapped into $A$. Normality of $A$ places
it in $A$. The injective comparison just proved therefore gives
$x\otimes1=a\otimes(1/f)$ in
$Q(R)\otimes_RR_{\mathfrak p}\cong Q(R)_{\mathfrak p}$ for
some original $a,f\in R$, $f\notin\mathfrak p$.
This last isomorphism sends $z\otimes(a/u)$ to $az/u$
and has inverse $z/u\mapsto z\otimes(1/u)$; tensor balancing
and the fraction equality relation verify both maps and composites.
Hence $fx-a$ becomes zero in that localization, so there is
$f'\notin\mathfrak p$ with $f'(fx-a)=0$ in $Q(R)$.
The equality $f'fx=f'a$ has its right side in the embedded $R$,
so $ff'\in I$ and $ff'\notin\mathfrak p$.
As this holds for every prime, $I$ cannot be proper; otherwise
a maximal ideal containing it would contradict this property.
Thus $1\in I$ and the original $x$ lies in $R$.
If $R=0$, its total quotient ring is zero and the assertion
holds without any choice of a prime.
\end{proof}

\begin{lemma}
\label{editorial-source-lemma-localization-normal-ring}
A localization of a normal ring is a normal ring.
\end{lemma}

\begin{proof}
Let $S\subset R$ be the original multiplicative subset and
let $\mathfrak q$ be a prime of $S^{-1}R$. Its contraction
$\mathfrak p$ to $R$ is prime and disjoint from $S$, and
$\mathfrak q=S^{-1}\mathfrak p$. The maps
$$
(S^{-1}R)_{\mathfrak q}\longrightarrow R_{\mathfrak p},
\quad\frac{a/s}{b/t}\longmapsto\frac{at}{sb},
\qquad
R_{\mathfrak p}\longrightarrow(S^{-1}R)_{\mathfrak q},
\quad a/u\longmapsto\frac{a/1}{u/1}
$$
are defined because $s,t\in S$ and $b\notin\mathfrak p$,
while $u\notin\mathfrak p$ in the second formula. The original
maps from $R$ send each required denominator to a unit, hence
extend to these localizations; on fractions the displayed formulas
prove that the two extensions are inverse. Normality of $R$
makes $R_{\mathfrak p}$ a normal domain, proving the assertion
at every prime of $S^{-1}R$. If this localization is zero,
there are no primes to test and it is normal by definition.
\end{proof}

\begin{lemma}
\label{editorial-source-lemma-polynomial-ring-normal}
Let $R$ be a normal ring. Then $R[x]$ is a normal ring.
More generally, the polynomial ring over $R$ in any set of
indeterminates is normal.
\end{lemma}

\begin{proof}
Let $\mathfrak q\subset R[x]$ be prime and let
$\mathfrak p$ be its contraction to the original $R$.
The ring $R_{\mathfrak p}[x]$ is a normal domain by
Lemma \ref{editorial-source-lemma-polynomial-domain-normal}. It is canonically
$(R\setminus\mathfrak p)^{-1}R[x]$: a finite polynomial with
fraction coefficients is represented by multiplying all their
original denominators to obtain one common denominator, and
the inverse sends $F/s$ to the polynomial whose coefficients
are the corresponding fractions. These maps preserve coefficient
equalities and are inverse. The extended ideal
$\widetilde{\mathfrak q}=(R\setminus\mathfrak p)^{-1}\mathfrak q$
is prime, with contraction $\mathfrak q$. The exact local maps are
$$
(R[x])_{\mathfrak q}\longrightarrow
(R_{\mathfrak p}[x])_{\widetilde{\mathfrak q}},
\quad F/G\longmapsto(F/1)/(G/1),
$$
with inverse $(F/s)/(G/t)\mapsto tF/(sG)$.
Here $s,t\in R\setminus\mathfrak p$ and $G\notin\mathfrak q$,
so every displayed denominator is valid. The original coefficient
maps invert the indicated sets, proving well-definedness by the
localization relations, and the fraction formulas verify both
composites. Lemma \ref{editorial-source-lemma-localize-normal-domain} makes this
local ring a normal domain, proving the one-variable assertion.

Induction proves the assertion for any finite set of variables.
For an arbitrary set $J$, retain all its original labeled variables.
The polynomial ring $R[x_j:j\in J]$ is the directed colimit of
$R[x_j:j\in F]$ over the finite subsets $F\subset J$:
every polynomial has finite support, and coefficient equality
already holds in a finite set containing the variables of both
polynomials. The transition maps keep each original variable and
coefficient. Lemma \ref{editorial-source-lemma-colimit-normal-ring} applies to
these normal rings and gives the full assertion. Its proof below
does not use this polynomial-ring assertion.
\end{proof}

\begin{lemma}
\label{editorial-source-lemma-finite-product-normal}
A finite product of normal rings is normal. More generally,
for any set $I$, the product $\prod_{i\in I}R_i$ is normal
if and only if every original factor $R_i$ is normal.
Every ultraproduct of normal rings is normal as well.
\end{lemma}

\begin{proof}
It suffices to show that the product of two normal rings, say $R$ and $S$, is
normal. By Lemma \ref{algebra-lemma-disjoint-decomposition} the prime ideals of
$R\times S$ are of the form $\mathfrak{p}\times S$ and $R\times
\mathfrak{q}$, where $\mathfrak{p}$ and $\mathfrak{q}$ are primes of $R$
and $S$ respectively. Localization yields 
$(R\times S)_{\mathfrak{p}\times S}=R_{\mathfrak{p}}$ which is a normal domain
by assumption. Explicitly the map sends
$(a,b)/(s,t)$ to $a/s$, and its inverse sends $a/s$ to
$(a,0)/(s,1)$. Both are defined because $s\notin\mathfrak p$.
The idempotent $(1,0)$ is a unit in this localization, so $(0,1)$
is zero there, verifying the inverse identities. The corresponding
maps at $R\times\mathfrak q$ use the other coordinate, proving
the finite assertion without any assumption on the factor maps.
\medskip\noindent
For the arbitrary-product assertion we first prove an algebraic
characterization. A ring $B$ is normal if and only if, for every
$n\geq1$ and all $a,b,c_0,\ldots,c_{n-1}\in B$, it satisfies
$$
a^n+\sum_{j=0}^{n-1}c_ja^jb^{n-j}=0
\quad\Longrightarrow\quad a=bd\text{ for some }d\in B.
\qquad(\mathcal D_n)
$$
Here $b$ may be zero or a zero divisor.
Suppose first that $B$ is normal and the displayed equation holds.
At a prime $\mathfrak p$, write $a_{\mathfrak p},b_{\mathfrak p}$
for the original images in $B_{\mathfrak p}$. If
$b_{\mathfrak p}=0$, the equation gives $a_{\mathfrak p}^n=0$,
hence $a_{\mathfrak p}=0$ since this localization is a domain.
If $b_{\mathfrak p}\ne0$, division by $b_{\mathfrak p}^n$
gives the monic equation
$$
(a_{\mathfrak p}/b_{\mathfrak p})^n+
\sum_{j=0}^{n-1}c_{j,\mathfrak p}
 (a_{\mathfrak p}/b_{\mathfrak p})^j=0.
$$
Integral closedness then gives
$a_{\mathfrak p}\in b_{\mathfrak p}B_{\mathfrak p}$ in both cases.
This implies $a\in bB$: otherwise the ideal
$J=\{r\in B:ra\in bB\}$ is proper and lies in a maximal ideal
$\mathfrak p$. Membership in the localized principal ideal gives
$s\notin\mathfrak p$ with $sa\in bB$ (clear the denominator and
retain the further multiplier witnessing equality in the localization).
This contradicts $s\in J\subset\mathfrak p$ and proves
$(\mathcal D_n)$.

Conversely suppose all $(\mathcal D_n)$ hold. Condition
$(\mathcal D_2)$ with $b=0$ and both coefficients zero gives
$z^2=0\Rightarrow z=0$. If $ab=0$, apply $(\mathcal D_2)$
to the exact equation $a^2-a(a+b)=0$ to get
$a=(a+b)d$ for some $d\in B$. Put
$z=(1-d)a=db$. Then
$z^2=((1-d)a)(db)=d(1-d)ab=0$, so $z=0$ and
$$
(1-d)a=0,\qquad db=0,\qquad (1-d)+d=1.
$$
These identities show that every $B_{\mathfrak p}$ is a domain.
Indeed a zero product $(a/u)(b/v)=0$ has a witness
$w\notin\mathfrak p$ with $wab=0$ in $B$. Apply the identities
to $wa$ and $b$. At least one of $d,1-d$ is outside
$\mathfrak p$. If $d$ is outside, $b/v=0$ in the localization;
if $1-d$ is outside, $a/u=0$ since $w$ is also a unit there.
The localization is nonzero because $\mathfrak p$ is proper.

To prove this domain integrally closed, retain a fraction
$y=(a_0/s)/(b_0/t)$ with $s,t\notin\mathfrak p$ and
$b_0/t\ne0$. Set $a=a_0t$ and $b=b_0s$ in the original $B$,
so $b/1\ne0$ and $y=(a/1)/(b/1)$ with these factors recorded.
Suppose its original monic equation is
$$
y^n+\sum_{j=0}^{n-1}(r_j/t_j)y^j=0,
\qquad r_j\in B,\quad t_j\notin\mathfrak p.
$$
Put $u=\prod_{j=0}^{n-1}t_j$ and
$c_j=r_j\prod_{k\ne j}t_k$, retaining every denominator factor.
Then $r_j/t_j=c_j/u$ in the localization. Multiplication by
$(b/1)^n(u/1)^n$ says that the element
$$
F=(ua)^n+\sum_{j=0}^{n-1}c_j u^{n-1-j}(ua)^jb^{n-j}
$$
has zero image in $B_{\mathfrak p}$. Therefore some
$v\notin\mathfrak p$ satisfies $vF=0$ in the original ring.
The full equation
$$
(vua)^n+\sum_{j=0}^{n-1}c_j u^{n-1-j}(vua)^j(vb)^{n-j}
=v^nF=0
$$
holds there, because $n\geq1$. Condition $(\mathcal D_n)$ gives
$vua=vbd$ for some $d\in B$. Since $v/1$ and $u/1$ are units
and $b/1\ne0$ in the localized domain, the original fraction
satisfies $y=d/u\in B_{\mathfrak p}$. This proves the characterization,
including the zero ring, whose prime-local condition is vacuous.

Now let $R=\prod_{i\in I}R_i$ with every $R_i$ normal.
An equation as in $(\mathcal D_n)$ in $R$ means precisely
$$
a_i^n+\sum_{j=0}^{n-1}c_{j,i}a_i^jb_i^{n-j}=0
\quad\text{in every original }R_i.
$$
Choose $d_i\in R_i$ with $a_i=b_id_i$ using the proved
characterization, and put $d=(d_i)_{i\in I}$. It gives $a=bd$
in the original product. Thus $R$ satisfies every $(\mathcal D_n)$
and is normal. Only the equation degree and its finitely many
coefficients are common; no characteristic or finiteness condition
on $I$ or on the factors is used.

We describe the exact prime and localization maps as well,
including the primes of an infinite product that are not obtained
from a single coordinate. For $E\subseteq I$, let $e_E\in R$
have coordinate $1$ on $E$ and $0$ on its complement. For a
prime $\mathfrak p\subset R$, the family
$$
\mathcal U=\{E\subseteq I:e_E\notin\mathfrak p\}
$$
is an ultrafilter, that is, a family closed under finite intersections
and enlargement, containing $I$ but not the empty set, and containing
exactly one of each subset and its complement. These facts follow
respectively from
$$
e_Ee_F=e_{E\cap F},\quad
e_E=e_Ee_F\ (E\subseteq F),\quad
e_I=1,\quad e_\varnothing=0,
$$
and from $e_Ee_{I\setminus E}=0$,
$e_E+e_{I\setminus E}=1$, using primality.
Set
$$
K_{\mathcal U}=\{a=(a_i):\{i:a_i=0\}\in\mathcal U\},
\qquad A_{\mathcal U}=R/K_{\mathcal U},
\qquad\pi:R\to A_{\mathcal U}.
$$
The set $K_{\mathcal U}$ is an ideal: intersections of zero sets
control sums, and multiplication by any tuple enlarges a zero set.
The quotient $A_{\mathcal U}$ is the ultraproduct, here defined by
this exact ideal. Also $K_{\mathcal U}\subseteq\mathfrak p$:
if $E=\{i:a_i=0\}\in\mathcal U$, then $e_Ea=0$ and
$e_E\notin\mathfrak p$, so primality gives $a\in\mathfrak p$.
Therefore $\mathfrak q=\mathfrak p/K_{\mathcal U}$ is a prime
of $A_{\mathcal U}$, and the map
$$
R_{\mathfrak p}\longrightarrow (A_{\mathcal U})_{\mathfrak q},
\qquad a/s\longmapsto\pi(a)/\pi(s)
$$
is defined, because $s\notin\mathfrak p$ if and only if
$\pi(s)\notin\mathfrak q$. It is surjective by lifting each
numerator and denominator through the quotient. It is injective:
if $\pi(a)/\pi(s)=0$, choose $t\notin\mathfrak p$ with
$ta\in K_{\mathcal U}$. For
$E=\{i:t_ia_i=0\}\in\mathcal U$, one has
$e_Eta=0$ and $e_Et\notin\mathfrak p$. This is the original
localization witness that $a/s=0$. The inverse map is consequently
$\pi(a)/\pi(s)\mapsto a/s$, independent of the chosen lifts by
the same argument applied to their difference. Both composites fix
the original fractions.

Conversely a prime of an ultraproduct contracts to a prime of
the product and recovers its ultrafilter: the class of $e_E$ is
$1$ when $E\in\mathcal U$ and $0$ when its complement is in
$\mathcal U$. Thus, as sets, all product primes are uniquely the
pairs of an ultrafilter and a prime of its ultraproduct, with the
displayed localization isomorphism for each pair.

The ultraproduct of normal rings is itself normal. If a displayed
$(\mathcal D_n)$ equation holds in $A_{\mathcal U}$, it holds
coordinatewise on a set $E\in\mathcal U$ by the definition of
the quotient ideal. For $i\in E$ choose $d_i$ with $a_i=b_id_i$
in $R_i$; set $d_i=0$ elsewhere. Their classes satisfy
$\pi(a)=\pi(b)\pi(d)$. The characterization proves normality
of $A_{\mathcal U}$, so the localization isomorphism also gives
the product result through its actual primes.

When every factor is a nonzero normal domain, these maps have
a stronger description. The ultraproduct is then a domain:
the zero set of a product is the union of the zero sets of its
two factors, so membership in an ultrafilter forces one of those
zero sets to belong to it. Its identity is nonzero because the
coordinate identities are nonzero. If $\pi(a)/\pi(b)$ is integral
in its fraction field with $\pi(b)\ne0$, then the sets
$$
T=\{i:b_i\ne0\},\qquad
E=\{i:a_i^n+\sum_{j=0}^{n-1}c_{j,i}a_i^jb_i^{n-j}=0\}
$$
belong to $\mathcal U$. On $E\cap T$, the original fraction
$a_i/b_i$ satisfies its monic equation over $R_i$ and belongs
to $R_i$. Choose that value for $d_i$ there and zero elsewhere.
Then $\pi(a)=\pi(b)\pi(d)$, proving integral closedness directly.
Moreover in this domain-factor case
$\ker(R\to R_{\mathfrak p})=K_{\mathcal U}$.
One inclusion was proved by the idempotent witnesses above.
For the other, if $ta=0$ with $t\notin\mathfrak p$, then
$t\notin K_{\mathcal U}$, so $\{i:t_i\ne0\}\in\mathcal U$;
on that set the domains force $a_i=0$, giving $a\in K_{\mathcal U}$.

Finally normality of the product implies normality of each factor.
For $\mathfrak r\in\operatorname{Spec}(R_i)$, take
$\mathfrak p=\operatorname{pr}_i^{-1}(\mathfrak r)$.
The localization map
$$
R_{\mathfrak p}\longrightarrow(R_i)_{\mathfrak r},
\qquad a/s\longmapsto a_i/s_i
$$
has inverse obtained by using the numerator tuple supported at $i$
and the denominator tuple equal to $s_i$ at $i$ and to $1$ elsewhere.
Multiplication by $e_{\{i\}}\notin\mathfrak p$ kills any
differences outside that coordinate, which verifies both composites
and independence of fraction representatives. Hence these local
rings are isomorphic, proving normality of the factor.
Zero-ring factors cause no exception: if $J=\{i:R_i\ne0\}$,
then $e_J=1_R$, so any ultrafilter obtained from a product prime
contains $J$. An ultrafilter concentrated on zero factors gives
the zero ultraproduct and contributes no primes. The empty product
is the zero ring and satisfies the same conclusion.
\end{proof}

\begin{lemma}
\label{editorial-source-lemma-characterize-reduced-ring-normal}
Let $R$ be a ring. Assume $R$ is reduced and has finitely many
minimal primes. Then the following are equivalent:
\begin{enumerate}
\item $R$ is a normal ring,
\item $R$ is integrally closed in its total ring of fractions, and
\item $R$ is a finite product of normal domains.
\end{enumerate}
\end{lemma}

\begin{proof}
The implications (1) $\Rightarrow$ (2) and
(3) $\Rightarrow$ (1) hold in general,
see Lemmas \ref{editorial-source-lemma-normal-ring-integrally-closed} and
\ref{editorial-source-lemma-finite-product-normal}.

\medskip\noindent
Let $\mathfrak p_1, \ldots, \mathfrak p_n$ be the minimal primes of $R$.
By Lemmas \ref{algebra-lemma-reduced-ring-sub-product-fields} and
\ref{algebra-lemma-total-ring-fractions-no-embedded-points} we have
$Q(R) = R_{\mathfrak p_1} \times \ldots \times R_{\mathfrak p_n}$, and
by Lemma \ref{algebra-lemma-minimal-prime-reduced-ring} each factor is a field.
Let $e_i = (0, \ldots, 0, 1, 0, \ldots, 0)$ be the $i$th idempotent
of $Q(R)$.

\medskip\noindent
Suppose $R$ is integrally closed in its original $Q(R)$.
Each $e_i$ satisfies $e_i^2-e_i=0$, hence lies in $R$.
The orthogonality identities and $\sum_i e_i=1$ give inverse maps
$$
R\longrightarrow\prod_{i=1}^n e_iR,\quad
r\longmapsto(e_ir)_i,
\qquad (r_i)_i\longmapsto\sum_{i=1}^n r_i.
$$
The factor $e_iR$ has identity $e_i$ and embeds in the original
field $R_{\mathfrak p_i}$ by the $i$th coordinate of $Q(R)$.
The kernel of $R\to e_iR$ is exactly $\mathfrak p_i$.
Indeed the kernel of $R\to R_{\mathfrak p_i}$ is
$\mathfrak p_i$: the localization is a field, whose maximal ideal
$\mathfrak p_iR_{\mathfrak p_i}$ is zero; conversely an equation
$ur=0$ with $u\notin\mathfrak p_i$ forces $r\in\mathfrak p_i$
by primality. Multiplication by $e_i$ keeps precisely this
coordinate, and injectivity of $R\to Q(R)$ identifies its
vanishing with $e_ir=0$ in $R$. Therefore
$R/\mathfrak p_i\cong e_iR$ by $r+\mathfrak p_i\mapsto e_ir$,
with inverse $e_ir\mapsto r+\mathfrak p_i$.
The maps
$$
\operatorname{Frac}(R/\mathfrak p_i)\longrightarrow R_{\mathfrak p_i},
\quad\frac{a+\mathfrak p_i}{b+\mathfrak p_i}
\longmapsto\frac{a/1}{b/1},
$$
and $a/s\mapsto(a+\mathfrak p_i)/(s+\mathfrak p_i)$ are inverse:
the denominators avoid $\mathfrak p_i$, and the original fraction
equality witnesses give equality in the domain $R/\mathfrak p_i$.
This proves the asserted field-of-fractions comparison.

Each subring $R/\mathfrak p_i\subset R_{\mathfrak p_i}$ is
integrally closed, as also follows from
Lemma \ref{editorial-source-lemma-finite-product-integral-closure}.
Directly, if $z_i$ in that field satisfies a monic equation of
degree $d\geq1$ over $e_iR$, form $z\in Q(R)$ with $i$th
coordinate $z_i$ and all others zero. Regard each coefficient
as the element of $R$ supported by $e_i$. The same monic
polynomial vanishes on $z$: at the other coordinates it is
evaluated at zero with all coefficients zero, including its
constant coefficient. Thus $z$ is integral over $R$ and belongs
to $R$ by hypothesis. Its $i$th coordinate lies in $e_iR$,
so $z_i$ belongs to the stated subring. This proves that each
$R/\mathfrak p_i$ is a normal domain and gives (3).
If $R=0$, there are no minimal primes and (3) is the empty
product, so the same conclusion holds directly.
\end{proof}

\begin{lemma}
\label{editorial-source-lemma-colimit-normal-ring}
Let $(R_i,\varphi_{ii'})$ be a system of rings indexed by a nonempty
directed preordered set $I$
(Categories, Definition \ref{categories-definition-directed-system}).
If each $R_i$ is normal, so is
$R = \colim_i R_i$.
\end{lemma}

\begin{proof}
Write $\lambda_i:R_i\to R$ for the original colimit maps.
If $R=0$, the conclusion is vacuous. Otherwise fix a prime
$\mathfrak p\subset R$ and set
$\mathfrak p_i=\lambda_i^{-1}(\mathfrak p)$.
These are proper primes, and
$\varphi_{ij}^{-1}(\mathfrak p_j)=\mathfrak p_i$ because
$\lambda_j\varphi_{ij}=\lambda_i$.
Thus the localized transition maps are defined by
$$
\psi_{ij}:(R_i)_{\mathfrak p_i}\longrightarrow(R_j)_{\mathfrak p_j},
\qquad a_i/s_i\longmapsto
\varphi_{ij}(a_i)/\varphi_{ij}(s_i).
$$
Indeed denominators outside $\mathfrak p_i$ stay outside
$\mathfrak p_j$, and an equality witness
$u_i(a_it_i-b_is_i)=0$ with $u_i\notin\mathfrak p_i$
maps to the corresponding equality witness in $R_j$.
Put $A_i=(R_i)_{\mathfrak p_i}$ and $L=\colim_i A_i$,
with maps $\rho_i:A_i\to L$.

The element rule for a directed colimit says that every element
has a representative at one stage, and equality of two representatives
holds at a common later stage. To recall why, construct the underlying
set from pairs $(i,a_i)$, identifying two pairs when their images
agree at some common stage. Directedness makes this relation
transitive; addition and multiplication are computed at a common
stage and are independent of that stage by the same rule.
This construction satisfies the universal property of the ring
colimit, since every compatible family of maps sends equivalent
pairs to the same value. In particular a zero equality may need
a later stage; no transition map is assumed injective.

We prove the original localization comparison with both maps:
$$
\alpha:L\longrightarrow R_{\mathfrak p},\qquad
\rho_i(a_i/s_i)\longmapsto
\lambda_i(a_i)/\lambda_i(s_i),
$$
and
$$
\beta:R_{\mathfrak p}\longrightarrow L,\qquad
a/s\longmapsto\rho_i(a_i/s_i),
$$
where numerator and denominator are represented at one common
stage $i$. Since $s\notin\mathfrak p$, every such representative
$s_i$ is outside $\mathfrak p_i$. For well-definedness of the
second map, an equality $a/s=b/t$ has a witness
$u\notin\mathfrak p$ with $u(at-bs)=0$ in $R$.
Represent all five elements at one common stage. The equality
rule gives a later stage $j$ with
$u_j(a_jt_j-b_js_j)=0$ in $R_j$.
All three denominators $u_j,s_j,t_j$ are outside
$\mathfrak p_j$, so $a_j/s_j=b_j/t_j$ in $A_j$.
The same argument handles any change of representative.
Addition and multiplication use the ordinary fraction formulas
at a common stage. The first map is defined by the compatible
localized maps to $R_{\mathfrak p}$, whose fraction relations
are respected in exactly the same way.
The composites fix every $a/s$ and every generator
$\rho_i(a_i/s_i)$, respectively. Hence
$$
R_{\mathfrak p}\cong\colim_i(R_i)_{\mathfrak p_i}.
$$

Each $A_i$ is a nonzero normal domain by the hypothesis.
The colimit $L$ is nonzero: otherwise $1=0$ would hold at
some later stage, contradicting this property of $A_i$.
It is a domain even though the transition maps may have kernels.
For $xy=0$ in $L$, represent $x,y$ at one stage and then
pass to a later stage where their product is zero.
The domain at that stage forces one representative to be zero,
so one of the original $x,y$ is zero in $L$.
Thus $A=R_{\mathfrak p}$ is a nonzero domain.

Retain $a,b\in A$, $b\ne0$, and an original monic equation
$$
z^n+\sum_{k=0}^{n-1}c_kz^k=0,
\qquad z=a/b,\quad c_k\in A,\quad n\geq1.
$$
Multiplication by the full denominator power $b^n$ gives
$a^n+\sum_{k=0}^{n-1}c_ka^kb^{n-k}=0$ in $A$.
Choose a common stage $i$ representing each of the finitely
many localized elements as
$$
a=\frac{\lambda_i(a_i)}{\lambda_i(u_i)},\qquad
b=\frac{\lambda_i(b_i)}{\lambda_i(v_i)},\qquad
c_k=\frac{\lambda_i(c_{k,i})}{\lambda_i(w_{k,i})},
$$
with $u_i,v_i,w_{k,i}\notin\mathfrak p_i$.
Keep the products
$$
W_i=\prod_{k=0}^{n-1}w_{k,i},\qquad
W_{k,i}=\prod_{\ell\ne k}w_{\ell,i},\qquad
D_i=u_i^nv_i^nW_i
$$
and the full numerator
$$
H_i=a_i^nv_i^nW_i+
\sum_{k=0}^{n-1}c_{k,i}a_i^kb_i^{n-k}
 u_i^{n-k}v_i^kW_{k,i}.
$$
In $A_i$, without omitting any denominator factor, one has
$$
\frac{H_i}{D_i}
=\left(\frac{a_i}{u_i}\right)^n+
\sum_{k=0}^{n-1}\frac{c_{k,i}}{w_{k,i}}
 \left(\frac{a_i}{u_i}\right)^k
 \left(\frac{b_i}{v_i}\right)^{n-k}.
$$
The image of this fraction in $A$ is zero. Therefore some
$t\in R\setminus\mathfrak p$ satisfies
$t\lambda_i(H_i)=0$ in the original $R$.
Represent $t$ at a stage above $i$, then apply the equality rule
again to obtain a stage $q$ where $t_qH_q=0$ in $R_q$.
Here every entry of $H_q$ is the image of the corresponding
original entry of $H_i$ under the indicated transition map:
$$
H_q=a_q^nv_q^n\prod_{\ell=0}^{n-1}w_{\ell,q}
+\sum_{k=0}^{n-1}c_{k,q}a_q^kb_q^{n-k}u_q^{n-k}v_q^k
 \prod_{\ell\ne k}w_{\ell,q}.
$$
The witness $t_q$ and all the denominator factors are outside
$\mathfrak p_q$. Thus $t_q$ and $D_q$ are units in $A_q$,
and the complete displayed equation for $H_q/D_q$ is zero there.
The element $b_q/v_q$ is nonzero in $A_q$, since its image is
the original nonzero $b$ in $A$. Consequently
$$
z_q=\frac{a_q/u_q}{b_q/v_q}\in\operatorname{Frac}(A_q)
$$
is defined, and division by $(b_q/v_q)^n$ gives its monic equation
$$
z_q^n+\sum_{k=0}^{n-1}(c_{k,q}/w_{k,q})z_q^k=0.
$$
Normality gives $z_q=h_q\in A_q$. Transport the equality
$a_q/u_q=(b_q/v_q)h_q$ in $A_q$ through its actual map to $A$.
It becomes $a=bh$ for the image $h\in A$ of $h_q$.
Since $A$ is a domain and $b\ne0$, the original $z=a/b$ equals
$h$ and belongs to $A$. This proves integral closedness of $A$
and hence normality of $R$.
No map from $\operatorname{Frac}(A_q)$ to
$\operatorname{Frac}(A)$ has been assumed: it need not exist
when the stage map has a kernel. The equality transported is
the displayed equality inside $A_q$ itself.

The divisibility characterization in the arbitrary-product proof
gives a second direct verification. Represent the finitely many
elements of a $(\mathcal D_n)$ equation in $R$ at one stage,
then pass to a later stage where that whole equation holds.
The normal ring at that stage supplies $a_i=b_id_i$;
its image gives the required $a=bd$ in $R$. Thus all
$(\mathcal D_n)$ hold, in agreement with the localization proof.
Neither proof requires antisymmetry of the directed preorder
or injectivity of the transition maps.
\end{proof}










\subsection{Editorial treatment of Algebra lines 8347--8357}

\hypertarget{algebra-context-090}{}

\noindent\href{https://github.com/stacks/stacks-project/blob/a04446e57ec1fbc252a871afcec7752fb2807b14/algebra.tex\#L8347}{Original source, lines 8347--8357}. Complete original and reviewed TeX passages accompany this note. Every intervention is separately recorded; the following treatment is editorial.


\begin{definition}
\label{editorial-source-definition-integral-over-ideal}
Let $\varphi : R \to S$ be a ring map.
Let $I \subset R$ be an ideal.
We say an element $g \in S$ is
{\it integral over $I$} if
there exists a monic
polynomial $P = x^d + \sum_{j=0}^{d-1} a_j x^j$, with $d\geq1$,
with coefficients $a_j \in I^{d-j}$ such
that $P^\varphi(g) = 0$ in $S$.
\end{definition}


\subsection{Editorial treatment of Algebra lines 8365--8396}

\hypertarget{algebra-context-091}{}

\noindent\href{https://github.com/stacks/stacks-project/blob/a04446e57ec1fbc252a871afcec7752fb2807b14/algebra.tex\#L8365}{Original source, lines 8365--8396}. Complete original and reviewed TeX passages accompany this note. Every intervention is separately recorded; the following treatment is editorial.


\begin{lemma}
\label{editorial-source-lemma-characterize-integral-ideal}
Let $\varphi : R \to S$ be a ring map.
Let $I \subset R$ be an ideal.
Let $A = \bigoplus_{h\geq0} I^ht^h \subset R[t]$ be the
subring of the polynomial ring
generated by $R \oplus It \subset R[t]$.
An element $s \in S$ is integral over $I$ if
and only if the element $st \in S[t]$
is integral over $A$, with the coefficient map induced by $\varphi$.
More generally, for every integer $k\geq1$, the element $s$ is
integral over $I^k$ if and only if $st^k$ is integral over $A$.
\end{lemma}

\begin{proof}
Let $\widetilde\varphi:R[t]\to S[t]$ send
$\sum_h r_ht^h$ to $\sum_h\varphi(r_h)t^h$, and use its
restriction to $A$. Its additivity and multiplicativity follow
coefficient by coefficient from those of $\varphi$. The displayed
description of $A$ is the subring generated by $R\oplus It$:
products of $h$ elements of $It$ generate $I^ht^h$ additively,
and products of its homogeneous summands have the required powers
$I^hI^\ell=I^{h+\ell}$. Here $I^0=R$.

We prove the assertion for every $k\geq1$, retaining $k=1$ as
the first assertion. Suppose $st^k$ is integral over $A$, and
keep an original monic equation
$$
P=x^d+\sum_{j=0}^{d-1}a_jx^j,\qquad
a_j=\sum_{h\geq0}a_{jh}t^h,\quad a_{jh}\in I^h,
\qquad P^{\widetilde\varphi}(st^k)=0.
$$
Each coefficient sum has finite support. The full equation is
$$
s^dt^{kd}+\sum_{j=0}^{d-1}\sum_{h\geq0}
  \varphi(a_{jh})s^jt^{h+kj}=0.
$$
For every $\ell\geq0$ its coefficient of $t^\ell$ is
$$
\delta_{\ell,kd}s^d+
\sum_{\substack{0\leq j<d,\ h\geq0\\h+kj=\ell}}
  \varphi(a_{jh})s^j=0,
$$
where $\delta_{\ell,kd}$ equals $1$ for $\ell=kd$ and $0$
otherwise. Thus none of the other coefficient equations is discarded.
Taking $\ell=kd$ gives
$$
s^d+\sum_{j=0}^{d-1}\varphi(a_{j,k(d-j)})s^j=0,
\qquad a_{j,k(d-j)}\in I^{k(d-j)}=(I^k)^{d-j}.
$$
This proves integrality over $I^k$. In the original $k=1$
notation, $a_j''=a_{j,d-j}$ and $a_j'=a_j''t^{d-j}$,
so the degree-$d$ equation is exactly
$(st)^d+\sum_{j=0}^{d-1}\varphi(a_j'')t^{d-j}(st)^j=0$.

Conversely retain a monic $P=x^d+\sum_{j=0}^{d-1}a_jx^j$
with $a_j\in I^{k(d-j)}$ and $P^\varphi(s)=0$.
Every coefficient of
$$
Q=x^d+\sum_{j=0}^{d-1}(a_jt^{k(d-j)})x^j
$$
belongs to $A$, and the exact equation
$$
Q^{\widetilde\varphi}(st^k)
=t^{kd}\left(s^d+\sum_{j=0}^{d-1}\varphi(a_j)s^j\right)=0
$$
proves the converse, without cancellation of $t$ or any
injectivity assumption on $\varphi$.
\end{proof}


\subsection{Editorial treatment of Algebra lines 8398--8419}

\hypertarget{algebra-context-092}{}

\noindent\href{https://github.com/stacks/stacks-project/blob/a04446e57ec1fbc252a871afcec7752fb2807b14/algebra.tex\#L8398}{Original source, lines 8398--8419}. Complete original and reviewed TeX passages accompany this note. Every intervention is separately recorded; the following treatment is editorial.


\begin{lemma}
\label{editorial-source-lemma-integral-over-ideal-is-submodule}
Let $\varphi : R \to S$ be a ring map.
Let $I \subset R$ be an ideal.
The set of elements of $S$ which are integral
over $I$ forms an $R$-submodule of $S$.
Furthermore, if $s \in S$ is integral over
$R$, and $s'$ is integral over $I$, then
$ss'$ is integral over $I$.
\end{lemma}

\begin{proof}
Use $A$ and its map to $S[t]$ from
Lemma \ref{editorial-source-lemma-characterize-integral-ideal}. The elements of
$S[t]$ integral over $A$ form a subring by
Lemma \ref{editorial-source-lemma-integral-closure-is-ring}. If $u,v\in S$
are integral over $I$, the elements $ut$ and $vt$ belong to
that subring, so do $(u+v)t=ut+vt$ and $(-u)t=-ut$.
The same characterization gives closure under addition and negatives;
$0$ is integral over $I$ by the monic equation $x=0$.

If $s\in S$ is integral over $R$, its original monic equation
over $R$ is also a monic equation for the constant polynomial
$s\in S[t]$ over $A$, since the degree-zero part of $A$ is
$R$. For $s'$ integral over $I$, the element $s't$ is integral
over $A$. Their product $s(s't)=(ss')t$ is therefore integral
over $A$, and the characterization gives $ss'$ integral over
$I$, which proves the second assertion. Finally every
$\varphi(r)$ for $r\in R$ is integral over $R$ by the equation
$x-r$, so this product assertion supplies multiplication by
every original $R$-scalar and proves the submodule assertion.
\end{proof}


\subsection{Editorial treatment of Algebra lines 8421--8430}

\hypertarget{algebra-context-093}{}

\noindent\href{https://github.com/stacks/stacks-project/blob/a04446e57ec1fbc252a871afcec7752fb2807b14/algebra.tex\#L8421}{Original source, lines 8421--8430}. Complete original and reviewed TeX passages accompany this note. Every intervention is separately recorded; the following treatment is editorial.


\begin{lemma}
\label{editorial-source-lemma-integral-integral-over-ideal}
Suppose $\varphi : R \to S$ is integral.
Suppose $I \subset R$ is an ideal.
Then every element of $IS$ is integral over $I$.
\end{lemma}

\begin{proof}
An element of $IS$ is a finite sum
$\sum_{i=1}^r\varphi(a_i)s_i$ with $a_i\in I$ and $s_i\in S$.
Each $\varphi(a_i)$ is integral over $I$ by the degree-one
equation $x-a_i$. Every $s_i$ is integral over $R$ by the
hypothesis on $\varphi$. The product assertion and then addition
in Lemma \ref{editorial-source-lemma-integral-over-ideal-is-submodule} give
integrality of the entire sum over $I$. The empty sum is $0$,
already covered by its degree-one equation.
\end{proof}

\begin{lemma}
\label{editorial-source-lemma-normal-principal-ideals-integrally-closed}
A ring $R$ is normal if and only if every principal ideal
$bR$, for $b\in R$, is integrally closed in $R$ in the sense
of Definition \ref{editorial-source-definition-integral-over-ideal}.
This includes $b=0$ and zero divisors $b$.
\end{lemma}

\begin{proof}
For every integer $e\geq1$, the equality $(bR)^e=b^eR$
follows in both directions by multiplying generators and by
writing $b^er$ as a product of $e$ elements of $bR$.
Consequently an element $a\in R$ is integral over $bR$ exactly
when, for some $n\geq1$ and $c_0,\ldots,c_{n-1}\in R$,
$$
a^n+\sum_{j=0}^{n-1}c_ja^jb^{n-j}=0.
$$
Indeed each original coefficient of a defining equation lies in
$(bR)^{n-j}=b^{n-j}R$ and can be written $c_jb^{n-j}$;
conversely these displayed coefficients have the required ideal
powers. This uses no cancellation or uniqueness of $c_j$.
Also every $a\in bR$ is integral over $bR$ by $x-a$.
Thus every principal ideal is integrally closed exactly when
every displayed equation implies $a=bd$ for some $d\in R$.
This is precisely the family of conditions $(\mathcal D_n)$,
for all $n\geq1$, proved equivalent to normality in the proof of
Lemma \ref{editorial-source-lemma-finite-product-normal}. Both implications of
that characterization hold for arbitrary rings, including the
zero ring and the cases where $b$ is a zero divisor; applying
them proves the assertion.
\end{proof}


\subsection{Editorial treatment of Algebra lines 8432--8496}

\hypertarget{algebra-context-094}{}

\noindent\href{https://github.com/stacks/stacks-project/blob/a04446e57ec1fbc252a871afcec7752fb2807b14/algebra.tex\#L8432}{Original source, lines 8432--8496}. Complete original and reviewed TeX passages accompany this note. Every intervention is separately recorded; the following treatment is editorial.


\begin{lemma}
\label{editorial-source-lemma-polynomials-divide}
Let $K$ be a field. Let $n, m \in \mathbf{N}$ and
$a_0, \ldots, a_{n - 1}, b_0, \ldots, b_{m - 1} \in K$.
If the polynomial $x^n + a_{n - 1}x^{n - 1} + \ldots + a_0$
divides the polynomial $x^m + b_{m - 1} x^{m - 1} + \ldots + b_0$
in $K[x]$ then
\begin{enumerate}
\item $a_0, \ldots, a_{n - 1}$ are integral over any subring
$R_0$ of $K$ containing the elements $b_0, \ldots, b_{m - 1}$, and
\item each $a_i$ lies in $\sqrt{(b_0, \ldots, b_{m-1})R}$
for any subring $R \subset K$ containing the elements
$a_0, \ldots, a_{n - 1}, b_0, \ldots, b_{m - 1}$.
\end{enumerate}
\end{lemma}

\begin{proof}
Write
$$
Q=x^n+\sum_{j=0}^{n-1}a_jx^j,\qquad
P=x^m+\sum_{i=0}^{m-1}b_ix^i.
$$
Since both are monic and $Q$ divides $P$ over the field $K$,
we have $m\geq n$. Choose a splitting field $L/K$ and retain
the roots with their multiplicities:
$$
P=\prod_{i=1}^m(x-\beta_i),\qquad \beta_i\in L.
$$
See Fields, Section \ref{fields-section-splitting-fieds}.
Unique factorization in $L[x]$ gives an indexed subset
$J\subseteq\{1,\ldots,m\}$ of cardinality $n$ with
$Q=\prod_{i\in J}(x-\beta_i)$. Equal roots remain distinct
indices in this expression. Every $\beta_i$ is integral over
$R_0$ by the original $P$, and every original coefficient is
$$
a_j=(-1)^{n-j}
\sum_{\substack{E\subseteq J\\ |E|=n-j}}
\prod_{i\in E}\beta_i\qquad(0\leq j<n).
$$
The subring property in Lemma \ref{editorial-source-lemma-integral-closure-is-ring}
proves (1), retaining every sign and repeated-root contribution.

For (2), first consider $m=n$. The monic quotient has degree
zero and is $1$, so $a_i=b_i$ for every $i$, which gives the
assertion directly. Now let $e=m-n>0$ and write the original
monic quotient as
$$
C=x^e+\sum_{k=0}^{e-1}c_kx^k,\qquad P=QC.
$$
Put $a_n=c_e=1$. Comparing the coefficient of $x^{n+k}$
for $k=e-1,e-2,\ldots,0$ gives the exact recurrence
$$
c_k=b_{n+k}-
\sum_{i=\max(0,n+k-e)}^{n-1}a_i c_{n+k-i}.
$$
Every index $n+k-i$ in this sum is between $k+1$ and $e$.
Thus descending induction shows $c_k\in R$ for all $k$,
without negative coefficient indices or assumptions on $n>1$.
Set $H=\sqrt{(b_0,\ldots,b_{m-1})R}$ and
$\overline R=R/H$. This quotient is reduced. If it is the
zero ring, the required membership in $H=R$ already holds.
Otherwise the complete coefficient image of $P=QC$ is
$$
x^m+\sum_{i=0}^{m-1}\overline b_i x^i=x^m
=\left(x^n+\sum_{j=0}^{n-1}\overline a_jx^j\right)
 \left(x^e+\sum_{k=0}^{e-1}\overline c_kx^k\right).
$$
At a minimal prime $\mathfrak p$ of $\overline R$, the ring
$\overline R_{\mathfrak p}$ is a field by
Lemma \ref{algebra-lemma-minimal-prime-reduced-ring}. In its polynomial
ring unique factorization makes the first monic factor
$f=x^n+\sum_{j=0}^{n-1}\overline a_jx^j$ equal to $x^n$:
every irreducible factor divides $x$, its degree is $n$, and
its leading coefficient fixes the unit factor to be $1$.
For each $j$ choose $s_j\notin\mathfrak p$ with
$s_j\overline a_j=0$, using the equality in this localization,
and take $s=\prod_{j=0}^{n-1}s_j$.
Then there exists an $s\in\overline R$, $s\notin\mathfrak p$,
such that $s(f-x^n)=0$, with this explicit common choice.
Primality gives $\overline a_j\in\mathfrak p$ for every $j$.
As this holds at every minimal prime and their intersection
is zero in the reduced ring, every $\overline a_j$ is zero
(Lemma \ref{algebra-lemma-reduced-ring-sub-product-fields}).
Hence every $a_j$ belongs to $H$, proving (2).
\end{proof}

\begin{lemma}
\label{editorial-source-lemma-weighted-monic-factor-coefficients}
Let $R\subseteq K$ be a subring of a field and let $I\subseteq R$
be an ideal. Suppose
$$
P=x^m+\sum_{i=0}^{m-1}b_ix^i,\qquad b_i\in I^{m-i},
\qquad Q=x^n+\sum_{j=0}^{n-1}a_jx^j\in K[x]
$$
are monic, $1\leq n\leq m$, and $Q$ divides $P$.
For each $0\leq j<n$, the coefficient $a_j$ is integral over
$I^{n-j}$. More precisely, with $N=\binom{m}{n}$, it satisfies
an equation
$$
a_j^N+\sum_{\ell=0}^{N-1}C_{j,\ell}a_j^\ell=0,
\qquad C_{j,\ell}\in I^{(n-j)(N-\ell)}.
$$
\end{lemma}

\begin{proof}
In a splitting field $L/K$ retain the original indexed roots,
including multiplicities, so
$P=\prod_{r=1}^m(x-\beta_r)$. As in the proof of
Lemma \ref{editorial-source-lemma-polynomials-divide}, there is an indexed subset
$J\subseteq\{1,\ldots,m\}$ of cardinality $n$ with
$Q=\prod_{r\in J}(x-\beta_r)$. Fix $j$ and put $w=n-j$.
Let $\mathcal T$ be the set of all $n$-element subsets of
$\{1,\ldots,m\}$, so $|\mathcal T|=N$. For each $T\in\mathcal T$
define the full coefficient
$$
r_{j,T}=(-1)^w
\sum_{\substack{U\subseteq T\\|U|=w}}\prod_{r\in U}\beta_r.
$$
In particular $a_j=r_{j,J}$. Form
$$
H_j(Y)=\prod_{T\in\mathcal T}(Y-r_{j,T})
=Y^N+\sum_{\ell=0}^{N-1}C_{j,\ell}Y^\ell,
$$
where every coefficient, sign and repeated contribution is
$$
C_{j,\ell}=(-1)^{N-\ell}
\sum_{\substack{\mathcal A\subseteq\mathcal T\\|\mathcal A|=N-\ell}}
\prod_{T\in\mathcal A}r_{j,T}.
$$
Then $H_j(a_j)=0$.

Replace the $\beta_r$ temporarily by independent variables
$Z_r$. The same expression for $C_{j,\ell}$ is an integer
polynomial, symmetric in the $m$ variables and homogeneous of
total degree $D=w(N-\ell)$. It is an integer polynomial in
the elementary symmetric polynomials $e_1,\ldots,e_m$, each
monomial having weighted degree $D$ when $e_s$ has weight $s$.
Here is a proof retaining the weights. In lexicographic order
$Z_1>\cdots>Z_m$, the leading exponent vector
$(\lambda_1,\ldots,\lambda_m)$ of a nonzero symmetric
homogeneous polynomial satisfies
$\lambda_1\geq\cdots\geq\lambda_m$: an inversion would give
a larger monomial after interchanging the corresponding
variables, contradicting maximality. With $\lambda_{m+1}=0$,
the polynomial
$$
\prod_{s=1}^m e_s^{\lambda_s-\lambda_{s+1}}
$$
has that same leading monomial with coefficient $1$ and
weighted degree
$\sum_{s=1}^m s(\lambda_s-\lambda_{s+1})
=\sum_{s=1}^m\lambda_s=D$.
Subtract the integer multiple matching the leading coefficient.
The result remains symmetric and homogeneous of degree $D$,
with a smaller leading monomial. Only finitely many monomials
have total degree $D$, so this procedure terminates. It proves
an expression with integers $\kappa_{j,\ell,\nu}$:
$$
C_{j,\ell}
=\sum_{\substack{\nu\in\mathbf Z_{\geq0}^m\\
                  \sum_{s=1}^m s\nu_s=D}}
\kappa_{j,\ell,\nu}
\prod_{s=1}^m e_s(\beta_1,\ldots,\beta_m)^{\nu_s}.
$$
The original coefficients of $P$ give
$e_s(\beta_1,\ldots,\beta_m)=(-1)^s b_{m-s}$, so the
complete substituted expression is
$$
C_{j,\ell}
=\sum_{\substack{\nu\in\mathbf Z_{\geq0}^m\\
                  \sum_{s=1}^m s\nu_s=D}}
\kappa_{j,\ell,\nu}
\prod_{s=1}^m\bigl((-1)^s b_{m-s}\bigr)^{\nu_s}.
$$
Each product belongs to
$\prod_s(I^s)^{\nu_s}=I^D=I^{w(N-\ell)}$.
This proves the required equation over $R$. All identities
were proved over the integers before evaluation, so no
separability or characteristic restriction is present.

The exact Rees comparison gives a second proof of the
integrality conclusion. Put
$A=\bigoplus_{h\geq0}I^ht^h\subseteq R[t]$ and map it to
$L[t]$ by the given inclusion. Every $t\beta_r$ satisfies
$$
(t\beta_r)^m+
\sum_{i=0}^{m-1}(b_it^{m-i})(t\beta_r)^i=t^mP(\beta_r)=0.
$$
Its coefficients belong to $A$, so all these elements are
integral over $A$. The original signed identity
$$
a_jt^{n-j}=(-1)^{n-j}
\sum_{\substack{U\subseteq J\\|U|=n-j}}
\prod_{r\in U}(t\beta_r)
$$
and Lemma \ref{editorial-source-lemma-integral-closure-is-ring} show that this
element is integral over $A$. The weight-$n-j$ assertion of
Lemma \ref{editorial-source-lemma-characterize-integral-ideal} gives integrality
of the original $a_j$ over $I^{n-j}$.
\end{proof}

\begin{lemma}
\label{editorial-source-lemma-integral-ideal-minimal-polynomial-criterion}
Let $R$ be a domain, $K=\operatorname{Frac}(R)$, $I\subseteq R$
an ideal, and $g$ an algebraic element in a field extension
$L/K$. Write its monic minimal polynomial as
$Q=x^n+\sum_{j=0}^{n-1}a_jx^j\in K[x]$.
Then $g$ is integral over $I$ if and only if every $a_j$ is
integral over $I^{n-j}$. If $R$ is normal, this is equivalent
to $a_j\in\overline{I^{n-j}}\subseteq R$ for every $j$,
where the bar denotes integral closure of the ideal inside $R$.
\end{lemma}

\begin{proof}
If $g$ is integral over $I$, retain a witnessing equation
$P(g)=0$ with $P=x^m+\sum_{i=0}^{m-1}b_ix^i$ and
$b_i\in I^{m-i}$. Division by $Q$ in $K[x]$ gives a
remainder of degree less than $n$ that vanishes at $g$, hence
is zero by minimality. Thus $Q\mid P$, and
Lemma \ref{editorial-source-lemma-weighted-monic-factor-coefficients} gives
the asserted coefficient integrality.

Conversely suppose every coefficient has that integrality.
Use $A=\bigoplus_{h\geq0}I^ht^h\subseteq L[t]$ with the
original coefficient inclusion. By the weighted Rees
characterization every $h_j=a_jt^{n-j}$ is integral over $A$.
The algebra $B=A[h_0,\ldots,h_{n-1}]$ is finite over $A$:
if the chosen monic equation for $h_j$ has degree $d_j$,
the products $\prod_jh_j^{e_j}$ with $0\leq e_j<d_j$
generate it as an $A$-module, using those original monic
equations to reduce each higher power. The element $gt$
satisfies the full monic equation
$$
(gt)^n+\sum_{j=0}^{n-1}h_j(gt)^j=t^nQ(g)=0.
$$
Thus $B[gt]$ is generated over $A$ by all
$\bigl(\prod_jh_j^{e_j}\bigr)(gt)^r$ with $0\leq r<n$.
To see directly that this gives integrality over $A$, list
these generators as $v_1,\ldots,v_M$, including $1$.
Multiplication by $gt$ has relations
$(gt)v_i=\sum_k d_{ik}v_k$ with $d_{ik}\in A$.
Multiplying the vector equation by the adjugate of
$XI-(d_{ik})$ at $X=gt$ shows that the monic polynomial
$\det(XI-(d_{ik}))$ annihilates every $v_i$, including $1$.
It is therefore an integral equation for $gt$ over $A$.
The weight-one Rees characterization makes $g$ integral over
$I$, proving the converse. This argument does not require
the displayed module generators to be linearly independent.

If $R$ is normal, each integral equation for $a_j$ over an
ideal is in particular an integral equation over $R$.
Since $a_j\in K$, integral closedness of the normal domain
gives $a_j\in R$. The definition of integral closure of
$I^{n-j}$ inside $R$ now gives exactly the stated equivalence.
\end{proof}

\begin{example}
\label{editorial-source-example-weighted-coefficients-need-integral-closure}
Normality does not place these coefficients in the ideal
powers themselves. Let $k$ be any field,
$R=k[u,v]$, $K=k(u,v)$ and $I=(u^2,v^2)$.
The polynomial ring $R$ is a normal domain by
Lemma \ref{editorial-source-lemma-polynomial-domain-normal}, applied twice to $k$.
Choose a root $g$ of $x^2-u^3v$ in a field extension of $K$.
This quadratic is irreducible: in $k(u)[v]$ the exponent of
the factor $v$ defines an integer order on nonzero rational
functions by numerator order minus denominator order.
Multiplication adds this order, as follows by removing the
largest powers of $v$ and observing that the remaining
polynomials have nonzero values at $v=0$. Thus every nonzero
square in $K$ has even order, whereas $u^3v$ has order $1$.
It cannot be a square, so the quadratic has no root in $K$
and is irreducible, including in characteristic two.
Consequently $Q=x^2-u^3v$ is the minimal polynomial of $g$,
and $S=R[g]$ in the chosen extension field is a domain.
Yet the original equation
$$
g^4-u^6v^2=0,\qquad u^6v^2=(u^2)^3v^2\in I^4
$$
makes $g$ integral over $I$; its other three coefficients
are zero and belong to the required powers. The constant
coefficient $a_0=-u^3v$ of $Q$ is not in
$I^2=(u^4,u^2v^2,v^4)$: the monomial $u^3v$ is divisible
by none of these generators, and every polynomial multiple
of a generator is a sum of divisible monomials. Its distinct
monomial coefficient therefore cannot be obtained in their
sum. However
$$
a_0^2-u^6v^2=0,\qquad -u^6v^2\in I^4=(I^2)^2,
$$
so $a_0\in\overline{I^2}\setminus I^2$.
Even weight one need not give actual membership: $h=uv\in R$
satisfies $h^2-u^2v^2=0$ with $u^2v^2\in I^2$, but the
constant coefficient $-uv$ of its minimal polynomial $x-uv$
is outside $I$ by the same monomial argument.
\end{example}

\begin{lemma}
\label{editorial-source-lemma-integrally-closed-powers-monic-factors}
For a normal domain $R$, $K=\operatorname{Frac}(R)$ and an
ideal $I\subseteq R$, the following are equivalent:
\begin{enumerate}
\item Every positive power $I^w$ is integrally closed in $R$.
\item For every monic $P=x^m+\sum_{i=0}^{m-1}b_ix^i$ with
$b_i\in I^{m-i}$, every monic divisor
$Q=x^n+\sum_{j=0}^{n-1}a_jx^j\in K[x]$ of positive degree
has $a_j\in I^{n-j}$ for every $j$.
\item Every element integral over $I$ in a field extension of
$K$ has minimal-polynomial coefficients $a_j\in I^{n-j}$,
where $n$ is its minimal-polynomial degree.
\end{enumerate}
\end{lemma}

\begin{proof}
The weighted coefficient lemma gives (1)$\Rightarrow$(2):
its equations first give $a_j\in R$ by normality and then
$a_j\in I^{n-j}$ by integral closedness of that power.
Minimal-polynomial division as above gives (2)$\Rightarrow$(3).
For (3)$\Rightarrow$(2), factor the given $Q$ in $K[x]$
into monic irreducible factors, repeated with their original
multiplicities. A root of each factor in a splitting field
is a root of $P$, hence is integral over $I$. Condition (3)
gives the required powers for each irreducible factor.
Write the factors as
$Q_r=\sum_{j=0}^{d_r}c_{r,j}x^j$ for $1\leq r\leq q$,
with $c_{r,d_r}=1\in I^0$ and
$c_{r,j}\in I^{d_r-j}$ for $j<d_r$.
The full coefficient of $x^j$ in their product is
$$
\sum_{\substack{0\leq j_r\leq d_r\ (1\leq r\leq q)\\
                 j_1+\cdots+j_q=j}}
\prod_{r=1}^q c_{r,j_r}
\ \in I^{\sum_{r=1}^qd_r-j}=I^{n-j}.
$$
This retains the multiplicities and proves (2).

Finally assume (2), fix $w\geq1$, and take
$c\in\overline{I^w}\subseteq R$. Keep its defining equation
$$
c^d+\sum_{\ell=0}^{d-1}r_\ell c^\ell=0,
\qquad r_\ell\in I^{w(d-\ell)}.
$$
Set $F(Y)=Y^d+\sum_{\ell=0}^{d-1}r_\ell Y^\ell$.
The complete difference identity is
$$
F(Y)-F(c)=(Y-c)\left(
\sum_{k=0}^{d-1}Y^{d-1-k}c^k+
\sum_{\ell=1}^{d-1}r_\ell
   \sum_{k=0}^{\ell-1}Y^{\ell-1-k}c^k\right).
$$
Since $F(c)=0$, substituting $Y=x^w$ proves that $x^w-c$
divides
$P=x^{wd}+\sum_{\ell=0}^{d-1}r_\ell x^{w\ell}$.
The coefficient at $x^{w\ell}$ belongs to $I^{wd-w\ell}$;
every omitted coefficient is zero and is retained in its
required power. Condition (2) applied to this divisor gives
$-c\in I^w$, hence $c\in I^w$. The reverse inclusion
$I^w\subseteq\overline{I^w}$ follows from the degree-one
equation for each of its elements. This proves (1).
\end{proof}

\begin{remark}
\label{editorial-source-remark-weighted-coefficient-consequences}
Under the normal-domain hypotheses, the original coefficient
$a_j\in\overline{I^{n-j}}$ lies in $\sqrt I$: reduce its
defining equation modulo $\sqrt I$, where all nonleading
coefficients vanish; its positive leading power is zero, so
it belongs to that radical. If $I$ is integrally closed,
the same defining coefficients belong to $I^{N-\ell}$
because $(n-j)(N-\ell)\geq N-\ell$, and therefore $a_j$
is integral over $I$ and belongs to $I$. If $I$ is radical,
the radical conclusion already gives membership in $I$.
These claims do not assert membership in $I^{n-j}$ without
the corresponding integral-closedness condition.

The exponent $n-j$ cannot be increased in general.
For $R=k[u]$, $I=(u)$ and any $w\geq1$, take
$P=Q=x^w-u^w$. Its constant coefficient $a_0=-u^w$
belongs to $I^w$. If it were integral over $I^{w+1}$,
an equation of degree $d\geq1$ would have leading term
$(-u^w)^d$ of $u$-order $wd$, whereas each other term
$c_\ell(-u^w)^\ell$ would have order at least
$(w+1)(d-\ell)+w\ell=wd+d-\ell>wd$.
The coefficient of $u^{wd}$ in their sum would remain
$(-1)^d\ne0$, a contradiction in every characteristic.
\end{remark}


\subsection{Editorial treatment of Algebra lines 8498--8515}

\hypertarget{algebra-context-095}{}

\noindent\href{https://github.com/stacks/stacks-project/blob/a04446e57ec1fbc252a871afcec7752fb2807b14/algebra.tex\#L8498}{Original source, lines 8498--8515}. Complete original and reviewed TeX passages accompany this note. Every intervention is separately recorded; the following treatment is editorial.


\begin{lemma}
\label{editorial-source-lemma-minimal-polynomial-normal-domain}
Let $R \subset S$ be an inclusion of domains.
Assume $R$ is normal. Let $g \in S$ be integral
over $R$. Then the minimal polynomial of $g$
has coefficients in $R$.
\end{lemma}

\begin{proof}
The original inclusion $R\subset S$ of domains gives an
injective map $K=\operatorname{Frac}(R)\to\operatorname{Frac}(S)$,
sending $a/b$ to the same fraction; denominators stay nonzero,
and equality of fractions is preserved in both directions by
cross multiplication. Regard $g$ in this fixed extension field.
Choose its original monic equation
$P=x^m+\sum_{i=0}^{m-1}b_ix^i\in R[x]$, with $P(g)=0$,
and its monic minimal polynomial
$Q=x^n+\sum_{j=0}^{n-1}a_jx^j\in K[x]$.
Division in $K[x]$ gives $P=QC+D$ with $\deg D<n$ or $D=0$.
Evaluation at $g$ gives $D(g)=0$, so minimality forces $D=0$.
Part (1) of Lemma \ref{editorial-source-lemma-polynomials-divide} makes every
$a_j$ integral over $R$. Each already belongs to the original
$K$, so normality gives $a_j\in R$.
\end{proof}


\subsection{Editorial treatment of Algebra lines 8517--8572}

\hypertarget{algebra-context-096}{}

\noindent\href{https://github.com/stacks/stacks-project/blob/a04446e57ec1fbc252a871afcec7752fb2807b14/algebra.tex\#L8517}{Original source, lines 8517--8572}. Complete original and reviewed TeX passages accompany this note. Every intervention is separately recorded; the following treatment is editorial.


\begin{proposition}
\label{editorial-source-proposition-going-down-normal-integral}
Let $R \subset S$ be an inclusion of domains.
Assume $R$ is normal and $S$ integral over $R$.
Let $\mathfrak p \subset \mathfrak p' \subset R$
be primes. Let $\mathfrak q'$ be a prime of $S$
with $\mathfrak p' = R \cap \mathfrak q'$.
Then there exists a prime $\mathfrak q$
with $\mathfrak q \subset \mathfrak q'$
such that $\mathfrak p = R \cap \mathfrak q$. In other words:
the going down property holds for $R \to S$, see
Definition \ref{algebra-definition-going-up-down}.
\end{proposition}

\begin{proof}
Keep the stated primes and the original inclusion, and put
$B=S_{\mathfrak q'}$. It suffices to prove
$\mathfrak pB\cap R=\mathfrak p$, as in Lemma \ref{editorial-source-lemma-in-image}.
Here is the precise prime construction supplied by that equality.
The ring
$$
F=(R\setminus\mathfrak p)^{-1}(B/\mathfrak pB)
$$
is nonzero: if $1=0$ in $F$, some $r\in R\setminus\mathfrak p$
has zero image in $B/\mathfrak pB$, contradicting the equality.
A prime of $F$ pulls back through the quotient and localization
maps to a prime of $B$ containing $\mathfrak pB$ and avoiding
the image of $R\setminus\mathfrak p$. Its inverse image
$\mathfrak q$ under $S\to B$ has contraction $\mathfrak p$
in $R$. Every element of $S\setminus\mathfrak q'$ is a unit
in $B$, so none belongs to $\mathfrak q$ and
$\mathfrak q\subseteq\mathfrak q'$. This proves the desired
conclusion once the equality is established.

The inclusion $\mathfrak p\subseteq\mathfrak pB\cap R$ is
immediate. For the reverse inclusion choose an original
$z\in R$ whose image belongs to $\mathfrak pB$. Write that
image as a finite sum $\sum_{i=1}^r r_i h_i/g_i$ with
$r_i\in\mathfrak p$, $h_i\in S$ and $g_i\notin\mathfrak q'$.
Take
$$
g=\prod_{i=1}^r g_i,\qquad
y=\sum_{i=1}^r r_i h_i\prod_{k\ne i}g_k\in\mathfrak pS.
$$
Then $g\notin\mathfrak q'$ and $z=y/g$ in $B$. Since $S$
is a domain, the map $S\to B$ is injective, and the exact
identity $zg=y$ holds in $S$ itself. For the empty sum use
$g=1$ and $y=0$.

\medskip\noindent
By Lemma \ref{editorial-source-lemma-integral-integral-over-ideal}
there exists a monic polynomial
$P = x^m + b_{m-1} x^{m-1} + \ldots + b_0$
with $b_i\in\mathfrak p^{m-i}\subseteq\mathfrak p$
for every $0\leq i<m$, such that $P(y)=0$.
These are the full coefficient powers in the original definition
of integrality over $\mathfrak p$.

\medskip\noindent
By Lemma \ref{editorial-source-lemma-minimal-polynomial-normal-domain}, the
minimal polynomial of $g$ over $K=\operatorname{Frac}(R)$ has
coefficients in $R$. Write it as
$Q=x^n+\sum_{j=0}^{n-1}a_jx^j$. If every $a_j$ belonged to
$\mathfrak p$, its original equation would give
$$
g^n=-\sum_{j=0}^{n-1}a_jg^j\in\mathfrak pS
\subseteq\mathfrak p'S\subseteq\mathfrak q'.
$$
Since $n\geq1$ and $\mathfrak q'$ is prime, this would give
$g\in\mathfrak q'$, a contradiction. Therefore at least one
of the original $a_j$ is outside $\mathfrak p$.

\medskip\noindent
If $z=0$, the desired membership $z\in\mathfrak p$ already
holds. Hence assume $z\ne0$, so $z$ is a unit in the original
$K$, and retain $y=zg$. The coefficient-preserving substitution
map $K[x]\to K[x]$, $H(x)\mapsto H(zx)$, has inverse
$H(x)\mapsto H(x/z)$. Set
$$
Q'(x)=z^nQ(x/z)
=x^n+\sum_{j=0}^{n-1}z^{n-j}a_jx^j.
$$
This is monic of degree $n$ and satisfies
$Q'(y)=z^nQ(g)=0$. If $H(y)=0$ for $H\in K[x]$, then
$H(zg)=0$, so minimality of $Q$ gives $H(zx)=Q(x)D(x)$
for some $D\in K[x]$. The inverse substitution yields
$$
H(x)=Q(x/z)D(x/z)=Q'(x)\bigl(z^{-n}D(x/z)\bigr).
$$
Thus $Q'$ divides every annihilating polynomial of $y$ and is
its monic minimal polynomial; in particular $Q'$ divides the
original $P$. All coefficients of $P$ and $Q'$ belong to $R$.
Part (2) of Lemma \ref{editorial-source-lemma-polynomials-divide} gives
$$
z^ja_{n-j}\in\sqrt{(b_0,\ldots,b_{m-1})R}
\subseteq\mathfrak p,\qquad 1\leq j\leq n.
$$
The full powers $b_i\in\mathfrak p^{m-i}$ also yield the
stronger statement that each $z^ja_{n-j}$ is integral over
$\mathfrak p^j$, by
Lemma \ref{editorial-source-lemma-weighted-monic-factor-coefficients}.
Choose $a_{n-j}\notin\mathfrak p$ from the previous paragraph.
Primality first gives $z^j\in\mathfrak p$ and then
$z\in\mathfrak p$, since $j\geq1$. The required contraction
equality, and therefore the prime constructed above, follow.
\end{proof}


\subsection{Editorial treatment of Algebra lines 8587--9307}

\hypertarget{algebra-context-097}{}

\noindent\href{https://github.com/stacks/stacks-project/blob/a04446e57ec1fbc252a871afcec7752fb2807b14/algebra.tex\#L8587}{Original source, lines 8587--9307}. Complete original and reviewed TeX passages accompany this note. Every intervention is separately recorded; the following treatment is editorial.


\subsubsection{Flat modules and flat ring maps}
\label{editorial-source-section-flat}

\noindent
One often used result is that if $M = \colim_{i\in \mathcal{I}} M_i$
is a colimit of $R$-modules and if $N$ is an $R$-module then
$$
M \otimes N
=
\colim_{i\in \mathcal{I}} M_i \otimes_R N,
$$
see Lemma \ref{editorial-source-lemma-tensor-products-commute-with-limits}.
This property is usually expressed by saying
that {\it $\otimes$ commutes with colimits}.
Another often used result is that if $0 \to N_1 \to N_2 \to N_3 \to 0$
is an exact sequence and if $M$ is any $R$-module, then
$$
M \otimes_R N_1
\to
M \otimes_R N_2
\to
M \otimes_R N_3
\to
0
$$
is still exact, see Lemma \ref{algebra-lemma-tensor-product-exact}.
Both of these properties tell us that the functor
$N \mapsto M \otimes_R N$ {\it is right exact}.
See Categories, Section \ref{categories-section-exact-functor}
and Homology, Section \ref{homology-section-functors}.
An $R$-module $M$ is flat if $N \mapsto N \otimes_R M$ is also left exact,
i.e., if it is exact. Here is the precise definition.

\begin{definition}
\label{editorial-source-definition-flat}
Let $R$ be a ring.
\begin{enumerate}
\item An $R$-module $M$ is called {\it flat} if whenever
$N_1 \to N_2 \to N_3$ is an exact sequence of $R$-modules
the sequence $M \otimes_R N_1 \to M \otimes_R N_2 \to M \otimes_R N_3$
is exact as well.
\item An $R$-module $M$ is called {\it faithfully flat} if the
complex of $R$-modules
$N_1 \to N_2 \to N_3$ is exact if and only if
the sequence $M \otimes_R N_1 \to M \otimes_R N_2 \to M \otimes_R N_3$
is exact.
\item A ring map $R \to S$ is called {\it flat} if
$S$ is flat as an $R$-module.
\item A ring map $R \to S$ is called {\it faithfully flat} if
$S$ is faithfully flat as an $R$-module.
\end{enumerate}
\end{definition}

\noindent
Here is an example of how you can use the flatness condition.

\begin{lemma}
\label{editorial-source-lemma-flat-intersect-ideals}
Let $R$ be a ring. Let $I, J \subset R$ be ideals. Let $M$ be a flat
$R$-module. Then $IM\cap JM=(I\cap J)M$.
More generally, if $A_1,\ldots,A_r$ are submodules of an
$R$-module $N$, their tensors with $M$ identify with
submodules of $N\otimes_R M$, and
$$
\left(\bigcap_{i=1}^r A_i\right)\otimes_R M
\ \cong\ \bigcap_{i=1}^r(A_i\otimes_R M)
$$
under these original inclusions, for every finite $r\geq1$.
\end{lemma}

\begin{proof}
In the original ideal case use the exact sequence
$$
0\longrightarrow I\cap J\longrightarrow R
\longrightarrow R/I\oplus R/J,
$$
whose last map is $r\mapsto(r+I,r+J)$. Tensoring with $M$
preserves exactness. The maps
$$
R\otimes_R M\longrightarrow M,\quad r\otimes m\longmapsto rm,
\qquad
(R/I)\otimes_R M\longrightarrow M/IM,\quad
(r+I)\otimes m\longmapsto rm+IM
$$
are isomorphisms with inverses $m\mapsto1\otimes m$ and
$m+IM\mapsto(1+I)\otimes m$ respectively. The latter is
well-defined since each $a m$ with $a\in I$ maps to
$(a+I)\otimes m=0$; their composites fix the original tensor
generators and residue classes. The same formula applies to
$J$. Tensor product respects the indicated finite direct sum
by the component projections and inclusions. Thus the exact
sequence becomes
$$
0\longrightarrow (I\cap J)\otimes_R M\longrightarrow M
\longrightarrow M/IM\oplus M/JM.
$$
The middle kernel is $IM\cap JM$. The first map sends
$a\otimes m$ to $am$, so its image is exactly $(I\cap J)M$.
This proves the original equality.

For the general statement, tensoring $0\to A_i\to N$ with
$M$ proves that each original map $A_i\otimes_R M\to N\otimes_R M$
is injective. For two submodules $A,B\subseteq N$, use
$$
0\longrightarrow A\cap B\longrightarrow N
\longrightarrow (N/A)\oplus(N/B),\qquad
n\longmapsto(n+A,n+B).
$$
After tensoring, exactness identifies
$(A\cap B)\otimes_R M$ with the kernel of the two quotient
maps. Right exactness for $A\to N\to N/A\to0$ identifies
the kernel of the first quotient map with the image of
$A\otimes_R M$; the sequence for $B$ identifies the other
kernel. The kernel of the pair is their intersection, proving
the displayed isomorphism for two submodules. Its map is
induced by their original inclusions, so induction, intersecting
the first $r-1$ submodules with $A_r$, proves the finite case.
For $r=1$ it is the identity under the given inclusion.
\end{proof}

\begin{lemma}
\label{editorial-source-lemma-colimit-flat}
Let $R$ be a ring. Let $\{M_i, \varphi_{ii'}\}$ be a directed system of
flat $R$-modules. Then $\colim_i M_i$ is a flat $R$-module.
\end{lemma}

\begin{proof}
Write $M=\colim_i M_i$ with its original structure maps
$\lambda_i:M_i\to M$. For each module $N$, the natural
isomorphism
$$
\colim_i(N\otimes_R M_i)\longrightarrow N\otimes_R M,
\qquad [n\otimes m_i]\longmapsto n\otimes\lambda_i(m_i)
$$
is Lemma \ref{editorial-source-lemma-tensor-products-commute-with-limits};
it commutes with every map of $N$ because both composites
send the displayed generator to the tensor of its original
image. Given an exact complex $N_1\to N_2\to N_3$, flatness
of each $M_i$ makes its tensor complex exact at every stage.
Lemma \ref{algebra-lemma-directed-colimit-exact} gives exactness of
their directed colimit, and the displayed natural isomorphisms
identify it with the tensor complex for $M$. This proves
flatness without assuming that any transition map is injective.
\end{proof}

\begin{example}
\label{editorial-source-example-flat-infinite-intersections}
The finite-intersection assertion above cannot be extended to
arbitrary families of ideals for all flat modules. Take
$R=\mathbf Z$, $M=\mathbf Q$ and $I_n=2^n\mathbf Z$ for
$n\geq1$. The module $\mathbf Q$ is flat by
Lemma \ref{editorial-source-lemma-colimit-flat}: it is the colimit of copies
of $\mathbf Z$ indexed by $n\geq1$, with transition from
$n$ to $n+1$ multiplication by $n+1$ and structure map
$a\mapsto a/n!$. These maps are compatible. Every rational
$a/b$, $b\ne0$, has such a representation since $|b|$ divides
$n!$ for $n\geq|b|$. Equality $a/n!=a'/m!$ is equivalent
after a common stage $k\geq n,m$ to
$a(k!/n!)=a'(k!/m!)$, which is exactly the colimit relation.
Each copy of $\mathbf Z$ is flat because its tensor functor
is the identity via $a\otimes n\mapsto an$ and inverse
$n\mapsto1\otimes n$.

An integer divisible by every $2^n$ must be zero: for a
nonzero integer $a$, choose $n$ with $2^n>|a|$; it cannot
be a nonzero multiple of $2^n$. Thus $\bigcap_{n\geq1}I_n=0$.
On the other hand $I_nM=\mathbf Q$, since every original
$q\in\mathbf Q$ equals $2^n(q/2^n)$. Therefore
$$
\left(\bigcap_{n\geq1}I_n\right)M=0
\quad\hbox{whereas}\quad
\bigcap_{n\geq1}(I_nM)=\mathbf Q.
$$
\end{example}

\begin{lemma}
\label{editorial-source-lemma-composition-flat}
A composition of (faithfully) flat ring maps is
(faithfully) flat.
If $R \to R'$ is (faithfully) flat, and $M'$ is a
(faithfully) flat $R'$-module, then $M'$ is a
(faithfully) flat $R$-module.
\end{lemma}

\begin{proof}
For every original $R$-module $N$ there are mutually inverse
maps
$$
\begin{aligned}
\theta_N:M'\otimes_{R'}(R'\otimes_R N)&\longrightarrow M'\otimes_R N,
&m'\otimes(r'\otimes n)&\longmapsto r'm'\otimes n,\\
\eta_N:M'\otimes_R N&\longrightarrow M'\otimes_{R'}(R'\otimes_R N),
&m'\otimes n&\longmapsto m'\otimes(1\otimes n).
\end{aligned}
$$
The first respects $R'$-balancing because multiplication by
an $R'$-scalar can move between $m'$ and $r'$. It respects
$R$-balancing in the inner tensor by the restricted $R$-action
on $M'$. For the inverse, the equality
$(rm')\otimes(1\otimes n)=m'\otimes(\varphi(r)\otimes n)
=m'\otimes(1\otimes rn)$ proves $R$-balancing.
The composite $\theta_N\eta_N$ fixes $m'\otimes n$, and
$\eta_N\theta_N$ sends a generator to
$r'm'\otimes(1\otimes n)=m'\otimes(r'\otimes n)$.
Both maps commute with each homomorphism of $N$ on those
generators, so they are natural. Here $\varphi:R\to R'$
is the original ring map.

If $R'$ is flat over $R$ and $M'$ flat over $R'$, apply
their tensor functors successively to an exact complex
$N_1\to N_2\to N_3$. Both steps preserve exactness, and
$\theta_{N_i}$ identifies the resulting complex with
$M'\otimes_R N_1\to M'\otimes_R N_2\to M'\otimes_R N_3$.
This proves flatness over $R$. If both hypotheses are faithful,
an arbitrary original complex is exact if and only if the
first tensor complex is exact, and this in turn holds if
and only if the second tensor complex is exact. The same
natural maps prove faithful flatness over $R$. Taking $M'$
to be the target ring of a second ring map gives the first
assertion, for flat and faithfully flat maps alike.
\end{proof}

\begin{lemma}
\label{editorial-source-lemma-flat}
Let $M$ be an $R$-module. The following are equivalent:
\begin{enumerate}
\item
\label{editorial-source-item-flat}
$M$ is flat over $R$.
\item
\label{editorial-source-item-injective}
for every injection of $R$-modules $N \subset N'$
the map $N \otimes_R M \to N'\otimes_R M$ is injective.
\item
\label{editorial-source-item-f-ideal}
for every ideal $I \subset R$ the map
$I \otimes_R M \to R \otimes_R M = M$ is injective.
\item
\label{editorial-source-item-ffg-ideal}
for every finitely generated ideal $I \subset R$
the map $I \otimes_R M \to R \otimes_R M = M$ is injective.
\end{enumerate}
\end{lemma}

\begin{proof}
The successive implications
$(\ref{editorial-source-item-flat})\Rightarrow(\ref{editorial-source-item-injective})
\Rightarrow(\ref{editorial-source-item-f-ideal})\Rightarrow(\ref{editorial-source-item-ffg-ideal})$
follow directly: apply flatness to $0\to N\to N'$, then
specialize to inclusions $I\subseteq R$, and then to finitely
generated ideals. We prove
$(\ref{editorial-source-item-ffg-ideal})\Rightarrow(\ref{editorial-source-item-flat})$.
Suppose that $N_1 \to N_2 \to N_3$ is exact.
Let $K = \Ker(N_2 \to N_3)$ and $Q = \Im(N_2 \to N_3)$.
Then we get maps
$$
N_1 \otimes_R M \to
K \otimes_R M \to
N_2 \otimes_R M \to
Q \otimes_R M \to
N_3 \otimes_R M
$$
Observe that the first and third arrows are surjective. Thus if we show
that the second and fourth arrows are injective, then we are
done\footnote{Here is the argument in more detail:
Assume that we know that the second and fourth arrows are
injective. Lemma \ref{algebra-lemma-tensor-product-exact} (applied
to the exact sequence $K \to N_2 \to Q \to 0$) yields that
the sequence $K \otimes_R M \to N_2 \otimes_R M \to
Q \otimes_R M \to 0$ is exact. Hence,
$\Ker \left(N_2 \otimes_R M \to Q \otimes_R M\right)
= \Im \left(K \otimes_R M \to N_2 \otimes_R M\right)$.
Since
$\Im \left(K \otimes_R M \to N_2 \otimes_R M\right)
= \Im \left(N_1 \otimes_R M \to N_2 \otimes_R M\right)$
(due to the surjectivity of $N_1 \otimes_R M \to
K \otimes_R M$) and
$\Ker \left(N_2 \otimes_R M \to Q \otimes_R M\right)
= \Ker \left(N_2 \otimes_R M \to N_3 \otimes_R M\right)$
(due to the injectivity of $Q \otimes_R M \to
N_3 \otimes_R M$), this becomes
$\Ker \left(N_2 \otimes_R M \to N_3 \otimes_R M\right)
= \Im \left(N_1 \otimes_R M \to N_2 \otimes_R M\right)$,
which shows that the functor $- \otimes_R M$ is exact,
whence $M$ is flat.}.
Hence it suffices to show that $- \otimes_R M$ transforms
injective $R$-module maps into injective $R$-module maps.

\medskip\noindent
Assume $K \to N$ is an injective $R$-module map and
let $x \in \Ker(K \otimes_R M \to N \otimes_R M)$.
We have to show that $x$ is zero.
The $R$-module $K$ is the union of its finite
$R$-submodules; hence, $K \otimes_R M$ is
the colimit of $R$-modules of the form
$K_i \otimes_R M$ where $K_i$ runs over all finite
$R$-submodules of $K$
(because tensor product commutes with colimits).
Thus, for some $i$ our $x$ comes from an element
$x_i \in K_i \otimes_R M$. Thus we may assume that $K$
is a finite $R$-module. Assume this. We regard the
injection $K \to N$ as an inclusion, so that
$K \subset N$.

\medskip\noindent
The $R$-module $N$ is the union of its finite
$R$-submodules that contain $K$. Hence, $N \otimes_R M$
is the colimit of $R$-modules of the form
$N_i \otimes_R M$ where $N_i$ runs over all finite
$R$-submodules of $N$ that contain $K$
(again since tensor product commutes with colimits).
Notice that this is a colimit over a directed system
(since the sum of two finite submodules of $N$ is
again finite).
Hence, (by Lemma \ref{algebra-lemma-zero-directed-limit})
the element $x \in K \otimes_R M$ maps to
zero in at least one of these $R$-modules
$N_i \otimes_R M$ (since $x$ maps to zero
in $N \otimes_R M$).
Thus we may assume $N$ is a finite $R$-module.

\medskip\noindent
Assume $N$ is a finite $R$-module. Choose a surjection
$R^{\oplus n}\to N$, put $L$ equal to its kernel, and let
$L'$ be the inverse image of the original $K\subseteq N$.
Then $L\subseteq L'\subseteq R^{\oplus n}$,
$N=R^{\oplus n}/L$ and $K=L'/L$, under the quotient maps
just specified. For any submodule $G\subseteq R^{\oplus n}$,
the map $G\otimes_R M\to M^{\oplus n}$ sends
$(g_1,\ldots,g_n)\otimes m$ to $(g_1m,\ldots,g_nm)$.
Suppose these maps are injective for $G=L,L'$, and denote
them by $\beta$ and $\beta'$ respectively. Tensoring the
original inclusion gives the map in the direction
$$
\alpha:L\otimes_R M\longrightarrow L'\otimes_R M,
\qquad \ell\otimes m\longmapsto\ell\otimes m,
\qquad \beta'\alpha=\beta.
$$
It is injective since $\beta$ is injective. Right exactness
identifies the original $K\otimes_R M$ with
$(L'\otimes_R M)/\alpha(L\otimes_R M)$ and $N\otimes_R M$
with $M^{\oplus n}/\beta(L\otimes_R M)$.
The induced quotient map is injective: if
$v\in L'\otimes_R M$ has $\beta'(v)=\beta(u)$ for some
$u\in L\otimes_R M$, then
$\beta'(v-\alpha(u))=0$, hence $v=\alpha(u)$ by injectivity
of $\beta'$. This proves the required injection using the
original inclusions, without a reverse map from $L'$ to $L$.

\medskip\noindent
Thus it suffices to show that $L \otimes_R M \to M^{\oplus n}$
is injective when $L \subset R^{\oplus n}$ is an $R$-submodule.
We do this by induction on $n$. If $n=0$, both $L$ and
$M^{\oplus n}$ are zero, so the assertion holds. The base
case $n=1$ is proved below.
For the induction step assume $n > 1$ and set
$L'=L\cap(R\oplus0^{\oplus n-1})$. It is isomorphic to
the ideal $\{a\in R:(a,0,\ldots,0)\in L\}$ by the first
coordinate map, with inverse $a\mapsto(a,0,\ldots,0)$.
The last $n-1$ coordinate projection on $L$ has kernel $L'$,
so it induces the injection
$L''=L/L'\to R^{\oplus n-1}$ sending the class of
$(a_1,\ldots,a_n)$ to $(a_2,\ldots,a_n)$.
It is well-defined since elements of $L'$ have all these
coordinates zero, and injective by the kernel equality.
We obtain a diagram
$$
\xymatrix{
&
L' \otimes_R M \ar[r] \ar[d] &
L \otimes_R M \ar[r] \ar[d] &
L'' \otimes_R M \ar[r] \ar[d] &
0 \\
0 \ar[r] &
M \ar[r] &
M^{\oplus n} \ar[r] &
M^{\oplus n - 1} \ar[r] & 0
}
$$
The left vertical map is induced by the first coordinate,
the right by the last $n-1$ coordinates, and the lower maps
are the corresponding inclusion and projection. The upper
row is right exact, the lower row is exact, and the squares
commute on every tensor generator. The left and right
vertical maps are injective by the base case and induction.
If $v\in L\otimes_R M$ maps to zero in $M^{\oplus n}$,
its image in $L''\otimes_R M$ maps to zero on the right and
is therefore zero. Lift $v$ to $u\in L'\otimes_R M$ using
upper-row exactness. Its image in $M$ maps to zero in
$M^{\oplus n}$ and hence is zero, since the lower left map
is injective. Left vertical injectivity now gives $u=0$ and
then $v=0$, proving the middle injectivity.

\medskip\noindent
The base case of the induction above is when $L \subset R$ is an ideal.
In other words, we have to show that $I \otimes_R M \to M$ is injective
for any ideal $I$ of $R$. We know this is true when $I$ is finitely
generated. However, $I = \bigcup I_\alpha$ is the union of the
finitely generated ideals $I_\alpha$ contained in it. In other words,
$I=\colim_\alpha I_\alpha$, over the directed set of finitely
generated subideals; the sum of two such subideals contains
both. Tensor-colimit compatibility gives
$I\otimes_R M=\colim_\alpha(I_\alpha\otimes_R M)$.
If an element of this colimit maps to zero in $M$, choose
a representative $v_\alpha\in I_\alpha\otimes_R M$.
Its image in $M$ is already zero. The assumed injectivity
at that stage gives $v_\alpha=0$, so the original colimit
element is zero. This proves the base case for every ideal,
and completes the stated criterion.
\end{proof}

\begin{lemma}
\label{editorial-source-lemma-colimit-rings-flat}
Let $\{R_i, \varphi_{ii'}\}$ be a system of rings over the directed set $I$.
Let $R = \colim_i R_i$.
\begin{enumerate}
\item If $M$ is an $R$-module such that $M$ is flat as an $R_i$-module
for all $i$, then $M$ is flat as an $R$-module.
\item For $i \in I$ let $M_i$ be a flat $R_i$-module and
for $i' \geq i$ let $f_{ii'} : M_i \to M_{i'}$ be a $\varphi_{ii'}$-linear
map such that $f_{i' i''} \circ f_{i i'} = f_{i i''}$. Then
$M = \colim_{i \in I} M_i$ is a flat $R$-module.
\end{enumerate}
\end{lemma}

\begin{proof}
Write $\lambda_i:R_i\to R$ for the original ring maps and
$\mu_i:M_i\to M$ for the module structure maps. Part (1)
is the instance of part (2) with every $M_i=M$, acted on
through $\lambda_i$, and all underlying transition maps the
identity. They are semilinear because
$\lambda_{i'}\varphi_{ii'}=\lambda_i$.

We first specify the colimit constructions in part (2).
Only a nonempty directed indexing set is needed; none of
the arguments uses antisymmetry of its order. For a system
of abelian groups $B_i$ with transition maps $b_{ij}$,
the colimit is the direct sum of the $B_i$ modulo the
subgroup generated by the original relations
$\iota_i(x)-\iota_j(b_{ij}(x))$ for $i\leq j$.
Every element has a representative at a single stage:
move its finitely many summands to a common upper stage
and add their images there. Two representatives $x_i,x_j$
are equal in the colimit exactly when their images agree
at some common upper stage. To prove the nontrivial
direction, express their difference as a finite sum of
the generating relations and move every stage occurring
in that finite expression to one upper stage; compatibility
kills each relation there. The converse follows from the
defining transition relations. In particular a stage
element represents zero exactly when its image is zero
at a later stage. This description permits noninjective
transition maps throughout.

For $r=\lambda_i(r_i)$ and $m=\mu_j(m_j)$ choose $k\geq i,j$
and define the original scalar action by
$$
rm=\mu_k\bigl(\varphi_{ik}(r_i)f_{jk}(m_j)\bigr).
$$
Changing either representative or the chosen upper stage
does not change this element: move the relevant equalities
to one further stage and use semilinearity and compatibility.
The distributive laws and associativity of the action follow
by moving all finitely many entries of each identity to one
stage and applying its module axiom. The unit acts as the
identity because $[f_{ik}(m_i)]=[m_i]$.

Here is the varying-base tensor comparison needed below.
Let $A_i$ be $R_i$-modules with compatible semilinear maps
$h_{ij}:A_i\to A_j$, and put $A=\colim_i A_i$ with the
same scalar construction and structure maps $\alpha_i$.
There are transition maps
$$
\tau_{ij}:A_i\otimes_{R_i}M_i\longrightarrow A_j\otimes_{R_j}M_j,
\qquad a_i\otimes m_i\longmapsto h_{ij}(a_i)\otimes f_{ij}(m_i).
$$
They are well-defined: the two images of the balancing
relation $r_ia_i\otimes m_i=a_i\otimes r_im_i$ agree by
$$
\varphi_{ij}(r_i)h_{ij}(a_i)\otimes f_{ij}(m_i)
=h_{ij}(a_i)\otimes\varphi_{ij}(r_i)f_{ij}(m_i).
$$
Let $T=\colim_i(A_i\otimes_{R_i}M_i)$ with maps $\theta_i$.
The first comparison is
$$
F:T\longrightarrow A\otimes_R M,\qquad
\theta_i\left(\sum_{\ell=1}^q a_{i,\ell}\otimes m_{i,\ell}\right)
\longmapsto\sum_{\ell=1}^q
\alpha_i(a_{i,\ell})\otimes\mu_i(m_{i,\ell}).
$$
Each stage formula respects balancing and every transition
relation, so it defines a map of abelian groups. The scalar
formula above makes it $R$-linear by moving a scalar and
the tensor representative to one common stage.
For the inverse, define an additive balanced pairing on
representatives by
$$
B\bigl(\alpha_i(a_i),\mu_j(m_j)\bigr)
=\theta_k\bigl(h_{ik}(a_i)\otimes f_{jk}(m_j)\bigr),
\qquad k\geq i,j.
$$
Its value is independent of representatives and $k$:
move any two eventual equalities of representatives and
both choices of $k$ to one further stage, where the tensors
are identical. Additivity is checked at a common stage for
the two summands. For $r=\lambda_\ell(r_\ell)$, move
$r,a,m$ to one stage $k$. Then
$$
B(ra,m)=\theta_k\bigl(\varphi_{\ell k}(r_\ell)a_k\otimes m_k\bigr)
=\theta_k\bigl(a_k\otimes\varphi_{\ell k}(r_\ell)m_k\bigr)
=B(a,rm).
$$
Thus $B$ induces an $R$-linear map $G:A\otimes_R M\to T$.
On a pure tensor $FG$ gives
$\alpha_k(h_{ik}(a_i))\otimes\mu_k(f_{jk}(m_j))
=\alpha_i(a_i)\otimes\mu_j(m_j)$.
On a stage tensor $GF$ gives its image at a common later
stage, which is the same class $\theta_i(a_i\otimes m_i)$.
Finite sums of these generators prove both composites are
the identity. Hence the exact comparison, with both maps
specified, is
$$
\colim_i(A_i\otimes_{R_i}M_i)
\ \cong\ (\colim_i A_i)\otimes_R(\colim_i M_i).
$$

For completeness, stagewise injections induce an injection
on these directed colimits. Indeed, for a compatible family
of injections $u_i:X_i\to Y_i$, if $[x_i]$ maps to zero,
then at a later stage $j$ compatibility gives $u_j(x_j)=0$.
Stagewise injectivity gives $x_j=0$, hence $[x_i]=0$.
This uses no injectivity of the transition maps.

Now retain the original finitely generated ideal
$\mathfrak a=(a_1,\ldots,a_n)\subseteq R$. Choose one stage
$i$ containing representatives $a_{\ell,i}$ of all its
generators, and put
$$
\mathfrak a'=(a_{1,i},\ldots,a_{n,i})\subseteq R_i,
\qquad
\mathfrak a_j=\mathfrak a'R_j
=(\varphi_{ij}(a_{1,i}),\ldots,\varphi_{ij}(a_{n,i}))
\quad(j\geq i).
$$
The tail $j\geq i$ is cofinal: every original representative,
and every eventual equality, can be moved to a stage above
both its original stage and $i$. It has the same ring and
module colimits. The map
$$
\delta:\colim_{j\geq i}\mathfrak a_j\longrightarrow\mathfrak a,
\qquad [x_j]\longmapsto\lambda_j(x_j)
$$
is surjective. For $x=\sum_{\ell=1}^n r_\ell a_\ell$,
choose a stage $j\geq i$ representing all $r_\ell$ and use
$x_j=\sum_{\ell=1}^n r_{\ell,j}\varphi_{ij}(a_{\ell,i})$.
It is injective: if $\lambda_j(x_j)=0$ in $R$, some
$k\geq j$ has $\varphi_{jk}(x_j)=0$ in $R_k$; because
$\mathfrak a_k$ is an actual ideal in that ring, the image
is zero in $\mathfrak a_k$ itself. This proves that the
inverse formula $x\mapsto[x_j]$ is independent of its
expression and representatives; equal expressions differ
by an element eventually zero in the original ring system.
Both maps respect the scalar action and fix the represented
elements in their composites.

Apply the proved tensor comparison to $A_j=\mathfrak a_j$.
It identifies
$$
\colim_{j\geq i}(\mathfrak a_j\otimes_{R_j}M_j)
\ \cong\ \mathfrak a\otimes_R M.
$$
For every $j$, flatness of $M_j$ makes
$\mathfrak a_j\otimes_{R_j}M_j\to M_j$ injective. Its
formula is $x_j\otimes m_j\mapsto x_jm_j$, and compatibility
of these maps is exactly semilinearity of $f_{jk}$.
Their colimit is injective by the element argument above.
Under the displayed comparison it is the original map
$\mathfrak a\otimes_R M\to M$, $a\otimes m\mapsto am$.
Lemma \ref{editorial-source-lemma-flat} now proves flatness of $M$.
The zero ideal, including a presentation with no generators,
has zero tensor product and requires no exception.

The comparison also proves the full injection criterion
directly. For an injection $N'\to N$ of $R$-modules, regard
both as constant underlying systems with restricted
$R_i$-actions and identity transition maps. Each
$N'\otimes_{R_i}M_i\to N\otimes_{R_i}M_i$ is injective.
Their colimit comparison is exactly
$N'\otimes_R M\to N\otimes_R M$, hence it is injective.
This gives a second proof with all original modules retained.

The explicitly displayed composition law for $f_{ij}$ also
suffices if identities on the diagonal were not stipulated.
Indeed $e_i=f_{ii}$ is an idempotent $R_i$-linear map and
$$
M_i=\operatorname{Im}(e_i)\oplus\operatorname{Ker}(e_i),
\qquad m_i=e_i(m_i)+(m_i-e_i(m_i)).
$$
The intersection of the two summands is zero since $e_i$
is the identity on its image. Each summand of a flat module
is flat: its tensor map on any injection is a direct summand
of the injective tensor map for $M_i$. Compatibility gives
$f_{ij}=e_jf_{ij}=f_{ij}e_i$, so the restrictions to
$\operatorname{Im}(e_i)$ form a genuine directed system.
The colimit by all the original relations is isomorphic to
the colimit of these images. The maps send $[m_i]$ to
$[e_i(m_i)]$ and an image representative to $[m_i]$ in the
original colimit; they respect transitions by the displayed
identities and are inverse because $[e_i(m_i)]=[m_i]$ is
one of the original relations. Thus no extra identity
hypothesis is needed for the flatness conclusion.

\medskip\noindent
We record exactly what is additionally needed for faithful
flatness. For a maximal ideal $\mathfrak m\subset R$ set
$\mathfrak m_i=\lambda_i^{-1}(\mathfrak m)$. Each is a
proper prime, though it need not be maximal in $R_i$, and
$\varphi_{ij}^{-1}(\mathfrak m_j)=\mathfrak m_i$.
Consequently the induced maps $R_i/\mathfrak m_i\to R_j/\mathfrak m_j$
are injective. The comparison
$$
\colim_i(R_i/\mathfrak m_i)\longrightarrow R/\mathfrak m,
\qquad [r_i+\mathfrak m_i]\longmapsto\lambda_i(r_i)+\mathfrak m
$$
is surjective by representatives and injective because an
element mapping to zero already has $r_i\in\mathfrak m_i$.
Its inverse sends any representative $\lambda_i(r_i)+\mathfrak m$
to $[r_i+\mathfrak m_i]$; the same zero test proves independence
of the choice. Applying the varying-base tensor comparison
gives
$$
M/\mathfrak mM\ \cong\
\colim_i(M_i/\mathfrak m_iM_i).
$$
Here the stage quotient-tensor isomorphism sends
$(r_i+\mathfrak m_i)\otimes u_i$ to
$r_iu_i+\mathfrak m_iM_i$ and has inverse
$u_i+\mathfrak m_iM_i\mapsto(1+\mathfrak m_i)\otimes u_i$.
The same formulas identify the receiving quotient over $R$.
Balancing makes both maps well-defined, and both composites
fix the indicated generators. Thus the original transition
on the right is
$u_i+\mathfrak m_iM_i\mapsto f_{ij}(u_i)+\mathfrak m_jM_j$,
and the receiving map is
$u_i+\mathfrak m_iM_i\mapsto\mu_i(u_i)+\mathfrak mM$.
The eventual-zero rule now proves
$$
\mu_i(u_i)+\mathfrak mM=0
\quad\Longleftrightarrow\quad
f_{ij}(u_i)\in\mathfrak m_jM_j\text{ for some }j\geq i.
$$

A flat module is faithfully flat exactly when all these
maximal-ideal quotients are nonzero. Here is the full
argument, using the definition by exactness reflection.
Faithful flatness detects nonzero modules by applying that
definition to $0\to N\to0$, and in particular detects
$N=R/\mathfrak m$. Conversely suppose all the quotients
are nonzero. Given $0\ne n\in N$, let
$J=\operatorname{Ann}_R(n)$, choose a maximal ideal
$\mathfrak m\supseteq J$, and retain the injection
$R/J\to N$, $r+J\mapsto rn$. The module $M/JM$ surjects
onto the nonzero $M/\mathfrak mM$, so it is nonzero.
Flatness makes $(R/J)\otimes_R M=M/JM\to N\otimes_R M$
injective, proving $N\otimes_R M\ne0$ for every nonzero $N$.

To deduce exactness reflection, take an original complex
$N_1\xrightarrow{d_1}N_2\xrightarrow{d_2}N_3$ and write
$J=\operatorname{Im}(d_1)\subseteq K=\operatorname{Ker}(d_2)$.
Flatness embeds $J\otimes_R M$ and $K\otimes_R M$ in
$N_2\otimes_R M$. The exact sequences defining image and
kernel identify them with the tensor image and tensor
kernel, respectively: the map $N_1\to J$ is surjective,
and $0\to K\to N_2\to N_3$ stays exact. Right exactness
for $J\to K\to K/J\to0$ then identifies
$$
(K/J)\otimes_R M\ \cong\
\operatorname{Ker}(d_2\otimes\mathrm{id}_M)\big/
\operatorname{Im}(d_1\otimes\mathrm{id}_M).
$$
These are the maps induced by the original inclusions and
quotient $K\to K/J$, so the identification respects those
original submodules. If the tensor complex is exact, this
quotient is zero. Nonzero-module detection gives $K/J=0$,
hence the original complex is exact. This proves the criterion.

Combining it with the exact quotient comparison proves:
$M$ is faithfully flat over $R$ if and only if, for every
maximal $\mathfrak m\subset R$, there are an index $i$ and
$u_i\in M_i$ such that
$$
f_{ij}(u_i)\notin\mathfrak m_jM_j
\quad\hbox{for every }j\geq i.
$$
This quantifier keeps the same original element at every
later stage. It cannot be replaced by merely having a
nonzero quotient at each stage.

For example take constant $R_i=k$ and $M_i=k$ over a
nonzero field, indexed by $i\in\mathbf Z_{\geq0}$, with
identity ring maps and $f_{ij}=0$ for $i<j$, $f_{ii}$ the
identity. Each stage module is faithfully flat, but every
element dies at the next stage, so $R=k$ and $M=0$.
Even injective module transitions do not suffice. Take
$R_i=M_i=\mathbf Z$ with identity ring maps and
$f_{ij}(a)=2^{j-i}a$. The map $[a\text{ at }i]\mapsto a/2^i$
identifies $M$ with $\mathbf Z[1/2]$: it is surjective by
the description of its elements, and equality of two
fractions is exactly equality after multiplication to a
common power-of-two denominator. Every stage is faithfully
flat and every transition injective, but $M/2M=0$ since
$a/2^i=2(a/2^{i+1})$, so $M$ is not faithfully flat.

A useful sufficient condition follows with the actual
residue maps retained. Suppose every $M_i$ is faithfully
flat and, for each maximal $\mathfrak m\subset R$, there
is a tail $i\geq i_0$ on which all maps
$M_i/\mathfrak m_iM_i\to M_j/\mathfrak m_jM_j$ are injective.
Faithful flatness at $i_0$ makes its quotient nonzero because
$R_{i_0}/\mathfrak m_{i_0}\ne0$. A nonzero element there
remains nonzero at every later stage, which is precisely
the necessary and sufficient condition just proved.

One condition on the original transitions that guarantees
this is purity of the scalar-extended maps
$$
g_{ij}:R_j\otimes_{R_i}M_i\longrightarrow M_j,
\qquad r_j\otimes u_i\longmapsto r_jf_{ij}(u_i).
$$
Here purity means that tensoring this map with every
$R_j$-module gives an injection. The formula is balanced
by semilinearity, so $g_{ij}$ is an $R_j$-linear map.
The contraction equality gives an injection of
$R_i$-modules $R_i/\mathfrak m_i\to R_j/\mathfrak m_j$.
Flatness of $M_i$ and then purity give injections
$$
(R_i/\mathfrak m_i)\otimes_{R_i}M_i
\longrightarrow (R_j/\mathfrak m_j)\otimes_{R_i}M_i
\longrightarrow (R_j/\mathfrak m_j)\otimes_{R_j}M_j.
$$
For the second map the associativity comparison sends
$(s_j+\mathfrak m_j)\otimes(r_j\otimes u_i)$ to
$(s_jr_j+\mathfrak m_j)\otimes u_i$ and has inverse
$(s_j+\mathfrak m_j)\otimes u_i\mapsto
(s_j+\mathfrak m_j)\otimes(1\otimes u_i)$; balancing and
these formulas verify both composites. Under the quotient
identifications, the composite injection is exactly
$u_i+\mathfrak m_iM_i\mapsto f_{ij}(u_i)+\mathfrak m_jM_j$.
Thus faithful flatness of the stages together with these
pure transition maps suffices.

If $R=0$, its unital action forces $M=0$. The only unital
module over that ring is zero, so its tensor functor
preserves and reflects exactness; it is faithfully flat in
the stated sense. The maximal-ideal criterion gives the
same result because there are no maximal ideals.
\end{proof}

\begin{lemma}
\label{editorial-source-lemma-flat-base-change}
Suppose that $M$ is (faithfully) flat over $R$, and that $R \to R'$
is a ring map. Then $M \otimes_R R'$ is (faithfully) flat over $R'$.
\end{lemma}

\begin{proof}
For every original $R'$-module $N$, define
$$
\begin{aligned}
\theta_N:N\otimes_{R'}(R'\otimes_R M)&\longrightarrow N\otimes_R M,
&n\otimes(r'\otimes m)&\longmapsto r'n\otimes m,\\
\eta_N:N\otimes_R M&\longrightarrow N\otimes_{R'}(R'\otimes_R M),
&n\otimes m&\longmapsto n\otimes(1\otimes m).
\end{aligned}
$$
The first formula respects the outer $R'$-balancing and the
inner $R$-balancing by the original $R'$-action on $N$.
The inverse respects $R$-balancing because
$(rn)\otimes(1\otimes m)=n\otimes(\varphi(r)\otimes m)
=n\otimes(1\otimes rm)$. Both composites fix the tensor
generators, using $r'n\otimes(1\otimes m)=n\otimes(r'\otimes m)$.
The formulas also commute with every $R'$-linear map of $N$.
The flip $M\otimes_R R'\to R'\otimes_R M$, $m\otimes r'\mapsto r'\otimes m$,
identifies the module in the statement with this one and is
its own inverse after reversing the factors. Exactness of a
complex of $R'$-modules is equivalent to exactness of its
underlying $R$-module complex, since kernels and images are
the same subsets. Apply flatness, or preservation and
reflection of exactness, for the original $M$ over $R$ and
then use the displayed natural maps. This proves both claims.
\end{proof}

\begin{lemma}
\label{editorial-source-lemma-flatness-descends}
Let $R \to R'$ be a faithfully flat ring map.
Let $M$ be a module over $R$, and set $M' = R' \otimes_R M$.
Then $M$ is flat over $R$ if and only if $M'$ is flat over $R'$.
\end{lemma}

\begin{proof}
By Lemma \ref{editorial-source-lemma-flat-base-change} we see that if $M$ is flat
then $M'$ is flat. For the converse, suppose that $M'$ is flat.
Let $N_1 \to N_2 \to N_3$ be an exact sequence of $R$-modules.
We want to show that $N_1 \otimes_R M \to N_2 \otimes_R M \to N_3 \otimes_R M$
is exact. We know that
$N_1 \otimes_R R' \to N_2 \otimes_R R' \to N_3 \otimes_R R'$ is
exact, because $R \to R'$ is flat. Flatness of $M'$ implies that
$N_1 \otimes_R R' \otimes_{R'} M'
\to N_2 \otimes_R R' \otimes_{R'} M'
\to N_3 \otimes_R R' \otimes_{R'} M'$ is exact.
The exact natural comparison for every original $N$ is
$$
\begin{split}
(N\otimes_R R')\otimes_{R'}(R'\otimes_R M)
&\longrightarrow (N\otimes_R M)\otimes_R R',\\
(n\otimes a)\otimes(b\otimes m)&\longmapsto(n\otimes m)\otimes ab.
\end{split}
$$
Its inverse sends $(n\otimes m)\otimes c$ to
$(n\otimes c)\otimes(1\otimes m)$. All $R$-balancing
relations follow from the original coefficient map $R\to R'$;
the $R'$-balancing relation moves a scalar between $a$ and
$b$ and preserves their product $ab$. These observations
make both maps well-defined. One composite fixes
$(n\otimes m)\otimes c$. The other gives
$(n\otimes ab)\otimes(1\otimes m)
=(n\otimes a)\otimes(b\otimes m)$ by $R'$-balancing.
Their formulas commute with every map of $N$. Thus the
exact complex above is the tensor with $R'$ of the original
complex $N_1\otimes_R M\to N_2\otimes_R M\to N_3\otimes_R M$.
Faithful flatness of $R'$ over $R$ reflects exactness, so
this original complex is exact, as required.
\end{proof}

\begin{lemma}
\label{editorial-source-lemma-flatness-descends-more-general}
Let $R$ be a ring. Let $S \to S'$ be a flat map of $R$-algebras.
Let $M$ be a module over $S$, and set $M' = S' \otimes_S M$.
\begin{enumerate}
\item If $M$ is flat over $R$, then $M'$ is flat over $R$.
\item If $S \to S'$ is faithfully flat, then $M$ is flat
over $R$ if and only if $M'$ is flat over $R$.
\end{enumerate}
\end{lemma}

\begin{proof}
Let $f:N\to N'$ be an injection of $R$-modules and put
$K=\Ker(f\otimes\mathrm{id}_M)$. The map
$f\otimes\mathrm{id}_M:N\otimes_R M\to N'\otimes_R M$
is $S$-linear, with the $S$-action on its $M$ factor, so
$K$ is an $S$-module. Tensoring its defining exact sequence
with the flat $S$-module $S'$ identifies $K\otimes_S S'$
with the kernel after that tensor operation. For every $N$
the precise comparison is
$$
\begin{split}
(N\otimes_R M)\otimes_S S'&\longrightarrow
N\otimes_R(S'\otimes_S M),\\
(n\otimes m)\otimes s'&\longmapsto n\otimes(s'\otimes m),
\end{split}
$$
with inverse $n\otimes(s'\otimes m)\mapsto(n\otimes m)\otimes s'$.
The $S$-balancing relations agree on the two sides by moving
the original $S$-scalar between $m$ and $s'$; the $R$-balancing
relations agree through $R\to S\to S'$. The composites fix
the generators, and the formulas commute with $f$. Consequently
the kernel identification is exactly
$$
K\otimes_S S'\ \cong\
\Ker(N\otimes_R M'\longrightarrow N'\otimes_R M').
$$
If $M$ is flat over $R$, then $K=0$, so the right kernel
vanishes for every injection $f$ and Lemma \ref{editorial-source-lemma-flat}
gives flatness of $M'$ over $R$. Conversely, if $M'$ is
flat, then $K\otimes_S S'=0$. When $S'$ is faithfully flat
over $S$, apply reflection of exactness to the complex
$0\to K\to0$: its tensor complex is exact, hence $K=0$.
The same criterion now gives flatness of $M$ over $R$.
\end{proof}

\begin{lemma}
\label{editorial-source-lemma-flat-permanence}
Let $R \to S$ be a ring map. Let $M$ be an $S$-module.
If $M$ is flat as an $R$-module and faithfully flat as an $S$-module,
then $R\to S$ is flat.
More generally, whenever $M$ is faithfully flat over $S$,
the map $R\to S$ is flat (respectively faithfully flat)
if and only if $M$ is flat (respectively faithfully flat)
as an $R$-module.
\end{lemma}

\begin{proof}
For every original $R$-module $N$, the comparison
$$
(N\otimes_R S)\otimes_S M\longrightarrow N\otimes_R M,
\qquad (n\otimes s)\otimes m\longmapsto n\otimes sm
$$
has inverse $n\otimes m\mapsto(n\otimes1)\otimes m$.
Balancing over $S$ preserves the product $sm$, and balancing
over $R$ is compatible with the restricted action on $M$.
Both composites fix the generators; for one of them use
$(n\otimes1)\otimes sm=(n\otimes s)\otimes m$.
These maps are natural in $N$.

If $M$ is flat over $R$, an exact complex
$N_1\to N_2\to N_3$ stays exact after tensoring with $M$.
The comparison identifies this with the result of first
tensoring with $S$ and then with $M$ over $S$. Faithful
flatness over $S$ reflects exactness, so the first tensor
complex is exact. This proves the original conclusion that
$R\to S$ is flat. The converse under the same faithful
$S$-module hypothesis follows from
Lemma \ref{editorial-source-lemma-composition-flat}.

If $M$ is faithfully flat over $R$, take any original complex
$N_1\to N_2\to N_3$ whose tensor with $S$ is exact. Tensoring
that complex with the flat $S$-module $M$ and using the
comparison shows its tensor with $M$ over $R$ is exact.
Faithful flatness over $R$ reflects this to exactness of the
original complex. Together with flatness just proved, this
makes $R\to S$ faithfully flat. Conversely, if $R\to S$
is faithfully flat, composition with the faithfully flat
$S$-module $M$ gives faithful flatness over $R$, again by
Lemma \ref{editorial-source-lemma-composition-flat}.
\end{proof}

\noindent
Let $R$ be a ring.
Let $M$ be an $R$-module.
Let $\sum_{i=1}^n f_ix_i=0$ be a finite relation in $M$,
with $f_i\in R$, $x_i\in M$, and $n\geq0$.
We say the relation $\sum f_i x_i$
is {\it trivial} if there exist an integer $m \geq 0$,
elements $y_j \in M$, $j = 1, \ldots, m$, and elements $a_{ij} \in R$,
$i = 1, \ldots, n$, $j = 1, \ldots, m$ such that
$$
x_i = \sum\nolimits_j a_{ij} y_j, \forall i,
\quad\text{and}\quad
0 = \sum\nolimits_i f_ia_{ij}, \forall j.
$$

\begin{lemma}[Equational criterion of flatness]
\label{editorial-source-lemma-flat-eq}
A module $M$ over $R$ is flat if and only if
every relation in $M$ is trivial.
\end{lemma}

\begin{proof}
Assume $M$ is flat and let $\sum f_i x_i = 0$ be a relation in $M$.
Let $I = (f_1, \ldots, f_n)$, and let
$K = \Ker(R^n \to I, (a_1, \ldots, a_n) \mapsto \sum_i a_i f_i)$.
So we have the short exact sequence
$0 \to K \to R^n \to I \to 0$. Then $\sum f_i \otimes x_i$
is an element of $I \otimes_R M$ which maps
to zero in $R \otimes_R M = M$. By flatness
$\sum f_i \otimes x_i$ is zero in $I \otimes_R M$.
Thus there exists an element of $K \otimes_R M$ mapping
to $\sum e_i \otimes x_i \in R^n \otimes_R M$ where $e_i$
is the $i$th basis element of $R^n$.
Write this element as a finite sum
$\sum_{j=1}^m k_j\otimes y_j$, and write each original
$k_j\in K\subseteq R^n$ as $\sum_{i=1}^n a_{ij}e_i$.
Membership in $K$ gives $\sum_{i=1}^n f_i a_{ij}=0$
for every $j$. Under the coordinate isomorphism
$R^n\otimes_R M\to M^n$, $e_i\otimes y$ maps to the
vector having $y$ in coordinate $i$ and zero elsewhere;
its inverse sends $(x_1,\ldots,x_n)$ to
$\sum_i e_i\otimes x_i$. Comparing all coordinates of
the original lifted tensor gives
$x_i=\sum_{j=1}^m a_{ij}y_j$ for every $i$, exactly the
two defining conditions. The sum may be empty; when $n=0$
there are no coordinates and the empty relation is handled
with $m=0$.

\medskip\noindent
Assume every relation is trivial, let $I$
be a finitely generated ideal, and let $x = \sum f_i \otimes x_i$
be an element of $I \otimes_R M$ mapping to zero in $R \otimes_R M = M$.
This just means exactly that $\sum f_i x_i$ is a relation in
$M$. And the fact that it is trivial implies easily that
$x$ is zero, because
$$
x
=
\sum f_i \otimes x_i
=
\sum f_i \otimes \left(\sum a_{ij}y_j\right)
=
\sum \left(\sum f_i a_{ij}\right) \otimes y_j
=
0
$$
The sums are finite, so the middle equality follows by
additivity and the original balancing relation
$f_i\otimes a_{ij}y_j=(f_i a_{ij})\otimes y_j$ in
$I\otimes_R M$. Thus the map $I\otimes_R M\to M$ is
injective for every finitely generated ideal $I$, and
Lemma \ref{editorial-source-lemma-flat} proves that $M$ is flat.
\end{proof}

\begin{lemma}
\label{editorial-source-lemma-flat-finite-simultaneous-relations}
An $R$-module $M$ is flat if and only if it has the following
property. For every matrix $(f_{\ell i})\in\operatorname{Mat}_{r\times n}(R)$
and elements $x_1,\ldots,x_n\in M$ satisfying
$\sum_{i=1}^n f_{\ell i}x_i=0$ for every $1\leq\ell\leq r$,
there are an integer $m\geq0$, elements $y_1,\ldots,y_m\in M$
and coefficients $a_{ij}\in R$ such that
$$
x_i=\sum_{j=1}^m a_{ij}y_j\quad(1\leq i\leq n),\qquad
\sum_{i=1}^n f_{\ell i}a_{ij}=0
\quad(1\leq\ell\leq r,\ 1\leq j\leq m).
$$
Here $r,n$ may be zero; all matrices and sums are finite.
\end{lemma}

\begin{proof}
Keep the original matrix and let
$\psi:R^n\to R^r$ send the $i$th basis vector to the
column $(f_{1i},\ldots,f_{ri})$. Put $K=\Ker\psi$.
Flatness makes
$K\otimes_R M\to R^n\otimes_R M\to R^r\otimes_R M$
exact. The coordinate tensor isomorphisms in the preceding
proof identify its last map with the original matrix
acting on $M^n$. Thus the given vector is the image of
a finite sum $\sum_{j=1}^m k_j\otimes y_j$.
Write $k_j=\sum_i a_{ij}e_i$. All $r$ original equations
for $k_j\in K$ are precisely the second displayed family,
and coordinate equality gives the first family. No finite
generation assumption on $K$ is used. Conversely, the case
$r=1$ gives the equational criterion
of Lemma \ref{editorial-source-lemma-flat-eq}, hence flatness. When $r=0$
one may take $m=n$, $a_{ij}=\delta_{ij}$ and $y_j=x_j$;
when $n=0$ take $m=0$. These choices verify the boundary
cases with the original empty matrices retained.
\end{proof}

\begin{lemma}
\label{editorial-source-lemma-flat-tor-zero}
Suppose that $R$ is a ring, that $0\to M''\to M'\to M\to0$
is a short exact sequence, and that $N$ is an $R$-module. If $M$ is flat
then $N \otimes_R M'' \to N \otimes_R M'$ is injective, i.e., the
sequence
$$
0 \to N \otimes_R M'' \to N \otimes_R M' \to N \otimes_R M \to 0
$$
is a short exact sequence.
\end{lemma}

\begin{proof}
Retain a surjection $q:F=R^{(I)}\to N$ with kernel
$K$ and inclusion $j:K\to F$. Write $u:M''\to M'$ and
$v:M'\to M$ for the original maps. The coordinate
isomorphism $L\otimes_R F\to L^{(I)}$ for any $L$ sends
$\ell\otimes\sum_i r_ie_i$ to $(r_i\ell)_i$ and has
inverse $(\ell_i)_i\mapsto\sum_i\ell_i\otimes e_i$;
all supports are finite. The rows and columns of the
following diagram are induced by $u,v,j,q$:
$$
\begin{matrix}
 & & 0 & & 0 & & 0 & & \\
 & & \uparrow & & \uparrow & & \uparrow & & \\
 & & M''\otimes_R N & \to & M' \otimes_R N & \to & M \otimes_R N & \to & 0 \\
 & & \uparrow & & \uparrow & & \uparrow & & \\
0 & \to & (M'')^{(I)} & \to & (M')^{(I)} & \to & M^{(I)} & \to & 0 \\
 & & \uparrow & & \uparrow & & \uparrow & & \\
 & & M''\otimes_R K & \to & M' \otimes_R K & \to & M \otimes_R K & \to & 0 \\
 & & & & & & \uparrow & & \\
 & & & & & & 0 & &
\end{matrix}
$$
Each column ending in zero is right exact by tensoring
$K\to F\to N\to0$. The bottom and top rows are right
exact by tensoring $M''\to M'\to M\to0$. The middle row
is short exact coordinate by coordinate, including
surjectivity since a finitely supported tuple has finitely
many lifts. The bottom right vertical map is injective
because $M$ is flat and $K\to F$ is injective.
Every square commutes on its tensor generators.

Here is the full chase proving the remaining injection.
Let $x\in M''\otimes_R N$ map to zero in $M'\otimes_R N$.
Lift it to $a\in(M'')^{(I)}$. The image $u^{(I)}(a)$
has zero image in $M'\otimes_R N$, so it is the image
of some $b\in M'\otimes_R K$. The image of $b$ in
$M\otimes_R K$ maps to zero in $M^{(I)}$, because
$v^{(I)}u^{(I)}(a)=0$. Bottom right injectivity forces
that image to be zero. Bottom-row exactness supplies
$c\in M''\otimes_R K$ mapping to $b$.
Its image $d\in(M'')^{(I)}$ satisfies
$u^{(I)}(d)=u^{(I)}(a)$ by commutativity. Middle-row
injectivity gives $d=a$. As $d$ comes from
$M''\otimes_R K$, its image in $M''\otimes_R N$ is zero,
so $x=0$. This proves the asserted injection.
Finally the flips $L\otimes_R N\to N\otimes_R L$,
$\ell\otimes n\mapsto n\otimes\ell$, are mutual inverses
and commute with $u,v$, so the diagram conclusion is
exactly the sequence in the statement.
\end{proof}

\begin{lemma}
\label{editorial-source-lemma-flat-quotient-cyclic-tensor-criterion}
In the original short exact sequence
$0\to M''\xrightarrow{u}M'\xrightarrow{v}M\to0$,
assume $M'$ is flat. The following are equivalent:
\begin{enumerate}
\item $M$ is flat.
\item For every $R$-module $N$, the map
$N\otimes_R M''\to N\otimes_R M'$ is injective.
\item For every finitely generated ideal $I\subseteq R$,
the map $M''/IM''\to M'/IM'$ induced by $u$ is injective.
\end{enumerate}
\end{lemma}

\begin{proof}
The first implication is Lemma \ref{editorial-source-lemma-flat-tor-zero}.
The second follows by taking $N=R/I$ and using the
quotient-tensor maps $(r+I)\otimes m\mapsto rm+IM$
and $m+IM\mapsto(1+I)\otimes m$, on each original module.

Assume the third condition and retain any finite original
relation $\sum_{i=1}^n f_ix_i=0$ in $M$.
Choose lifts $\widetilde x_i\in M'$ with
$v(\widetilde x_i)=x_i$, and set $I=(f_1,\ldots,f_n)$.
The element $\sum_i f_i\widetilde x_i$ has zero image
under $v$, so it equals $u(w)$ for a unique $w\in M''$.
Its class in $M'/IM'$ is zero. Injectivity of the specified
quotient map gives $w\in IM''$.
Write $w=\sum_{\ell=1}^q c_\ell z_\ell$ with
$c_\ell\in I$, $z_\ell\in M''$, and express each original
$c_\ell=\sum_{i=1}^n f_i d_{i\ell}$.
Putting $w_i=\sum_{\ell=1}^q d_{i\ell}z_\ell$ gives the
exact equality $w=\sum_i f_iw_i$.
The corrected lifts
$\widehat x_i=\widetilde x_i-u(w_i)$ therefore satisfy
$$
\sum_i f_i\widehat x_i
=u(w)-u\left(\sum_i f_iw_i\right)=0,
\qquad v(\widehat x_i)=x_i.
$$
By flatness of the original $M'$ and
Lemma \ref{editorial-source-lemma-flat-eq}, there are $y_j\in M'$ and
$a_{ij}\in R$ with
$\widehat x_i=\sum_j a_{ij}y_j$ and
$\sum_i f_i a_{ij}=0$ for every $j$.
Applying $v$ gives $x_i=\sum_j a_{ij}v(y_j)$ with the
same coefficients and the same vanishing equations.
Thus every original relation in $M$ is trivial and
$M$ is flat by that criterion. The empty relation
requires no lifts and is already trivial.
\end{proof}


\begin{lemma}
\label{editorial-source-lemma-flat-ses}
Suppose that $0 \to M' \to M \to M'' \to 0$ is
a short exact sequence of $R$-modules.
If $M'$ and $M''$ are flat so is $M$.
If $M$ and $M''$ are flat so is $M'$.
\end{lemma}

\begin{proof}
We will use the criterion that a module $N$ is flat if for
every ideal $I \subset R$ the map $N \otimes_R I \to N$ is injective,
see Lemma \ref{editorial-source-lemma-flat}.
Consider an ideal $I \subset R$.
Consider the diagram
$$
\begin{matrix}
0 & \to & M' & \to & M & \to & M'' & \to & 0 \\
& & \uparrow & & \uparrow & & \uparrow & & \\
& & M'\otimes_R I & \to & M \otimes_R I & \to & M''\otimes_R I & \to & 0
\end{matrix}
$$
The vertical maps are $m\otimes a\mapsto am$ for the
original ideal $I$, so both squares commute on generators.
The upper row is the given short exact sequence and the
lower row is right exact. For the first assertion suppose
$M'$ and $M''$ are flat, and let $x\in M\otimes_R I$
map to zero in $M$. Its image in $M''\otimes_R I$ maps
to zero in $M''$ and hence is zero by flatness of $M''$.
Lower-row exactness lifts $x$ to $y\in M'\otimes_R I$.
The image of $y$ in $M'$ maps to zero in $M$, so it is
zero by upper-row injectivity. Flatness of $M'$ now gives
$y=0$, hence $x=0$. Lemma \ref{editorial-source-lemma-flat} proves that
$M$ is flat.

For the second assertion suppose $M$ and $M''$ are flat.
Lemma \ref{editorial-source-lemma-flat-tor-zero}, applied with quotient
$M''$ and tensor module $I$, makes the lower left map
injective. If $y\in M'\otimes_R I$ maps to zero in $M'$,
its image $x\in M\otimes_R I$ maps to zero in $M$.
Flatness of $M$ gives $x=0$, and lower left injectivity
gives $y=0$. The same ideal criterion proves flatness
of the original $M'$.
\end{proof}

\begin{lemma}
\label{editorial-source-lemma-easy-ff}
Let $R$ be a ring.
Let $M$ be an $R$-module.
The following are equivalent:
\begin{enumerate}
\item $M$ is faithfully flat, and
\item $M$ is flat and for all $R$-module homomorphisms $\alpha : N \to N'$
we have $\alpha = 0$ if and only if $\alpha \otimes \text{id}_M = 0$.
\end{enumerate}
\end{lemma}

\begin{proof}
Faithful flatness includes flatness. It also implies that
$L\otimes_R M=0$ forces $L=0$, by reflecting exactness
of $0\to L\to0$. For an original map $\alpha:N\to N'$,
retain $K=\Ker\alpha$. Flatness makes
$0\to K\otimes_R M\to N\otimes_R M\to N'\otimes_R M$
exact. If $\alpha\otimes\mathrm{id}_M=0$, the first
injection is also surjective. Right exactness of tensoring
$K\to N\to N/K\to0$ then gives $(N/K)\otimes_R M=0$.
The vanishing reflection just proved gives $N/K=0$, so
$K=N$ and $\alpha=0$. Conversely a zero map tensors
to a zero map by its formula on generators. This proves
the first implication with the original kernel retained.

Assume the second condition. Flatness already preserves
exactness. To prove reflection, take an original complex
$N_1\xrightarrow{d_1}N_2\xrightarrow{d_2}N_3$ whose
tensor complex is exact, and put $J=\operatorname{Im}d_1$.
For each $x\in\Ker d_2$, keep the original map
$\alpha:R\to N_2/J$, $r\mapsto rx+J$.
For every $m\in M$, the element $x\otimes m$ belongs to
$\Ker(d_2\otimes\mathrm{id}_M)$, hence to the image of
$d_1\otimes\mathrm{id}_M$. The quotient map $N_2\to N_2/J$
tensored with $M$ kills that image, so
$(x+J)\otimes m=0$. Therefore on every tensor generator
$$
(\alpha\otimes\mathrm{id}_M)(r\otimes m)
=(rx+J)\otimes m=(x+J)\otimes rm=0.
$$
The assumed reflection of zero morphisms gives $\alpha=0$,
and evaluation at $1$ gives $x\in J$. Thus
$\Ker d_2\subseteq\operatorname{Im}d_1$; the reverse
inclusion holds because the original sequence is a complex.
It is exact, proving faithful flatness.
\end{proof}

\begin{lemma}
\label{editorial-source-lemma-ff}
\begin{slogan}
A flat module is faithfully flat if and only if it has nonzero fibers.
\end{slogan}
Let $M$ be a flat $R$-module.
The following are equivalent:
\begin{enumerate}
\item $M$ is faithfully flat,
\item for every nonzero $R$-module $N$, the tensor product $M \otimes_R N$
is nonzero,
\item for all $\mathfrak p \in \Spec(R)$
the tensor product $M \otimes_R \kappa(\mathfrak p)$ is nonzero, and
\item for all maximal ideals $\mathfrak m$ of $R$
the tensor product $M \otimes_R \kappa(\mathfrak m) = M/{\mathfrak m}M$
is nonzero.
\end{enumerate}
\end{lemma}

\begin{proof}
Assume $M$ faithfully flat and $N \not = 0$. By Lemma \ref{editorial-source-lemma-easy-ff}
the nonzero map $1 : N \to N$ induces a nonzero map
$M \otimes_R N \to M \otimes_R N$, so $M \otimes_R N \not = 0$.
Thus (1) implies (2). The implications (2) $\Rightarrow$ (3) $\Rightarrow$ (4)
are immediate.

\medskip\noindent
Assume (4), and retain a complex
$N_1\xrightarrow{d_1}N_2\xrightarrow{d_2}N_3$ whose
tensor complex is exact. Write
$K=\Ker d_2$, $J=\operatorname{Im}d_1\subseteq K$,
and keep its original cohomology $H=K/J$.
Flatness embeds $J\otimes_R M$ and $K\otimes_R M$ in
$N_2\otimes_R M$. Tensoring the surjection $N_1\to J$
identifies the first image with
$\operatorname{Im}(d_1\otimes\mathrm{id}_M)$; tensoring
$0\to K\to N_2\to N_3$ identifies the second image with
$\operatorname{Ker}(d_2\otimes\mathrm{id}_M)$.
The tensor complex is exact, so these submodules coincide.
Right exactness applied to $J\to K\to H\to0$ gives
$H\otimes_R M=0$ with these exact original maps.

Take any $x\in H$ and set $I=\{f\in R:fx=0\}$.
The map $R/I\to H$, $f+I\mapsto fx$, is well-defined
and injective by the definition of $I$.
Flatness makes $(R/I)\otimes_R M\to H\otimes_R M=0$
injective. The quotient-tensor isomorphism
$(f+I)\otimes m\mapsto fm+IM$ with inverse
$m+IM\mapsto(1+I)\otimes m$ therefore gives $M/IM=0$.
If $I$ were proper, choose a maximal ideal
$\mathfrak m\supseteq I$. The quotient map
$M/IM\to M/\mathfrak mM$, $m+IM\mapsto m+\mathfrak mM$,
is well-defined and surjective. Its source is zero, forcing
$M/\mathfrak mM=0$, contrary to (4). Thus $I=R$,
$x=1x=0$, and $H=0$. This proves exactness reflection.
For the zero ring all modules are zero and all prime and
maximal-ideal conditions are vacuous, so the same equivalence
also covers that case.
\end{proof}

\begin{lemma}
\label{editorial-source-lemma-ff-rings}
Let $R \to S$ be a flat ring map.
The following are equivalent:
\begin{enumerate}
\item $R \to S$ is faithfully flat,
\item the induced map on $\Spec$ is surjective, and
\item any closed point $x \in \Spec(R)$ is
in the image of the map $\Spec(S) \to \Spec(R)$.
\end{enumerate}
\end{lemma}

\begin{proof}
For the original prime $\mathfrak p\subset R$, the fibre
ring $S\otimes_R\kappa(\mathfrak p)$ is nonzero exactly
when there is a prime of $S$ contracting to $\mathfrak p$,
by Remark \ref{algebra-remark-fundamental-diagram} and
Lemma \ref{editorial-source-lemma-in-image}. The prime correspondence is
through the quotient by $\mathfrak p$ followed by inversion
of the original $R\setminus\mathfrak p$; its primes
contract to primes of $S$ with precisely that contraction.
Thus condition (3) of Lemma \ref{editorial-source-lemma-ff} says exactly
that every prime is in the spectrum image. Its condition
(4) says exactly that every maximal ideal is in the image,
and maximal ideals are the closed points of $\Spec(R)$.
Since $S$ is flat over $R$, that lemma proves all three
equivalences asserted here.
\end{proof}

\begin{lemma}
\label{editorial-source-lemma-local-flat-ff}
A flat local ring homomorphism of local rings is faithfully flat.
\end{lemma}

\begin{proof}
Write $\varphi:(R,\mathfrak m)\to(S,\mathfrak n)$ for
the original local homomorphism. Its defining local
condition is $\varphi^{-1}(\mathfrak n)=\mathfrak m$.
Thus the unique closed point of $\Spec(R)$ is the image
of the prime $\mathfrak n$ of $S$. Flatness and
Lemma \ref{editorial-source-lemma-ff-rings} give faithful flatness.
\end{proof}

\noindent
Flatness meshes well with localization.

\begin{lemma}
\label{editorial-source-lemma-flat-localization}
Let $R$ be a ring. Let $S \subset R$ be a multiplicative subset.
\begin{enumerate}
\item The localization $S^{-1}R$ is a flat $R$-algebra.
\item If $M$ is an $S^{-1}R$-module, then $M$ is a flat $R$-module
if and only if $M$ is a flat $S^{-1}R$-module.
\item Suppose $M$ is an $R$-module. Then
$M$ is a flat $R$-module if and only if $M_{\mathfrak p}$ is a flat
$R_{\mathfrak p}$-module for all primes $\mathfrak p$ of $R$.
\item Suppose $M$ is an $R$-module. Then $M$ is a flat $R$-module if
and only if $M_{\mathfrak m}$ is a flat
$R_{\mathfrak m}$-module for all maximal ideals $\mathfrak m$ of $R$.
\item Suppose $R \to A$ is a ring map, $M$ is an $A$-module,
and $g_1, \ldots, g_m \in A$ are elements generating the unit
ideal of $A$. Then $M$ is flat over $R$ if and only if each localization
$M_{g_i}$ is flat over $R$.
\item Suppose $R \to A$ is a ring map, and $M$ is an $A$-module.
Then $M$ is a flat $R$-module if and only if the localization
$M_{\mathfrak q}$ is a flat $R_{\mathfrak p}$-module
(with $\mathfrak p$ the prime of $R$ lying under $\mathfrak q$)
for all primes $\mathfrak q$ of $A$.
\item Suppose $R \to A$ is a ring map, and $M$ is an $A$-module.
Then $M$ is a flat $R$-module if and only if the localization
$M_{\mathfrak m}$ is a flat $R_{\mathfrak p}$-module
(with $\mathfrak p$ the inverse image of $\mathfrak m$ under $R\to A$)
for all maximal ideals $\mathfrak m$ of $A$.
\end{enumerate}
\end{lemma}

\begin{proof}
Retain every original coefficient map; in (5)--(7) write it
as $\varphi:R\to A$. All contractions below are inverse
images under this map, which need not be injective.

We first recall the exact fraction facts used in the proof.
For a ring $B$, a multiplicative set $W\subseteq B$ and a
$B$-module $L$, the equality $\ell/w=\ell'/w'$ means
that $v(w'\ell-w\ell')=0$ for some original $v\in W$.
Thus $\ell/w=0$ means that some $v\in W$ kills $\ell$.
If $L'\to L$ is injective, this equality for an element
of $L'$ holds in $L'$ whenever it holds in $L$, so
localization preserves injections. For a map $u:L\to N$,
an element $\ell/w$ in the localized kernel has
$v u(\ell)=0$ for some $v\in W$. Then $v\ell\in\Ker u$
and $\ell/w=(v\ell)/(vw)$, proving that localization
preserves kernels. Its image consists exactly of fractions
$u(\ell)/w$, and a surjection remains surjective by lifting
the numerator. These facts prove exactness with the
original fraction witnesses retained.

For later zero detection, if $0\ne\ell\in L$ its
annihilator is a proper ideal of $B$ and is contained in
a maximal ideal $\mathfrak m$. Then $\ell/1\ne0$ in
$L_{\mathfrak m}$, since an element outside $\mathfrak m$
killing $\ell$ would belong to that annihilator. Hence
$L=0$ if all its maximal localizations vanish, as in
Lemma \ref{algebra-lemma-characterize-zero-local}.

For an original $R$-module $N$, the $A$-action on
$N\otimes_R M$ is $a(n\otimes m)=n\otimes am$.
For every multiplicative $W\subseteq A$, there are
mutually inverse maps
$$
\begin{split}
W^{-1}(N\otimes_R M)&\longrightarrow N\otimes_R W^{-1}M,\\
\frac{\sum_{\ell=1}^d n_\ell\otimes m_\ell}{w}
&\longmapsto\sum_{\ell=1}^d n_\ell\otimes\frac{m_\ell}{w},
\end{split}
$$
and $n\otimes(m/w)\mapsto(n\otimes m)/w$.
For the forward map start with $n\otimes m\mapsto
n\otimes(m/1)$, which is $A$-linear and $R$-balanced;
every element of $W$ acts invertibly on its target, so
it extends to the localization with the stated formula.
For the inverse, if $v(w'm-wm')=0$, then
$v(w'(n\otimes m)-w(n\otimes m'))=0$ by the original
$A$-action. This proves independence of the fraction.
The $R$-balancing relation is
$rn\otimes m=n\otimes\varphi(r)m$, and is preserved by
the formula. Additivity follows by common denominators.
Both composites fix the displayed generators and their
finite sums. The maps are natural in $N$ and $M$, since
they apply their homomorphisms only to the numerators.
Consequently, if $M$ is flat over $R$, then $W^{-1}M$
is flat over $R$: localize the injection obtained by
tensoring any injection of $R$-modules with $M$ and use
these natural isomorphisms.

For a prime $\mathfrak q\subset A$, put
$\mathfrak p=\varphi^{-1}(\mathfrak q)$. The map and action are
$$
R_{\mathfrak p}\longrightarrow A_{\mathfrak q},\quad
\frac rs\longmapsto\frac{\varphi(r)}{\varphi(s)},
\qquad
\frac rs\cdot\frac ma=\frac{\varphi(r)m}{\varphi(s)a}.
$$
All denominators are valid because $s\notin\mathfrak p$
is equivalent to $\varphi(s)\notin\mathfrak q$.
The exact two-sided comparison for every original $N$ is
$$
(N\otimes_R M)_{\mathfrak q}
\ \cong\ N_{\mathfrak p}\otimes_{R_{\mathfrak p}}M_{\mathfrak q},
\qquad
\frac{n\otimes m}{a}\longmapsto\frac n1\otimes\frac ma,
$$
with inverse
$$
\frac ns\otimes\frac ma\longmapsto
\frac{n\otimes m}{\varphi(s)a}.
$$
If $n/s=n'/s'$, choose $u\notin\mathfrak p$ with
$u(s'n-sn')=0$. The two proposed inverse images have
common denominator $\varphi(s)\varphi(s')a$ and numerator
$(s'n-sn')\otimes m$. Multiplication by
$\varphi(u)\notin\mathfrak q$ kills it, by $R$-balancing.
If $m/a=m'/a'$, choose $b\notin\mathfrak q$ with
$b(a'm-am')=0$. Using common denominator $\varphi(s)aa'$
gives numerator $n\otimes(a'm-am')$, killed by $b$.
For $r/t\in R_{\mathfrak p}$ the two balancing images are
$$
\frac{rn\otimes m}{\varphi(t)\varphi(s)a}
=\frac{n\otimes\varphi(r)m}{\varphi(s)\varphi(t)a},
$$
so the inverse is well-defined and additive.
The forward map is the localization of
$n\otimes m\mapsto(n/1)\otimes(m/1)$; elements outside
$\mathfrak q$ act invertibly on its target.
One composite fixes $(n\otimes m)/a$, while the other
sends $(n/s)\otimes(m/a)$ to
$(n/1)\otimes(m/(\varphi(s)a))$, which equals the
original tensor by balancing with $1/s$.
The formulas prove naturality for the original maps.

We also retain the exact ideal descent used in the source.
For any ideal $I\subseteq R_{\mathfrak p}$, put
$J=\{r\in R:r/1\in I\}$. Then $J_{\mathfrak p}=I$:
if $j\in J$, its fraction $j/s=(1/s)(j/1)$ is in $I$;
if $a/s\in I$, then $a/1=(s/1)(a/s)\in I$, so $a\in J$.
The localized map $J_{\mathfrak p}\to R_{\mathfrak p}$ is
injective by localization of $J\hookrightarrow R$.
For specified finite generators $a_\ell/s_\ell$ of $I$,
the smaller specified ideal $J_0=(a_1,\ldots,a_n)$ also
has $(J_0)_{\mathfrak p}=I$, because each $s_\ell/1$ is
a unit and both identities above apply to that generator.
This does not assert finite generation of the contraction $J$.

We now prove (1). Keep the original multiplicative set
$S\subseteq R$ and put $B=S^{-1}R$.
For every $N$ the maps
$$
N\otimes_R B\longrightarrow S^{-1}N,\quad
n\otimes(r/s)\longmapsto rn/s,
\qquad
n/s\longmapsto n\otimes(1/s)
$$
are inverse. To check the inverse on fraction equality,
if $u(s'n-sn')=0$ then its difference of images is
$$
(s'n-sn')\otimes\frac1{ss'}
=u(s'n-sn')\otimes\frac1{uss'}=0.
$$
The forward map respects equal fractions for the same
reason and respects the original $R$-balancing relation.
The two composites fix the generators, since
$rn\otimes(1/s)=n\otimes(r/s)$.
Localization preserves injections, so these natural maps
and Lemma \ref{editorial-source-lemma-flat} prove that $B$ is flat over $R$.

For (2), if $M$ is flat over $B$, then (1) and
Lemma \ref{editorial-source-lemma-composition-flat} imply flatness over $R$.
The receiving tensor comparison can also be written as
$$
(S^{-1}N)\otimes_B M\longrightarrow N\otimes_R M,
\qquad (n/s)\otimes m\longmapsto n\otimes(1/s)m,
$$
with inverse $n\otimes m\mapsto(n/1)\otimes m$.
For example if $u(s'n-sn')=0$, the difference between
the first images is
$(s'n-sn')\otimes(1/(ss'))m
=u(s'n-sn')\otimes(1/(uss'))m=0$.
Balancing by $r/t\in B$ agrees on both sides after
moving $r$ and the inverse action of $t$ to the $M$ factor.
Both composites are identities by these same relations.

Conversely let $M$ be a $B$-module flat over $R$.
For any $B$-module $P$, the quotient map
$P\otimes_R M\to P\otimes_B M$ has inverse
$p\otimes_B m\mapsto p\otimes_R m$.
To verify the additional balancing needed for this inverse,
for each original $s\in S$ use
$$
(1/s)p\otimes_R m
=(1/s)p\otimes_R s((1/s)m)
=p\otimes_R(1/s)m.
$$
Together with $R$-balancing this gives balancing by every
$r/s\in B$. The composites fix all generators.
An injection of $B$-modules is an injection of underlying
$R$-modules, so these natural maps and flatness over $R$
give the injection criterion over $B$. This proves (2).

Next prove (6) and (7), retaining the original ideal maps.
Suppose $M_{\mathfrak m}$ is flat over $R_{\mathfrak p}$
for each maximal $\mathfrak m\subset A$, where
$\mathfrak p=\varphi^{-1}(\mathfrak m)$.
For an ideal $I\subseteq R$, consider the original
$A$-linear map
$$
\eta:I\otimes_R M\longrightarrow M,
\qquad a\otimes m\longmapsto\varphi(a)m,
\qquad K=\Ker\eta.
$$
At $\mathfrak m$ the comparison already proved gives
$$
(I\otimes_R M)_{\mathfrak m}
\cong I_{\mathfrak p}\otimes_{R_{\mathfrak p}}M_{\mathfrak m}.
$$
It identifies $\eta_{\mathfrak m}$ with the multiplication
map $(a/s)\otimes(m/t)\mapsto\varphi(a)m/(\varphi(s)t)$,
so that map is injective by the asserted local flatness
and the inclusion $I_{\mathfrak p}\hookrightarrow R_{\mathfrak p}$.
Exactness of localization gives $K_{\mathfrak m}=0$.
Zero detection at all maximal ideals of $A$ gives $K=0$,
and Lemma \ref{editorial-source-lemma-flat} proves flatness over $R$.

Conversely suppose $M$ is flat over $R$ and retain any
prime $\mathfrak q\subset A$ with
$\mathfrak p=\varphi^{-1}(\mathfrak q)$.
For an ideal $I\subseteq R_{\mathfrak p}$ take its original
contraction $J\subseteq R$ above, so $J_{\mathfrak p}=I$.
Flatness makes $J\otimes_R M\to M$ injective. Localizing
at $\mathfrak q$ preserves its injectivity, and the exact
comparison identifies it with
$I\otimes_{R_{\mathfrak p}}M_{\mathfrak q}\to M_{\mathfrak q}$.
The ideal criterion proves flatness over $R_{\mathfrak p}$.
This proves necessity at every prime and sufficiency at
all maximal ideals. Since the former condition includes
the latter, (6) and (7) both follow. A maximal ideal of
$A$ may contract to a nonmaximal prime of $R$; no step
replaces the original contracted prime by a maximal one.

Taking $A=R$ and $\varphi=\mathrm{id}_R$ in (6) and (7)
gives (3) and (4), respectively, with the same module $M$
and the same prime and maximal localizations.

For (5) we prove the stronger arbitrary-family statement:
if $(g_\lambda)_{\lambda\in\Lambda}$ is any family in $A$
such that for every maximal $\mathfrak m$ with
$M_{\mathfrak m}\ne0$ some $g_\lambda\notin\mathfrak m$,
then $M$ is flat over $R$ if and only if every
$M_{g_\lambda}$ is flat over $R$.

First, for each original $g\in A$, multiplication by $g$
gives a directed system of the original $R$-module $M$,
indexed by $n\in\mathbf Z_{\geq0}$. Its colimit is
$$
\colim(M\xrightarrow{g}M\xrightarrow{g}\cdots)
\longrightarrow M_g,\qquad [m\text{ at }n]\longmapsto m/g^n.
$$
Compatibility is $gm/g^{n+1}=m/g^n$.
Every fraction occurs. A fraction $m/g^n$ is zero exactly
when some $g^k m=0$, which is exactly eventual vanishing
of its original representative. This proves injectivity
as well as surjectivity; the inverse sends a fraction to
that class, independently of its representative by the
same equality test. Lemma \ref{editorial-source-lemma-colimit-flat} then
proves that $M_g$ is flat whenever $M$ is flat.

For the converse retain any injection of $R$-modules
$N'\hookrightarrow N$ and its original $A$-module kernel
$K=\Ker(N'\otimes_R M\to N\otimes_R M)$.
If every $M_{g_\lambda}$ is flat over $R$, localization
and the first tensor comparison identify the localized
map with
$N'\otimes_R M_{g_\lambda}\to N\otimes_R M_{g_\lambda}$,
which is injective. Thus $K_{g_\lambda}=0$ for each $\lambda$.
At a maximal ideal $\mathfrak m$ with $M_{\mathfrak m}=0$,
the same comparison gives
$(N'\otimes_R M)_{\mathfrak m}=N'\otimes_R M_{\mathfrak m}=0$,
so $K_{\mathfrak m}=0$. At any other maximal ideal choose
$g_\lambda\notin\mathfrak m$ by the hypothesis.
For every $k\in K$, the equality $K_{g_\lambda}=0$ gives
$g_\lambda^e k=0$ for some $e\geq0$, which may depend
on $k$. This power is outside $\mathfrak m$, so $k/1=0$
in $K_{\mathfrak m}$. Thus $K_{\mathfrak m}=0$ at every
maximal ideal, and zero detection gives $K=0$.
The injection criterion proves flatness over $R$.

The stated finite unit-ideal hypothesis implies this cover:
if every $g_i$ belonged to a maximal ideal, so would their
generated ideal, contrary to $(g_1,\ldots,g_m)=A$.
This proves (5), including both directions.

The stronger covering hypothesis can equivalently be
stated on all primes $\mathfrak q$ with $M_{\mathfrak q}\ne0$.
Indeed choose $x\in M$ with $x/1\ne0$ at such a prime.
Then $\operatorname{Ann}_A(x)\subseteq\mathfrak q$ by
the exact fraction-zero test. At every maximal
$\mathfrak m\supseteq\mathfrak q$, the same inclusion
shows $x/1\ne0$ in $M_{\mathfrak m}$. A cover element
outside $\mathfrak m$ is also outside $\mathfrak q$.
The reverse implication specializes the prime condition
to maximal ideals. No finite generation of $M$ or of the
family is required.

This really weakens the unit-ideal requirement. For a
nonzero ring $R$, take $A=R\times R$ with the diagonal
map from $R$, and an $R$-module $M$ viewed as an $A$-module
through the first projection. The element $g=(1,0)$ acts
as the identity on $M$, so $M_g=M$ by the maps
$m/g^n\mapsto m$ and $m\mapsto m/1$.
The element $1-g=(0,1)$ acts as zero. If a prime contains
$g$, it cannot contain $1-g$, which becomes an invertible
annihilator of the localized module; hence that localization
is zero. Thus $D(g)$ covers every nonzero stalk of $M$,
although $gA=R\times0$ is a proper ideal of $A$.

Finally all assertions include zero cases. If $0$ lies
in a multiplicative set, its localized ring and modules
are zero by the fraction rule; a unital module over that
ring is zero. In particular $g=0$ or a nilpotent $g$ gives
$M_g=0$ and an empty principal open. If $A=0$, then $M=0$
and its prime and maximal conditions are empty. If the
cover family is empty, its hypothesis says all maximal
stalks of $M$ are zero, so $M=0$ by zero detection and the
equivalence is valid. The zero module is flat by the
definition in every case.
\end{proof}

\begin{lemma}
\label{editorial-source-lemma-flat-going-down}
Let $R \to S$ be flat. Let $\mathfrak p \subset \mathfrak p'$
be primes of $R$. Let $\mathfrak q' \subset S$ be a prime of $S$
mapping to $\mathfrak p'$. Then there exists a prime
$\mathfrak q \subset \mathfrak q'$ mapping to $\mathfrak p$.
\end{lemma}

\begin{proof}
Let $\varphi:R\to S$ be the original map and retain the
stated primes. The localization map
$$
\varphi':R_{\mathfrak p'}\longrightarrow S_{\mathfrak q'},
\qquad r/t\longmapsto\varphi(r)/\varphi(t)
$$
is defined because $t\notin\mathfrak p'$ implies
$\varphi(t)\notin\mathfrak q'$. Its maximal-ideal
contraction is $\mathfrak p'R_{\mathfrak p'}$, so it is
local, and Lemma \ref{editorial-source-lemma-flat-localization} makes it
flat. It is faithfully flat by Lemma \ref{editorial-source-lemma-local-flat-ff}.
The ideal $\mathfrak pR_{\mathfrak p'}$ is prime and has
contraction $\mathfrak p$ in $R$ by prime localization.
Lemma \ref{editorial-source-lemma-ff-rings} supplies a prime
$\widetilde{\mathfrak q}\subset S_{\mathfrak q'}$ whose
contraction under $\varphi'$ is that ideal. Define
$$
\mathfrak q=\{s\in S:s/1\in\widetilde{\mathfrak q}\}.
$$
It is prime as an inverse image of a prime. An element of
$S\setminus\mathfrak q'$ becomes a unit in $S_{\mathfrak q'}$
and cannot belong to this inverse image, so
$\mathfrak q\subseteq\mathfrak q'$. Finally
$$
r\in\varphi^{-1}(\mathfrak q)
\Longleftrightarrow
r/1\in\mathfrak pR_{\mathfrak p'}
\Longleftrightarrow r\in\mathfrak p.
$$
This proves the exact required contraction, without
assuming the original flat map is injective.
\end{proof}

\noindent
The property of $R \to S$ described in the lemma is called the
``going down property''. See Definition \ref{algebra-definition-going-up-down}.

\begin{lemma}
\label{editorial-source-lemma-colimit-faithfully-flat}
Let $R$ be a ring. Let $\{S_i, \varphi_{ii'}\}$ be a directed system of
faithfully flat $R$-algebras. Then $S = \colim_i S_i$ is a faithfully flat
$R$-algebra.
\end{lemma}

\begin{proof}
The underlying $R$-module colimit is flat by
Lemma \ref{editorial-source-lemma-colimit-flat}. For every maximal ideal
$\mathfrak m\subset R$, faithful flatness at each stage
makes $S_i/\mathfrak mS_i$ nonzero, by Lemma \ref{editorial-source-lemma-ff}.
The canonical ring comparison is
$$
\colim_i(S_i/\mathfrak mS_i)\longrightarrow S/\mathfrak mS,
\qquad [s_i+\mathfrak mS_i]\longmapsto[s_i]+\mathfrak mS.
$$
It is surjective by representatives. If $[s_i]\in\mathfrak mS$,
write its original expression as a finite sum
$[s_i]=\sum_{\ell=1}^q a_\ell[t_{j_\ell}]$ with
$a_\ell\in\mathfrak m$. Move all terms to a common stage
and then to a further stage where this equality holds.
There the image of $s_i$ belongs to the actual ideal
$\mathfrak mS_k$, so its quotient class is eventually zero.
This proves injectivity and gives the inverse by choosing
any original representative; equality of two choices is
handled by the same kernel argument.

If this colimit were the zero ring, the class of $1$ at
one stage would become zero at some later stage. Every
transition is a unital $R$-algebra map, so the image is
the original $1$ of that later quotient, contradicting
its nonzero ring structure. Hence $S/\mathfrak mS\ne0$
for every maximal ideal, and Lemma \ref{editorial-source-lemma-ff} gives
faithful flatness. This is the specific reason algebra
colimits preserve faithful flatness although arbitrary
module colimits need not: the same element $1$ survives
in every receiving residue quotient. For $R=0$ all
unital $R$-algebras and their modules are zero, and the
faithfully-flat conclusion holds by the definition.
\end{proof}








\subsection{Editorial treatment of Algebra lines 9328--9344}

\hypertarget{algebra-context-098}{}

\noindent\href{https://github.com/stacks/stacks-project/blob/a04446e57ec1fbc252a871afcec7752fb2807b14/algebra.tex\#L9328}{Original source, lines 9328--9344}. Complete original and reviewed TeX passages accompany this note. Every intervention is separately recorded; the following treatment is editorial.


\begin{lemma}
\label{editorial-source-lemma-support-zero}
\begin{slogan}
A module over a ring has empty support if and only if it is the trivial module.
\end{slogan}
Let $R$ be a ring. Let $M$ be an $R$-module. Then
$$
M = (0) \Leftrightarrow \text{Supp}(M) = \emptyset.
$$
\end{lemma}

\begin{proof}
If $M=0$, all its localizations are zero, so its support
is empty. Conversely, for any original nonzero $m\in M$,
the ideal $\{r\in R:rm=0\}$ is proper and lies in a
maximal ideal $\mathfrak p$. The fraction $m/1$ is nonzero
in $M_{\mathfrak p}$: if it were zero, a witness
$s\notin\mathfrak p$ would satisfy $sm=0$, contradicting
membership of every such annihilator in $\mathfrak p$.
Thus the support contains a maximal ideal, as also proved
in Lemma \ref{algebra-lemma-characterize-zero-local}.
\end{proof}


\subsection{Editorial treatment of Algebra lines 9363--9382}

\hypertarget{algebra-context-099}{}

\noindent\href{https://github.com/stacks/stacks-project/blob/a04446e57ec1fbc252a871afcec7752fb2807b14/algebra.tex\#L9363}{Original source, lines 9363--9382}. Complete original and reviewed TeX passages accompany this note. Every intervention is separately recorded; the following treatment is editorial.


\begin{lemma}
\label{editorial-source-lemma-annihilator-flat-base-change}
Let $R \to S$ be a flat ring map. Let $M$ be an $R$-module and
$m \in M$. Then $\text{Ann}_R(m) S = \text{Ann}_S(m \otimes 1)$.
If $M$ is a finite $R$-module, then
$\text{Ann}_R(M)S=\text{Ann}_S(M\otimes_R S)$.
The latter conclusion also holds if finitely many elements
$m_1,\ldots,m_n\in M$ have
$\operatorname{Ann}_R(M)=\bigcap_{i=1}^n\operatorname{Ann}_R(m_i)$,
without requiring that they generate $M$.
\end{lemma}

\begin{proof}
Write the original ring map as $\varphi:R\to S$ and retain
$I=\operatorname{Ann}_R(m)$. The exact sequence
$$
0\longrightarrow I\longrightarrow R\longrightarrow M,
\qquad r\longmapsto rm,
$$
remains exact after tensoring with the flat $R$-module $S$.
Under $R\otimes_RS\to S$, $r\otimes s\mapsto\varphi(r)s$,
the last map becomes $s\mapsto m\otimes s=(m\otimes1)s$.
The image of $I\otimes_RS\to S$, $a\otimes s\mapsto\varphi(a)s$,
is exactly the extended ideal $IS$. Its equality with the
kernel therefore proves
$$
\operatorname{Ann}_R(m)S=\operatorname{Ann}_S(m\otimes1).
$$

For the original finite generating family $m_1,\ldots,m_n$,
retain $I_i=\operatorname{Ann}_R(m_i)$. An element kills $M$
exactly when it kills every generator. The elements
$m_i\otimes1$ generate $M\otimes_RS$, so the full comparison is
$$
\begin{split}
\operatorname{Ann}_S(M\otimes_RS)
&=\bigcap_{i=1}^n\operatorname{Ann}_S(m_i\otimes1)\\
&=\bigcap_{i=1}^n(I_iS)
=\left(\bigcap_{i=1}^n I_i\right)S
=\operatorname{Ann}_R(M)S.
\end{split}
$$
The finite-intersection equality is
Lemma \ref{editorial-source-lemma-flat-intersect-ideals}.

For the stronger assertion, suppose only that the original
finite list detects the annihilator:
$$
J=\operatorname{Ann}_R(M)=\bigcap_{i=1}^n I_i.
$$
These elements need not generate $M$. The sequence
$$
0\longrightarrow J\longrightarrow R\longrightarrow M^n,
\qquad r\longmapsto(rm_1,\ldots,rm_n),
$$
is exact. Tensoring and using the coordinate maps
$M^n\otimes_RS\simeq(M\otimes_RS)^n$ gives
$JS=\bigcap_i\operatorname{Ann}_S(m_i\otimes1)$.
Consequently every element annihilating the whole tensor
module belongs to $JS$. Conversely each finite sum
$\sum_j\varphi(a_j)s_j$ with $a_j\in J$ kills every tensor:
$$
\left(\sum_j\varphi(a_j)s_j\right)(m\otimes s)
=\sum_j(a_jm)\otimes s_js=0.
$$
This proves equality for the weaker hypothesis. For $n=0$
the empty intersection is $R$, forcing $M=0$; both module
annihilators are then the whole rings, as required.
For example, any nonempty free module, even with an infinite
basis, has its annihilator detected by one original basis
vector: $re=0$ implies $r=0$ by that coordinate.

\medskip\noindent
For an arbitrary module the same element calculation gives
the exact formula
$$
\operatorname{Ann}_S(M\otimes_RS)
=\bigcap_{m\in M}\bigl(\operatorname{Ann}_R(m)S\bigr),
$$
since all the elements $m\otimes1$ generate the tensor module.
Extension of ideals need not commute with this infinite
intersection, even for a faithfully flat ring map.

Retain the explicit example
$$
R=\mathbf Z,\qquad
M=\bigoplus_{n\geq1}\mathbf Z/2^n\mathbf Z,\qquad
S=\mathbf Z\times\mathbf Q,
\qquad \varphi(a)=(a,a).
$$
The rational numbers are the localization at the nonzero
integers. For every module $N$ the tensor-localization map is
$N\otimes_{\mathbf Z}\mathbf Q\to
(\mathbf Z\setminus\{0\})^{-1}N$,
$n\otimes a/b\mapsto an/b$, with inverse $n/b\mapsto n\otimes1/b$.
Localization preserves injections: a fraction mapping to zero
has a nonzero integer witness killing its numerator, and
injectivity gives the same relation in the original module.
Thus $\mathbf Q$ is flat, as also established in
Example \ref{editorial-source-example-flat-infinite-intersections}.
The exact tensor comparison
$$
N\otimes_{\mathbf Z}(\mathbf Z\times\mathbf Q)
\longrightarrow N\oplus(N\otimes_{\mathbf Z}\mathbf Q),
\qquad n\otimes(a,q)\longmapsto(an,n\otimes q)
$$
has inverse
$$
\left(n,\sum_j n_j\otimes q_j\right)
\longmapsto n\otimes(1,0)+\sum_j n_j\otimes(0,q_j).
$$
Both are balanced and their composites are identities on
the displayed generators. This proves flatness of $S$ and
faithfulness: the tensor of a homomorphism $h$ has first
component $h$, so it is zero only if $h=0$.

The original annihilator is
$$
\operatorname{Ann}_{\mathbf Z}(M)
=\bigcap_{n\geq1}2^n\mathbf Z=0,
$$
because no nonzero integer is divisible by every $2^n$.
Each original $m\in M$ has finite support, hence is killed
by some $2^N$. Therefore
$m\otimes q=(2^Nm)\otimes(q/2^N)=0$ in
$M\otimes_{\mathbf Z}\mathbf Q$. The preceding comparison gives
$M\otimes_{\mathbf Z}S\simeq M\oplus0$, with the action
$(a,q)(m,0)=(am,0)$. It follows exactly that
$$
\operatorname{Ann}_{\mathbf Z}(M)S=0,
\qquad
\operatorname{Ann}_S(M\otimes_{\mathbf Z}S)=0\times\mathbf Q\ne0.
$$
Every finite family in $M$ is killed by a common nonzero
power of $2$, so no such family detects its zero annihilator.
This same module has support $\{(2)\}$: its first summand
has nonzero localization at $(2)$, while at every other
prime $2$ is a unit and every element is killed by a power
of $2$. Thus its support is closed but is strictly smaller
than $V(\operatorname{Ann}_{\mathbf Z}(M))=\Spec(\mathbf Z)$.

The counterexample even has a finitely presented faithfully
flat algebra version. Put $S'=\mathbf Z\times\mathbf Z[1/2]$.
The identical localization and direct-sum arguments prove
faithful flatness and give annihilator $0\times\mathbf Z[1/2]$.
More explicitly,
$$
\mathbf Z[T]/(2T^2-T)\longrightarrow S',
\qquad f(T)\longmapsto(f(0),f(1/2)),
$$
is an isomorphism. The ideals $(T)$ and $(2T-1)$ are
comaximal since $2T-(2T-1)=1$; hence their intersection
is their product $(2T^2-T)$. The Chinese remainder map
is surjective with this kernel. Its factors are $\mathbf Z$
and $\mathbf Z[1/2]$: in the second quotient $2T=1$, and
the inverse to evaluation at $1/2$ sends $a/2^j$ to $aT^j$.
This proves the asserted presentation with the original
integer factors retained.

Finally, faithful flatness does always imply the contraction
identity
$$
\varphi^{-1}\bigl(\operatorname{Ann}_S(M\otimes_RS)\bigr)
=\operatorname{Ann}_R(M).
$$
Multiplication by $r$ on $M$ becomes multiplication by
$\varphi(r)$ on its tensor module, and a faithfully flat
tensor functor detects whether this homomorphism is zero.
Thus the counterexample fails extension, while contraction
still recovers the original annihilator. For $M=0$ or the
zero ring, all the annihilator statements retain these meanings.
\end{proof}


\subsection{Editorial treatment of Algebra lines 9384--9410}

\hypertarget{algebra-context-100}{}

\noindent\href{https://github.com/stacks/stacks-project/blob/a04446e57ec1fbc252a871afcec7752fb2807b14/algebra.tex\#L9384}{Original source, lines 9384--9410}. Complete original and reviewed TeX passages accompany this note. Every intervention is separately recorded; the following treatment is editorial.


\begin{lemma}
\label{editorial-source-lemma-support-closed}
Let $R$ be a ring and let $M$ be an $R$-module.
If $M$ is finite, then $\text{Supp}(M)$ is closed.
More precisely, if $I = \text{Ann}(M)$ is the annihilator of $M$, then
$V(I) = \text{Supp}(M)$.
\end{lemma}

\begin{proof}
We will show that $V(I) = \text{Supp}(M)$.

\medskip\noindent
Suppose $\mathfrak p \in \text{Supp}(M)$. Then $M_{\mathfrak p} \not = 0$.
Choose an element $m \in M$ whose image in $M_\mathfrak p$ is nonzero.
Then the annihilator of $m$ is contained in $\mathfrak p$ by construction
of the localization $M_\mathfrak p$. Hence
a fortiori $I = \text{Ann}(M)$ must be contained in $\mathfrak p$.

\medskip\noindent
If $M=0$, then $I=R$ and both sets are empty.
Otherwise retain generators $x_1,\ldots,x_r$ with $r\geq1$.
Suppose $\mathfrak p\notin\operatorname{Supp}(M)$.
Each $x_i/1$ is zero in $M_{\mathfrak p}$, so
Lemma \ref{algebra-lemma-localization-colimit} gives a common
$f\notin\mathfrak p$ for which all are zero in $M_f$.
For each original generator choose $n_i\geq1$ with
$f^{n_i}x_i=0$, and keep $n=\max\{n_1,\ldots,n_r\}$.
The complete equality
$f^nx_i=f^{n-n_i}(f^{n_i}x_i)=0$ shows that $f^n$ kills
every generator and therefore $f^nM=0$.
Thus $f^n\in I$ while $f^n\notin\mathfrak p$, proving
$\mathfrak p\notin V(I)$ and the required equality.
\end{proof}

\begin{lemma}
\label{editorial-source-lemma-support-finite-annihilator-detection}
Let $M$ be an arbitrary $R$-module. If finitely many
$m_1,\ldots,m_n\in M$ satisfy
$\operatorname{Ann}_R(M)=\bigcap_{i=1}^n\operatorname{Ann}_R(m_i)$,
then $\operatorname{Supp}(M)=V(\operatorname{Ann}_R(M))$.
\end{lemma}

\begin{proof}
Retain the original submodule $N=\sum_{i=1}^n Rm_i\subseteq M$.
An element kills $N$ exactly when it kills each $m_i$,
so $\operatorname{Ann}_R(N)=\operatorname{Ann}_R(M)$ by
the stated hypothesis. Since $N$ is finite, the preceding
lemma gives $\operatorname{Supp}(N)=V(\operatorname{Ann}_R(M))$.
Localizing the actual injection $N\hookrightarrow M$
is injective, so $\operatorname{Supp}(N)\subseteq\operatorname{Supp}(M)$.
For every $M$, the first half of the preceding proof
gives $\operatorname{Supp}(M)\subseteq V(\operatorname{Ann}_R(M))$.
The inclusions prove equality without replacing the
original $M$ by $N$. If the list is empty its intersection
of annihilators is $R$, hence $M=0$, which also satisfies
the conclusion. For example an arbitrary nonempty free
module is covered by this hypothesis: one original basis
vector already has annihilator zero, although the basis
may be infinite.
\end{proof}


\subsection{Editorial treatment of Algebra lines 9412--9444}

\hypertarget{algebra-context-101}{}

\noindent\href{https://github.com/stacks/stacks-project/blob/a04446e57ec1fbc252a871afcec7752fb2807b14/algebra.tex\#L9412}{Original source, lines 9412--9444}. Complete original and reviewed TeX passages accompany this note. Every intervention is separately recorded; the following treatment is editorial.


\begin{lemma}
\label{editorial-source-lemma-support-base-change}
Let $R \to R'$ be a ring map and let $M$ be a finite $R$-module.
Then $\text{Supp}(M \otimes_R R')$ is the inverse image of
$\text{Supp}(M)$.
\end{lemma}

\begin{proof}
Write the original map as $\varphi:R\to R'$ and the induced
spectrum map as $f:\Spec(R')\to\Spec(R)$.
Fix an original prime $\mathfrak p'\subset R'$ and its
contraction $\mathfrak p=\varphi^{-1}(\mathfrak p')$.
The local map is
$$
R_{\mathfrak p}\longrightarrow R'_{\mathfrak p'},
\qquad r/t\longmapsto\varphi(r)/\varphi(t).
$$
All denominators are valid, because $t\notin\mathfrak p$
implies $\varphi(t)\notin\mathfrak p'$. Its inverse image
of $\mathfrak p'R'_{\mathfrak p'}$ is
$\mathfrak pR_{\mathfrak p}$: a fraction in the former
ideal has numerator in $\mathfrak p'$, and contraction
of that numerator is the stated condition in $R$.

For any original module $M$, not yet assumed finite, the
exact localization-tensor comparison is
$$
(M\otimes_RR')_{\mathfrak p'}
\longrightarrow M_{\mathfrak p}\otimes_{R_{\mathfrak p}}R'_{\mathfrak p'},
\qquad \frac{m\otimes s}{u}\longmapsto\frac m1\otimes\frac su.
$$
Its inverse is
$$
\frac mr\otimes\frac su
\longmapsto\frac{m\otimes s}{\varphi(r)u},
\qquad r\notin\mathfrak p,\quad u\notin\mathfrak p'.
$$
The forward map is induced by the balanced map to the
right-hand module and then by its invertible action of
every $u\notin\mathfrak p'$. For the inverse, if
$v(r'm-rm')=0$ with $v\notin\mathfrak p$, multiplying
the corresponding numerator difference by
$\varphi(v)\notin\mathfrak p'$ gives zero. This checks
the equivalence relation in $M_{\mathfrak p}$. If
$w(u's-us')=0$ with $w\notin\mathfrak p'$, the same
$w$ kills the tensor numerator difference, checking the
other localization relation. Additivity uses common
denominators and tensor additivity. For $a/t\in R_{\mathfrak p}$,
the two balanced expressions map respectively to
$$
\frac{am\otimes s}{\varphi(t)\varphi(r)u}
=\frac{m\otimes\varphi(a)s}{\varphi(r)\varphi(t)u}.
$$
This proves balancing with every original factor. Composing
in one direction gives the original fraction immediately;
in the other, $m/1\otimes s/(\varphi(r)u)=m/r\otimes s/u$.
Thus the maps are inverse. In particular, for every $M$
and every $R\to R'$,
$$
\operatorname{Supp}_{R'}(M\otimes_RR')
\subseteq f^{-1}\operatorname{Supp}_R(M).
$$

The original residue-field map is the injective map
$$
\iota:\kappa(\mathfrak p)\longrightarrow\kappa(\mathfrak p'),
\qquad
\frac{r+\mathfrak p}{t+\mathfrak p}
\longmapsto
\frac{\varphi(r)+\mathfrak p'}{\varphi(t)+\mathfrak p'}.
$$
It exists and is injective because $R/\mathfrak p\to
R'/\mathfrak p'$ is injective and every displayed
nonzero denominator stays nonzero. There are exact maps
$$
\begin{split}
(M\otimes_RR')\otimes_{R'}\kappa(\mathfrak p')
&\longrightarrow M\otimes_R\kappa(\mathfrak p'),\\
(m\otimes s)\otimes\lambda
&\longmapsto m\otimes\overline s\lambda,
\end{split}
$$
with inverse $m\otimes\lambda\mapsto(m\otimes1)\otimes\lambda$,
and
$$
\begin{split}
(M\otimes_R\kappa(\mathfrak p))
\otimes_{\kappa(\mathfrak p)}\kappa(\mathfrak p')
&\longrightarrow M\otimes_R\kappa(\mathfrak p'),\\
(m\otimes\alpha)\otimes\beta
&\longmapsto m\otimes\iota(\alpha)\beta,
\end{split}
$$
with inverse $m\otimes\beta\mapsto(m\otimes1)\otimes\beta$.
The tensor balancing relations show independence and give
both composites on the displayed generators. Extension
along the field inclusion detects a nonzero vector space:
if $B$ is a basis, the map from the direct sum of copies
of $\kappa(\mathfrak p')$ indexed by $B$, sending its
$b$-coordinate $\lambda$ to $b\otimes\lambda$, is an
isomorphism, with inverse given by the finite basis
expansion of each vector. Thus it is zero exactly when
$B$ is empty.

For completeness, for any local ring $(A,\mathfrak m)$
and module $N$, the map
$N\otimes_A A/\mathfrak m\to N/\mathfrak mN$,
$n\otimes\overline a\mapsto an+\mathfrak mN$, has inverse
$n+\mathfrak mN\mapsto n\otimes1$; balancing and the
relations $\mathfrak mN=0$ in both targets prove the claim.
Applying this after the localization comparisons gives the
original fibre quotient
$$
M\otimes_R\kappa(\mathfrak p)
\simeq M_{\mathfrak p}/\mathfrak pM_{\mathfrak p},
$$
and the corresponding quotient for $M\otimes_RR'$ at
$\mathfrak p'$. Explicitly,
$m\otimes(a/t+\mathfrak pR_{\mathfrak p})$ maps to
$am/t+\mathfrak pM_{\mathfrak p}$, with inverse
$m/t+\mathfrak pM_{\mathfrak p}\mapsto
m\otimes(1/t+\mathfrak pR_{\mathfrak p})$.

Now restore the original finiteness hypothesis. The module
$M_{\mathfrak p}$ is finite over the local ring, so
Nakayama's Lemma \ref{algebra-lemma-NAK} says that it is nonzero
exactly when this fibre quotient is nonzero. If
$\mathfrak p\in\operatorname{Supp}_R(M)$, the field
extension and the displayed fibre maps show that the
original quotient
$$
(M\otimes_RR')_{\mathfrak p'}/
\mathfrak p'(M\otimes_RR')_{\mathfrak p'}
$$
is nonzero. Therefore its module is nonzero. Together
with the already proved reverse inclusion, this proves
the stated lemma, keeping both original primes and all maps.

\medskip\noindent
There are two further exact forms. First, for an arbitrary
$M$ put
$$
F_R(M)=\{\mathfrak p\in\Spec(R):
M\otimes_R\kappa(\mathfrak p)\ne0\}.
$$
The field-extension argument above did not use finiteness
or flatness. Hence for every ring map it proves
$$
F_{R'}(M\otimes_RR')=f^{-1}F_R(M).
$$
The localization comparison always gives
$F_R(M)\subseteq\operatorname{Supp}_R(M)$.
Consequently localization support commutes with every
base change if and only if
$$
F_R(M)=\operatorname{Supp}_R(M).
$$
For sufficiency, a prime in the inverse image of support
is in the inverse image of $F_R(M)$, hence in the new
fibre set and therefore in the new support. The general
reverse inclusion was proved above. For necessity, apply
the asserted support equality to the particular original
map $R\to\kappa(\mathfrak p)$. The field's unique prime
contracts to $\mathfrak p$, and its module is nonzero
exactly when $M\otimes_R\kappa(\mathfrak p)\ne0$.
This proves both directions of the criterion.

This criterion strictly extends finite generation. An
arbitrary direct sum of finite modules satisfies it.
Localization commutes with direct sums by the maps
$(m_i)_i/s\mapsto(m_i/s)_i$ and their inverse obtained
from the finite sum of the localized inclusions.
For coordinates $m_i/s_i$ indexed by a finite set $E$,
the inverse has denominator $\prod_{i\in E}s_i$ and
$i$-th numerator $m_i\prod_{j\in E,\,j\ne i}s_j$.
Tensor products commute with these direct sums by
$(m_i)_i\otimes c\mapsto(m_i\otimes c)_i$, with inverse
given by the sum of the tensor products of the original
inclusions. Balancing verifies these maps, and both
composites are identities on elements of finite support.
A direct sum is nonzero exactly when one summand is;
Nakayama for each finite summand therefore proves the
criterion for the sum. For a nonzero ring, a free module
on an infinite set is not finite: any finite collection
of vectors involves only finitely many basis coordinates,
so cannot generate a basis vector outside their union.
It supplies the strictness example.

\medskip\noindent
Second, if the original map $R\to R'$ is flat, support
commutes with this base change for every $M$. The local
map $R_{\mathfrak p}\to R'_{\mathfrak p'}$ is flat:
tensor with $R'$ preserves injections, and subsequent
localization at $\mathfrak p'$ preserves them. For any
$R_{\mathfrak p}$-module $N$, the comparison
$N\otimes_RR'_{\mathfrak p'}\simeq
N\otimes_{R_{\mathfrak p}}R'_{\mathfrak p'}$
sends each original elementary tensor to itself.
Its inverse is balanced over $R_{\mathfrak p}$ since
every $t\notin\mathfrak p$ acts invertibly in both factors;
thus the comparison follows with no injectivity assumption
on the ring map.

A flat local ring map $(A,\mathfrak m)\to(B,\mathfrak n)$
detects nonzero arbitrary modules. Choose $x\ne0$ in such
a module $N$ and put $J=\operatorname{Ann}_A(x)$.
The map $A/J\hookrightarrow N$, $a+J\mapsto ax$, is
injective. Flatness gives $B/JB\hookrightarrow N\otimes_AB$.
Here $J\subseteq\mathfrak m$ and
$JB\subseteq\mathfrak mB\subseteq\mathfrak n$, so
$B/JB\ne0$. This proves faithful flatness of the local
map, and applies to the original localization map above.
Its tensor comparison proves
$$
M_{\mathfrak p}\ne0
\quad\Longleftrightarrow\quad
(M\otimes_RR')_{\mathfrak p'}\ne0
$$
for every $M$, proving the asserted flat-base-change extension.

Flatness of the module alone does not imply the universal
support criterion. Retain the example
$R=\mathbf Z$, $M=\mathbf Q$, $R'=\mathbf F_p$ for a
prime integer $p$, with the quotient map. The module
$\mathbf Q$ is flat by its exact localization comparison,
proved in Example \ref{editorial-source-example-flat-infinite-intersections}.
At every prime of $\mathbf Z$, its localization is still
$\mathbf Q$: the map sending an original fraction $q/s$
to the same rational number is inverse to $q\mapsto q/1$.
Thus its localization support is all of $\Spec(\mathbf Z)$.
But
$$
q\otimes\overline a=(q/p)\otimes p\overline a=0
$$
in $\mathbf Q\otimes_{\mathbf Z}\mathbf F_p$, so the
new support is empty, whereas the inverse image of the
old support is the nonempty spectrum of $\mathbf F_p$.
The fibre set is exactly $\{(0)\}$: all nonzero prime
fibres vanish as just proved, while multiplication gives
$\mathbf Q\otimes_{\mathbf Z}\mathbf Q\simeq\mathbf Q$,
with inverse $q\mapsto q\otimes1$. For a fraction $a/b$,
$q\otimes a/b=(qa/b)\otimes1$ because $b$ is invertible
on both sides of the tensor, proving the inverse formula.
Moreover the annihilator of $\mathbf Q$ is detected by
its element $1$. Thus finite annihilator detection does
not imply the universal support criterion. Empty spectra
and zero modules satisfy all the stated identities as well.
\end{proof}


\subsection{Editorial treatment of Algebra lines 9446--9458}

\hypertarget{algebra-context-102}{}

\noindent\href{https://github.com/stacks/stacks-project/blob/a04446e57ec1fbc252a871afcec7752fb2807b14/algebra.tex\#L9446}{Original source, lines 9446--9458}. Complete original and reviewed TeX passages accompany this note. Every intervention is separately recorded; the following treatment is editorial.


\begin{lemma}
\label{editorial-source-lemma-support-element}
Let $R$ be a ring, let $M$ be an $R$-module, and let $m \in M$.
Then $\mathfrak p \in V(\text{Ann}(m))$ if and only if
$m$ does not map to zero in $M_\mathfrak p$.
\end{lemma}

\begin{proof}
If $\operatorname{Ann}_R(m)\subseteq\mathfrak p$ and
$m/1=0$ in the original $M_{\mathfrak p}$, there would
be $s\notin\mathfrak p$ with $sm=0$. This contradicts
$s\in\operatorname{Ann}_R(m)\subseteq\mathfrak p$.
Conversely, if $\operatorname{Ann}_R(m)$ is not contained
in $\mathfrak p$, take an original annihilator
$s\notin\mathfrak p$. It becomes a unit at $\mathfrak p$
and kills $m/1$, so $m/1=0$.
The exact cyclic comparison underlying this argument is
the injection
$$
R/\operatorname{Ann}_R(m)\longrightarrow M,\qquad
r+\operatorname{Ann}_R(m)\longmapsto rm,
$$
whose image is the actual submodule $Rm$ and whose
inverse on that image sends $rm$ to the indicated class.
Equality of two such expressions is exactly an
annihilator relation, proving independence. Localization
gives the injection $(r+\operatorname{Ann}_R(m))/s\mapsto rm/s$
into the original $M_{\mathfrak p}$. Thus the cyclic
module comparison preserves, rather than replaces, $M$.
\end{proof}

\begin{lemma}
\label{editorial-source-lemma-support-arbitrary-cyclic-union}
For every original $R$-module $M$,
$$
\operatorname{Supp}(M)=\bigcup_{m\in M}V(\operatorname{Ann}_R(m)).
$$
Consequently supports of arbitrary modules are stable
under specialization, and every subset of $\Spec(R)$
stable under specialization is the support of an $R$-module.
\end{lemma}

\begin{proof}
A localization $M_{\mathfrak p}$ is nonzero exactly when
some original $m\in M$ has $m/1\ne0$: a nonzero fraction
$m/s$ has nonzero numerator image since $s/1$ is a unit.
Lemma \ref{editorial-source-lemma-support-element} gives the displayed union.
If $\mathfrak p\subseteq\mathfrak q$ and
$\operatorname{Ann}_R(m)\subseteq\mathfrak p$, the same
annihilator is contained in $\mathfrak q$. Hence support
is stable under specialization with the original element
as its witness.

Localization commutes with arbitrary direct sums by the
maps
$$
(\bigoplus_{\lambda\in\Lambda}M_\lambda)_{\mathfrak p}
\longrightarrow\bigoplus_{\lambda\in\Lambda}(M_\lambda)_{\mathfrak p},
\qquad (m_\lambda)_\lambda/s\longmapsto(m_\lambda/s)_\lambda.
$$
For finitely many nonzero coordinates indexed by a set
$E\subseteq\Lambda$, its inverse sends
$(m_\lambda/s_\lambda)_{\lambda\in E}$ to the vector
with numerator coordinate
$m_\lambda\prod_{\mu\in E,\mu\ne\lambda}s_\mu$
and common denominator $\prod_{\mu\in E}s_\mu$.
All factors are outside $\mathfrak p$.
Equivalently the inverse is the sum of the localized
original inclusions $M_\lambda\to\bigoplus M_\lambda$,
which proves independence of representations; these
formulas verify both composites. Thus the support of
a direct sum is the union of the supports of its summands.

For a specialization-stable subset $T\subseteq\Spec(R)$,
take the specified module
$M_T=\bigoplus_{\mathfrak p\in T}R/\mathfrak p$.
Each summand has support $V(\mathfrak p)$ by
Lemma \ref{editorial-source-lemma-support-closed}, since its annihilator
is exactly $\mathfrak p$. Their union is $T$: it contains
each original $\mathfrak p\in T$, and every prime
containing one of them belongs to $T$ by specialization
stability. The empty subset gives the zero direct sum.
\end{proof}


\subsection{Editorial treatment of Algebra lines 9460--9484}

\hypertarget{algebra-context-103}{}

\noindent\href{https://github.com/stacks/stacks-project/blob/a04446e57ec1fbc252a871afcec7752fb2807b14/algebra.tex\#L9460}{Original source, lines 9460--9484}. Complete original and reviewed TeX passages accompany this note. Every intervention is separately recorded; the following treatment is editorial.


\begin{lemma}
\label{editorial-source-lemma-support-finite-presentation-constructible}
Let $R$ be a ring and let $M$ be an $R$-module.
If $M$ is a finitely presented $R$-module, then $\text{Supp}(M)$ is a
closed subset of $\Spec(R)$ whose complement is quasi-compact.
\end{lemma}

\begin{proof}
Retain the original finite presentation
$$
R^{\oplus m}\xrightarrow{A}R^{\oplus n}
\xrightarrow{\pi}M\longrightarrow0,
$$
where $A=(a_{ri})$ is the original $n\times m$ matrix.
Write $e_1,\ldots,e_n$ for the original standard basis
of $R^{\oplus n}$ and $a_i$ for column $i$ of $A$.
For each sorted subset $S=\{i_1<\cdots<i_n\}$ of
$\{1,\ldots,m\}$, set
$$
A_S=(a_{i_1},\ldots,a_{i_n}),\qquad
\Delta_S=\det(A_S),\qquad
I=(\Delta_S:|S|=n).
$$
Thus $I$ is exactly the original ideal of maximal minors.

Fix $\mathfrak p\in\Spec(R)$ and retain
$K=\kappa(\mathfrak p)$.
Applying $R\to R_{\mathfrak p}\to K$ to every matrix
entry gives $A_{\mathfrak p}$ and $A(\mathfrak p)$.
The localized presentation is exact. In detail, the last
map sends $v/s$ to $\pi(v)/s$ and is surjective by lifting
the original numerator. If $\pi(v)/s=0$, there is
$u\notin\mathfrak p$ with $u\pi(v)=0$. Then $uv=Aw$
for some $w\in R^{\oplus m}$, and the full equality
$$
\frac vs=\frac{uv}{us}
=A_{\mathfrak p}\left(\frac w{us}\right)
$$
proves its kernel statement. No injectivity of
$R\to R_{\mathfrak p}$ is assumed.
The fibre quotient comparison is
$$
K^n/\operatorname{im}A(\mathfrak p)
\longrightarrow M\otimes_RK,
\qquad
[(\lambda_1,\ldots,\lambda_n)]
\longmapsto\sum_{r=1}^n\pi(e_r)\otimes\lambda_r.
$$
Its inverse sends $m\otimes\lambda$ to
$[\lambda\overline v]$ for a lift $\pi(v)=m$.
Changing $v$ by $Aw$ changes this vector by
$\lambda A(\mathfrak p)\overline w$, which is in the
matrix image. Additivity and $R$-balancing follow from
the original linear map $\pi$. Both composites are
identities on the displayed generators. Hence
$$
\dim_K(M\otimes_RK)=n-\operatorname{rank}_K A(\mathfrak p).
$$

For $n\geq1$, a minor $\Delta_S\notin\mathfrak p$
gives the explicit right inverse
$$
A_{\mathfrak p}
\left(E_S\frac{\operatorname{adj}(A_S)_{\mathfrak p}}
{\Delta_S/1}\right)=I_n,
$$
where $E_S:R^n\to R^m$ sends its $\ell$-th basis
vector to the $i_\ell$-th original basis vector.
Indeed $AE_S=A_S$ and the adjugate identity gives the
displayed product. Thus $M_{\mathfrak p}=0$.
If $M_{\mathfrak p}=0$, the localized map is surjective;
preimages of the original basis vectors reduce to a
right inverse for $A(\mathfrak p)$, so the fibre is zero.
Conversely, if the fibre is zero, choose
$C\in\operatorname{Mat}_{m\times n}(K)$ with
$A(\mathfrak p)C=I_n$. Cauchy--Binet gives
$$
1=\sum_{\substack{S\subseteq\{1,\ldots,m\}\\|S|=n}}
\overline{\Delta_S}\det(C_{S,\{1,\ldots,n\}}).
$$
At least one original minor therefore lies outside
$\mathfrak p$. This proves, with explicit maps, all
the original equivalences
$$
M_{\mathfrak p}\ne0
\quad\Longleftrightarrow\quad
M\otimes_RK\ne0
\quad\Longleftrightarrow\quad
\operatorname{rank}_K A(\mathfrak p)<n.
$$
In particular $\operatorname{Supp}(M)=V(I)$.

If $n=0$, the target is zero and $M=0$. There is one
$0\times0$ minor indexed by the empty subset, with
determinant $1$, so $I=R$ and both sets are empty.
If $m<n$, there are no maximal minors, so $I=0$.
Every fibre then has dimension at least $n-m>0$, giving
support all of $\Spec(R)$. For the zero ring there are
no primes; every module in the presentation is zero,
and the same support identities hold.

There are only finitely many original maximal minors,
and the complement is the exact finite union
$$
\Spec(R)\setminus\operatorname{Supp}(M)
=\bigcup_{\substack{S\subseteq\{1,\ldots,m\}\\|S|=n}}D(\Delta_S).
$$
Each $D(f)$ is quasi-compact. To recall the proof,
localization identifies $\Spec(R_f)$ with $D(f)$:
contraction is one map and
$\mathfrak p\mapsto\{r/f^a:r\in\mathfrak p,\ a\geq0\}$
is its inverse. The prime localization correspondence
gives these inverse maps, and $D(r/f^a)$ corresponds
to $D(r)\cap D(f)$, proving that they are homeomorphisms.
For any ring $B$, refine a cover of $\Spec(B)$ by
basic opens $D(b_\lambda)$. If their ideal were proper,
a maximal ideal containing it would be missed. Hence
$1=\sum_{j=1}^t c_jb_{\lambda_j}$ for finitely many
original cover elements, and those basic opens cover.
This proves quasi-compactness of every spectrum and
of $D(f)$. A finite union of quasi-compact subspaces
is quasi-compact, by restricting any cover to each
of the finitely many subspaces and joining their finite
subcovers. This proves the original assertion. The
empty union is allowed; for $n=0$ the union is $D(1)$.

\medskip\noindent
The presentation also proves the more precise ideal
bounds. Put $H=\operatorname{Ann}_R(M)$. For $n\geq1$,
$$
I\subseteq H,\qquad H^n\subseteq I.
$$
To prove the first, retain the original column indices
and every cofactor sign in the adjugate identity:
$$
\Delta_Se_r
=\sum_{\ell=1}^n(-1)^{r+\ell}
\det A_{\{1,\ldots,n\}\setminus\{r\},\,
S\setminus\{i_\ell\}}\ a_{i_\ell}.
$$
The remaining rows and columns have inherited increasing
order. In coordinate $k$ the right side is the cofactor
expansion along row $r$ of the matrix with that row
replaced by row $k$. It is $\Delta_S$ if $k=r$ and
zero if $k\ne r$, since then two rows coincide.
Applying $\pi$ proves $\Delta_S\pi(e_r)=0$ for all
$r$. These original elements generate $M$, so
$\Delta_S\in H$ and $I\subseteq H$.

For the second bound take arbitrary original
$t_1,\ldots,t_n\in H$. Since $\pi(t_je_j)=0$,
choose $b_j\in R^m$ with $Ab_j=t_je_j$.
Let $B=(b_{ij})$ be the $m\times n$ matrix with these
columns. Then $AB=\operatorname{diag}(t_1,\ldots,t_n)$.
The full ordered-column determinant expansion is
$$
t_1\cdots t_n=\det(AB)
=\sum_{u_1=1}^m\cdots\sum_{u_n=1}^m
\left(\prod_{j=1}^n b_{u_j,j}\right)
\det(a_{u_1},\ldots,a_{u_n}).
$$
Terms with repeated indices vanish because the determinant
has repeated columns. For distinct indices let
$S=\{i_1<\cdots<i_n\}$ be their sorted set and let
$u_j=i_{\sigma(j)}$. The determinant is
$\operatorname{sgn}(\sigma)\Delta_S$. Thus the full sum is
$$
\sum_{\substack{S=\{i_1<\cdots<i_n\}\\S\subseteq\{1,\ldots,m\}}}
\Delta_S
\left(\sum_{\sigma\in\mathfrak S_n}
\operatorname{sgn}(\sigma)\prod_{j=1}^n b_{i_{\sigma(j)},j}\right).
$$
For $\tau=\sigma^{-1}$ the inner product is
$\prod_{\ell=1}^n b_{i_\ell,\tau(\ell)}$, with the same
sign, so the exact Cauchy--Binet identity is
$$
t_1\cdots t_n
=\sum_{\substack{S\subseteq\{1,\ldots,m\}\\|S|=n}}
\Delta_S\det(B_{S,\{1,\ldots,n\}}).
$$
This proves that every product of $n$ possibly different
elements of $H$ is in $I$. Those products generate $H^n$,
so it proves the ideal-power statement, not only the
statement about individual $n$-th powers. If $m<n$
the sum is empty, giving $H^n=0$. If $n=0$, both $H$
and $I$ are $R$ and the convention $H^0=R$ retains the
two containments. They also hold over the zero ring.
For $n\geq1$ they imply $\sqrt I=\sqrt H$: one inclusion
uses $I\subseteq H$; for $h\in H$, the other uses
$h^n\in H^n\subseteq I$, followed by radicals.
Thus the exact support formula is also
$V(I)=V(\operatorname{Ann}_R(M))$.

The exponent $n$ is sharp uniformly. For a field $k$
take the original presentation $A=xI_n$ over $R=k[x]$,
with $m=n\geq1$. It gives
$$
M=(R/(x))^n,\qquad H=(x),\qquad I=(x^n).
$$
Here an element kills all basis quotient classes exactly
when it belongs to $(x)$, and the single maximal minor
is $x^n$. Thus $H^n=I$; for $n\geq2$ the polynomial
$x^{n-1}$ belongs to $H^{n-1}$ and not to $(x^n)$.
Also $I\subsetneq H$ for these $n$.

\medskip\noindent
There is an exact weaker condition for the topological
conclusion. For any original module $N$ and specified
ideal $L$ with $\operatorname{Supp}(N)=V(L)$, its
complement is quasi-compact if and only if
$$
\sqrt L=\sqrt{(f_1,\ldots,f_t)}
\quad\hbox{for finitely many }f_1,\ldots,f_t\in L.
$$
Indeed $D(L)=\bigcup_{f\in L}D(f)$, so quasi-compactness
gives a finite subcover of this specific cover. Taking
complements gives $V(L)=V((f_1,\ldots,f_t))$ and hence
the radical equality. Conversely that equality gives the
same finite union of quasi-compact basic opens.
The empty list is allowed, with generated ideal zero.
The radical characterization used here follows directly:
every prime containing $J$ contains $\sqrt J$. If
$a\notin\sqrt J$, the ring $(R/J)_a$ is nonzero, since
otherwise some original $a^N$ would be in $J$.
A maximal ideal of this localization contracts to a prime
containing $J$ and avoiding $a$. Consequently $\sqrt J$
is exactly the intersection of the primes containing $J$,
which proves the asserted equivalence of radicals and
closed sets without a finiteness assumption on $J$.

Finite presentation is not necessary for this topology.
Take the specified example
$$
R=k[x,y_1,y_2,\ldots],\qquad
I=(x^2,xy_i:i\geq1),\qquad M=R/I,
$$
where $k$ is a field. Here $I\subseteq(x)$ and $x^2\in I$,
while $R/(x)=k[y_1,y_2,\ldots]$ is a domain: two
nonzero polynomials belong to a polynomial ring in
finitely many variables over $k$, where their product
is nonzero. Thus $\sqrt I=(x)$.
The original cyclic quotient has support $V(I)$.
If $a\in I\setminus\mathfrak p$, the unit $a/1$
makes $(R/I)_{\mathfrak p}$ zero; if $I\subseteq\mathfrak p$,
this quotient maps onto the nonzero residue field.
Therefore the support is $V(x)$ and its complement
is $D(x)=D(x^2)$, a quasi-compact basic open. In the
weaker criterion one can take the actual original
element $f_1=x^2\in I$; there is no replacement of
the original ideal by its radical.

To prove that $I$ is not finitely generated, suppose
$I=(h_1,\ldots,h_r)$ and choose $N$ containing all
indices of variables $y_i$ occurring in any $h_j$.
Define $\rho:R\to k[x,z]$ by fixing $x$, sending
$y_{N+1}$ to $z$, and sending every other $y_i$ to
zero. Define also $\rho_0:R\to k[x]$ by fixing $x$
and sending all the $y_i$ to zero. Since $y_{N+1}$
does not occur in $h_j$, one has
$\rho(h_j)=\rho_0(h_j)\in(x^2)$: the last membership
follows by applying $\rho_0$ to the original generators
$x^2,xy_i$ of $I$. Every combination of the proposed
finite generators therefore maps into $(x^2)$.
But $xy_{N+1}\in I$ maps to $xz\notin(x^2)$.
Indeed its coefficient of $x^1$ over $k[z]$ is the
nonzero polynomial $z$, whereas every multiple of
$x^2$ has that coefficient zero. This is a contradiction.

Finally, for every ring and ideal, the cyclic module
$R/I$ is finitely presented exactly when $I$ is
finitely generated. A finite generating list gives
the presentation $R^t\to R\to R/I\to0$ with that
list as the first matrix. Conversely choose a finite
presentation with surjection $\pi:R^b\to R/I$ and
finitely generated kernel $K=(k_1,\ldots,k_t)$.
Let $q:R\to R/I$ be the original quotient. Choose
$v\in R^b$ with $\pi(v)=1+I$, and an $R$-linear
lift $f:R^b\to R$ of $\pi$ by choosing representatives
for the images of the original basis vectors.
Thus $qf=\pi$, $f(v)-1\in I$, and $f(k_j)\in I$.
For each original $a\in I$, the vector $av$ lies
in $K$, so $av=\sum_j c_jk_j$ for some coefficients.
Retaining both original terms gives the exact identity
$$
a=f(av)+a(1-f(v))
=\sum_{j=1}^t c_jf(k_j)+a(1-f(v)).
$$
Every listed generator is in $I$, and the identity
proves the reverse inclusion in
$$
I=(f(k_1),\ldots,f(k_t),1-f(v)).
$$
This covers $R/I=0$ as well, including $b=0$, where
$1-f(v)=1$ supplies the unit. Applying it to the
specified infinitely generated ideal proves that the
example's original $M$ is not finitely presented,
despite its quasi-compact support complement.
\end{proof}


\subsection{Editorial treatment of Algebra lines 9486--9508}

\hypertarget{algebra-context-104}{}

\noindent\href{https://github.com/stacks/stacks-project/blob/a04446e57ec1fbc252a871afcec7752fb2807b14/algebra.tex\#L9486}{Original source, lines 9486--9508}. Complete original and reviewed TeX passages accompany this note. Every intervention is separately recorded; the following treatment is editorial.


\begin{lemma}
\label{editorial-source-lemma-support-quotient}
Let $R$ be a ring, let $M$ be an $R$-module, and let $I\subseteq R$ be an ideal.
\begin{enumerate}
\item If $M$ is finite then the support
of $M/IM$ is $\text{Supp}(M) \cap V(I)$.
\item If $N \subset M$, then $\text{Supp}(N) \subset
\text{Supp}(M)$.
\item If $Q$ is a quotient module of $M$ then $\text{Supp}(Q) \subset
\text{Supp}(M)$.
\item If $0 \to N \to M \to Q \to 0$ is a short exact sequence
then $\text{Supp}(M) = \text{Supp}(Q) \cup \text{Supp}(N)$.
\end{enumerate}
\end{lemma}

\begin{proof}
Localization preserves the original injections, surjections
and short exact sequences. For (2), the map
$N_{\mathfrak p}\to M_{\mathfrak p}$ is injective, so
the first module can be nonzero only if the second is.
For (3), the map $M_{\mathfrak p}\to Q_{\mathfrak p}$
is surjective, so the same conclusion holds for a nonzero
quotient. In (4), the exact localized sequence shows
$M_{\mathfrak p}=0$ if and only if both
$N_{\mathfrak p}=0$ and $Q_{\mathfrak p}=0$, which is
precisely the stated union formula.

For (1), retain the exact quotient comparison
$$
(M/IM)_{\mathfrak p}\longrightarrow M_{\mathfrak p}/IM_{\mathfrak p},
\qquad (m+IM)/s\longmapsto m/s+IM_{\mathfrak p}.
$$
Its inverse sends $m/s+IM_{\mathfrak p}$ to $(m+IM)/s$.
The original quotient map localizes to a surjection with
kernel $(IM)_{\mathfrak p}=IM_{\mathfrak p}$: the equality
follows by writing finite sums of products over a common
denominator. This proves well-definedness and both
composites of the displayed maps. If $M_{\mathfrak p}=0$,
the quotient is zero. If some $a\in I$ lies outside
$\mathfrak p$, it is a unit in $R_{\mathfrak p}$ and
$IM_{\mathfrak p}=M_{\mathfrak p}$, so the quotient is
again zero. Conversely, if $M_{\mathfrak p}\ne0$ and
$I\subseteq\mathfrak p$, then $IR_{\mathfrak p}$ lies
in the maximal ideal of the local ring. The module
$M_{\mathfrak p}$ is finite, so Nakayama's
Lemma \ref{algebra-lemma-NAK} prevents
$M_{\mathfrak p}/IM_{\mathfrak p}$ from being zero.
These cases prove the exact intersection formula.
\end{proof}


\subsection{Editorial treatment of Algebra lines 9566--9586}

\hypertarget{algebra-context-105}{}

\noindent\href{https://github.com/stacks/stacks-project/blob/a04446e57ec1fbc252a871afcec7752fb2807b14/algebra.tex\#L9566}{Original source, lines 9566--9586}. Complete original and reviewed TeX passages accompany this note. Every intervention is separately recorded; the following treatment is editorial.


\begin{lemma}
\label{editorial-source-lemma-open-going-down}
Let $R \to S$ be a ring map. If the induced map
$\varphi : \Spec(S) \to \Spec(R)$ is open, then
$R \to S$ satisfies going down.
\end{lemma}

\begin{proof}
Retain $\mathfrak p\subseteq\mathfrak p'$ and the original
$\mathfrak q'$ lying over $\mathfrak p'$.
For $g\notin\mathfrak q'$, the open set $D(g)$ contains
$\mathfrak q'$, so its image is an open neighbourhood of
$\mathfrak p'$. This image contains a standard open
$D(f)$ with $f\notin\mathfrak p'$. Then
$f\notin\mathfrak p$, so it also contains $\mathfrak p$.
The fibre criterion in Lemma \ref{editorial-source-lemma-in-image} gives
$S_g\otimes_R\kappa(\mathfrak p)\ne0$ for every such $g$.

Here is the precise directed localization argument. Index
finite subsets $E\subseteq S\setminus\mathfrak q'$ by
inclusion and set $g_E=\prod_{g\in E}g$, with $g_\varnothing=1$.
For $E\subseteq F$ the map is
$$
S_{g_E}\longrightarrow S_{g_F},\qquad
\frac{s}{g_E^n}\longmapsto
\frac{s\,g_{F\setminus E}^{\,n}}{g_F^n}.
$$
It is the original localization map, since all factors
in the product $g_F$ become units. The map from their
colimit to $S_{\mathfrak q'}$ sends each fraction to the
same original fraction. It is surjective: $s/t$ comes
from $E=\{t\}$. If $s/g_E^n$ maps to zero, there is
$u\notin\mathfrak q'$ with $us=0$. At the stage
$F=E\cup\{u\}$, $u$ is a unit, so the image of $s$
and then that fraction is zero. This proves injectivity
and the representative formula for the inverse. It is
the comparison of Lemma \ref{algebra-lemma-localization-colimit}.

Tensor-colimit compatibility gives the unital ring
comparison
$$
S_{\mathfrak q'}\otimes_R\kappa(\mathfrak p)
\cong\colim_E(S_{g_E}\otimes_R\kappa(\mathfrak p)),
$$
sending a stage tensor $s/g_E^n\otimes c$ to the tensor
of the same fraction and $c$; the inverse follows by
representing the finitely many fractions at a common
stage, as in Lemma \ref{editorial-source-lemma-tensor-products-commute-with-limits}.
Every stage ring is nonzero by the previous paragraph,
because $g_E\notin\mathfrak q'$. Their colimit is nonzero:
if its unit were zero, it would become zero at a later
stage, but every transition is unital and that later
ring is nonzero.

Choose a prime of this nonzero fibre ring and pull it
back along $S\to S_{\mathfrak q'}\to
S_{\mathfrak q'}\otimes_R\kappa(\mathfrak p)$ to a prime
$\mathfrak q$ of $S$. Every element outside
$\mathfrak q'$ maps to a unit, so
$\mathfrak q\subseteq\mathfrak q'$.
The contraction to the field $\kappa(\mathfrak p)$ is
its zero prime, and the original map $R\to\kappa(\mathfrak p)$
has kernel $\mathfrak p$. Thus the contraction of
$\mathfrak q$ to $R$ is exactly $\mathfrak p$.

In fact this argument proves the exact weaker criterion:
$R\to S$ satisfies going down if and only if the image
of every principal open $D(g)\subseteq\Spec(S)$ is stable
under generalization. For sufficiency that stability
supplies $\mathfrak p$ in each image containing
$\mathfrak p'$, and the remaining argument just given
applies without openness. Conversely, if going down
holds and $\mathfrak q'\in D(g)$ lies over
$\mathfrak p'$, every $\mathfrak p\subseteq\mathfrak p'$
has a lift $\mathfrak q\subseteq\mathfrak q'$.
Since $g\notin\mathfrak q'$, also $g\notin\mathfrak q$,
so the lift still belongs to the original $D(g)$.
This proves necessity and both directions of the criterion.

Openness is stronger than this criterion. The flat map
$\mathbf Z\to\mathbf Q$ satisfies going down by
Lemma \ref{editorial-source-lemma-flat-going-down}, but the image of its
whole spectrum is $\{(0)\}$, which is not open in
$\Spec(\mathbf Z)$. Indeed any open neighbourhood of
$(0)$ contains $D(n)$ for some nonzero integer $n$.
A prime divisor $p$ of $|n|+1>1$ does not divide $n$,
so $(p)\in D(n)$ and $(p)\ne(0)$. This proves the
failure of openness. Flatness of $\mathbf Q$ was proved
in Example \ref{editorial-source-example-flat-infinite-intersections}.
\end{proof}


\subsection{Editorial treatment of Algebra lines 9588--9603}

\hypertarget{algebra-context-106}{}

\noindent\href{https://github.com/stacks/stacks-project/blob/a04446e57ec1fbc252a871afcec7752fb2807b14/algebra.tex\#L9588}{Original source, lines 9588--9603}. Complete original and reviewed TeX passages accompany this note. Every intervention is separately recorded; the following treatment is editorial.


\begin{lemma}
\label{editorial-source-lemma-going-up-down-specialization}
Let $R \to S$ be a ring map.
\begin{enumerate}
\item $R \to S$ satisfies going down if and only if
generalizations lift along the map $\Spec(S) \to \Spec(R)$,
see Topology, Definition \ref{topology-definition-lift-specializations}.
\item $R \to S$ satisfies going up if and only if
specializations lift along the map $\Spec(S) \to \Spec(R)$,
see Topology, Definition \ref{topology-definition-lift-specializations}.
\end{enumerate}
\end{lemma}

\begin{proof}
For the original map $\varphi:R\to S$, the spectrum map
sends a prime $\mathfrak q$ to $\varphi^{-1}(\mathfrak q)$.
The point of $\mathfrak p$ specializes to the point of
$\mathfrak p'$ exactly when $\mathfrak p\subseteq\mathfrak p'$,
because the closure of the point of $\mathfrak p$ is
$V(\mathfrak p)$. The same statement holds for primes of $S$.
Thus, given the original $\mathfrak p\subseteq\mathfrak p'$
and a prime $\mathfrak q'$ over $\mathfrak p'$, lifting
the generalization is precisely finding
$\mathfrak q\subseteq\mathfrak q'$ with contraction
$\mathfrak p$. This is the stated going-down condition,
in both directions. Given instead a prime $\mathfrak q$
over $\mathfrak p$, lifting the specialization is precisely
finding $\mathfrak q'\supseteq\mathfrak q$ with contraction
$\mathfrak p'$, the going-up condition, again in both
directions. These identifications agree exactly with
Topology, Definition \ref{topology-definition-lift-specializations}.
\end{proof}


\subsection{Editorial treatment of Algebra lines 9605--9615}

\hypertarget{algebra-context-107}{}

\noindent\href{https://github.com/stacks/stacks-project/blob/a04446e57ec1fbc252a871afcec7752fb2807b14/algebra.tex\#L9605}{Original source, lines 9605--9615}. Complete original and reviewed TeX passages accompany this note. Every intervention is separately recorded; the following treatment is editorial.


\begin{lemma}
\label{editorial-source-lemma-going-up-down-composition}
Suppose $R \to S$ and $S \to T$ are ring maps satisfying
going down. Then so does $R \to T$. Similarly for going up.
\end{lemma}

\begin{proof}
Write the original maps as $\varphi:R\to S$ and
$\psi:S\to T$. For going down, take
$\mathfrak p\subseteq\mathfrak p'$ in $R$ and a prime
$\mathfrak r'$ of $T$ contracting to $\mathfrak p'$.
Put $\mathfrak q'=\psi^{-1}(\mathfrak r')$.
Going down for $\varphi$ gives
$\mathfrak q\subseteq\mathfrak q'$ contracting to
$\mathfrak p$. Going down for $\psi$ then gives
$\mathfrak r\subseteq\mathfrak r'$ with
$\psi^{-1}(\mathfrak r)=\mathfrak q$. Its original
composite contraction is
$(\psi\circ\varphi)^{-1}(\mathfrak r)
=\varphi^{-1}(\mathfrak q)=\mathfrak p$.

For going up, take the same chain in $R$ and a prime
$\mathfrak r$ of $T$ contracting to $\mathfrak p$.
Set $\mathfrak q=\psi^{-1}(\mathfrak r)$.
Going up for $\varphi$ supplies
$\mathfrak q'\supseteq\mathfrak q$ contracting to
$\mathfrak p'$, and going up for $\psi$ supplies
$\mathfrak r'\supseteq\mathfrak r$ contracting to
$\mathfrak q'$. The same inverse-image composition gives
$(\psi\circ\varphi)^{-1}(\mathfrak r')=\mathfrak p'$.
These are the precise prime maps underlying
Lemma \ref{editorial-source-lemma-going-up-down-specialization} and
Topology, Lemma \ref{topology-lemma-lift-specialization-composition}.
\end{proof}


\subsection{Editorial treatment of Algebra lines 9617--9649}

\hypertarget{algebra-context-108}{}

\noindent\href{https://github.com/stacks/stacks-project/blob/a04446e57ec1fbc252a871afcec7752fb2807b14/algebra.tex\#L9617}{Original source, lines 9617--9649}. Complete original and reviewed TeX passages accompany this note. Every intervention is separately recorded; the following treatment is editorial.


\begin{lemma}
\label{editorial-source-lemma-image-stable-specialization-closed}
Let $R \to S$ be a ring map. Let $T \subset \Spec(R)$
be the image of $\Spec(S)$. If $T$ is stable under specialization,
then $T$ is closed. More precisely, for every ring map
$\varphi:R\to S$, without a specialization hypothesis,
$$
\overline T=V(\ker\varphi)
=\{\mathfrak p:\text{there is }\mathfrak p'\in T
\text{ with }\mathfrak p'\subseteq\mathfrak p\}.
$$
\end{lemma}

\begin{proof}
We retain both original proofs and all original rings.
Write the map as $\varphi:R\to S$ and put $I=\ker\varphi$.

\medskip\noindent
First proof. Let $\mathfrak p$ be a point of $\overline T$.
Every $f\notin\mathfrak p$ gives $D(f)\cap T\ne\varnothing$.
The prime localization correspondence identifies this
intersection with the image of $\Spec(S_{\varphi(f)})$,
so $S_{\varphi(f)}\ne0$.
The directed localization comparison with $S_{\mathfrak p}$
can be written without suppressing the original coefficient map.
Index finite subsets $E\subseteq R\setminus\mathfrak p$
by inclusion and put $f_E=\prod_{f\in E}f$.
For $E\subseteq F$, the transition is
$$
S_{\varphi(f_E)}\longrightarrow S_{\varphi(f_F)},\qquad
\frac{s}{\varphi(f_E)^n}\longmapsto
\frac{s\varphi(f_{F\setminus E})^n}{\varphi(f_F)^n}.
$$
The colimit map to $S_{\mathfrak p}=(R\setminus\mathfrak p)^{-1}S$
sends a fraction to the same fraction. It is surjective
because $s/\varphi(t)$ comes from $E=\{t\}$.
If an original numerator $s$ maps to zero, some
$u\notin\mathfrak p$ has $\varphi(u)s=0$; the numerator
then vanishes at the stage $E\cup\{u\}$ where
$\varphi(u)$ is a unit. Thus the map is also injective.
Every stage is nonzero, including $E=\varnothing$,
because $f_E\notin\mathfrak p$; their common unit
cannot die in the colimit, since that would make it
zero at a nonzero later stage. Hence $S_{\mathfrak p}\ne0$.

Choose a prime of this nonzero ring and contract it
to a prime $\mathfrak q\subset S$. All
$\varphi(t)$, $t\notin\mathfrak p$, are units in
$S_{\mathfrak p}$, so
$\mathfrak p'=\varphi^{-1}(\mathfrak q)\subseteq\mathfrak p$.
This is a point of the original image $T$.
Thus every point of $\overline T$ is a specialization
of a point of $T$. The reverse containment follows
because $\overline T$ is closed and contains $T$.

The same localization argument identifies the closed set
exactly. Every prime in $T$ contains $I$, so
$\overline T\subseteq V(I)$. Conversely, for
$I\subseteq\mathfrak p$, the ring $S_{\mathfrak p}$
cannot be zero: the equality $1/1=0$ would give
$t\notin\mathfrak p$ with $\varphi(t)1=0$,
contradicting $t\in I\subseteq\mathfrak p$.
Its prime again contracts to a point of $T$ contained
in $\mathfrak p$. This proves both displayed equalities.
If $T$ is stable under specialization, they imply
$T=V(I)$ and hence the asserted closedness.

\medskip\noindent
Second proof. Retain the original quotient map
$q:R\to R/I$ and the induced injective map
$\overline\varphi:R/I\to S$, $r+I\mapsto\varphi(r)$.
The spectrum map for $q$ is the homeomorphism onto
$V(I)$ given by contraction, with inverse
$\mathfrak p\mapsto\mathfrak p/I$.
For each $\mathfrak p\in V(I)$ choose a minimal prime
$\overline{\mathfrak p}_0\subseteq\mathfrak p/I$
of $R/I$. Existence follows from Zorn's lemma: the
intersection of a chain of primes contained in
$\mathfrak p/I$ is prime and is a lower bound.
Indeed if $ab$ belongs to that intersection and $a$
is absent from one chain member, primality forces
$b$ into that member and all smaller members, and
then into all larger members as well.
Lemma \ref{algebra-lemma-injective-minimal-primes-in-image}
gives a prime $\mathfrak q\subset S$ whose contraction
along $\overline\varphi$ is $\overline{\mathfrak p}_0$.
Along the original map its contraction is
$$
\varphi^{-1}(\mathfrak q)
=q^{-1}(\overline{\mathfrak p}_0)\subseteq\mathfrak p.
$$
Thus every point of $V(I)$ is again a specialization
of an original image point; no original ring has been
replaced in the comparison. If $S=0$, then $I=R$,
and all three sets in the statement are empty.
\end{proof}


\subsection{Editorial treatment of Algebra lines 9651--9679}

\hypertarget{algebra-context-109}{}

\noindent\href{https://github.com/stacks/stacks-project/blob/a04446e57ec1fbc252a871afcec7752fb2807b14/algebra.tex\#L9651}{Original source, lines 9651--9679}. Complete original and reviewed TeX passages accompany this note. Every intervention is separately recorded; the following treatment is editorial.


\begin{lemma}
\label{editorial-source-lemma-going-up-closed}
Let $R \to S$ be a ring map. The following are equivalent:
\begin{enumerate}
\item Going up holds for $R \to S$, and
\item the map $\Spec(S) \to \Spec(R)$ is closed.
\end{enumerate}
\end{lemma}

\begin{proof}
Write the original map as $\varphi:R\to S$ and its
spectrum map as $f$. Suppose first that $f$ is closed.
Given $\mathfrak p\subseteq\mathfrak p'$ and a prime
$\mathfrak q$ of $S$ with $\varphi^{-1}(\mathfrak q)=\mathfrak p$,
the closed set $V(\mathfrak q)$ contains $\mathfrak q$.
Its image is closed and contains $\mathfrak p$, so it
contains $\mathfrak p'$. Thus there exists
$\mathfrak q'\in V(\mathfrak q)$ with
$\varphi^{-1}(\mathfrak q')=\mathfrak p'$.
This says $\mathfrak q\subseteq\mathfrak q'$ and proves
going up. These are the exact prime maps in
Topology, Lemma \ref{topology-lemma-closed-open-map-specialization}.

Conversely suppose going up holds. For an original
ideal $I\subseteq S$, the image of $V(I)$ is stable
under specialization. Given a prime $\mathfrak q$
containing $I$ and a specialization of its contraction,
going up supplies a containing prime $\mathfrak q'$;
it still contains the original $I$ and maps to that
specialization. The quotient map $S\to S/I$ identifies
$\Spec(S/I)$ with $V(I)$ by
$\overline{\mathfrak q}\mapsto\{s:s+I\in\overline{\mathfrak q}\}$,
with inverse $\mathfrak q\mapsto\mathfrak q/I$.
The composite $R\to S/I$ has kernel $\varphi^{-1}(I)$.
Applying Lemma \ref{editorial-source-lemma-image-stable-specialization-closed}
to this composite proves not only that the image is
closed but also the exact formula
$$
f(V(I))=V(\varphi^{-1}(I)).
$$
This proves closedness for every original closed set.
The quotient map itself satisfies going up: a larger
prime of $S$ above one containing $I$ also contains
$I$, so its quotient supplies the lift. Together
with Lemma \ref{editorial-source-lemma-going-up-down-composition}, this
also verifies the composite-map argument in the original proof.

More generally, for any ring map and any ideal $I$
the same quotient comparison proves
$$
\overline{f(V(I))}=V(\varphi^{-1}(I)).
$$
Consequently going up is equivalent to the unclosed
image equality for every $I$. Necessity was proved
above; sufficiency implies every closed-set image is
closed, so the first paragraph supplies the prime lift.
\end{proof}


\subsection{Editorial treatment of Algebra lines 9681--9705}

\hypertarget{algebra-context-110}{}

\noindent\href{https://github.com/stacks/stacks-project/blob/a04446e57ec1fbc252a871afcec7752fb2807b14/algebra.tex\#L9681}{Original source, lines 9681--9705}. Complete original and reviewed TeX passages accompany this note. Every intervention is separately recorded; the following treatment is editorial.


\begin{lemma}
\label{editorial-source-lemma-constructible-stable-specialization-closed}
Let $R$ be a ring. Let $E \subset \Spec(R)$ be a constructible subset.
\begin{enumerate}
\item If $E$ is stable under specialization, then $E$ is closed.
\item If $E$ is stable under generalization, then $E$ is open.
\end{enumerate}
\end{lemma}

\begin{proof}
First proof. By Lemma \ref{editorial-source-lemma-constructible}, retain
a finite expression
$$
E=\bigcup_{i=1}^a\left(D(f_i)\cap
V(g_{i1},\ldots,g_{ib_i})\right).
$$
For $J_i=(g_{i1},\ldots,g_{ib_i})$, the original
quotient and localization maps give
$\Spec((R/J_i)_{f_i})$ the image $D(f_i)\cap V(J_i)$.
Contraction is the map; the inverse on this subset
sends $\mathfrak p$ to $(\mathfrak p/J_i)_{f_i}$.
The finite product $A=\prod_i(R/J_i)_{f_i}$ therefore
has spectrum image $E$ under the diagonal original
map $R\to A$, as in Lemma \ref{editorial-source-lemma-constructible-is-image}.
Indeed the coordinate idempotents sum to $1$ and
have pairwise zero products, so every prime selects
exactly one factor; this identifies its contraction
with that factor's contraction. The empty product
is the zero ring and gives the empty set.
If $E$ is stable under specialization, the preceding
image lemma proves it closed.

The complement $E^c$ is constructible by
Topology, Lemma \ref{topology-lemma-constructible}.
If $E$ is stable under generalization, $E^c$ is
stable under specialization: a specialization from
a point outside $E$ cannot enter $E$, since the
reverse generalization would then enter $E$ too.
Applying the first assertion to $E^c$ proves it
closed, hence $E$ open.
Each original basic term $D(f_i)\cap V(J_i)$ is
quasi-compact, being the image of the spectrum of
$(R/J_i)_{f_i}$. Therefore every constructible subset
of $\Spec(R)$ is quasi-compact by this finite union.
The complement is constructible as well. Thus in (1)
the closed set has quasi-compact complement, and in
(2) the open set is quasi-compact.

\medskip\noindent
Second proof, including its stronger scope. By
Lemma \ref{algebra-lemma-spec-spectral}, $X=\Spec(R)$ is
spectral. Its constructible topology has the
quasi-compact opens and their complements as a
subbasis; these sets are therefore both open and
closed in that topology. The topology is
quasi-compact Hausdorff by Topology, Lemma
\ref{topology-lemma-constructible-hausdorff-quasi-compact}.
Let $C\subseteq X$ be closed in the constructible
topology, which is weaker as a hypothesis than
being constructible, and let $x\in\overline C$
for the original topology. For every quasi-compact
open neighbourhood $U$ of $x$, the set $U\cap C$
is nonempty and is closed in the constructible
topology. Finite intersections of these $U$ are
again quasi-compact open neighbourhoods, so this
family has the finite intersection property.
Compactness gives $y\in C$ lying in every such $U$.
Since these sets form a neighbourhood basis of $x$,
$x$ is a specialization of $y$. Thus
$$
\overline C=\{x:\text{some }y\in C\text{ specializes to }x\}.
$$
If $C$ is stable under specialization it is closed
in the original topology. By taking complements,
any subset open in the constructible topology
and stable under generalization is open in the
original topology. This gives the full stronger
form of Topology, Lemma
\ref{topology-lemma-constructible-stable-specialization-closed},
as well as both original assertions. The stronger
scope does not assert that these sets or their
complements are constructible or quasi-compact in
all the cases just stated.
\end{proof}


\subsection{Editorial treatment of Algebra lines 9707--9737}

\hypertarget{algebra-context-111}{}

\noindent\href{https://github.com/stacks/stacks-project/blob/a04446e57ec1fbc252a871afcec7752fb2807b14/algebra.tex\#L9707}{Original source, lines 9707--9737}. Complete original and reviewed TeX passages accompany this note. Every intervention is separately recorded; the following treatment is editorial.


\begin{proposition}
\label{editorial-source-proposition-fppf-open}
Let $R \to S$ be flat and of finite presentation.
Then $\Spec(S) \to \Spec(R)$ is open.
More generally this holds for any ring map $R \to S$ of
finite presentation which satisfies going down.
\end{proposition}

\begin{proof}
If the original map $\varphi:R\to S$ is flat, it
satisfies going down by Lemma \ref{editorial-source-lemma-flat-going-down}.
Thus to prove the proposition it suffices to retain
finite presentation and going down.
For an original $g\in S$, write a finite presentation
$S=R[x_1,\ldots,x_n]/(h_1,\ldots,h_m)$ and choose
an original polynomial $G(x_1,\ldots,x_n)$ representing
$g$. Then the localization has the finite presentation
$$
S_g\simeq
R[x_1,\ldots,x_n,t]/(h_1,\ldots,h_m,tG-1).
$$
The map sends $x_i$ to its original class and $t$
to $1/g$. The inverse sends $s/g^d$ to the class
of $P(x)t^d$ for a polynomial lift $P$ of $s$.
Independence follows from the original relations
$h_j$ and the relation $tG=1$; these maps are
inverse on the algebra generators and fractions.
Hence Chevalley's Theorem \ref{editorial-source-theorem-chevalley}
applies to this particular $R\to S_g$ and shows
that the image $E_g$ of $D(g)$ is constructible.

Going down shows that $E_g$ is stable under
generalization. Given $\mathfrak q'\in D(g)$ over
$\mathfrak p'$ and $\mathfrak p\subseteq\mathfrak p'$
in $R$, take $\mathfrak q\subseteq\mathfrak q'$
over $\mathfrak p$. Since $g\notin\mathfrak q'$,
also $g\notin\mathfrak q$, and the lift remains
in the original $D(g)$. This is also the exact
going-down comparison for $S\to S_g$ and the
composition lemma. The preceding constructible-set
lemma now makes $E_g$ open. Every original open
of $\Spec(S)$ is a union of such $D(g)$ and images
commute with unions, proving openness of the map.

The proof only needs constructibility of the
images of all principal opens. Under that weaker
assumption, going down is equivalent to openness:
the proof just given gives one direction, and
Lemma \ref{editorial-source-lemma-open-going-down} gives the other.
In particular this is an equivalence for maps of
finite presentation. Finite presentation is not
necessary: for a field $k$, the map
$k\to k[x_i:i\geq1]$ has open spectrum map,
because the target spectrum $\Spec(k)$ has one
point. The algebra is not even of finite type:
finitely many polynomials use finitely many
variables and cannot generate an omitted original
indeterminate, as evaluation fixing those variables
and changing that indeterminate verifies.

For the flat finite-presentation case the map is
universally open. Indeed for every original
$R$-algebra $A$, the base change has presentation
$$
S\otimes_RA\simeq
A[x_1,\ldots,x_n]/(h_1^A,\ldots,h_m^A),
$$
where $h_j^A$ retains every coefficient through
the original map $R\to A$. The map sends
$s\otimes a$ to $a$ times a polynomial representative
of $s$; the inverse sends the original variables
to $x_i\otimes1$ and each $a$ to $1\otimes a$.
The displayed relations and tensor balancing make
these maps inverse. Flatness is preserved because
for an injection of $A$-modules $L\hookrightarrow L'$,
the original comparison
$$
L\otimes_A(S\otimes_RA)\simeq L\otimes_RS,
\qquad l\otimes(s\otimes a)\longmapsto al\otimes s,
$$
with inverse $l\otimes s\mapsto l\otimes(s\otimes1)$,
identifies its tensor map with an injection by
flatness of $S$ over $R$. Thus the proved proposition
applies after every base change. This last assertion
uses flatness, not merely the original going-down
hypothesis.
\end{proof}


\subsection{Editorial treatment of Algebra lines 9739--9772}

\hypertarget{algebra-context-112}{}

\noindent\href{https://github.com/stacks/stacks-project/blob/a04446e57ec1fbc252a871afcec7752fb2807b14/algebra.tex\#L9739}{Original source, lines 9739--9772}. Complete original and reviewed TeX passages accompany this note. Every intervention is separately recorded; the following treatment is editorial.


\begin{lemma}
\label{editorial-source-lemma-same-image}
Let $k$ be a field, and let $R$, $S$ be $k$-algebras.
Let $S' \subset S$ be a sub $k$-algebra, and let $f \in S' \otimes_k R$.
In the commutative diagram
$$
\xymatrix{
\Spec((S \otimes_k R)_f) \ar[rd] \ar[rr] & &
\Spec((S' \otimes_k R)_f) \ar[ld] \\
& \Spec(R) &
}
$$
the images of the diagonal arrows are the same.
\end{lemma}

\begin{proof}
Retain the original inclusion $i:S'\hookrightarrow S$,
the original $k$-algebra $R$, and the specified
$f\in S'\otimes_kR$. For a prime $\mathfrak p\subset R$
put $K=\kappa(\mathfrak p)$. The coefficient maps
$R\to K$ give elements $f_{S'}\in S'\otimes_kK$
and $f_S\in S\otimes_kK$, related by $i\otimes1$.
For either $B=S'$ or $B=S$, the exact comparison is
$$
((B\otimes_kR)_f)\otimes_RK
\longrightarrow (B\otimes_kK)_{f_B},
\qquad
\frac{\sum_j b_j\otimes r_j}{f^d}\otimes\lambda
\longmapsto
\frac{\sum_j b_j\otimes\overline{r_j}\lambda}{f_B^d}.
$$
Here the $f$ in the $B=S$ expression denotes its
image under the original inclusion, not a different
choice of element. The inverse sends
$$
\frac{b\otimes\lambda}{f_B^d}
\longmapsto \frac{b\otimes1}{f^d}\otimes\lambda
$$
and extends additively. To check these maps, first use
the balanced isomorphism
$(B\otimes_kR)\otimes_RK\simeq B\otimes_kK$,
$(b\otimes r)\otimes\lambda\mapsto b\otimes\overline r\lambda$,
whose inverse is $b\otimes\lambda\mapsto(b\otimes1)\otimes\lambda$.
On localizing, $f$ maps to $f_B$ and is inverted in
both targets. The universal property of localization
therefore extends the two inverse algebra maps to
the displayed fraction maps. In particular it checks
all zero-divisor witnesses as well as tensor balancing;
their composites are identities on the original
algebra generators and on $1/f$.

For a ring $C$ and $h\in C$, one has $C_h=0$
exactly when $h$ is nilpotent. Indeed $1/1=0$
holds exactly when $h^d\cdot1=0$ for some
$d\geq0$; in the case $d=0$ the ring is zero
and $h^1=0$ also holds. Conversely any such positive
power makes $h$ an invertible nilpotent, forcing the
localized unit to vanish. Thus the fibre criterion
of Lemma \ref{editorial-source-lemma-in-image} identifies the image
of each diagonal arrow with the primes for which
its original $f_B$ is not nilpotent.

Since $k$ is a field, tensoring the injection
$S'\hookrightarrow S$ with the $k$-vector space $K$
is injective. Explicitly, choose a $k$-basis of $K$.
Under the direct-sum identifications of both tensor
modules, the map is the direct sum of copies of the
original injection; hence its kernel is zero.
Consequently
$$
f_{S'}^d=0\quad\Longleftrightarrow\quad f_S^d=0
\qquad(d\geq1).
$$
The forward implication is functoriality; the reverse
uses this injection. Nonnilpotence is therefore the
same in both original fibres. The fibre criterion
and the fraction comparisons prove equality of the
two spectrum images. Independently, one inclusion
is also the factorization through the horizontal
arrow: every prime of the larger tensor algebra
contracts to a prime of the smaller one, with the
same contraction to $R$. The reverse inclusion is
the nonnilpotence argument just supplied.

\medskip\noindent
This proof gives a more general assertion with its
exact sufficient hypothesis. Let $A$ be any base
ring, let $B\to C$ be a map of $A$-algebras, and
let $R$ be an $A$-algebra. Suppose that for every
original $\mathfrak p\in\Spec(R)$ the map
$$
B\otimes_A\kappa(\mathfrak p)
\longrightarrow C\otimes_A\kappa(\mathfrak p)
$$
is injective. Then for every $f\in B\otimes_AR$,
the two original principal-localization spectrum
images in $\Spec(R)$ agree. All the preceding
fraction maps work with $A$ in place of $k$,
and injectivity of these particular maps supplies
exactly the nilpotence reflection used in the proof.
In particular the hypothesis holds when the
$A$-module map $B\to C$ is pure, meaning that
$B\otimes_AL\to C\otimes_AL$ is injective for
every $A$-module $L$. It suffices, more weakly,
that each of the displayed fibre maps reflect
nilpotence of the particular image of $f$.

An arbitrary injection over a general base ring
does not suffice. Take $A=R=B=\mathbf Z$,
$C=\mathbf Q$, and $f=1$. The spectrum image
for $B$ is all of $\Spec(\mathbf Z)$, whereas
the image for $C$ is $\{(0)\}$. At a nonzero
prime $(p)$ the receiving map of fibres is
$\mathbf F_p\to0$, since
$q\otimes\overline a=(q/p)\otimes p\overline a=0$.
This identifies the failed residue injection and
the failed image equality for the same original
maps and element. Finally, nonzero elements cannot
replace nonnilpotent elements in the proof: over
$k[\epsilon]/(\epsilon^2)$ the nonzero element
$\epsilon$ has empty principal open. Its square
is zero and its localization is the zero ring.
\end{proof}


\subsection{Editorial treatment of Algebra lines 9774--9791}

\hypertarget{algebra-context-113}{}

\noindent\href{https://github.com/stacks/stacks-project/blob/a04446e57ec1fbc252a871afcec7752fb2807b14/algebra.tex\#L9774}{Original source, lines 9774--9791}. Complete original and reviewed TeX passages accompany this note. Every intervention is separately recorded; the following treatment is editorial.


\begin{lemma}
\label{editorial-source-lemma-map-into-tensor-algebra-open}
Let $k$ be a field.
Let $R$ and $S$ be $k$-algebras.
The map $\Spec(S \otimes_k R) \to \Spec(R)$
is open.
\end{lemma}

\begin{proof}
Retain the original $k$-algebras $S,R$ and write the
specified tensor as a finite sum
$$
f=\sum_{i=1}^d s_i\otimes r_i\in S\otimes_kR.
$$
Set $S'=k[s_1,\ldots,s_d]\subseteq S$, keeping every
original coefficient $s_i$. Tensoring $S'\hookrightarrow S$
with $R$ is injective because $R$ is a $k$-vector space,
so the same displayed tensor is an element of
$S'\otimes_kR$ with its specified image $f$.
The $k$-algebra $S'$ is finitely presented: it is
a quotient of $k[x_1,\ldots,x_d]$, and this polynomial
ring is Noetherian by Lemma \ref{editorial-source-lemma-Noetherian-permanence}.
Hence the kernel of the original map $x_i\mapsto s_i$
has finitely many generators $h_1,\ldots,h_m$.
This also follows from
Lemma \ref{editorial-source-lemma-Noetherian-finite-type-is-finite-presentation}.

The base change retains the presentation
$$
S'\otimes_kR\simeq R[x_1,\ldots,x_d]/(h_1^R,\ldots,h_m^R),
$$
where every coefficient of $h_j^R$ is the image of
the corresponding coefficient of $h_j$ in $R$.
The map sends $s_i\otimes1$ to the class of $x_i$
and $1\otimes r$ to $r$, and its inverse sends the
original polynomial generators and coefficients back
to these tensors. The relations and balancing prove
both maps are well defined and inverse.
The algebra is flat over $R$. To see this directly,
choose a $k$-basis $(b_\lambda)$ of $S'$.
The map
$$
\bigoplus_\lambda R\longrightarrow S'\otimes_kR,
\qquad (a_\lambda)_\lambda\longmapsto
\sum_\lambda b_\lambda\otimes a_\lambda
$$
is an $R$-module isomorphism; its inverse is the
finite basis expansion in the first tensor factor.
Tensoring a free module preserves injections by
the corresponding direct-sum maps. Thus
Proposition \ref{editorial-source-proposition-fppf-open} applies to
$R\to S'\otimes_kR$ and makes the image of the
original principal open $D(f)$ in this spectrum
open. Lemma \ref{editorial-source-lemma-same-image}, with the exact
same $f$, identifies it with the image of $D(f)$
in $\Spec(S\otimes_kR)$. Principal opens form a
basis and images commute with unions, so the
original spectrum map is open, as claimed.

This also proves universal openness of the
original map. For an arbitrary original $R$-algebra
$A$, its base change is identified with
$\Spec(S\otimes_kA)\to\Spec(A)$ through the map
$$
(S\otimes_kR)\otimes_RA\longrightarrow S\otimes_kA,
\qquad (s\otimes r)\otimes a\longmapsto s\otimes\overline r a,
$$
whose inverse is $s\otimes a\mapsto(s\otimes1)\otimes a$.
Both maps respect the original algebra products
and balancing, and their composites fix all
elementary tensors. The result already proved
applies to this $k$-algebra $A$, giving openness
after every base change, without finite generation
of the full algebra $S$.

\medskip\noindent
One can make the principal-open image more explicit
without confusing nonzero and nonnilpotent elements.
For the fixed original list $s_1,\ldots,s_d$, there
exists an integer $N\geq1$, independent of $R$
and $r_1,\ldots,r_d$, such that the image of $D(f)$
is the union of the coefficient opens of the full
power $f^N$ in a $k$-basis of $S'$.
Here is a proof with all coefficients retained.

Put $A_0=k[T_1,\ldots,T_d]$ and
$F=\sum_{i=1}^d s_i\otimes T_i\in S'\otimes_kA_0$.
Choose a $k$-basis $(b_\lambda)_{\lambda\in\Lambda}$
of $S'$. For each ordered tuple $(i_1,\ldots,i_n)$,
write the complete finite expansion
$$
s_{i_1}\cdots s_{i_n}
=\sum_{\lambda\in\Lambda}
\mu_{i_1,\ldots,i_n,\lambda}b_\lambda,
\qquad \mu_{i_1,\ldots,i_n,\lambda}\in k.
$$
Then, with all ordered tuples and their repeated
terms retained, the coefficients of $F^n$ are
$$
\begin{split}
F^n&=\sum_{\lambda\in\Lambda}b_\lambda\otimes C_{n,\lambda}(T),\\
C_{n,\lambda}(T)
&=\sum_{i_1=1}^d\cdots\sum_{i_n=1}^d
\mu_{i_1,\ldots,i_n,\lambda}T_{i_1}\cdots T_{i_n}.
\end{split}
$$
For fixed $n$ only finitely many of these polynomials
are nonzero. Put $J_n=(C_{n,\lambda}:\lambda\in\Lambda)
\subseteq A_0$ and $Z_n=V(J_n)$. Extension of the
basis to every residue field shows that
$\mathfrak p\in Z_n$ exactly when the original
$F$ has $n$-th power zero in
$S'\otimes_k\kappa(\mathfrak p)$. Hence
$Z_n\subseteq Z_{n+1}$. The fibre criterion and
the tensor-localization maps of the preceding
lemma give
$$
\Spec(A_0)\setminus
\operatorname{image}\bigl(D(F)\to\Spec(A_0)\bigr)
=\bigcup_{n\geq1}Z_n.
$$
Call this complement $Z$. The proven openness
result, applied to the original algebra $S'$
and base $A_0$, says that $Z$ is closed.

The ring $A_0$ is Noetherian, so its closed set
$Z$ has finitely many irreducible components
$V(\mathfrak p_1),\ldots,V(\mathfrak p_e)$ with
their indicated generic points. This follows
by applying Lemma \ref{editorial-source-lemma-Noetherian-irreducible-components}
to the quotient defining $Z$. Each generic point
lies in some $Z_{n_j}$, because their union is
$Z$. Since $Z_{n_j}$ is closed, it contains
the full component $V(\mathfrak p_j)$.
Thus for $N=\max(1,n_1,\ldots,n_e)$,
$$
Z=Z_N.
$$
If $Z$ is empty, use $N=1$: every $Z_n$ is
already empty. This is a finite-component argument;
it does not infer stabilization of increasing
closed sets merely from Noetherianity.

For any original $R$ and $r_i$, use the original
map $A_0\to R$, $T_i\mapsto r_i$.
At a prime $\mathfrak p\subset R$, its contraction
$\mathfrak p_0\subset A_0$ induces an injective
field map $\kappa(\mathfrak p_0)\to\kappa(\mathfrak p)$.
The map
$S'\otimes_k\kappa(\mathfrak p_0)\to
S'\otimes_k\kappa(\mathfrak p)$ is injective by
extension of vector spaces, and carries the full
original $F$ to $f$. Therefore nilpotence is
equivalent in these two fibres, and by $Z=Z_N$
it is equivalent to the vanishing of the $N$-th
power. Its coefficients are exactly
$C_{N,\lambda}(r_1,\ldots,r_d)$, with the same
ordered-tuple formula. It follows that
$$
\operatorname{image}\bigl(D(f)\subseteq
\Spec(S\otimes_kR)\to\Spec(R)\bigr)
=\bigcup_{\lambda\in\Lambda}
D\bigl(C_{N,\lambda}(r_1,\ldots,r_d)\bigr).
$$
Only finitely many original coefficient polynomials
contribute, giving a finite union of basic opens
and hence a quasi-compact open image. The equality
for $S$ uses the same-image lemma after the equality
for its specified subalgebra $S'$. None of these
comparisons replaces $f$ by its linear nonzero
coefficient test: the required test uses the
proved power $N$. For instance in
$S'=k[\epsilon]/(\epsilon^2)$ with $f=\epsilon\otimes1$,
the first power is nonzero and has a unit coefficient,
while the actual principal open is empty; the
second power vanishes, giving the correct empty
union. The case $d=0$ has $F=f=0$ and can use
$N=1$, and zero tensor algebras give empty images
under the same formulas.
\end{proof}


\subsection{Editorial treatment of Algebra lines 9796--9845}

\hypertarget{algebra-context-114}{}

\noindent\href{https://github.com/stacks/stacks-project/blob/a04446e57ec1fbc252a871afcec7752fb2807b14/algebra.tex\#L9796}{Original source, lines 9796--9845}. Complete original and reviewed TeX passages accompany this note. Every intervention is separately recorded; the following treatment is editorial.


\begin{lemma}
\label{editorial-source-lemma-unique-prime-over-localize-below}
Let $R \to S$ be a ring map.
Let $\mathfrak p \subset R$ be a prime.
Assume that
\begin{enumerate}
\item there exists a unique prime $\mathfrak q \subset S$ lying over
$\mathfrak p$, and
\item either
\begin{enumerate}
\item going up holds for $R \to S$, or
\item going down holds for $R \to S$ and there is at most one prime
of $S$ above every prime of $R$.
\end{enumerate}
\end{enumerate}
Then $S_{\mathfrak p} = S_{\mathfrak q}$.
\end{lemma}

\begin{proof}
Write the original ring map as $\varphi:R\to S$.
The rings in the conclusion are the specified
localizations
$$
S_{\mathfrak p}=\varphi(R\setminus\mathfrak p)^{-1}S,
\qquad S_{\mathfrak q}=(S\setminus\mathfrak q)^{-1}S.
$$
Since $\varphi^{-1}(\mathfrak q)=\mathfrak p$,
the first multiplicative set is contained in the
second. Thus there is an original localization map
$j:S_{\mathfrak p}\to S_{\mathfrak q}$, sending
$s/\varphi(r)$ to the same fraction.
No injectivity of $\varphi$ or of these localization
maps is assumed.

A prime of $S_{\mathfrak p}$ corresponds exactly
to a prime $\mathfrak q'\subset S$ whose contraction
$\mathfrak p'=\varphi^{-1}(\mathfrak q')$ is contained
in $\mathfrak p$. A prime of $S_{\mathfrak q}$
corresponds exactly to a prime
$\mathfrak q'\subseteq\mathfrak q$. Both statements
follow from contraction and extension under
localization: the corresponding prime consists of
the fractions with numerator in $\mathfrak q'$,
and these maps are inverse because the denominator
set is disjoint from that prime.

Assume first (1) and (2)(a). For such a
$\mathfrak q'$ over $\mathfrak p'\subseteq\mathfrak p$,
going up supplies $\mathfrak q''\supseteq\mathfrak q'$
whose contraction is $\mathfrak p$. Uniqueness over
$\mathfrak p$ makes $\mathfrak q''=\mathfrak q$,
so $\mathfrak q'\subseteq\mathfrak q$.
Assume instead (1) and (2)(b). Going down applied
to $\mathfrak q$ gives
$\mathfrak q''\subseteq\mathfrak q$ over
$\mathfrak p'$. Uniqueness over $\mathfrak p'$
makes $\mathfrak q''=\mathfrak q'$, with the same
conclusion. This argument uses uniqueness only
over the primes $\mathfrak p'\subseteq\mathfrak p$;
the global uniqueness assumption in (2)(b) can
therefore be weakened to precisely those primes.

It follows in either case that for every
$u\in S\setminus\mathfrak q$, the element $u/1$
belongs to no prime of $S_{\mathfrak p}$: every
such prime contracts to a prime contained in
$\mathfrak q$. Hence $u/1$ is a unit. Indeed a
nonunit generates a proper ideal lying in a
maximal ideal, which is prime, a contradiction.
Write its inverse as $v/\varphi(r)$ with
$r\notin\mathfrak p$. The equality
$(u/1)(v/\varphi(r))=1$ means that for some
$t\notin\mathfrak p$,
$$
\varphi(t)(uv-\varphi(r))=0\quad\hbox{in }S.
$$
Thus the inverse to $j$ has the full fraction
formula
$$
k:S_{\mathfrak q}\longrightarrow S_{\mathfrak p},
\qquad
\frac su\longmapsto\frac{sv}{\varphi(r)},
$$
where $v,r$ satisfy that original inverse identity.
Independence of this choice follows from uniqueness
of the inverse of $u/1$. Independence of fraction
representatives can also be checked explicitly:
if $w(u's-us')=0$ for $w\notin\mathfrak q$,
then $w/1,u/1,u'/1$ are all units in
$S_{\mathfrak p}$, so multiplication by their
inverses gives $s/u=s'/u'$ there. Addition and
multiplication follow by the same common-denominator
identities. Consequently $k$ is the localization
extension of the original map $S\to S_{\mathfrak p}$.
Both composites fix $S$ and all the specified inverse
denominators, hence are identities. This proves
$S_{\mathfrak p}\simeq S_{\mathfrak q}$ through
the original canonical map, giving the equality
notation in the statement its exact meaning.

In fact for any original $\varphi$, $\mathfrak p$
and $\mathfrak q$ with
$\varphi^{-1}(\mathfrak q)=\mathfrak p$, the
following conditions are equivalent:
$$
\begin{gathered}
j:S_{\mathfrak p}\longrightarrow S_{\mathfrak q}
\text{ is an isomorphism};\\
u/1\text{ is a unit in }S_{\mathfrak p}
\text{ for every }u\notin\mathfrak q;\\
\varphi^{-1}(\mathfrak q')\subseteq\mathfrak p
\Longrightarrow\mathfrak q'\subseteq\mathfrak q
\text{ for every prime }\mathfrak q'\subset S.
\end{gathered}
$$
The first implies the second because these elements
are units in $S_{\mathfrak q}$. The second implies
the third since a prime of $S_{\mathfrak p}$ cannot
contain a unit; contraction then avoids every
element outside $\mathfrak q$. The third implies
the second by the maximal-ideal argument, and the
second supplies the proved inverse fraction map.
This is the exact criterion used by both cases
of the original lemma. It is specific to the
specified further localization; a bijective
spectrum map alone has not been substituted for
an isomorphism of rings.
\end{proof}


\subsection{Editorial treatment of Algebra lines 9851--9882}

\hypertarget{algebra-context-115}{}

\noindent\href{https://github.com/stacks/stacks-project/blob/a04446e57ec1fbc252a871afcec7752fb2807b14/algebra.tex\#L9851}{Original source, lines 9851--9882}. Complete original and reviewed TeX passages accompany this note. Every intervention is separately recorded; the following treatment is editorial.


\begin{lemma}
\label{editorial-source-lemma-going-down-flat-module}
Let $R \to S$ be a ring map. Let $N$ be a finite $S$-module flat over $R$.
Endow $\text{Supp}(N) \subset \Spec(S)$ with the induced topology.
Then generalizations lift along $\text{Supp}(N) \to \Spec(R)$.
\end{lemma}

\begin{proof}
Write the original map as $\varphi:R\to S$. The
statement means the following. Given original primes
$\mathfrak p\subseteq\mathfrak p'\subset R$ and
$\mathfrak q'\in\operatorname{Supp}_S(N)$ with
$\varphi^{-1}(\mathfrak q')=\mathfrak p'$, there
must be a prime $\mathfrak q\subseteq\mathfrak q'$
such that $\varphi^{-1}(\mathfrak q)=\mathfrak p$
and $\mathfrak q\in\operatorname{Supp}_S(N)$.
The contraction condition on $\mathfrak q'$ is
part of the lifting statement and is retained
explicitly here. For the induced topology on
support, specialization between two of its points
is the original prime inclusion: closure in a
subspace is the intersection with the original
closure. Thus these prime conditions are exactly
the required generalization lift.

Put
$$
A=R_{\mathfrak p'},\qquad C=S_{\mathfrak q'},
\qquad E=N_{\mathfrak q'},\qquad
\mathfrak m=\mathfrak p'R_{\mathfrak p'}.
$$
Keep the local ring map
$A\to C$, $r/t\mapsto\varphi(r)/\varphi(t)$.
Its denominators are valid because
$t\notin\mathfrak p'$ implies
$\varphi(t)\notin\mathfrak q'$.
The module $E$ is flat over $A$, by
Lemma \ref{editorial-source-lemma-flat-localization}. One may also
see the original comparison directly: tensor an
injection of $A$-modules with the $R$-flat $N$
and then localize the resulting $S$-modules at
$\mathfrak q'$. Exactness of localization preserves
the injection, and the maps
$$
(L\otimes_RN)_{\mathfrak q'}\longrightarrow L\otimes_RE,
\qquad (l\otimes n)/s\longmapsto l\otimes n/s
$$
and their inverses identify the result with
tensoring by $E$. Here every original element of
$R\setminus\mathfrak p'$ is invertible on both
$L$ and $E$, so balancing already factors uniquely
through $A$. This gives the same $A$-flatness.

The module $E$ is nonzero by the chosen support
point, and is finite over the local ring $C$.
Nakayama's Lemma \ref{algebra-lemma-NAK} gives
$E/\mathfrak q' E\ne0$. Since the image of
$\mathfrak m$ lies in $\mathfrak q'C$, the quotient
map
$$
E/\mathfrak mE\longrightarrow E/\mathfrak q'E,
\qquad e+\mathfrak mE\longmapsto e+\mathfrak q'E
$$
is surjective. Consequently
$E\otimes_A\kappa(\mathfrak p')=E/\mathfrak mE\ne0$.
The quotient-tensor identification sends
$e\otimes\overline a$ to $ae+\mathfrak mE$,
with inverse $e+\mathfrak mE\mapsto e\otimes1$.

The flat $A$-module $E$ is therefore faithfully
flat. For a direct verification, if an arbitrary
$A$-module $L$ has an element $l\ne0$, set
$H=\operatorname{Ann}_A(l)\subseteq\mathfrak m$.
The injection $A/H\hookrightarrow L$,
$a+H\mapsto al$, becomes the injection
$E/HE\hookrightarrow L\otimes_AE$ by flatness.
The source surjects onto $E/\mathfrak mE\ne0$,
so $L\otimes_AE\ne0$. This proves nonzero-module
detection and, together with flatness, faithful
flatness, in agreement with Lemma \ref{editorial-source-lemma-ff}.
In particular for the nonzero $A$-module
$K=\kappa(\mathfrak p)$, defined using the original
chain $\mathfrak p\subseteq\mathfrak p'$, one has
$$
P=E\otimes_AK\ne0,
\qquad B=C\otimes_AK.
$$

Here are the full localization and quotient maps
in this fibre. Let $W=R\setminus\mathfrak p$,
$D=W^{-1}C$ and $Q=W^{-1}E$, using the images of
$W$ under the original map. Then
$$
B\simeq D/\mathfrak pD,\qquad
P\simeq Q/\mathfrak pQ.
$$
The ring comparison sends
$$
c\otimes\frac{a+\mathfrak p}{w+\mathfrak p}
\longmapsto\frac{\varphi(a)c}{\varphi(w)}+\mathfrak pD,
$$
where coefficients in $C$ are interpreted via its
original map from $R$. Its inverse sends
$c/\varphi(w)+\mathfrak pD$ to
$c\otimes(1/(w+\mathfrak p))$.
The module comparison uses the same formulas
with $e\in E$ in place of $c$. The original
relations of localization are respected because
each image of $w\in W$ becomes a unit in $K$.
The original ideal $\mathfrak p$ is killed in
both targets. Thus the universal properties of
localization and quotient, followed by tensor
balancing over $A$, give these maps. Their
composites fix every original coefficient,
element of $E$, and inverse denominator, proving
the assertions. These expressions retain the
possibly noninjective maps and all denominators.

Since $N$ is finite over $S$, the localized
$E$ is finite over $C$, and its original finite
generators tensor to generators of $P$ over $B$.
Put $J=\operatorname{Ann}_B(P)$, the original
annihilator ideal in the source proof.
Because $P\ne0$, $J$ is proper. Choose a prime
$\mathfrak b$ of $B$ containing $J$.
The finite-module support equality of
Lemma \ref{editorial-source-lemma-support-closed} gives
$P_{\mathfrak b}\ne0$.
Equivalently this prime is obtained from a prime
of the original quotient $B/J$ by inverse image.
Write
$$
\theta:S\longrightarrow C\longrightarrow B,
\qquad s\longmapsto (s/1)\otimes1,
\qquad \mathfrak q=\theta^{-1}(\mathfrak b).
$$
It is prime. Every $s\notin\mathfrak q'$ is a
unit in $C$ and hence in $B$, so
$\mathfrak q\subseteq\mathfrak q'$.
Every $r\in\mathfrak p$ maps to zero, while
every $r\notin\mathfrak p$ maps to the nonzero
scalar $r+\mathfrak p$ in the field $K$, which
is a unit in the nonzero ring $B$. The map
$K\to B$ is unital and injective, since a
nonzero field has no proper nonzero kernel.
Thus
$\varphi^{-1}(\mathfrak q)=\mathfrak p$ exactly.

To prove support membership through the original
module, use the balanced isomorphism
$$
N\otimes_SB\longrightarrow P,
\qquad n\otimes(c\otimes a)\longmapsto(cn/1)\otimes a,
$$
with inverse $(n/s)\otimes a\mapsto
n\otimes((1/s)\otimes a)$ for $s\notin\mathfrak q'$.
All such $s$ are inverted in $C$, so localization
and tensor balancing make these maps independent
of representatives; both composites fix the
original elements and scalar tensors. Localizing
at $\mathfrak b$ gives the further comparison
$$
P_{\mathfrak b}\simeq
N_{\mathfrak q}\otimes_{S_{\mathfrak q}}B_{\mathfrak b}.
$$
The map on an elementary fraction sends
$(n\otimes b)/u$ to $n/1\otimes b/u$.
Its inverse sends $n/s\otimes b/u$ to
$(n\otimes b)/(\theta(s)u)$, where
$s\notin\mathfrak q$ implies
$\theta(s)\notin\mathfrak b$. These formulas
respect both localization witness relations and
the scalar balancing, and are inverse on the
generators. Since the left side is nonzero,
$N_{\mathfrak q}$ is nonzero. This proves the
required lift in the original support.

\medskip\noindent
The exact places where finiteness was used give
a weaker sufficient hypothesis. Suppose now that
$N$ is any $S$-module flat over $R$ and that for
every original $\mathfrak q'\in\operatorname{Supp}_S(N)$,
with $\mathfrak p'=\varphi^{-1}(\mathfrak q')$, one has
$$
N_{\mathfrak q'}\otimes_{R_{\mathfrak p'}}
\kappa(\mathfrak p')\ne0.
$$
The argument above still proves $A$-flatness of
$E$; the displayed hypothesis replaces the
Nakayama step, and gives its faithful flatness.
Hence the same original fibre $P$ is nonzero
for every $\mathfrak p\subseteq\mathfrak p'$.
It need not be finite. Choose an original
$z\ne0$ in $P$, and a maximal ideal
$\mathfrak b\supseteq\operatorname{Ann}_B(z)$.
Then $z/1\ne0$ in $P_{\mathfrak b}$: a witness
$u\notin\mathfrak b$ with $uz=0$ would contradict
the containment of its annihilator. The same
prime contractions and comparison maps now give
the lift. This replaces the finite-module use
of $\operatorname{Ann}_B(P)$ by the exact cyclic
support of the chosen original element.

In particular it suffices that the localization
support equal the residue-fibre set over $S$:
$$
\operatorname{Supp}_S(N)
=\{\mathfrak q:N\otimes_S\kappa(\mathfrak q)\ne0\}.
$$
At a support point this supplies
$E/\mathfrak q'E\ne0$, and the same quotient
surjection supplies the required nonzero fibre
over $R_{\mathfrak p'}$. Every arbitrary direct
sum of finite $S$-modules satisfies this equality,
as proved in the support base-change comparison;
when that sum is also $R$-flat, the lifting
conclusion therefore follows without finite
generation of $N$. This is a sufficient hypothesis,
not a claim that such fibre nonvanishing is
necessary for every possible support lift.

\medskip\noindent
There is also an openness consequence. Suppose the
original algebra $S$ is finitely presented over $R$,
and the original $S$-module $N$ is finitely presented
and $R$-flat. Then the map on its support is open,
and the corresponding support map remains open
after every base change $R\to A_1$.
To prove this with the specified objects, retain
a finite presentation $S^b\xrightarrow{H}S^a\to N\to0$.
Let $J_H$ be the ideal generated by the original
$a\times a$ minors of $H$, using determinant $1$
in size zero. The full presentation calculation in
Lemma \ref{editorial-source-lemma-support-finite-presentation-constructible}
gives $\operatorname{Supp}_S(N)=V(J_H)$, with the
quotient-spectrum homeomorphism
$\Spec(S/J_H)\to V(J_H)$ given by contraction.
The ideal $J_H$ is finitely generated. Appending
polynomial representatives of its finitely many
specified generators to a finite presentation of
$S$ over $R$ proves that $S/J_H$ is finitely
presented over $R$. The going-down result just
proved lifts every prime generalization inside
$V(J_H)$; under the quotient homeomorphism this
is exactly going down for $R\to S/J_H$.
Proposition \ref{editorial-source-proposition-fppf-open} therefore
proves openness of the original support map.

For the arbitrary original base change $R\to A_1$,
put $S_1=S\otimes_RA_1$ and $N_1=N\otimes_RA_1$.
Tensoring the displayed presentation and its
coefficient maps gives
$S_1^b\xrightarrow{H_1}S_1^a\to N_1\to0$,
where every entry of $H_1$ is the original entry
of $H$ tensored with $1$. The full determinant
polynomial, with all its signs, shows that its
minor ideal is $J_HS_1$. Thus
$$
\operatorname{Supp}_{S_1}(N_1)=V(J_HS_1),
\qquad
(S/J_H)\otimes_RA_1\simeq S_1/J_HS_1.
$$
The last map sends $(s+J_H)\otimes a$ to
$s\otimes a+J_HS_1$; its inverse is the
original quotient map followed by tensoring.
The ideal generators are killed on both sides,
so these are inverse maps with the indicated
formulas. The algebra $S_1$ is finitely presented
over $A_1$ by the base-change presentation.
The module $N_1$ is $A_1$-flat: for an injection
of $A_1$-modules $L\hookrightarrow L'$, the maps
$$
L\otimes_{A_1}N_1\simeq L\otimes_RN,
\qquad l\otimes(n\otimes a)\longmapsto al\otimes n,
$$
with inverse $l\otimes n\mapsto l\otimes(n\otimes1)$,
identify the tensor map with an injection by
the original $R$-flatness of $N$.
Applying the proven support argument to these
original base-change objects proves openness
for every $A_1$. Equivalently, the quotient map
$\Spec(S/J_H)\to\Spec(R)$ is universally open,
even though the argument did not require
$S/J_H$ itself to be flat over $R$.

Flatness alone, and even flatness together with
finite annihilator detection, does not suffice.
Retain the original ring and module
$$
R=S=k[x,y]/(xy),\qquad \varphi=\operatorname{id}_R,
\qquad N=R_x,
$$
where $k$ is a field and $x,y$ denote their
specified quotient classes. The module is flat
by exactness of localization. Its annihilator
is exactly $(y)$, and is detected by $1\in R_x$.
Indeed $y$ kills the entire localization since
$xy=0$ and $x$ is inverted. Conversely, if
$r\cdot1=0$ in $R_x$, then $x^nr=0$ for some
$n\geq0$. Modulo $(y)$ this says
$x^n\overline r=0$ in the domain $k[x]$,
so $r\in(y)$. The same condition is equivalent
to $r$ killing every element of $R_x$ because
all its elements are fractions with powers
of $x$ as denominators.

The exact comparison
$$
R_x\longrightarrow k[x,x^{-1}],\qquad
r/x^n\longmapsto r(x,0)/x^n,
$$
is an isomorphism, with inverse given by the
original class of $x$ and its inverse; $y$ is
zero in $R_x$ by $xy=0$. Thus as an $R$-module,
$N$ has support $V(y)$. To verify both directions,
if $y\notin\mathfrak a$, its unit image in
$R_{\mathfrak a}$ kills $N_{\mathfrak a}$,
so that localization is zero. If
$y\in\mathfrak a$, the prime
$\overline{\mathfrak a}=\mathfrak a/(y)$
in $k[x]$ identifies $N_{\mathfrak a}$ with
the localization of $k[x,x^{-1}]$ at the
images of $k[x]\setminus\overline{\mathfrak a}$.
These images are all nonzero in a domain,
so the localized ring is nonzero. The comparison
sends each original fraction to the same
fraction after $y=0$; the inverse follows from
the two specified localization maps.
Its fibre set, however, is exactly $D(x)$:
$$
N\otimes_R\kappa(\mathfrak a)
\simeq\kappa(\mathfrak a)_{\overline x}.
$$
This is zero if $x\in\mathfrak a$ and is the
nonzero field otherwise, with fraction maps
$r/x^n\otimes c\mapsto\overline r c/\overline x^n$
and the tensor inverse. For primes avoiding
$x$, the relation $xy=0$ forces $y$ into the
prime, consistently with $D(x)\subseteq V(y)$.

Take the original primes
$\mathfrak p=(x)\subset\mathfrak p'=(x,y)$
and $\mathfrak q'=\mathfrak p'$. The latter
is in $V(y)=\operatorname{Supp}(N)$, but the
former is not, because $y\notin(x)$.
Since the spectrum map for $\varphi$ is the
identity, the only possible lift over
$\mathfrak p$ is $\mathfrak p$ itself, which
is outside support. This proves the failure
of generalization lifting despite flatness
and finite annihilator detection. At
$\mathfrak q'$ the required base residue fibre
is zero, exactly identifying the failed
hypothesis in the extended proof.
\end{proof}


\subsection{Editorial treatment of Algebra lines 9933--9942}

\hypertarget{algebra-context-116}{}

\noindent\href{https://github.com/stacks/stacks-project/blob/a04446e57ec1fbc252a871afcec7752fb2807b14/algebra.tex\#L9933}{Original source, lines 9933--9942}. Complete original and reviewed TeX passages accompany this note. Every intervention is separately recorded; the following treatment is editorial.


\begin{lemma}
\label{editorial-source-lemma-subextensions-are-separable}
Let $K/k$ be a separable field extension.
For any subextension $K/K'/k$ the field
extension $K'/k$ is separable.
\end{lemma}

\begin{proof}
Let $k\subseteq L\subseteq K'$ be any original
intermediate field finitely generated over $k$.
Then $L$ is also an intermediate field of the
original $K/k$. Separability of $K/k$, with
exactly the given definition, supplies a
transcendence basis of $L/k$ over which $L$
is separable algebraic. Thus every such $L$
is separably generated over $k$, which is
precisely separability of $K'/k$. No finite
generation of $K'$ or $K$ is required.
\end{proof}


\subsection{Editorial treatment of Algebra lines 9944--9959}

\hypertarget{algebra-context-117}{}

\noindent\href{https://github.com/stacks/stacks-project/blob/a04446e57ec1fbc252a871afcec7752fb2807b14/algebra.tex\#L9944}{Original source, lines 9944--9959}. Complete original and reviewed TeX passages accompany this note. Every intervention is separately recorded; the following treatment is editorial.


\begin{lemma}
\label{editorial-source-lemma-generating-finitely-generated-separable-field-extensions}
Let $K/k$ be a separably generated and finitely generated
field extension.
Set $r = \text{trdeg}_k(K)$. Then there exist elements
$x_1, \ldots, x_{r + 1}$ of $K$ such that
\begin{enumerate}
\item $x_1, \ldots, x_r$ is a transcendence basis of $K$ over $k$,
\item $K = k(x_1, \ldots, x_{r + 1})$, and
\item $x_{r + 1}$ is separable over $k(x_1, \ldots, x_r)$.
\end{enumerate}
\end{lemma}

\begin{proof}
Choose a separating transcendence basis of the
original extension. Since $K/k$ is finitely
generated, it has finite size $r$; enumerate
the chosen elements as $x_1,\ldots,x_r$.
Put $F=k(x_1,\ldots,x_r)$, retaining this
original subfield of $K$. If
$K=k(z_1,\ldots,z_m)$ is a finite generating
list, each $z_j$ is algebraic over $F$,
and $K=F(z_1,\ldots,z_m)$. The successive
finite algebraic degrees therefore prove
that $K/F$ is finite. It is separable by
the chosen separating basis. Fields,
Lemma \ref{fields-lemma-primitive-element}
supplies an original $x_{r+1}\in K$ with
$K=F(x_{r+1})$. Its minimal polynomial over
$F$ is separable because the extension is.
This proves all three assertions with the
same original elements. If $K=F$, take
$x_{r+1}=0$, whose minimal polynomial $T$
is separable. If $r=0$, the basis is empty,
$F=k$, and the same finite-extension
argument applies without any transcendental
generators.
\end{proof}


\subsection{Editorial treatment of Algebra lines 9961--10032}

\hypertarget{algebra-context-118}{}

\noindent\href{https://github.com/stacks/stacks-project/blob/a04446e57ec1fbc252a871afcec7752fb2807b14/algebra.tex\#L9961}{Original source, lines 9961--10032}. Complete original and reviewed TeX passages accompany this note. Every intervention is separately recorded; the following treatment is editorial.


\begin{lemma}
\label{editorial-source-lemma-make-separably-generated}
Let $K/k$ be a finitely generated field extension.
There exists a diagram
$$
\xymatrix{
K \ar[r] & K' \\
k \ar[u] \ar[r] & k' \ar[u]
}
$$
where $k'/k$, $K'/K$ are finite purely inseparable field
extensions such that $K'/k'$ is a separably generated field extension.
\end{lemma}

\begin{proof}
Choose the original transcendence basis
$x_1,\ldots,x_r$ of $K$ over $k$, and put
$F=k(x_1,\ldots,x_r)$. Its size $r$ is finite
because the original field extension is finitely
generated. If $K=k(z_1,\ldots,z_m)$ is a specified
finite generating family, then every $z_j$ is
algebraic over $F$ and $K=F(z_1,\ldots,z_m)$.
Successively adjoining these original elements
therefore makes $K/F$ finite. In characteristic
zero every irreducible algebraic polynomial has
nonzero derivative and is separable, so this
basis already separates $K/k$ and one can take
$k'=k$, $K'=K$.

Suppose now that the characteristic is $p>0$.
Let $K/K_{sep}/F$ be the original separable-first
tower from Fields, Lemma \ref{fields-lemma-separable-first}.
Thus $K_{sep}/F$ is finite separable and
$K/K_{sep}$ is finite purely inseparable. Retain
$d=[K:K_{sep}]$. If $d=1$, the original basis
is separating and again no extension is needed.
We construct the strict induction step for $d>1$.

Choose an original $\eta\in K\setminus K_{sep}$.
Pure inseparability gives a least $h\geq1$
such that $\eta^{p^h}\in K_{sep}$. Set
$$
\beta=\eta^{p^{h-1}},\qquad
\alpha=\beta^p=\eta^{p^h}.
$$
Then $\beta\notin K_{sep}$ and
$\alpha\in K_{sep}$. The polynomial $T^p-\alpha$
is irreducible over $K_{sep}$, since a $p$-th
root of $\alpha$ there would equal $\beta$
in the ambient field, contradicting this choice;
see Fields, Lemma \ref{fields-lemma-take-pth-root}.
Hence $[K_{sep}(\beta):K_{sep}]=p$.
Keep the original monic minimal polynomial
$$
P(T)=T^n+a_1T^{n-1}+\cdots+a_n\in F[T]
$$
of $\alpha$. It is separable since
$\alpha\in K_{sep}$; in particular $P'(\alpha)\ne0$.

Here is the omitted finite coefficient construction.
For every original $a_i$, choose complete polynomial
numerator and denominator expressions
$$
a_i=\frac{A_i(x_1,\ldots,x_r)}{B_i(x_1,\ldots,x_r)},
\qquad A_i,B_i\in k[X_1,\ldots,X_r],\quad B_i\ne0.
$$
No common factors are canceled and all coefficients
and summands of these chosen polynomials are kept.
Let $\mathcal C\subset k$ be the finite set of
all their coefficients. Work in one fixed
algebraic closure $\Omega$ of the original $K$,
and choose the unique $p$-th roots there of
the elements of $\mathcal C$ and of the $x_j$.
Set
$$
k'=k(c^{1/p}:c\in\mathcal C),\qquad
y_j=x_j^{1/p},\qquad
F'=k'(y_1,\ldots,y_r).
$$
Each adjoined constant satisfies $U^p-c=0$.
Thus $k'/k$ is finite purely inseparable,
every element of $k'$ has $p$-th power in $k$,
and $[k':k]\leq p^{|\mathcal C|}$.
The last assertion about powers follows for
polynomials and fractions from the characteristic
$p$ identity $(a+b)^p=a^p+b^p$.
The elements $y_1,\ldots,y_r$ are algebraically
independent over $k'$: if
$\sum_\nu d_\nu y^\nu=0$ is a nonzero polynomial
relation, raising it to the $p$-th power gives
$\sum_\nu d_\nu^p x^\nu=0$ over $k$.
At least one coefficient remains nonzero because
Frobenius is injective on a field, contradicting
the original algebraic independence of the $x_j$.

Write the full original expansions as
$A_i(X)=\sum_\nu a_{i,\nu}X^\nu$ and
$B_i(X)=\sum_\nu b_{i,\nu}X^\nu$, with finite
support and nonnegative integer multi-indices.
Their coefficient-root polynomials are
$$
\widetilde A_i(Y)=\sum_\nu a_{i,\nu}^{1/p}Y^\nu,
\qquad
\widetilde B_i(Y)=\sum_\nu b_{i,\nu}^{1/p}Y^\nu.
$$
The denominator is nonzero at the algebraically
independent original $y_j$, and the full identities
are
$$
\widetilde A_i(y)^p=A_i(x),\qquad
\widetilde B_i(y)^p=B_i(x),\qquad
\left(\frac{\widetilde A_i(y)}{\widetilde B_i(y)}\right)^p=a_i.
$$
Put $b_i=\widetilde A_i(y)/\widetilde B_i(y)\in F'$.
This retains each original numerator, denominator,
monomial and coefficient and proves the required
$p$-th power property for every $a_i$, including zero.
The field $F'$ contains $F$ by the original
inclusions and the identities $x_j=y_j^p$.
Moreover $(F')^p\subseteq F$ and
$[F':F]\leq p^{|\mathcal C|+r}$.

Let $L=KF'\subseteq\Omega$ be the actual compositum.
It gives the original commutative diagram
$$
\xymatrix{
K\ar[r]&L\\
k(x_1,\ldots,x_r)\ar[u]\ar[r]&
k'(x_1^{1/p},\ldots,x_r^{1/p})\ar[u].
}
$$
Every arrow is the indicated inclusion in $\Omega$.
The extension $L/K$ is finite purely inseparable
and $L^p\subseteq K$, because the only new
generators are the displayed constant roots
and coordinate roots. In particular
$[L:K]\leq p^{|\mathcal C|+r}$.
The original field $L$ is also finite over $F'$,
as $K$ is finite over $F$.

Put $C=K_{sep}F'$. A finite collection of separable
generators of $K_{sep}/F$ remains separable after
extension to $F'$: their minimal polynomials over
the larger base divide their original separable
polynomials. Thus $C/F'$ is finite separable.
The extension $L/C$ is purely inseparable. To
check this on every element, take finite
generators of $K/K_{sep}$, choose a common
$p$-power placing all of them in $K_{sep}$,
and apply that power to any rational expression
in these generators with coefficients in $C$.
All resulting powers lie in $C$.
It follows that the original separable-first
field $L_{sep}$ over $F'$ is exactly $C$.
Indeed all of $C$ is separable over $F'$, while
an element of $L$ separable over $F'$ is also
separable over $C$, and is purely inseparable
over $C$ by what was just proved. Its minimal
polynomial over $C$ is then both separable and
purely inseparable, hence linear. This proves
the identification $L_{sep}=K_{sep}F'$ without
an assumption about linear disjointness.

The chosen $\beta$ already belongs to $L_{sep}$.
To verify this with the original polynomial, put
$$
Q(T)=T^n+b_1T^{n-1}+\cdots+b_n\in F'[T].
$$
Then $Q(\beta)^p=P(\alpha)=0$, so
$Q(\beta)=0$. Keeping all derivative factors gives
$$
\begin{split}
Q'(\beta)&=n\beta^{n-1}
+\sum_{i=1}^{n-1}(n-i)b_i\beta^{n-i-1},\\
Q'(\beta)^p&=n\alpha^{n-1}
+\sum_{i=1}^{n-1}(n-i)a_i\alpha^{n-i-1}
=P'(\alpha)\ne0.
\end{split}
$$
Here each integer coefficient is its element of
the prime field, whose $p$-th power is itself;
terms divisible by $p$ are retained with their
actual zero coefficients. Thus the minimal
polynomial of $\beta$ over $F'$ is separable:
it divides $Q$, and differentiating that
factorization at $\beta$ shows that its
derivative is nonzero. Consequently
$\beta\in L_{sep}$. Equivalently, Fields,
Lemma \ref{fields-lemma-pth-root} applies to
the original $\alpha$ and $P$ inside the separable
extension $L_{sep}/F'$; its root $\alpha'$
satisfies $(\alpha')^p=\alpha=\beta^p$ and
therefore $\alpha'=\beta$ by injectivity of
Frobenius in $\Omega$. This proves the original
root identification, including its coefficient map.

All entries of the original induction tower are
therefore the actual indicated subfields:
$$
\xymatrix{
K\ar[r]&L\\
K_{sep}(\beta)\ar[r]\ar[u]&L_{sep}\ar[u]\\
K_{sep}\ar[r]\ar[u]&L_{sep}\ar@{=}[u]\\
k(x_1,\ldots,x_r)\ar[r]\ar[u]&
k'(x_1^{1/p},\ldots,x_r^{1/p})\ar[u]\\
k\ar[r]\ar[u]&k'\ar[u].
}
$$
The decrease can be bounded exactly. Since
$K_{sep}(\beta)\subseteq L_{sep}$ and
$L=KL_{sep}$, a basis of $K$ over
$K_{sep}(\beta)$ spans $L$ over $L_{sep}$.
For completeness, its span is the image of
$L_{sep}\otimes_{K_{sep}(\beta)}K\to L$.
This image is a finite-dimensional algebra over
$L_{sep}$, is a domain as a subring of $L$,
and contains both original fields. Any nonzero
element has invertible multiplication on this
finite-dimensional vector space, so the image
is itself a field. It is therefore the
compositum $L$, proving the asserted spanning.
The original tower law now gives
$$
[L:L_{sep}]
\leq[K:K_{sep}(\beta)]
=\frac{[K:K_{sep}]}{[K_{sep}(\beta):K_{sep}]}
=\frac d p<d.
$$
This verifies the strict induction step without
replacing it by an unproved degree comparison.
Applying the same construction to $L/k'$ with
the retained basis $y_1,\ldots,y_r$ completes
the induction. Compositions of the finite
purely inseparable extensions remain finite
and purely inseparable: choose common powers
for finite generating families at successive
stages and compose those powers. Thus the
final inclusions $k\subseteq k''$ and
$K\subseteq L'$ give the required original
diagram, with $k''$ and $L'$ used as the
final $k'$ and $K'$ in the statement.

\medskip\noindent
The same proof provides finite quantitative
information. The original purely inseparable
degree is $d=p^e$ for some integer $e\geq0$.
Indeed adjoining a purely inseparable algebraic
element has degree a power of $p$: its irreducible
polynomial has only one root in an algebraic
closure, and repeated removal of zero derivatives
writes it as a separable irreducible polynomial
in $T^{p^a}$; the separable polynomial must be
linear because there is only one root. A finite
generating tower and the degree law give $d=p^e$.
Every nonterminal step divides the bound for
the remaining inseparable degree by at least $p$.
Consequently the construction terminates after
some $s\leq e$ steps.

Write the actual fields at these steps as
$k_h\subseteq K_h$ for $0\leq h\leq s$,
with $k_0=k$ and $K_0=K$, and retain the
separating-coordinate candidates
$x_j^{1/p^h}$ in the fixed $\Omega$.
At step $h$, let $c_h$ be the number of
coefficients in the finite set used for its
chosen complete rational numerator and
denominator expressions. The construction gives
$$
[k_{h+1}:k_h]\leq p^{c_h},\qquad
[K_{h+1}:K_h]\leq p^{c_h+r},\qquad
k_{h+1}^{,p}\subseteq k_h,\qquad
K_{h+1}^{,p}\subseteq K_h.
$$
Iterating these actual tower maps yields
$$
[k_s:k]\leq p^{\sum_{h=0}^{s-1}c_h},\qquad
[K_s:K]\leq p^{\sum_{h=0}^{s-1}(c_h+r)},\qquad
k_s^{,p^s}\subseteq k,\qquad K_s^{,p^s}\subseteq K.
$$
The terminal extension
$K_s/k_s(x_1^{1/p^s},\ldots,x_r^{1/p^s})$
is separable by the termination criterion.
These are bounds for the original finite
construction, with all constants and coordinate
root adjunctions retained. They do not assert
that the coefficient counts $c_h$ depend only
on $d$, nor that the original transcendence
basis separates $K/k$ before these extensions.
When $r=0$ the coordinate list is empty; when
$s=0$ the sums are empty, the bounds are $1$,
and the original fields already give the result.
\end{proof}


\subsection{Editorial treatment of Algebra lines 10037--10288}

\hypertarget{algebra-context-119}{}

\noindent\href{https://github.com/stacks/stacks-project/blob/a04446e57ec1fbc252a871afcec7752fb2807b14/algebra.tex\#L10037}{Original source, lines 10037--10288}. Complete original and reviewed TeX passages accompany this note. Every intervention is separately recorded; the following treatment is editorial.


\subsubsection{Geometrically reduced algebras}
\label{editorial-source-section-geometrically-reduced}

\noindent
The main result on geometrically reduced algebras is
Lemma \ref{editorial-source-lemma-geometrically-reduced-finite-purely-inseparable-extension}.
We suggest the reader skip to the lemma after reading the definition.

\begin{definition}
\label{editorial-source-definition-geometrically-reduced}
Let $k$ be a field. Let $S$ be a $k$-algebra.
We say $S$ is {\it geometrically reduced over $k$}
if for every field extension $K/k$ the
$K$-algebra $K \otimes_k S$ is reduced.
\end{definition}

\noindent
Let $k$ be a field and let $S$ be a reduced $k$-algebra.
To check that $S$ is geometrically reduced it will suffice
to check that $\overline{k} \otimes_k S$ is reduced (where
$\overline{k}$ denotes the algebraic closure of $k$).
In fact it is enough to check this for finite purely inseparable
field extensions $k'/k$. See
Lemma \ref{editorial-source-lemma-geometrically-reduced-finite-purely-inseparable-extension}.

\begin{lemma}
\label{editorial-source-lemma-subalgebra-separable}
Elementary properties of geometrically reduced algebras.
Let $k$ be a field. Let $S$ be a $k$-algebra.
\begin{enumerate}
\item If $S$ is geometrically reduced over $k$ so is every
$k$-subalgebra.
\item If all finitely generated $k$-subalgebras of $S$ are
geometrically reduced, then $S$ is geometrically reduced.
\item A directed colimit of geometrically reduced $k$-algebras
is geometrically reduced.
\item If $S$ is geometrically reduced over $k$, then any localization
of $S$ is geometrically reduced over $k$.
\end{enumerate}
\end{lemma}

\begin{proof}
Fix an arbitrary original field extension $K/k$.
For (1), if $S'\subseteq S$ is the specified
$k$-subalgebra, the map
$K\otimes_kS'\to K\otimes_kS$ is injective.
Indeed choosing a $k$-basis of $K$ identifies it
with the direct sum of copies of the original
injection. If an element of the first tensor
ring is nilpotent, its image in the reduced
second ring is zero, and injectivity makes the
original element zero. This proves (1).

For (2), retain an original nilpotent element
$z=\sum_{i=1}^n\lambda_i\otimes s_i$ of
$K\otimes_kS$, with $z^d=0$ for some $d\geq1$.
The same sum belongs to $K\otimes_kS_0$ for
the original finite subalgebra
$S_0=k[s_1,\ldots,s_n]\subseteq S$.
The map into $K\otimes_kS$ is injective by
the same basis argument, so the identical
$d$-th power is zero already in this subalgebra.
It is reduced by the hypothesis on $S_0$;
thus the sum is zero and so is its original
image $z$. This proves (2), including the
empty sum, which is zero from the outset.

For (3), write the original system as
$(S_i,\varphi_{ij})$ and its directed colimit
as $S$. The comparison
$$
\mathop{\rm colim}_i(K\otimes_kS_i)
\longrightarrow K\otimes_kS,
\qquad [\lambda\otimes s_i]\longmapsto
\lambda\otimes[s_i]
$$
is an isomorphism. Its inverse sends a finite
sum of tensors to representatives at a common
stage. Independence follows because any finite
list of equalities of colimit elements holds
at a later common stage; the tensor balancing
and additivity relations are finite such lists.
These formulas give both composites on the
original elementary tensors. If the class of
$z_i\in K\otimes_kS_i$ has $d$-th power zero,
then at some later stage $j$ the image of
$z_i^d$ is zero. The image of $z_i$ has
$d$-th power zero in the reduced ring
$K\otimes_kS_j$, and is therefore zero there.
Thus its original colimit class is zero.
No injectivity of the transition maps is used.
This proves the second and third properties
with the actual tensor-colimit comparison.

For (4), let $U\subseteq S$ be the original
multiplicative set and let $U_K$ be its image
$\{1\otimes u:u\in U\}$ in $K\otimes_kS$.
There are inverse algebra maps
$$
K\otimes_kU^{-1}S\longrightarrow
U_K^{-1}(K\otimes_kS),\qquad
\lambda\otimes s/u\longmapsto
\frac{\lambda\otimes s}{1\otimes u},
$$
and
$$
\frac{\sum_i\lambda_i\otimes s_i}{1\otimes u}
\longmapsto\sum_i\lambda_i\otimes s_i/u.
$$
Products of denominators still have the displayed
form. The original balancing maps send every
denominator to a unit, so localization gives
well-defined homomorphisms; their composites
fix both tensor factors and the inverses of
the original denominators. A localization of
a reduced ring $B$ is reduced: if
$(b/u)^d=0$, there is $v$ in the multiplicative
set with $vb^d=0$. Then
$(vb)^d=v^{d-1}(vb^d)=0$, so $vb=0$ in $B$.
This is a zero witness for the original fraction
$b/u$. If zero is inverted, the localization
is the zero ring and the same conclusion holds.
Apply this to the reduced $K\otimes_kS$.
Since $K$ was arbitrary, all four assertions
hold for the original algebras.
\end{proof}

\begin{lemma}
\label{editorial-source-lemma-geometrically-reduced-permanence}
Let $k$ be a field.
If $R$ is geometrically reduced over $k$,
and $S \subset R$ is a multiplicative subset, then the localization
$S^{-1}R$ is geometrically reduced over $k$.
If $R$ is geometrically reduced over $k$, then $R[x]$ is geometrically
reduced over $k$.
\end{lemma}

\begin{proof}
For the original multiplicative set $S\subseteq R$,
the preceding proof gives the exact maps
$$
K\otimes_kS^{-1}R\simeq
S_K^{-1}(K\otimes_kR),\qquad
\lambda\otimes r/s\longmapsto
\frac{\lambda\otimes r}{1\otimes s},
$$
where $S_K=\{1\otimes s:s\in S\}$.
Its inverse sends a fraction with numerator
$\sum_i\lambda_i\otimes r_i$ and denominator
$1\otimes s$ to $\sum_i\lambda_i\otimes r_i/s$.
That proof verifies the original balancing,
fraction witnesses and reducedness of this
localization, including the zero localization.

For the original polynomial ring the comparison is
$$
K\otimes_kR[x]\longrightarrow(K\otimes_kR)[x],
\qquad \lambda\otimes\sum_i r_i x^i
\longmapsto\sum_i(\lambda\otimes r_i)x^i.
$$
Its inverse sends each coefficient tensor
times $x^i$ to $\lambda\otimes r_i x^i$
and extends additively, retaining every
original coefficient and exponent. The
monomial bases verify both composites,
balancing and multiplication.
If a nonzero polynomial over a reduced
ring $B$ has highest nonzero coefficient
$b_d$, its $n$-th power has coefficient
$b_d^n$ at the original degree $nd$.
That coefficient is nonzero, since $b_d$
is nonzero and $B$ reduced. Thus a nonzero
polynomial cannot be nilpotent. Applying
this to $B=K\otimes_kR$ proves reducedness
for every original $K/k$, as required.

The same comparison and induction apply
to any finite list of original variables.
A polynomial in an arbitrary variable
set uses only a finite sublist; its
nilpotence relation therefore holds in
that finite polynomial subring by the
injective monomial inclusion. Thus the
polynomial assertion holds for arbitrary
sets of indeterminates. It is an equivalence:
the constant inclusion
$K\otimes_kR\hookrightarrow(K\otimes_kR)[x_i]$
is injective, with evaluation at all
$x_i=0$ a left inverse. Hence geometric
reducedness of the polynomial algebra
implies that of the original coefficient
algebra. Further localization of these
original polynomial rings is covered
by the proved localization maps.
\end{proof}

\noindent
In the proofs of the following lemmas we will repeatedly use
the following observation: Suppose that $R' \subset R$ and
$S' \subset S$ are inclusions of $k$-algebras.
Then the map $R' \otimes_k S' \to R \otimes_k S$
is injective. More explicitly it factors as
$$
R'\otimes_kS'\longrightarrow R\otimes_kS'
\longrightarrow R\otimes_kS.
$$
The first map is injective by choosing a
$k$-basis of $S'$ and applying the original
injection $R'\to R$ in each coordinate.
The second is injective by choosing a
$k$-basis of $R$ and applying $S'\to S$
in each coordinate. Both maps send every
original elementary tensor to its indicated
image and are algebra homomorphisms, so
their composite is the stated injection.

\begin{lemma}
\label{editorial-source-lemma-limit-argument}
Let $k$ be a field. Let $R$, $S$ be $k$-algebras.
\begin{enumerate}
\item If $R \otimes_k S$ is nonreduced, then there exist
finitely generated subalgebras $R' \subset R$,
$S' \subset S$ such that $R' \otimes_k S'$ is not reduced.
\item If $R \otimes_k S$ contains a nonzero zerodivisor, then there exist
finitely generated subalgebras $R' \subset R$,
$S' \subset S$ such that $R' \otimes_k S'$ contains a nonzero zerodivisor.
\item If $R \otimes_k S$ contains a nontrivial idempotent, then there exist
finitely generated subalgebras $R' \subset R$,
$S' \subset S$ such that $R' \otimes_k S'$ contains a nontrivial idempotent.
\end{enumerate}
\end{lemma}

\begin{proof}
For (1), retain a nonzero nilpotent original
$z=\sum_{i=1}^n x_i\otimes y_i$ with
$z^d=0$, $d\geq1$. Let $R'$ be the
original subalgebra generated by all $x_i$
and $S'$ that generated by all $y_i$.
The same full sum defines $z'\in R'\otimes_kS'$.
Its image is $z$, and the preceding
injection gives $(z')^d=0$. It is nonzero,
since its image is nonzero. Thus this
original tensor subalgebra is not reduced.

For (2), retain both original witnesses
$z\ne0$ and $w\ne0$ with $zw=0$.
Write each as a finite tensor sum, and
generate $R'$ and $S'$ using all the
coefficients from both sums. Their
preimages $z',w'$ are nonzero, and
$z'w'=0$ by the same tensor injection.
Thus $z'$ is a nonzero zerodivisor
in the stated finite subalgebra.
Keeping only the coefficients of $z$
would not by itself supply its required
nonzero annihilating witness.

For (3), retain an original idempotent
$e$ with $e^2=e$, $e\ne0$ and $e\ne1$.
Use all the coefficients of a finite
tensor expression of $e$. The tensor
injection is unital, so its preimage
$e'$ satisfies $(e')^2=e'$ and has
$e'\ne0,1$. This proves (3).

More generally any fixed finite collection
of original tensor elements and polynomial
equations and inequations over $k$ among
them can be witnessed in such finitely
generated $R',S'$. Include every coefficient
of the finite tensor expressions for all
the specified elements. Each original
equation descends because the algebra
map is injective, and each inequation
is preserved because its evaluated image
is nonzero. This assertion concerns the
complete specified finite list and keeps
all witnesses; no infinite collection
is replaced by a finite one.
\end{proof}

\begin{lemma}
\label{editorial-source-lemma-geometrically-reduced-any-reduced-base-change}
Let $k$ be a field.
Let $S$ be a geometrically reduced $k$-algebra.
Let $R$ be any reduced $k$-algebra.
Then $R \otimes_k S$ is reduced.
\end{lemma}

\begin{proof}
We first retain the original finite-subalgebra
argument and its maps. Let
$z=\sum_{i=1}^n r_i\otimes s_i$ be nilpotent
in the original $R\otimes_kS$, and let
$R_0=k[r_1,\ldots,r_n]\subseteq R$.
It is reduced as a subring of the reduced
original $R$ and Noetherian by
Lemma \ref{editorial-source-lemma-Noetherian-permanence}.
The map $R_0\otimes_kS\to R\otimes_kS$ is
injective by the vector-space tensor comparison,
so the same original sum has its power zero
already in $R_0\otimes_kS$.

The original minimal-prime localizations of
$R_0$ are fields $F_j=(R_0)_{\mathfrak p_j}$,
and there are finitely many of them. The maps
$$
R_0\longrightarrow\prod_j F_j,\qquad
r\longmapsto(r/1)_j
$$
are injective, by
Lemmas \ref{editorial-source-lemma-Noetherian-irreducible-components}
and \ref{algebra-lemma-minimal-prime-reduced-ring}.
More explicitly each localization has just its
zero prime, since it is reduced and has a
unique prime. Its kernel on $R_0$ is
$\mathfrak p_j$: a numerator outside that
prime becomes a unit, and a numerator inside
it maps into the zero maximal ideal. The
intersection of the minimal primes is the
nilradical, which is zero. This proves the
injection without replacing the original
$R_0$ by its product of fields.
It agrees with the original total-fraction
comparison from
Lemma \ref{algebra-lemma-total-ring-fractions-no-embedded-points}:
a fraction $r/t$ with $t$ a nonzerodivisor
maps to $(r/1)/(t/1)$ in every factor.
That lemma identifies the total fraction ring
with this finite product; the argument here
uses its indicated injection and maps.

Tensoring the injection with the $k$-vector
space $S$ preserves injectivity. For this
finite product the exact comparison is
$$
\left(\prod_j F_j\right)\otimes_kS
\longrightarrow\prod_j(F_j\otimes_kS),\qquad
(a_j)_j\otimes s\longmapsto(a_j\otimes s)_j.
$$
Its inverse sends $(w_j)_j$ to
$\sum_j(\iota_j\otimes1)(w_j)$, where
$\iota_j:F_j\to\prod_jF_j$ is the original
coordinate inclusion as a $k$-linear map.
The finite sum proves both inverse identities
and preserves multiplication, since distinct
coordinate factors multiply to zero.
Each $F_j\otimes_kS$ is reduced by geometric
reducedness of $S$. A product of reduced rings
is reduced, because nilpotence is checked
coordinate by coordinate. Thus the image of
the original $z$ in this product is zero.
Both injections show $z=0$ in the original
$R\otimes_kS$. If $R_0=0$, the original sum
is already zero; the empty product causes
no exception. This proves the lemma.

\medskip\noindent
The same result yields an exact formula with
an arbitrary original $k$-algebra $A$, which
need not be reduced. Write
$\mathcal N=\sqrt{(0)}\subseteq A$ for its
original nilradical. Then
$$
\sqrt{(0)}_{\,A\otimes_kS}
=\operatorname{image}\bigl(\mathcal N\otimes_kS
\longrightarrow A\otimes_kS\bigr)
=\mathcal N(A\otimes_kS).
$$
The tensor map is injective, by flatness over
the field. Its image is the displayed ideal:
multiplying a tensor $n\otimes s$ by
$a\otimes t$ gives $an\otimes st$, still
with $an\in\mathcal N$, and every generator
$(n\otimes1)(1\otimes s)$ is such a tensor.
Every original finite element
$w=\sum_{i=1}^m n_i\otimes s_i$ of this image
is nilpotent. Choose $d_i\geq1$ with
$n_i^{d_i}=0$ and retain
$D=1+\sum_{i=1}^m(d_i-1)$.
For $m\geq1$ its full multinomial expansion is
$$
w^D=\sum_{\substack{a_1+\cdots+a_m=D\\a_i\geq0}}
\frac{D!}{a_1!\cdots a_m!}
\left(\prod_{i=1}^m n_i^{a_i}\right)
\otimes\left(\prod_{i=1}^m s_i^{a_i}\right)=0.
$$
For each tuple some $a_i\geq d_i$; otherwise
their sum would be at most $D-1$. Thus every
displayed term is zero with its full integer
multinomial coefficient retained. These are
integer coefficients acting in the ring,
not divisions by possibly noninvertible
factorials. The empty sum is zero.

Conversely, exactness of tensoring gives the
quotient comparison
$$
(A\otimes_kS)/\operatorname{image}(\mathcal N\otimes_kS)
\longrightarrow(A/\mathcal N)\otimes_kS,
\qquad [a\otimes s]\longmapsto(a+\mathcal N)\otimes s.
$$
The inverse is induced by the same elementary
formula in the reverse direction, independent
of the representative because the difference
has numerator in $\mathcal N$. The original
quotient $A/\mathcal N$ is reduced, so the
proved lemma makes this tensor quotient reduced.
Every nilpotent in $A\otimes_kS$ consequently
maps to zero there, proving the reverse
inclusion and the exact nilradical identity.

If $S\ne0$, the map $A\to A\otimes_kS$,
$a\mapsto a\otimes1$, is injective: extend
$1\in S$ to a $k$-basis and take the linear
functional $\ell:S\to k$ with $\ell(1)=1$
and zero on the other basis elements.
Then $1\otimes\ell$ is a left inverse.
It follows that for this geometrically
reduced nonzero $S$, the original $A$ is
reduced if and only if $A\otimes_kS$ is
reduced. For $S=0$ the tensor is always
zero, so this last converse is not asserted;
the nilradical formula itself still holds.
\end{proof}

\begin{lemma}
\label{editorial-source-lemma-separable-extension-preserves-reducedness}
Let $k$ be a field.
Let $S$ be a reduced $k$-algebra.
Let $K/k$ be either a separable field extension,
or a separably generated field extension.
Then $K \otimes_k S$ is reduced.
\end{lemma}

\begin{proof}
We first prove the assertion for an original
finitely generated and separably generated
extension $K/k$ and an original domain $D$
over $k$. Choose the original elements
$x_1,\ldots,x_r,\alpha=x_{r+1}$ from
Lemma \ref{editorial-source-lemma-generating-finitely-generated-separable-field-extensions}.
Put $E=k(x_1,\ldots,x_r)$, so that
$K=E(\alpha)$, and retain the monic separable
minimal polynomial $P(T)\in E[T]$ of $\alpha$.
If $K=E$, the choice $\alpha=0$, $P(T)=T$
gives the same proof.

The tensor map
$$
k[x_1,\ldots,x_r]\otimes_kD
\longrightarrow D[x_1,\ldots,x_r],\qquad
\left(\sum_\nu c_\nu x^\nu\right)\otimes d
\longmapsto\sum_\nu c_\nu d x^\nu
$$
is an isomorphism. Its inverse sends
$\sum_\nu d_\nu x^\nu$ to
$\sum_\nu x^\nu\otimes d_\nu$; the original
monomial bases prove balancing and both
composites. With
$W=k[x_1,\ldots,x_r]\setminus\{0\}$,
the original fraction map gives
$$
A=E\otimes_kD\simeq W^{-1}D[x_1,\ldots,x_r],
\qquad (g/h)\otimes d\longmapsto dg/h.
$$
The inverse sends $\sum_\nu d_\nu x^\nu/h$
to $\sum_\nu(x^\nu/h)\otimes d_\nu$.
The universal property of inverting the
same original $h\in W$ verifies these maps,
including equality witnesses for fractions.
Every element of $W$ remains nonzero in the
polynomial ring over the domain $D$, because
the original coefficient map $k\to D$ is
injective. Thus $A$ is a domain. Retain its
fraction field $F=\operatorname{Frac}(A)$.

There is an exact original algebra comparison
$$
A[T]/(P)\longrightarrow K\otimes_kD,
\qquad T\longmapsto\alpha\otimes1,
$$
with coefficients $e\otimes d$ sent to the
same original tensor. The powers
$1,\alpha,\ldots,\alpha^{n-1}$, where
$n=\deg P$, form an $E$-basis of $K$.
Consequently both sides are free $A$-modules
with these corresponding bases, proving the
map an isomorphism. Its inverse sends
$\gamma\otimes d$, for the unique expansion
$\gamma=\sum_{i=0}^{n-1}e_i\alpha^i$, to
$\sum_{i=0}^{n-1}(e_i\otimes d)T^i$ modulo
the original $P$. Multiplication is preserved
because both sides impose exactly $P(\alpha)=0$.
The map $A[T]/(P)\to F[T]/(P)$ is injective:
monic division leaves a unique remainder
of degree less than $n$, and its coefficients
cannot become zero in the fraction field
unless they were zero in $A$.

Separability supplies polynomials $U,V\in E[T]$
with the full identity $UP+VP'=1$.
The original coefficient maps carry this
identity into $F[T]$. Thus $P$ has no repeated
irreducible factor there. Write its actual
factorization $P=\prod_{j=1}^t P_j$ into
distinct monic irreducible polynomials over
$F$. The Chinese remainder comparison sends
$$
F[T]/(P)\longrightarrow\prod_{j=1}^t F[T]/(P_j),
\qquad [H]\longmapsto([H])_j.
$$
For the inverse, put $H_j=P/P_j$ and choose
$U_j$ with $U_jH_j\equiv1\pmod{P_j}$.
A tuple $([G_j])_j$ maps to
$[\sum_jG_jU_jH_j]$ modulo $P$. It is
independent of lifts and composes to the
identity, since each term is $G_j$ in its
own factor and zero in the others; a
polynomial vanishing modulo all the coprime
$P_j$ is divisible by their original
product $P$. Each quotient is a field.
Hence $F[T]/(P)$ is reduced, and the
proved injections show $K\otimes_kD$
is reduced for the original $K$ and $D$.

Now suppose $K/k$ is separable according to
the definition, with no finiteness hypothesis,
and retain an original nilpotent
$z=\sum_{i=1}^m\lambda_i\otimes s_i$
in $K\otimes_kS$. Let
$K_0=k(\lambda_1,\ldots,\lambda_m)\subseteq K$
and $S_0=k[s_1,\ldots,s_m]\subseteq S$.
By definition $K_0/k$ is separably generated
and finitely generated. The original $S_0$
is reduced and Noetherian. The tensor map
$K_0\otimes_kS_0\to K\otimes_kS$ is injective,
so the same original sum is nilpotent already
in the smaller tensor ring.
As in the preceding reduced-base-change
proof, retain the injection
$S_0\to\prod_j (S_0)_{\mathfrak p_j}$
through its finitely many minimal-prime
fields, with every map $s\mapsto s/1$.
This is the original total-fraction comparison
of Lemmas \ref{algebra-lemma-minimal-prime-reduced-ring}
and \ref{algebra-lemma-total-ring-fractions-no-embedded-points}.
Tensoring with $K_0$ is injective, and the
finite-product tensor map is an isomorphism
with inverse given by the finite sum of
coordinate inclusions. Each factor
$K_0\otimes_k(S_0)_{\mathfrak p_j}$ is
reduced by the domain case just proved.
Thus the original sum is zero in their
product and, by injectivity, zero in
$K\otimes_kS$. This proves the separable case.
If $S_0=0$ the tensor was already zero.

Finally suppose $K/k$ is separably generated,
possibly by an infinite original separating
basis $(x_i)_{i\in I}$. We do not assume
the later theorem that this implies separability.
Put $E=k(x_i:i\in I)$ and take any finite
list $\lambda_1,\ldots,\lambda_m\in K$.
Each $\lambda_j$ has a monic separable
minimal polynomial $P_j\in E[T]$, with a
Bezout identity $U_jP_j+V_jP_j'=1$ in $E[T]$.
Every coefficient in this finite collection
of complete polynomials and identities is
a rational function in finitely many of
the original $x_i$. Choose a finite
$J\subseteq I$ containing all those indices.
The same polynomials and identities then
belong to $E_J[T]$, where
$E_J=k(x_i:i\in J)$, since its inclusion
into $E$ is injective. Each $\lambda_j$
is algebraic and separable over $E_J$:
its minimal polynomial divides the same
monic $P_j$, which has no repeated roots
by the retained Bezout identity.
Therefore
$$
K_{J,\lambda}=k(x_i:i\in J)(\lambda_1,\ldots,\lambda_m)
$$
is finite separable over $E_J$, and is a
finitely generated, separably generated
subextension of the original $K/k$.
The original algebraic independence of the
$x_i$ shows that the selected sublist is
indeed a transcendence basis for this field.

These subfields form a directed family:
for two finite lists and two finite index
sets, include their unions and the finitely
many coefficient indices required for the
combined list. Enlarging the rational base
retains separability because each minimal
polynomial still divides its specified
separable polynomial. Their union is $K$,
since every original element is included
by taking it in a one-element list.
Every nilpotent tensor in $K\otimes_kS$
therefore belongs to one of these original
finite subextensions tensored with $S$.
Its power relation holds there by injectivity,
and that tensor algebra is reduced by the
finite separably generated argument above
and the finite-subalgebra argument for $S$.
Hence the original tensor is zero. This
also gives explicitly the directed-colimit
description in the original proof.

As a consequence, either hypothesis on
$K/k$ makes $K$ geometrically reduced over
$k$: apply the result with an arbitrary
receiving field as the reduced algebra $S$
and use the tensor flip
$\lambda\otimes s\mapsto s\otimes\lambda$,
whose inverse is the same flip. Moreover,
for every original $k$-algebra $B$, the
nilradical formula of the preceding lemma
now gives the exact equality
$$
\sqrt{(0)}_{\,K\otimes_kB}
=\operatorname{image}\bigl(K\otimes_k\sqrt{(0)}_B
\longrightarrow K\otimes_kB\bigr).
$$
Thus this result preserves the entire
nilradical under the original scalar
extension, not only reducedness when
that ideal happens to vanish.
\end{proof}

\begin{lemma}
\label{editorial-source-lemma-generic-points-geometrically-reduced}
Let $k$ be a field and let $S$ be a $k$-algebra. Assume that
$S$ is reduced and that $S_{\mathfrak p}$ is geometrically
reduced for every minimal prime $\mathfrak p$ of $S$.
Then $S$ is geometrically reduced.
\end{lemma}

\begin{proof}
Retain the original localizations $S_{\mathfrak p}$
at every minimal prime. Since $S$ is reduced,
so is each localization. Its sole prime is
$\mathfrak pS_{\mathfrak p}$, and the intersection
of its primes is its zero nilradical; hence
$\mathfrak pS_{\mathfrak p}=0$ and this local
ring is a field. The original map
$S\to S_{\mathfrak p}$ has kernel exactly
$\mathfrak p$: elements in it map to zero,
whereas an element outside it maps to a unit
in this nonzero field. The intersection of
the minimal primes of $S$ is zero, so
$$
S\longrightarrow\prod_{\mathfrak p\text{ minimal}}S_{\mathfrak p},
\qquad s\longmapsto(s/1)_{\mathfrak p}
$$
is injective, as in
Lemma \ref{algebra-lemma-reduced-ring-sub-product-fields}.

For any original field extension $K/k$, tensoring
this map with $K$ preserves injectivity. The
second map in the source proof is
$$
\Theta:\left(\prod_{\mathfrak p}S_{\mathfrak p}\right)\otimes_kK
\longrightarrow\prod_{\mathfrak p}(S_{\mathfrak p}\otimes_kK),
\qquad (a_{\mathfrak p})_{\mathfrak p}\otimes\lambda
\longmapsto(a_{\mathfrak p}\otimes\lambda)_{\mathfrak p}.
$$
To prove its injectivity, write any original
tensor as a finite sum
$w=\sum_{j=1}^r a_j\otimes\lambda_j$ with
$k$-linearly independent $\lambda_j\in K$.
Such an expression is obtained by expanding
the finitely many original scalar coefficients
in a basis of their span, retaining the
resulting sums in each $a_j$.
Extend these $\lambda_j$ to a basis of $K$.
If $\Theta(w)=0$, then for every original
$\mathfrak p$ the finite basis expansion
$\sum_j(a_j)_{\mathfrak p}\otimes\lambda_j=0$
forces $(a_j)_{\mathfrak p}=0$ for each $j$.
Thus every original $a_j$ is zero and $w=0$.
The map is an algebra map since it is
coordinatewise on elementary products.

Each factor $S_{\mathfrak p}\otimes_kK$ is
reduced by its assumed geometric reducedness.
The product is reduced, so the two original
injections imply that $S\otimes_kK$ is reduced.
The tensor flip identifies this with the
ordering in the definition, and proves the
lemma for every $K$. If $S=0$ the index set
is empty and all the displayed rings are
zero, which is reduced.

The converse also holds: if $S$ is geometrically
reduced, taking $K=k$ shows that $S$ is reduced,
and Lemma \ref{editorial-source-lemma-subalgebra-separable} shows
that each of its original localizations is
geometrically reduced. Thus the original
conditions are an equivalence.

The same proof applies more generally whenever
there is a specified injective $k$-algebra map
$S\hookrightarrow\prod_{i\in I}B_i$ and each
$B_i$ is geometrically reduced. The tensor
map into the product of the base changes
is injective by the identical finite-basis
argument, so $S$ is geometrically reduced.
In particular, an arbitrary product of
geometrically reduced $k$-algebras is
geometrically reduced. Conversely if such
a product is geometrically reduced, fix
an original index $i$. Decomposing it as
the two factors $B_i\times\prod_{j\ne i}B_j$
and tensoring uses the finite-product
isomorphism. The first factor $B_i\otimes_kK$
is reduced for every $K$, because a nilpotent
there would give a nilpotent with zero
second coordinate in the reduced product.
Hence each $B_i$ is geometrically reduced.
The empty product is the zero algebra and
satisfies the assertion as well.

Injectivity of $\Theta$ must not be strengthened
to surjectivity without further hypotheses.
For example take every $B_i=k$, indexed by
$i\geq1$, and $K=k(t)$. The sequence
$(t^i)_{i\geq1}\in\prod_iK$ is outside the
image of $(\prod_i k)\otimes_kK$.
Indeed a finite tensor expression uses a
single finite list $\lambda_1,\ldots,\lambda_r$
of elements of $K$, so every coordinate
of its image belongs to their fixed finite
$k$-linear span. The powers $t^i$ are
$k$-linearly independent, since a nonzero
polynomial in the transcendental $t$ cannot
vanish. They cannot all lie in that span.
The original proof and its extensions use
only the established injection.
\end{proof}

\begin{lemma}
\label{editorial-source-lemma-separable-algebraic-diagonal}
Let $k'/k$ be a separable algebraic extension.
Then there exists a multiplicative subset $S \subset k' \otimes_k k'$
such that the multiplication map $k' \otimes_k k' \to k'$
is identified with $k' \otimes_k k' \to S^{-1}(k' \otimes_k k')$.
\end{lemma}

\begin{proof}
First let the original $k'/k$ be finite
separable, and choose its original primitive
element $\alpha$ using Fields,
Lemma \ref{fields-lemma-primitive-element}.
Retain its monic separable minimal polynomial
$P(x)\in k[x]$, of degree $n\geq1$.
The original tensor map, with both sums
included, is
$$
\psi:k'[x]/(P)\longrightarrow k'\otimes_kk',
\qquad
\left[\sum_i\alpha_i x^i\right]
\longmapsto\sum_i\alpha_i\otimes\alpha^i.
$$
The relation $P(1\otimes\alpha)=0$ makes this
well defined. Both rings are free over the
first $k'$ with bases corresponding to
$1,x,\ldots,x^{n-1}$ and
$1\otimes1,1\otimes\alpha,\ldots,1\otimes\alpha^{n-1}$.
The latter basis is obtained by tensoring
the original $k$-basis of the second $k'$.
Thus $\psi$ is an isomorphism. Its inverse
sends $a\otimes b$, with the unique original
expansion $b=\sum_{i=0}^{n-1}c_i\alpha^i$,
$c_i\in k$, to $[\sum_{i=0}^{n-1}ac_i x^i]$.
This also verifies its balancing and its
product law through the same polynomial $P$.
Under this map multiplication
$\mu:k'\otimes_kk'\to k'$ corresponds to
evaluation at $x=\alpha$.

Retain the factorization $P(x)=(x-\alpha)Q(x)$
in the original $k'[x]$. Differentiating gives
$P'(\alpha)=Q(\alpha)\ne0$ by separability.
For $B=k'[x]/(P)$ and the multiplicative set
$S'$ generated by the class of $Q$, the
evaluation map extends to
$$
B_Q\longrightarrow k',\qquad
\frac{[H(x)]}{[Q(x)]^d}\longmapsto
\frac{H(\alpha)}{Q(\alpha)^d}.
$$
Its inverse sends $a\in k'$ to $a/1$.
Indeed $(x-\alpha)Q=0$ in $B$; once the
original $Q$ is a unit, $x=\alpha$.
Consequently every displayed fraction is
the image of its stated scalar value,
and both composites are identities. The
fraction map is independent of lifts and
zero-divisor witnesses by the localization
universal property. Transporting this exact
set $S'$ under $\psi$ gives the required
original multiplicative set $S$.

There is also an explicit original direct
factor. Put $q=Q(\alpha)\in k'\setminus\{0\}$
and retain the polynomial $H$ defined by
$Q(x)-q=(x-\alpha)H(x)$. In $B$, set
$e=[Q(x)]/q$, where the scalar denominator
is the original invertible $q$.
Then the full identities are
$$
e^2-e=\frac{Q(x)(Q(x)-q)}{q^2}
=\frac{P(x)H(x)}{q^2}=0,
\qquad (x-\alpha)e=\frac{P(x)}q=0,
\qquad e(\alpha)=1.
$$
Thus $e$ is an idempotent and
$\ker(B\to k')=(x-\alpha)B=(1-e)B$.
For the equality of ideals, one inclusion
uses $1-e=-(x-\alpha)H/q$; the other uses
$(x-\alpha)=(x-\alpha)(1-e)$.
The map $eB\to k'$ is evaluation, with
inverse $a\mapsto ae$. Also $B_Q=B_e$,
because $Q=qe$ with $q$ already a unit.
These formulas retain the original $Q$
and its full scalar factor. If $n=1$,
then $Q=1$, $H=0$, $e=1$, and the kernel
is zero, so the same formulas apply.

Now retain an arbitrary separable algebraic
extension $k'/k$. Its finite intermediate
fields $k_i$ form a directed family: a
compositum of two finite subextensions
inside $k'$ is again finite and separable.
Every finite set of original coefficients
of $k'$ lies in some $k_i$. The injections
$$
k_i\otimes_kk_i\longrightarrow k'\otimes_kk'
$$
are injective by tensoring the two original
field inclusions over the field $k$.
Their images have union $k'\otimes_kk'$,
since every tensor is a finite sum of
elementary tensors. For each $i$, use the
actual multiplicative set $S_i$ and
idempotent $e_i$ constructed above for
$k_i/k$, and let $S$ be the multiplicative
closure of their images $\bigcup_iS_i$.
The original multiplication map $\mu$
sends every element of $S$ to a nonzero
element of $k'$, because this holds for
each original generator $Q_i$ and their
finite products. It therefore extends to
$S^{-1}(k'\otimes_kk')\to k'$.

If an original $z\in\ker\mu$ is represented
at a stage $k_i\otimes_kk_i$, its multiplication
there is zero: its image in the containing
field $k'$ is zero and $k_i\hookrightarrow k'$
is injective. The finite-stage proof makes
$z$ zero on inverting that stage's $Q_i$;
indeed $Q_i$ annihilates its diagonal kernel.
It is therefore zero in the full original
localization. Multiplication is surjective
already before localization, since
$\mu(a\otimes1)=a$ for every $a\in k'$.
The localized map is consequently bijective:
a fraction mapping to zero has numerator
in $\ker\mu$, which is killed as just proved.
Its inverse is explicitly
$a\mapsto(a\otimes1)/1$.
For a general original fraction $z/s$,
the difference from the image of
$\mu(z)/\mu(s)$ has numerator
$(\mu(s)\otimes1)z-(\mu(z)\otimes1)s$,
whose multiplication is zero; its denominator
is the original product $(\mu(s)\otimes1)s$.
This verifies the second composite and
keeps all scalar and localization factors.

\medskip\noindent
The diagonal kernel also gives a sharp finite
generation consequence. For this original
separable algebraic extension, the ideal
$I=\ker(k'\otimes_kk'\to k')$ is finitely
generated if and only if $k'/k$ is finite.
The finite case follows from the explicit
idempotent generator $1-e$ above.
For the converse, suppose
$I=(z_1,\ldots,z_m)$, and choose a finite
original subfield $k_i$ containing every
coefficient of all these tensors.
Every $z_j$ is in the finite diagonal
kernel, so $e_i z_j=0$ by its proved
idempotent formula. Therefore $e_iI=0$.
Since $\mu(e_i)=1$, also $1-e_i\in I$,
and the exact equality is
$I=(1-e_i)(k'\otimes_kk')$.

The quotient by this particular ideal is
canonically $k'\otimes_{k_i}k'$. In detail,
$\ker(k_i\otimes_kk_i\to k_i)$ is generated
by $a\otimes1-1\otimes a$, $a\in k_i$:
modulo these relations the two original
copies of every $a$ coincide, and multiplication
and $a\mapsto a\otimes1$ are inverse maps.
Extending these relations to $k'\otimes_kk'$
imposes precisely the balancing over $k_i$.
The map sends the original $a\otimes_k b$
to $a\otimes_{k_i}b$, and its inverse is
the same elementary formula into the
quotient; the just-stated relations prove
well-definedness and both composites.
Since the extended finite kernel is exactly
$(1-e_i)$, the full kernel assumption would
make multiplication
$k'\otimes_{k_i}k'\to k'$ an isomorphism.
If $k'\ne k_i$, choose an original
$b\in k'\setminus k_i$ and extend the
$k_i$-independent list $1,b$ to a basis
of $k'$. In the induced basis of the
second tensor factor, the element
$1\otimes_{k_i}b-b\otimes_{k_i}1$ is nonzero,
but its multiplication is zero. This
contradicts that isomorphism. Hence
$k'=k_i$ and the original extension is
finite. The argument includes an empty
generating list for the zero ideal.
\end{proof}

\begin{lemma}
\label{editorial-source-lemma-geometrically-reduced-over-separable-algebraic}
Let $k'/k$ be a separable algebraic field extension.
Let $A$ be an algebra over $k'$. Then $A$ is geometrically
reduced over $k$ if and only if it is geometrically reduced over $k'$.
\end{lemma}

\begin{proof}
First assume that the original $A$ is geometrically
reduced over $k'$, and let $K/k$ be any original
field extension. The ring $K\otimes_kk'$ is
reduced by
Lemma \ref{editorial-source-lemma-separable-extension-preserves-reducedness},
applied to the separable field extension $k'/k$
and the reduced $k$-algebra $K$, followed by
the tensor flip. Treat it as a $k'$-algebra
through its original second factor. Then
Lemma \ref{editorial-source-lemma-geometrically-reduced-any-reduced-base-change}
over $k'$ shows that
$(K\otimes_kk')\otimes_{k'}A$ is reduced.
The exact comparison is
$$
(K\otimes_kk')\otimes_{k'}A\longrightarrow K\otimes_kA,
\qquad (\lambda\otimes b)\otimes a\longmapsto\lambda\otimes ba,
$$
with inverse $\lambda\otimes a\mapsto
(\lambda\otimes1)\otimes a$.
Balancing over the original $k'$ moves $b$
into $a$, and balancing over $k$ is the
original tensor relation. These formulas
verify both composites and multiplication.
Hence $K\otimes_kA$ is reduced for every
$K/k$, proving this direction.

Conversely assume that the original $A$ is
geometrically reduced over $k$, and let $K/k'$
be an arbitrary original field extension.
Retain the rings
$$
H=k'\otimes_kk',\qquad B=K\otimes_kA.
$$
The original $H$-algebra structure on $B$
is the map
$$
\chi:H\longrightarrow B,\qquad
b\otimes c\longmapsto b\otimes c,
$$
where $b$ enters through $k'\to K$ and
$c$ through $k'\to A$. Give $k'$ the
$H$-algebra structure of multiplication
$\mu(b\otimes c)=bc$.
The source's exact comparison, now including
its original actions and punctuation, is
$$
B\otimes_Hk'\longrightarrow K\otimes_{k'}A,
\qquad (\lambda\otimes a)\otimes c
\longmapsto c\lambda\otimes a.
$$
It is balanced over $H$: multiplying the
first tensor by $b\otimes d$ gives
$cb\lambda\otimes da= cbd\lambda\otimes a$
over $k'$, equal to multiplying the last
scalar by $bd$. Its inverse sends
$\lambda\otimes_{k'}a$ to
$(\lambda\otimes_ka)\otimes1$.
For its $k'$-balancing, the difference
$c\lambda\otimes a-\lambda\otimes ca$
is the action of $c\otimes1-1\otimes c$
on $\lambda\otimes a$, and that element
of $H$ is killed by $\mu$. The inverse
is therefore well defined. Both composites
are identities, since the scalar $c$ can
be moved from the last factor by the
original element $c\otimes1\in H$.

Lemma \ref{editorial-source-lemma-separable-algebraic-diagonal}
gives an actual multiplicative set $U\subseteq H$
and the original identification $k'\simeq U^{-1}H$
through $\mu$. Thus the preceding tensor
ring is the localization $\chi(U)^{-1}B$.
The full comparison sends
$$
b\otimes\frac hu\longmapsto
\frac{b\chi(h)}{\chi(u)},
$$
with inverse $b/\chi(u)\mapsto b\otimes1/u$.
Balancing and the same original inverse
denominators prove both maps and composites.
Since $K$ is a field extension of $k$,
the hypothesis makes $B=K\otimes_kA$
reduced. Its localization is reduced by
the zero-witness proof already given.
Therefore the original $K\otimes_{k'}A$
is reduced for every $K/k'$, completing
the converse without identifying a general
quotient of a reduced ring with a reduced ring.

\medskip\noindent
The forward implication has a wider scope:
it holds if the original $k'/k$ is any
separable or separably generated extension,
with no algebraicity assumption. The first
paragraph uses only that
$K\otimes_kk'$ is reduced for every $K/k$,
and the preceding separable-extension lemma
proves exactly this in either case. Its
two displayed original tensor maps remain
unchanged. The reverse implication does
not hold for arbitrary separably generated
extensions, as the following example shows.

Let $k$ be a field of characteristic $p>0$,
take $k'=k(t)$ and $A=k(u)$ with the
original embedding $k'\to A$ given by
$t\mapsto u^p$, where $u$ is transcendental
over $k$. The field $k'/k$ is separably
generated using its basis $t$, and $A$
is geometrically reduced over $k$.
Indeed for any field $L/k$,
$$
L\otimes_kk(u)\simeq
(k[u]\setminus\{0\})^{-1}L[u]
$$
through $\lambda\otimes f(u)/g(u)\mapsto
\lambda f(u)/g(u)$, with inverse obtained
by expanding each numerator in its
original powers of $u$. The specified
nonzero denominators stay nonzero over
$L$, so this is a domain.

The polynomial $T^p-t$ is irreducible
over the original $k(t)$. There is no
$p$-th root of $t$ in $k(t)$: for nonzero
polynomials $f,g\in k[t]$, the equation
$f(t)^p=tg(t)^p$ would imply
$p\operatorname{ord}_t(f)=1+p\operatorname{ord}_t(g)$,
which is impossible. These orders are the
actual smallest powers of $t$ with nonzero
coefficients; the product law follows by
multiplying the corresponding first
nonzero coefficients. No common factor
of the original $f,g$ is removed.
Fields, Lemma \ref{fields-lemma-take-pth-root}
therefore gives irreducibility and the
original degree $[k(u):k(u^p)]=p$.
For the receiving field $K=A$ over $k'$
the original tensor map gives
$$
A\otimes_{k'}A\simeq A[T]/(T^p-u^p),
\qquad \sum_{i=0}^{p-1}a_i\otimes u^i
\longleftrightarrow\left[\sum_{i=0}^{p-1}a_iT^i\right].
$$
More precisely the first $A$ supplies
the coefficients and $1\otimes u$
corresponds to $T$; the original basis
$1,u,\ldots,u^{p-1}$ of the second
factor proves the isomorphism and its
inverse on all tensors. The element
$$
z=1\otimes u-u\otimes1
$$
is nonzero: its polynomial representative
$T-u$ has degree $1<p$ and nonzero leading
coefficient, so it is not a multiple of
the original monic $T^p-u^p$.
But its full binomial power is
$$
(T-u)^p=\sum_{j=0}^p
\binom pj T^j(-u)^{p-j}=T^p-u^p=0
$$
in this quotient, since each middle
integer coefficient is divisible by $p$.
Thus $A$ is not geometrically reduced
over the original $k'$, despite being
geometrically reduced over $k$ and
despite $k'/k$ being separably generated.
This locates the role of algebraicity
in the converse diagonal-localization proof.

Separability cannot simply be dropped
from the original algebraic assertion
either. Take $k=\mathbf F_p(t)$,
$k'=\mathbf F_p(u)$ with the same
$t\mapsto u^p$, and $A=k'$.
The identity algebra $A$ is geometrically
reduced over $k'$, because every field
base change is that field. Its tensor
with the receiving field $k'$ over $k$
is the same nonreduced ring
$k'[T]/(T^p-u^p)$ with the same original
nonzero nilpotent. Hence it is not
geometrically reduced over $k$.
\end{proof}








\subsection{Editorial treatment of Algebra lines 10295--10350}

\hypertarget{algebra-context-120}{}

\noindent\href{https://github.com/stacks/stacks-project/blob/a04446e57ec1fbc252a871afcec7752fb2807b14/algebra.tex\#L10295}{Original source, lines 10295--10350}. Complete original and reviewed TeX passages accompany this note. Every intervention is separately recorded; the following treatment is editorial.


\begin{lemma}
\label{editorial-source-lemma-mini-separability}
Let $k$ be a field of characteristic $p > 1$. Let $K/k$
be a field extension generated by $x_1, \ldots, x_{n + 1} \in K$ such that
\begin{enumerate}
\item $\{x_1, \ldots, x_n\}$ is a transcendence base of $K/k$,
\item for every $k$-linearly independent subset
$\{a_1, \ldots, a_m\}$ of $K$ the set
$\{a^p_1, \ldots, a_m^p\}$ is $k$-linearly independent.
\end{enumerate}
Then there is $1 \leq j \leq n+1$ such that
$\{ x_1, \ldots, \widehat{x}_j, \ldots, x_{n+1}\}$
is a separating transcendence base for $K / k$.
\end{lemma}

\begin{proof}
Retain the original generators $x_1,\ldots,x_{n+1}$.
Since $x_{n+1}$ is algebraic over $k(x_1,\ldots,x_n)$,
clearing the nonzero polynomial denominators in one of its
algebraic equations gives a nonzero
$F\in k[X_1,\ldots,X_{n+1}]$ with $F(x_1,\ldots,x_{n+1})=0$.
Choose such an $F$ of least total degree $D$. A nonzero
constant cannot vanish, so $D\geq1$. If $F=GH$ with both
factors of positive degree, evaluation in the field $K$
makes at least one factor zero. That factor has degree
less than $D$, a contradiction. Thus $F$ is irreducible.

Write the full polynomial as
$$
F=\sum_{\alpha\in A}\lambda_\alpha
X_1^{\alpha_1}\cdots X_{n+1}^{\alpha_{n+1}},
\qquad 0\neq\lambda_\alpha\in k,
$$
where $A$ is its finite set of distinct exponent tuples,
all of whose coordinates are integers at least zero.
Suppose every coordinate of every $\alpha\in A$ were
divisible by $p$. Put
$b_\alpha=x_1^{\alpha_1/p}\cdots x_{n+1}^{\alpha_{n+1}/p}$.
The family $(b_\alpha^p)_{\alpha\in A}$ is dependent,
since $\sum_{\alpha\in A}\lambda_\alpha b_\alpha^p=0$.
Hypothesis (2) also preserves independence of indexed
finite families: an independent family has distinct
entries, and Frobenius is injective on $K$, so its
$p$th powers still have distinct entries and hypothesis
(2) applies to the corresponding sets. Its contrapositive
therefore makes $(b_\alpha)_{\alpha\in A}$ dependent,
including the case where two evaluated monomials agree.
Choose $c_\alpha\in k$, not all zero, such that
$\sum_\alpha c_\alpha b_\alpha=0$. The polynomial
$$
G=\sum_{\alpha\in A}c_\alpha
X_1^{\alpha_1/p}\cdots X_{n+1}^{\alpha_{n+1}/p}
$$
is nonzero because its formal monomials are distinct.
It vanishes at the same original generators and has
degree at most $D/p<D$, a contradiction. Consequently
some exponent of some $X_i$ in $F$ is not divisible
by $p$, equivalently $\partial F/\partial X_i\neq0$.

Fix any such $i$, and retain the original coefficient
polynomials in
$$
F=\sum_{r=0}^h A_r(X_1,\ldots,\widehat X_i,\ldots,X_{n+1})X_i^r.
$$
There is $r\geq1$ with $p\nmid r$ and $A_r\neq0$.
Every nonzero $A_r$ with $r\geq1$ has total degree
at most $D-r<D$. Hence its value at the other original
generators is nonzero: otherwise that same coefficient
polynomial, viewed in all $n+1$ variables, would be a
smaller nonzero relation. In particular the polynomial
$$
P(T)=F(x_1,\ldots,x_{i-1},T,x_{i+1},\ldots,x_{n+1})
\in L[T],\qquad
L=k(x_1,\ldots,\widehat x_i,\ldots,x_{n+1}),
$$
is nonconstant and has $P(x_i)=0$. Thus $K=L(x_i)$
is finite algebraic over $L$. A transcendence basis of
$L/k$, chosen among its $n$ displayed generators by
Fields, Lemma \ref{fields-lemma-transcendence-degree},
is also a transcendence basis of $K/k$: algebraicity
in a tower is transitive. Its size must therefore be
$n$. All the displayed generators of $L$ are consequently
algebraically independent over $k$.

Evaluation thus identifies the polynomial ring
$B=k[X_1,\ldots,\widehat X_i,\ldots,X_{n+1}]$
with $k[x_1,\ldots,\widehat x_i,\ldots,x_{n+1}]$,
and identifies its fraction field with the original $L$.
The irreducible polynomial $F\in B[X_i]$ is primitive:
a nonunit common divisor of its coefficients would
factor $F$, with the other factor still of positive
$X_i$-degree. Gauss' lemma now makes the unchanged
polynomial $P(T)$ irreducible in $L[T]$.
Its derivative is the full sum
$$
P'(T)=\sum_{r=1}^h r A_r(x_1,\ldots,\widehat x_i,\ldots,x_{n+1})T^{r-1}.
$$
The previously chosen term is nonzero, so $P'\neq0$.
Irreducibility and $\deg P'<\deg P$ give
$\gcd(P,P')=1$. Hence $x_i$ is separable over $L$,
and $K/L$ is separable. This proves the assertion,
also when $n=0$ and the displayed basis of $L=k$ is empty.

The proof gives an exact criterion for every omitted
coordinate, not merely existence of one. For this same
minimal-degree $F$, omission of $x_j$ gives a separating
transcendence basis if and only if
$\partial F/\partial X_j\neq0$. Sufficiency was proved
for every such index. For necessity, suppose the other
$n$ elements form a separating basis and put
$L_j=k(x_1,\ldots,\widehat x_j,\ldots,x_{n+1})$.
The polynomial $F$ must depend on $X_j$, since otherwise
it would be a relation among that basis. Its coefficient
evaluation is injective, and the same primitive Gauss
argument makes its image $P_j(T)$ irreducible over $L_j$.
If $c_j$ is its actual leading coefficient and $M_j(T)$
is the monic minimal polynomial of $x_j$ over $L_j$,
then $P_j(T)=c_jM_j(T)$ with $c_j\neq0$. Separability
gives $M_j'\neq0$, so $P_j'=c_jM_j'\neq0$.
Injectivity of coefficient evaluation yields
$\partial F/\partial X_j\neq0$, as required.

There is also an exact relation-ring comparison for
these original generators. Evaluation induces
$$
k[X_1,\ldots,X_{n+1}]/(F)
\longrightarrow k[x_1,\ldots,x_{n+1}].
$$
It is surjective. If an original polynomial $H$ is in
its kernel, view $H$ in $B[X_i]$ for the index already
chosen. Over $L$, its image is divisible by the
minimal polynomial of $x_i$, hence by the original
$P=c_iM_i$ as $c_i$ is a unit of $L$. Primitive
divisibility in Gauss' lemma gives $F\mid H$ already
in $B[X_i]$. The kernel is therefore exactly $(F)$,
and the displayed map is an isomorphism with the
original evaluation map as its defining morphism.
\end{proof}


\subsection{Editorial treatment of Algebra lines 10359--10407}

\hypertarget{algebra-context-121}{}

\noindent\href{https://github.com/stacks/stacks-project/blob/a04446e57ec1fbc252a871afcec7752fb2807b14/algebra.tex\#L10359}{Original source, lines 10359--10407}. Complete original and reviewed TeX passages accompany this note. Every intervention is separately recorded; the following treatment is editorial.


\begin{lemma}
\label{editorial-source-lemma-characterize-separable-field-extensions}
Let $k$ be a field of characteristic $p > 0$.
Let $K/k$ be a field extension.
The following are equivalent:
\begin{enumerate}
\item $K$ is separable over $k$,
\item for every $k$-linearly independent subset $\{a_1, \ldots, a_m\}$
of $K$ the set $\{a^p_1, \ldots, a_m^p\}$ is $k$-linearly independent,
\item the ring $K \otimes_k k^{1/p}$ is reduced, and
\item $K$ is geometrically reduced over $k$.
\end{enumerate}
\end{lemma}

\begin{proof}
The implication (1) $\Rightarrow$ (4) is Lemma
\ref{editorial-source-lemma-separable-extension-preserves-reducedness}
applied to every field extension of $k$ as a reduced
$k$-algebra. The symmetry map $a\otimes b\mapsto b\otimes a$
has itself as inverse. The implication (4) $\Rightarrow$ (3)
then follows by using the specified field $k^{1/p}$.

For (3) $\Rightarrow$ (2), retain the original tensor
ring $B=K\otimes_k k^{1/p}$ and the original map
$$
m:B\longrightarrow K,\qquad
\sum_i\lambda_i\otimes\mu_i\longmapsto
\sum_i\lambda_i^p\mu_i^p.
$$
This is a ring homomorphism. Additivity and multiplicativity
follow from Frobenius in characteristic $p$, and balancing
follows from
$(c\lambda)^p\mu^p=\lambda^p(c\mu)^p$ for $c\in k$.
Every $\mu_i^p$ belongs to the original $k$, so the
displayed sum belongs to $K$. For every original
$z=\sum_i\lambda_i\otimes\mu_i$ one has
$$
z^p=\sum_i\lambda_i^p\otimes\mu_i^p
=\left(\sum_i\lambda_i^p\mu_i^p\right)\otimes1
=m(z)\otimes1.
$$
Here all mixed multinomial terms have coefficients
divisible by the prime $p$. The map
$K\to B$, $a\mapsto a\otimes1$, is injective, as
extension of scalars from one field to another is
faithful on vector spaces. If $B$ is reduced, $m(z)=0$
therefore implies $z^p=0$ and then $z=0$.

Let $a_1,\ldots,a_m$ be $k$-linearly independent,
and suppose $\sum_{i=1}^m c_i a_i^p=0$ with $c_i\in k$.
Let $\mu_i\in k^{1/p}$ be the unique root with
$\mu_i^p=c_i$. Then
$m(\sum_i a_i\otimes\mu_i)=0$, so the tensor is zero.
The elements $a_i\otimes1$ are independent over
$k^{1/p}$: extend the original finite family to a
$k$-basis of $K$ and extend that basis by scalars.
Consequently every $\mu_i=0$ and every $c_i=0$.
This proves (2). The argument uses the displayed
ring homomorphism, not an assertion that $m$ is
$k$-linear: in fact $m(cz)=c^p m(z)$ for the
original scalar action of $c\in k$.

For (2) $\Rightarrow$ (1), property (2) is inherited
by every intermediate field of the original $K/k$.
It suffices, by the definition of separability, to
prove that each finitely generated such field is
separably generated. Work with one of these fields,
denoted by $K$ in the remainder of this argument.
Choose a transcendence basis $x_1,\ldots,x_d$ and
retain $K'=k(x_1,\ldots,x_d)\subseteq K$. Since
$K/k$ is finitely generated and $K/K'$ is algebraic,
it is finite by Fields, Lemma
\ref{fields-lemma-algebraic-finitely-generated}.
Among these original bases choose one for which
the positive integer $[K:K']_i$ is least. Such a
least value exists because the set of bases is
nonempty and these values are positive integers.

If $K/K'$ is separable, this basis is separating.
Otherwise choose $x_{d+1}\in K$ not separable over
$K'$, and retain the intermediate field
$E=K'(x_{d+1})$. Then $[E:K']_i>1$.
The original $E/k$ is generated by the displayed
$d+1$ elements, has $x_1,\ldots,x_d$ as a
transcendence basis, and inherits (2) from $K/k$.
Apply Lemma \ref{editorial-source-lemma-mini-separability} to $E/k$,
with its integer $n$ equal to this $d$. It supplies
$1\leq j\leq d+1$ such that
$$
K''=k(x_1,\ldots,\widehat x_j,\ldots,x_{d+1})
\subseteq E
$$
has the indicated generators as a transcendence
basis and $E/K''$ is finite separable. They are
also a transcendence basis of $K/k$, since
$K/E$ is finite algebraic. The two finite towers
are $K/E/K'$ and $K/E/K''$. Fields, Lemma
\ref{fields-lemma-multiplicativity-all-degrees}
gives the full comparison
$$
[K:K']_i=[K:E]_i[E:K']_i,
\qquad
[K:K'']_i=[K:E]_i[E:K'']_i=[K:E]_i.
$$
Thus $[K:K'']_i<[K:K']_i$, contradicting the choice
of basis. This proves separable generation of
every finitely generated intermediate field, and
hence (1) for the original extension, with no
finite-generation assumption on that extension.
When $d=0$, the same argument uses an empty basis
and $K''=k$; no omitted index is out of range.

The same original tensor maps give a stronger
description even when the equivalent conditions
fail. Choose an algebraic closure $\Omega$ of $K$
and embed $k^{1/p}$ there by its unique $p$th roots.
Put $C=Kk^{1/p}\subseteq\Omega$ and
$D=K^p k\subseteq K$, retaining both composita
with their indicated embeddings. Multiplication
defines a ring map
$$
\nu:B\longrightarrow C,\qquad
\sum_i\lambda_i\otimes\mu_i\longmapsto
\sum_i\lambda_i\mu_i.
$$
Its image is the subring of finite such sums.
For a nonzero element $w$ of this subring,
$w^p=\sum_i\lambda_i^p\mu_i^p\in K\setminus\{0\}$,
so $w^{-1}=w^{p-1}/w^p$ lies in the same subring.
It is consequently a field containing $K$ and
$k^{1/p}$, and is exactly $C$. Hence $\nu$ is
surjective. Frobenius on $C$ maps into $K$ and
sends these original sums to $\sum_i\lambda_i^p\mu_i^p$.
Its image is exactly $D$: it is a field containing
$K^p$ and $k$, and each image lies in the field
generated by them. Frobenius on a field is injective,
so it is a ring isomorphism $C\longrightarrow D$.
The original $m$ is precisely its composite with
$\nu$, followed by the inclusion $D\subseteq K$.

Since the target of $\nu$ is a field, every
nilpotent of $B$ belongs to $\ker\nu$. Conversely,
if $\nu(z)=0$, then $m(z)=\nu(z)^p=0$, and the
displayed tensor identity gives $z^p=0$. Therefore
$$
\operatorname{Nil}(B)=\ker\nu=\ker m
=\{z\in B:z^p=0\}.
$$
The exact quotient maps identify
$B/\operatorname{Nil}(B)$ with $C$ by $[z]\mapsto\nu(z)$,
and with $D$ by $[z]\mapsto m(z)$. The first is a
$K$-algebra isomorphism for the original action;
the second sends the scalar $a\in K$ to $a^p\in D$.
Neither map replaces the original tensor ring or
changes its scalar structure silently. In particular
$B$ has a unique prime ideal, namely its nilradical:
primes contain all nilpotents, and its reduced
quotient is the nonzero field $C$. The equivalent
conditions of the lemma can also be expressed as
injectivity of $\nu$, or equivalently as the original
map $K\otimes_k k^{1/p}\to Kk^{1/p}$ being an
isomorphism of $K$-algebras.

Elementwise annihilation by the $p$th power in this
description does not assert that the $p$th ideal
power of the nilradical is zero. For an exact
example, take algebraically independent $u,v$
over $\mathbf F_p$, and retain
$$
s=u^p,\quad t=v^p,\quad
k=\mathbf F_p(s,t),\quad K=\mathbf F_p(u,v).
$$
Then $k^{1/p}=K$ inside an algebraic closure of $K$.
Indeed every rational function in $s,t$ has the
unique root obtained by using the same numerator
and denominator in $u,v$, since Frobenius fixes
each coefficient in $\mathbf F_p$. The elements
$u^a v^b$ for $0\leq a,b<p$ form a $k$-basis of
$K$. Independence follows on clearing the original
denominators in any proposed relation: the exponent
pairs in its polynomial terms have distinct residues
modulo $p$. They span because powers can be reduced
using $u^p=s$, $v^p=t$, and every nonzero polynomial
denominator $q(u,v)$ has inverse
$q(u,v)^{p-1}/q(u,v)^p$, with $q(u,v)^p\in k$.

The original tensor comparison is consequently
$$
B=K\otimes_k K\simeq
K[U,V]/(U^p-s,V^p-t),\qquad
1\otimes u\longleftrightarrow U,\quad
1\otimes v\longleftrightarrow V.
$$
It fixes $K$ through the left tensor factor.
Monic division gives the basis $U^aV^b$ for
$0\leq a,b<p$ on the polynomial side; the stated
basis of the right tensor factor proves the map
and its inverse on every basis vector. Put
$\epsilon=U-u$, $\eta=V-v$. The full binomial
expansions give $\epsilon^p=U^p-u^p=0$ and
$\eta^p=V^p-v^p=0$, because every intermediate
binomial coefficient is divisible by $p$.
Evaluation at $U=u,V=v$ identifies the quotient
by $(\epsilon,\eta)$ with $K$, and all elements
of this ideal have $p$th power zero. Thus it is
the nilradical $N$. Its power $N^{2p-1}$ is zero:
each generating monomial has either epsilon or
eta exponent at least $p$. But
$\epsilon^{p-1}\eta^{p-1}\neq0$, since in the
original basis its $U^{p-1}V^{p-1}$ coefficient
is $1$ and all other terms have smaller exponent
in at least one coordinate. Hence the ideal
nilpotence index is exactly $2p-1$, and in
particular $N^p\neq0$ for every prime $p$.
\end{proof}


\subsection{Editorial treatment of Algebra lines 10409--10417}

\hypertarget{algebra-context-122}{}

\noindent\href{https://github.com/stacks/stacks-project/blob/a04446e57ec1fbc252a871afcec7752fb2807b14/algebra.tex\#L10409}{Original source, lines 10409--10417}. Complete original and reviewed TeX passages accompany this note. Every intervention is separately recorded; the following treatment is editorial.


\begin{lemma}
\label{editorial-source-lemma-separably-generated-separable}
A separably generated field extension is separable.
\end{lemma}

\begin{proof}
In characteristic $p>0$, Lemma
\ref{editorial-source-lemma-separable-extension-preserves-reducedness}
applies directly to the original separably generated
extension $K/k$ and to every reduced $k$-algebra.
In particular $K\otimes_k L$ is reduced for every
field extension $L/k$. Thus $K/k$ is geometrically
reduced, and Lemma
\ref{editorial-source-lemma-characterize-separable-field-extensions}
makes this original extension separable.

In characteristic zero, every finitely generated
intermediate field $E/k$ has a finite transcendence
basis $y_1,\ldots,y_d$, and $E$ is finite algebraic
over $k(y_1,\ldots,y_d)$. Each irreducible polynomial
of positive degree over this field has nonzero
derivative: its leading derivative coefficient is
the original nonzero leading coefficient multiplied
by its positive integer degree. The derivative has
smaller degree, so irreducibility makes their greatest
common divisor equal to a unit. Every algebraic
element of $E$ is consequently separable over the
chosen rational function field. Thus $E/k$ is
separably generated. This holds for every such $E$,
which is precisely the definition of separability
of $K/k$. In fact this argument proves that every
field extension in characteristic zero is separable,
without requiring that the entire extension be
separably generated.
\end{proof}


\subsection{Editorial treatment of Algebra lines 10424--10478}

\hypertarget{algebra-context-123}{}

\noindent\href{https://github.com/stacks/stacks-project/blob/a04446e57ec1fbc252a871afcec7752fb2807b14/algebra.tex\#L10424}{Original source, lines 10424--10478}. Complete original and reviewed TeX passages accompany this note. Every intervention is separately recorded; the following treatment is editorial.


\begin{lemma}
\label{editorial-source-lemma-geometrically-reduced-finite-purely-inseparable-extension}
Let $k$ be a field. Let $S$ be a $k$-algebra.
The following are equivalent:
\begin{enumerate}
\item $k' \otimes_k S$ is reduced for every finite
purely inseparable extension $k'$ of $k$,
\item if $\operatorname{char}(k)=p>0$, then
$k^{1/p}\otimes_k S$ is reduced; if
$\operatorname{char}(k)=0$, then $S$ is reduced,
\item $k^{perf} \otimes_k S$ is reduced, where $k^{perf}$ is the
perfect closure of $k$,
\item $\overline{k} \otimes_k S$ is reduced, where $\overline{k}$ is the
algebraic closure of $k$, and
\item $S$ is geometrically reduced over $k$.
\end{enumerate}
\end{lemma}

\begin{proof}
We treat characteristic zero explicitly. Every field
extension in that characteristic is separable, as
proved in Lemma \ref{editorial-source-lemma-separably-generated-separable}.
Every finite purely inseparable extension is the
identity extension and $k^{perf}=k$. Thus (1), (2)
and (3) each say that the original $S$ is reduced.
If it is reduced, Lemma
\ref{editorial-source-lemma-separable-extension-preserves-reducedness}
makes $L\otimes_k S$ reduced for every field
extension $L/k$. This proves (5), and hence (4).
Conversely (4) implies that $S$ is reduced through
the injection $s\mapsto1\otimes s$ into
$\overline k\otimes_k S$. This proves all equivalences
in characteristic zero without using undefined
$p$th-root notation.

Now suppose $\operatorname{char}(k)=p>0$. Use an
algebraic closure $\overline k$ to contain the
specified perfect closure and $k^{1/p}$. A finite
purely inseparable extension $k'/k$ has a unique
$k$-embedding into $\overline k$, and its image lies
in $k^{perf}$, since each element has a $p$-power
in $k$. Extension of scalars by the $k$-vector
space $S$ preserves these field injections.
Consequently their original maps
$a\otimes s\mapsto a\otimes s$ prove
(5) $\Rightarrow$ (4) $\Rightarrow$ (3) $\Rightarrow$ (2)
and (3) $\Rightarrow$ (1). A subring of a reduced
ring is reduced, which is the exact assertion
used at each injection.

For (1) $\Rightarrow$ (5), let $K/k$ be any original
field extension and let $z=\sum_{i=1}^m a_i\otimes s_i$
be a nilpotent of $K\otimes_k S$. Retain the
finitely generated intermediate field
$K_0=k(a_1,\ldots,a_m)\subseteq K$. The map
$K_0\otimes_k S\to K\otimes_k S$ is injective,
so the same displayed tensor is already nilpotent
in $K_0\otimes_k S$. Apply Lemma
\ref{editorial-source-lemma-make-separably-generated} to this $K_0/k$.
It gives the original commutative field diagram
$$
\xymatrix{
K_0\ar[r]&K_0'\\
k\ar[u]\ar[r]&k'\ar[u]
}
$$
where $k'/k$ and $K_0'/K_0$ are finite purely
inseparable and $K_0'/k'$ is separably generated.
By (1), $S'=k'\otimes_k S$ is reduced. The original
comparison map is
$$
K_0'\otimes_{k'}S'\longrightarrow K_0'\otimes_k S,
\qquad b\otimes(c\otimes s)\longmapsto bc\otimes s.
$$
Its inverse sends $b\otimes s$ to
$b\otimes(1\otimes s)$. Balancing follows from
the displayed commutative field diagram, and
both composites are the identity on generating
pure tensors. Lemma
\ref{editorial-source-lemma-separable-extension-preserves-reducedness}
makes the left side reduced, so the right side
is reduced. The injection
$K_0\otimes_k S\to K_0'\otimes_k S$ now makes
the original $z$ zero. Its image in $K\otimes_k S$
is the tensor we started with. Thus every
nilpotent in that original tensor ring is zero,
proving (5).

For (2) $\Rightarrow$ (5), inject $S$ into
$k^{1/p}\otimes_k S$ by $s\mapsto1\otimes s$;
hence $S$ is reduced. Let $\mathfrak p$ be any
minimal prime of $S$ and put $W=S\setminus\mathfrak p$.
There are inverse original localization maps
$$
k^{1/p}\otimes_k S_{\mathfrak p}
\longrightarrow (1\otimes W)^{-1}(k^{1/p}\otimes_k S),
\qquad a\otimes(s/w)\longmapsto (a\otimes s)/(1\otimes w),
$$
and
$$
\frac{\sum_i a_i\otimes s_i}{1\otimes w}
\longmapsto\sum_i a_i\otimes(s_i/w).
$$
The first comes from inverting exactly $W$ in
the $S$-factor; the second comes from inverting
exactly $1\otimes W$ in the tensor ring. These
universal properties prove well-definedness,
and the formulas show both composites are
identities, with the same original denominators.
The localized tensor ring is reduced by (2).
Since $S_{\mathfrak p}$ is a field by Lemma
\ref{algebra-lemma-minimal-prime-reduced-ring}, Lemma
\ref{editorial-source-lemma-characterize-separable-field-extensions}
applied to this original field says it is
geometrically reduced over $k$. Lemma
\ref{editorial-source-lemma-generic-points-geometrically-reduced}
then proves that $S$ is geometrically reduced.
This also covers the zero ring: it has no
minimal primes and all its scalar extensions
are zero.

There are two useful exact forms of the test in
positive characteristic. First, it suffices in
(1) to test only finite purely inseparable
extensions $l/k$ with $l^p\subseteq k$.
Necessity follows from (5). Conversely every
finite list of coefficients in $k^{1/p}$ lies
in a field $l=k(b_1,\ldots,b_r)$ with $b_i^p\in k$.
This extension is finite of degree at most $p^r$,
and every element of $l$ has its $p$th power in
$k$, by applying Frobenius to its full rational
expression. If a tensor in $k^{1/p}\otimes_k S$
is nilpotent, capture its finite coefficient
list in this original $l$. Injectivity of
$l\otimes_k S\to k^{1/p}\otimes_k S$ places the
same nilpotence equation in the smaller tensor
ring. The restricted test makes it zero there
and therefore zero in the original ring.
Thus the restricted tests imply (2) and all
five conditions.

Second, the field's linear-independence criterion
extends to the original arbitrary $k$-algebra $S$:
geometric reducedness is equivalent to the
following condition. Every finite $k$-linearly
independent family $s_1,\ldots,s_r$ in $S$ has
$k$-linearly independent $p$th powers. To prove
this, retain $A=S\otimes_k k^{1/p}$ and the ring
homomorphism
$$
m_S:A\longrightarrow S,\qquad
\sum_i s_i\otimes b_i\longmapsto\sum_i s_i^p b_i^p.
$$
Balancing uses $(cs)^p b^p=s^p(cb)^p$, and the
full Frobenius identity gives
$z^p=m_S(z)\otimes1$. The original map
$S\to A$ is injective, so $\ker m_S$ is exactly
the kernel of the $p$th-power map on $A$.
That kernel is zero exactly when $A$ is reduced:
for the nontrivial direction, if $z$ is nilpotent,
choose $p^e$ at least a nilpotence exponent and
apply injectivity of Frobenius $e$ times to
$z^{p^e}=0$.

If $m_S$ is injective and
$\sum_i c_i s_i^p=0$, take the unique roots
$b_i^p=c_i$ in $k^{1/p}$. Injectivity makes
$\sum_i s_i\otimes b_i=0$, and extension of the
independent family to a vector-space basis
forces every $b_i=0$, hence every $c_i=0$.
Conversely any original $z\in A$ can be written
$z=\sum_i s_i\otimes b_i$ with $s_i$ independent:
choose a basis of the finite-dimensional span
of its original first-factor coefficients and
rewrite those coefficients in that basis,
using tensor balancing. If the asserted
independence of $p$th powers holds and $m_S(z)=0$,
then every $b_i^p=0$ and hence every $b_i=0$.
Thus $m_S$ is injective. We have proved the
equivalence with reducedness of $A$, which is
(2) after the explicit tensor-factor symmetry.
No domain, finiteness, or Noetherian hypothesis
on the original $S$ is used.
\end{proof}


\subsection{Editorial treatment of Algebra lines 10482--10604}

\hypertarget{algebra-context-124}{}

\noindent\href{https://github.com/stacks/stacks-project/blob/a04446e57ec1fbc252a871afcec7752fb2807b14/algebra.tex\#L10482}{Original source, lines 10482--10604}. Complete original and reviewed TeX passages accompany this note. Every intervention is separately recorded; the following treatment is editorial.


\subsubsection{Perfect fields}
\label{editorial-source-section-perfect-fields}

\noindent
Here is the definition.

\begin{definition}
\label{editorial-source-definition-perfect}
Let $k$ be a field. We say $k$ is {\it perfect}
if every field extension of $k$ is separable over $k$.
\end{definition}

\begin{lemma}
\label{editorial-source-lemma-perfect}
A field $k$ is perfect if and only if it is a field of characteristic $0$
or a field of characteristic $p > 0$ such that every element has a $p$th
root.
\end{lemma}

\begin{proof}
In characteristic zero every field extension is
separable, by the characteristic-zero argument in
Lemma \ref{editorial-source-lemma-separably-generated-separable}.
This is precisely the stated definition of perfect.

Suppose now that $\operatorname{char}(k)=p>0$.
If $k$ is perfect, take any $a\in k$ and its unique
$p$th root $b$ in an algebraic closure. The original
extension $k(b)/k$ is purely inseparable, since
$b^p=a$, and is separable by perfectness. A minimal
polynomial of $b$ divides $T^p-a=(T-b)^p$ over the
algebraic closure. Thus it has only one distinct
root, whereas separability makes all its roots
distinct. Its degree is therefore one, so $b\in k$.
This includes $a=0$, whose root is zero.

Conversely, suppose every $a\in k$ has a $p$th root
in $k$. Frobenius is injective on a field, so those
roots are unique and the specified extension
$k^{1/p}/k$ is the identity. For every original
field extension $K/k$ the tensor map
$$
K\otimes_k k\longrightarrow K,\qquad
\lambda\otimes a\longmapsto a\lambda
$$
has inverse $\lambda\mapsto\lambda\otimes1$.
It identifies $K\otimes_k k^{1/p}$ with the reduced
field $K$. Lemma
\ref{editorial-source-lemma-characterize-separable-field-extensions}
then makes every original $K/k$ separable, so $k$
is perfect.

In particular, in positive characteristic it
suffices to test the separability of the individual
extensions $k(a^{1/p})/k$. If $a$ has no root in $k$,
then $T^p-a$ is irreducible: a proper monic factor
over $k$ would be $(T-b)^d$ over the algebraic
closure for some $1\leq d<p$, whose coefficient
of $T^{d-1}$ is $-db$. Since $d\neq0$ in $k$,
that coefficient would put $b$ in $k$, a contradiction.
Thus every nontrivial one of these tests has degree
exactly $p$ and is inseparable. This proves the
test and its exact degree without assuming a
general extension is finitely generated.
\end{proof}

\begin{lemma}
\label{editorial-source-lemma-make-separable}
Let $K/k$ be a finitely generated field extension.
There exists a diagram
$$
\xymatrix{
K \ar[r] & K' \\
k \ar[u] \ar[r] & k' \ar[u]
}
$$
where $k'/k$, $K'/K$ are finite purely inseparable field
extensions such that $K'/k'$ is a separable field extension.
In this situation we can assume that $K' = k'K$ is the compositum,
and also that $K' = (k' \otimes_k K)_{red}$.
\end{lemma}

\begin{proof}
Lemma \ref{editorial-source-lemma-make-separably-generated} supplies
the original diagram, with its upper-right field
temporarily denoted by $K'_0$, in which $k'/k$
and $K'_0/K$ are finite purely inseparable and
$K'_0/k'$ is separably generated. By Lemma
\ref{editorial-source-lemma-separably-generated-separable}, the last
extension is separable. Retain both original
embeddings into $K'_0$ and let $L=k'K\subseteq K'_0$
be their compositum. Since $k'\subseteq L\subseteq K'_0$,
Lemma \ref{editorial-source-lemma-subextensions-are-separable} makes
$L/k'$ separable. Also $L/K$ is finite purely
inseparable, being an intermediate extension of
$K'_0/K$. Thus choosing the new upper-right field
to be $K'=L$, with its actual inclusion into
$K'_0$, proves the compositum assertion.

For the tensor assertion, first treat characteristic
zero. A purely inseparable extension is then the
identity, so $k'=k$ and $K'_0=K$. The multiplication
map $k\otimes_k K\to K$, $a\otimes b\mapsto ab$,
is an isomorphism with inverse $b\mapsto1\otimes b$.
Its nilradical is zero, and it is the required map.

In characteristic $p>0$, put $B=k'\otimes_k K$.
Because the original extension $k'/k$ is finite
purely inseparable, there is a single $q=p^e$
with $(k')^q\subseteq k$: choose finitely many
field generators, take the largest of their
inseparability exponents, and apply Frobenius
to their full rational expressions. Retain
the multiplication map
$$
\mu:B\longrightarrow K'_0,\qquad
\sum_i a_i\otimes b_i\longmapsto\sum_i a_i b_i.
$$
It is balanced because the original field diagram
commutes. For every original $z=\sum_i a_i\otimes b_i$
one has the full identity
$$
z^q=\sum_i a_i^q\otimes b_i^q
=1\otimes\left(\sum_i a_i^q b_i^q\right),
\qquad
\mu(z)^q=\sum_i a_i^q b_i^q\in K.
$$
The equality for $q$ follows by iterating the
$p$th-power binomial identity; each intermediate
mixed coefficient is divisible by $p$. The map
$K\to B$, $b\mapsto1\otimes b$, is injective.
If $\mu(z)=0$, the coefficient sum in $K$ is
zero, and hence $z^q=0$. Conversely every
nilpotent maps to zero in the field $K'_0$.
Therefore $\ker\mu=\operatorname{Nil}(B)$.

The image of $\mu$ is the subring of finite sums
of products $a_i b_i$. It is a field: for a
nonzero such sum $w$, its $q$th power is a
nonzero element of $K$, and
$w^{-1}=w^{q-1}/w^q$ belongs to the same subring.
The image contains $k'$ and $K$, so it is exactly
the original $L=k'K$. Consequently the original
map gives the canonical $K$-algebra isomorphism
$$
(k'\otimes_k K)_{red}\longrightarrow L,
\qquad [\sum_i a_i\otimes b_i]\longmapsto\sum_i a_i b_i.
$$
Its inverse sends a finite sum in $L$ to the
class of the corresponding tensor; if two such
sums are equal, their difference lies in the
kernel just calculated, proving independence
of the representation. Thus both asserted
choices of $K'$ agree through this exact map.

This last comparison needs neither finite
generation of $K/k$ nor finiteness of the purely
inseparable extension in the first tensor factor.
Indeed let $l/k$ be any algebraic purely
inseparable extension and embed it, by its unique
roots, in an algebraic closure of the original
$K$. Every individual tensor in $l\otimes_k K$
uses finitely many coefficients $a_i\in l$,
so there is a $q=p^e$ for that finite list with
all $a_i^q\in k$. The same two displayed power
identities prove that the multiplication kernel
is exactly the nilradical, and the same inverse
formula for each nonzero image proves that the
image is the original compositum $lK$.
Thus $(l\otimes_k K)_{red}\simeq lK$ with the
same explicit map in this full generality.
Only the exponent may now depend on the tensor.

Finite generation cannot, however, be discarded
from the earlier assertion that one finite
purely inseparable base extension makes the
whole receiving extension separable. Take
$$
k=\mathbf F_p(t),\qquad
K=\bigcup_{r\geq0}\mathbf F_p(t^{1/p^r})
$$
inside a fixed algebraic closure, with $t$
transcendental and all root inclusions retained.
Every finite purely inseparable $k'/k$ has a
unique image in $K$: a root of an original
rational function $A(t)/B(t)$ of order $p^r$
is $A(t^{1/p^r})/B(t^{1/p^r})$, since its
coefficients belong to $\mathbf F_p$.
The elements of $k'$ have
$p^e$th powers in $k$ for some finite $e$.
The element $t^{1/p^{e+1}}$ is not in $k'$:
otherwise its $p^e$th power $t^{1/p}$ would
belong to $k$. The latter is impossible since
a rational function in $t$ cannot have $p$th
power $t$. To see this directly, writing such
a function as $A(t)/B(t)$ would give
$A(t)^p=tB(t)^p$; the exponents on the left
are divisible by $p$ and those on the right
are congruent to $1$ modulo $p$, so a nonzero
polynomial equality is impossible.
Thus $K/k'$ is a nontrivial purely inseparable
extension and cannot be separable. Its
compositum with the original $k'$ is still
$K$. Moreover any purported larger receiving
field separable over $k'$ would have the
subextension $K/k'$ separable, again impossible.
This proves the necessity of some restriction
on the original $K/k$ for the finite-base-change
existence assertion, while retaining the
unrestricted tensor-quotient comparison.
\end{proof}

\begin{lemma}
\label{editorial-source-lemma-perfection}
\begin{slogan}
Every field has a unique perfect closure.
\end{slogan}
For every field $k$ there exists a purely inseparable extension
$k'/k$ such that $k'$ is perfect. The field extension
$k'/k$ is unique up to unique isomorphism.
\end{lemma}

\begin{proof}
In characteristic zero $k$ itself is perfect by
Lemma \ref{editorial-source-lemma-perfect}, and a purely inseparable
extension must be the identity. This gives
existence and uniqueness, including the specified
embedding of $k$.

Let $\operatorname{char}(k)=p>0$, and fix an
algebraic closure $\Omega$ of the original $k$.
For every integer $n\geq0$ put
$$
E_n=\{b\in\Omega:b^{p^n}\in k\}.
$$
The Frobenius identities show that this is a
field: sums and products have $p^n$th powers
in $k$, as does the inverse of a nonzero
element. One has $E_0=k$ and $E_n\subseteq E_{n+1}$.
Every $a\in k$ has a unique $p^n$th root in
$\Omega$, because $T^{p^n}-a$ has a root there
and the difference of any two roots has zero
$p^n$th power in a field. Thus $E_n$ has exactly
properties (a) and (b) of the stated construction.

To compare this with the original abstract
construction, let $k_n$ denote a copy of the
underlying field $k$, with structural map
$i_n:k\to k_n$, $a\mapsto a^{p^n}$.
The transition map $f_n:k_n\to k_{n+1}$ is
$r\mapsto r^p$, and
$f_n(i_n(a))=a^{p^{n+1}}=i_{n+1}(a)$.
The map
$$
\theta_n:k_n\longrightarrow E_n,\qquad
r\longmapsto r^{1/p^n}
$$
is a field isomorphism. Uniqueness of roots
proves additivity and multiplicativity;
surjectivity follows from the definition of
$E_n$, and injectivity follows from Frobenius.
It respects the original $k$-structure since
$\theta_n(i_n(a))=a$, and it respects the
transitions since
$\theta_{n+1}(r^p)=\theta_n(r)$.
Accordingly the union of the actual subfields
$E_n$ realizes the directed colimit of the
original copies with their Frobenius transition
maps. No literal inclusion of a copy with its
different structural map is being assumed.

Set $k'=\bigcup_{n\geq0}E_n$. This is a field,
and every one of its elements has a $p$-power
in the original $k$, so $k'/k$ is purely
inseparable. If $b\in E_n$, its unique $p$th
root in $\Omega$ lies in $E_{n+1}$, since its
$p^{n+1}$th power equals $b^{p^n}\in k$.
Thus Frobenius on $k'$ is surjective, and
Lemma \ref{editorial-source-lemma-perfect} makes $k'$ perfect.

For uniqueness, let $L/k$ be any perfect purely
inseparable extension, with its original
structural embedding of $k$. Every $a\in k$
has a unique $p^n$th root in $L$. Define
$\phi_n:E_n\to L$ by sending $b$ to that
root of $b^{p^n}\in k$. Frobenius and root
uniqueness prove that $\phi_n$ preserves
addition, multiplication, $1$, and the given
embedding of $k$. It is injective as a
nonzero field homomorphism. These maps agree
on the inclusions $E_n\subseteq E_{n+1}$,
again by root uniqueness, and give
$\phi:k'\to L$. Every $c\in L$ has some
$c^{p^n}\in k$, so it is the image of the
unique corresponding root in $E_n$.
Hence $\phi$ is surjective. Any $k$-homomorphism
from $k'$ to $L$ must take every such root
to that same unique root, so it equals $\phi$.
This proves uniqueness up to unique
$k$-isomorphism.

The construction proves the stronger universal
property with its maps: for every perfect
field $P$ and every field homomorphism
$\iota:k\to P$, there is exactly one field
homomorphism $\iota^{perf}:k'\to P$ extending
$\iota$. On $b\in E_n$ it is the unique
$p^n$th root in $P$ of $\iota(b^{p^n})$.
The preceding addition, multiplication and
compatibility proofs apply verbatim; they
did not use pure inseparability of the
receiving field for existence, only for
surjectivity. This also proves functoriality:
two composites that agree on $k$ send every
root to the same unique root. In
characteristic zero the assertion is just
the original map $\iota$ on $k'=k$.
\end{proof}

\begin{definition}
\label{editorial-source-definition-perfection}
Let $k$ be a field. The field extension $k'/k$ of Lemma \ref{editorial-source-lemma-perfection}
is called the {\it perfect closure} of $k$. Notation $k^{perf}/k$.
\end{definition}

\noindent
Let $l/k$ be any algebraic purely inseparable
extension. In positive characteristic, send an
original $a\in l$, with $a^{p^n}\in k$, to
the unique root of $T^{p^n}-a^{p^n}$ in
$k^{perf}$. The value is independent of the
chosen exponent, and a common exponent for
two elements proves additivity and
multiplicativity by Frobenius. This defines
the unique $k$-embedding $l\to k^{perf}$.
In characteristic zero $l=k$ and this is
the identity. Retain this actual embedding.
The extension $k^{perf}/l$ is purely
inseparable because every element has a
$p$-power in $k\subseteq l$, and its upper
field is perfect. It is therefore a
perfect closure of $l$. More explicitly,
the universal property just proved gives
maps $k^{perf}\to l^{perf}$ extending
$k\to l^{perf}$ and $l^{perf}\to k^{perf}$
extending the displayed embedding of $l$.
The first map fixes the original $l$
under these embeddings, since roots of
its elements' powers in $k$ are unique.
Both composites consequently fix the
base fields and every adjoining root,
so are identities. These are the
specified inverse isomorphisms of
perfect closures, not merely an
abstract equality of field types.

\begin{lemma}
\label{editorial-source-lemma-perfect-reduced}
Let $k$ be a perfect field.
Any reduced $k$-algebra is geometrically reduced over $k$.
Let $R$, $S$ be $k$-algebras.
Assume both $R$ and $S$ are reduced.
Then the $k$-algebra $R \otimes_k S$ is reduced.
\end{lemma}

\begin{proof}
By Lemma \ref{editorial-source-lemma-perfect}, either $k$ has
characteristic zero or its Frobenius map is
surjective. In characteristic zero the test
in Lemma
\ref{editorial-source-lemma-geometrically-reduced-finite-purely-inseparable-extension}
is precisely reducedness of the original $S$.
In positive characteristic $k^{1/p}=k$ through
the actual root embedding, and the same test
is reducedness of $k\otimes_k S\simeq S$,
with maps $a\otimes s\mapsto as$ and
$s\mapsto1\otimes s$. Thus every reduced
$k$-algebra is geometrically reduced. Applying
Lemma \ref{editorial-source-lemma-geometrically-reduced-any-reduced-base-change}
to the original reduced $R$ and geometrically
reduced $S$ proves that $R\otimes_k S$ is
reduced, including either zero-ring case.

The hypothesis on $k$ is also necessary for
this conclusion to hold for every pair of
reduced $k$-algebras, or even for every pair
of fields over $k$. If $k$ is not perfect,
Lemma \ref{editorial-source-lemma-perfect} gives characteristic
$p>0$ and an $a\in k$ with no $p$th root in
$k$. Retain $L=k(b)$, $b^p=a$. As proved in
Lemma \ref{editorial-source-lemma-perfect}, $T^p-a$ is
irreducible, so $1,b,\ldots,b^{p-1}$ is a
$k$-basis of $L$. The original map
$$
L[T]/(T^p-a)\longrightarrow L\otimes_k L,
\qquad T\longmapsto1\otimes b
$$
fixes $L$ through the left factor and is an
isomorphism by the corresponding bases.
Its inverse sends $c\otimes\sum_i d_i b^i$
to $\sum_i cd_i T^i$, for $d_i\in k$.
The element $\delta=1\otimes b-b\otimes1$
has $\delta^p=1\otimes a-a\otimes1=0$,
but is nonzero because its inverse image
$T-b$ has degree $1<p$ in the monic
quotient. Consequently the tensor of the
two original reduced fields is not reduced.
The same example shows that $L$ is not
geometrically reduced. Thus perfectness is
equivalent both to every reduced $k$-algebra
being geometrically reduced and to every
tensor product of two reduced $k$-algebras
being reduced.

For arbitrary original $k$-algebras $R,S$
over this perfect field, their nilradicals
$N_R=\operatorname{Nil}(R)$ and
$N_S=\operatorname{Nil}(S)$ give the exact
formula in $T=R\otimes_k S$
$$
\operatorname{Nil}(T)=J,
\qquad
J=\operatorname{im}(N_R\otimes_k S\to T)
+\operatorname{im}(R\otimes_k N_S\to T).
$$
This is a sum of ideals, with no assertion
that the sum is direct. The tensor quotient
map $T\to (R/N_R)\otimes_k(S/N_S)$ is
surjective and kills $J$. Its inverse
after passage to $T/J$ sends
$[r]\otimes[s]$ to $[r\otimes s]$.
Changing either representative changes
the tensor by an element of the indicated
ideal, so this is well-defined; balancing
and both composites follow on pure tensors.
Thus the kernel is exactly $J$.

Every element of $J$ is a finite sum
$z=z_1+\cdots+z_m$ of tensors having at
least one nilpotent factor. Choose positive
integers $d_i$ with $z_i^{d_i}=0$, and put
$D=1+\sum_{i=1}^m(d_i-1)$. Its full
multinomial expansion is
$$
z^D=
\sum_{\substack{a_1+\cdots+a_m=D\\a_i\geq0}}
\frac{D!}{a_1!\cdots a_m!}
z_1^{a_1}\cdots z_m^{a_m}=0.
$$
Indeed at least one $a_i\geq d_i$ in each
summand. The coefficients are the original
integers; no factorial is inverted in the
receiving ring. The empty sum is zero.
Hence $J\subseteq\operatorname{Nil}(T)$.
The proved quotient is reduced by the
first part, applied to $R/N_R$ and $S/N_S$.
Every nilpotent of $T$ therefore maps to
zero in that quotient and belongs to $J$,
proving the reverse inclusion and the
claimed formula with its exact quotient
isomorphism.
\end{proof}












\subsection{Editorial treatment of Algebra lines 10614--10630}

\hypertarget{algebra-context-125}{}

\noindent\href{https://github.com/stacks/stacks-project/blob/a04446e57ec1fbc252a871afcec7752fb2807b14/algebra.tex\#L10614}{Original source, lines 10614--10630}. Complete original and reviewed TeX passages accompany this note. Every intervention is separately recorded; the following treatment is editorial.


\begin{lemma}
\label{editorial-source-lemma-surjective-locally-nilpotent-kernel}
Let $\varphi : R \to S$ be a surjective map with locally nilpotent kernel.
Then $\varphi$ induces a homeomorphism of spectra and isomorphisms
on residue fields. For any ring map $R \to R'$ the ring map
$R' \to R' \otimes_R S$ is surjective with locally nilpotent kernel.
\end{lemma}

\begin{proof}
Put $I=\ker\varphi$ and retain the isomorphism
$R/I\to S$, $[r]\mapsto\varphi(r)$, supplied
by surjectivity. Every prime of $R$ contains
$I$, since each element of $I$ is nilpotent.
The inverse maps on primes are therefore
$$
\mathfrak q\longmapsto\varphi^{-1}(\mathfrak q),
\qquad \mathfrak p\longmapsto\varphi(\mathfrak p).
$$
The second is a prime ideal because its
quotient is the original domain $R/\mathfrak p$.
Surjectivity and $I\subseteq\mathfrak p$
show both composites are identities.
Moreover $D(\varphi(r))$ corresponds to
$D(r)$ for every original $r\in R$;
these sets form bases, so the bijection
and its inverse are continuous.

For corresponding primes $\mathfrak p$
and $\mathfrak q$, the residue-field map is
the explicit isomorphism
$$
\operatorname{Frac}(R/\mathfrak p)
\longrightarrow\operatorname{Frac}(S/\mathfrak q),
\qquad
\frac{r+\mathfrak p}{t+\mathfrak p}
\longmapsto
\frac{\varphi(r)+\mathfrak q}{\varphi(t)+\mathfrak q},
\quad t\notin\mathfrak p.
$$
The inverse lifts numerator and denominator
along $\varphi$; different lifts differ by
$I\subseteq\mathfrak p$, so the fractions
are unchanged. A nonzero denominator on
either side remains nonzero on the other.

For an arbitrary original map $f:R\to R'$,
there are inverse algebra maps
$$
R'\otimes_R S\longrightarrow R'/IR',\qquad
a\otimes\varphi(r)\longmapsto[af(r)],
\qquad
R'/IR'\longrightarrow R'\otimes_R S,\quad
[a]\longmapsto a\otimes1.
$$
Changing a lift $r$ of an element of $S$
adds an element of $I$, and balancing
uses the same $f(r)$, so the first map
is well-defined. The second kills $IR'$
because $f(i)\otimes1=1\otimes\varphi(i)=0$
for $i\in I$. The formulas prove both
composites are identities. Thus the
original map $R'\to R'\otimes_R S$ is
surjective with kernel exactly $IR'$.

Each element of that kernel is a finite
sum $\sum_{j=1}^m a_j f(i_j)$ with $i_j\in I$.
Choose $d_j\geq1$ with $i_j^{d_j}=0$.
In the full multinomial expansion of its
$D$th power, for $D=1+\sum_j(d_j-1)$,
every term contains some $f(i_j)^{d_j}$,
and is zero. Thus $IR'$ is locally
nilpotent, without a uniform exponent
being assumed for the whole ideal.
Applying the first part to this actual
quotient proves the asserted base-change
conclusions, with the same prime and
residue-field maps.

The kernel hypothesis is also necessary
for a surjective ring map to induce a
homeomorphism onto the entire spectrum:
its image is $V(I)$, so surjectivity on
primes requires $I$ to lie in every
prime, equivalently in the nilradical.
In addition, the given quotient induces
an exact isomorphism of reduced rings
$R_{red}\to S_{red}$. Indeed if
$\varphi(r)$ is nilpotent, then $r^n\in I$
for some $n$, and $(r^n)^m=0$ for some
$m$, so $r$ is nilpotent. The converse
is immediate under a ring homomorphism.
The preimage of $\operatorname{Nil}(S)$
is consequently exactly
$\operatorname{Nil}(R)$, proving the
stated reduced-quotient isomorphism and
the same assertion after each base change.
\end{proof}


\subsection{Editorial treatment of Algebra lines 10632--10715}

\hypertarget{algebra-context-126}{}

\noindent\href{https://github.com/stacks/stacks-project/blob/a04446e57ec1fbc252a871afcec7752fb2807b14/algebra.tex\#L10632}{Original source, lines 10632--10715}. Complete original and reviewed TeX passages accompany this note. Every intervention is separately recorded; the following treatment is editorial.


\begin{lemma}
\label{editorial-source-lemma-powers-field}
\begin{reference}
\cite[Lemma 3.1.6]{Alper-adequate}
\end{reference}
Let $k'/k$ be a field extension. The following are equivalent
\begin{enumerate}
\item for each $x \in k'$ there exists an $n > 0$ such that $x^n \in k$, and
\item $k' = k$ or $k$ and $k'$ have characteristic $p > 0$ and
either $k'/k$ is a purely inseparable extension or
$k$ and $k'$ are algebraic extensions of $\mathbf{F}_p$.
\end{enumerate}
\end{lemma}

\begin{proof}
Each case of (2) implies (1). For $k'=k$
take $n=1$. For a purely inseparable
extension in characteristic $p>0$, each
original element has a $p$-power in $k$.
If both fields are algebraic over
$\mathbf F_p$, a nonzero $x\in k'$
lies in the finite field $\mathbf F_p(x)$,
whose multiplicative group is finite;
some positive power of $x$ is $1\in k$.
For $x=0$ take $n=1$.

Assume (1) and retain the original extension
$k'/k$. Each $x\in k'$ is algebraic over
$k$, being a root of $T^n-x^n\in k[T]$.
Let $E$ be its maximal separable
subextension from Fields, Lemma
\ref{fields-lemma-separable-first}.
In positive characteristic, if $E=k$
then $k'/k$ is purely inseparable and
we have finished. In characteristic
zero $E=k'$, so $E=k$ is precisely the
identity case. Otherwise choose
$x\in E\setminus k$ and keep the finite
separable intermediate field $L=k(x)$.
Property (1) is inherited by $L/k$.
It is enough to show that $k$ has
positive characteristic and is algebraic
over its prime field: the original
$k'/k$ is already algebraic, and then
$k'$ is algebraic over the same prime
field by transitivity. This proves the
reduction for infinite $k'/k$ as well.

First let $z$ be any element of $L\setminus k$.
There are $k$-embeddings $\sigma,\tau$ of
$k(z)$ into an algebraic closure $\overline k$
with $A=\sigma(z)\neq B=\tau(z)$, because
its separable minimal polynomial has
degree at least two. This is the
embedding count of Fields, Lemma
\ref{fields-lemma-count-embeddings}.
Both $z$ and $z+1$ are nonzero, since
neither belongs to $k$. Choose the
original positive integers $n,m$ with
$z^n,(z+1)^m\in k$. Their images give
$A^n=B^n$ and $(A+1)^m=(B+1)^m$.
Thus, with all denominators nonzero,
$$
\zeta=A/B,\qquad \zeta'=(A+1)/(B+1),
\qquad \zeta^n=1,\quad(\zeta')^m=1.
$$
One has $\zeta'\neq1$ because $A\neq B$.
The full identity
$$
\zeta'(B+1)=A+1=\zeta B+1
$$
gives $B(\zeta'-\zeta)=1-\zeta'$.
Since the right side is nonzero,
$\zeta'-\zeta\neq0$, and hence
$$
B=\frac{1-\zeta'}{\zeta'-\zeta}.
$$
Both roots of unity are algebraic over
the prime field, so this formula makes
$B$ algebraic over that field. The
embedding $\tau$ fixes the prime field;
pulling back an annihilating polynomial
through the injective $\tau$ proves that
the original $z$ is algebraic over it.
This applies to every $z\in L\setminus k$.
For any $a\in k$, both $x$ and $x+a$
lie outside $k$ and are algebraic over
the prime field, so their difference
$a=(x+a)-x$ is algebraic. It follows
that $k$, and then the original $k'$,
are algebraic over the prime field.

It remains to exclude characteristic
zero. Use the fixed original $x\in L\setminus k$
and two embeddings as above, with
$A=\sigma(x)$, $B=\tau(x)$, and the
same $\zeta,\zeta'$. Both $A=\zeta B$
and $B=(1-\zeta')/(\zeta'-\zeta)$
belong to the number field
$F=\mathbf Q(\zeta,\zeta')$.
This field contains only finitely many
roots of unity, for the following exact
reason. Put $d=[F:\mathbf Q]$. The monic
minimal polynomial of a root of unity
in $F$ has degree $e\leq d$ and integer
coefficients: it divides some $T^N-1$
in $\mathbf Q[T]$, and the primitive
Gauss lemma gives a factor in $\mathbf Z[T]$
whose leading coefficient divides $1$.
Every complex root $\xi_i$ of that
minimal polynomial is also an $N$th
root of unity and has absolute value
$1$. Its coefficient of $T^{e-j}$ is
the full signed sum
$$
(-1)^j\sum_{1\leq i_1<\cdots<i_j\leq e}
\xi_{i_1}\cdots\xi_{i_j},
$$
so its absolute value is at most
$\binom ej$. There are finitely many
integer coefficient lists with these
bounds and $e\leq d$, and each resulting
nonzero polynomial has finitely many
roots. This proves the required
finiteness inside $F$.

For every original rational $y$, one
has $x+y\notin k$ and $B+y\neq0$.
Property (1) gives a positive exponent
for $x+y$, so
$$
\zeta_y=\frac{A+y}{B+y}
=\frac{\zeta B+y}{B+y}\in F
$$
is a root of unity, and $\zeta_y\neq1$.
Infinitely many rational $y$ and only
finitely many roots of unity in $F$
give distinct $y,y'$ with
$\zeta_y=\zeta_{y'}$. Subtracting the
two original identities gives
$$
y-y'=(A+y)-(A+y')
=\zeta_y\big((B+y)-(B+y')\big)
=\zeta_y(y-y').
$$
Since $y-y'\neq0$, this forces
$\zeta_y=1$, a contradiction. Thus the
nontrivial separable case has positive
characteristic, and the already proved
algebraicity gives exactly (2).

The argument also determines when a
single positive exponent can work for
all elements. Such an exponent exists
if and only if $k'=k$, or the fields
have characteristic $p>0$ and either
$k'/k$ is purely inseparable of bounded
exponent, or $k'$ is a finite field.
The converse is immediate with exponent
$1$, a common $p$-power, or $|k'|-1$,
respectively. For necessity retain a
single original exponent $N$ for all
elements. If a nontrivial separable
intermediate field exists, choose
$A=\sigma(x)\neq B=\tau(x)$ as above.
For every $a\in k$, the quotient
$(A+a)/(B+a)$ is an $N$th root of unity.
These quotients are distinct for
distinct $a$: cross multiplication
gives the full identity
$$
(A+a)(B+b)-(A+b)(B+a)=(A-B)(b-a).
$$
Their denominators are nonzero because
$x\notin k$. Hence $|k|\leq N$, as a
nonzero degree-$N$ polynomial has at
most $N$ roots. The field $k$ is finite.
Every element of $k'$ is then among
the roots of the finitely many original
polynomials $T^N-c$, $c\in k$. For
$c=0$ the only root is zero; for
each of the $|k|-1$ nonzero values
there are at most $N$ roots. Thus
$|k'|\leq1+N(|k|-1)$, and $k'$ is finite.

In the remaining nonidentity case
$k'/k$ is purely inseparable in
characteristic $p$. Write the original
$N=p^e m$ with $p\nmid m$. For every
nonzero $x\in k'$, put $u=x^{p^e}$.
It is purely inseparable over $k$ and
satisfies $u^m=x^N\in k$. The polynomial
$T^m-x^N$ has nonzero derivative at
each root, since $m$ and each root
are nonzero. Hence $u$ is also
separable over $k$, and belongs to
$k$. The zero element satisfies the
same conclusion directly. Therefore
$(k')^{p^e}\subseteq k$, giving a
uniform inseparability exponent from
the full original integer $N$.

For finite original fields the fixed-exponent
condition has an exact form. Put $q=|k|$
and $q'=|k'|$. The multiplicative group
of a finite field is cyclic, as follows
directly from the root bound. For a
finite subgroup $G$ of a field's
multiplicative group, let $h$ be the
least common multiple of its element
orders. For every prime $\ell$ dividing
$h$, choose an element whose order
has the largest $\ell$-power $\ell^{a_\ell}$
occurring in $h$, and raise that element
to the prime-to-$\ell$ part of its
order. The result has order exactly
$\ell^{a_\ell}$. The product of these
elements has order $h$: if its $m$th
power is $1$, raising to
$h/\ell^{a_\ell}$ shows
$\ell^{a_\ell}\mid m$ for each $\ell$,
since $h/\ell^{a_\ell}$ is coprime
to $\ell$. Conversely its $h$th power
is $1$. All elements of $G$ are
roots of $T^h-1$, so $|G|\leq h$;
the element of order $h$ gives
$|G|\geq h$. Thus it generates $G$.
For the trivial group the empty
product has order $1$, giving the
same conclusion.

Choose a generator $g$ of the original
$(k')^\times$, of order $q'-1$.
Inside $k'$, the nonzero roots of
$T^{q-1}-1$ are exactly $k^\times$:
all its $q-1$ elements are roots,
and a degree-$q-1$ polynomial has
no additional ones. Therefore
$$
\begin{split}
(\forall x\in k')\ x^N\in k
&\quad\Longleftrightarrow\quad g^N\in k^\times\\
&\quad\Longleftrightarrow\quad g^{N(q-1)}=1\\
&\quad\Longleftrightarrow\quad (q'-1)\mid N(q-1).
\end{split}
$$
The zero element satisfies the condition
for every positive $N$. If $d=[k':k]$,
then $q'=q^d$, by counting the coordinate
vectors of a $k$-basis. Retaining the
full factorization
$$
q'-1=q^d-1=(q-1)(1+q+\cdots+q^{d-1}),
$$
the displayed divisibility is equivalently
$(1+q+\cdots+q^{d-1})\mid N$. The least
positive common exponent is consequently
the full sum $1+q+\cdots+q^{d-1}$.
This also proves the sharp cardinality
bound above: equality holds for this
least exponent, with every contribution
from the original zero and nonzero
elements retained.
\end{proof}


\subsection{Editorial treatment of Algebra lines 10717--10767}

\hypertarget{algebra-context-127}{}

\noindent\href{https://github.com/stacks/stacks-project/blob/a04446e57ec1fbc252a871afcec7752fb2807b14/algebra.tex\#L10717}{Original source, lines 10717--10767}. Complete original and reviewed TeX passages accompany this note. Every intervention is separately recorded; the following treatment is editorial.


\begin{lemma}
\label{editorial-source-lemma-powers}
Let $\varphi : R \to S$ be a ring map. If
\begin{enumerate}
\item for any $x \in S$ there exists $n > 0$ such that
$x^n$ is in the image of $\varphi$, and
\item $\Ker(\varphi)$ is locally nilpotent,
\end{enumerate}
then $\varphi$ induces a homeomorphism on spectra and induces residue
field extensions satisfying the equivalent conditions of
Lemma \ref{editorial-source-lemma-powers-field}.
\end{lemma}

\begin{proof}
Write $I=\ker\varphi$. First use only
condition (1). Every original $x\in S$
has an equation $x^n-\varphi(y)=0$
for some $n>0$ and $y\in R$. This is
a monic equation over the actual image
$R/I$, so $R/I\to S$ is an integral
injection. Lemma
\ref{editorial-source-lemma-integral-overring-surjective}
therefore makes $\Spec(S)\to\Spec(R/I)$
surjective. Under the original quotient
map, $\Spec(R/I)$ identifies with
the closed subset $V(I)\subseteq\Spec(R)$.

To prove injectivity, suppose primes
$\mathfrak q,\mathfrak q'$ have the
same contraction $\mathfrak p$ and
$x\in\mathfrak q\setminus\mathfrak q'$.
Choose the original equation
$x^n=\varphi(y)$ from (1). It gives
$y\in\mathfrak p$, and therefore
$x^n\in\mathfrak q'$, contradicting
primality and $x\notin\mathfrak q'$.
Interchanging the two primes proves
equality. For this same equation,
the resulting bijection sends $D(x)$
exactly onto $D(y)\cap V(I)$: a prime
contains $x$ if and only if its
contraction contains $y$. Hence it
is open for the subspace topology
on $V(I)$, and is a homeomorphism
onto that closed subset, since the
original map on spectra is continuous.
Condition (2) says precisely that
$I\subseteq\operatorname{Nil}(R)$,
so $V(I)=\Spec(R)$ and gives the
homeomorphism in the statement.

Retain corresponding original primes
$\mathfrak q$ and $\mathfrak p$.
The map $R/\mathfrak p\to S/\mathfrak q$
is injective and thus extends to the
specified residue-field embedding.
Let $x$ in the receiving residue
field be represented by $y/z$ with
$y,z\in S$, $z\notin\mathfrak q$.
Choose original $a,b\in R$ and
positive $n,m$ with
$y^n=\varphi(a)$, $z^m=\varphi(b)$.
Then $b\notin\mathfrak p$, because
its image is $z^m\notin\mathfrak q$.
The full identity in $S_{\mathfrak q}$ is
$$
(y/z)^{nm}
=\frac{(y^n)^m}{(z^m)^n}
=\frac{\varphi(a)^m}{\varphi(b)^n}
=\varphi_{\mathfrak p,\mathfrak q}(a^m/b^n).
$$
Passing to residue fields proves
that the original $x^{nm}$ belongs
to the image of $\kappa(\mathfrak p)$.
Lemma \ref{editorial-source-lemma-powers-field} now
applies to these exact residue fields.
This part did not require (2).

Thus (1) alone gives a homeomorphism
onto $V(\ker\varphi)$ and the same
residue-field conclusion; under (1),
(2) is equivalent to surjectivity
on all primes of $R$. No universal
homeomorphism follows from these two
conditions alone. For an exact
counterexample take the original
inclusion
$$
R=\mathbf F_2\longrightarrow
S=\mathbf F_2[\alpha]/(\alpha^2+\alpha+1).
$$
The polynomial has no root in
$\mathbf F_2$, so $S$ is a field
with elements $0,1,\alpha,\alpha^2$,
and $\alpha^3=1$. Hence every
element of $S$ has cube in $R$,
and the kernel is zero. After the
original base change $R\to S$,
the tensor ring has the comparison
$$
S\otimes_R S\simeq S[T]/(T^2+T+1),
\qquad1\otimes\alpha\longleftrightarrow T,
$$
fixing $S$ through the first factor.
The bases $1,\alpha$ and $1,T$
prove this map and its inverse.
Retain the two distinct roots
$\alpha,\alpha^2$ and the full
factorization
$T^2+T+1=(T-\alpha)(T-\alpha^2)$.
Evaluation gives
$$
S[T]/(T^2+T+1)\longrightarrow S\times S,
\qquad[H]\longmapsto(H(\alpha),H(\alpha^2)).
$$
Its inverse sends $(a,b)$ to
$$
\left[
a\frac{T-\alpha^2}{\alpha-\alpha^2}
+b\frac{T-\alpha}{\alpha^2-\alpha}
\right].
$$
The original denominators are nonzero
in $S$; evaluating verifies both
coordinates, and a polynomial zero
at both roots is divisible by the
original product of the two distinct
linear factors. These are inverse
ring maps. The base-changed spectrum
therefore has two points mapping
to the single point of $\Spec(S)$,
and cannot be a homeomorphism.
\end{proof}


\subsection{Editorial treatment of Algebra lines 10769--10811}

\hypertarget{algebra-context-128}{}

\noindent\href{https://github.com/stacks/stacks-project/blob/a04446e57ec1fbc252a871afcec7752fb2807b14/algebra.tex\#L10769}{Original source, lines 10769--10811}. Complete original and reviewed TeX passages accompany this note. Every intervention is separately recorded; the following treatment is editorial.


\begin{lemma}
\label{editorial-source-lemma-2-3-ring-map}
Let $\varphi : R \to S$ be a ring map. Assume
\begin{enumerate}
\item[(a)] $S$ is generated as an $R$-algebra by elements $x$ such
that $x^2, x^3 \in \varphi(R)$, and
\item[(b)] $\Ker(\varphi)$ is locally nilpotent,
\end{enumerate}
Then $\varphi$ induces isomorphisms on residue fields and
a homeomorphism of spectra. For any ring map $R \to R'$
the ring map $R' \to R' \otimes_R S$ also satisfies (a) and (b).
\end{lemma}

\begin{proof}
For every original generator $x$, choose $a_x,b_x\in R$
with $\varphi(a_x)=x^2$ and $\varphi(b_x)=x^3$.
The monic equation $T^2-a_x$ proves integrality of
$x$ over $R$. Every element of the generated algebra
involves finitely many such generators, so the
integral-closure subring property proves that
$R\to S$ is integral. Put $I=\ker\varphi$.
The integral injection $R/I\to S$ is surjective
on spectra by Lemma \ref{editorial-source-lemma-integral-overring-surjective}.
Since $I$ is locally nilpotent, $V(I)=\Spec(R)$.
Thus the original map on spectra is surjective,
and is closed by Lemmas \ref{editorial-source-lemma-integral-going-up}
and \ref{editorial-source-lemma-going-up-closed}.

Let $\mathfrak q$ lie over $\mathfrak p$, and use
the original injection
$\kappa(\mathfrak p)\to\kappa(\mathfrak q)$.
If $\alpha$ is the image of $x$ in the receiving
field, then $\alpha^2=\overline a_x$ and
$\alpha^3=\overline b_x$, where the bars are the
images of these original coefficients in
$\kappa(\mathfrak p)$. If $\overline a_x=0$,
then $\alpha=0$. Otherwise
$\alpha=\overline b_x/\overline a_x$.
Consequently every generator belongs to the
image of the original base residue field.
Since $\kappa(\mathfrak q)=\operatorname{Frac}(S/\mathfrak q)$
is generated as a field by those images and
$\kappa(\mathfrak p)$, this residue-field map
is an isomorphism.

If $\mathfrak q'$ also lies over $\mathfrak p$,
compose the maps $S\to\kappa(\mathfrak q)$ and
$S\to\kappa(\mathfrak q')$ with the respective
inverses of the base residue-field isomorphisms.
Both resulting maps to $\kappa(\mathfrak p)$
agree on $R$ and send each original $x$ to
the same displayed value, zero or
$\overline b_x/\overline a_x$.
They agree on every polynomial in the
generators, hence on $S$. Their kernels
are respectively $\mathfrak q$ and
$\mathfrak q'$, so these primes are equal.
The original spectrum map is a continuous
closed bijection, hence a homeomorphism.

Retain any original $f:R\to R'$ and put
$S'=R'\otimes_R S$. It is generated over
$R'$ by $1\otimes x$, with
$$
(1\otimes x)^2=f(a_x)\otimes1,
\qquad (1\otimes x)^3=f(b_x)\otimes1.
$$
To verify the new kernel without assuming
$f$ flat, let $\mathfrak p'\subset R'$
contract to $\mathfrak p\subset R$.
Choose an original $\mathfrak q$ above
$\mathfrak p$, as is possible by surjectivity.
The tensor ring
$$
C=\kappa(\mathfrak p')\otimes_{\kappa(\mathfrak p)}
\kappa(\mathfrak q)
$$
is nonzero: a vector-space basis of either
nonzero field remains a basis after extending
scalars. The original maps give
$S'\to C$, $r'\otimes s\mapsto\overline r'\otimes\overline s$.
Take a prime of $C$. Its quotient is a
nonzero domain, and the induced unital
map from the field $\kappa(\mathfrak p')$
is injective. The inverse image of that
prime in $S'$ therefore contracts to
exactly $\mathfrak p'$ in $R'$.
Thus $\Spec(S')\to\Spec(R')$ is surjective,
and its kernel lies in every prime of
$R'$, hence consists of nilpotents.
This proves both base-change hypotheses;
the already proved argument gives the
same homeomorphism and residue-field
isomorphisms after every such base change.

The exponents $2,3$ can more generally
be replaced, separately for each original
generator $x$, by a finite nonempty list
of positive integers $d_1,\ldots,d_r$
with greatest common divisor $1$, provided
all $x^{d_i}$ belong to $\varphi(R)$.
One positive exponent supplies the monic
integral equation. Choose integers $u_i$
with $\sum_i u_i d_i=1$. For a nonzero
residue $\alpha$ of $x$, every $\alpha^{d_i}$
is a nonzero element of the base residue
field, and the full identity
$$
\alpha=\prod_{i=1}^r(\alpha^{d_i})^{u_i}
$$
puts it in that field, with negative
powers interpreted through its actual
field inverses. If any power is zero,
then $\alpha=0$. This formula determines
the same residue at every prime over
the given base prime. The preceding
integrality, kernel, and generator-map
arguments therefore prove the same
conclusions and their preservation by
arbitrary base change for these lists.

The greatest-common-divisor condition
cannot be weakened to an arbitrary
positive common divisor. For a list
with greatest common divisor $d>1$,
take the original field inclusion
$\mathbf Q(u^d)\to\mathbf Q(u)$, with
$u$ transcendental. The powers in the
list all belong to the smaller field.
The elements $1,u,\ldots,u^{d-1}$ are
independent over $\mathbf Q(u^d)$:
clearing the original denominators
leaves polynomial monomials with
distinct exponent residues modulo $d$.
They span the algebra generated by $u$
using $u^d\in\mathbf Q(u^d)$; that
finite-dimensional domain is a field,
since multiplication by a nonzero
element is injective and hence
surjective on its finite-dimensional
vector space. Thus it is the original
$\mathbf Q(u)$ and the residue-field
extension has degree exactly $d$,
not degree one.

Even under the original $2,3$ hypotheses,
an injective map can acquire a nonzero
nilpotent kernel after base change.
For any field $k$ let
$$
R=k[a,b]/(b^2-a^3)\longrightarrow S=k[t],
\qquad a\longmapsto t^2,\quad b\longmapsto t^3.
$$
This is injective: monic division in
$b$ writes an element uniquely as
$f(a)+b g(a)$, and its image
$f(t^2)+t^3g(t^2)$ has disjoint even
and odd exponent supports. Also $S$
is generated by the original $t$ with
$t^2,t^3$ in the image. Base change
by $R\to R/(a)=k[b]/(b^2)$ gives
$$
k[b]/(b^2)\longrightarrow k[t]/(t^2),
\qquad b\longmapsto t^3=0.
$$
Its kernel is exactly the nonzero
ideal $(b)$, whose square is zero.
The quotient-tensor comparison
$(R/(a))\otimes_R S\simeq S/(t^2)$
sends $[r]\otimes s$ to $[\varphi(r)s]$
and has inverse $[s]\mapsto1\otimes s$.
Thus this is the actual base-changed
map, with its nilpotent kernel fully
identified.
\end{proof}


\subsection{Editorial treatment of Algebra lines 10813--10845}

\hypertarget{algebra-context-129}{}

\noindent\href{https://github.com/stacks/stacks-project/blob/a04446e57ec1fbc252a871afcec7752fb2807b14/algebra.tex\#L10813}{Original source, lines 10813--10845}. Complete original and reviewed TeX passages accompany this note. Every intervention is separately recorded; the following treatment is editorial.


\begin{lemma}
\label{editorial-source-lemma-help-with-powers}
Let $p$ be a prime number. Let $n, m > 0$ be two integers. There exists
an integer $a$ such that
$(x + y)^{p^a}, p^a(x + y) \in \mathbf{Z}[x^{p^n}, p^nx, y^{p^m}, p^my]$.
\end{lemma}

\begin{proof}
Retain the original subring
$$
A=\mathbf Z[x^{p^n},p^n x,y^{p^m},p^m y]
\subseteq\mathbf Z[x,y].
$$
We calculate exactly which integer
coefficients can occur at each original
monomial $x^i y^j$. Write uniquely
$$
i=q p^n+r,\quad 0\leq r<p^n,
\qquad j=q' p^m+r',\quad 0\leq r'<p^m.
$$
In any product of the four displayed
generators contributing to this
monomial, the exponent on $p^n x$
is at least $r$, and the exponent
on $p^m y$ is at least $r'$.
Thus its coefficient is divisible
by $p^{nr+mr'}$. Conversely the
product
$$
(x^{p^n})^q(p^n x)^r
(y^{p^m})^{q'}(p^m y)^{r'}
=p^{nr+mr'}x^i y^j
$$
belongs to $A$. Since distinct formal
monomials are independent over
$\mathbf Z$, a polynomial belongs
to $A$ exactly when each of its
coefficients is divisible by this
corresponding integer. This proves
necessity as well as sufficiency
of the source divisibility test.

For a nonnegative integer $a$,
$p^a(x+y)$ belongs to $A$ exactly
when $a\geq b=\max(n,m)$: the
coefficient tests for $x$ and $y$
are respectively divisibility by
$p^n$ and $p^m$. In that case its
full expression in the generators is
$$
p^a(x+y)=p^{a-n}(p^n x)+p^{a-m}(p^m y).
$$
Now keep $a\geq b$ and $N=p^a$.
For $0<i<N$, the original binomial
coefficient has the exact valuation
$$
v_p\binom Ni=a-v_p(i).
$$
Indeed
$\binom Ni=(N/i)\binom{N-1}{i-1}$,
and the second factor is not divisible
by $p$. To verify the latter assertion,
reduce the full polynomial identity
$(1+T)^N=1+T^N$ modulo $p$ and divide
by $1+T$ in $\mathbf F_p[T]$.
It gives
$$
(1+T)^{N-1}=\sum_{h=0}^{N-1}(-T)^h
\quad\hbox{in }\mathbf F_p[T].
$$
Multiplication by $1+T$ verifies this
identity: the last sign is positive
for odd $p$, and signs coincide for
$p=2$. Every coefficient is nonzero,
giving the valuation assertion with
the original integer coefficient
still retained.

The two endpoint terms $i=0,N$
already belong to $A$. For an
intermediate $i$, put $t=i\bmod p^b$
with $0\leq t<p^b$. If $t=0$, both
remainders $r,r'$ are zero. If $t>0$,
then $v_p(i)=v_p(t)$ and
$$
r=t\bmod p^n,\qquad
r'=(-t)\bmod p^m,
$$
where both residues are chosen in
their original nonnegative ranges.
Every $1\leq t<p^b$ occurs among
$1\leq i<N$. Consequently the least
possible $a$ for both assertions
is exactly
$$
\max\left(b,
\max_{1\leq t<p^b}
\left(v_p(t)+n(t\bmod p^n)+m((-t)\bmod p^m)\right)\right).
$$
This is an exact test on the
original integer coefficients,
not merely a sufficient estimate.

It has the following explicit value:
$$
a_{min}=
\begin{cases}
n p^n+n-1,&n=m,\\
b(p^b-1)+c,&n\neq m,\quad
b=\max(n,m),\ c=\min(n,m).
\end{cases}
$$
For $n=m$, each nonzero residue $t$
has $r+r'=p^n$, and its valuation
is at most $n-1$, with equality at
$t=p^{n-1}$. This gives the first
value, which is at least $b$.
For $n<m$, write $t=q p^n+r$.
If $r>0$, its full weight is
$$
m p^m-m q p^n-(m-n)r+v_p(r).
$$
For every positive $r$ one has
$v_p(r)\leq r-1$: if $v=v_p(r)$,
then $r\geq p^v\geq v+1$, the
last inequality following by
induction on $v$. Since $m-n\geq1$,
this gives
$v_p(r)\leq(m-n)(r-1)$.
The weight is therefore at most
$m p^m-(m-n)$, and equality is
attained at the original $t=1$.
If $r=0$ and $t>0$, its weight is
$m(p^m-t)+v_p(t)$, at most
$m p^m-(m-1)t-1$. Since $t\geq1$,
this is at most $m p^m-m$, and
hence at most $m p^m-(m-n)$.
Thus the maximum is exactly
$m(p^m-1)+n$. For $m<n$, the
automorphism of $\mathbf Z[x,y]$
exchanging the original $x$ and
$y$ takes $A$ to the corresponding
subring with $n,m$ exchanged, and
fixes both $(x+y)^{p^a}$ and
$p^a(x+y)$. Applying this actual
inverse-to-itself map proves the
other unequal case. These values
are all at least $b$.

For every $a\geq a_{min}$, the full
desired expression is therefore
$$
(x+y)^{p^a}=
\sum_{\substack{i,j\geq0\\i+j=p^a}}
\frac{\binom{p^a}{i,j}}{p^{nr+mr'}}
(x^{p^n})^q(p^n x)^r
(y^{p^m})^{q'}(p^m y)^{r'}.
$$
Each displayed quotient is an
integer by the proved divisibility;
no power of $p$ is inverted in the
ring. All original binomial
coefficients and all four generator
factors are retained. If $b\leq a<a_{min}$,
the maximizing residue just described
gives an actual monomial whose
coefficient fails the necessary
divisibility. If $a<b$, the linear
expression already fails. This
proves minimality.

In particular the source choice
$a_0=n p^n+m p^m+n+m$ works.
For $n=m$, its difference from
$a_{min}$ is $n p^n+n+1$.
For $n\neq m$, with the same
$b,c$, the difference is $c p^c+2b$.
Both are positive. Thus the original
sufficient choice is preserved and
verified, alongside the sharp bound.
\end{proof}


\subsection{Editorial treatment of Algebra lines 10847--10865}

\hypertarget{algebra-context-130}{}

\noindent\href{https://github.com/stacks/stacks-project/blob/a04446e57ec1fbc252a871afcec7752fb2807b14/algebra.tex\#L10847}{Original source, lines 10847--10865}. Complete original and reviewed TeX passages accompany this note. Every intervention is separately recorded; the following treatment is editorial.


\begin{lemma}
\label{editorial-source-lemma-p-ring-map-field}
Let $k'/k$ be a field extension. Let $p$ be a prime number.
The following are equivalent
\begin{enumerate}
\item $k'$ is generated as a field extension of $k$ by elements
$x$ such that there exists an $n > 0$ with $x^{p^n} \in k$ and
$p^nx \in k$, and
\item $k = k'$ or the characteristic of $k$
and $k'$ is $p$ and $k'/k$ is purely inseparable.
\end{enumerate}
\end{lemma}

\begin{proof}
The original field embedding preserves $1$, hence
the characteristic of $k$ and $k'$ is the
same. If that characteristic differs from
the specified prime $p$, the image of
$p^n$ is a nonzero field element and
therefore a unit of $k$. Thus
$x=(p^n x)/p^n\in k$ for every generator
in (1). Those generators generate $k'/k$,
so $k'=k$.

If the common characteristic is $p$,
each such generator has a $p$-power
in the original $k$. Every element
of the generated field is a rational
expression in finitely many original
generators. Choose an exponent at
least all their witnessing exponents.
Applying Frobenius to that full
rational expression raises its
coefficients and each generator to
elements of $k$; its denominator
remains nonzero in the field.
Thus every element of $k'$ has a
$p$-power in $k$, proving pure
inseparability without a finiteness
assumption on the whole extension.

Conversely, if $k'=k$, each element
itself satisfies (1) with $n=1$.
If the extension is purely
inseparable in characteristic $p$,
each original element has a
$p^e$th power in $k$. Choose
$n=\max(1,e)$; then its $p^n$th
power is still in $k$ and
$p^n x=0\in k$. The set of all
original elements can serve as
the generating set, proving (1).
The proof also shows exactly
which condition is needed for
fields: outside characteristic
$p$ the integer-multiple condition
alone forces the identity extension;
in characteristic $p$ it is automatic,
and the power condition is sufficient.
\end{proof}


\subsection{Editorial treatment of Algebra lines 10867--10910}

\hypertarget{algebra-context-131}{}

\noindent\href{https://github.com/stacks/stacks-project/blob/a04446e57ec1fbc252a871afcec7752fb2807b14/algebra.tex\#L10867}{Original source, lines 10867--10910}. Complete original and reviewed TeX passages accompany this note. Every intervention is separately recorded; the following treatment is editorial.


\begin{lemma}
\label{editorial-source-lemma-p-ring-map}
Let $\varphi : R \to S$ be a ring map. Let $p$ be a prime number. Assume
\begin{enumerate}
\item[(a)] $S$ is generated as an $R$-algebra by elements $x$ such
that there exists an $n > 0$ with $x^{p^n} \in \varphi(R)$ and
$p^nx \in \varphi(R)$, and
\item[(b)] $\Ker(\varphi)$ is locally nilpotent,
\end{enumerate}
Then $\varphi$ induces a homeomorphism of spectra and induces
residue field extensions satisfying the equivalent conditions
of Lemma \ref{editorial-source-lemma-p-ring-map-field}. For any ring map $R \to R'$
the ring map $R' \to R' \otimes_R S$ also satisfies (a) and (b).
\end{lemma}

\begin{proof}
Retain the actual image $A=\varphi(R)\subseteq S$.
Let $T$ consist of the original $x\in S$ for
which some positive $n$ satisfies
$x^{p^n}\in A$ and $p^n x\in A$.
Every element of $A$ lies in $T$, by taking
$n=1$. If $x\in T$, then $-x\in T$ with
the same $n$, since both required expressions
change only by the respective integer signs.
For $r\in R$,
$$
(\varphi(r)x)^{p^n}=\varphi(r)^{p^n}x^{p^n}\in A,
\qquad
p^n\varphi(r)x=\varphi(r)(p^n x)\in A.
$$
Thus this verification never assumes that
the original $\varphi$ is injective.
If $x,y\in T$ have the original witnesses
$n,m$, then
$$
(xy)^{p^{n+m}}
=(x^{p^n})^{p^m}(y^{p^m})^{p^n}\in A,
\qquad
p^{n+m}xy=(p^n x)(p^m y)\in A.
$$
For addition, apply the actual evaluation
homomorphism $\mathbf Z[X,Y]\to S$,
$X\mapsto x$, $Y\mapsto y$, to Lemma
\ref{editorial-source-lemma-help-with-powers}. Its four
original generators map to
$x^{p^n},p^n x,y^{p^m},p^m y\in A$.
The displayed full coefficient formulas
in that lemma give both required expressions
for $x+y$ in $A$, with its specified
integer exponent. Therefore $T$ is an
$R$-subalgebra of $S$. Hypothesis (a)
places every generating element in $T$,
so $T=S$.

In particular every original element
of $S$ has a positive power in $A$.
Together with the nil-kernel hypothesis,
Lemma \ref{editorial-source-lemma-powers} proves that
the original map on spectra is a
homeomorphism. Let $\mathfrak q$ lie
over $\mathfrak p$. The field
$\kappa(\mathfrak q)$ is generated over
$\kappa(\mathfrak p)$ by the residues
of the original algebra generators:
the images generate $S/\mathfrak q$
as an $R/\mathfrak p$-algebra, and
fractions give its residue field.
Each such residue still satisfies
both original power and integer-multiple
conditions. Lemma \ref{editorial-source-lemma-p-ring-map-field}
therefore proves precisely the asserted
residue-field classification.

For any original $R\to R'$, the
algebra $S'=R'\otimes_R S$ has the
generators $1\otimes x$. If the
original preimages of $x^{p^n}$ and
$p^n x$ are $a,b\in R$, then
$$
(1\otimes x)^{p^n}=f(a)\otimes1,
\qquad p^n(1\otimes x)=f(b)\otimes1.
$$
Thus (a) is preserved with exactly
the same $n$ for each generator.
Surjectivity on spectra is preserved
as well: for $\mathfrak p'\subset R'$
contracting to $\mathfrak p$, choose
$\mathfrak q\subset S$ above
$\mathfrak p$ and use the nonzero
field tensor
$\kappa(\mathfrak p')\otimes_{\kappa(\mathfrak p)}\kappa(\mathfrak q)$.
The map $r'\otimes s\mapsto
\overline r'\otimes\overline s$
and any prime of that tensor have
inverse image in $S'$ contracting
to exactly $\mathfrak p'$, since
a field injects into every nonzero
unital algebra over it. This is
the same explicit fibre construction
as in Lemma \ref{editorial-source-lemma-2-3-ring-map}.
Every element of the new kernel
therefore lies in every prime of
$R'$ and is nilpotent, proving (b).
The theorem applies to this actual
base-changed map, giving a
homeomorphism and the stated
residue-field extensions after
every base change.

The integer-multiple hypothesis
also gives an exact global conclusion
after inverting the original prime
integer $p$. The map
$$
R[1/p]\longrightarrow S[1/p]
$$
is surjective: each original
generator with $p^n x=\varphi(b)$
is the image of the actual fraction
$b/p^n$. Its kernel is exactly
$I[1/p]$, for $I=\ker\varphi$.
Indeed if $r/p^j$ maps to zero,
there is an integer $h\geq0$
with $p^h\varphi(r)=0$, and then
$p^h r\in I$ and
$r/p^j=(p^h r)/p^{j+h}$.
The reverse containment is immediate.
Thus
$S[1/p]\simeq R[1/p]/I[1/p]$
through this original fraction map,
and the kernel remains locally
nilpotent. In particular the
localized map is an isomorphism
if the original $\varphi$ is
injective. This statement also
covers the case that inverting
$p$ makes both rings zero.
\end{proof}


\subsection{Editorial treatment of Algebra lines 10912--10966}

\hypertarget{algebra-context-132}{}

\noindent\href{https://github.com/stacks/stacks-project/blob/a04446e57ec1fbc252a871afcec7752fb2807b14/algebra.tex\#L10912}{Original source, lines 10912--10966}. Complete original and reviewed TeX passages accompany this note. Every intervention is separately recorded; the following treatment is editorial.


\begin{lemma}
\label{editorial-source-lemma-radicial}
Let $\varphi : R \to S$ be a ring map. Assume
\begin{enumerate}
\item $\varphi$ induces an injective map of spectra,
\item $\varphi$ induces purely inseparable residue field extensions.
\end{enumerate}
Then for any ring map $R \to R'$ properties (1) and (2) are true for
$R' \to R' \otimes_R S$.
\end{lemma}

\begin{proof}
Retain the original base-change algebra
$S'=R'\otimes_R S$ and its two structural
maps. Let $\mathfrak p'\subset R'$ contract
to $\mathfrak p\subset R$, and set
$K=\kappa(\mathfrak p)$ and
$K'=\kappa(\mathfrak p')$ with the original
field embedding $K\to K'$. The original
fibre algebra is $F=K\otimes_R S$.
By Remark \ref{algebra-remark-fundamental-diagram},
its primes correspond to the primes
of $S$ over $\mathfrak p$, with their
original residue-field maps.

If there is no such prime, $F$ has
no prime and is zero, since every
nonzero unital ring has a maximal
ideal. The new fibre is zero as
well by the explicit comparison
below. Otherwise there is exactly
one original prime $\mathfrak q$
above $\mathfrak p$. Let $\mathfrak r$
be its corresponding prime in $F$.
It is the only prime of $F$, hence
is both its nilradical and its
unique maximal ideal. Consequently
$F/\mathfrak r$ is a field, and
the canonical map $F\to\kappa(\mathfrak r)$
is surjective with kernel $\mathfrak r$.
The original fibre correspondence
identifies this field with
$E=\kappa(\mathfrak q)$ over $K$.
By hypothesis $E/K$ is purely
inseparable. This proves the
needed nil-kernel quotient without
assuming that $F$ itself is reduced.

There is an original fibre isomorphism
$$
K'\otimes_{R'} S'\longrightarrow
K'\otimes_K F,\qquad
c\otimes(r'\otimes s)\longmapsto
c\overline r'\otimes(1\otimes s).
$$
Its inverse sends
$c\otimes(a\otimes s)$ to
$ca\otimes(1\otimes s)$, using the
specified embedding $K\to K'$.
Balancing follows from the original
commutative scalar diagram, and
these formulas compose to the
identity on every pure tensor.
Write $H=K'\otimes_K F$.
Lemma \ref{editorial-source-lemma-surjective-locally-nilpotent-kernel}
applied to $F\to E$ gives a
surjection
$H\to K'\otimes_K E$ with locally
nilpotent kernel.

Embed the purely inseparable $E$
in an algebraic closure of $K'$
by the unique roots compatible
with $K$. The full tensor comparison
proved in Lemma \ref{editorial-source-lemma-make-separable}
gives
$$
(K'\otimes_K E)_{red}\simeq C=K'E,
\qquad [\sum_i c_i\otimes e_i]
\longmapsto\sum_i c_i e_i.
$$
It applies also when $E/K$ is
infinite; each individual tensor
uses only finitely many root
exponents. The field $C/K'$ is
purely inseparable, because every
finite sum or quotient in the
compositum has a sufficiently high
$p$-power in $K'$ in positive
characteristic; in characteristic
zero $E=K$ and $C=K'$.
Both kernels in the composite
$H\to K'\otimes_K E\to C$ are nil.
Its full kernel is nil as well:
if an element maps to zero, some
power lies in the first nil kernel,
and a further power is zero.
Thus $H_{red}\simeq C$. The fibre
$H$ has exactly one prime and its
residue field is the original $C$.
The fibre correspondence proves
both required base-change properties.

In particular we obtain the exact
prime and residue field, rather
than only their uniqueness. Whenever
$\mathfrak q$ exists, the unique
$\mathfrak q'\subset S'$ over
$\mathfrak p'$ is the kernel of
$$
S'\longrightarrow C,\qquad
r'\otimes s\longmapsto
\overline r'\,\overline s,
$$
with the two factors embedded in
$K'$ and $E$ as above. Its contraction
to $S$ is $\mathfrak q$ and its
contraction to $R'$ is $\mathfrak p'$.
The fraction field of its image
is exactly $C$: it contains the
fractions giving $K'$ and $E$,
and is contained in their compositum.
This identifies $\kappa(\mathfrak q')$
with $C$ through the original maps.
It also proves that the image on
spectra after base change is exactly
the inverse image of the original
image on spectra.

Conversely, these two original
hypotheses are necessary for
injectivity on spectra after every
base change. Injectivity before
base change follows by choosing
$R'=R$. For an original
$\mathfrak q$ over $\mathfrak p$,
use the same $K,E,F$ above; injectivity
makes $F\to E$ a surjective nil-kernel
map. If $E/K$ were not purely
inseparable, the ring $E\otimes_K E$
would have at least two primes, as
we now prove with the original maps.
Its multiplication map onto $E$
has a prime kernel containing every
$z\otimes1-1\otimes z$.
If some $z\in E$ is transcendental
over $K$, no positive power of
this difference is zero: in its
full binomial expansion, the
coefficient of $1\otimes z^n$
is $(-1)^n$, and the powers
$1,z,\ldots,z^n$ remain independent
after extension from $K$ to $E$.

Otherwise $E/K$ is algebraic but
not purely inseparable. The
separable-first theorem supplies
$z\in E\setminus K$ separable over
$K$. Let $P(T)$ be its original
monic minimal polynomial, of degree
at least two. The ring
$E\otimes_K K(z)\simeq E[T]/(P)$
is reduced because $P$ is separable,
by the original polynomial and
Bezout argument in Lemma
\ref{editorial-source-lemma-separable-extension-preserves-reducedness}.
The element corresponding to $T-z$
is nonzero, by monic division,
and so is not nilpotent. The map
$E\otimes_K K(z)\to E\otimes_K E$
is injective by extension of a
vector-space basis. Thus the same
diagonal difference is not nilpotent
in the larger tensor ring either.
In both cases some prime avoids
that difference, while the original
multiplication kernel contains it.
They are two distinct primes.

Now take the actual base change
$R\to E$ through $R\to K\to E$.
Its algebra is
$E\otimes_R S\simeq E\otimes_K F$.
The surjection from this algebra
to $E\otimes_K E$ pulls the two
distinct primes back to distinct
primes, both over the unique prime
of the field $E$. This contradicts
universal injectivity. Hence every
original residue-field extension
must be purely inseparable, proving
the converse and the exact
characterization.
\end{proof}


\subsection{Editorial treatment of Algebra lines 10968--10986}

\hypertarget{algebra-context-133}{}

\noindent\href{https://github.com/stacks/stacks-project/blob/a04446e57ec1fbc252a871afcec7752fb2807b14/algebra.tex\#L10968}{Original source, lines 10968--10986}. Complete original and reviewed TeX passages accompany this note. Every intervention is separately recorded; the following treatment is editorial.


\begin{lemma}
\label{editorial-source-lemma-radicial-integral}
Let $\varphi : R \to S$ be a ring map. Assume
\begin{enumerate}
\item $\varphi$ is integral,
\item $\varphi$ induces an injective map of spectra,
\item $\varphi$ induces purely inseparable residue field extensions.
\end{enumerate}
Then $\varphi$ induces a homeomorphism from $\Spec(S)$ onto a closed
subset of $\Spec(R)$ and for any ring map
$R \to R'$ properties (1), (2), (3) are true for $R' \to R' \otimes_R S$.
\end{lemma}

\begin{proof}
Put $I=\ker\varphi$. The original integral
injection $R/I\to S$ gives image $V(I)$
on spectra by lying over. More generally,
for every ideal $J\subseteq S$, the
map $R/\varphi^{-1}(J)\to S/J$ is an
integral injection, since its monic
equations are the images of the
original ones. Its spectrum is
surjective, so the image of $V(J)$
is exactly $V(\varphi^{-1}(J))$.
Thus the original spectrum map is
closed. Its assumed injectivity
makes it a homeomorphism onto
the closed subset $V(I)$: closed
subsets of the domain have closed
images in this subspace, which
proves continuity of the inverse.

For any original $R\to R'$, write
$S'=R'\otimes_R S$. Each original
monic equation for $s\in S$ maps
to the same monic equation for
$1\otimes s$ with coefficients
mapped into $R'$. Those elements
generate $S'$ over $R'$, so
integrality is preserved, including
infinite generating families by
the finite-support subring argument.
Lemma \ref{editorial-source-lemma-radicial} supplies
injectivity on spectra and pure
inseparability of residue fields
for this exact base-changed map.
The preceding closed-image proof
therefore applies to it as well.

There is also an exact radical
comparison for its kernel $I'$:
$$
\sqrt{I'}=\sqrt{IR'}.
$$
Indeed the original image is
$V(I)$, and the explicit fibre
construction in Lemma
\ref{editorial-source-lemma-radicial} identifies the
new image with its inverse image
in $\Spec(R')$, namely $V(IR')$.
Integrality also identifies the
new image with $V(I')$. Equality
of these closed sets gives the
stated equality of radicals.
The same formula in fact holds
for every integral ring map:
for a prime of $R'$ containing
$IR'$, its contraction contains
$I$, so lying over supplies a
prime of $S$ above that contraction;
the nonzero tensor of the two
original residue fields supplies
a prime after base change. This
proves equality of the images
without requiring uniqueness.
Equality of the kernels themselves
need not hold: the explicit cusp
base change in Lemma
\ref{editorial-source-lemma-2-3-ring-map} has $I=0$
and $I'=(b)\neq0$ in $k[b]/(b^2)$,
while both radicals are $(b)$.
All comparisons concern the
original extended ideal $IR'$.
\end{proof}


\subsection{Editorial treatment of Algebra lines 10988--11003}

\hypertarget{algebra-context-134}{}

\noindent\href{https://github.com/stacks/stacks-project/blob/a04446e57ec1fbc252a871afcec7752fb2807b14/algebra.tex\#L10988}{Original source, lines 10988--11003}. Complete original and reviewed TeX passages accompany this note. Every intervention is separately recorded; the following treatment is editorial.


\begin{lemma}
\label{editorial-source-lemma-radicial-integral-bijective}
Let $\varphi : R \to S$ be a ring map. Assume
\begin{enumerate}
\item $\varphi$ is integral,
\item $\varphi$ induces a bijective map of spectra,
\item $\varphi$ induces purely inseparable residue field extensions.
\end{enumerate}
Then $\varphi$ induces a homeomorphism on spectra and for any ring map
$R \to R'$ properties (1), (2), (3) are true for $R' \to R' \otimes_R S$.
\end{lemma}

\begin{proof}
By Lemma \ref{editorial-source-lemma-radicial-integral}, the
original integral injective map on spectra
is a homeomorphism onto its image.
Its assumed bijectivity makes that image
the entire $\Spec(R)$. For any original
$R\to R'$, the same lemma preserves
integrality, injectivity on spectra,
and purely inseparable residue fields.
Surjectivity on spectra is preserved
by the original nonzero residue-field
tensor construction of Lemma
\ref{editorial-source-lemma-surjective-spec-radical-ideal}.
Thus the receiving map is bijective
on spectra and still integral, hence
a homeomorphism by the closed-image
argument, with exactly the asserted
residue fields. Equivalently, for
an integral map already satisfying
the injectivity and residue-field
hypotheses, full universal surjectivity
is the condition $\ker\varphi$
locally nilpotent: its image is
the exact $V(\ker\varphi)$, and
the same holds after every base
change by the radical formula.
\end{proof}


\subsection{Editorial treatment of Algebra lines 11005--11061}

\hypertarget{algebra-context-135}{}

\noindent\href{https://github.com/stacks/stacks-project/blob/a04446e57ec1fbc252a871afcec7752fb2807b14/algebra.tex\#L11005}{Original source, lines 11005--11061}. Complete original and reviewed TeX passages accompany this note. Every intervention is separately recorded; the following treatment is editorial.


\begin{lemma}
\label{editorial-source-lemma-universally-bijective}
Let $\varphi : R \to S$ be a ring map such that
\begin{enumerate}
\item the kernel of $\varphi$ is locally nilpotent, and
\item $S$ is generated as an $R$-algebra by elements $x$
such that there exist $n > 0$ and a polynomial $P(T) \in R[T]$
whose image in $S[T]$ is $(T - x)^n$.
\end{enumerate}
Then $\Spec(S) \to \Spec(R)$ is a homeomorphism and $R \to S$
induces purely inseparable extensions of residue fields.
Moreover, conditions (1) and (2) remain true on arbitrary base change.
\end{lemma}

\begin{proof}
If $S=0$, the nil-kernel hypothesis
makes $1\in R$ nilpotent, so $R=0$.
Every unital base change of the zero
ring is again zero, and all assertions
follow. Hence assume $S\neq0$.
Retain the original map $\varphi:R\to S$
and its kernel $I$. For every original
generator $x$ retain the given positive
integer $n$ and the entire original
polynomial $P(T)=\sum_{j=0}^d c_jT^j$,
whose image is $(T-x)^n$. Coefficient
comparison gives $d\geq n$,
$c_n-1\in I$, and $c_j\in I$ for
$j>n$. The polynomial
$Q(T)=T^n+\sum_{j=0}^{n-1}c_jT^j$
is a monic annihilating polynomial
for $x$. Indeed the full difference
$$
P(T)-Q(T)=(c_n-1)T^n+
\sum_{j=n+1}^d c_jT^j
$$
has all coefficients in the original
kernel, so $Q$ has the same image
$(T-x)^n$. This retains every
coefficient of $P$ and proves
integrality without presuming that
the given lift itself is monic.
All generators, hence $S$, are
integral over $R$. The integral
injection $R/I\to S$ is surjective
on spectra. Since $I$ is nil,
$V(I)=\Spec(R)$, so the original
spectrum map is surjective; it is
also closed by integrality.

We first prove the global scalar
identity used in the fibre argument.
Write
$$
e_j=\binom nj x^j
=\varphi((-1)^j c_{n-j})\in\varphi(R),
\qquad 1\leq j\leq n.
$$
All signs here are the actual signs
in the given $(T-x)^n$. For
$1\leq h\leq n$, the full Newton
recurrence is
$$
n x^h=
\sum_{j=1}^{h-1}(-1)^{j-1}e_j\,n x^{h-j}
+(-1)^{h-1}h e_h.
$$
It follows directly on substituting
the original $e_j$: Pascal's identity
telescopes to
$\sum_{j=0}^{h-1}(-1)^j\binom nj
=(-1)^{h-1}\binom{n-1}{h-1}$,
and $n\binom{n-1}{h-1}=h\binom nh$.
These two integer identities give
exactly the coefficient $n$ on the
right, with every displayed summand
retained. Induction on $h$, starting
with $n x=e_1$, proves $n x^h\in\varphi(R)$.
Explicit preimages are supplied by
the same recurrence in $R$, taking
$E_j=(-1)^j c_{n-j}$ and
$$
U_h=\sum_{j=1}^{h-1}(-1)^{j-1}E_j U_{h-j}
+(-1)^{h-1}h E_h.
$$
The empty sum for $h=1$ gives
$U_1=E_1$, and $\varphi(U_h)=n x^h$.

Fix a prime integer $p$ and write
the original $n=p^e m$, with
$p\nmid m$, retaining $q=p^e$.
One has the exact integer identity
$$
\gcd\left(n,\binom nq\right)=m.
$$
For a prime $\ell\neq p$ dividing
$n$, the identity
$\binom nq=(n/q)\binom{n-1}{q-1}$
shows its $\ell$-adic valuation is
at least that of $n$, since
$q$ is prime to $\ell$.
For $p$, the factor $n/q=m$
is prime to $p$, and
$\binom{n-1}{q-1}$ is nonzero
modulo $p$. In fact
$(1+T)^{qm}=(1+T^q)^m$ modulo
$p$, so dividing by $1+T$ shows
that its degree-$q-1$ coefficient
is $(-1)^{q-1}=1$ in
$\mathbf F_p$. This also treats
$q=1$. These valuations prove
the displayed greatest common
divisor, including every prime
factor of the original $n$.
Choose integers $A,B$ with
$An+B\binom nq=m$. Then
$$
\varphi(AU_q+BE_q)
=A n x^q+B\binom nq x^q=m x^q.
$$
Thus the source assertion
$m x^{p^e}\in\varphi(R)$ is
indeed true globally, not merely
after taking a residue field.
The recurrence and this integer
identity supply its missing proof.

Now let $\mathfrak p\subset R$
and retain $K=\kappa(\mathfrak p)$
and the original fibre
$C=K\otimes_R S$. It is nonzero
by surjectivity on spectra. Thus
the map from the field $K$ to $C$
is injective. If $K$ has
characteristic zero, each generator
$y=1\otimes x$ satisfies
$n y=-\overline c_{n-1}$.
The original positive $n$ is
invertible in $K$, so
$y=-n^{-1}\overline c_{n-1}$
belongs to the image of $K$.
All generators do, and hence
$K\to C$ is an isomorphism.

If $K$ has characteristic $p>0$,
use the original decomposition
$n=q m$, $q=p^e$, for each
generator separately. The full
fibre polynomial is
$$
(T-y)^n=(T^q-y^q)^m
=\sum_{j=0}^m\binom mj(-1)^j
T^{q(m-j)}y^{qj}.
$$
Its coefficient of $T^{n-q}$
is $-m y^q=\overline c_{n-q}$.
Since $m$ is a unit of $K$,
$y^q=-m^{-1}\overline c_{n-q}$
is in the original image of $K$.
This is the direct fibre formula;
the global preimage proved above
gives the same conclusion before
reducing coefficients.
If $e=0$, use the exponent $1$
in Lemma \ref{editorial-source-lemma-p-ring-map};
otherwise use $e$. In either
case a positive $p$-power of $y$
is in $K$, and the corresponding
integer multiple is zero in
characteristic $p$. Thus that
lemma applies to $K\to C$, whose
kernel is zero. The fibre has
one prime and its residue field
is purely inseparable over $K$.
The same conclusion holds in
characteristic zero by the
already proved isomorphism.
Remark \ref{algebra-remark-fundamental-diagram}
now gives bijectivity on spectra
and the asserted original residue
extensions. A continuous closed
bijection is a homeomorphism.

Under any original $f:R\to R'$,
the image polynomial
$P_f(T)=\sum_j f(c_j)T^j$ has
image $(T-1\otimes x)^n$ in
$(R'\otimes_R S)[T]$, with all
coefficients and the same $n$
retained. These elements generate
the base-changed algebra, so
condition (2) persists. For
condition (1), the original
surjectivity on spectra persists
under base change by the explicit
nonzero residue-field tensor
construction in Lemma
\ref{editorial-source-lemma-2-3-ring-map}. Hence
the new kernel lies in every
prime and is locally nilpotent.
Both hypotheses, and therefore
all conclusions, hold for the
actual base-changed map.

The scalar identities give more
global information for a finite
original generating list
$x_1,\ldots,x_r$ with integers
$n_1,\ldots,n_r$. Put
$N=\prod_{i=1}^r n_i$, taking
$N=1$ for the empty list.
Since $n_i x_i$ has the
specified preimage $U_{i,1}$,
the map $R[1/N]\to S[1/N]$
is surjective, sending the
original fraction
$U_{i,1}/n_i=((N/n_i)U_{i,1})/N$
to $x_i$. Its kernel is exactly
$I[1/N]$: the fraction-witness
proof is the same as for
inverting $p$ in Lemma
\ref{editorial-source-lemma-p-ring-map}, replacing
each power of $p$ there by the
same power of this original $N$.
Thus it is an explicit nil-kernel
quotient, and is an isomorphism
if $\varphi$ was injective.

For a fixed prime $p$, retain
instead $n_i=p^{e_i}m_i$ and
$M=\prod_i m_i$. Over $R[1/M]$,
the global preimages above give
both
$$
x_i^{p^{e_i}}=
\frac{\varphi(A_iU_{i,p^{e_i}}+B_iE_{i,p^{e_i}})}{m_i},
\qquad
p^{e_i}x_i=\frac{\varphi(U_{i,1})}{m_i}.
$$
When $e_i=0$ these put $x_i$
itself in the image; one may
then choose exponent $1$.
For $e_i>0$ these are the
original two $p$-power conditions.
The kernel remains $I[1/M]$
and is nil. Thus the whole
localized map satisfies Lemma
\ref{editorial-source-lemma-p-ring-map}, with the
exact root exponents and all
prime-to-$p$ denominators retained.
\end{proof}


\subsection{Editorial treatment of Algebra lines 11082--11103}

\hypertarget{algebra-context-136}{}

\noindent\href{https://github.com/stacks/stacks-project/blob/a04446e57ec1fbc252a871afcec7752fb2807b14/algebra.tex\#L11082}{Original source, lines 11082--11103}. Complete original and reviewed TeX passages accompany this note. Every intervention is separately recorded; the following treatment is editorial.


\begin{lemma}
\label{editorial-source-lemma-flat-fibres-irreducible}
Let $R \to S$ be a ring map. Assume
\begin{enumerate}
\item[(a)] $\Spec(R)$ is irreducible,
\item[(b)] $R \to S$ is flat,
\item[(c)] $R \to S$ is of finite presentation,
\item[(d)] the fibre rings $S \otimes_R \kappa(\mathfrak p)$
have irreducible spectra for a dense collection of primes $\mathfrak p$ of $R$.
\end{enumerate}
Then $\Spec(S)$ is irreducible.
This is true more generally with (b) $+$ (c)
replaced by ``the map $\Spec(S) \to \Spec(R)$ is open''.
\end{lemma}

\begin{proof}
Write $\varphi:R\to S$ for the original
ring map. Put $X=\Spec(S)$, $Y=\Spec(R)$
and retain its spectrum map $f:X\to Y$. Proposition
\ref{editorial-source-proposition-fppf-open} makes $f$ open
under (b) and (c). We prove the stated
more general open-map assertion directly.
Let $D\subseteq Y$ be the given dense
set of points with irreducible fibres.
The irreducible space $Y$ is nonempty,
so $D$ is nonempty. Each fibre at a
point of $D$ is nonempty; consequently
$X$ is nonempty as well.

Let $U,V$ be nonempty open subsets of
$X$. The original images $f(U),f(V)$
are nonempty open subsets of the
irreducible $Y$, so their intersection
is nonempty and open. Density supplies
$y\in D\cap f(U)\cap f(V)$.
The subsets $U\cap f^{-1}(y)$ and
$V\cap f^{-1}(y)$ are nonempty open
subsets of the irreducible fibre.
Their intersection is nonempty,
so $U\cap V\neq\emptyset$.
This is the nonempty-space criterion
for irreducibility, proving the lemma
and supplying the nonemptiness step
in the cited Topology, Lemma
\ref{topology-lemma-irreducible-on-top}.

The same proof needs less than
openness: it is enough that the
image of every nonempty open
$U\subseteq X$ contain a nonempty
open subset of $Y$. Choose such
subsets inside $f(U)$ and $f(V)$
and perform the same intersection
argument with $D$. The whole proof
also works with the actual subspace
$f(X)\subseteq Y$ as codomain,
provided it is irreducible, the
specified interior condition is
relative to that subspace, and
the good fibres form a dense
subset of it. These are statements
about the original map with its
codomain restricted, not an
assumption that the image equals $Y$.

For flat ring maps there is a
different exact criterion without
finite presentation. Under (a),
let $\mathfrak p=\operatorname{Nil}(R)$
be the unique minimal prime of
the original $R$. Put
$C=S\otimes_R\kappa(\mathfrak p)$.
The original fibre comparison is
$$
C\simeq(R\setminus\mathfrak p)^{-1}(S/\mathfrak pS),
\qquad
s\otimes(\overline a/\overline b)
\longmapsto\varphi(a)s/\varphi(b),
\quad b\notin\mathfrak p.
$$
The inverse sends $\overline s/\varphi(b)$
to $s\otimes(1/\overline b)$;
quotient and localization balancing
prove well-definedness and both
composites. Primes of $C$ therefore
correspond exactly to original
primes of $S$ contracting to
$\mathfrak p$.

If $\mathfrak q$ is a minimal
prime of $S$, going down for
flat maps (Lemma \ref{editorial-source-lemma-flat-going-down}) lifts the inclusion
$\mathfrak p\subseteq\varphi^{-1}(\mathfrak q)$
to a prime contained in
$\mathfrak q$. Minimality forces
equality, so its contraction is
exactly $\mathfrak p$. Its fibre
prime is minimal, since a smaller
fibre prime would contract to
a smaller prime of $S$.
Conversely, if a minimal fibre
prime corresponds to $\mathfrak q$,
choose a minimal prime
$\mathfrak q_0\subseteq\mathfrak q$
of $S$. Its contraction is
$\mathfrak p$ by the previous
argument. Its fibre prime is
contained in the given minimal
one, so both are equal, and
$\mathfrak q_0=\mathfrak q$.
Thus these original fibre maps
give a bijection between the
minimal primes of $S$ and $C$.
If either ring is zero, both
sets are empty; if $S$ is nonzero,
it has a minimal prime and hence
$C$ is nonzero. Consequently,
under (a) and flatness alone,
$\Spec(S)$ is irreducible if and
only if this original generic
fibre has irreducible spectrum.

Dense irreducible fibres cannot
replace the generic-fibre condition
after dropping the openness
assumption. Retain the original
diagonal map
$$
\mathbf Z\longrightarrow\mathbf Q\times\mathbf Z,
\qquad n\longmapsto(n,n).
$$
It is flat: its underlying module
is the finite direct sum of the
flat localization $\mathbf Q$
and the free module $\mathbf Z$.
For every prime integer $\ell$,
its fibre is $\mathbf F_\ell$.
Indeed $\mathbf Q\otimes_{\mathbf Z}\mathbf F_\ell=0$,
since $1\otimes1=\ell^{-1}\otimes\ell=0$,
and the second factor gives
$\mathbf Z\otimes\mathbf F_\ell\simeq\mathbf F_\ell$.
These closed primes are dense:
every nonempty $D(n)$, $n\neq0$,
contains a prime dividing $|n|+1$
and hence not dividing $n$.
But the generic fibre is
$\mathbf Q\times\mathbf Q$,
through the multiplication map
$q\otimes(a,z)\mapsto(qa,qz)$
and its inverse
$(a,b)\mapsto a\otimes(1,0)+b\otimes(0,1)$.
The equality on the first factor
uses the original localization
identity $\mathbf Q\otimes_{\mathbf Z}\mathbf Q\simeq\mathbf Q$.
Both the total spectrum and its
generic fibre are reducible.
The nonempty open component
$\Spec(\mathbf Q)$ has image
$\{(0)\}\subseteq\Spec(\mathbf Z)$,
which has empty interior by the
same $D(n)$ argument. This
identifies exactly the failed
topological hypothesis.
\end{proof}


\subsection{Editorial treatment of Algebra lines 11105--11143}

\hypertarget{algebra-context-137}{}

\noindent\href{https://github.com/stacks/stacks-project/blob/a04446e57ec1fbc252a871afcec7752fb2807b14/algebra.tex\#L11105}{Original source, lines 11105--11143}. Complete original and reviewed TeX passages accompany this note. Every intervention is separately recorded; the following treatment is editorial.


\begin{lemma}
\label{editorial-source-lemma-separably-closed-irreducible}
Let $k$ be a separably closed field.
Let $R$, $S$ be $k$-algebras. If $R$, $S$ have a unique
minimal prime, so does $R \otimes_k S$.
\end{lemma}

\begin{proof}
Choose the original perfect closure
$\overline k/k$ of Definition
\ref{editorial-source-definition-perfection}. It is
algebraically closed in this situation.
Indeed take an algebraic closure
$\Omega/k$. Its maximal separable
subextension is $k$, because $k$
is separably closed. Fields, Lemma
\ref{fields-lemma-separable-first}
therefore makes $\Omega/k$ purely
inseparable. The algebraically
closed field $\Omega$ is perfect,
so uniqueness of perfect closure
gives the unique $k$-isomorphism
$\overline k\simeq\Omega$.
In characteristic zero both
extensions are already the identity.

For each original $k$-algebra $A$,
the map $A\to A\otimes_k\overline k$
is injective by extension of a
vector-space basis. In positive
characteristic the receiving algebra
is generated by $1\otimes\lambda$,
each with a $p$-power in the image
of $k$, and the corresponding
integer multiple is zero.
Lemma \ref{editorial-source-lemma-p-ring-map} gives
a homeomorphism on spectra; in
characteristic zero the map is
the identity under multiplication.
Apply this to $R$, $S$ and the
original $R\otimes_k S$.
The scalar comparison is the
isomorphism
$$
(R\otimes_k S)\otimes_k\overline k
\longrightarrow
(R\otimes_k\overline k)\otimes_{\overline k}
(S\otimes_k\overline k),
\qquad
(r\otimes s)\otimes c\longmapsto
(r\otimes c)\otimes(s\otimes1).
$$
Its inverse sends
$(r\otimes a)\otimes(s\otimes b)$
to $(r\otimes s)\otimes ab$.
The original scalar actions prove
balancing and both composites.
The two changed factors still
have a unique minimal prime.
It therefore suffices to prove
the assertion over the actual
algebraically closed field
$\overline k$, and then descend
through these homeomorphisms.

Work now over an algebraically
closed field $k$, keeping the
algebras then under consideration.
For arbitrary $R,S$, the tensor
quotient
$$
R\otimes_k S\longrightarrow R_{red}\otimes_k S_{red}
$$
is surjective with kernel
the sum of the two ideals
generated by the original factor
nilradicals. The inverse of its
induced quotient isomorphism
sends $[r]\otimes[s]$ to the
class of $r\otimes s$.
Each element of this kernel is
a finite sum of nilpotent pure
tensors; if their exponents are
$d_1,\ldots,d_m$, the full
multinomial expansion at exponent
$1+\sum_i(d_i-1)$ is zero.
Thus the kernel is locally
nilpotent and Lemma
\ref{editorial-source-lemma-surjective-locally-nilpotent-kernel}
identifies the spectra.
Since the original factors have
unique minimal primes, their
reductions are nonzero domains.
We may consequently prove the
remaining assertion for these
two domains, without losing the
nilpotent information in the
original quotient map.

For domains $R,S$ over this
algebraically closed $k$, their
tensor ring is nonzero, as is
seen by extending the vectors
$1\in R$ and $1\in S$ to bases.
It is reduced by Lemma
\ref{editorial-source-lemma-perfect-reduced}.
If it were not a domain, there
would be nonzero original tensors
$z,w$ with $zw=0$. Each uses
finitely many coefficients from
$R$ and $S$. Choose the subalgebras
$R_0,S_0$ generated by all those
coefficients. They are domains
of finite type over the
algebraically closed field $k$.
The map $R_0\otimes_k S_0\to R\otimes_k S$
is injective by applying the
two vector-space injections
successively. Hence the same
tensors $z,w$ are nonzero already
in $R_0\otimes_k S_0$, and
their product is zero there.

The original $R_0$-algebra
$R_0\otimes_k S_0$ is flat:
a $k$-basis of $S_0$ gives a
free $R_0$-module basis.
It is finitely presented since
$S_0$ has a finite presentation
over $k$ by the Hilbert basis
theorem, and the same generators
and relations give its scalar
extension over $R_0$.
For each maximal ideal
$\mathfrak m\subset R_0$, the
Nullstellensatz (Theorem \ref{editorial-source-theorem-nullstellensatz}) gives the
original scalar isomorphism
$k\to R_0/\mathfrak m$.
The fibre map
$$
(R_0\otimes_k S_0)\otimes_{R_0}R_0/\mathfrak m
\longrightarrow S_0,
\qquad(r\otimes s)\otimes\overline a
\longmapsto\overline r\,\overline a\,s
$$
uses that isomorphism; its
inverse is $s\mapsto(1\otimes s)\otimes1$.
These fibres are irreducible.
The maximal ideals are dense
because $R_0$ is a finite type
algebra over a field and is
Jacobson, by Lemma
\ref{algebra-lemma-finite-type-field-Jacobson}.
Lemma \ref{editorial-source-lemma-flat-fibres-irreducible}
therefore makes the tensor's
spectrum irreducible. It is
also reduced by perfectness,
so its nilradical is the zero
prime and it is a domain.
This contradicts the original
nonzero $z,w$ with product zero.
Thus the arbitrary domain tensor
is a domain, and the preceding
quotient and field-extension
maps give the result for the
original $R,S$ over the original
separably closed field.
\end{proof}


\subsection{Editorial treatment of Algebra lines 11145--11204}

\hypertarget{algebra-context-138}{}

\noindent\href{https://github.com/stacks/stacks-project/blob/a04446e57ec1fbc252a871afcec7752fb2807b14/algebra.tex\#L11145}{Original source, lines 11145--11204}. Complete original and reviewed TeX passages accompany this note. Every intervention is separately recorded; the following treatment is editorial.


\begin{lemma}
\label{editorial-source-lemma-geometrically-irreducible}
Let $k$ be a field.
Let $R$ be a $k$-algebra.
The following are equivalent
\begin{enumerate}
\item for every field extension $k'/k$ the
spectrum of $R \otimes_k k'$ is irreducible,
\item for every finite separable field extension $k'/k$ the
spectrum of $R \otimes_k k'$ is irreducible,
\item the spectrum of $R \otimes_k \overline{k}$ is irreducible
where $\overline{k}$ is the separable algebraic closure of $k$, and
\item the spectrum of $R \otimes_k \overline{k}$ is irreducible
where $\overline{k}$ is the algebraic closure of $k$.
\end{enumerate}
\end{lemma}

\begin{proof}
Condition (1) implies (2) directly.
Under (2), taking the identity field
extension shows that the original
$R$ has nonempty spectrum, so $R\neq0$.
Let $k_s$ be a separable algebraic
closure of $k$. Its finite
subextensions $E/k$ are separable,
are directed under compositum,
and have union $k_s$.
Every original finite list of
coefficients in $k_s$ belongs to
one such $E$, so
$R\otimes_k k_s=\bigcup_E R\otimes_k E$
through the actual injective
coefficient maps.

This tensor is nonzero by a
vector-space basis argument.
Suppose $\mathfrak q_1,\mathfrak q_2$
are minimal primes in it.
For each original finite $E$,
the map $R\otimes_k E\to R\otimes_k k_s$
is flat: its module comparison
is scalar extension by the field
$k_s/E$. Going down makes the
contractions of both primes
minimal. Indeed a smaller prime
under a contraction would lift
to a smaller prime under the
given minimal one. By (2), the
ring $R\otimes_k E$ has only
one minimal prime, so these
contractions agree. Each element
of $R\otimes_k k_s$ lies in some
one of the displayed finite
tensors; membership in either
$\mathfrak q_i$ is therefore
determined by equal contractions.
Thus $\mathfrak q_1=\mathfrak q_2$.
A nonzero ring has a minimal
prime, so this proves (3),
including the required existence.

For (3) $\Rightarrow$ (4), choose
an algebraic closure $\overline k$
containing the specified $k_s$.
The extension $\overline k/k_s$
is purely inseparable, or is the
identity in characteristic zero,
by the separable-first theorem.
The map
$R\otimes_k k_s\to(R\otimes_k k_s)\otimes_{k_s}\overline k$
is a homeomorphism on spectra
by Lemma \ref{editorial-source-lemma-p-ring-map}
in positive characteristic, and
is the identity in characteristic
zero. The latter tensor is
identified with $R\otimes_k\overline k$
by
$(r\otimes a)\otimes b\mapsto r\otimes ab$,
whose inverse is
$r\otimes b\mapsto(r\otimes1)\otimes b$.
This proves (4) for these original
field embeddings.

For (4) $\Rightarrow$ (1), keep
an arbitrary original extension
$L/k$. A common receiving field
must contain compatible copies
over $k$; construct it as follows.
The tensor $L\otimes_k\overline k$
is nonzero by extension of a
vector-space basis. Choose a
prime $\mathfrak a$ in it and put
$$
F=\operatorname{Frac}((L\otimes_k\overline k)/\mathfrak a).
$$
The maps $a\mapsto[a\otimes1]$
and $b\mapsto[1\otimes b]$ embed
$L$ and $\overline k$ into this
nonzero field and agree on the
original $k$ by tensor balancing.
Thus they provide the needed
commutative field diagram, not
merely unrelated abstract copies.
The scalar comparison
$$
(R\otimes_k\overline k)\otimes_{\overline k}F
\longrightarrow R\otimes_k F,
\qquad(r\otimes a)\otimes b\longmapsto r\otimes ab
$$
has inverse $r\otimes b\mapsto(r\otimes1)\otimes b$.
Its source has irreducible
spectrum by Lemma
\ref{editorial-source-lemma-separably-closed-irreducible},
applied over the algebraically
closed $\overline k$ to the
original $R\otimes_k\overline k$
and the field $F$.

The map $B=R\otimes_k L\to R\otimes_k F$
is scalar extension by $F/L$.
It is faithfully flat: for any
nonzero $B$-module $M$, the
comparison
$M\otimes_B(B\otimes_L F)\simeq M\otimes_L F$
sends $m\otimes(b\otimes a)$ to
$bm\otimes a$, with inverse
$m\otimes a\mapsto m\otimes(1\otimes a)$;
extension of a vector-space
basis preserves nonzeroness
and exact sequences. More
explicitly, over each prime
$\mathfrak b$ of $B$ the
nonzero ring $\kappa(\mathfrak b)\otimes_L F$
has a prime. Its inverse image
is a prime of $B\otimes_L F$
contracting to $\mathfrak b$.
Hence the spectrum map is
surjective. The continuous
surjective image of the
nonempty irreducible spectrum
just proved is irreducible,
giving (1) for the original
$R\otimes_k L$.

Two further forms follow with
the same maps. First, in (2)
it suffices to test finite
Galois extensions. Every finite
separable $E/k$ embeds in a
finite separable normal splitting
field $N/k$: take a primitive
element for $E/k$ and adjoin
all roots of its separable
minimal polynomial in an
algebraic closure. These are
finitely many algebraic roots,
all separable, so the resulting
extension is finite and
separable; each embedding
permutes the roots, so it is
normal. The field extension
$N/E$ makes
$R\otimes_k E\to R\otimes_k N$
faithfully flat by the same
module comparison above.
Irreducibility of the latter
spectrum therefore descends
to the former, proving (2).

Second, for every fixed field
extension $L/k$, the original
$R$ is geometrically irreducible
over $k$ if and only if
$R\otimes_k L$ is geometrically
irreducible over $L$.
The forward direction uses
$(R\otimes_k L)\otimes_L F\simeq R\otimes_k F$
for every field $F/L$.
For the reverse direction and
any original field $E/k$, use
the same prime-of-a-tensor
construction to place $L$ and
$E$ compatibly in a field $F$.
The hypothesis over $L$ makes
$R\otimes_k F$ irreducible,
and the faithfully flat map
$R\otimes_k E\to R\otimes_k F$
descends irreducibility as
above. Every scalar comparison
is the displayed tensor map
and its displayed inverse.
\end{proof}


\subsection{Editorial treatment of Algebra lines 11220--11232}

\hypertarget{algebra-context-139}{}

\noindent\href{https://github.com/stacks/stacks-project/blob/a04446e57ec1fbc252a871afcec7752fb2807b14/algebra.tex\#L11220}{Original source, lines 11220--11232}. Complete original and reviewed TeX passages accompany this note. Every intervention is separately recorded; the following treatment is editorial.


\begin{lemma}
\label{editorial-source-lemma-separably-closed-irreducible-implies-geometric}
Let $k$ be a field.
Let $R$ be a $k$-algebra.
If $k$ is separably closed then $R$ is
geometrically irreducible over $k$ if and only if the
spectrum of $R$ is irreducible.
\end{lemma}

\begin{proof}
If $R$ is geometrically irreducible,
take the identity field extension in
Definition \ref{algebra-definition-geometrically-irreducible}.
It gives irreducibility of $\Spec(R)$.
Conversely suppose $\Spec(R)$ is
irreducible and $k$ is separably closed.
For every finite separable extension
$E/k$, each original element of $E$
has a separable irreducible minimal
polynomial over $k$. Separably closedness
forces that polynomial to be linear,
so the given embedding $k\to E$ is
surjective and hence an isomorphism.
The map $R\otimes_k E\to R$ sends
$r\otimes a$ to $r\iota^{-1}(a)$,
where $\iota:k\to E$ is that given
embedding; its inverse is $r\mapsto r\otimes1$.
Thus every finite separable test has
irreducible spectrum. Lemma
\ref{editorial-source-lemma-geometrically-irreducible}
gives geometric irreducibility over
the original $k$. The proof imposes
neither perfectness nor algebraic
closedness on that separably closed
field. If $R=0$, both irreducibility
assertions are false, in accordance
with the nonemptiness convention.
\end{proof}


\subsection{Editorial treatment of Algebra lines 11234--11254}

\hypertarget{algebra-context-140}{}

\noindent\href{https://github.com/stacks/stacks-project/blob/a04446e57ec1fbc252a871afcec7752fb2807b14/algebra.tex\#L11234}{Original source, lines 11234--11254}. Complete original and reviewed TeX passages accompany this note. Every intervention is separately recorded; the following treatment is editorial.


\begin{lemma}
\label{editorial-source-lemma-subalgebra-geometrically-irreducible}
Let $k$ be a field. Let $S$ be a $k$-algebra.
\begin{enumerate}
\item If $S$ is geometrically irreducible over $k$ so is every
$k$-subalgebra.
\item If all finitely generated $k$-subalgebras of $S$ are
geometrically irreducible, then $S$ is geometrically irreducible.
\item A directed colimit of geometrically irreducible $k$-algebras
is geometrically irreducible.
\end{enumerate}
\end{lemma}

\begin{proof}
For a nonzero ring $A$, its spectrum
is irreducible exactly when its
nilradical is prime, by Lemma
\ref{editorial-source-lemma-irreducible}. Equivalently,
whenever $ab$ is nilpotent, at
least one of $a,b$ is nilpotent.
We keep nonzeroness in this
criterion; the zero ring has
empty spectrum and is not
irreducible.

For (1), retain the original
subalgebra $S'\subseteq S$ and
any original field extension
$L/k$. The map
$S'\otimes_k L\to S\otimes_k L$
is injective by extension of
a vector-space basis. For any
injective ring map $A\to B$,
the nilradical of $A$ is exactly
the inverse image of that of
$B$: a power becomes zero in
$B$ exactly when it was zero
in $A$. Here the receiving
nilradical is a proper prime,
so its inverse image is also
a proper prime. The source is
nonzero because its $1$ remains
the nonzero $1$ of the target.
The criterion proves that
$S'\otimes_k L$ is irreducible,
for every such $L$, proving (1).
This argument works for any
injective $k$-algebra map, with
its actual image subalgebra.

For (3), let $(S_i)$ be a
nonempty directed system of
geometrically irreducible
$k$-algebras with their original
unital transition maps; they
need not be injective. Put
$S=\mathop{\rm colim}_i S_i$.
For each original $L/k$ there
are inverse maps
$$
S\otimes_k L\longrightarrow
\mathop{\rm colim}_i(S_i\otimes_k L),
\qquad [i,s]\otimes a\longmapsto[i,s\otimes a],
$$
and
$$
[i,\sum_j s_j\otimes a_j]
\longmapsto\sum_j[i,s_j]\otimes a_j.
$$
All expressions use finitely
many elements, so a common
later index verifies balancing,
independence of representatives,
and both composites. Let $A_i$
denote the original $S_i\otimes_k L$.
Each is nonzero with prime
nilradical. Their unital colimit
$A$ is nonzero: an equality
$1=0$ there would hold at a
later stage, contradicting
$1\neq0$ in that $A_i$.

Suppose original elements
$a,b\in A$ satisfy $(ab)^n=0$.
Choose representatives at a
common stage $i$. The zero
equality holds at some later
stage $j$. Its nilradical is
prime, so the image of $a$
or of $b$ is nilpotent at
that stage. The corresponding
element in the original
colimit is then nilpotent.
This proves that the nilradical
of $A$ is prime, and proves
(3) after the explicit tensor
comparison. In particular no
injection assumption on the
transition maps is used.
If an empty colimit is allowed
in the category of $k$-algebras,
it is its initial object $k$,
which also satisfies the
assertion.

For (2), the original finitely
generated $k$-subalgebras of $S$
form a directed system under
inclusion: adjoining the union
of two finite lists gives a
common member. Every original
element of $S$ occurs in such
a subalgebra, so their colimit
is exactly $S$. Apply (3).
The hypothesis excludes $S=0$,
since its finitely generated
subalgebra generated by the
image of $k$ would itself be
zero and would not satisfy
the stated hypothesis.

The same nilradical criterion
also proves localization
stability with its necessary
nonemptiness qualification.
Assume $S$ is geometrically irreducible.
Let $W\subseteq S$ be a
multiplicative subset and
assume $W^{-1}S\neq0$.
Since $S$ is geometrically
irreducible, its nilradical
$N$ is prime; $W$ avoids $N$,
because a nilpotent becoming
invertible would make the
localization zero. For every
original $L/k$, the injection
$S\to S\otimes_k L$ shows
that no image of an element
of $W$ is nilpotent. The
prime nilradical of that
tensor therefore survives
localization and gives a
nonzero ring with one minimal
prime. Explicitly the nilradical
of a localization $B_W$ is
$W^{-1}\operatorname{Nil}(B)$:
if $(b/w)^n=0$, some $v\in W$
has $v b^n=0$, and then
$(vb)^n=v^{n-1}(v b^n)=0$,
while $b/w=(vb)/(vw)$.
The converse follows by taking
powers of a nilpotent numerator.
The original tensor-localization
maps are
$$
(W^{-1}S)\otimes_k L
\longrightarrow (W\otimes1)^{-1}(S\otimes_k L),
\qquad(s/w)\otimes a\longmapsto(s\otimes a)/(w\otimes1),
$$
and the inverse sends
$\sum_i(s_i\otimes a_i)/(w\otimes1)$
to $\sum_i(s_i/w)\otimes a_i$.
The two localization universal
properties prove these maps
well-defined and inverse.
Thus every nonzero localization
of the original $S$ remains
geometrically irreducible.
\end{proof}


\subsection{Editorial treatment of Algebra lines 11256--11290}

\hypertarget{algebra-context-141}{}

\noindent\href{https://github.com/stacks/stacks-project/blob/a04446e57ec1fbc252a871afcec7752fb2807b14/algebra.tex\#L11256}{Original source, lines 11256--11290}. Complete original and reviewed TeX passages accompany this note. Every intervention is separately recorded; the following treatment is editorial.


\begin{lemma}
\label{editorial-source-lemma-geometrically-irreducible-any-base-change}
Let $k$ be a field.
Let $S$ be a geometrically irreducible $k$-algebra.
Let $R$ be any $k$-algebra.
The map
$$
\Spec(R \otimes_k S) \longrightarrow \Spec(R)
$$
induces a bijection on irreducible components.
\end{lemma}

\begin{proof}
Write $A=R\otimes_k S$ and retain
$f:\Spec(A)\to\Spec(R)$. The
original $S$ is nonzero, since
geometric irreducibility includes
nonemptiness. If $R=0$, then
$A=0$ and both sets of components
are empty, so the assertion
holds. Otherwise, for every
$\mathfrak p\in\Spec(R)$ the
original quotient and fibre maps
give
$$
A/\mathfrak pA\simeq(R/\mathfrak p)\otimes_k S
\longrightarrow\kappa(\mathfrak p)\otimes_k S.
$$
The quotient isomorphism sends
$[r\otimes s]$ to $\overline r\otimes s$
and has inverse
$\overline r\otimes s\mapsto[r\otimes s]$.
The arrow is injective, since
$R/\mathfrak p\to\kappa(\mathfrak p)$
is an injection of $k$-vector
spaces and tensoring by $S$
preserves it. Its target has
prime nilradical by the given
geometric irreducibility of $S$.
The inverse image of this
nilradical is precisely the
nilradical of the source, by
injectivity, and is therefore
a proper prime. In particular
$A/\mathfrak pA$ has irreducible
nonempty spectrum.

Consequently the original ideal
$$
\sigma(\mathfrak p)=\sqrt{\mathfrak pA}
$$
is prime, and $f^{-1}(V(\mathfrak p))$
is the irreducible closed subset
$V(\sigma(\mathfrak p))$.
Its contraction is exactly
$\mathfrak p$. To verify this,
the injection
$R/\mathfrak p\to(R/\mathfrak p)\otimes_k S$
can be split as a linear map
by choosing a $k$-linear
functional $S\to k$ taking
$1$ to $1$. Hence an element
of $R/\mathfrak p$ becomes
nilpotent in that tensor
exactly when it was nilpotent
in the original domain, that
is, exactly when it is zero.
This proves the contraction
claim with the given scalar
map.

The map $R\to A$ is flat: a
$k$-basis of $S$ gives a free
$R$-module basis of $A$.
Going down therefore sends
every minimal prime of $A$
to a minimal prime of $R$.
If $\mathfrak p$ is minimal
in $R$, choose a minimal
$\mathfrak q_0\subseteq\sigma(\mathfrak p)$
of $A$. Its contraction is
a minimal prime contained
in $\mathfrak p$, hence is
exactly $\mathfrak p$.
It contains $\mathfrak pA$,
so it contains
$\sqrt{\mathfrak pA}=\sigma(\mathfrak p)$.
Thus equality holds, proving
that $\sigma(\mathfrak p)$
is minimal. Conversely, for
any minimal $\mathfrak q$
of $A$, put
$\mathfrak p=\mathfrak q\cap R$.
Then $\sigma(\mathfrak p)\subseteq\mathfrak q$,
so minimality makes them equal.
These actual inverse maps on
minimal primes prove the
bijection on components via
Lemma \ref{editorial-source-lemma-irreducible}.

The construction gives more:
$\sigma$ is a continuous section
of $f$ on the entire spectra,
selecting the generic point
of each original fibre.
For every $\mathfrak p$, its
prime is contained in every
prime of $A$ over $\mathfrak p$,
and its contraction equals
$\mathfrak p$, so it is the
unique generic fibre prime.
For an original $a\in A$ one
has the exact equality
$$
\sigma^{-1}(D(a))=f(D(a)).
$$
If a prime above $\mathfrak p$
does not contain $a$, the
smaller $\sigma(\mathfrak p)$
does not contain it. Conversely
if $a\notin\sigma(\mathfrak p)$,
this prime itself is a point
of $D(a)$ above $\mathfrak p$.
Lemma \ref{editorial-source-lemma-map-into-tensor-algebra-open}
makes the right side open,
proving continuity. Since
$f\sigma$ is the identity,
the section is a topological
embedding, with inverse the
restriction of $f$ to its
image. That image is dense:
every nonempty basic open
$D(a)$ has nonempty image,
and the displayed equality
puts a section point in it.
Also
$$
\overline{\{\sigma(\mathfrak p)\}}
=V(\sqrt{\mathfrak pA})
=f^{-1}(V(\mathfrak p)).
$$
Restricting the two continuous
inverse maps to the minimal
primes proves a homeomorphism
of their subspaces in the
respective spectra, strengthening
the stated component bijection.

Finally this section is
compatible with arbitrary
original $k$-algebra base change
$\theta:R\to R'$. If $\mathfrak p'$
contracts to $\mathfrak p$,
the induced map of fields
$\kappa(\mathfrak p)\to\kappa(\mathfrak p')$
gives an injection
$$
\kappa(\mathfrak p)\otimes_k S
\longrightarrow\kappa(\mathfrak p')\otimes_k S.
$$
The inverse image of the
nilradical on the right is
the nilradical on the left,
since injective ring maps
reflect nilpotence. The two
quotient-fibre descriptions
above therefore show that
$\sigma_{R'}(\mathfrak p')$
contracts to
$\sigma_R(\mathfrak p)$ in
$R\otimes_k S$. This proves
the stated compatibility
through the original tensor
map $r\otimes s\mapsto\theta(r)\otimes s$.
The section is a continuous
map of these spectra; no
extra residue-field
isomorphisms are asserted.
\end{proof}


\subsection{Editorial treatment of Algebra lines 11296--11320}

\hypertarget{algebra-context-142}{}

\noindent\href{https://github.com/stacks/stacks-project/blob/a04446e57ec1fbc252a871afcec7752fb2807b14/algebra.tex\#L11296}{Original source, lines 11296--11320}. Complete original and reviewed TeX passages accompany this note. Every intervention is separately recorded; the following treatment is editorial.


\begin{lemma}
\label{editorial-source-lemma-field-extension-geometrically-irreducible}
Let $K/k$ be a field extension. If $k$ is algebraically closed in $K$, then
$K$ is geometrically irreducible over $k$.
\end{lemma}

\begin{proof}
We prove the assertion under the weaker
hypothesis that every element of the
original $K$ separably algebraic over
$k$ belongs to $k$. The stated relative
algebraic closedness implies this.
By Lemma \ref{editorial-source-lemma-geometrically-irreducible}
it suffices to treat each original
finite separable extension $k'/k$.
Choose its primitive element $\alpha$
using Fields, Lemma
\ref{fields-lemma-primitive-element},
and retain its full minimal polynomial
$$
P=T^d+a_1T^{d-1}+\cdots+a_d\in k[T].
$$
The map $K[T]/(P)\to K\otimes_k k'$
sends $\sum_{j=0}^{d-1}b_jT^j$ to
$\sum_{j=0}^{d-1}b_j\otimes\alpha^j$.
Its inverse sends $b\otimes Q(\alpha)$
to the class of $bQ(T)$.
The original basis $1,\alpha,\ldots,\alpha^{d-1}$
and monic division by $P$ show that
these maps are well-defined inverses.

Suppose the original polynomial has
a nontrivial factorization $P=P_1P_2$
in $K[T]$. Set $r=\deg(P_1)$ and
$c=\operatorname{lc}(P_1)\in K^\times$.
Then $1\leq r\leq d-1$ and
$\operatorname{lc}(P_2)=c^{-1}$.
In an algebraic closure $\Omega_K$ of
$K$, the original separable $P$ has
distinct roots $\beta_1,\ldots,\beta_d$.
There is a subset $J$ of cardinality
$r$ such that the complete original
factors are
$$
P_1=c\prod_{i\in J}(T-\beta_i),\qquad
P_2=c^{-1}\prod_{i\notin J}(T-\beta_i).
$$
Define the auxiliary polynomials
$Q_1=c^{-1}P_1$ and $Q_2=cP_2$,
retaining the exact identity
$P=(cQ_1)(c^{-1}Q_2)=Q_1Q_2$.
For $0\leq j\leq r$, the coefficient
of $T^{r-j}$ in the original $P_1$ is
$$
c(-1)^j\sum_{\substack{H\subseteq J\\|H|=j}}
\prod_{i\in H}\beta_i.
$$
For $P_2$ the corresponding coefficient
has factor $c^{-1}$ and the same signed
sum over subsets of the complementary
root set. Empty products give the
original leading coefficients $c,c^{-1}$.

Each $\beta_i$ is separably algebraic
over $k$ because $P$ is separable.
The finite compositum $k(\beta_1,\ldots,\beta_d)$
is separable over $k$, by Fields,
Lemma \ref{fields-lemma-separable-permanence}
and permanence under enlarging the
base field. Consequently every displayed
coefficient ratio, namely each coefficient
of $Q_1$ and $Q_2$, is separably algebraic
over $k$. These ratios also lie in
the original $K$ by their definitions,
so the weakened hypothesis puts them
in $k$. Both $Q_1,Q_2$ thus belong to
$k[T]$ and have positive degree.
Their product is the original $P$,
contradicting its irreducibility over $k$.
It follows that $K[T]/(P)$ is a field;
in particular its spectrum is irreducible.
All finite separable tests now give
geometric irreducibility as asserted.

The leading factors cannot be discarded
from the coefficient argument in Lemma
\ref{editorial-source-lemma-polynomials-divide}, whose
statement concerns monic polynomials.
For example, over the original
$\mathbf Q(u)$ one has
$$
T^2-1=\bigl(u(T-1)\bigr)\bigl(u^{-1}(T+1)\bigr).
$$
The coefficients $u$ and $u^{-1}$
are not algebraic over $\mathbf Q$;
a nonzero polynomial equation for
$u^{-1}$ would, after multiplication
by its full power of $u$, give a
nonzero polynomial equation for $u$.
This example concerns the coefficient
inference for an arbitrary monic
product, not the irreducible-polynomial
hypothesis of the lemma. The proof
above keeps $c,c^{-1}$ and proves
exactly the required claim about
their coefficient ratios. It also
proves the stronger conclusion that
$K\otimes_k k'$ is a field for every
finite separable original $k'/k$.
\end{proof}


\subsection{Editorial treatment of Algebra lines 11322--11338}

\hypertarget{algebra-context-143}{}

\noindent\href{https://github.com/stacks/stacks-project/blob/a04446e57ec1fbc252a871afcec7752fb2807b14/algebra.tex\#L11322}{Original source, lines 11322--11338}. Complete original and reviewed TeX passages accompany this note. Every intervention is separately recorded; the following treatment is editorial.


\begin{lemma}
\label{editorial-source-lemma-geometrically-irreducible-transitive}
Let $K/k$ be a geometrically irreducible field extension.
Let $S$ be a geometrically irreducible $K$-algebra.
Then $S$ is geometrically irreducible over $k$.
\end{lemma}

\begin{proof}
Let $k'/k$ be any original finite
separable extension. By Fields,
Lemma \ref{fields-lemma-primitive-element},
write $k'=k(\alpha)$ with monic
separable minimal polynomial $P$.
The inverse tensor-polynomial maps
from Lemma
\ref{editorial-source-lemma-field-extension-geometrically-irreducible}
give $K'=K\otimes_k k'\simeq K[T]/(P)$.
Since $P$ is separable, there are
$U,V\in k[T]$ with $UP+VP'=1$.
The same full identity holds in
$K[T]$. Thus $P$ has no repeated
irreducible factor there: a repeated
factor would divide $P$ and $P'$
and hence divide $1$.

Factor $P=\prod_{i=1}^m Q_i$ into
its distinct monic irreducible
factors in $K[T]$. The leading
coefficient of $P$ is $1$, so this
product has exactly the original
leading coefficient. The Chinese
remainder map
$$
K[T]/(P)\longrightarrow\prod_{i=1}^m K[T]/(Q_i)
$$
sends the class of $F$ to its residue
classes. For each $i$, Bezout gives
$A_i\prod_{j\ne i}Q_j+B_iQ_i=1$;
the inverse sends $([F_i])_i$ to
$[\sum_i F_i A_i\prod_{j\ne i}Q_j]$.
These formulas give both composites
by reduction modulo every $Q_i$.
Each factor is a nonzero field,
so the spectrum has exactly $m$
points and is irreducible exactly
when $m=1$. The hypothesis on
$K/k$ therefore makes $K'$ a field.
It is finite over $K$ with basis
$1,T,\ldots,T^{\deg(P)-1}$, and
separable since its irreducible
polynomial $P$ has the displayed
Bezout identity with its derivative.

There are inverse original maps
$$
S\otimes_K(K\otimes_k k')\longrightarrow S\otimes_k k',
\qquad s\otimes(a\otimes b)\longmapsto as\otimes b,
$$
and $s\otimes b\mapsto s\otimes(1\otimes b)$.
The $K$-balancing relation proves
both composites are the identity.
Geometric irreducibility of $S/K$
applies to the actual field $K'/K$,
so this tensor has irreducible
spectrum. Every finite separable
$k'/k$ passes the test, and Lemma
\ref{editorial-source-lemma-geometrically-irreducible}
proves the result over the original $k$.
\end{proof}


\subsection{Editorial treatment of Algebra lines 11340--11375}

\hypertarget{algebra-context-144}{}

\noindent\href{https://github.com/stacks/stacks-project/blob/a04446e57ec1fbc252a871afcec7752fb2807b14/algebra.tex\#L11340}{Original source, lines 11340--11375}. Complete original and reviewed TeX passages accompany this note. Every intervention is separately recorded; the following treatment is editorial.


\begin{lemma}
\label{editorial-source-lemma-geometrically-irreducible-base-change-transcendental}
Let $K/k$ be a field extension. The following are equivalent
\begin{enumerate}
\item $K$ is geometrically irreducible over $k$, and
\item the induced extension $K(t)/k(t)$ of purely transcendental extensions
is geometrically irreducible.
\end{enumerate}
\end{lemma}

\begin{proof}
Assume (1). Let $\Omega$ be an
algebraic closure of the original
$k(t)$. The algebra $K\otimes_k\Omega$
has irreducible spectrum by (1).
Put $U=k[t]\setminus\{0\}$ and
$V=K[t]\setminus\{0\}$, keeping
the same indeterminate $t$.
There is an isomorphism
$$
K\otimes_k k(t)\longrightarrow U^{-1}K[t],
\qquad a\otimes(P(t)/Q(t))\longmapsto aP(t)/Q(t).
$$
The inverse sends
$(\sum_i a_it^i)/Q(t)$ to
$\sum_i a_i\otimes(t^i/Q(t))$.
The quotient relations for the
original denominators verify
independence of representatives,
and the formulas give both composites.
Localizing further at the images
of all of $V$ gives precisely the
original fraction field $K(t)$.

The two original scalar maps
$$
(K\otimes_k k(t))\otimes_{k(t)}\Omega
\longrightarrow K\otimes_k\Omega,
\qquad(a\otimes b)\otimes c\longmapsto a\otimes bc
$$
and $a\otimes c\mapsto(a\otimes1)\otimes c$
are inverse. Localization commutes
with these tensors: the fraction
$b/v\otimes c$ goes to
$(b\otimes c)/(v\otimes1)$, with
inverse determined by inverting
the same original $v$. Hence
$K(t)\otimes_{k(t)}\Omega$ is a
localization of $K\otimes_k\Omega$
at exactly those images of $V$.
It is nonzero because it is a
tensor product of two nonzero
fields over $k(t)$; extending
$1$ to vector-space bases proves
this directly. A nonzero
localization of a ring with
prime nilradical again has prime
nilradical, as proved in Lemma
\ref{editorial-source-lemma-subalgebra-geometrically-irreducible}.
Thus its spectrum is irreducible.
The algebraic-closure test of Lemma
\ref{editorial-source-lemma-geometrically-irreducible}
proves (2).

Conversely, assume (2) and retain
any original extension $k'/k$.
Put $A=K\otimes_k k'$ and let $W$
be the multiplicative subset of
$A[t]$ generated by the images of
all nonzero polynomials of $K[t]$
and of $k'[t]$. Both coefficient
maps into $A$ are injective:
extension of a vector-space basis
over $k$ proves this. More strongly,
the leading coefficient of each
such polynomial is a unit of $A$,
with inverse given by the inverse
of the original nonzero field
coefficient in the same tensor
factor. Multiplication by a
polynomial with unit leading
coefficient is injective on $A[t]$:
if $H$ is nonzero of degree $n$
and leading coefficient $h_n$,
the leading coefficient of its
product with a degree-$d$
polynomial of leading coefficient
$c$ is exactly $c h_n\neq0$.
Products of these injective
multiplication maps remain
injective. Thus every element
of $W$ is a nonzerodivisor and
$A\to A[t]\to W^{-1}A[t]$ is
injective, without assuming $A$
reduced or a domain.

The full receiving algebra is
exactly
$$
K(t)\otimes_{k(t)}k'(t)\simeq W^{-1}A[t].
$$
To construct the map from left
to right, use the original
coefficient embeddings and send
$$
\frac{F(t)}{f(t)}\otimes\frac{G(t)}{g(t)}
\longmapsto\frac{F(t)G(t)}{f(t)g(t)},
\qquad
f\in K[t]\setminus\{0\},\quad
g\in k'[t]\setminus\{0\}.
$$
The field-fraction relations hold
after precisely these denominators
are inverted. The two images of
every rational function in $k(t)$
agree, so the map is balanced over
$k(t)$. For the inverse, send
$a\otimes b\in A$ to $a\otimes b$
in the left-hand tensor and send
$t$ to $t\otimes1=1\otimes t$.
The image of each original
$f\in K[t]\setminus\{0\}$ is
the unit $f(t)\otimes1$, and
each original nonzero $g\in k'[t]$
gives the unit $1\otimes g(t)$.
The universal property therefore
extends this map to $W^{-1}A[t]$.
The composites fix $A$, $t$ and
each specified inverse denominator,
so they are identities on their
generating sets. In particular,
the injection just proved is
the original tensor-constant map
used in the source argument.

By (2), this receiving tensor
has nonempty irreducible spectrum,
so its nilradical is a proper
prime. Its inverse image under
the injection from $A$ is exactly
$\operatorname{Nil}(A)$, because
an injection reflects the vanishing
of every power. Hence $A$ is
nonzero with prime nilradical,
and its spectrum is irreducible.
This works for every original
$k'/k$, proving (1).

Iteration gives the same equivalence
after adjoining any specified finite
list of common independent
indeterminates. At each step the
proved one-variable equivalence
applies to the actual preceding
field extension and retains all
its polynomial denominators.
The empty list is the identity
extension, so it also satisfies
the assertion.
\end{proof}


\subsection{Editorial treatment of Algebra lines 11377--11389}

\hypertarget{algebra-context-145}{}

\noindent\href{https://github.com/stacks/stacks-project/blob/a04446e57ec1fbc252a871afcec7752fb2807b14/algebra.tex\#L11377}{Original source, lines 11377--11389}. Complete original and reviewed TeX passages accompany this note. Every intervention is separately recorded; the following treatment is editorial.


\begin{lemma}
\label{editorial-source-lemma-geometrically-irreducible-add-transcendental}
Let $K/L/M$ be a tower of fields with $L/M$ geometrically irreducible.
Let $x \in K$ be transcendental over $L$. Then $L(x)/M(x)$ is geometrically
irreducible.
\end{lemma}

\begin{proof}
Keep the original inclusions
$M\subseteq L\subseteq K$ and
the original $x\in K$.
Transcendence over $L$ makes
evaluation at $x$ injective on
$L[t]$, and hence on its subring
$M[t]$. In particular no original
nonzero denominator vanishes.
The maps
$$
\lambda:L(t)\longrightarrow L(x),\quad
\frac{P(t)}{Q(t)}\longmapsto\frac{P(x)}{Q(x)},
\qquad
\mu:M(t)\longrightarrow M(x)
$$
with the same evaluation formula
are injective field maps. They
are surjective because their
images are fields containing the
original coefficient fields and
$x$, respectively. Their inverses
replace $x$ by $t$ in a rational
expression; injectivity proves
independence of that expression.
The square with the original
inclusions commutes because the
formulas agree on every $M$
coefficient, on $t$, and on
every inverse denominator.

For any original field $E/M(x)$,
use its given embedding composed
with $\mu$ to view it over $M(t)$.
The map
$$
L(t)\otimes_{M(t)}E\longrightarrow L(x)\otimes_{M(x)}E,
\qquad a\otimes b\longmapsto\lambda(a)\otimes b
$$
is well-defined by that commuting
square, with inverse
$c\otimes b\mapsto\lambda^{-1}(c)\otimes b$.
Lemma
\ref{editorial-source-lemma-geometrically-irreducible-base-change-transcendental}
makes the left-hand spectrum
irreducible from the hypothesis
on $L/M$. This proves the assertion
for every given $E/M(x)$.
Conversely the same tensor maps,
using $\mu^{-1}$ to transport a
field over $M(t)$, show that
geometric irreducibility of
$L(x)/M(x)$ implies that of
$L(t)/M(t)$. The converse in the
preceding lemma then gives that
of $L/M$. Thus, under the stated
transcendence hypothesis, the
original conclusion is in fact
equivalent to geometric
irreducibility of $L/M$.
\end{proof}


\subsection{Editorial treatment of Algebra lines 11391--11425}

\hypertarget{algebra-context-146}{}

\noindent\href{https://github.com/stacks/stacks-project/blob/a04446e57ec1fbc252a871afcec7752fb2807b14/algebra.tex\#L11391}{Original source, lines 11391--11425}. Complete original and reviewed TeX passages accompany this note. Every intervention is separately recorded; the following treatment is editorial.


\begin{lemma}
\label{editorial-source-lemma-geometrically-irreducible-separable-elements}
Let $K/k$ be a field extension. The following are equivalent
\begin{enumerate}
\item $K/k$ is geometrically irreducible, and
\item every element $\alpha \in K$ separably algebraic over $k$ is
in $k$.
\end{enumerate}
\end{lemma}

\begin{proof}
Assume (1) and let the original
$\alpha\in K$ be separably algebraic
over $k$. Put $k'=k(\alpha)$ and
let $P\in k[T]$ be its monic
minimal polynomial of degree $d$.
Choose an algebraic closure
$\overline k$ of $k$ and retain
the distinct roots
$\alpha_1,\ldots,\alpha_d$ of $P$
in it. The original inclusion
$k'\subseteq K$ and Lemma
\ref{editorial-source-lemma-subalgebra-geometrically-irreducible}
make $k'/k$ geometrically irreducible.
The tensor-polynomial comparison
and the following explicit Chinese
remainder maps give
$$
k'\otimes_k\overline k\simeq
\overline k[T]/(P)\simeq\prod_{i=1}^d\overline k.
$$
The first sends $Q(\alpha)\otimes b$
to $[bQ(T)]$; its inverse sends
$[\sum_j b_jT^j]$ to
$\sum_j\alpha^j\otimes b_j$.
The second sends $[F]$ to
$(F(\alpha_i))_i$; its inverse is
$$
(b_1,\ldots,b_d)\longmapsto
\left[\sum_{i=1}^d b_i
\frac{\prod_{j\ne i}(T-\alpha_j)}
{\prod_{j\ne i}(\alpha_i-\alpha_j)}\right].
$$
All displayed denominators are
nonzero because the roots are
distinct. The interpolation
polynomials evaluate to the
coordinate vectors. Conversely,
a polynomial vanishing at all
these roots is divisible by
their full product $P$, proving
the other composite. These are
the embedding maps of Fields,
Lemma
\ref{fields-lemma-finite-separable-tensor-alg-closed},
with their inverse now explicit.
The product has $d$ distinct
maximal and minimal primes,
the kernels of its projections.
It has irreducible spectrum
only if $d=1$. Therefore the
original $\alpha$ belongs to $k$,
proving (2).

Assume (2) and retain the original
subfield $k'\subseteq K$ consisting
of all elements algebraic over $k$.
It is a field: a finite list of
such elements generates a finite
field extension of $k$, containing
their sums, products and inverses
of nonzero elements. Furthermore
$k'$ is algebraically closed in
$K$. Indeed if $\beta\in K$ is
algebraic over $k'$, the finitely
many coefficients of a polynomial
equation for $\beta$ belong to a
finite extension $E/k$ in $k'$.
Then $E(\beta)/E$ is finite, so
$E(\beta)/k$ is finite and
$\beta\in k'$ by the definition.
Lemma
\ref{editorial-source-lemma-field-extension-geometrically-irreducible}
therefore makes $K/k'$ geometrically
irreducible, using the full relative
algebraic closedness just proved.

By Fields, Lemma
\ref{fields-lemma-separable-first},
the algebraic extension $k'/k$
is purely inseparable over its
maximal separable subextension.
Condition (2) says precisely
that this subextension is $k$.
In characteristic zero this
forces $k'=k$. In characteristic
$p>0$, for every original $L/k$
the map $L\to L\otimes_k k'$
is injective by extension of
a vector-space basis. Its
receiving algebra is generated
over $L$ by $1\otimes a$ with
$a\in k'$, and each such
generator has some $p$-power
in the image of $L$, while
its corresponding integer
$p$-power multiple is zero.
Lemma \ref{editorial-source-lemma-p-ring-map}
thus identifies its spectrum
with the one-point $\Spec(L)$.
This proves geometric
irreducibility of $k'/k$,
including infinite extensions
with no common exponent assumed.
Lemma
\ref{editorial-source-lemma-geometrically-irreducible-transitive}
now proves (1) for the original
tower $K/k'/k$.

The criterion has an exact
tensor-field form: (1) and (2)
are also equivalent to
$K\otimes_k E$ being a field
for every separable algebraic
extension $E/k$. To prove it
from (2), each finite separable
subextension $E_i/k$ gives a
field $K\otimes_k E_i$ by the
coefficient proof of Lemma
\ref{editorial-source-lemma-field-extension-geometrically-irreducible}.
These fields form a directed
system of injective original
maps whose union is
$K\otimes_k E$: every tensor
has finitely many coefficients.
A nonzero element belongs to
one of those fields and has
its inverse there. Their union
is therefore a field. The
converse follows at once by
applying the finite separable
tests of Lemma
\ref{editorial-source-lemma-geometrically-irreducible}.

For any compatible embeddings
of the original $K$ and $E$
into a common field $F$, the
multiplication map
$$
K\otimes_k E\longrightarrow KE\subseteq F,
\qquad a\otimes b\longmapsto ab
$$
is consequently an isomorphism:
it is a unital map from a field,
so is injective; its image is
a field containing both original
embedded fields and consists of
their finite tensor sums, hence
equals their compositum.
For finite $E/k$, extension of
its original $k$-basis gives
the exact degree equality
$[KE:K]=[E:k]$. More generally,
every original finite linearly
independent family over $k$
in $E$ remains independent
over $K$: its tensor extension
is independent, and the proved
multiplication map is injective.
This identifies the precise
linear-disjointness consequence
without replacing the original
field embeddings.
\end{proof}


\subsection{Editorial treatment of Algebra lines 11427--11443}

\hypertarget{algebra-context-147}{}

\noindent\href{https://github.com/stacks/stacks-project/blob/a04446e57ec1fbc252a871afcec7752fb2807b14/algebra.tex\#L11427}{Original source, lines 11427--11443}. Complete original and reviewed TeX passages accompany this note. Every intervention is separately recorded; the following treatment is editorial.


\begin{lemma}
\label{editorial-source-lemma-make-geometrically-irreducible}
Let $K/k$ be a field extension. Consider the subextension $K/k'/k$ consisting
of elements separably algebraic over $k$. Then $K$ is geometrically irreducible
over $k'$. If $K/k$ is a finitely generated field extension, then
$[k' : k] < \infty$.
\end{lemma}

\begin{proof}
The original set $k'$ is a subfield
by Fields, Lemma \ref{fields-lemma-separable-elements}.
If $a\in K$ is separably algebraic
over $k'$, transitivity of separable
algebraic extensions (Fields, Lemma
\ref{fields-lemma-separable-permanence})
makes $a$ separably algebraic over
$k$. Hence $a\in k'$ by its
definition. Lemma
\ref{editorial-source-lemma-geometrically-irreducible-separable-elements}
now proves geometric irreducibility
of the original $K/k'$.

For completeness, retain the full
degree argument of Fields, Lemma
\ref{fields-lemma-algebraic-closure-in-finitely-generated}.
Choose a finite transcendence basis
$x_1,\ldots,x_r$ of the finitely
generated extension $K/k$ and put
$$
F=k(x_1,\ldots,x_r),\qquad n=[K:F]<\infty.
$$
Here $r=0$ is allowed, with $F=k$.
The remaining members of a finite
generating set for $K/k$ are
algebraic over $F$, so adjoining
them successively proves the
displayed finiteness. Let $E/k$
be any original finite algebraic
subextension of $K/k$, and choose
a basis $e_1=1,e_2,\ldots,e_d$ of
$E/k$, where $d=[E:k]$.

The original $x_i$ remain
algebraically independent over $E$.
Here is a direct proof retaining
every coefficient. For nonzero
$P\in E[T_1,\ldots,T_r]$, let $M_P$
be the matrix of multiplication
by $P$ on the free $k[T_1,\ldots,T_r]$-module
$E[T_1,\ldots,T_r]$ in the basis
$e_1,\ldots,e_d$. This multiplication
is injective because the latter
polynomial ring is a domain.
After localization to $k(T_1,\ldots,T_r)$,
the matrix is an injective square
linear map on a finite-dimensional
vector space, hence
$N_P=\det(M_P)\neq0$ in $k[T_1,\ldots,T_r]$.
Write $H_P\in E[T_1,\ldots,T_r]$
for the element whose coordinate
column is $\operatorname{adj}(M_P)(1,0,\ldots,0)^t$.
The full adjugate identity gives
$$
P H_P=N_P.
$$
If $P(x_1,\ldots,x_r)=0$, this
identity would imply
$N_P(x_1,\ldots,x_r)=0$, contradicting
the original independence over $k$.
This proves the asserted independence
over $E$, also for $r=0$.

Consequently the original evaluation
map induces inverse isomorphisms
$$
E\otimes_k F\longrightarrow E(x_1,\ldots,x_r),
\qquad a\otimes h(x)\longmapsto a h(x).
$$
To specify the inverse, write a
numerator in the basis $e_i$ and
retain each original denominator
$Q(x)$ with $Q\in E[T_1,\ldots,T_r]\setminus\{0\}$.
Use exactly the identity
$Q^{-1}=H_Q/N_Q$ just proved,
then expand the numerator $P H_Q$
in the same basis. This gives
its unique expression
$\sum_i e_i f_i(x)$ with $f_i\in F$,
and the inverse sends it to
$\sum_i e_i\otimes f_i(x)$.
Uniqueness follows by multiplying
all denominators in $k[T_1,\ldots,T_r]$
and comparing polynomial coefficients
in the basis $e_i$; no denominator
vanishes at the original $x$.
Thus these are well-defined inverse
maps and
$$
[E:k]=[E(x_1,\ldots,x_r):F]\leq[K:F]=n.
$$

In fact this bounds the full
algebraic constant field $k_a$
inside $K$, not only $k'$.
The degrees of its finite
subextensions lie among the
integers $1,\ldots,n$, so choose
one, $E_0$, of maximal degree.
For any $a\in k_a$, the original
$E_0(a)/k$ is finite and has
degree at least $[E_0:k]$.
Maximality and degree multiplicativity
give $[E_0(a):E_0]=1$. Hence every
$a$ lies in $E_0$, proving
$k_a=E_0$ and the precise bound
$$
[k':k]\leq[k_a:k]\leq[K:k(x_1,\ldots,x_r)].
$$

There is also an exact description
of the algebraic intermediate
fields $k\subseteq L\subseteq K$
over which $K$ is geometrically
irreducible: they are exactly
those containing $k'$.
For necessity, every $a\in k'$
is separably algebraic over $L$,
because its minimal polynomial
over $L$ divides its original
separable polynomial over $k$.
Geometric irreducibility of $K/L$
and the preceding criterion force
$a\in L$. For sufficiency, suppose
$k'\subseteq L$ and $L/k$ algebraic.
The full algebraic constant field
$k_a/k'$ is purely inseparable
by Fields, Lemma
\ref{fields-lemma-separable-first}.
If $a\in K$ is separably algebraic
over $L$, it is algebraic over
$k$, hence belongs to $k_a$.
Thus it is also purely inseparable
over $L$: some original $p$-power
lies in $k'\subseteq L$ in positive
characteristic, and in characteristic
zero $k_a=k'$. An element both
separable and purely inseparable
has degree one, since its minimal
polynomial is separable and divides
a polynomial with only one distinct
root. Therefore $a\in L$, and the
criterion proves sufficiency.
The restriction $L/k$ algebraic
is part of this characterization;
it is not removed.
\end{proof}


\subsection{Editorial treatment of Algebra lines 11445--11479}

\hypertarget{algebra-context-148}{}

\noindent\href{https://github.com/stacks/stacks-project/blob/a04446e57ec1fbc252a871afcec7752fb2807b14/algebra.tex\#L11445}{Original source, lines 11445--11479}. Complete original and reviewed TeX passages accompany this note. Every intervention is separately recorded; the following treatment is editorial.


\begin{lemma}
\label{editorial-source-lemma-Galois-orbit}
Let $K/k$ be an extension of fields.
Let $\overline{k}/k$ be a separable algebraic closure.
Then $\text{Gal}(\overline{k}/k)$ acts transitively on the
primes of $\overline{k} \otimes_k K$.
\end{lemma}

\begin{proof}
Retain the original separable closure
$\overline k$, put $G=\operatorname{Gal}(\overline k/k)$,
and let $K/k'/k$ be the original
subextension of Lemma
\ref{editorial-source-lemma-make-geometrically-irreducible}.
Write
$$
B=\overline k\otimes_k k',\qquad
A=\overline k\otimes_k K.
$$
Both tensors are nonzero by extension
of vector-space bases over $k$.
As $\overline k/k$ is algebraic,
these are integral algebras over
the respective fields $k'$ and
$K$: each generator $a\otimes1$
satisfies its original monic
polynomial over $k$. Thus all
their primes are maximal by
Lemma \ref{editorial-source-lemma-integral-no-inclusion},
and are also minimal since
there are no strict inclusions
between primes over a field.
The two original comparison maps
$$
B\otimes_{k'}K\longrightarrow A,
\qquad(a\otimes b)\otimes c\longmapsto a\otimes bc,
$$
and $a\otimes c\mapsto(a\otimes1)\otimes c$
are inverse by $k'$-balancing.
Since $K/k'$ is geometrically
irreducible, Lemma
\ref{editorial-source-lemma-geometrically-irreducible-any-base-change}
therefore identifies the spectra
of $A$ and $B$, including their
topologies by its minimal-prime
homeomorphism. The map is
contraction, and commutes with
the original automorphisms
$g\otimes1$ for $g\in G$.

For a prime $\mathfrak p$ of $B$,
its residue field is separably
algebraic over the original
$\overline k$. Indeed $B$ is
generated over $\overline k$ by
$1\otimes b$ for $b\in k'$,
and each of their images satisfies
the original separable polynomial
of $b$ over $k$. Minimal
polynomials over $\overline k$
divide these separable polynomials.
Every element of the field
$B/\mathfrak p$ belongs to a
finite subextension generated
by finitely many such images,
so that extension is separable.
As $\overline k$ is separably
closed, its structural embedding
into $B/\mathfrak p$ is an
isomorphism. Its inverse composed
with $b\mapsto[1\otimes b]$
defines a unique embedding
$\sigma:k'\to\overline k$ over $k$.
Consequently
$$
\mathfrak p=\mathfrak p_\sigma
=\ker(\mu_\sigma),\qquad
\mu_\sigma:B\to\overline k,\quad
\sum_i a_i\otimes b_i\longmapsto\sum_i a_i\sigma(b_i).
$$
Conversely this map is surjective
because $a=\mu_\sigma(a\otimes1)$,
so its kernel is a prime. The
structural field isomorphism
above proves uniqueness of
$\sigma$ from its kernel.
This gives the stated bijection
with $\operatorname{Hom}_k(k',\overline k)$.
Separability, rather than
maximality alone, is what
identifies these residue fields
with the separably closed field.

Given embeddings $\sigma,\tau$,
extend the isomorphism
$\tau\sigma^{-1}:\sigma(k')\to\tau(k')$
to a $k$-embedding of $\overline k$
in itself. Here are the needed
extension details, including
infinite $k'/k$. Order partial
extensions by inclusion and
take unions on chains. If a
maximal partial extension has
domain $L\neq\overline k$, choose
$a\in\overline k\setminus L$.
Its minimal polynomial over
$L$ divides a separable
polynomial over $k$, hence is
separable. Transporting its
coefficients gives a separable
polynomial over the image field,
which has a root in the
separably closed $\overline k$.
Sending $a$ to that root extends
the embedding to $L(a)$, a
contradiction. Thus the domain
is all of $\overline k$.
Any $k$-embedding of this field
in itself is surjective: each
$b\in\overline k$ lies in a
finite Galois subextension $N/k$,
obtained by adjoining all roots
of its separable minimal polynomial.
The embedding permutes those
roots, sends $N$ into itself,
and its injective $k$-linear
map on the finite-dimensional
$N$ is surjective. In particular
$b$ is in its image. We have
therefore constructed $g\in G$
with $g\sigma=\tau$.

Use the action on ideals
$\mathfrak p\mapsto(g\otimes1)(\mathfrak p)$.
For every original $z\in B$,
$$
\mu_{g\sigma}((g\otimes1)z)=g(\mu_\sigma(z)).
$$
It follows that
$(g\otimes1)(\mathfrak p_\sigma)=\mathfrak p_{g\sigma}$.
This proves transitivity on
$\Spec(B)$ and, by equivariant
contraction, on the original
$\Spec(A)$.

The correspondence also identifies
the actual prime and residue
field of $A$. Give $\overline k$
its $k'$-structure through
$\sigma$. The extension
$\overline k/\sigma(k')$ is
separable algebraic. The tensor-field
form of Lemma
\ref{editorial-source-lemma-geometrically-irreducible-separable-elements}
applied to $K/k'$ therefore makes
$$
C_\sigma=\overline k\otimes_{k',\sigma}K
$$
a field. The original quotient
comparison sends the class of
$a\otimes x$ to $a\otimes x$
and gives an isomorphism
$A/\mathfrak p_\sigma A\simeq C_\sigma$;
its inverse sends the same pure
tensor to its class. The additional
$k'$-balancing relations are
exactly the extended kernel
of $\mu_\sigma$, so both maps
are well-defined. Hence the
corresponding prime is exactly
$\mathfrak q_\sigma=\mathfrak p_\sigma A$,
and its residue field is
$C_\sigma$. No further radical
or residue-field identification
is needed or assumed.

For a fixed $\sigma$, its
stabilizer in $G$ is exactly
the subgroup fixing $\sigma(k')$
pointwise, by the injectivity
of the embedding-to-prime
correspondence. Thus the orbit
is the corresponding coset
set. If $K/k$ is finitely
generated, the preceding lemma
makes $k'/k$ finite separable;
Fields, Lemma
\ref{fields-lemma-separable-equality}
then gives exactly $[k':k]$
primes of $A$. Since this
finite spectrum has every
point closed, each point is
also open: its complement is
a finite union of closed
points. It is therefore a
finite discrete orbit. No
finite-degree count is asserted
for an infinite $k'/k$.
\end{proof}


\subsection{Editorial treatment of Algebra lines 11487--11528}

\hypertarget{algebra-context-149}{}

\noindent\href{https://github.com/stacks/stacks-project/blob/a04446e57ec1fbc252a871afcec7752fb2807b14/algebra.tex\#L11487}{Original source, lines 11487--11528}. Complete original and reviewed TeX passages accompany this note. Every intervention is separately recorded; the following treatment is editorial.


\begin{lemma}
\label{editorial-source-lemma-separably-closed-connected}
Let $k$ be a separably closed field.
Let $R$, $S$ be $k$-algebras. If $\Spec(R)$ and
$\Spec(S)$ are connected, then so is
$\Spec(R \otimes_k S)$.
\end{lemma}

\begin{proof}
The source convention for connectedness
includes nonemptiness (Topology,
Definition \ref{topology-definition-connected-components}).
Thus the original $R$ and $S$ are
nonzero, and Lemma
\ref{algebra-lemma-characterize-spec-connected}
applies to them. Their tensor is
also nonzero, since extending $1$
in each factor to a $k$-basis
shows $1\otimes1\neq0$.

Suppose this tensor had a nontrivial
idempotent $e$. A finite tensor
expression for this original $e$
uses coefficients in finite type
$k$-subalgebras $R_0\subseteq R$
and $S_0\subseteq S$. These are
nonzero and have no nontrivial
idempotents, since their inclusions
in the original factors preserve
idempotents and reflect $0$ and $1$.
They therefore have connected spectra.
The map $R_0\otimes_k S_0\to R\otimes_k S$
is injective by applying the two
vector-space injections successively.
It follows that the same $e$ is
a nontrivial idempotent in this
finite type tensor: its equation
$e^2-e=0$ is reflected by that
injection. Thus it suffices to
prove connectedness for the actual
finite type subalgebras, with
the original $e$ and all its
coefficients retained.

For the finite type assertion,
use $R,S$ for the two algebras
under consideration. They are
Noetherian and have finite
nonzero numbers $n,m$ of
irreducible components by Lemma
\ref{editorial-source-lemma-Noetherian-irreducible-components}.
We induct on $n+m$. For $n=m=1$,
Lemma \ref{editorial-source-lemma-separably-closed-irreducible}
gives irreducibility, hence
connectedness, of the tensor.
If $n=1$ and $m>1$, interchange
the factors through the inverse
maps $r\otimes s\mapsto s\otimes r$
and $s\otimes r\mapsto r\otimes s$.
It remains to handle $n>1$.

Here are the component-selection
details of Topology, Lemma
\ref{topology-lemma-remove-irreducible-connected}.
Form the finite graph whose
vertices are the original
irreducible components and whose
edges record nonempty intersection.
It is connected: otherwise the
unions over two parts with no
edges between them would be
disjoint nonempty closed sets
covering the spectrum, a separation.
Choose a spanning tree, obtained
by adding adjacent new vertices
until every vertex is reached.
A leaf exists by taking an
endpoint of a longest simple
path in this finite tree.
Remove that leaf. The remaining
tree is connected, so the union
of its connected vertex spaces
is connected: add them in tree
order, each meeting the existing
union. This proves the required
choice without changing any
irreducible component.

Let $T=V(\mathfrak p)$ be the
chosen component and let $T'$
be the union of the other
$n-1$ components. Take the
original ideal $I$ to be the
intersection of their minimal
primes. Then $T'=V(I)$, by
finite union of closed sets,
and its irreducible components
are those same $n-1$ components:
their defining primes are
pairwise incomparable. Both
$T$ and $T'$ are connected
and nonempty. They meet,
since otherwise these two
closed sets would separate
the original connected spectrum.
Their intersection is exactly
$V(\mathfrak p+I)$, so
$R/(\mathfrak p+I)\neq0$.

Under the original spectrum map,
the inverse images of $T$ and
$T'$ are respectively
$$
\Spec((R/\mathfrak p)\otimes_k S),\qquad
\Spec((R/I)\otimes_k S).
$$
For each of the ideals $J$
involved, the quotient comparison
sends $[r\otimes s]$ to
$[r]\otimes s$ and has inverse
$[r]\otimes s\mapsto[r\otimes s]$,
identifying $(R\otimes_k S)/J(R\otimes_k S)$
with $(R/J)\otimes_k S$.
The induction parameters for
these two inverse images are
$1+m$ and $(n-1)+m$, both
smaller than $n+m$; hence
both are connected. Their
intersection is the inverse
image of $T\cap T'$, namely
$$
\Spec((R/(\mathfrak p+I))\otimes_k S).
$$
It is nonempty because both
original tensor factors are
nonzero vector spaces and
their unit tensor is nonzero.
The two inverse images cover
the spectrum of $R\otimes_k S$.
A union of two connected sets
with nonempty intersection is
connected: a separation would
put each in one part, and
the common point forces the
same part for both.
This finishes the induction
and excludes the original
nontrivial idempotent $e$.
\end{proof}


\subsection{Editorial treatment of Algebra lines 11530--11558}

\hypertarget{algebra-context-150}{}

\noindent\href{https://github.com/stacks/stacks-project/blob/a04446e57ec1fbc252a871afcec7752fb2807b14/algebra.tex\#L11530}{Original source, lines 11530--11558}. Complete original and reviewed TeX passages accompany this note. Every intervention is separately recorded; the following treatment is editorial.


\begin{lemma}
\label{editorial-source-lemma-geometrically-connected}
Let $k$ be a field.
Let $R$ be a $k$-algebra.
The following are equivalent
\begin{enumerate}
\item for every field extension $k'/k$ the
spectrum of $R \otimes_k k'$ is connected, and
\item for every finite separable field extension $k'/k$ the
spectrum of $R \otimes_k k'$ is connected.
\end{enumerate}
\end{lemma}

\begin{proof}
Condition (1) implies (2).
Assume (2); its identity field
extension says that $\Spec(R)$
is connected and in particular
nonempty, so $R\neq0$.
Every scalar extension considered
below is nonzero by extension
of a vector-space basis.
For these nonzero rings, Lemma
\ref{algebra-lemma-characterize-spec-connected}
identifies connectedness with
absence of nontrivial idempotents.
The nonzero condition cannot be
omitted from that criterion:
the zero ring has no nontrivial
idempotents but has empty,
nonconnected spectrum under
Topology, Definition
\ref{topology-definition-connected-components}.

Let $k_s/k$ be the original
separable algebraic closure.
Every finite list of its
elements lies in a finite
separable subextension $E/k$.
The original coefficient maps
give the injective directed union
$$
R\otimes_k k_s=\bigcup_E R\otimes_k E.
$$
An idempotent in this union
lies in one such tensor, and
its equation is already true
there by injectivity. By (2)
it equals $0$ or $1$ there,
hence in the union. This
proves connectedness over $k_s$.

For any original field $L/k$,
choose a prime $\mathfrak a$
of the nonzero ring $L\otimes_k k_s$
and put
$$
F=\operatorname{Frac}((L\otimes_k k_s)/\mathfrak a).
$$
Its original field embeddings
$a\mapsto[a\otimes1]$ and
$b\mapsto[1\otimes b]$ are
injective and agree on $k$.
The scalar comparison maps
$$
(R\otimes_k k_s)\otimes_{k_s}F\longrightarrow R\otimes_k F,
\qquad(r\otimes a)\otimes b\longmapsto r\otimes ab,
$$
and $r\otimes b\mapsto(r\otimes1)\otimes b$
are inverse by balancing.
The source has connected
spectrum by Lemma
\ref{editorial-source-lemma-separably-closed-connected},
since both $R\otimes_k k_s$
and the field $F$ do.
The injection
$R\otimes_k L\to(R\otimes_k L)\otimes_L F\simeq R\otimes_k F$
reflects nontrivial idempotents,
again by a vector-space basis
and the displayed scalar maps.
Thus $R\otimes_k L$ is nonzero
and has no nontrivial idempotents.
It is connected, proving (1).

The same argument proves that
connectedness after just the
separable closure is an equivalent
test. One may instead use the
algebraic closure $\overline k$:
the extension $\overline k/k_s$
is purely inseparable, or the
identity in characteristic zero.
The original injection
$R\otimes_k k_s\to R\otimes_k\overline k$
is a homeomorphism on spectra
by Lemma \ref{editorial-source-lemma-p-ring-map}
and the same scalar comparison,
so it preserves and reflects
connectedness. Also finite Galois
extensions suffice in (2): every
finite separable $E/k$ embeds
in a finite Galois $N/k$, and
the injection $R\otimes_k E\to R\otimes_k N$
reflects nontrivial idempotents
and nonzeroness.

Finally, for every original $L/k$,
geometric connectedness of $R/k$
is equivalent to geometric
connectedness of $(R\otimes_k L)/L$.
The forward direction uses
the displayed scalar comparison
for any $F/L$. For the reverse
direction and any $E/k$, construct
a common receiving field for
$E$ and $L$ by the same
prime-quotient construction.
Geometric connectedness over
$L$ gives connectedness of
$R\otimes_k F$, and its
injected subalgebra $R\otimes_k E$
is nonzero and has no nontrivial
idempotents. This proves the
reverse implication with the
original scalar actions.
\end{proof}


\subsection{Editorial treatment of Algebra lines 11573--11585}

\hypertarget{algebra-context-151}{}

\noindent\href{https://github.com/stacks/stacks-project/blob/a04446e57ec1fbc252a871afcec7752fb2807b14/algebra.tex\#L11573}{Original source, lines 11573--11585}. Complete original and reviewed TeX passages accompany this note. Every intervention is separately recorded; the following treatment is editorial.


\begin{lemma}
\label{editorial-source-lemma-separably-closed-connected-implies-geometric}
Let $k$ be a field.
Let $R$ be a $k$-algebra.
If $k$ is separably closed then $R$ is
geometrically connected over $k$ if and only if the
spectrum of $R$ is connected.
\end{lemma}

\begin{proof}
Geometric connectedness implies
connectedness by the identity field
extension in Definition
\ref{algebra-definition-geometrically-connected}.
Conversely, if $\Spec(R)$ is
connected, then $R\neq0$ under
the source nonemptiness convention.
For any finite separable extension
$E/k$, separably closedness of $k$
makes the given embedding $\iota:k\to E$
an isomorphism: each element has
a separable minimal polynomial,
which must be linear over $k$.
The inverse maps
$R\otimes_k E\to R$, $r\otimes a\mapsto r\iota^{-1}(a)$,
and $r\mapsto r\otimes1$ show
that all the finite separable
tests have connected spectrum.
Lemma \ref{editorial-source-lemma-geometrically-connected}
proves geometric connectedness.
For $R=0$, neither side is
connected, so the same equivalence
holds without changing the definition.
\end{proof}


\subsection{Editorial treatment of Algebra lines 11587--11604}

\hypertarget{algebra-context-152}{}

\noindent\href{https://github.com/stacks/stacks-project/blob/a04446e57ec1fbc252a871afcec7752fb2807b14/algebra.tex\#L11587}{Original source, lines 11587--11604}. Complete original and reviewed TeX passages accompany this note. Every intervention is separately recorded; the following treatment is editorial.


\begin{lemma}
\label{editorial-source-lemma-subalgebra-geometrically-connected}
Let $k$ be a field. Let $S$ be a $k$-algebra.
\begin{enumerate}
\item If $S$ is geometrically connected over $k$ so is every
$k$-subalgebra.
\item If all finitely generated $k$-subalgebras of $S$ are
geometrically connected, then $S$ is geometrically connected.
\item A directed colimit of geometrically connected $k$-algebras
is geometrically connected.
\end{enumerate}
\end{lemma}

\begin{proof}
Use throughout the nonzero-ring
criterion of Lemma
\ref{algebra-lemma-characterize-spec-connected}.
For (1), the original inclusion
$S'\subseteq S$ gives an injection
$S'\otimes_k L\to S\otimes_k L$
for every field $L/k$ by extension
of a vector-space basis. Its
target is nonzero, and the
unital source is therefore
nonzero as well. Any nontrivial
idempotent of the source would
remain nontrivial in the target,
which is impossible. This proves (1).

For (3), retain the original
nonempty directed system $(S_i)$
and its unital transition maps;
no injectivity is required.
For every $L/k$, the maps
$$
(\mathop{\rm colim}_i S_i)\otimes_k L
\longrightarrow\mathop{\rm colim}_i(S_i\otimes_k L),
\qquad[i,s]\otimes a\longmapsto[i,s\otimes a]
$$
and $[i,\sum_j s_j\otimes a_j]\mapsto\sum_j[i,s_j]\otimes a_j$
are well-defined inverses:
all representatives and balancing
equations involve finitely many
terms and can be compared at
one common later stage.
Each original $S_i\otimes_k L$
is nonzero. The colimit is
nonzero, since $1=0$ there
would hold at a later stage.
If $e$ is an idempotent in
the colimit, represent it at
a stage $i$. Its equation
$e^2-e=0$ holds at some later
stage $j$, making that later
representative idempotent.
It is $0$ or $1$ there by
connectedness, and hence is
$0$ or $1$ in the colimit.
This proves connectedness after
each $L/k$, and gives (3).

For (2), the finite type
$k$-subalgebras form a directed
system under adjoining the union
of two finite lists of generators,
and their union is the original
$S$. Apply (3). The hypotheses
exclude $S=0$, since its
subalgebra generated by the
image of $k$ would already
be zero and would fail
geometric connectedness.

Unlike geometric irreducibility,
geometric connectedness need
not survive a nonzero localization.
For the exact example
$$
S=k[x,y]/(xy),\qquad s=x+y,
$$
each field extension gives
two affine lines meeting at
their origin. More explicitly,
$V(x)$ and $V(y)$ cover its
spectrum, are irreducible
because their quotients are
$L[y]$ and $L[x]$, and meet
in the nonempty $\Spec(L)$.
Their union is therefore
connected. But in the original
$S_s$ the element
$$
e=\frac{x}{x+y},\qquad
e^2-e=\frac{-xy}{(x+y)^2}=0
$$
is a nontrivial idempotent.
The map
$$
S_s\longrightarrow k[x,x^{-1}]\times k[y,y^{-1}],
\qquad\frac{h(x,y)}{(x+y)^n}
\longmapsto\left(\frac{h(x,0)}{x^n},\frac{h(0,y)}{y^n}\right)
$$
sends it to $(1,0)$, so it
is neither $0$ nor $1$.
This map is in fact an
isomorphism with the following
explicit inverse. In $eS_s$,
whose unit is $e$, send $x$
to the original $x=ex$ and
$x^{-1}$ to $e/s$; their
product is $e$ because
$x(e/s)=x^2/s^2=e$.
In $(1-e)S_s$, send $y$ to
$y=(1-e)y$ and $y^{-1}$ to
$(1-e)/s$; their product is
$1-e$. Add the images of
the two Laurent polynomials.
Orthogonality of $e,1-e$
makes this a ring map.
The composites fix both
component units and Laurent
generators. On the original
side they fix $x,y$ and
$s^{-1}$, since
$e/s+(1-e)/s=s^{-1}$.
Thus both composites are
identities, proving the full
product description and the
failure of localization stability.
\end{proof}


\subsection{Editorial treatment of Algebra lines 11610--11649}

\hypertarget{algebra-context-153}{}

\noindent\href{https://github.com/stacks/stacks-project/blob/a04446e57ec1fbc252a871afcec7752fb2807b14/algebra.tex\#L11610}{Original source, lines 11610--11649}. Complete original and reviewed TeX passages accompany this note. Every intervention is separately recorded; the following treatment is editorial.


\begin{lemma}
\label{editorial-source-lemma-geometrically-connected-any-base-change}
Let $k$ be a field.
Let $S$ be a geometrically connected $k$-algebra.
Let $R$ be any $k$-algebra.
The map
$$
R \longrightarrow R \otimes_k S
$$
induces a bijection on idempotents, and the map
$$
\Spec(R \otimes_k S) \longrightarrow \Spec(R)
$$
induces a bijection on connected components.
\end{lemma}

\begin{proof}
Put $A=R\otimes_k S$ and retain
$f:X=\Spec(A)\to Y=\Spec(R)$.
Geometric connectedness of the
original $S$ includes nonemptiness,
so $S\neq0$. If $R=0$, both
spectra are empty and both
rings have one idempotent,
so the two asserted bijections
hold. Assume $R\neq0$.
For every $\mathfrak p\in Y$,
the original fibre comparison
$$
A\otimes_R\kappa(\mathfrak p)\longrightarrow
\kappa(\mathfrak p)\otimes_k S,
\qquad(r\otimes s)\otimes a\longmapsto ra\otimes s
$$
has inverse $a\otimes s\mapsto(1\otimes s)\otimes a$.
The receiving spectrum is
connected by the hypothesis
on $S$ and is nonempty.
Hence $f$ is surjective with
connected fibres. It is open
without finite type assumptions,
by Lemma
\ref{editorial-source-lemma-map-into-tensor-algebra-open}.

Let $C\subseteq X$ be open
and closed. Connectedness
of each fibre implies that
the fibre is contained either
in $C$ or in its complement:
otherwise their intersections
would separate the fibre.
The images $f(C)$ and $f(X\setminus C)$
are therefore disjoint open
sets covering $Y$. In
particular $f(C)$ is open
and closed, and
$$
C=f^{-1}(f(C)).
$$
Conversely, inverse images of
open and closed subsets are
open and closed, and surjectivity
gives $f(f^{-1}(D))=D$ for
every such $D\subseteq Y$.
This proves a bijection of
the original open-and-closed
subsets in both directions.

By Lemma \ref{algebra-lemma-disjoint-decomposition},
the map $e\mapsto D(e)$ is
a bijection between idempotents
of each ring and its open-and-closed
subsets. Under the original
ring map $R\to A$, the
identity $f^{-1}(D(e))=D(e\otimes1)$
therefore proves exactly the
claimed idempotent bijection.
The inverse takes an idempotent
$a\in A$ to the unique
idempotent $e\in R$ such
that $D(e)=f(D(a))$.
These maps also preserve
complements and intersections,
represented by $1-e$ and
$ee'$, so give an isomorphism
of the corresponding Boolean
algebras. They commute with
arbitrary original $k$-algebra
base change: the image of
$e\otimes1$ is $\theta(e)\otimes1$,
and uniqueness of descent
gives the inverse compatibility.

We prove the component assertion
directly as well. For any
connected component $D$ of
$Y$, the restriction
$f^{-1}(D)\to D$ is open,
surjective and has the same
connected fibres. Openness
holds because for every open
$U\subseteq X$,
$$
f(U\cap f^{-1}(D))=f(U)\cap D.
$$
If $f^{-1}(D)$ were separated
into disjoint nonempty open
and closed parts, the same
fibre argument would give
a separation of $D$ by
their open images. Thus
$f^{-1}(D)$ is connected
and nonempty. It is maximal
connected: any larger connected
subset would have connected
image contained in a component
of $Y$, which must be $D$
because it meets $D$.
Conversely each component
of $X$ has connected image
in some component $D$ of
$Y$ and hence equals this
connected inverse image.
This proves the claimed
bijection and the exact
inverse-image description,
as in Topology, Lemma
\ref{topology-lemma-connected-fibres-connected-components}.

With the quotient topologies
on the sets of components,
the bijection is a homeomorphism.
Write $q_X,q_Y$ for the two
quotient maps and $h$ for
the induced bijection.
The relation $h q_X=q_Y f$
proves continuity by the
quotient property. If $V$
is open in the component
space of $X$, then
$$
q_Y^{-1}(h(V))=f(q_X^{-1}(V))
$$
by the exact component
correspondence. The right
side is open because $f$
is open, proving that $h$
is open and hence a
homeomorphism. None of this
requires individual components
to be open in their spectra.
\end{proof}


\subsection{Editorial treatment of Algebra lines 11673--11683}

\hypertarget{algebra-context-154}{}

\noindent\href{https://github.com/stacks/stacks-project/blob/a04446e57ec1fbc252a871afcec7752fb2807b14/algebra.tex\#L11673}{Original source, lines 11673--11683}. Complete original and reviewed TeX passages accompany this note. Every intervention is separately recorded; the following treatment is editorial.


\begin{lemma}
\label{editorial-source-lemma-geometrically-integral}
Let $k$ be a field.
Let $S$ be a $k$-algebra.
In this case $S$ is geometrically integral over $k$ if and only if
$S$ is geometrically irreducible as well as geometrically reduced over $k$.
\end{lemma}

\begin{proof}
For any original field extension $L/k$,
put $A=S\otimes_k L$. If $A$ is
a domain, it is nonzero and reduced,
and its zero ideal is prime. Its
spectrum is therefore irreducible.
Conversely, if $A$ is reduced and
its spectrum is irreducible, then
it is nonzero and its nilradical
is prime by Lemma \ref{editorial-source-lemma-irreducible}.
Reducedness makes that nilradical
the zero ideal, so the zero ideal
is prime and $A$ is a domain.
Apply this exact ring equivalence
for every original $L/k$ to obtain
both directions of the assertion.
Nonemptiness comes from irreducibility;
reducedness alone would allow the
zero ring and would not suffice.
\end{proof}


\subsection{Editorial treatment of Algebra lines 11685--11702}

\hypertarget{algebra-context-155}{}

\noindent\href{https://github.com/stacks/stacks-project/blob/a04446e57ec1fbc252a871afcec7752fb2807b14/algebra.tex\#L11685}{Original source, lines 11685--11702}. Complete original and reviewed TeX passages accompany this note. Every intervention is separately recorded; the following treatment is editorial.


\begin{lemma}
\label{editorial-source-lemma-characterize-geometrically-integral}
Let $k$ be a field. Let $S$ be a $k$-algebra.
The following are equivalent
\begin{enumerate}
\item $S$ is geometrically integral over $k$,
\item for every finite extension $k'/k$ of fields
the ring $S \otimes_k k'$ is a domain,
\item $S \otimes_k \overline{k}$ is a domain
where $\overline{k}$ is the algebraic closure of $k$.
\end{enumerate}
\end{lemma}

\begin{proof}
Condition (1) implies (2) by
restricting the field extensions.
Assume (2). In particular its
identity extension makes $S$
a nonzero domain. Choose the
original algebraic closure
$\overline k$ and consider its
finite subextensions $E/k$.
They are directed under compositum
and contain every finite list
of coefficients of a tensor.
The original injective coefficient
maps therefore give
$$
S\otimes_k\overline k=\bigcup_E S\otimes_k E.
$$
Each ring in this union is a
domain by (2). The union is
nonzero, since the maps preserve
the nonzero unit. For two
nonzero elements, choose one
stage containing both. Their
product there is nonzero and
remains so by injectivity.
This proves (3), including
nonemptiness and all zero-divisor
conditions.

Assume (3), and let $L/k$
be any original extension.
Choose a prime $\mathfrak a$
of the nonzero field tensor
$L\otimes_k\overline k$ and
form the compatible receiving
field
$$
F=\operatorname{Frac}((L\otimes_k\overline k)/\mathfrak a).
$$
The maps $a\mapsto[a\otimes1]$
and $b\mapsto[1\otimes b]$
embed the two original fields
and agree on $k$. Write
$B=S\otimes_k\overline k$.
This is a nonzero domain
by (3). Since $\overline k$
is perfect, Lemma
\ref{editorial-source-lemma-perfect-reduced}
makes $B\otimes_{\overline k}F$
reduced. Since $\overline k$
is separably closed and both
$B$ and $F$ have irreducible
spectra, Lemma
\ref{editorial-source-lemma-separably-closed-irreducible}
makes its spectrum irreducible.
Thus it is a domain by Lemma
\ref{editorial-source-lemma-geometrically-integral}'s
underlying ring criterion.

For either original subfield
$H=L$ or $H=\overline k$ of
the receiving $F$, the maps
$$
(S\otimes_k H)\otimes_H F\longrightarrow S\otimes_k F,
\qquad(s\otimes a)\otimes b\longmapsto s\otimes ab
$$
and $s\otimes b\mapsto(s\otimes1)\otimes b$
are inverse by balancing.
They identify the domain just
proved with $S\otimes_k F$.
The original scalar map
$S\otimes_k L\to(S\otimes_k L)\otimes_L F$
is injective by extension of
an $L$-vector-space basis.
Its source is therefore a
nonzero subring of that domain,
so is a domain. This proves
(1) for every original $L/k$.

The same receiving-field maps
prove geometric integrality
is preserved and reflected
by every field extension
$L/k$. Preservation follows
from the displayed comparison
for all $F/L$. For reflection,
suppose $S\otimes_k L$ is
geometrically integral over $L$.
For any $E/k$, construct a
compatible common field for
$E$ and $L$ as above. The
hypothesis makes $S\otimes_k F$
a domain, and the injection
$S\otimes_k E\to S\otimes_k F$
then makes the original
$S\otimes_k E$ a domain.
This is precisely geometric
integrality of $S/k$.
The zero algebra satisfies
none of (1)--(3), as each
condition includes a nonzero
domain after scalar extension.
\end{proof}


\subsection{Editorial treatment of Algebra lines 11704--11718}

\hypertarget{algebra-context-156}{}

\noindent\href{https://github.com/stacks/stacks-project/blob/a04446e57ec1fbc252a871afcec7752fb2807b14/algebra.tex\#L11704}{Original source, lines 11704--11718}. Complete original and reviewed TeX passages accompany this note. Every intervention is separately recorded; the following treatment is editorial.


\begin{lemma}
\label{editorial-source-lemma-geometrically-integral-any-integral-base-change}
Let $k$ be a field. Let $S$ be a geometrically integral $k$-algebra.
Let $R$ be a $k$-algebra and an integral domain. Then $R \otimes_k S$
is an integral domain.
\end{lemma}

\begin{proof}
Let $F=\operatorname{Frac}(R)$.
The original injection $R\to F$
gives an injection
$$
R\otimes_k S\longrightarrow F\otimes_k S
$$
because tensoring $k$-vector spaces
preserves injections. The receiving
ring is a domain by geometric
integrality of $S/k$. The source
is nonzero: extending $1$ to
bases of the two original
nonzero factors shows
$1\otimes1\neq0$. Hence it
is a domain, as asserted.
The corresponding localization
maps retain the original
denominators $W=R\setminus\{0\}$:
$$
F\otimes_k S\longrightarrow W^{-1}(R\otimes_k S),
\qquad(r/w)\otimes s\longmapsto(r\otimes s)/(w\otimes1),
$$
and the inverse sends
$(\sum_i r_i\otimes s_i)/(w\otimes1)$
to $\sum_i(r_i/w)\otimes s_i$.
The fraction equivalence relations
and tensor balancing make these
well-defined inverses. The
displayed injection is exactly
the original constant-fraction
map under this comparison.

The argument gives precise
ideal consequences for any
original $k$-algebra $R$, now
without a domain hypothesis.
Write $A=R\otimes_k S$ and
let $I\subseteq R$ be an
arbitrary ideal. The maps
$[r\otimes s]\mapsto[r]\otimes s$
and $[r]\otimes s\mapsto[r\otimes s]$
identify
$$
A/IA\simeq(R/I)\otimes_k S.
$$
Choose a $k$-linear functional
$\ell:S\to k$ with $\ell(1)=1$.
The map $\operatorname{id}\otimes\ell$
splits $R/I\to(R/I)\otimes_k S$
as a linear map, proving
$IA\cap R=I$ for every $I$.
If $I$ is prime, $R/I$ is
a domain and the argument
above makes $A/IA$ a domain,
so $IA$ is prime. Conversely,
if $IA$ is prime, its
contraction $I$ is prime.
This proves the exact equivalence
for the original extended ideal,
including the exclusion of
the unit ideal from prime ideals.

There is also the full radical
comparison
$$
\sqrt{IA}=(\sqrt I)A.
$$
To prove it, keep the quotient
map
$$
(R/I)\otimes_k S\longrightarrow(R/\sqrt I)\otimes_k S.
$$
Its kernel is the ideal
generated by the original
$\sqrt I/I$ factor, by the
same explicit quotient maps.
Each element of this kernel
is a finite sum of tensors
$r_j\otimes s_j$ with $r_j$
nilpotent. If $r_j^{d_j}=0$,
its power of degree
$N=1+\sum_j(d_j-1)$ vanishes:
in the full multinomial sum
$$
\sum_{i_1+\cdots+i_m=N}
\frac{N!}{i_1!\cdots i_m!}
\prod_{j=1}^m(r_j\otimes s_j)^{i_j},
$$
some $i_j\geq d_j$ in every
summand. The displayed coefficient
is an integer; no factorial
is inverted in the receiving ring.
The target is reduced by Lemma
\ref{editorial-source-lemma-geometrically-reduced-any-reduced-base-change},
since $S$ is geometrically reduced
and $R/\sqrt I$ is reduced.
Thus this kernel is exactly
the nilradical of $A/IA$.
Taking its inverse image in
$A$ proves the radical formula.

In particular for every prime
$\mathfrak p$ of $R$ the ideal
$\mathfrak pA$ is already prime,
so the continuous section of
Lemma
\ref{editorial-source-lemma-geometrically-irreducible-any-base-change}
has the exact value
$\sigma(\mathfrak p)=\mathfrak pA$.
Its contraction, topology,
base-change compatibility and
minimal-prime correspondence
are those already proved for
the original tensor map;
the stronger integrality
hypothesis removes no terms
from the original ideal.
\end{proof}


\subsection{Editorial treatment of Algebra lines 11723--12149}

\hypertarget{algebra-context-157}{}

\noindent\href{https://github.com/stacks/stacks-project/blob/a04446e57ec1fbc252a871afcec7752fb2807b14/algebra.tex\#L11723}{Original source, lines 11723--12149}. Complete original and reviewed TeX passages accompany this note. Every intervention is separately recorded; the following treatment is editorial.


\subsubsection{Valuation rings}
\label{editorial-source-section-valuation-rings}

\noindent
Here are some definitions.

\begin{definition}
\label{editorial-source-definition-valuation-ring}
Valuation rings.
\begin{enumerate}
\item Let $K$ be a field. Let $A$, $B$ be local rings contained
in $K$. We say that $B$ {\it dominates} $A$ if $A \subset B$
and $\mathfrak m_A = A \cap \mathfrak m_B$.
\item Let $A$ be a ring. We say $A$ is a {\it valuation ring}
if $A$ is a local domain and if $A$ is maximal
for the relation of domination among local rings contained in
the fraction field of $A$.
\item Let $A$ be a valuation ring with fraction field $K$.
If $R \subset K$ is a subring of $K$, then we say $A$
is {\it centered} on $R$ if $R \subset A$.
\end{enumerate}
\end{definition}

\noindent
With this definition a field is a valuation ring.

\begin{lemma}
\label{editorial-source-lemma-dominate}
Let $K$ be a field. Let $A \subset K$ be a local subring.
Then there exists a valuation ring with fraction field $K$
dominating $A$.
\end{lemma}

\begin{proof}
Consider the set of local subrings
$B\subseteq K$ dominating the
original $A$, ordered by domination.
It is nonempty since it contains
$A$. The relation is transitive:
if $B$ dominates $A$ and $C$
dominates $B$, then
$A\cap\mathfrak m_C=A\cap\mathfrak m_B=\mathfrak m_A$.
Antisymmetry follows from the
inclusions of the actual subrings.

For a nonempty chain $(B_i)$,
put $B=\bigcup_i B_i$ and
$\mathfrak n=\bigcup_i\mathfrak m_{B_i}$.
The maximal ideals are compatible
under contraction and inclusion.
A common chain member containing
any finite list of elements proves
that $B$ is a ring and
$\mathfrak n$ is an ideal.
It is proper because no member
contains $1$. If $b\in B\setminus\mathfrak n$,
it is a unit in any member
containing it, hence in $B$.
If $b\in\mathfrak n$, an inverse
in $B$ would put both elements
in a common later member and
contradict properness of its
maximal ideal. Thus $B$ is
local with maximal ideal
$\mathfrak n$, and
$B_i\cap\mathfrak n=\mathfrak m_{B_i}$
for every $i$. It dominates
all members and the original
$A$. The empty chain has
upper bound $A$. Zorn's lemma
therefore gives a maximal
member $V$.

We prove that its actual
fraction field $F=\operatorname{Frac}(V)$
equals $K$. Otherwise choose
$t\in K\setminus F$.
If $t$ is transcendental over
$F$, the evaluation map embeds
the original polynomial ring
$V[T]$ in $K$, with $T\mapsto t$.
The ideal $(\mathfrak m_V,t)$
of $V[t]$ is maximal, since
its quotient is $V/\mathfrak m_V$.
Its localization embeds in
$K$ by the original fraction
map and dominates $V$:
the maximal ideal contracts
to $\mathfrak m_V$. It contains
$t\notin V$, contradicting
maximality of $V$.

If $t$ is algebraic over
$F$, retain its monic equation
$$
t^d+c_1t^{d-1}+\cdots+c_d=0,
\qquad c_i\in F.
$$
Choose a nonzero $a\in V$
that is a product of denominators
for all original $c_i$, and
write $b_i=ac_i\in V$.
The original element $u=at$
satisfies the full equation
$$
u^d+b_1u^{d-1}+a b_2u^{d-2}
+\cdots+a^{d-1}b_d=0.
$$
This is exactly $a^d$ times
the displayed original equation,
using $a^{i-1}b_i=a^ic_i$;
every coefficient is retained.
Hence $V[u]$ is finite integral
over $V$, generated as a module
by $1,u,\ldots,u^{d-1}$.
By Lemma
\ref{editorial-source-lemma-integral-overring-surjective},
there is a prime $\mathfrak q$
of $V[u]$ contracting to
$\mathfrak m_V$. The original
fraction map embeds $V[u]_{\mathfrak q}$
in $K$, with maximal-ideal
contraction $\mathfrak m_V$.
It dominates $V$ and contains
$u\notin V$, since $u\in V$
would imply $t=u/a\in F$.
This is again a contradiction.
Thus $F=K$.

Every local subring of $K$
dominating $V$ also dominates
$A$. Maximality in the chosen
set therefore makes $V$ maximal
for domination in its own
fraction field $K$, exactly
the definition of a valuation
ring. This proves the lemma.

In particular, for any original
subring $R\subseteq K$ and
prime $\mathfrak p\subset R$,
apply the result to the
actual local subring $R_{\mathfrak p}$
of $K$. The resulting $V$
has fraction field $K$ and
$$
\mathfrak m_V\cap R
=(\mathfrak m_V\cap R_{\mathfrak p})\cap R
=\mathfrak pR_{\mathfrak p}\cap R
=\mathfrak p.
$$
Thus the center can be
prescribed as any given
prime of the original subring.
\end{proof}

\begin{lemma}
\label{editorial-source-lemma-valuation-ring-normal}
Let $A$ be a valuation ring.
Then $A$ is a normal domain.
\end{lemma}

\begin{proof}
Let $K=\operatorname{Frac}(A)$, retaining
the original fraction embedding. Suppose
$x\in K$ is integral over $A$.
The original $A'=A[x]\subseteq K$
is finite integral over $A$: a
monic equation of degree $d$ for
$x$ expresses all its powers in
the span of $1,x,\ldots,x^{d-1}$.
Lemma \ref{editorial-source-lemma-integral-overring-surjective}
gives a prime $\mathfrak m'$ of
$A'$ contracting to $\mathfrak m_A$.
The actual fraction map
$A'_{\mathfrak m'}\to K$, $a/s\mapsto a/s$,
is injective because $A'$ is a
domain and $s\notin\mathfrak m'$
is nonzero. The maximal ideal
of this localization contracts to
$\mathfrak m'$ in $A'$ and then
to $\mathfrak m_A$ in $A$.
Thus it dominates the original
$A$. The valuation-ring maximality
in this same $K$ gives
$A=A'_{\mathfrak m'}$, so $x\in A$.
Every element of the fraction field
integral over $A$ therefore belongs
to $A$, proving normality.
In particular any integral overring
$A\subseteq D\subseteq K$ equals
$A$, by applying this proof to
each of its original elements.
\end{proof}


\begin{lemma}
\label{editorial-source-lemma-valuation-ring-x-or-x-inverse}
Let $A$ be a valuation ring with maximal ideal $\mathfrak m$ and
fraction field $K$.
Let $x \in K^*$. Then either $x \in A$ or $x^{-1} \in A$ or both.
\end{lemma}

\begin{proof}
It suffices to treat the case
$x\notin A$; then $x\neq0$,
since $0\in A$. Retain the
original subring $A'=A[x]\subseteq K$.
If a prime $\mathfrak q$ of
$A'$ contracted to $\mathfrak m$,
the original fraction map would
embed $A'_{\mathfrak q}$ in $K$
as a local ring dominating
$A$. Maximality of the valuation
ring would give $A=A'_{\mathfrak q}$,
contrary to $x\notin A$.
There is therefore no such
prime. Every prime containing
$\mathfrak m A'$ would contract
to the maximal ideal $\mathfrak m$,
so $\mathfrak m A'$ is the
unit ideal. A finite expression
for its unit, expanded in
the original generator $x$,
gives
$$
1=\sum_{i=0}^d t_i x^i,\qquad t_i\in\mathfrak m.
$$
Here $d\geq1$, since $d=0$
would imply $1=t_0\in\mathfrak m$.
With $y=x^{-1}$, multiplication
by the full $x^{-d}$ gives
$$
(1-t_0)y^d-\sum_{i=1}^d t_i y^{d-i}=0.
$$
The index starts at $1$:
the original $t_0$ contribution
is already in $1-t_0$.
That coefficient is a unit
of the local $A$, so the
monic polynomial
$$
Q(T)=T^d-\sum_{i=1}^d(1-t_0)^{-1}t_iT^{d-i}
$$
annihilates $y$. The complete
original polynomial is
$(1-t_0)Q(T)$; its unit
factor and all coefficients
are kept in the comparison.
Thus $y$ is integral over
$A$ and belongs to $A$ by
Lemma \ref{editorial-source-lemma-valuation-ring-normal}.
This proves the assertion.

The inverse is in fact in
$\mathfrak m$: if $y$ were
a unit of $A$, then
$x=y^{-1}$ would belong to
$A$, contrary to the case
under consideration. Conversely,
the inverse of a nonzero
element of $\mathfrak m$
cannot belong to $A$, since
their product would put $1$
in $\mathfrak m$. Therefore
$$
K\setminus A=\{a^{-1}:0\neq a\in\mathfrak m\},
$$
and a nonzero $x$ and its
inverse both belong to $A$
exactly when $x\in A^*$.
The zero element belongs to
$A$ independently and is
never assigned an inverse.
\end{proof}

\begin{lemma}
\label{editorial-source-lemma-x-or-x-inverse-valuation-ring}
Let $A \subset K$ be a subring of a field $K$ such that for all
$x \in K^*$ either $x \in A$ or $x^{-1} \in A$ or both.
Then $A$ is a valuation ring with fraction field $K$.
\end{lemma}

\begin{proof}
The original $A$ is a nonzero
domain, as a unital subring
of the field $K$. The given
alternative proves that its
fraction field is all of $K$:
for nonzero $x\in K$, either
$x=x/1$ with $x\in A$ or
$x=1/x^{-1}$ with $x^{-1}\in A$.
These expressions are the
inverse of the canonical
injection $\operatorname{Frac}(A)\to K$.
The zero element is $0/1$.

If $A=K$, it is a local
ring with zero maximal ideal.
Otherwise it is not a field,
since a subfield equals its
own fraction field. Choose
a maximal ideal $\mathfrak m$
of $A$. If another maximal
ideal $\mathfrak m'$ were
distinct, their incomparability
would give
$$
x\in\mathfrak m\setminus\mathfrak m',\qquad
y\in\mathfrak m'\setminus\mathfrak m.
$$
Both original elements are
nonzero. If $x/y\in A$,
then $x=y(x/y)\in\mathfrak m'$,
a contradiction. If $y/x\in A$,
then $y=x(y/x)\in\mathfrak m$,
also a contradiction. The
given alternative excludes
two maximal ideals. Thus
$A$ is local, including
the field case.

Let the original local subring
$A'\subseteq K$ dominate $A$.
If $x\in A'\setminus A$,
then $x\neq0$ and the
hypothesis gives $x^{-1}\in A$.
It is not a unit of $A$,
since otherwise $x\in A$.
Hence $x^{-1}\in\mathfrak m_A$
and domination puts it in
$\mathfrak m_{A'}$. Multiplying
by the original $x\in A'$
would put $1$ in that ideal.
This contradiction gives
$A'=A$, proving the valuation
ring assertion in the exact
fraction field $K$.

This criterion also proves
that all ideals of a valuation
ring are linearly ordered
by inclusion. For nonzero
$a,b\in A$, applying the
alternative to $a/b$ gives
either $a=bc$ or $b=ac$
with $c\in A$. The zero
cases are immediate divisibility
equalities. Thus principal
ideals are comparable. If
ideals $I,J$ satisfy
$I\not\subseteq J$, choose
$a\in I\setminus J$.
For every $b\in J$, the
alternative $a=bc$ is
impossible because it puts
$a$ in $J$. The other
alternative gives $b\in(a)\subseteq I$;
this includes $b=0$.
Consequently $J\subseteq I$.
In particular all original
prime ideals are in one
inclusion chain. Conversely,
comparability of the original
principal ideals in a nonzero
domain gives the fraction
alternative in its own
fraction field, and the proof
above makes it a valuation ring.
\end{proof}

\begin{lemma}
\label{editorial-source-lemma-colimit-valuation-rings}
\begin{slogan}
Valuation rings are stable under filtered direct limits
\end{slogan}
Let $I$ be a directed set. Let $(A_i, \varphi_{ij})$
be a system of valuation rings over $I$.
Then $A = \colim A_i$ is a valuation ring.
\end{lemma}

\begin{proof}
The directed set is nonempty,
and all transition maps are
the original unital ring maps;
they need not be injective
or local. The colimit $A$
is nonzero: an equality
$1=0$ in it would hold
at a later stage, impossible
in a valuation domain $A_j$.
For $a,b\in A$ with $ab=0$,
choose representatives $a_i,b_i$
at one stage. At a later
stage their product is zero.
That stage is a domain,
so one of the two images
is zero and the corresponding
original $a$ or $b$ is zero.
Thus $A$ is a domain.

For arbitrary $a,b\in A$,
choose such representatives.
If either representative is
zero, one divides the other.
Otherwise Lemma
\ref{editorial-source-lemma-valuation-ring-x-or-x-inverse}
in $\operatorname{Frac}(A_i)$
gives either $a_i=b_i c_i$
or $b_i=a_i c_i$ with
$c_i\in A_i$. The same
original equality holds in
$A$ under the structure map.
For a fraction $a/b$ with
$b\neq0$, this proves that
$a/b\in A$ or, if $a\neq0$,
its inverse belongs to $A$.
Lemma
\ref{editorial-source-lemma-x-or-x-inverse-valuation-ring}
now proves the assertion.

There is an exact fraction-field
description even for the original
noninjective system. Put
$I_i=\ker(A_i\to A)$ and
$B_i=A_i/I_i$. Since $B_i$
injects in the domain $A$,
$I_i$ is prime and $B_i$
is a nonzero domain. The
divisibility equalities in
$A_i$ descend to $B_i$,
so the same criterion makes
$B_i$ a valuation ring.
The induced maps $B_i\to B_j$
are injective: their kernels
are zero because an element
vanishing in $B_j$ has zero
image in $A$, hence was
already in $I_i$.
These are the actual image
subrings of $A$, with
$A=\bigcup_i B_i$.

The field maps
$\operatorname{Frac}(B_i)\to\operatorname{Frac}(B_j)$
send each original $a_i/b_i$
to the ratio of its images;
denominators remain nonzero
by injectivity. They give
the isomorphism
$$
\mathop{\rm colim}_i\operatorname{Frac}(B_i)
\longrightarrow\operatorname{Frac}(A),
\qquad[i,a_i/b_i]\longmapsto a_i/b_i.
$$
Every fraction on the right
has both coefficients in
one common $B_i$, proving
surjectivity. Equality of two
fractions is the equality
of their cross products;
it holds in a common stage
because those stages inject
in $A$. This proves injectivity
and gives the inverse by
choosing a stage for the
two original coefficients.
No fraction-field map is
asserted for an original
transition map before its
kernel has been accounted for.

For a representative $a_i$,
its image in $A$ is a unit
exactly when its image is
a unit at some later stage.
One direction preserves the
original inverse equation.
For the other, an inverse
in $A$ has a representative
at a common stage, and
the equation giving product
$1$ holds at a later stage.
Hence the exact maximal-ideal
criterion is that every
later image of $a_i$ belong
to the corresponding maximal
ideal. A union of the
original maximal ideals need
not work: for the two-stage
system $V\to\operatorname{Frac}(V)$
with a nonfield valuation
ring $V$, each nonzero
element of $\mathfrak m_V$
becomes a unit at the
second stage. The colimit
is the field, with zero
maximal ideal, as the exact
criterion predicts.
\end{proof}

\begin{lemma}
\label{editorial-source-lemma-valuation-ring-cap-field}
Let $L/K$ be an extension of fields. If $B \subset L$
is a valuation ring, then $A = K \cap B$ is a valuation ring.
\end{lemma}

\begin{proof}
Retain the original fields $K\subseteq L$
and the original $B\subseteq L$.
Put $F=\operatorname{Frac}(B)\subseteq L$
and $E=K\cap F$. The ring in
the statement is exactly
$A=K\cap B=E\cap B$; these
are equal subrings of the
original $L$. For every
nonzero $x\in E$, the
valuation alternative in $B$
gives $x\in B$ or $x^{-1}\in B$.
Both $x,x^{-1}$ lie in $E$,
so the alternative holds in
$A$. Lemma
\ref{editorial-source-lemma-x-or-x-inverse-valuation-ring}
makes $A$ a valuation ring
whose exact fraction field
is $E=K\cap F$.
The canonical fraction map
$\operatorname{Frac}(A)\to E$
sends $a/b$ to the same
fraction in $L$. Its inverse
sends a nonzero $x$ to
$x/1$ if $x\in A$, or
to $1/x^{-1}$ otherwise;
the two possible expressions
agree when both apply.
This proves the fraction-field
identification without replacing
either original ambient field.

An original nonzero $a\in A$
is a unit of $A$ precisely
when it is a unit of $B$.
Indeed, if $a^{-1}\in B$,
then $a^{-1}\in K$ as well,
so it belongs to $A$.
The converse follows from
the inclusion. Hence
$$
\mathfrak m_A=A\cap\mathfrak m_B.
$$
The resulting residue map
$A/\mathfrak m_A\to B/\mathfrak m_B$
sends $[a]$ to $[a]$ and
is injective by this exact
kernel calculation. In particular
$B$ dominates the original
intersection $A$.

This also proves extension
of any given valuation ring
$A$ of a field $K$ to
any prescribed extension $L/K$.
Lemma \ref{editorial-source-lemma-dominate}
provides a valuation ring
$B\subseteq L$ of fraction
field $L$ dominating $A$.
Its intersection $A'=K\cap B$
has fraction field $K$ by
the proved formula and is
local with maximal ideal
$A'\cap\mathfrak m_B$.
It contains $A$ and its
maximal ideal contracts to
$\mathfrak m_A$, so it
dominates $A$. Maximality
of the original valuation
ring $A$ in $K$ forces
$A'=A$. Thus the extension
restricts to the original
ring and residue embedding
exactly, not merely to some
valuation ring containing it.
\end{proof}

\begin{lemma}
\label{editorial-source-lemma-valuation-ring-cap-field-finite}
Let $L/K$ be an algebraic extension of fields. If $B \subset L$
is a valuation ring with fraction field $L$ and not a field, then
$A = K \cap B$ is a valuation ring and not a field.
\end{lemma}

\begin{proof}
By Lemma \ref{editorial-source-lemma-valuation-ring-cap-field},
the original $A=K\cap B$ is
a valuation ring, and its
fraction field is exactly $K$
because $\operatorname{Frac}(B)=L$.
If $A$ were a field, it
would equal its fraction field,
so $A=K\subseteq B$.
Every $b\in B$ is algebraic
over $K$, hence integral over
this field by its monic
minimal polynomial. Thus
$K\subseteq B$ is integral.
Lemma \ref{editorial-source-lemma-integral-no-inclusion}
forbids a proper inclusion
of primes in $B$ over
the sole prime of $K$.
But the local domain $B$
has $(0)\subsetneq\mathfrak m_B$
because it is not a field,
a contradiction.

Equivalently the contradiction
has a direct inverse formula.
For nonzero $b\in B$, its
original monic minimal polynomial
$$
b^d+c_1b^{d-1}+\cdots+c_d=0
$$
has $c_d\neq0$ in $K$,
and yields
$$
b^{-1}=-c_d^{-1}
\left(b^{d-1}+\sum_{i=1}^{d-1}c_i b^{d-1-i}\right)\in B.
$$
For $d=1$ the parenthesis
is $1$. Thus every nonzero
$b$ would be a unit,
again making $B$ a field.

The algebraicity hypothesis
cannot be dropped. For any
field $k$, retain
$K=k$, $L=k(t)$ and
$B=k[t]_{(t)}$. Its fraction
field is the original $L$.
For a nonzero original
fraction $P(t)/Q(t)$, write
$P=t^rP_0$, $Q=t^sQ_0$
with $P_0(0),Q_0(0)\neq0$,
retaining both powers and
both residual factors.
If $r\geq s$, the fraction
$t^{r-s}P_0/Q_0$ lies in
$B$; if $r<s$, its inverse
$t^{s-r}Q_0/P_0$ lies in
$B$. The fraction criterion
therefore makes $B$ a
valuation ring. It is not
a field: evaluation at $t=0$
on this localization has
proper kernel $(t)B$, which
contains its nonzero $t$.
Nevertheless $K\cap B=k$
is a field. This identifies
the exact failure without
the original algebraicity
assumption.
\end{proof}

\begin{lemma}
\label{editorial-source-lemma-make-valuation-rings}
Let $A$ be a valuation ring. For any prime ideal $\mathfrak p \subset A$ the
quotient $A/\mathfrak p$ is a valuation ring. The same is true for the
localization $A_\mathfrak p$ and in fact any nonzero localization of $A$.
\end{lemma}

\begin{proof}
For the original prime
$\mathfrak p\subset A$, the
quotient $A/\mathfrak p$ is
a nonzero domain. Any two
nonzero residue classes have
original representatives $a,b$
outside $\mathfrak p$.
The valuation alternative in
$\operatorname{Frac}(A)$ gives
$a=bc$ or $b=ac$ for
some original $c\in A$.
The same equality in the
quotient gives its fraction
alternative, so Lemma
\ref{editorial-source-lemma-x-or-x-inverse-valuation-ring}
makes $A/\mathfrak p$ a
valuation ring. Its maximal
ideal is $\mathfrak m_A/\mathfrak p$:
a class outside that ideal
has an inverse inherited
from $A$. Its fraction field
is the original residue
field $\kappa(\mathfrak p)$,
through the inverse maps
$\overline a/\overline b\mapsto[a/b]$
and $[a/b]\mapsto\overline a/\overline b$
between $\operatorname{Frac}(A/\mathfrak p)$
and $A_{\mathfrak p}/\mathfrak pA_{\mathfrak p}$,
where $b\notin\mathfrak p$.

Let $W\subseteq A$ be an
original multiplicative subset.
If $0\in W$, Lemma
\ref{algebra-lemma-localization-zero}
gives $W^{-1}A=0$, which
is not a domain and hence
not a valuation ring.
This is an allowed localization
under Definition
\ref{algebra-definition-multiplicative-subset}.
If $0\notin W$, the
fraction map embeds $W^{-1}A$
in $K=\operatorname{Frac}(A)$:
equality of fractions is
exactly the original cross-product
equality, since nonzero elements
of $W$ do not annihilate
anything in the domain.
It contains $A$ and has
fraction field $K$. The
alternative $x\in A$ or
$x^{-1}\in A$ for every
nonzero $x\in K$ therefore
holds in $W^{-1}A$, making
it a valuation ring. This
includes $A_{\mathfrak p}$,
whose denominator set avoids zero.

More generally every original
overring $A\subseteq B\subseteq K$
is exactly $A_{\mathfrak q}$
for one prime $\mathfrak q$
of $A$. The same fraction
alternative first proves that
$B$ is a valuation ring
with fraction field $K$.
Put $\mathfrak q=A\cap\mathfrak m_B$.
The original fraction map
$A_{\mathfrak q}\to B$ is
defined because every denominator
outside $\mathfrak q$ is
a unit of $B$, and is
injective inside $K$.
For $b\in B$, if $b\in A$
it is already in $A_{\mathfrak q}$.
Otherwise $b\neq0$ and
$b^{-1}\in A$ by the
original valuation criterion.
It is a unit of $B$,
so $b^{-1}\notin\mathfrak q$
and $b=1/b^{-1}\in A_{\mathfrak q}$.
This proves equality of the
original subrings. The inverse
prime assignment is contraction
of the maximal ideal, since
$\mathfrak qA_{\mathfrak q}\cap A=\mathfrak q$.
Thus primes and overrings
correspond bijectively, reversing
inclusion: larger primes invert
fewer original elements.
Conversely an inclusion of overrings
preserves their units, so contraction
of their maximal ideals reverses
that inclusion as well.

For the given $W$, this
prime is explicitly
$$
\mathfrak q=\{a\in A:(a)\cap W=\emptyset\}.
$$
Indeed $a/1$ is a unit
in $W^{-1}A$ precisely when
$ab=w$ for some $b\in A$
and $w\in W$. An inverse
$b/w$ gives this exact
equality by cancelling the
nonzero localization witness
in the domain, and the
equality conversely gives
the inverse. This proves
the formula from the maximal
ideal of the valuation ring.
The prime is disjoint from
$W$ and contains every prime
disjoint from $W$: if
$a$ is in such a prime,
no multiple $ab$ can be
in $W$. It is therefore
the largest such prime,
and the canonical fraction
map identifies $W^{-1}A$
with exactly $A_{\mathfrak q}$.
\end{proof}

\begin{lemma}
\label{editorial-source-lemma-stack-valuation-rings}
Let $A'$ be a valuation ring with residue field $K$.
Let $A$ be a valuation ring with fraction field $K$.
Then
$C = \{\lambda \in A' \mid \lambda \bmod \mathfrak m_{A'} \in A\}$
is a valuation ring.
\end{lemma}

\begin{proof}
Keep the original residue map
$\pi:A'\to K$ and write
$\mathfrak p=\mathfrak m_{A'}$.
The original $C=\pi^{-1}(A)$
is a subring of $A'$,
and $\mathfrak p\subseteq C$.
The map $C/\mathfrak p\to A$
sending $[c]$ to $\pi(c)$
is an isomorphism: surjectivity
follows by lifting each $a\in A$
through $\pi$, and its kernel
is exactly $\mathfrak p$.
The inverse sends $a$ to
the class of any such lift;
different lifts differ in
$\mathfrak p$. In particular
$\mathfrak p$ is a prime
of the original $C$.

There is a more precise
comparison than just equality
of fraction fields:
$$
C_{\mathfrak p}\longrightarrow A',\qquad c/t\longmapsto c t^{-1}
$$
is an isomorphism inside
$F=\operatorname{Frac}(A')$.
It is defined and injective
because $t\in C\setminus\mathfrak p$
is a nonzero unit of $A'$.
For surjectivity, take the
original $b\in A'$ and
write its residue in
$K=\operatorname{Frac}(A)$
as $\pi(b)=a/d$ with
$a,d\in A$ and $d\neq0$;
for zero residue take $a=0,d=1$.
Choose $t\in A'$ lifting
this original $d$. Then
$t\in C\setminus\mathfrak p$
and $tb\in C$ because
$\pi(tb)=a$. Thus the
inverse sends $b$ to
$(tb)/t$. Independence of
the lift and denominator
follows from the equality
of cross products in the
original field $F$. Both
composites are identities.
It follows that
$\operatorname{Frac}(C)=F$,
including the case $A'$ is
a field and $\mathfrak p=0$.

For nonzero $x\in F$, if
$x\notin A'$, Lemma
\ref{editorial-source-lemma-valuation-ring-x-or-x-inverse}
gives $x^{-1}\in\mathfrak m_{A'}\subseteq C$.
If $x\in\mathfrak m_{A'}$,
then $x\in C$. In the
remaining case $x$ is a
unit of $A'$, so its
nonzero residue satisfies
$\pi(x)\in A$ or
$\pi(x)^{-1}\in A$ by the
valuation alternative in the
original residue field $K$.
Because $\pi(x^{-1})=\pi(x)^{-1}$,
this says $x\in C$ or
$x^{-1}\in C$. The criterion
of Lemma
\ref{editorial-source-lemma-x-or-x-inverse-valuation-ring}
therefore proves the assertion.

The maximal ideal and residue
field are also exact:
$$
\mathfrak m_C=\pi^{-1}(\mathfrak m_A),\qquad
C/\mathfrak m_C\simeq A/\mathfrak m_A.
$$
Indeed an element $c$ outside
that inverse image has
$\pi(c)\in A^*$, so it
is a unit of $A'$ and
$\pi(c^{-1})=\pi(c)^{-1}\in A$,
whence $c^{-1}\in C$.
An element inside cannot be
a unit, by applying $\pi$
to an inverse equation.
The residue isomorphism is
the map $[c]\mapsto[\pi(c)]$,
with inverse given by a
lift as above.

Finally the original prime
chain of $C$ is described
by these two exact maps.
Primes contained in $\mathfrak p$
correspond to primes of
$C_{\mathfrak p}=A'$ by
extension and contraction.
Primes containing $\mathfrak p$
correspond to primes of
$C/\mathfrak p=A$ by
quotient and inverse image.
All primes of the valuation
ring $C$ are comparable,
as proved in Lemma
\ref{editorial-source-lemma-x-or-x-inverse-valuation-ring},
so these two portions cover
its spectrum. They meet
at the single original prime
$\mathfrak p$, corresponding
to $\mathfrak m_{A'}$ on
the localization side and
to $(0)$ on the quotient
side. This proves the full
prime-chain composition.
\end{proof}

\begin{lemma}
\label{editorial-source-lemma-find-valuation-rings}
Let $A$ be a normal domain with fraction field $K$.
\begin{enumerate}
\item For every $x \in K$, $x \not \in A$ there exists a valuation ring
$A \subset V \subset K$ with fraction field $K$ such that $x \not \in V$.
\item If $A$ is local, we can moreover choose $V$ which dominates $A$.
\end{enumerate}
In other words, $A$ is the intersection of all valuation rings in $K$
containing $A$ and if $A$ is local, then $A$ is the intersection of
all valuation rings in $K$ dominating $A$.
\end{lemma}

\begin{proof}
For the given normal domain,
$x\in K\setminus A$ is
not integral over $A$.
We prove the construction
using exactly this nonintegrality
condition; it will also give
the integral-closure statement
for arbitrary domains below.
Since $0\in A$, the original
$x$ is nonzero. Put
$B=A[x^{-1}]\subseteq K$.
If $x\in B$, a finite
expression
$x=\sum_{i=0}^d a_i x^{-i}$
would give the full monic
equation
$$
x^{d+1}-\sum_{i=0}^d a_i x^{d-i}=0,
$$
contradicting nonintegrality.
Thus $x^{-1}$ is not a
unit of $B$, and
$x^{-1}B$ is a proper
ideal. Choose a prime
$\mathfrak p$ containing it
and apply Lemma
\ref{editorial-source-lemma-dominate} to the
original local subring
$B_{\mathfrak p}\subseteq K$.
It gives a valuation ring
$V$ of fraction field $K$
dominating $B_{\mathfrak p}$.
The element $x^{-1}$ belongs
to $\mathfrak m_V$, so
$x\notin V$: otherwise
their product would put
$1$ in $\mathfrak m_V$.
Since $A\subseteq B\subseteq V$,
this proves (1).

Suppose now the original
$A$ is local. Retain
$$
J=x^{-1}B+\mathfrak m_A B.
$$
If $J=B$, finite expressions
in the original $x^{-1}$,
with zero padding to a
common degree $d\geq0$,
give
$$
1=\left(\sum_{i=0}^d a_i x^{-i}\right)x^{-1}
+\sum_{i=0}^d a'_i x^{-i},
\qquad a_i\in A,\quad a'_i\in\mathfrak m_A.
$$
Multiplying by the full
$x^{d+1}$ and keeping the
two original coefficient lists
gives
$$
(1-a'_0)x^{d+1}
-\sum_{i=0}^d a_i x^{d-i}
-\sum_{i=1}^d a'_i x^{d+1-i}=0.
$$
For $d=0$ the last sum
is empty and the equation
is $(1-a'_0)x-a_0=0$.
In every case $u=1-a'_0$
is a unit of $A$.
The monic polynomial
$$
Q(T)=T^{d+1}
-\sum_{i=0}^d u^{-1}a_i T^{d-i}
-\sum_{i=1}^d u^{-1}a'_i T^{d+1-i}
$$
annihilates $x$, and the
complete polynomial originally
obtained is exactly $uQ(T)$.
Thus $x$ would be integral,
a contradiction. It follows
that the original $J$ is
proper. Choose a prime
$\mathfrak p\supseteq J$ and
then a valuation ring $V$
of fraction field $K$
dominating $B_{\mathfrak p}$,
as in the first construction.

The contraction $\mathfrak p\cap A$
contains $\mathfrak m_A$ and
is proper, so it equals
$\mathfrak m_A$. Localization
gives
$\mathfrak m_{B_{\mathfrak p}}\cap A=\mathfrak m_A$,
and domination gives
$\mathfrak m_V\cap A=\mathfrak m_A$.
Hence this original $V$
dominates $A$ and still
excludes $x$ because
$x^{-1}\in\mathfrak m_V$.
This proves (2). Each
asserted intersection equals
$A$: inclusion of $A$
is built into all its
members, and the construction
excludes every original element
of $K\setminus A$ from
at least one member.

More generally let $A$
be any domain with original
fraction field $K$, and
let $\widetilde A$ be its
integral closure in $K$.
Every $x\in\widetilde A$
belongs to every valuation
ring $A\subseteq V\subseteq K$:
its original monic equation
over $A$ is the same
monic equation over $V$,
and $V$ is normal by
Lemma \ref{editorial-source-lemma-valuation-ring-normal}.
If $x\notin\widetilde A$,
the construction above applies
because $x$ is not integral
over $A$, and gives a
member excluding it. Therefore
$$
\widetilde A=\bigcap_{A\subseteq V\subseteq K\text{ valuation}}V.
$$
For local $A$, exactly the
second construction gives
the same equality with the
intersection restricted to
valuation rings dominating $A$.

For a prescribed prime
$\mathfrak p\subset A$, the
valuation rings containing
$A$ with center
$\mathfrak m_V\cap A=\mathfrak p$
are precisely those dominating
the original $A_{\mathfrak p}$.
Indeed every $a\in A\setminus\mathfrak p$
is a unit of $V$, giving
the actual inclusion of
fractions $A_{\mathfrak p}\to V$,
and its maximal ideal
contracts to $\mathfrak pA_{\mathfrak p}$.
The converse follows by
contracting once more to
$A$. Their intersection
is thus exactly the
integral closure of
$A_{\mathfrak p}$ in $K$,
by the proved local formula.
\end{proof}

\noindent
A {\it totally ordered abelian group} is a pair $(\Gamma, \geq)$
consisting of an abelian group $\Gamma$ endowed with a total
ordering $\geq$ such that $\gamma \geq \gamma' \Rightarrow
\gamma + \gamma'' \geq \gamma' + \gamma''$ for all
$\gamma, \gamma', \gamma'' \in \Gamma$.

\begin{lemma}
\label{editorial-source-lemma-valuation-group}
Let $A$ be a valuation ring with field of fractions $K$.
Set $\Gamma = K^*/A^*$ (with group law written additively).
For $\gamma, \gamma' \in \Gamma$
define $\gamma \geq \gamma'$ if and only if
$\gamma - \gamma'$ is in the image of $A - \{0\} \to \Gamma$.
Then $(\Gamma, \geq)$ is a totally ordered abelian group.
\end{lemma}

\begin{proof}
Retain the original quotient
$\Gamma=K^*/A^*$ and write
$v(x)=[x]$ additively, so
$v(xy)=v(x)+v(y)$ by the
quotient group law. For original
$x,y\in K^*$, the proposed
order is exactly
$$
v(x)\geq v(y)\quad\Longleftrightarrow\quad x/y\in A.
$$
Indeed $v(x/y)$ is in the
image of $A\setminus\{0\}$
precisely when $x/y$ differs
from such an element by
a unit of $A$. Multiplication
by that unit and its inverse
proves the equivalence.
If the representatives are
changed to $xu,yw$ with
$u,w\in A^*$, the original
ratio becomes $(x/y)(u/w)$.
Its membership in $A$ is
unchanged, so the order is
well-defined on the quotient.

Reflexivity follows from $x/x=1$.
If $x/y,y/x\in A$, their
product is $1$, so $x/y$
is a unit and $v(x)=v(y)$;
this is antisymmetry. If
$x/y$ and $y/z$ belong
to $A$, their exact product
$x/z$ does too, proving
transitivity. For every pair,
the valuation alternative for
$x/y$ gives $x/y\in A$
or $y/x\in A$, proving
totality. Finally multiplying
both original representatives
by the same $z\in K^*$
does not change their ratio:
$(xz)/(yz)=x/y$. This
proves translation invariance
of the order and the
claimed ordered group structure.

Every totally ordered abelian
group is torsion-free. If
$\gamma>0$, repeated strict
translation invariance gives
$n\gamma>0$ for every integer
$n\geq1$; if $\gamma<0$,
apply the same argument to
$-\gamma$. Hence $n\gamma=0$
with $n\geq1$ forces
$\gamma=0$. In particular
this holds for the original
$\Gamma$.

The order recovers the
original ring and maximal
ideal exactly:
$$
A=\{0\}\cup\{x\in K^*:v(x)\geq0\},\qquad
\mathfrak m_A=\{0\}\cup\{x\in K^*:v(x)>0\}.
$$
The first equality is the
ratio criterion with $y=1$.
For $x\in A\setminus\{0\}$,
the quotient definition says
$v(x)=0$ exactly when
$x\in A^*$, giving the
second equality by localness.
Also $v(a/b)=v(a)-v(b)$
for all original nonzero
$a,b\in K$. The value
of $0$ is not introduced
into this group. If $A=K$,
the quotient is the zero
group, with its unique order.
Conversely, if this quotient
is zero, every nonzero
element of $K$ is a unit
of $A$, so $A=K$.
\end{proof}

\begin{definition}
\label{editorial-source-definition-value-group}
Let $A$ be a valuation ring.
\begin{enumerate}
\item The totally ordered abelian group $(\Gamma, \geq)$ of
Lemma \ref{editorial-source-lemma-valuation-group} is called the
{\it value group} of the valuation ring $A$.
\item The map $v : A - \{0\} \to \Gamma$ and also $v : K^* \to \Gamma$ is
called the {\it valuation} associated to $A$.
\item The valuation ring $A$ is called a {\it discrete valuation ring}
if $\Gamma \cong \mathbf{Z}$.
\end{enumerate}
\end{definition}

\noindent
Suppose the original ordered group
$\Gamma$ is isomorphic as a group
to $\mathbf Z$, and choose an
initial group isomorphism
$\psi:\Gamma\to\mathbf Z$.
Put $\gamma=\psi^{-1}(1)\neq0$.
Let $\varepsilon=1$ if
$\gamma>0$, and $\varepsilon=-1$
if $\gamma<0$. Then
$$
\phi=\varepsilon\psi:\Gamma\longrightarrow\mathbf Z
$$
is an order isomorphism with
the usual order on the
receiving integers. Indeed
$\varepsilon\gamma>0$, and
every original element is
$n\gamma$ for a unique
$n\in\mathbf Z$. Repeated
addition and translation
invariance give
$n\gamma\geq0$ precisely when
$\varepsilon n\geq0$, which
is precisely $\phi(n\gamma)\geq0$.
The same applies to differences,
proving both order implications.
Any other group isomorphism
is $\delta\psi$ for
$\delta\in\{1,-1\}$, since
the only generators of
$\mathbf Z$ are $1,-1$.
Requiring the inverse image
of $1$ to be positive
forces $\delta=\varepsilon$.
This is the precise unique
order-compatible isomorphism.
The original $\Gamma$ and
its value map $v$ are
retained; their integer comparison
is $\phi\circ v$, with
$v=\phi^{-1}\circ(\phi\circ v)$.

Write $\gamma_0=\phi^{-1}(1)>0$.
Surjectivity of the original
quotient $v:K^*\to\Gamma$
gives an element $\pi\in K^*$
with $v(\pi)=\gamma_0$.
The positive-cone description
puts $\pi$ in the maximal
ideal of $A$. For every
original nonzero $a\in A$,
put $n=\phi(v(a))\geq0$.
The exact original element
$u=a/\pi^n$ has value
$v(a)-n\gamma_0=0$, so
$u\in A^*$ and
$$
a=u\pi^n.
$$
Both the unit $u$ and
the original exponent are
part of this equality;
neither is discarded. For
the chosen $\pi$, the
integer $n$ is unique by
applying $\phi v$, and
then $u=a/\pi^n$ is
unique as well. Another
element $\pi'$ of the
same least positive value
has $\pi'/\pi\in A^*$,
which is the exact comparison
between such choices.

Every nonzero ideal $I$
has the form $(\pi^n)$
for exactly one $n\geq0$.
Choose the least integer
among $\phi(v(a))$ for
the original nonzero $a\in I$,
and choose such an $a$.
The full equality $a=u\pi^n$
with $u$ a unit puts
$\pi^n=u^{-1}a$ in $I$.
For every nonzero $b\in I$,
the inequality $v(b)\geq n\gamma_0$
puts $b/\pi^n$ in $A$,
and a zero member contributes
the equality $0=\pi^n0$.
Thus all of $I$ is
contained in $(\pi^n)$,
giving equality. Distinct
$n$ give distinct ideals
by their least original
values. The zero ideal
is retained as a separate
case and is not assigned
an infinite value in $\Gamma$.
In particular the maximal
ideal is $(\pi)$, every
ideal is principal, and
the only prime ideals are
$(0)$ and $(\pi)$: for
$n>1$, the product
$\pi\pi^{n-1}$ lies in
$(\pi^n)$ although neither
factor does. This proves
the ideal and prime descriptions
from the original ordered
quotient without suppressing
any unit factors.

\begin{lemma}
\label{editorial-source-lemma-properties-valuation}
Let $A$ be a valuation ring. The valuation $v : A -\{0\} \to \Gamma_{\geq 0}$
has the following properties:
\begin{enumerate}
\item $v(a) = 0 \Leftrightarrow a \in A^*$,
\item $v(ab) = v(a) + v(b)$,
\item $v(a + b) \geq \min(v(a), v(b))$ provided $a + b \not = 0$.
\end{enumerate}
\end{lemma}

\begin{proof}
The quotient map to the
original $\Gamma=K^*/A^*$
has kernel $A^*$, giving
(1), and its group law
gives (2). For (3), retain
the original nonzero $a,b$
and assume $a+b\neq0$.
Interchanging their roles
if necessary, suppose
$v(a)\leq v(b)$. The
order proved in Lemma
\ref{editorial-source-lemma-valuation-group}
then gives $b/a\in A$.
The exact equality
$$
\frac{a+b}{a}=1+\frac ba
$$
puts the nonzero left-hand
element in $A$, so
$v(a+b)-v(a)\geq0$.
This proves (3), retaining
both summands and their
original ratio.

If $v(a)<v(b)$, that
ratio $b/a$ belongs to
$\mathfrak m_A$. Thus
$1+b/a$ is a unit:
if it were in the maximal
ideal, subtracting $b/a$
would put $1$ there.
Consequently $a+b\neq0$
and $v(a+b)=v(a)$.
The reversed strict inequality
gives the symmetric equality.
This entire proof also
applies to nonzero $a,b\in K$
outside $A$: the comparison
still puts the same ratio
in $A$. The original value
map on $K^*$ therefore
satisfies the same inequality,
with equality whenever the
two original values differ.

For a finite original list
$a_1,\ldots,a_n\in K$, keep
all its terms, including
zeros. If at least one
term is nonzero, choose
$a_j$ of least value among
the nonzero terms. Every
ratio $a_i/a_j$ belongs
to $A$, with a zero
term giving the ratio $0$.
Hence if the original sum
$a=\sum_i a_i$ is nonzero,
then $a/a_j=\sum_i a_i/a_j$
is a nonzero element of
$A$, proving $v(a)\geq v(a_j)$.
If exactly one term has
the least value, every
other ratio is in
$\mathfrak m_A$, including
the zeros. The full sum
of ratios is $1$ plus
an element of that ideal,
hence a unit. In that
case the sum is nonzero
and has exactly the least
original value. For an
all-zero or empty list,
the sum is $0$ and no
value in $\Gamma$ is
assigned to it.
\end{proof}

\begin{lemma}
\label{editorial-source-lemma-characterize-valuation-ring}
Let $A$ be a ring. The following are equivalent
\begin{enumerate}
\item $A$ is a valuation ring,
\item $A$ is a local domain and every finitely generated
ideal of $A$ is principal.
\end{enumerate}
\end{lemma}

\begin{proof}
Suppose first that $A$ is
a valuation ring. Retain
the complete finite generator
list $f_1,\ldots,f_n$.
If the list is empty
or all generators are zero,
its ideal is $(0)$ and
is principal. Otherwise choose
an original nonzero $f_i$
whose value is least among
the nonzero generators, and
put $h=f_i$. For each
nonzero $f_j$, the comparison
$v(f_j)\geq v(h)$ gives
the original ratio $c_j=f_j/h\in A$.
For $f_j=0$, set $c_j=0$.
Then the exact equalities
$f_j=hc_j$ for every $j$
give $(f_1,\ldots,f_n)\subseteq(h)$,
while $h$ itself is in
the original list, giving
the reverse inclusion.
No value is taken at
a zero generator.

Conversely, assume the original
$A$ is a local domain
with every finitely generated
ideal principal. An original
fraction $f/g$ with $g\neq0$
belongs to $A$ if $f=0$.
For $f,g\neq0$, write
$(f,g)=(h)$. Then $h\neq0$,
and there are original
$a,b,c,d\in A$ with
$$
f=ah,\qquad g=bh,\qquad h=cf+dg.
$$
Substitution gives
$h=(ca+db)h$. The domain
property and $h\neq0$ permit
cancellation, giving $ca+db=1$.
Both $a,b$ cannot belong
to the maximal ideal,
so at least one is a
unit. If $a$ is a
unit, the original fraction
$g/f=b/a$ belongs to $A$;
if $b$ is a unit,
$f/g=a/b$ belongs to $A$.
This proves the nonzero
fraction alternative in the
actual $\operatorname{Frac}(A)$.
Lemma
\ref{editorial-source-lemma-x-or-x-inverse-valuation-ring}
therefore makes $A$ a
valuation ring.

The converse proof uses
only two-generated ideals.
Thus in (2) it suffices
to require that every
two-generated ideal be principal;
the forward proof then
recovers the assertion for
every original finite list.
The local hypothesis cannot
be discarded: every ideal
of $\mathbf Z$ is principal,
since for a nonzero ideal
its least positive element
divides all its elements
by division with remainder,
and the zero ideal is
$(0)$. But neither $2/3$
nor $3/2$ is an integer,
so the exact fraction
criterion shows that $\mathbf Z$
is not a valuation ring.
\end{proof}

\begin{lemma}
\label{editorial-source-lemma-valuation-valuation-ring}
Let $(\Gamma, \geq)$ be a totally ordered abelian group.
Let $K$ be a field. Let $v : K^* \to \Gamma$ be a homomorphism
of abelian groups such that $v(a + b) \geq \min(v(a), v(b))$ for
$a, b \in K$ with $a, b, a + b$ not zero. Then
$$
A =
\{
x \in K \mid x = 0 \text{ or } v(x) \geq 0
\}
$$
is a valuation ring with value group $\Im(v) \subset \Gamma$,
with maximal ideal
$$
\mathfrak m =
\{
x \in K \mid x = 0 \text{ or } v(x) > 0
\}
$$
and with group of units
$$
A^* =
\{
x \in K^* \mid v(x) = 0
\}.
$$
\end{lemma}

\begin{proof}
Keep the original ordered group
$\Gamma$ and homomorphism
$v:K^*\to\Gamma$, without
assuming it is surjective.
The homomorphism identities give
$v(1)=0$ and $v(x^{-1})=-v(x)$.
Also $2v(-1)=v(1)=0$.
A totally ordered abelian group
has no nonzero element killed
by a positive integer, by
the strict-order argument in
Lemma \ref{editorial-source-lemma-valuation-group}.
Hence $v(-1)=0$, and
$v(-x)=v(x)$ for nonzero $x$.

The original set $A$ contains
$0,1$ and is closed under
additive inverses. It is
closed under multiplication:
products involving zero are
zero, and otherwise the
sum of two nonnegative
original values is nonnegative.
For addition, a zero summand
causes no problem. If both
summands and their sum are
nonzero, the given inequality
puts the sum in $A$;
if their sum is zero,
it belongs to $A$ by
definition. Thus $A$ is
a unital subring of $K$.
For every nonzero $x\in K$,
totality gives $v(x)\geq0$
or $v(x)<0$. In the
second case $v(x^{-1})>0$,
so $x^{-1}\in A$.
The criterion of Lemma
\ref{editorial-source-lemma-x-or-x-inverse-valuation-ring}
therefore proves that $A$
is a valuation ring whose
fraction field is exactly $K$.

A nonzero original $x$
is a unit of $A$ exactly
when both $v(x)$ and
$-v(x)$ are nonnegative,
which by antisymmetry is
equivalent to $v(x)=0$.
Thus its nonunits are
precisely $0$ and the
elements of positive value,
giving the asserted maximal
ideal. Directly, that set
is a proper ideal: addition
uses the minimum inequality
unless the sum is zero,
and multiplying a positive
value by an element of
$A$ adds a nonnegative
value. It excludes $1$,
and every element outside
it has the inverse just
described. This also proves
maximality explicitly.

Let $q:K^*\to\Delta=K^*/A^*$
be the original associated
value-group quotient. The map
$$
\overline v:\Delta\longrightarrow\operatorname{Im}(v),
\qquad[x]\longmapsto v(x)
$$
is well-defined because
$A^*=\ker(v)$, injective
by the same kernel equality,
and surjective by the
definition of the image.
Its inverse sends a value
$\gamma=v(x)$ to $[x]$;
different original choices
$x,y$ have $v(x/y)=0$
and differ by a unit of
$A$, proving independence.
Moreover
$$
[x]\geq[y]\quad\Longleftrightarrow\quad
x/y\in A\quad\Longleftrightarrow\quad
v(x)-v(y)\geq0.
$$
Thus it is an order
isomorphism with exactly
the asserted subgroup of
the original $\Gamma$.
If $i:\operatorname{Im}(v)\hookrightarrow\Gamma$
is its original inclusion,
then $i\overline vq=v$;
the ambient codomain and
the original value map
have both been retained.

Finally two such original
maps $v:K^*\to\Gamma$ and
$w:K^*\to\Lambda$ define
the same valuation ring
exactly when there is an
order isomorphism
$\phi:\operatorname{Im}(v)\to\operatorname{Im}(w)$
with $\phi(v(x))=w(x)$
for every original $x$.
For a common ring, equality
of two $v$-values means
their ratio is a unit
of that ring, and hence
also has $w$-value zero.
This proves the formula
well-defined; the reverse
argument gives its inverse.
The same ratio criterion
proves that it preserves
and reflects order. Conversely
such an isomorphism preserves
the nonnegative cone, so
the two original sets
$\{0\}\cup v^{-1}(\Gamma_{\geq0})$
and $\{0\}\cup w^{-1}(\Lambda_{\geq0})$
are equal. Uniqueness follows
because every element of
the image is an original
value $v(x)$. No isomorphism
between unused parts of the
ambient groups is asserted.
\end{proof}

\noindent
Let $(\Gamma, \geq)$ be a totally ordered abelian group.
An {\it ideal of $\Gamma$} is a subset $I \subset \Gamma$ such
that all elements of $I$ are $\geq 0$ and $\gamma \in I$,
$\gamma' \geq \gamma$ implies $\gamma' \in I$. We say that such
an ideal is {\it prime} if $0 \not \in I$ and if
$\gamma + \gamma' \in I, \gamma, \gamma' \geq 0
\Rightarrow \gamma \in I \text{ or } \gamma' \in I$.

\begin{lemma}
\label{editorial-source-lemma-ideals-valuation-ring}
Let $A$ be a valuation ring.
Ideals in $A$ correspond $1 - 1$ with ideals of $\Gamma$.
This bijection is inclusion preserving, and maps prime
ideals to prime ideals.
\end{lemma}

\begin{proof}
Retain $K=\operatorname{Frac}(A)$, the original ordered quotient
$\Gamma=K^*/A^*$, and its quotient map $v:K^*\to\Gamma$.
For an ideal $J\subset A$ and an ideal $U\subset\Gamma$, define
$$
V(J)=\{v(x):x\in J,\ x\neq0\},\qquad
J(U)=\{0\}\cup\{x\in K^*:v(x)\in U\}.
$$
Every element of $V(J)$ is nonnegative. If $v(x)=\gamma\in V(J)$
and $\delta\geq\gamma$, choose $y\in K^*$ with $v(y)=\delta$.
Then $v(y/x)=\delta-\gamma\geq0$, so $y/x\in A$ and
$y=(y/x)x\in J$. Thus $V(J)$ is upward closed.
The same calculation shows that every representative of a value in
$V(J)$ lies in $J$; when the values are equal, $y/x$ is an actual unit.

Since $U\subset\Gamma_{\geq0}$, the set $J(U)$ lies in $A$.
It contains zero. If $x\neq0$ belongs to it, then
$v(-x)=v(x)\in U$. For nonzero $x,y\in J(U)$ with $x+y\neq0$,
the valuation inequality gives
$v(x+y)\geq\min\{v(x),v(y)\}\in U$, hence $x+y\in J(U)$.
If a summand or the sum is zero the closure follows directly from
the definition. For $a\in A$ and $x\in J(U)$, a zero factor gives
$ax=0$; otherwise $v(ax)=v(a)+v(x)\geq v(x)$, so $ax\in J(U)$.
This proves that $J(U)$ is an ideal.
Surjectivity of $v$ and the representative calculation give
$V(J(U))=U$ and $J(V(J))=J$. Both maps preserve inclusion.
In particular $J(\varnothing)=(0)$ and
$J(\Gamma_{\geq0})=A$; no value is assigned to the zero element.

The ideal $J$ is proper exactly when $0\notin V(J)$, since an
element of value zero is a unit. If $J$ is prime and
$\gamma,\delta\geq0$ with $\gamma+\delta\in V(J)$, choose
$x,y\in A\setminus\{0\}$ with these values. The representative
calculation puts $xy$ in $J$, so $x\in J$ or $y\in J$.
Consequently $\gamma\in V(J)$ or $\delta\in V(J)$.
Conversely, suppose $U$ is prime and $xy\in J(U)$ for $x,y\in A$.
If $xy=0$, the domain property gives a zero factor, which belongs
to $J(U)$. Otherwise the defining condition on
$v(x)+v(y)\in U$ puts one factor in $J(U)$.
Thus the correspondence preserves and reflects primality, including
the empty group ideal and the zero ring ideal.

The correspondence also describes the ideal operations exactly.
For every family $(J_\lambda)_{\lambda\in\Lambda}$, put
$U_\lambda=V(J_\lambda)$. Then
$$
V\left(\bigcap_{\lambda\in\Lambda}J_\lambda\right)
=\bigcap_{\lambda\in\Lambda}U_\lambda,
\qquad
V\left(\sum_{\lambda\in\Lambda}J_\lambda\right)
=\bigcup_{\lambda\in\Lambda}U_\lambda.
$$
The first equality follows by membership of each nonzero element.
For the second, a nonzero member of the sum is a finite sum of
members of the original ideals. Discarding no term from that
equality, list its nonzero terms. Their list is nonempty; its least
value belongs to one $U_\lambda$. The finite-sum inequality and
upward closure put the value of the sum in that same $U_\lambda$.
The reverse inclusion follows from each $J_\lambda$ being contained
in the sum. For an empty family the intersection is $A$ and its
value set is $\Gamma_{\geq0}$, while the sum is $(0)$ and its
value set is empty; the intersections on the right are taken
inside $\Gamma_{\geq0}$.
For two ideals with value sets $U,W$, one has
$$
V(J(U)J(W))=U+W
=\{\gamma+\delta:\gamma\in U,\ \delta\in W\},
$$
and
$$
V(\sqrt{J(U)})
=\{\gamma\in\Gamma_{\geq0}:n\gamma\in U
\text{ for some integer }n\geq1\}.
$$
Indeed $U+W$ is upward closed: if
$\eta\geq\gamma+\delta$, write
$\eta=\gamma+(\delta+\eta-\gamma-\delta)$ and use
$\delta+\eta-\gamma-\delta\geq\delta$.
Every value in $U+W$ is the value of a product of representatives
from the original ideals. A nonzero element of their product ideal
is a finite sum of such products, and the least-value argument
just used gives the reverse inclusion. If either ideal is zero,
both sides are empty. Finally, for nonzero $x\in A$, the equality
$v(x^n)=n v(x)$ proves the radical formula in both directions;
the original zero element lies in the radical separately.
\end{proof}

\begin{lemma}
\label{editorial-source-lemma-valuation-ring-Noetherian-discrete}
A valuation ring is Noetherian if and only if it is
a discrete valuation ring or a field.
\end{lemma}

\begin{proof}
Let $K=\operatorname{Frac}(A)$ and retain the original value map
$v:K^*\to\Gamma=K^*/A^*$. If $A$ is a field, its two ideals
are $(0)$ and $A$, so it is Noetherian.
If $A$ is a discrete valuation ring, let $\gamma_0>0$ generate
its original ordered group, and let
$\phi:\Gamma\to\mathbf Z$ be the order isomorphism
$\phi(n\gamma_0)=n$. Choose $\pi\in A$ with
$v(\pi)=\gamma_0$. The nonzero ideals, and only the nonzero
ideals, are
$$
I_n=\{0\}\cup\{x\in K^*:v(x)\geq n\gamma_0\}
=(\pi^n),\qquad n\geq0.
$$
To verify this directly, the values of a nonzero ideal under
$\phi$ form a nonempty upward closed subset of
$\mathbf Z_{\geq0}$, with a least element $n$.
For an original element $a$ of that value,
$u=a/\pi^n$ has value zero, so $u\in A^*$ and
$a=u\pi^n$. Every original nonzero $b$ of the ideal satisfies
$b/\pi^n\in A$, while $0=\pi^n0$. These equalities prove
that the ideal is $(\pi^n)$ and retain its original unit factor.
The remaining ideal $(0)$ is generated by zero. Thus every ideal
is finitely generated and $A$ is Noetherian. The map $v$ still
takes values in $\Gamma$; the integer-valued comparison is
$\phi\circ v$, with $v=\phi^{-1}\circ(\phi\circ v)$.

Conversely, suppose $A$ is Noetherian and is not a field.
Ascending chains of its ideals stabilize: their union is an ideal,
and finitely many generators of that union all occur in a single
member of the chain. By Lemma \ref{editorial-source-lemma-ideals-valuation-ring},
the same holds for the upward closed subsets of
$\Gamma_{\geq0}$. There is a positive element of $\Gamma$.
If there were no least positive element, starting from one such
element one could choose
$\gamma_1>\gamma_2>\gamma_3>\cdots>0$.
The corresponding ideals
$\gamma_i+\Gamma_{\geq0}$ would form a strictly ascending
chain: $\gamma_{i+1}$ belongs to its own set but not to the
preceding one. This is impossible. Denote the least positive
element by $\gamma_0$.

For any original $\gamma>0$, suppose that
$\gamma-n\gamma_0\geq0$ for every integer $n\geq0$.
These elements strictly decrease, and their upward closed sets
again give a strictly ascending chain of ideals, a contradiction.
Consequently there is a least integer $N\geq1$ such that
$\gamma-N\gamma_0<0$. Put
$r=\gamma-(N-1)\gamma_0$. Then
$0\leq r<\gamma_0$, so the least-positive property forces
$r=0$. Thus $\gamma=(N-1)\gamma_0$.
Zero is $0\gamma_0$, and applying the argument to $-\gamma$
expresses every negative element as a negative multiple of
$\gamma_0$. Distinct integer multiples are distinct, because
a positive multiple of $\gamma_0$ is positive. Therefore
$n\mapsto n\gamma_0$ is an order isomorphism
$\mathbf Z\to\Gamma$, and $A$ is a discrete valuation ring.
The argument uses the ascending chain condition twice; merely
having a least positive group element is not the assertion used
to conclude that all positive elements are its multiples.
\end{proof}


























\subsection{Editorial treatment of Algebra lines 12153--12184}

\hypertarget{algebra-context-158}{}

\noindent\href{https://github.com/stacks/stacks-project/blob/a04446e57ec1fbc252a871afcec7752fb2807b14/algebra.tex\#L12153}{Original source, lines 12153--12184}. Complete original and reviewed TeX passages accompany this note. Every intervention is separately recorded; the following treatment is editorial.


\begin{lemma}
\label{editorial-source-lemma-Noetherian-basic}
Let $R$ be a Noetherian ring.
Any finite $R$-module is of finite presentation.
Any submodule of a finite $R$-module is finite.
The ascending chain condition holds for $R$-submodules
of a finite $R$-module.
\end{lemma}

\begin{proof}
We first prove that every submodule $N$ of a finite $R$-module
$M$ is finite, by induction on a chosen number $n$ of generators
$x_1,\ldots,x_n$. If $n=0$, then $M=N=0$ and the empty list
generates $N$. If $n=1$, the surjection $R\to M$, $a\mapsto ax_1$,
has an ideal $I$ as kernel. The inverse image $J$ of $N$ is an
ideal containing $I$, and the induced map $J/I\to N$ is an
isomorphism. A finite generating list of $J$, which exists because
$R$ is Noetherian, maps to a generating list of $N$.

For $n>1$, let $M'$ be the submodule generated by
$x_1,\ldots,x_{n-1}$ and let $M''=M/M'$ with quotient map $q$.
The module $M''$ is generated by $q(x_n)$.
Set $N'=N\cap M'$ and $N''=q(N)$. The inclusion and restriction
of $q$ give the exact sequence
$$
0\longrightarrow N'\longrightarrow N\longrightarrow N''
\longrightarrow0.
$$
Both end modules are finite by induction. Choose generators
$u_1,\ldots,u_a$ of $N'$ and generators
$w_1,\ldots,w_b$ of $N''$, with lifts
$z_1,\ldots,z_b\in N$. If $z\in N$, write
$q(z)=\sum_{j=1}^b c_jw_j$. Then
$z-\sum_{j=1}^b c_jz_j\in N'$, so
$z=\sum_{i=1}^a d_i u_i+\sum_{j=1}^b c_jz_j$ for suitable
$d_i,c_j\in R$. This proves finite generation, including an
empty generating list at either end.
For a finite module $M$, choose a surjection $R^n\to M$.
Its kernel is a finite submodule of the finite module $R^n$.
Mapping a basis of $R^t$ to a finite generating list of this
kernel gives an exact presentation $R^t\to R^n\to M\to0$.

The last equivalence holds for every module $M$ over every ring
$R$, without a Noetherian or finite-generation hypothesis on
$R$ or $M$. If every submodule is finite and
$N_1\subset N_2\subset\cdots$ is an ascending chain,
its union $N$ is a submodule. Finitely many generators of $N$
all belong to one $N_j$, so $N=N_j=N_{j+1}=\cdots$.
Conversely, if a submodule $N$ is not finite, start with
$N_0=0$. Once $x_1,\ldots,x_j\in N$ have been chosen, their
span $N_j$ is not $N$, so choose $x_{j+1}\in N\setminus N_j$.
The strictly ascending chain $N_0\subsetneq N_1\subsetneq\cdots$
contradicts the ascending chain condition for submodules of $M$.
Thus that condition implies that every submodule is finite.
Applying the equivalence to the first part proves the final claim.
\end{proof}


\subsection{Editorial treatment of Algebra lines 12186--12218}

\hypertarget{algebra-context-159}{}

\noindent\href{https://github.com/stacks/stacks-project/blob/a04446e57ec1fbc252a871afcec7752fb2807b14/algebra.tex\#L12186}{Original source, lines 12186--12218}. Complete original and reviewed TeX passages accompany this note. Every intervention is separately recorded; the following treatment is editorial.


\begin{lemma}[Artin-Rees]
\label{editorial-source-lemma-Artin-Rees}
Suppose that $R$ is Noetherian. Let $I \subset R$ be an ideal.
Let $N \subset M$ be finite $R$-modules.
There exists a constant $c > 0$ such that
$I^n M \cap N  =  I^{n-c}(I^cM \cap N)$ for all $n \geq c$.
\end{lemma}

\begin{proof}
Keep the graded ring
$$
S=\bigoplus_{n\geq0}I^n,\qquad I^0=R,
$$
with its actual degree inclusions $\iota_n:I^n\to S$.
Thus $\iota_n(a)\iota_m(b)=\iota_{n+m}(ab)$.
Since $R$ is Noetherian, write $I=(f_1,\ldots,f_t)$.
The $R$-algebra map $R[X_1,\ldots,X_t]\to S$ sends $r\in R$
to $\iota_0(r)$ and sends a monomial with its coefficient to
$$
rX_1^{e_1}\cdots X_t^{e_t}
\longmapsto
\iota_{e_1+\cdots+e_t}
 (rf_1^{e_1}\cdots f_t^{e_t}).
$$
Every element of $I^n$ is a finite sum of these products of total
degree $n$, so the map is surjective. The polynomial ring, and
therefore $S$, is Noetherian by Lemma
\ref{editorial-source-lemma-Noetherian-permanence}. If $I=0$, the empty generator
list gives $S=R$ in degree zero and the same argument applies.

The graded $S$-module
$E=\bigoplus_{n\geq0}I^nM$ has action
$\iota_d(a)\jmath_n(x)=\jmath_{d+n}(ax)$, where $\jmath_n$
denotes its degree inclusion. If $x_1,\ldots,x_r$ generate
$M$, their degree-zero images generate $E$: each member of
$I^nM$ is a sum $\sum_{i=1}^r a_ix_i$ with $a_i\in I^n$.
The graded subgroup
$$
L=\bigoplus_{n\geq0}(I^nM\cap N)\subset E
$$
is an $S$-submodule, since
$I^d(I^nM\cap N)\subset I^{d+n}M\cap N$.
By Lemma \ref{editorial-source-lemma-Noetherian-basic}, $L$ is finite over $S$.
Take a finite generating list and replace it by the collection of
all its homogeneous components. Each component is in $L$, each
old generator is the sum of its components, and the generated
submodule is consequently still exactly $L$. Denote this finite
list by $\jmath_{d_j}(\xi_j)$ for $1\leq j\leq s$, where
$\xi_j\in I^{d_j}M\cap N$. The list may be empty when $L=0$.
Set
$$
c=\max\bigl(\{1\}\cup\{d_j:1\leq j\leq s\}\bigr)>0.
$$
For $n\geq c$, taking the degree-$n$ component of an expression
in these generators shows that every $x\in I^nM\cap N$ is
$$
x=\sum_{j=1}^s h_j\xi_j,\qquad h_j\in I^{n-d_j}.
$$
The degree inclusions above specify the graded expression from
which this equality in the original $M$ is obtained.
For each $j$, the ideal identity
$I^{n-d_j}=I^{n-c}I^{c-d_j}$ means that one can write
$h_j=\sum_{\ell=1}^{q_j}a_{j\ell}b_{j\ell}$ with
$a_{j\ell}\in I^{n-c}$ and $b_{j\ell}\in I^{c-d_j}$.
Each $b_{j\ell}\xi_j$ lies in $I^cM\cap N$, so the full
equality
$x=\sum_{j=1}^s\sum_{\ell=1}^{q_j}
a_{j\ell}(b_{j\ell}\xi_j)$ proves
$I^nM\cap N\subset I^{n-c}(I^cM\cap N)$.
The opposite inclusion holds because a product with a member of
$I^cM\cap N$ remains in $N$ and lies in $I^nM$.
If the generating list is empty then $L=0$, both modules in the
desired equality are zero, and the same choice $c=1$ is valid.

Every larger integer $d\geq c$ is also valid. Indeed, for $n\geq d$,
the equalities already proved at $n$ and at $d$ give
$$
I^nM\cap N=I^{n-c}(I^cM\cap N)
=I^{n-d}\bigl(I^{d-c}(I^cM\cap N)\bigr)
=I^{n-d}(I^dM\cap N).
$$
In particular, the two original filtrations on $N$ satisfy
$$
I^nN\subset I^nM\cap N\quad(n\geq0),\qquad
I^{k+c}M\cap N\subset I^kN\quad(k\geq0).
$$
The first containment follows from $N\subset M$; the second
uses the equality at $n=k+c$ and $I^cM\cap N\subset N$.
Thus the subspace $I$-adic topology from $M$ and the $I$-adic
topology on $N$ have cofinal neighbourhood systems, with the
displayed explicit index bound.
\end{proof}


\subsection{Editorial treatment of Algebra lines 12220--12238}

\hypertarget{algebra-context-160}{}

\noindent\href{https://github.com/stacks/stacks-project/blob/a04446e57ec1fbc252a871afcec7752fb2807b14/algebra.tex\#L12220}{Original source, lines 12220--12238}. Complete original and reviewed TeX passages accompany this note. Every intervention is separately recorded; the following treatment is editorial.


\begin{lemma}
\label{editorial-source-lemma-map-AR}
Suppose that $0 \to K \to M \xrightarrow{f} N$ is an
exact sequence of finitely generated modules
over a Noetherian ring $R$. Let $I \subset R$ be an ideal.
Then there exists a $c$ such that
$$
f^{-1}(I^nN) = K + I^{n-c}f^{-1}(I^cN)
\quad\text{and}\quad
f(M) \cap I^nN \subset f(I^{n - c}M)
$$
for all $n \geq c$.
\end{lemma}

\begin{proof}
Identify $K$ with its image in $M$, which is $\ker(f)$ by
exactness. The module $f(M)$ is finite over $R$.
Apply Lemma \ref{editorial-source-lemma-Artin-Rees} to $f(M)\subset N$, giving
an integer $c>0$ with
$$
f(M)\cap I^nN=I^{n-c}(f(M)\cap I^cN)\quad(n\geq c).
$$
Put $P=f^{-1}(I^cN)$. The restriction $f:P\to f(M)\cap I^cN$
is surjective: an element of the intersection already has an
inverse image in $M$, and every such inverse image lies in $P$.
For every integer $d\geq0$, linearity and lifting each member
of a finite sum give the exact equality
$f(I^dP)=I^d f(P)$.
Consequently, for $n\geq c$,
$$
f(M)\cap I^nN
=f\bigl(I^{n-c}P\bigr)
\subset f(I^{n-c}M).
$$

If $x\in f^{-1}(I^nN)$, the preceding equality expresses
$f(x)=\sum_{j=1}^t a_jf(p_j)$ for original coefficients
$a_j\in I^{n-c}$ and original elements $p_j\in P$.
Thus $x-\sum_{j=1}^t a_jp_j\in K$ and
$x\in K+I^{n-c}P$. Conversely, $f(K)=0$ and
$f(I^{n-c}P)\subset I^nN$, so every member of
$K+I^{n-c}P$ belongs to $f^{-1}(I^nN)$. This proves the
first equality as well as the second claimed containment.
The same maps give the more precise exact sequence
$$
0\longrightarrow K\cap I^{n-c}P
\longrightarrow I^{n-c}P
\xrightarrow{\ f\ }f(M)\cap I^nN
\longrightarrow0.
$$
Here the first arrow is the inclusion, its image is exactly the
kernel of the displayed restriction of the original $f$, and
surjectivity is the lifting calculation above. The kernel $K$
has not been assumed to vanish or omitted from the preimage formula.
\end{proof}


\subsection{Editorial treatment of Algebra lines 12240--12254}

\hypertarget{algebra-context-161}{}

\noindent\href{https://github.com/stacks/stacks-project/blob/a04446e57ec1fbc252a871afcec7752fb2807b14/algebra.tex\#L12240}{Original source, lines 12240--12254}. Complete original and reviewed TeX passages accompany this note. Every intervention is separately recorded; the following treatment is editorial.


\begin{lemma}[Krull's intersection theorem]
\label{editorial-source-lemma-intersect-powers-ideal-module-zero}
Let $R$ be a Noetherian local ring. Let $I \subset R$ be
a proper ideal. Let $M$ be a finite $R$-module.
Then $\bigcap_{n \geq 0} I^nM = 0$.
\end{lemma}

\begin{proof}
Let $N=\bigcap_{n\geq0}I^nM$. It is a submodule of the finite
module $M$, hence is finite by Lemma \ref{editorial-source-lemma-Noetherian-basic}.
For every $n\geq0$ one has $I^nM\cap N=N$.
Apply Lemma \ref{editorial-source-lemma-Artin-Rees} to $N\subset M$ and take
$n=c+1$, where $c>0$ is the integer it supplies. This gives
$$
N=I^{c+1}M\cap N=I(I^cM\cap N)=IN.
$$
The proper ideal $I$ lies in the maximal ideal $\mathfrak m$
of the local ring $R$. Hence
$N=IN\subset\mathfrak mN\subset N$, so
$N=\mathfrak mN$. Nakayama's Lemma \ref{algebra-lemma-NAK}, applied
to the finite module $N$, gives $N=0$.
\end{proof}


\subsection{Editorial treatment of Algebra lines 12256--12280}

\hypertarget{algebra-context-162}{}

\noindent\href{https://github.com/stacks/stacks-project/blob/a04446e57ec1fbc252a871afcec7752fb2807b14/algebra.tex\#L12256}{Original source, lines 12256--12280}. Complete original and reviewed TeX passages accompany this note. Every intervention is separately recorded; the following treatment is editorial.


\begin{lemma}
\label{editorial-source-lemma-intersection-powers-ideal-module}
Let $R$ be a Noetherian ring. Let $I \subset R$ be an ideal.
Let $M$ be a finite $R$-module. Let $N = \bigcap_n I^n M$.
\begin{enumerate}
\item For every prime $\mathfrak p$ containing $I$, there
exists $f \in R$, $f \not \in \mathfrak p$ such that $N_f = 0$.
\item If $I$ is contained in the Jacobson radical
of $R$, then $N = 0$.
\end{enumerate}
\end{lemma}

\begin{proof}
In fact there is one element $f\in1+I$ such that $fN=0$.
We prove this without a local hypothesis. The module $N$ is finite
by Lemma \ref{editorial-source-lemma-Noetherian-basic}. Applying
Lemma \ref{editorial-source-lemma-Artin-Rees} to $N\subset M$ at $n=c+1$
gives $N=IN$, since $I^nM\cap N=N$ for every $n\geq0$.
Choose generators $x_1,\ldots,x_t$ of $N$. Because $N=IN$,
there are actual coefficients $a_{ij}\in I$ with
$x_i=\sum_{j=1}^t a_{ij}x_j$ for $1\leq i\leq t$.
Put $A=(a_{ij})$ and $B=(\delta_{ij}-a_{ij})$.
The adjugate identity $\operatorname{adj}(B)B=\det(B)\,1_t$
applied to the column $(x_1,\ldots,x_t)^{\mathrm t}$ shows
that $f=\det(B)$ annihilates every $x_i$, hence annihilates $N$.
Retaining every coefficient and sign, this element is
$$
f=\sum_{\sigma\in\mathfrak S_t}\operatorname{sgn}(\sigma)
       \prod_{i=1}^t(\delta_{i,\sigma(i)}-a_{i,\sigma(i)}).
$$
Modulo $I$ the matrix $B$ is the identity matrix, so $f-1\in I$.
Equivalently $f=1-a$ for the specific element $a=1-f\in I$;
the determinant expression above records all contributions to $a$.
For an empty generating list, $N=0$ and the determinant of the
empty matrix is $f=1$, which has all the same properties.

For every prime $\mathfrak p$ containing $I$, the equality
$f\equiv1\pmod I$ implies $f\notin\mathfrak p$.
Since $fN=0$, inverting $f$ gives $N_f=0$. This proves (1)
with the same $f$ for all such primes. If $I$ is contained in the
Jacobson radical, $f=1-a$ is a unit: otherwise the proper ideal
$(f)$ would lie in a maximal ideal $\mathfrak m$, while
$a\in I\subset\mathfrak m$, contradicting $f+a=1$.
Thus $fN=0$ implies $N=0$, proving (2). If $R$ is the zero
ring then $M=N=0$ and both assertions hold directly.

The argument gives two exact descriptions of the same original
intersection. Let $W=1+I=\{1+b:b\in I\}$, a multiplicative
set because $(1+b)(1+d)=1+(b+d+bd)$ and $b+d+bd\in I$.
Then
$$
N=\{x\in M:(1-a)x=0\text{ for some }a\in I\}
=\ker\bigl(M\longrightarrow W^{-1}M\bigr).
$$
The inclusion from $N$ to the first set follows from the single
annihilator just constructed. Conversely, if $(1-a)x=0$,
then $x=ax$, and induction gives $x=a^nx\in I^nM$ for
every $n\geq0$. The localization kernel is exactly the set of
$x$ killed by some $w\in W$; writing $w=1-a$ gives the
second equality. This reasoning includes $I=R$, when $0\in W$,
$W^{-1}M=0$ and $N=M$, as well as $I=0$, when $W=\{1\}$
and $N=0$.

Moreover, $N$ is the largest submodule $L\subset M$ satisfying
$IL=L$. Indeed this equality implies $I^nL=L\subset I^nM$
for every $n$, so $L\subset N$, while $IN=N$ was proved above.
The quotient $M/N$ is $I$-adically separated:
$$
\bigcap_{n\geq0} I^n(M/N)=0.
$$
To prove this using the original quotient map $q:M\to M/N$,
note that $I^n(M/N)=q(I^nM)$ and $N\subset I^nM$.
Thus $q^{-1}(I^n(M/N))=I^nM+N=I^nM$.
A class belonging to every such image has any representative in
$\bigcap_n I^nM=N$ and is therefore zero. Finally, for every
$R$-linear map $h:M\to Q$ into an $I$-adically separated
module $Q$, one has $h(N)\subset\bigcap_n I^nQ=0$.
The original map $h$ therefore factors uniquely through $q$,
by $x+N\mapsto h(x)$. This proves the universal property of
this separated quotient with no finite-generation hypothesis on $Q$.
\end{proof}


\subsection{Editorial treatment of Algebra lines 12282--12294}

\hypertarget{algebra-context-163}{}

\noindent\href{https://github.com/stacks/stacks-project/blob/a04446e57ec1fbc252a871afcec7752fb2807b14/algebra.tex\#L12282}{Original source, lines 12282--12294}. Complete original and reviewed TeX passages accompany this note. Every intervention is separately recorded; the following treatment is editorial.


\begin{remark}
\label{editorial-source-remark-intersection-powers-ideal}
Lemma \ref{editorial-source-lemma-intersect-powers-ideal-module-zero} in particular
gives $\bigcap_{n\geq0}I^n=(0)$ for a proper ideal $I$ in a
Noetherian local ring. More generally let $R$ be Noetherian,
let $I\subset R$ be an ideal, and retain an original element
$f\in\bigcap_{n\geq0}I^n$. Apply the proof of
Lemma \ref{editorial-source-lemma-intersection-powers-ideal-module} with $M=R$.
It gives a single $g\in1+I$ annihilating the whole ideal
$\bigcap_{n\geq0}I^n$, so in particular $gf=0$.
Every prime containing $I$ avoids $g$, hence
$V(I)\subset D(g)$. In $R_g$, the equality $gf=0$ implies
$f/1=0$. Thus the original $f$, and indeed every member of the
intersection, is zero on the same principal open neighbourhood
$D(g)$ of $V(I)$, with the coefficient condition $g\equiv1\pmod I$.
When $I=R$, the closed set $V(I)$ is empty and $g=0$ is
allowed; no nonempty neighbourhood or nonzero denominator is
asserted in that case.
\end{remark}


\subsection{Editorial treatment of Algebra lines 12296--12323}

\hypertarget{algebra-context-164}{}

\noindent\href{https://github.com/stacks/stacks-project/blob/a04446e57ec1fbc252a871afcec7752fb2807b14/algebra.tex\#L12296}{Original source, lines 12296--12323}. Complete original and reviewed TeX passages accompany this note. Every intervention is separately recorded; the following treatment is editorial.


\begin{lemma}[Artin-Tate]
\label{editorial-source-lemma-Artin-Tate}
Let $R$ be a Noetherian ring. Let $S$ be a finitely
generated $R$-algebra. If $T \subset S$ is an $R$-subalgebra such
that $S$ is finitely generated as a $T$-module, then $T$ is of
finite type over $R$.
\end{lemma}

\begin{proof}
If $S=0$, then its subalgebra $T=0$ is the image of $R$ and
is of finite type over $R$. We may therefore assume $S\neq0$.
Choose original algebra generators $x_1,\ldots,x_n$ of $S$
over $R$ and original module generators $y_1,\ldots,y_m$ of
$S$ over $T$. Here $n$ may be zero, while $m\geq1$.
Choose coefficients $a_{ij},b_{ijk}\in T$ giving the full equations
$$
x_i=\sum_{j=1}^m a_{ij}y_j\quad(1\leq i\leq n),\qquad
y_i y_j=\sum_{k=1}^m b_{ijk}y_k\quad(1\leq i,j\leq m).
$$
Let $T'\subset T$ be the $R$-subalgebra generated by all these
coefficients. It is a quotient of a polynomial ring in finitely
many variables over the given $R$, and is Noetherian by
Lemma \ref{editorial-source-lemma-Noetherian-permanence}. This does not require
the structure map $R\to T'$ to be injective.
Form the algebra with every indicated relation
$$
S'=T'[Y_1,\ldots,Y_m]\big/
\left(Y_iY_j-\sum_{k=1}^m b_{ijk}Y_k:
1\leq i,j\leq m\right).
$$
Its underlying $T'$-module is generated by the classes of
$1,Y_1,\ldots,Y_m$. Indeed a monomial of total degree at least
two has a factor $Y_iY_j$. Replacing that pair by
$\sum_{k=1}^m b_{ijk}Y_k$ expresses it as a sum of monomials
of total degree one less, with their actual coefficients
$b_{ijk}$. Induction on total degree expresses every monomial,
and then every polynomial, in the asserted span. This argument
retains the constant generator $1$ and does not assume that one
of the chosen $y_j$ equals $1$.

The coefficient inclusion $T'\subset S$ and the assignments
$Y_j\mapsto y_j$ define a $T'$-algebra map $\rho:S'\to S$:
every displayed relation maps to its given zero expression in $S$.
For each $i$, the element $\sum_{j=1}^m a_{ij}Y_j$ maps to
the original $x_i$. The image also contains the image of $R$,
so it contains all of $S$, even if the algebra-generator list is
empty. Therefore $\rho$ is surjective and $S$ is finite over
$T'$. No assertion that $\rho$ is injective is needed or made.
Now $T\subset S$ is a $T'$-submodule. By
Lemma \ref{editorial-source-lemma-Noetherian-basic}, choose a finite generating
list $z_1,\ldots,z_q$ for this module. The exact equality
$$
T=R[a_{ij},b_{ijk},z_1,\ldots,z_q]
$$
uses all $a_{ij}$ for $1\leq i\leq n$, $1\leq j\leq m$,
all $b_{ijk}$ for $1\leq i,j,k\leq m$, and every $z_\ell$
for $1\leq\ell\leq q$, together with the image of $R$.
To verify it, every $t\in T$ is a sum
$\sum_{\ell=1}^q c_\ell z_\ell$ with $c_\ell\in T'$;
each $c_\ell$ is a polynomial in the original $a_{ij},b_{ijk}$
with coefficients from $R$. This gives containment in the
right-hand subalgebra. The reverse containment follows because
all its generators already lie in $T$. Hence $T$ is of finite
type over $R$, as required.
\end{proof}


\subsection{Editorial treatment of Algebra lines 12346--12699}

\hypertarget{algebra-context-165}{}

\noindent\href{https://github.com/stacks/stacks-project/blob/a04446e57ec1fbc252a871afcec7752fb2807b14/algebra.tex\#L12346}{Original source, lines 12346--12699}. Complete original and reviewed TeX passages accompany this note. Every intervention is separately recorded; the following treatment is editorial.


\subsubsection{Length}
\label{editorial-source-section-length}

\begin{definition}
\label{editorial-source-definition-length}
Let $R$ be a ring. For any $R$-module $M$
we define the {\it length} of $M$ over $R$ by the
formula
$$
\text{length}_R(M)
=
\sup
\{
n \in \mathbf Z_{\geq0}
\mid
\exists\ 0 = M_0 \subset M_1 \subset \ldots \subset M_n = M,
\text{ }M_i \not = M_{i + 1}\quad(0\leq i<n)
\}.
$$
\end{definition}

\noindent
The supremum is taken in $\mathbf Z_{\geq0}\cup\{\infty\}$.
For $M=0$ the length-zero chain is allowed, so its length is zero.
In other words it is the supremum of the lengths of chains
of submodules. There is an obvious notion of when a chain
of submodules is a refinement of another. This gives a
partial ordering on the collection of all chains of submodules,
with the smallest chain having the shape $0 = M_0 \subset M_1 = M$
if $M$ is not zero.
We note the obvious fact that if the length of
$M$ is finite, then every chain can be refined to a
maximal chain. But it is not as obvious that all maximal
chains have the same length (as we will see later).

\begin{lemma}
\label{editorial-source-lemma-finite-length-finite}
\begin{slogan}
Modules of finite length are finite.
\end{slogan}
Let $R$ be a ring.
Let $M$ be an $R$-module.
If $\text{length}_R(M) < \infty$ then $M$ is a finite $R$-module.
\end{lemma}

\begin{proof}
Write $\ell=\operatorname{length}_R(M)<\infty$.
If $M=0$, its empty generating list suffices. Otherwise start with
$N_0=0$. Whenever $N_j\neq M$, choose
$x_{j+1}\in M\setminus N_j$ and put
$N_{j+1}=N_j+Rx_{j+1}$. Each such step is strict.
There cannot be $\ell+1$ such steps: adjoining the final endpoint
$M$ if it is not already present would give a chain from zero
to $M$ with at least $\ell+1$ strict steps, contrary to the
definition of length. Therefore $N_j=M$ for some $j\leq\ell$,
and $x_1,\ldots,x_j$ generate the original $M$.

The same bound applies to every subquotient $N/P$ of $M$.
Indeed a strict chain in $N/P$ pulls back under the original
quotient map $N\to N/P$ to a strict chain from $P$ to $N$.
Adjoin zero if $P\neq0$ and $M$ if $N\neq M$. This is a
chain of submodules of $M$, so the original chain has at most
$\ell$ strict steps. Thus $N/P$ has length at most $\ell$
and is generated by at most its own length many elements by the
same construction. Moreover, no ascending or descending chain
of submodules of $M$ can have more than $\ell$ strict changes:
any finite selection of its strict changes, ordered increasingly
and completed with the endpoints, gives such a chain in $M$.
In particular both chain conditions hold for $M$.
\end{proof}

\begin{lemma}
\label{editorial-source-lemma-length-additive}
\begin{slogan}
Length is additive in short exact sequences.
\end{slogan}
If $0 \to M' \to M \to M'' \to 0$
is a short exact sequence of modules over $R$ then
the length of $M$ is the sum of the
lengths of $M'$ and $M''$.
\end{lemma}

\begin{proof}
Write the original sequence as
$$
0\longrightarrow M'\xrightarrow{u}M\xrightarrow{q}M''
\longrightarrow0.
$$
Lengths take values in $\mathbf Z_{\geq0}\cup\{\infty\}$,
with $a+\infty=\infty+a=\infty$ for every such $a$.
Given strict chains from zero to $M'$ and $M''$, of lengths
$a$ and $b$, respectively, first apply $u$ to the chain in
$M'$. Its last member is $u(M')=\ker(q)$.
Continue with the inverse images under $q$ of the nonzero steps
of the chain in $M''$. Injectivity of $u$ and surjectivity of
$q$ show that every indicated step is strict. This yields a
chain in $M$ of length $a+b$. If either length on the right is
infinite, finite chains of arbitrarily large length give
$\operatorname{length}_R(M)=\infty$. If both are finite, each
is an attained integer supremum, so choosing chains with those
lengths proves
$\operatorname{length}_R(M)\geq
\operatorname{length}_R(M')+\operatorname{length}_R(M'')$.
This construction includes a zero module with its length-zero chain.

Conversely, let $0=M_0\subsetneq\cdots\subsetneq M_n=M$
be a finite strict chain. Put
$M_i'=u^{-1}(M_i)$ and $M_i''=q(M_i)$, using the original
maps rather than silently identifying $M'$ with a different module.
Let $a$ and $b$ count only the strict changes in these two
nondecreasing sequences. Removing their repeated terms gives
chains of lengths $a$ and $b$ in $M'$ and $M''$.
At least one changes at each of the $n$ original steps.
For if $M_i'=M_{i+1}'$ and $M_i''=M_{i+1}''$, take
$x\in M_{i+1}$. There is $y\in M_i$ with $q(x)=q(y)$.
Then $x-y=u(z)$ for some $z\in M'$ and
$u(z)\in M_{i+1}$, so $z\in M_{i+1}'=M_i'$.
Consequently $u(z)\in M_i$ and $x=y+u(z)\in M_i$,
contradicting strictness. It follows that
$$
n\leq a+b\leq
\operatorname{length}_R(M')+\operatorname{length}_R(M'').
$$
Taking the supremum over the original chains proves the reverse
inequality, including infinite lengths. This proves additivity.
\end{proof}

\begin{lemma}
\label{editorial-source-lemma-length-infinite}
Let $R$ be a local ring with maximal ideal $\mathfrak m$.
If $M$ is an $R$-module and $\mathfrak m^n M \not = 0$ for all $n \geq 0$,
then $\text{length}_R(M) = \infty$. In other words, if $M$
has finite length then $\mathfrak m^nM = 0$ for some $n$.
\end{lemma}

\begin{proof}
Fix an integer $n\geq0$ such that $\mathfrak m^nM\neq0$.
A nonzero member of this module is a finite sum of products
$f_1\cdots f_nx$, with every $f_i\in\mathfrak m$ and $x\in M$.
At least one of these products is nonzero; choose it and put
$y_j=f_1\cdots f_jx$ for $0\leq j\leq n$, where $y_0=x$.
All $y_j$ are nonzero, since multiplying a zero earlier product
by the remaining original factors would give $y_n=0$.
The inclusions
$$
0\subsetneq Ry_n\subsetneq Ry_{n-1}\subsetneq\cdots
\subsetneq Ry_0\subset M
$$
are strict up to $Ry_0$. To verify any step with $1\leq j\leq n$,
suppose $Ry_{j-1}=Ry_j$. Then
$y_{j-1}=g y_j=g f_j y_{j-1}$ for some $g\in R$, hence
$(1-gf_j)y_{j-1}=0$. Since $gf_j\in\mathfrak m$, the original
factor $1-gf_j$ is a unit in the local ring. This forces
$y_{j-1}=0$, a contradiction. The first inclusion is strict
because $y_n\neq0$. Append $M$ if necessary; if $Ry_0=M$,
there is no extra final step. Thus
$\operatorname{length}_R(M)\geq n+1$.
For $n=0$ this uses the nonzero element $x$ and the empty product,
and gives the same bound.

If every $\mathfrak m^nM$ is nonzero, these bounds for all $n$
give infinite length. They also give the sharper quantitative
conclusion
$$
\mathfrak m^\ell M=0
\quad\text{if}\quad
\ell=\operatorname{length}_R(M)<\infty.
$$
Indeed a nonzero left-hand module would imply
$\ell\geq\ell+1$. When $\ell=0$, the definition of length
already gives $M=0$, so the displayed assertion remains valid.
\end{proof}

\begin{lemma}
\label{editorial-source-lemma-length-independent}
Let $R \to S$ be a ring map. Let $M$ be an $S$-module.
We always have $\text{length}_R(M) \geq \text{length}_S(M)$.
If $R \to S$ is surjective then equality holds.
\end{lemma}

\begin{proof}
Every $S$-submodule of the original $M$ is an $R$-submodule by
restriction of scalars. Hence every strict chain for the former
action gives the identical strict chain for the latter, proving
the inequality even for infinite length.
If $R\to S$ is surjective, then for any $R$-submodule $N\subset M$,
$s\in S$ and $x\in N$, choose $r\in R$ mapping to $s$.
The equality $sx=rx$ puts $sx$ in $N$. Thus the two collections
of submodules, and all their chains, agree exactly.

The proof gives a weaker sufficient hypothesis for equality:
surjectivity of the original composite
$R\to S/\operatorname{Ann}_S(M)$ suffices.
For then any $s\in S$ has an $r\in R$ such that their images
in this quotient agree; their difference annihilates every
original $x\in M$, so once again $sx=rx$.
The original actions are retained and their equality on each
element proves equality of the submodule lattices. When $M=0$,
both lattices have a single member and both lengths are zero.
\end{proof}

\begin{lemma}
\label{editorial-source-lemma-dimension-is-length}
Let $R$ be a ring with maximal ideal $\mathfrak m$.
Suppose that $M$ is an $R$-module with
$\mathfrak m M  =  0$. Then the length of $M$ as
an $R$-module agrees with its dimension as a vector space over
$R/\mathfrak m$, interpreting every infinite dimension as $\infty$.
The length is finite if and only if $M$ is a finite $R$-module.
\end{lemma}

\begin{proof}
Set $k=R/\mathfrak m$. The original action on $M$ factors
through $k$, since $\mathfrak mM=0$; explicitly
$(r+\mathfrak m)x=rx$, which is independent of the representative.
An $R$-submodule is exactly a $k$-vector subspace by
Lemma \ref{editorial-source-lemma-length-independent}.
If $\dim_k M=d<\infty$, a strict inclusion of vector subspaces
increases dimension by at least one. Thus every strict chain has
at most $d$ steps. Given a basis $x_1,\ldots,x_d$, the spans
of its first $i$ vectors, for $0\leq i\leq d$, give a chain
of exactly $d$ steps. Therefore the length is $d$, including $d=0$.
If the dimension is infinite, for each finite $n$ choose
$n$ independent vectors. Their initial spans form $n$ strict
steps, and adjoining $M$ as endpoint if needed gives length at
least $n$. The length is therefore $\infty$; it is this extended
integer value, not the particular infinite basis cardinal, that
is asserted here.
Finally, a finite $R$-generating list is exactly a finite
$k$-spanning list because the two actions just described have the
same spans. Such a list contains a basis after removing dependent
vectors, and a finite basis is itself a finite $R$-generating
list. This proves the final equivalence.
\end{proof}

\begin{lemma}
\label{editorial-source-lemma-length-localize}
Let $R$ be a ring. Let $M$ be an $R$-module. Let $S \subset R$ be
a multiplicative subset. Then
$\text{length}_R(M) \geq \text{length}_{S^{-1}R}(S^{-1}M)$.
\end{lemma}

\begin{proof}
Let $j:M\to S^{-1}M$ be the original map $x\mapsto x/1$.
For an $S^{-1}R$-submodule $L\subset S^{-1}M$, put
$j^{-1}(L)=\{x\in M:x/1\in L\}$. Its localization has image
exactly $L$. Indeed if $x/s\in L$, then $x/1=s(x/s)\in L$,
so $x\in j^{-1}(L)$; conversely $x/1\in L$ implies
$x/s=(1/s)(x/1)\in L$. Localization of the inclusion
$j^{-1}(L)\subset M$ is injective, so this identifies the actual
localized submodule with $L$.
Consequently a strict chain $0=L_0\subsetneq\cdots\subsetneq
L_n=S^{-1}M$ pulls back to a strict chain
$j^{-1}(L_0)\subsetneq\cdots\subsetneq j^{-1}(L_n)=M$:
equality of two adjacent inverse images would give equality of
their localized images. The first inverse image is $\ker(j)$,
which need not be zero. Adjoin zero when it is nonzero. This
produces at least $n$ strict steps from zero to $M$, proving
the length inequality. If $S^{-1}M=0$, its length is zero and
the inequality holds directly; no nonzero localization is assumed.
For finite length, the proof of Lemma \ref{editorial-source-lemma-simple-pieces}
below computes the exact loss: it is the number of original
composition factors whose maximal ideals meet $S$, counted
with their multiplicities.
\end{proof}

\begin{lemma}
\label{editorial-source-lemma-length-finite}
Let $R$ be a ring with finitely generated
maximal ideal $\mathfrak m$. (For example, this holds if $R$ is Noetherian.)
Suppose that $M$ is a finite $R$-module with
$\mathfrak m^n M  =  0$ for some integer $n\geq0$.
Then $\text{length}_R(M) < \infty$.
\end{lemma}

\begin{proof}
If $n=0$, then $M=\mathfrak m^0M=0$ and the claim follows.
Assume $n\geq1$ and write $k=R/\mathfrak m$.
Choose original generators $x_1,\ldots,x_r$ of $M$ and
$f_1,\ldots,f_t$ of $\mathfrak m$. For $i\geq0$, put
$$
E_i=\{(e_1,\ldots,e_t)\in\mathbf Z_{\geq0}^{,t}:
e_1+\cdots+e_t=i\}.
$$
The quotient $\mathfrak m^iM/\mathfrak m^{i+1}M$ is killed
by $\mathfrak m$ and has the following explicit surjection
of $k$-vector spaces:
$$
\begin{aligned}
k^{\{1,\ldots,r\}\times E_i}&\longrightarrow
\mathfrak m^iM/\mathfrak m^{i+1}M,\\
(\bar a_{j,e})_{j,e}&\longmapsto
\sum_{j=1}^r\sum_{e\in E_i}
a_{j,e} f_1^{e_1}\cdots f_t^{e_t}x_j
\pmod{\mathfrak m^{i+1}M}.
\end{aligned}
$$
Changing any representative $a_{j,e}$ by an element of
$\mathfrak m$ changes the indicated summand by a member of
$\mathfrak m^{i+1}M$, so the map is well defined. The products
in the specified ideal generators generate $\mathfrak m^i$;
expanding a member of $\mathfrak m^iM$ in these products and
the original $x_j$ proves surjectivity. Thus each layer has
finite dimension over $k$, and its length is that dimension by
Lemma \ref{editorial-source-lemma-dimension-is-length}. Applying additivity
successively to the original filtration by powers gives the exact formula
$$
\operatorname{length}_R(M)
=\sum_{i=0}^{n-1}\dim_k
 (\mathfrak m^iM/\mathfrak m^{i+1}M)
\leq r\sum_{i=0}^{n-1}\#E_i<\infty.
$$
Repeated adjacent powers contribute their actual zero-dimensional
quotients. If $t\geq1$, distributing $i$ positions among $t$
nonnegative exponents gives $\#E_i=\binom{t+i-1}{i}$: encode
the tuple by $i$ identical positions separated into $t$ blocks
by $t-1$ separators, choosing the $i$ positions among the
$t+i-1$ total positions. If $t=0$, there is one empty tuple
in $E_0$ and no tuple in $E_i$ for $i>0$.

In fact finite generation of $\mathfrak m$ can be replaced by
finite dimension of the original cotangent space
$\mathfrak m/\mathfrak m^2$ over $k$.
Choose lifts $f_1,\ldots,f_t\in\mathfrak m$ of a finite
spanning list and put $J=(f_1,\ldots,f_t)$.
Then $\mathfrak m=J+\mathfrak m^2$.
For $i\geq1$, expanding this actual ideal equality retains all terms:
$$
\mathfrak m^i
=(J+\mathfrak m^2)^i
=\sum_{h=0}^i J^{i-h}\mathfrak m^{2h}.
$$
Every term with $h\geq1$ lies in $\mathfrak m^{i+h}$,
hence in $\mathfrak m^{i+1}$. Conversely both $J^i$ and
$\mathfrak m^{i+1}$ lie in $\mathfrak m^i$. Therefore
$\mathfrak m^i=J^i+\mathfrak m^{i+1}$ and
$\mathfrak m^iM=J^iM+\mathfrak m^{i+1}M$.
All the terms with $h\geq1$ have zero image under this particular
quotient map; no equality $\mathfrak m=J$ has been assumed.
The same displayed coefficient map is consequently surjective
for every layer, with these original lifts $f_j$.
It proves the same finite length conclusion and the same explicit
bound whenever $M$ is finite and $\mathfrak m^nM=0$.
This argument covers $t=0$ as well: then
$\mathfrak m^iM=\mathfrak m^{i+1}M$ for $i\geq1$, so
nilpotence on $M$ forces all those layers and $\mathfrak mM$ to vanish.
\end{proof}

\begin{definition}
\label{editorial-source-definition-simple-module}
Let $R$ be a ring. Let $M$ be an $R$-module.
We say $M$ is {\it simple} if $M \not = 0$ and
every submodule of $M$ is either equal to $M$ or
to $0$.
\end{definition}

\begin{lemma}
\label{editorial-source-lemma-characterize-length-1}
Let $R$ be a ring. Let $M$ be an $R$-module.
The following are equivalent:
\begin{enumerate}
\item $M$ is simple,
\item $\text{length}_R(M) = 1$, and
\item $M \cong R/\mathfrak m$ for some maximal ideal
$\mathfrak m \subset R$.
\end{enumerate}
\end{lemma}

\begin{proof}
Let $\mathfrak m$ be a maximal ideal of $R$.
By Lemma \ref{editorial-source-lemma-dimension-is-length} the module
$R/\mathfrak m$ has length $1$. The equivalence of
the first two assertions is tautological.
Suppose that $M$ is simple. Choose $x \in M$, $x \not = 0$.
As $M$ is simple we have $M = R \cdot x$.
Let $I \subset R$ be the annihilator of $x$, i.e.,
$I = \{f \in R \mid fx = 0\}$. The map $R/I \to M$,
$f \bmod I \mapsto fx$ is an isomorphism, hence
$R/I$ is a simple $R$-module. Since $R/I \not = 0$ we see $I \not = R$.
Let $\mathfrak m$ be a maximal ideal containing $I$.
If $I \not = \mathfrak m$, then $\mathfrak m /I \subset R/I$
is a nontrivial submodule contradicting the simplicity
of $R/I$. Hence we see $I = \mathfrak m$ as desired.
\end{proof}

\begin{lemma}
\label{editorial-source-lemma-simple-pieces}
Let $R$ be a ring. Let $M$ be a finite length $R$-module.
Choose any maximal chain of submodules
$$
0 = M_0 \subset M_1 \subset M_2 \subset \ldots \subset M_n = M
$$
with $M_i \not = M_{i-1}$, $i = 1, \ldots, n$. Then
\begin{enumerate}
\item $n = \text{length}_R(M)$,
\item each $M_i/M_{i-1}$ is simple,
\item each $M_i/M_{i-1}$ is of the form
$R/\mathfrak m_i$ for some maximal ideal $\mathfrak m_i$,
\item given a maximal ideal $\mathfrak m \subset R$
we have
$$
\# \{i \mid \mathfrak m_i = \mathfrak m\}
=
\text{length}_{R_{\mathfrak m}} (M_{\mathfrak m}).
$$
\end{enumerate}
\end{lemma}

\begin{proof}
For each $i$, let $q_i:M_i\to M_i/M_{i-1}$ be the original
quotient map. If its nonzero target were not simple, a nonzero
proper submodule would have inverse image strictly between
$M_{i-1}$ and $M_i$. Inserting it would contradict maximality
of the given chain. Thus every quotient is simple, and
Lemma \ref{editorial-source-lemma-characterize-length-1} gives an isomorphism
$M_i/M_{i-1}\cong R/\mathfrak m_i$ and length one.
Applying Lemma \ref{editorial-source-lemma-length-additive} to the successive
short exact sequences gives $\operatorname{length}_R(M)=n$.
For $M=0$ the chain has no strict steps, the lists of quotients
and maximal ideals are empty, and this equality reads $0=0$.

For a maximal ideal $\mathfrak m$, localization gives the exact sequences
$$
0\longrightarrow(M_{i-1})_{\mathfrak m}
\longrightarrow(M_i)_{\mathfrak m}
\longrightarrow(M_i/M_{i-1})_{\mathfrak m}
\longrightarrow0.
$$
If $\mathfrak m_i\neq\mathfrak m$, choose
$a\in\mathfrak m_i\setminus\mathfrak m$; such an element exists
because distinct maximal ideals cannot contain one another.
It annihilates $R/\mathfrak m_i$ and is inverted at
$\mathfrak m$, so this quotient localizes to zero.
If $\mathfrak m_i=\mathfrak m$, the isomorphism with the
residue field of the original local ring is explicitly
$$
(R/\mathfrak m)_{\mathfrak m}
\longrightarrow R_{\mathfrak m}/\mathfrak mR_{\mathfrak m},
\qquad
(r+\mathfrak m)/s\longmapsto r/s+\mathfrak mR_{\mathfrak m}.
$$
Its inverse sends $r/s+\mathfrak mR_{\mathfrak m}$ to
$(r+\mathfrak m)/s$. The localization fraction relations hold
on both sides, and coefficients in $\mathfrak m$ act as zero,
so both maps are well defined; their composites are identities.
The resulting nonzero quotient has length one over $R_{\mathfrak m}$.
Additivity proves exactly the asserted count.

More generally let $S\subset R$ be any multiplicative subset.
For each factor $R/\mathfrak m_i$, its localization is zero
if $S\cap\mathfrak m_i\neq\varnothing$. Otherwise every
$s\in S$ has nonzero, hence invertible, image in the field
$R/\mathfrak m_i$, and
$$
S^{-1}(R/\mathfrak m_i)\longrightarrow R/\mathfrak m_i,
\qquad (r+\mathfrak m_i)/s\longmapsto
(r+\mathfrak m_i)(s+\mathfrak m_i)^{-1}
$$
is an isomorphism, with inverse $z\mapsto z/1$.
It identifies this localized module with the simple quotient
$S^{-1}R/S^{-1}\mathfrak m_i$.
Localizing every short exact sequence and adding the lengths gives
$$
\operatorname{length}_{S^{-1}R}(S^{-1}M)
=\#\{i: S\cap\mathfrak m_i=\varnothing\},
$$
and therefore the exact loss of length is
$\#\{i:S\cap\mathfrak m_i\neq\varnothing\}$.
This refines Lemma \ref{editorial-source-lemma-length-localize} for finite length.
For an arbitrary prime $\mathfrak p$, the same calculation with
$S=R\setminus\mathfrak p$ retains a factor precisely when
$\mathfrak m_i\subset\mathfrak p$, hence precisely when
$\mathfrak p=\mathfrak m_i$.
Thus the original support of $M$ is the finite set of distinct
$\mathfrak m_i$, and
$$
\operatorname{length}_R(M)
=\sum_{\mathfrak m\in\operatorname{Supp}_R(M)}
  \operatorname{length}_{R_{\mathfrak m}}(M_{\mathfrak m}).
$$
In particular the multiplicity of each simple module
$R/\mathfrak m$ is independent of the chosen maximal chain.
The ideal $\mathfrak m$ is determined by that module, since
its annihilator is exactly $\mathfrak m$; hence these counts
determine the multiset of simple factors, not only the number
of steps.

The two chain conditions also characterize finite length for
an arbitrary $R$-module. Necessity was proved in the proof of
Lemma \ref{editorial-source-lemma-finite-length-finite}.
Conversely assume both conditions. In every nonzero submodule
$L$, the ascending chain condition supplies a maximal proper
submodule: if none existed, starting from zero one could choose
an infinite strictly increasing sequence of proper submodules
of $L$. Successively choose such a maximal proper submodule
inside each nonzero module, starting from $M$.
The descending chain condition forces this sequence to reach
zero after finitely many strict steps. Reversing it gives a
finite chain whose quotients are simple. Each quotient has
length one by Lemma \ref{editorial-source-lemma-characterize-length-1}, so
additivity gives finite length for the original $M$.
\end{proof}

\begin{lemma}
\label{editorial-source-lemma-pushdown-module}
Let $A$ be a local ring with maximal ideal $\mathfrak m$.
Let $B$ be a semi-local ring with maximal ideals $\mathfrak m_i$,
$i = 1, \ldots, n$.
Suppose that $A \to B$ is a homomorphism such that each $\mathfrak m_i$
lies over $\mathfrak m$ and such that
$$
[\kappa(\mathfrak m_i) : \kappa(\mathfrak m)] < \infty.
$$
Let $M$ be a $B$-module of finite length.
Then
$$
\text{length}_A(M) = \sum\nolimits_{i = 1, \ldots, n}
[\kappa(\mathfrak m_i) : \kappa(\mathfrak m)]
\text{length}_{B_{\mathfrak m_i}}(M_{\mathfrak m_i}),
$$
in particular $\text{length}_A(M) < \infty$.
\end{lemma}

\begin{proof}
Choose a maximal chain of $B$-submodules of the original $M$
as in Lemma \ref{editorial-source-lemma-simple-pieces}. Each of its quotients is
isomorphic to $B/\mathfrak m_{i(j)}=\kappa(\mathfrak m_{i(j)})$.
The original map $A\to B$ induces an injective field map
$\kappa(\mathfrak m)=A/\mathfrak m\to B/\mathfrak m_i$:
its kernel before quotienting is exactly the contraction
$\mathfrak m_i\cap A=\mathfrak m$.
The $A$-action on this quotient is the action through that map.
By Lemma \ref{editorial-source-lemma-dimension-is-length}, its $A$-length is
exactly $[\kappa(\mathfrak m_i):\kappa(\mathfrak m)]$.
Additivity along the original chain, now regarded as a chain
of $A$-submodules, gives the sum of these degrees over its factors.
Lemma \ref{editorial-source-lemma-simple-pieces} identifies the number of factors
at $\mathfrak m_i$ with
$\operatorname{length}_{B_{\mathfrak m_i}}(M_{\mathfrak m_i})$.
Grouping the finite sum by those maximal ideals proves the
displayed formula and its finiteness. When $M=0$, the chain
is empty and all terms are zero, as required.

The assumptions need only concern the maximal ideals in the
support of $M$. More precisely, let $\alpha:A\to B$ be any
ring homomorphism and let $M$ have finite length over $B$.
The preceding composition-factor proof shows that
$Q=\operatorname{Supp}_B(M)$ is a finite set of maximal ideals.
For each $\mathfrak q\in Q$, put
$\mathfrak p_{\mathfrak q}=\alpha^{-1}(\mathfrak q)$.
Then $M$ has finite length over $A$ if and only if every
$\mathfrak p_{\mathfrak q}$ is maximal in $A$ and every
extension $\kappa(\mathfrak q)/\kappa(\mathfrak p_{\mathfrak q})$
has finite degree. Under these conditions the exact formula is
$$
\operatorname{length}_A(M)
=\sum_{\mathfrak q\in Q}
[\kappa(\mathfrak q):\kappa(\mathfrak p_{\mathfrak q})]
\operatorname{length}_{B_{\mathfrak q}}(M_{\mathfrak q}).
$$
To prove the criterion in both directions, fix one composition
factor $k=B/\mathfrak q$. The original map identifies
$D=A/\mathfrak p_{\mathfrak q}$ with a subring of the field $k$.
If $\mathfrak p_{\mathfrak q}$ is maximal, $D$ is its residue
field and the $A$-length of $k$ is its dimension over $D$ by
Lemma \ref{editorial-source-lemma-dimension-is-length}, finite exactly in the
stated finite-degree case.
If $\mathfrak p_{\mathfrak q}$ is not maximal, the domain $D$
is not a field. Choose a nonzero nonunit $a\in D$.
Its image is nonzero in $k$, and the original $A$-submodules
$$
D\subsetneq a^{-1}D\subsetneq a^{-2}D\subsetneq\cdots\subset k
$$
are strictly increasing. Indeed equality of two consecutive
members would give $a^{-(j+1)}=a^{-j}d$ for some $d\in D$,
and multiplying in $k$ by $a^{j+1}$ would yield $1=ad$ in
the embedded $D$, contradicting the choice of $a$.
Adjoining zero and the endpoint $k$ gives finite chains of
arbitrarily large length, so $\operatorname{length}_A(k)=\infty$.
Finally, the original finite composition chain in $M$ remains
an exact filtration of $A$-modules. Additivity shows that its
$A$-length is the sum of its factor lengths, and is finite
exactly when each occurring factor has finite $A$-length.
Every $\mathfrak q\in Q$ occurs with the strictly positive
multiplicity specified above. This proves necessity, sufficiency
and the formula without any local or semilocal hypothesis on
the whole rings. If $Q$ is empty, $M=0$ and the criterion is
vacuous with an empty sum equal to zero.
\end{proof}

\begin{lemma}
\label{editorial-source-lemma-pullback-module}
Let $A \to B$ be a flat local homomorphism of local rings.
Then for any $A$-module $M$ we have
$$
\text{length}_A(M) \text{length}_B(B/\mathfrak m_AB)
=
\text{length}_B(M \otimes_A B).
$$
Here products of lengths use $0\cdot\infty=\infty\cdot0=0$.
In particular, if $\text{length}_B(B/\mathfrak m_AB) < \infty$
then $M$ has finite length if and only if $M \otimes_A B$ has finite length.
\end{lemma}

\begin{proof}
Retain the original flat local homomorphism $\alpha:A\to B$.
It is faithfully flat by Lemma \ref{editorial-source-lemma-local-flat-ff}.
The original closed fibre
$F=B/\mathfrak m_AB$ is nonzero, because locality implies
$\mathfrak m_AB\subset\mathfrak m_B\neq B$.
Thus $d=\operatorname{length}_B(F)$ lies in
$\mathbf Z_{\geq1}\cup\{\infty\}$.
When $M=0$, also $M\otimes_A B=0$, and the formula reads
$0=0$ using $0\cdot\infty=0$ if needed.

For every strict chain $0=M_0\subsetneq\cdots\subsetneq M_n=M$,
flatness gives injections of the actual modules $M_i\otimes_A B$
into $M\otimes_A B$ and exact sequences
$$
0\longrightarrow M_{i-1}\otimes_A B
\longrightarrow M_i\otimes_A B
\longrightarrow (M_i/M_{i-1})\otimes_A B
\longrightarrow0.
$$
Every quotient on the right is nonzero by faithful flatness,
so every step remains strict. Therefore infinite length of $M$
implies infinite length after base change, as required by
$\infty\cdot d=\infty$ for the positive value $d$.

If $0<\ell=\operatorname{length}_A(M)<\infty$, choose a
composition chain of length $\ell$. Since $A$ is local, each
simple factor is isomorphic to the original $A/\mathfrak m_A$.
The exact sequences above identify each new factor with $F$.
In the tensor order being used, define
$\Phi:(A/\mathfrak m_A)\otimes_A B\to F$ and
$\Psi:F\to(A/\mathfrak m_A)\otimes_A B$ by
$$
\begin{aligned}
\Phi((a+\mathfrak m_A)\otimes b)&=\alpha(a)b+\mathfrak m_AB,\\
\Psi(b+\mathfrak m_AB)&=(1+\mathfrak m_A)\otimes b.
\end{aligned}
$$
The first map respects tensor balancing and annihilates the
indicated representatives from $\mathfrak m_A$.
For the second map, every summand $\alpha(a)b$ with
$a\in\mathfrak m_A$ maps to
$(a+\mathfrak m_A)\otimes b=0$. Thus both maps are well
defined, and tensor balancing verifies that their composites
are identities on each pure tensor and quotient class.
Additivity along the filtration obtained by tensoring gives
$\operatorname{length}_B(M\otimes_A B)=\ell d$,
also when $d=\infty$, since a positive finite number of
infinite summands has infinite sum.
Together with the zero and infinite cases this proves the formula.

In particular, without any advance finiteness assumption on $d$,
the exact finite-length criterion is
$$
\begin{aligned}
&\operatorname{length}_B(M\otimes_A B)<\infty\\
&\quad\Longleftrightarrow\quad
M=0\ \text{or}\quad
\bigl(\operatorname{length}_A(M)<\infty\ \text{and}\ d<\infty\bigr).
\end{aligned}
$$
The positivity of $d$, faithful flatness and the proved equality
give both directions. When $d<\infty$, this specializes to
the stated equivalence.
\end{proof}

\begin{lemma}
\label{editorial-source-lemma-pullback-transitive}
Let $A \to B \to C$ be flat local homomorphisms of local rings. Then
$$
\text{length}_B(B/\mathfrak m_A B)
\text{length}_C(C/\mathfrak m_B C)
=
\text{length}_C(C/\mathfrak m_A C)
$$
\end{lemma}

\begin{proof}
Write the original maps as $\alpha:A\to B$ and $\beta:B\to C$.
Apply Lemma \ref{editorial-source-lemma-pullback-module} to $B\to C$ and the
original module $B/\mathfrak m_AB$. The receiving module in
that formula is identified with $C/\mathfrak m_AC$ by the maps
$\Phi:(B/\mathfrak m_AB)\otimes_B C\to C/\mathfrak m_AC$
and $\Psi:C/\mathfrak m_AC\to(B/\mathfrak m_AB)\otimes_B C$:
$$
\begin{aligned}
\Phi((b+\mathfrak m_AB)\otimes c)&=\beta(b)c+\mathfrak m_AC,\\
\Psi(c+\mathfrak m_AC)&=(1+\mathfrak m_AB)\otimes c.
\end{aligned}
$$
Here $\mathfrak m_AC$ is generated by the actual elements
$\beta(\alpha(a))$ for $a\in\mathfrak m_A$.
The first map respects balancing and sends any member of
$\mathfrak m_AB$ to that ideal in $C$.
For the second map, a generator times $c$ maps to
$(\alpha(a)+\mathfrak m_AB)\otimes c=0$.
Thus both are well defined, and their composites are identities
by the balancing relation. The base-change length formula now
gives precisely the claimed equality. Each of the three original
closed fibres is nonzero because the maps are local. Their
lengths are therefore positive or infinite, so their product
has no zero-times-infinity ambiguity.
\end{proof}














\subsection{Editorial treatment of Algebra lines 12713--12721}

\hypertarget{algebra-context-166}{}

\noindent\href{https://github.com/stacks/stacks-project/blob/a04446e57ec1fbc252a871afcec7752fb2807b14/algebra.tex\#L12713}{Original source, lines 12713--12721}. Complete original and reviewed TeX passages accompany this note. Every intervention is separately recorded; the following treatment is editorial.


\begin{lemma}
\label{editorial-source-lemma-finite-dimensional-algebra}
Suppose $R$ is a finite dimensional algebra over a field.
Then $R$ is Artinian.
\end{lemma}

\begin{proof}
Let $k$ be the given field and let $d=\dim_k R<\infty$.
Every ideal is a $k$-vector subspace of the original $R$,
because scalar multiplication is multiplication by the image
of an element of $k$ in $R$. In any strict descending chain
of ideals, the dimensions strictly decrease among the integers
$0,\ldots,d$. Thus there are at most $d$ strict changes and
every descending chain stabilizes, proving that $R$ is Artinian.
The same argument bounds every strict increasing chain.
In particular, directly from the definition of length,
$\operatorname{length}_R(R)\leq d$.
If $R=0$, then $d=0$, its only ideal is zero, and these
assertions include its length-zero chain.
\end{proof}


\subsection{Editorial treatment of Algebra lines 12723--12736}

\hypertarget{algebra-context-167}{}

\noindent\href{https://github.com/stacks/stacks-project/blob/a04446e57ec1fbc252a871afcec7752fb2807b14/algebra.tex\#L12723}{Original source, lines 12723--12736}. Complete original and reviewed TeX passages accompany this note. Every intervention is separately recorded; the following treatment is editorial.


\begin{lemma}
\label{editorial-source-lemma-artinian-finite-nr-max}
If $R$ is Artinian then $R$ has only finitely many maximal ideals.
\end{lemma}

\begin{proof}
If there were infinitely many maximal ideals, choose distinct
$\mathfrak m_1,\mathfrak m_2,\ldots$ and put
$J_n=\bigcap_{i=1}^n\mathfrak m_i$ for $n\geq1$.
For every $i\leq n$, distinctness and maximality imply
$\mathfrak m_i+\mathfrak m_{n+1}=R$. Choose actual elements
$a_i\in\mathfrak m_i$ and $b_i\in\mathfrak m_{n+1}$ with
$a_i+b_i=1$. Their product
$a_1\cdots a_n$ belongs to every $\mathfrak m_i$ for
$i\leq n$, hence to $J_n$, and is congruent to $1$ modulo
$\mathfrak m_{n+1}$. It therefore does not belong to
$J_{n+1}=J_n\cap\mathfrak m_{n+1}$.
Consequently $J_1\supsetneq J_2\supsetneq\cdots$ is an
infinite strict descending chain of ideals, contradicting the
Artinian hypothesis. Thus the set of maximal ideals is finite.
For the zero ring the set is empty, which also satisfies the claim.
\end{proof}


\subsection{Editorial treatment of Algebra lines 12738--12755}

\hypertarget{algebra-context-168}{}

\noindent\href{https://github.com/stacks/stacks-project/blob/a04446e57ec1fbc252a871afcec7752fb2807b14/algebra.tex\#L12738}{Original source, lines 12738--12755}. Complete original and reviewed TeX passages accompany this note. Every intervention is separately recorded; the following treatment is editorial.


\begin{lemma}
\label{editorial-source-lemma-artinian-radical-nilpotent}
Let $R$ be Artinian. The Jacobson radical of $R$ is a nilpotent ideal.
\end{lemma}

\begin{proof}
Let $I$ be the original Jacobson radical. By the descending
chain condition applied to $I\supset I^2\supset\cdots$,
choose $n\geq1$ such that $I^n=I^{n+1}$.
Put $J=\operatorname{Ann}_R(I^n)=\{x\in R:xI^n=0\}$.
If $J=R$, then $1\in J$ gives $I^n=0$, as desired.
Otherwise the nonempty collection of ideals strictly containing
$J$ has a minimal member $J'$: absence of a minimal member
would let one choose an infinite strict descending chain in that
collection, contradicting the hypothesis.
The original quotient $J'/J$ is nonzero and has no nonzero
proper submodule. Indeed the inverse image of such a submodule
under $J'\to J'/J$ would be an ideal strictly between $J$ and
$J'$. Hence $J'/J$ is simple.
By Lemma \ref{editorial-source-lemma-characterize-length-1}, it is isomorphic
to $R/\mathfrak m$ for some maximal ideal $\mathfrak m$.
It follows that $\mathfrak mJ'\subset J$.

Since $I\subset\mathfrak m$, for every $a\in I$ and $y\in J'$
we have $ay\in J$. For every $b\in I^n$ this gives
$b(ay)=0$. An element of $I^{n+1}=II^n$ is a finite sum
$\sum_{j=1}^t a_jb_j$ with $a_j\in I$ and $b_j\in I^n$;
its product with $y$ is the full sum
$\sum_{j=1}^t b_j(a_jy)=0$.
Thus $I^{n+1}J'=0$. The chosen equality $I^{n+1}=I^n$
then gives $I^nJ'=0$, so every $y\in J'$ lies in
$\operatorname{Ann}_R(I^n)=J$. This contradicts $J'\supsetneq J$.
Therefore $J=R$ and $I^n=0$. No finite generation of $I^n$
or Noetherian property has been assumed in this proof.
\end{proof}


\subsection{Editorial treatment of Algebra lines 12757--12784}

\hypertarget{algebra-context-169}{}

\noindent\href{https://github.com/stacks/stacks-project/blob/a04446e57ec1fbc252a871afcec7752fb2807b14/algebra.tex\#L12757}{Original source, lines 12757--12784}. Complete original and reviewed TeX passages accompany this note. Every intervention is separately recorded; the following treatment is editorial.


\begin{lemma}
\label{editorial-source-lemma-product-local}
Any ring with finitely many maximal ideals and
locally nilpotent Jacobson radical is the product of its localizations
at its maximal ideals. Also, all primes are maximal.
\end{lemma}

\begin{proof}
Let the maximal ideals be the distinct ideals
$\mathfrak m_1,\ldots,\mathfrak m_n$, and retain
$I=\bigcap_{i=1}^n\mathfrak m_i$.
If $n=0$, the ring is zero, since every nonzero ring has a
maximal ideal. Its spectrum is empty and its map to the empty
product of rings is the isomorphism of zero rings. Assume $n\geq1$.
Every element of $I$ is nilpotent, so it belongs to every prime.
Since $\mathfrak m_1\cdots\mathfrak m_n\subset I$, a prime
$\mathfrak p$ must contain some $\mathfrak m_i$.
Indeed, otherwise one could choose $a_i\in\mathfrak m_i\setminus
\mathfrak p$ for each $i$, whose product would belong to
$\mathfrak p$ while all its factors avoid it, contradicting
primality. Maximality then gives $\mathfrak p=\mathfrak m_i$.
This proves the assertion about prime ideals.

We construct the precise product isomorphism. For each ordered
pair $i\neq j$, choose $a_{ij}\in\mathfrak m_j$ and
$b_{ij}\in\mathfrak m_i$ with $a_{ij}+b_{ij}=1$, and set
$c_i=\prod_{j\neq i}a_{ij}$. For $n=1$ this product is $1$.
The element $c_i$ has residue $1$ at $\mathfrak m_i$ and
residue $0$ at every other maximal ideal. Thus
$c_i(1-c_i)\in I$. Choose $N_i\geq1$ for which
$[c_i(1-c_i)]^{N_i}=0$. Define the actual element
$$
e_i=\sum_{k=N_i}^{2N_i-1}
\binom{2N_i-1}{k}c_i^k(1-c_i)^{2N_i-1-k}.
$$
All its coefficients and powers are part of this construction.
The binomial identity $(c_i+(1-c_i))^{2N_i-1}=1$ gives
$$
1-e_i=\sum_{k=0}^{N_i-1}
\binom{2N_i-1}{k}c_i^k(1-c_i)^{2N_i-1-k}.
$$
Every term of $e_i$ contains $c_i^{N_i}$, and every term of
$1-e_i$ contains $(1-c_i)^{N_i}$. Their product is zero by
the chosen nilpotence equation. Hence $e_i^2=e_i$.
The residues of $e_i$ at all maximal ideals are the same as
those of $c_i$: substituting $0$ gives zero, and substituting
$1$ leaves just the term with $k=2N_i-1$, equal to one.
For $i\neq j$, the product $e_ie_j$ is an idempotent in $I$.
An idempotent that is nilpotent is zero, since all its positive
powers equal itself. Thus $e_ie_j=0$.
The sum $\sum_i e_i$ is an idempotent, and
$1-\sum_i e_i$ is an idempotent in $I$, so it is zero too.
We have obtained $\sum_i e_i=1$ with pairwise orthogonal
idempotents, without assuming a uniform nilpotence exponent for $I$.

Put $R_i=e_iR$, viewed as a ring with identity $e_i$.
The maps
$$
R\longrightarrow\prod_{i=1}^n R_i,\quad r\longmapsto(e_ir)_i,
\qquad
\prod_{i=1}^n R_i\longrightarrow R,\quad (x_i)_i\longmapsto\sum_i x_i
$$
are inverse ring homomorphisms, by the orthogonality and sum
equations just proved. The first component map is surjective
and has kernel $(1-e_i)R$: $e_ir=0$ implies
$r=(1-e_i)r$, and the reverse implication follows from
$e_i(1-e_i)=0$.
Maximal ideals of $R_i$ correspond to maximal ideals of $R$
containing $(1-e_i)R$. By the computed residues, only
$\mathfrak m_i$ contains $1-e_i$.
Therefore $R_i$ is a nonzero local ring with unique maximal
ideal $e_i\mathfrak m_i$.

The comparison with the original localization is
$$
R_i\longrightarrow R_{\mathfrak m_i},\quad x\longmapsto x/1,
\qquad
R_{\mathfrak m_i}\longrightarrow R_i,\quad
r/s\longmapsto (e_ir)(e_is)^{-1}.
$$
Here $s\notin\mathfrak m_i$ implies
$e_is\notin e_i\mathfrak m_i$: reducing modulo
$\mathfrak m_i$ would otherwise give $s=0$ in its residue
field. Hence $e_is$ is a unit in the original ring $R_i$.
Also $e_i/1$ is an idempotent unit in $R_{\mathfrak m_i}$,
so it equals one there. The first map is consequently unital.
For the second map, a fraction relation
$t(s'r-sr')=0$ with $t\notin\mathfrak m_i$ becomes the
same relation multiplied by $e_i$ in $R_i$; the images of
$t,s,s'$ are units, so the two prescribed images agree.
The maps preserve sums and products by their fraction formulas.
Their composites fix $x\in R_i$ and each original $r/s$,
respectively. Thus both are isomorphisms.
Combining them proves that the canonical map in the assertion is
$$
R\longrightarrow\prod_i R_{\mathfrak m_i},\qquad r\longmapsto(r/1)_i,
$$
with inverse
$(r_i/s_i)_i\mapsto\sum_i(e_ir_i)(e_is_i)^{-1}$.
All terms of this sum are elements of their specified subrings
$e_iR\subset R$; the original denominators and their inverses
in those subrings are retained.

The same construction works for every original $R$-module $M$,
without a finite-generation hypothesis. The maps
$M\to\bigoplus_i e_iM$, $x\mapsto(e_ix)_i$, and
$(x_i)_i\mapsto\sum_i x_i$ are inverse module maps.
The maps $e_iM\to M_{\mathfrak m_i}$, $x\mapsto x/1$,
have inverses
$$
x/s\longmapsto(e_is)^{-1}e_ix.
$$
The same fraction relation and the unit $e_it$ verify that this
inverse is well defined; the composites are identities because
$e_i/1=1$ in the localization.
Consequently $M\to\prod_i M_{\mathfrak m_i}$ is the canonical
isomorphism, with the displayed component maps and inverse sum.
Using finite additivity of length and the surjections
$R\to e_iR$, Lemma \ref{editorial-source-lemma-length-independent} gives
$$
\operatorname{length}_R(M)
=\sum_{i=1}^n\operatorname{length}_{R_{\mathfrak m_i}}
 (M_{\mathfrak m_i}),
$$
also when one or more terms are infinite. In the zero-ring case
every unital module is zero and this is the empty sum.
\end{proof}


\subsection{Editorial treatment of Algebra lines 12786--12819}

\hypertarget{algebra-context-170}{}

\noindent\href{https://github.com/stacks/stacks-project/blob/a04446e57ec1fbc252a871afcec7752fb2807b14/algebra.tex\#L12786}{Original source, lines 12786--12819}. Complete original and reviewed TeX passages accompany this note. Every intervention is separately recorded; the following treatment is editorial.


\begin{lemma}
\label{editorial-source-lemma-artinian-finite-length}
A ring $R$ is Artinian if and only if it has finite length
as a module over itself. Any such ring $R$ is both Artinian and
Noetherian, any prime ideal of $R$ is a maximal ideal, and $R$ is equal
to the (finite) product of its localizations at its maximal ideals.
\end{lemma}

\begin{proof}
If the original module $R$ has finite length, the proof of
Lemma \ref{editorial-source-lemma-finite-length-finite} gives both chain conditions
for its submodules, which are exactly its ideals. The descending
condition says that $R$ is Artinian. The ascending condition says
that every ideal is finite by the arbitrary-module equivalence
proved in Lemma \ref{editorial-source-lemma-Noetherian-basic}, so $R$ is Noetherian.

Conversely, suppose $R$ is Artinian. Lemmas
\ref{editorial-source-lemma-artinian-finite-nr-max},
\ref{editorial-source-lemma-artinian-radical-nilpotent}, and
\ref{editorial-source-lemma-product-local} identify $R$ with the finite product
of its original localizations $R_{\mathfrak m_i}$ and prove
that every prime is maximal. The product comparison is the
canonical map $r\mapsto(r/1)_i$, with the inverse proved there.
Its component maps are surjective: in the preceding proof each
localization is identified with the quotient $e_iR$ of $R$.
Thus each $R_{\mathfrak m_i}$ is Artinian, since a descending
chain of its ideals pulls back to such a chain in $R$.
When $R=0$, the product is empty and the length is zero.
For nonzero $R$, the module decomposition and length formula
in that same proof show that it suffices to prove finite length
for each original local factor.

Consider one such local Artinian ring $C=R_{\mathfrak m_i}$,
with maximal ideal $\mathfrak n=\mathfrak m_i C$ and residue
field $k=C/\mathfrak n$. By the nilpotence lemma,
$\mathfrak n^N=0$ for some $N\geq1$.
Each layer $V_j=\mathfrak n^j/\mathfrak n^{j+1}$,
$0\leq j<N$, is a $k$-vector space via the original quotient
map, because $\mathfrak n V_j=0$.
This space must have finite dimension. If not, choose a basis
containing distinct elements $v_1,v_2,\ldots$.
For each $a\geq0$ let $W_a$ be the span of all basis elements
except $v_1,\ldots,v_a$. Then
$W_0\supsetneq W_1\supsetneq\cdots$ is a strict descending
chain of vector subspaces. Their inverse images under
$\mathfrak n^j\to V_j$ are ideals of $C$ containing
$\mathfrak n^{j+1}$, and strictness is preserved by the
surjectivity of that quotient map. This contradicts the
Artinian condition on $C$.
By Lemma \ref{editorial-source-lemma-dimension-is-length}, the length of each
layer is its finite $k$-dimension. Additivity along the original
filtration gives
$$
\operatorname{length}_C(C)
=\sum_{j=0}^{N-1}\dim_k(\mathfrak n^j/\mathfrak n^{j+1})<\infty.
$$
Summing these finite lengths over the original local factors
proves finite length of $R$ and completes all assertions.

There are quantitative consequences for the same ring.
Write $\ell_i=\operatorname{length}_{R_{\mathfrak m_i}}
(R_{\mathfrak m_i})$. Each is positive and finite, and
$\operatorname{length}_R(R)=\sum_i\ell_i$.
The bound proved in Lemma \ref{editorial-source-lemma-length-infinite} gives
$(\mathfrak m_iR_{\mathfrak m_i})^{\ell_i}=0$.
Under the canonical product map, the original Jacobson radical
$I$ maps onto $\prod_i\mathfrak m_iR_{\mathfrak m_i}$:
membership in each component maximal ideal is exactly membership
in the corresponding original maximal ideal, and the product
map is an isomorphism. Powers are componentwise, so for
$R\neq0$,
$$
I^{\max_i\ell_i}=0.
$$
For $R=0$ the radical is already zero and no maximum of an empty
set is used.
Finally, every finite $R$-module $M$ has finite length. If its
given generating list has $r$ members, the associated surjection
$R^r\to M$ and additivity of length give the precise bound
$$
\operatorname{length}_R(M)\leq
r\operatorname{length}_R(R).
$$
The reverse implication, from finite length to finite generation,
is Lemma \ref{editorial-source-lemma-finite-length-finite}. Thus over this
Artinian ring the two conditions on $M$ are equivalent, with
the displayed bound retaining the chosen number of generators.
\end{proof}


\subsection{Editorial treatment of Algebra lines 12846--12855}

\hypertarget{algebra-context-171}{}

\noindent\href{https://github.com/stacks/stacks-project/blob/a04446e57ec1fbc252a871afcec7752fb2807b14/algebra.tex\#L12846}{Original source, lines 12846--12855}. Complete original and reviewed TeX passages accompany this note. Every intervention is separately recorded; the following treatment is editorial.


\begin{lemma}
\label{editorial-source-lemma-composition-essentially-of-finite-type}
The class of ring maps which are essentially of finite type is
preserved under composition. Similarly for essentially of finite
presentation.
\end{lemma}

\begin{proof}
Let the original maps be $R\to S\to T$. Choose an $R$-algebra
$A$ of finite type and a multiplicative subset $U\subset A$,
together with the given $R$-algebra isomorphism $U^{-1}A\cong S$.
Choose a finite type $S$-algebra $B$, a multiplicative subset
$V\subset B$, and the given $S$-algebra isomorphism
$V^{-1}B\cong T$.
Write $B=S[X_1,\ldots,X_t]/H$ for an ideal $H$.
Let $H_0\subset A[X_1,\ldots,X_t]$ be its inverse image under
the coefficient map to $S[X_1,\ldots,X_t]$, and put
$C=A[X_1,\ldots,X_t]/H_0$. This is a finite type $R$-algebra.
The natural map $U^{-1}C\to B$ is an isomorphism.
To see the kernel assertion explicitly, every polynomial over
$U^{-1}A$ is $G/u$ with $G\in A[X_1,\ldots,X_t]$ and
$u\in U$, by taking a product of its finitely many coefficient
denominators. If $G/u\in H$, then $G/1=u(G/u)\in H$,
so $G\in H_0$. Conversely every polynomial in $H_0$ maps
into $H$. This proves that the localized ideal is exactly
$H$, and proves the asserted quotient isomorphism with its
original polynomial and fraction maps.

In either this construction or the finite-presentation construction
below, set
$$
W=\{c\in C:\text{the image of }c\text{ in the original }T
\text{ is a unit}\}.
$$
This is a multiplicative subset. The canonical map
$W^{-1}C\to T$ is an isomorphism as follows.
Every image of $u\in U$ in $C$ belongs to $W$, so
$C\to W^{-1}C$ extends to $B=U^{-1}C$.
For $v\in V$, write its original representative in this
localization as $d/u$ with $d\in C$ and $u\in U$.
Its image is a unit in $T$, as is the image of $u$, hence
the image of $d$ is a unit in $T$ and $d\in W$.
Thus $v$ maps to a unit of $W^{-1}C$, and there is a map
$T=V^{-1}B\to W^{-1}C$.
It sends a fraction represented as $(c/u)/(d/w)$, with
$c,d\in C$, $u,w\in U$ and $d/w\in V$, to
$cw/(ud)$. The preceding argument applied to this actual
$d/w$ gives $d\in W$, and also $u\in W$, so its
displayed denominator belongs to $W$.
Both maps are defined by the localization relations and their
composites fix $C$, the inverted images of $U$ and $V$,
and the inverted elements of $W$. They are therefore mutual
inverses. This proves the assertion for finite type.

For finite presentation choose $A$ of finite presentation over
$R$ and choose a finite presentation
$B=S[X_1,\ldots,X_t]/(F_1,\ldots,F_m)$.
The contraction $H_0$ above need not be finitely generated,
so instead perform the following finite construction, keeping
all actual coefficients and denominators.
For each $j$, write
$$
F_j=\sum_{\nu\in E_j}\frac{a_{j\nu}}{u_{j\nu}}X^\nu,
\qquad a_{j\nu}\in A,\quad u_{j\nu}\in U,
$$
where $E_j$ is the finite set of exponent tuples occurring in
the polynomial. Set
$$
d_j=\prod_{\nu\in E_j}u_{j\nu},\qquad
G_j=\sum_{\nu\in E_j}
a_{j\nu}\left(\prod_{\mu\in E_j\setminus\{\nu\}}
u_{j\mu}\right)X^\nu\in A[X_1,\ldots,X_t].
$$
The coefficient image of $G_j$ in $S[X_1,\ldots,X_t]$ is
exactly $(d_j/1)F_j$. The factor $d_j/1$ is retained and is
a unit; thus the ideals generated by the images of the $G_j$
and by the $F_j$ are equal. If $F_j=0$, use $E_j=\varnothing$,
$d_j=1$ and $G_j=0$.
Put $C=A[X_1,\ldots,X_t]/(G_1,\ldots,G_m)$.
Then $U^{-1}C\cong B$ by precisely these coefficient maps.
Moreover $C$ is finitely presented over $R$: if
$A=R[Y_1,\ldots,Y_a]/(P_1,\ldots,P_b)$, choose polynomial
lifts $\widetilde G_j\in R[Y_1,\ldots,Y_a,X_1,\ldots,X_t]$
of all the original coefficients in $G_j$.
The resulting presentation is
$$
C\cong R[Y_1,\ldots,Y_a,X_1,\ldots,X_t]
/(P_1,\ldots,P_b,\widetilde G_1,\ldots,\widetilde G_m).
$$
Indeed quotienting by the $P_i$ first gives the polynomial
ring over $A$, after which the images of the remaining listed
relations are exactly the $G_j$.
Using the same $W$ and the same inverse localization maps proves
$T\cong W^{-1}C$, as required. These arguments allow empty
variable and relation lists and multiplicative subsets containing
zero; in the latter case the corresponding localization is the
zero ring and the same universal maps apply.
\end{proof}


\subsection{Editorial treatment of Algebra lines 12857--12866}

\hypertarget{algebra-context-172}{}

\noindent\href{https://github.com/stacks/stacks-project/blob/a04446e57ec1fbc252a871afcec7752fb2807b14/algebra.tex\#L12857}{Original source, lines 12857--12866}. Complete original and reviewed TeX passages accompany this note. Every intervention is separately recorded; the following treatment is editorial.


\begin{lemma}
\label{editorial-source-lemma-base-change-essentially-of-finite-type}
The class of ring maps which are essentially of finite type is
preserved by base change. Similarly for essentially of finite
presentation.
\end{lemma}

\begin{proof}
Retain the original map $R\to S$ and a map $\rho:R\to R'$.
Write $S=U^{-1}A$ through a specified $R$-algebra isomorphism,
where $A$ has the required finite type or finite presentation
property. Put $A'=A\otimes_R R'$ and let $U'$ be the image
of $U$ under $a\mapsto a\otimes1$.
There is an isomorphism of the actual base-changed rings
$$
\Phi:S\otimes_R R'\longrightarrow (U')^{-1}A',
\qquad
(a/u)\otimes r'\longmapsto
\frac{a\otimes r'}{u\otimes1}.
$$
To construct it, the ring map $A\to(U')^{-1}A'$ sends
$a$ to $(a\otimes1)/1$ and inverts $U$, so extends to $S$.
Together with the map $R'\to(U')^{-1}A'$,
$r'\mapsto(1\otimes r')/1$, it induces the displayed map;
the two maps agree on the original $R$ by tensor balancing.
Conversely $A'\to S\otimes_R R'$ sends
$a\otimes r'$ to $(a/1)\otimes r'$ and sends $u\otimes1$
to a unit with inverse $(1/u)\otimes1$.
It therefore induces a map $\Psi:(U')^{-1}A'\to S\otimes_R R'$.
Every numerator is a finite sum of pure tensors and every
denominator in $U'$ is the image of an original $u\in U$;
on those fractions its formula is
$$
\Psi\left(\frac{\sum_{i=1}^q a_i\otimes r_i'}{u\otimes1}\right)
=\sum_{i=1}^q(a_i/u)\otimes r_i'.
$$
The construction from the localization maps proves independence
of both tensor expressions and fraction representatives.
The composites fix each pure tensor $(a/u)\otimes r'$ and
each displayed fraction, respectively, so they are identities.

If $a_1,\ldots,a_t$ generate $A$ as an $R$-algebra, their
images $a_i\otimes1$ generate $A'$ as an $R'$-algebra:
apply the original polynomial expression for each $a$ in any
finite sum of pure tensors, retaining its coefficients through
$\rho$. Thus finite type is preserved.
For finite presentation, write
$A=R[X_1,\ldots,X_t]/(P_1,\ldots,P_m)$.
The coefficient map gives
$$
A'\cong R'[X_1,\ldots,X_t]/(\rho(P_1),\ldots,\rho(P_m)),
$$
where for the full original polynomial
$P_j=\sum_\nu r_{j\nu}X^\nu$ one has
$\rho(P_j)=\sum_\nu\rho(r_{j\nu})X^\nu$.
Indeed $R[X_1,\ldots,X_t]\otimes_R R'$ maps isomorphically
to $R'[X_1,\ldots,X_t]$ by
$rX^\nu\otimes r'\mapsto\rho(r)r'X^\nu$, with inverse
$r'X^\nu\mapsto X^\nu\otimes r'$.
Tensoring the quotient presentation is right exact, so the
kernel of the resulting surjection is generated by the images
of the original ideal $(P_1,\ldots,P_m)$.
A finite sum of its polynomial multiples maps to the same
finite sum of multiples of $\rho(P_j)$, and each listed
$\rho(P_j)$ occurs, proving the asserted ideal equality.
This is a finite presentation of $A'$.
In either case the proved isomorphism $S\otimes_R R'\cong
(U')^{-1}A'$ establishes the required essential property.
It remains valid when a denominator maps to zero, since both
localization constructions and their inverse maps then yield
the zero ring.
\end{proof}


\subsection{Editorial treatment of Algebra lines 12868--12924}

\hypertarget{algebra-context-173}{}

\noindent\href{https://github.com/stacks/stacks-project/blob/a04446e57ec1fbc252a871afcec7752fb2807b14/algebra.tex\#L12868}{Original source, lines 12868--12924}. Complete original and reviewed TeX passages accompany this note. Every intervention is separately recorded; the following treatment is editorial.


\begin{lemma}
\label{editorial-source-lemma-essentially-of-finite-type-into-artinian-local}
Let $R \to S$ be a ring map. Assume $S$ is an Artinian local ring with
maximal ideal $\mathfrak m$. Then
\begin{enumerate}
\item $R \to S$ is finite if and only if $R \to S/\mathfrak m$ is finite,
\item $R \to S$ is of finite type if and only if $R \to S/\mathfrak m$
is of finite type.
\item $R \to S$ is essentially of finite type if and
only if the composition $R \to S/\mathfrak m$ is essentially
of finite type.
\end{enumerate}
\end{lemma}

\begin{proof}
Write $k=S/\mathfrak m$ and retain its quotient map from $S$.
If $S$ is finite as an $R$-module, its quotient $k$ is finite.
Conversely, suppose $k$ is generated over $R$ by
$\bar z_1,\ldots,\bar z_r$. By
Lemma \ref{editorial-source-lemma-artinian-finite-length}, the original $S$
has a composition chain
$$
0=I_0\subsetneq I_1\subsetneq\cdots\subsetneq I_\ell=S,
\qquad \ell=\operatorname{length}_S(S).
$$
Each $I_j$ is an ideal and each quotient $I_j/I_{j-1}$ is
isomorphic, as an $S$-module, to the original $k$, since $S$
is local. For every $j$, choose lifts $z_{j1},\ldots,z_{jr}\in I_j$
of the images of the specified $\bar z_a$ under such an
isomorphism $k\to I_j/I_{j-1}$.
These $r\ell$ elements generate $S$ over $R$.
In fact, for $x\in I_j$ its class is
$\sum_{a=1}^r r_a(z_{ja}+I_{j-1})$ for some original
$r_a\in R$, so $x-\sum_a r_a z_{ja}\in I_{j-1}$.
Induction on $j$ supplies the claimed expression with all the
chosen lifted generators. This proves (1), and gives the
quantitative bound $r\operatorname{length}_S(S)$ on a number
of $R$-generators, with the full original filtration retained.

For (2), a quotient of a finite type $R$-algebra is of finite
type: the images of the given algebra generators still generate.
Conversely, lift a finite algebra generating list of $k$ to
$f_1,\ldots,f_n\in S$ and form the original coefficient map
$R[X_1,\ldots,X_n]\to S$, $X_i\mapsto f_i$.
Its composite to $k$ is surjective, so $k$ is a cyclic module
over this polynomial ring. Part (1), applied with this ring as
the base, shows that $S$ is a finite module over it.
Choose its finite module generators $s_1,\ldots,s_q$.
Every $s\in S$ is a sum of these $s_j$ with coefficients
that are polynomials in the original $f_i$ and coefficients
from $R$. Therefore $f_1,\ldots,f_n,s_1,\ldots,s_q$ generate
$S$ as an $R$-algebra, proving (2).

For (3), a quotient is a finite type map, hence essentially of
finite type; the forward implication follows from
Lemma \ref{editorial-source-lemma-composition-essentially-of-finite-type}.
Conversely suppose a specified localization of
$D=R[X_1,\ldots,X_n]/H$ is isomorphic to the original field
$k$. Choose $f_i\in S$ lifting the corresponding images of
$X_i$ in $k$. Let $J$ be the kernel of
$R[X_1,\ldots,X_n]\to S$, $X_i\mapsto f_i$, and retain
the actual subalgebra
$A=R[X_1,\ldots,X_n]/J\subset S$ and its prime
$\mathfrak p=A\cap\mathfrak m$.
The subring $A/\mathfrak p\subset k$ is exactly the image
of $D$ in $k$, since both are generated by the same original
coefficient images and the same $f_i$ modulo $\mathfrak m$.
Every element of $k$ is a fraction of elements of this image,
with a nonzero denominator coming from the specified
multiplicative subset of $D$. Thus $k\subset
\operatorname{Frac}(A/\mathfrak p)$ in the given ambient field.
Conversely the field $k$ contains that domain and the inverse
of each of its nonzero elements, so
$\operatorname{Frac}(A/\mathfrak p)\subset k$.
These two original embeddings identify
$\kappa(\mathfrak p)=\operatorname{Frac}(A/\mathfrak p)$
with the original $k$.

Every $a\in A\setminus\mathfrak p$ has image outside
$\mathfrak m$ and hence is a unit in $S$. The inclusion
$A\subset S$ therefore extends to
$A_{\mathfrak p}\to S$, $a/u\mapsto a\,u^{-1}$,
using the actual inverse in $S$.
Its composite to $k$ is the residue map just identified and
is surjective. Part (1) shows that $S$ is finite over
$A_{\mathfrak p}$. Hence it is finite type over that ring.
The map $R\to A$ is finite type and $A\to A_{\mathfrak p}$
is a localization. Applying the proved composition lemma to
these original maps gives that $R\to S$ is essentially of
finite type, as required.
\end{proof}

